\documentclass[10pt,a4paper,twoside,bg=print,twocolumn,openany,nodeprecatedcode]{dndbook}
\usepackage[utf8]{inputenc}
\usepackage{tabulary}
\usepackage[normalem]{ulem}
\usepackage{color,soul}
\usepackage{float}
\usepackage{caption}
\usepackage{wrapfig}
\usepackage{tocloft}
\usepackage[toc]{multitoc}
\usepackage[hidelinks]{hyperref}
\raggedbottom
\extrafloats{2000}
\cftsetindents{section}{0em}{0em}
\cftsetindents{subsection}{0em}{0em}
\cftsetindents{subsubsection}{0.2em}{0em}
\setcounter{tocdepth}{1}
\renewcommand*{\multicolumntoc}{3}
\setlist[description]{nolistsep,listparindent=3pt,topsep=6pt}
\date{}
\begin{document}
\frontmatter
\tableofcontents
\mainmatter
\addcontentsline{toc}{chapter}{Introduction: Welcome to the Nine Hells}
\markboth{INTRODUCTION: WELCOME TO THE NINE HELLS}{INTRODUCTION: WELCOME TO THE NINE HELLS}
\chapter*{Introduction: Welcome to the Nine Hells}
\DndDropCapLine{T}{}his story takes place in the Nine Hells and sees the characters confront its ruler, \textbf{Asmodeus}. They venture into the Nine Hells to save the souls of their loved ones (or their own) and are aided by a powerful faction. This faction has its own reasons for wanting revenge on \textbf{Asmodeus}. With the help of the characters, the faction believes it can deal a devastating blow to him. In return for a favor, the faction provides characters with access to the Nine Hells and a means of traversing the different layers. The characters travel by barge down the River Styx, also known as the River of Blood, a cross-planar river that links the layers of the Nine Hells. Those who touch or drink from its waters risk forgetting their past lives permanently.


The characters try to avoid \textbf{Asmodeus}'s attention, believing themselves unwanted infiltrators who are there to undermine the Lord of the Nine Hells. Little do they know, they're here by his design. They're integral to one of \textbf{Asmodeus}'s plots, and he aims to recruit them. But to get them under his control, he needs to corrupt them. It took years of careful planning, and now, finally, he has the characters exactly where he wants them. They will be tested, with temptations and traps set which only the pure of heart can resist. Each time a character gives in to sin, a link is added to the chains that bind them to the Nine Hells, to \textbf{Asmodeus}. Can they save their loved ones without losing themselves in the process?

\section{The Nine Hells}
The Nine Hells is the ultimate plane of evil and the epitome of premeditated cruelty. The devils of the Nine Hells are bound to obey the laws of their superiors, and at the very top of the hierarchy is \textbf{Asmodeus}. The Nine Hells has nine layers. The first eight are ruled by archdevils who answer to \textbf{Asmodeus}. To reach the deepest layer of the Nine Hells, one must descend through all eight of the layers above it, in order. The most expeditious means of doing so is the River Styx, which plunges ever deeper as it flows from one layer to the next. Only the most courageous adventurers can withstand the torment and horror of that journey.

\begin{DndSidebar}[float=!b]{Sourcebook and Adventure}
In addition to being an adventure, this is a source book for the Nine Hells. This book provides useful information for each level of the Nine Hells and its denizens and lords.



\end{DndSidebar}

The River Styx This river bubbles with foul water, and along its banks are the remains of thousands of battles. There are captains who are skilled enough to command a ship through the unpredictable currents of the river. For a price, these pilots are willing to carry passengers into the lower depths of the Nine Hells. They're the only ones who know how to avoid the dangers of the River Styx, including the warships that patrol it.

These warships are built from infernal iron (bristling with spikes, blades, and chains) and can travel the waters of the River Styx. Weapons used by the infernal warships include flamethrowers, acid spewed from bronze tubes, and launchers that fire a dozen ballista bolts at once.

\section{Sins}
Sin must be at the core of any story that takes place in the Nine Hells. These Hells are based on Dante's Inferno, the famed 14th-century poem which depicted Hell as nine concentric circles of torment. Each circle of Hell tormented a different class of sinner, from the unbelievers of the 1st circle down to the great betrayers of the 9th circle.

The poem focused on describing the acknowledgement and rejection of sin. In this story, good people have committed grievous sins that will have consequences. Some of these people might be the player characters.

\section{What Are the Sins?}
Each layer of the Nine Hells embodies a particular sin, as described in the sections that follow. Feel free to add other sins, but be careful when you do. Some classical interpretations of sins shouldn't be used in any adventure. Others, such as lust or gluttony, should be handled very carefully.

\subsubsection{Sins to the Self}

\begin{description}
\item[Anger.]
The sin of anger fuels wrath, vengeance, and inner turmoil. It takes root in Avernus, the domain of the archdevil \textbf{Zariel}.

\item[Greed.]
This sin leads to a life of accumulating wealth and sharing very little—a sin nurtured in Minauros, the domain of the archdevil \textbf{Mammon}.

\item[Jealousy.]
The jealous are never satisfied with what they have. Phlegethos, the domain of the archdevils \textbf{Belial} and \textbf{Fierna}, is where jealousy thrives.

\item[Pride.]
Some say that this is the original sin which led to all other sins and is the sin that led to the fall of \textbf{Asmodeus} and thus is his domain.

\end{description}

\subsubsection{Sins to the Other}

\begin{description}
\item[Betrayal.]
Betraying a loved one, a god, or a sovereign can damn one's soul to Stygia, the domain of the archdevil \textbf{Levistus}.

\item[Deceit.]
Lying that leads to harm is, like Maladomini, the domain of the archdevil \textbf{Baalzebul}.

\item[Harm.]
Mutilating or torturing an enemy is wicked work that warms the ice-cold heart of \textbf{Mephistopheles}, the archdevil who presides over Cania.

\item[Murder.]
Killing an innocent or family member is evil business that comes home to roost in Dis, the domain of the archdevil \textbf{Dispater}.

\item[Oppression.]
Unjustly imprisoning enemies and denying their freedom is the devil's work and reflective of Malbolge, the domain of the archdevil \textbf{Glasya}.

\end{description}

More details on the different archdevils, including their stat blocks, can be found in appendix A, "Lords of the Nine."

\begin{DndSidebar}[float=!b]{An Unjust Universe}
Some players may be uncomfortable if a soul is consigned to the Nine Hells for a sin when the person so condemned lived a good life in aggregate. However, all the souls consigned to the Nine Hells in this adventure were deceived by \textbf{Asmodeus}. While they committed sins, that wasn't the reason they were consigned to the Nine Hells. However, \textbf{Asmodeus} ensures that the soul is punished as if the sin was the one that defined it.



\end{DndSidebar}

\section{Running the Adventure}
To run this adventure, you need the fifth edition rulebooks: \textit{Player's Handbook}, \textit{Dungeon Master's Guide}, \textit{Monster Manual}, and \textit{Monsters of the Multiverse}. Baldur's Gate: Descent into Avernus would also be a good source; Chains of Asmodeus can be run as a continuation of that story.

The stat blocks for most creatures in this adventure can be found in the \textit{Monster Manual}. When a creature's name appears in bold type, that is a visual cue for you to look up the creature's stat block in the \textit{Monster Manual}, unless the adventure's text refers you to appendix B (for monsters), appendix C (for NPCs), or \textit{Monsters of the Multiverse}.

Spells and equipment mentioned in the adventure are described in the \textit{Player's Handbook}. Magic items are described in the \textit{Dungeon Master's Guide}, unless the text directs you to an item's description in appendix D of this book.

\begin{DndReadAloud}{}
Text that appears in a box like this is meant to be read aloud or paraphrased for the players when their characters arrive at a location or under specific circumstance, as described in the text.

\end{DndReadAloud}

\subsection{Abbreviations}
The following abbreviations appear in this book:

AC = Armor Class

DC = Difficulty Class

XP = experience points

NPC = non-player character

DM = Dungeon Master

pp = platinum piece(s)

gp = gold piece(s)

sp = silver piece(s)

cp = copper piece(s)

ep = electrum piece(s)

\subsection{Character Advancement}
The adventure is meant for characters of 11th level, and the story sees them rise to 20th. Creating characters of this level requires a little more preparation compared to 1st-level characters. If you and your players previously played Baldur's Gate: Descent into Avernus, your players may choose to continue with their characters from that adventure. Or, if your players have other premade characters of the appropriate level, they can decide to embody those characters once again. Players with preexisting characters must still choose the soul that their character is dedicated to saving.

The "Great Heroes" section helps you and your players get started with either creating a new character or with the choices that need to be made for preexisting characters.

\subsection{Consider Your Players}
Playing \textit{Dungeons \& Dragons} is about having fun with your friends and exploring new worlds together. Before venturing out onto the River Styx, discuss the nature of the story with your players.

You will be guiding them on a journey through the Nine Hells and introducing them to the various devils that make this infernal domain their home. The story contains references to the suffering and torture of those souls who are damned to an eternity of torment because of their sins. Make sure that you're aware of any boundaries that may exist, as you don't wish to make the experience uncomfortable for your players. The characters might be on a perilous quest through the Nine Hells themselves, but their players should enjoy the experience. You may need to adjust the story to accommodate all players.

\begin{DndSidebar}[float=!b]{Previous Adventures in Avernus}
An adventure that is the perfect prequel to this one is Baldur's Gate: Descent into Avernus. If you and your players have been on this adventure before our story begins, you may be familiar with some of the setting, and the players' characters should be at the appropriate level. However, depending on the outcomes of your earlier encounters within the Nine Hells, you may need to adjust some parts of the story.



\end{DndSidebar}

\subsection{Agency and Replayability}
This adventure has been designed so that each time it is played it can have entirely different objectives and outcomes. It does this mainly through the following ways:


\begin{description}
\item[Souls.]
Each player must choose a soul to save. It might be the soul of a loved one or their own soul. The soul, and the sins associated with it, determine where \textbf{Asmodeus} has hidden it away to be tormented.

\item[Group Patron.]
The players must decide upon a group patron before the start of the adventure. There are three group patrons, one for good, one for evil, and one for neutrality. Each has their own objective.

\item[The River Styx.]
The river passes dozens of dangerous locations, but the characters visit only some of them.

\item[Temptations and Corruptions.]
\textbf{Asmodeus} tries to lure the characters into his grasp. Each time the characters go back to the Nine Hells they may find different temptations. Some may be easy to resist, others not so much. Giving in to temptation corrupts them.

\end{description}

\section{Story Overview}
The Lord of the Nine Hells, \textbf{Asmodeus}, is known to have countless plots, some that take centuries to come to fruition, while others he conceives of and executes within a few years. This is a story about one of the latter.

\textbf{Asmodeus} desires to corrupt and dominate some of the most powerful mortals in a generation. To further his designs, \textbf{Asmodeus} wants to convert these powerful mortals to his cause. To this end, \textbf{Asmodeus} has made pacts with many of the powerful and influential of the Forgotten Realms. This includes kings, queens, high priests, and powerful adventurers. The player characters are among those that he has targeted. When not able to convince such august personages, he has gone after their spouses, children, or friends. He comes to them when they're in need, providing a service with the understanding that if they die before they can return the favor, their soul is forfeit.

\textbf{Asmodeus} never intended to allow these infernal contracts to be fulfilled by their mortal signatories. Instead, he has given a list of names to the most skilled killers in the world.

These killers embark on a murder spree to kill all the people on the list, thus consigning their souls to the Nine Hells and to the corrupting influence of \textbf{Asmodeus} and his court of devils. Among the victims when the adventure begins are the spouses, children, or friends of the characters. One or more characters may also be among those who are murdered.

\section{Adventure Flowchart}
\section{The Nine Hells}
This overview summarizes the rulers of each layer of the Nine Hells, and the objectives characters might seek in each. These objectives include the souls of the characters, the souls of those they care about, and the items or people their group patron requires.

\subsection{Important Characters}
These are the important characters in the story:


\begin{description}
\item[Aeshma (the Unmaker).]
A lost soul that the Conclave of Halruaa needs destroyed to stop a dangerous threat. Nobody knows the Unmaker's true identity.

\item[Asmodeus.]
The Lord of the Nine Hells, and this adventure's primary antagonist.

\item[Barachiel.]
A Hellrider and one of the friends that \textbf{Ramius} seeks.

\item[Koh Tam.]
A priest serving the God of the Dead.

\item[Ramius.]
A contact for the Hellriders of Elturel group patron.

\item[Sarevok.]
High priest of Bhaal, and the contact for the Deathstalkers of Bhaal group patron.

\item[Tiax.]
\textbf{Koh Tam}'s sidekick and comedy relief.

\item[Zythan.]
High Diviner of Halruaa, and a contact for the Conclave of Halruaa group patron.

\end{description}

\subsection{Prologue}
The characters are approached by a representative of one of three group patrons, representing either good, neutrality, or evil. This representative acts as the quest giver and has their own motivation for wanting the characters to go into the Nine Hells. The three patrons are as follows:


\begin{description}
\item[Ramius, Hellrider of Elturel (Good).]
\textbf{Ramius} has a deep hatred for \textbf{Asmodeus} because of what the archdevil did to his home city and its patron angel. \textbf{Ramius}' lifelong friend was a victim of the recent murder spree. The paladin wants the characters to enter the Nine Hells to rescue his friend and other Celestials who have fallen victim to \textbf{Asmodeus}. \textbf{Ramius} and the Hellriders know the location of these beings trapped within the Nine Hells.

\item[Zythan, High Diviner of the Conclave of Halruaa (Neutral).]
\textbf{Zythan} has foreseen a dreadful consequence of \textbf{Asmodeus}'s plan. A dangerous soul known as the Unmaker has emerged and must be dealt with. \textbf{Zythan} and the mages of Halruaa have divined the location of powerful magic items that can help the characters in this quest. He wants the characters to gather these, to defeat the threat he has foreseen. These include some of the most famed items in the \textit{Dungeon Master's Guide}.

\item[Sarevok, High Priest of the Deathstalkers of Bhaal (Evil).]
The Deathstalkers of Bhaal feel uneasy about their alliance with \textbf{Asmodeus}. They fear he will betray Bhaal and that the divine power promised him will never be given. \textbf{Sarevok} will help the characters navigate the River Styx and rescue the souls of their loved ones, but only if they help Bhaal gain leverage against \textbf{Asmodeus} by stealing some of the most infamous and powerful artifacts of the Nine Hells. He aims to keep one, the others are for the characters to decide what to do with.

\end{description}

More information on the patrons can be found in the "Group Patron" section of this chapter.

\subsection{Koh Tam}
The group patron puts the characters in contact with \textbf{Koh Tam} (see chapter 1), a priest serving the God of the Dead. \textbf{Koh Tam} is an expert on the Nine Hells and knows how to navigate the River Styx. But first, he must learn where the souls that the characters seek might be found. To that end, he summons \textbf{Baalzebul} and invites the archdevil to possess him. The characters must deal with the devil to learn the location of the souls. Afterwards, they must exorcise \textbf{Koh Tam} of the devil's spirit.

\subsection{Through the Nine Hells}
\textbf{Koh Tam} knows how to navigate the River Styx and offers to guide the characters through the Nine Hells to help them find the souls they seek. He opens a portal to the first layer of the Nine Hells, Avernus.

The characters must descend through the layers of the Nine Hells along the treacherous River Styx. It is possible to move against the flow of the River Styx, but it is difficult. However, \textbf{Koh Tam}'s barge has a magical addition that allows him to use the \textit{Plane Shift} spell to travel to a layer of the Nine Hells above the current one.

The characters have two objectives if they're to save the souls of their loved ones or their own:


\begin{description}
\item[Finding the Lost Souls.]
The main objective of the characters is to find the souls that have been trapped by \textbf{Asmodeus}.

\item[Dealing with Asmodeus.]
Once the characters have found their lost loved ones, they must convince \textbf{Asmodeus} to allow them to be freed. They can do this by gaining leverage over him, by performing a service for him, or signing an infernal contract. Characters who are able to resist temptation within the Nine Hells might even find a loophole allowing them to escape, successfully avoiding all of \textbf{Asmodeus}'s snares.

\end{description}

\subsection{Finale: Asmodeus}
\textbf{Asmodeus} is a lawful evil deity who is genuinely one of the most powerful creatures in the universe. He likes to appear respectful and well-mannered when speaking with others, and he hates being wrong. In fact, sometimes he even pretends to know more than he does, or that he can predict what actions creatures will take before they take them, even though he can't. He enjoys making contracts with others, especially when their end of the bargain is their soul.

At the climax of the adventure, the characters may need to make a contract with \textbf{Asmodeus} to release the souls of their loved ones. The characters need as much leverage as possible when going into the contract. Without enough leverage, they may have to promise to do services for \textbf{Asmodeus}.

\section{Great Heroes}
This adventure is for high-level play. Players start with characters who are 11th level and are considered Masters of the Realm. Over the course of play, they go beyond 17th level to become Masters of the World. See the \textit{Dungeon Master's Guide} for more details on these tiers of play. There are three steps that players must complete before the start of the adventure:


\begin{itemize}
\item Create a new character, or choose an existing one.
\item Choose a soul to be saved.
\item Choose a group patron.
\end{itemize}

Players need to create a new high-level character unless they have a preexisting one they want to use. Any player who creates a new high-level character must go through the character creation process and choose a soul that their hero is dedicated to saving. A player with a preexisting character must still choose a soul that their character is dedicated to saving.

When choosing a soul to save, a player has the option of saving the soul of a loved one (detailed in the "Lost Souls" section) or saving their own soul (detailed in the "Souls of the Damned" section). While the lost souls option was designed to work with existing characters and the souls of the damned option was designed to work with newly created high-level characters, the DM can allow players to choose from either. A group can be made up of lost souls and souls of the damned. However, to promote group unity, a DM can insist upon one or the other.

After players have created their characters, they must choose a group patron. These group patrons are designed for three types of play that coincide with good, neutral, and evil alignments.

\section{New Characters}
Creating an 11th-level character uses the same character creation steps outlined in the \textit{Player's Handbook}. The process will take longer than usual, as a player must select spells, feats, abilities, and other aspects that a 1st-level character wouldn't have to. A new character starts with 7,500 gp, the standard equipment associated with their background, and magic items as described next.

\subsubsection{Magical Equipment}
Each new high-level character should be allowed to pick either 1 very rare magic item and 1 uncommon magic item or 2 rare magic items. These items should be selected from the \textit{Dungeon Master's Guide}.

\section{Saving Souls}
The following sections describe the souls available for each player to pick from. The chosen soul is the soul their character seeks to save. The player can either choose a lost soul, a loved one whose soul is trapped in the Nine Hells, or they can decide it's their character's soul that is damned. \textbf{Asmodeus} keeps these souls in phylacteries scattered around the Nine Hells. Phylacteries don't need to be in the layer associated with their sin, but usually are. When characters locate such a phylactery, they must bring it with them for the remainder of the adventure.

Only \textbf{Asmodeus} can break the infernal contract that binds a soul to its phylactery. To convince him to do this, characters must gain leverage over him, as described in later sections.

\section{Lost Souls}
These lost souls belong to the character's loved ones who have signed infernal contracts. When their phylactery is recovered, the lost soul can communicate with the character, providing advice, guidance, commentary, and even criticism. This type of phylactery also provides characters with a statistical or gameplay bonus.

The section that follows provides examples of lost souls that a player can choose.

\subsubsection{Brother (Anger)}
From the day he was born, your brother was a problem child—impulsive, aggressive and violent, though never toward you. He was constantly getting into fights—most of which he started—and you were constantly coming to his rescue and cleaning up his messes. But despite his flaws, you always knew he had your back. His temper grew worse as he grew older, and his fights often resulted in people getting hurt. He fell in with an unsavory crew, and you drifted apart. Now you've learned he was foolish enough to sign an infernal contract with \textbf{Asmodeus}, and his soul is trapped in the Nine Hells. You can't deny your brother deserves his fate... but he's still your brother, and you're not going to turn your back on him.


\begin{description}
\item[Phylactery Location.]
The War-Slough on Avernus, the first layer of the Nine Hells

\item[Phylactery Benefit.]
While you possess your brother's phylactery, you can channel your brother's rage. When you make your first attack on your turn, you can decide to attack recklessly. Doing so gives you advantage on melee weapon attack rolls using Strength during this turn, but attack rolls against you have advantage until your next turn.

\end{description}

\subsubsection{Bounty (Murder)}
Someone of your reputation and talent is always in high demand. If the price is right, you'll take on any job... but you don't work cheap. Still, this is the strangest adventure you've ever been hired for. The ruthless matron of a powerful noble family has disappeared, and you've been hired to find her. Your first assumption was that she'd either been kidnapped for ransom, murdered by a rival family, or killed by one of her own ambitious kin. But during your investigations, you discovered the shocking truth: decades ago, she signed an infernal contract to help advance her family's fortunes, and \textbf{Asmodeus} came to collect his payment. Tracking her soul through the Nine Hells wasn't part of the original deal, but the reward for bringing her back is too high to pass up, and you never leave a job unfinished.


\begin{description}
\item[Phylactery Location.]
The Agora of Floating Knives on Dis, the second layer of the Nine Hells

\item[Phylactery Benefit.]
While you possess your bounty's phylactery, you're more adroit at killing. When you miss with a melee weapon attack, you can choose to hit instead. Once you use this benefit, you can't use it again until you finish a long rest.

\end{description}

\subsubsection{Business Partner (Betrayal)}
For years, you and your business partner shared the fruits of your adventures in a mutually beneficial arrangement. Sure, they were always trying to squeeze a few extra drops from your cut on every mission, but you never thought they'd actually betray you. On your last adventure, however, they double-crossed you, stealing a rare and powerful magic item you'd spent years questing for. They disappeared with your treasure, leaving you high and dry. You searched for the precious item in vain, only to discover that your old partner's soul is trapped in the Nine Hells because of an infernal contract they signed with \textbf{Asmodeus}. The only way you'll ever get back what is rightfully yours is to track them down and bring them back to life so they can tell you where the treasure is hidden.


\begin{description}
\item[Phylactery Location.]
The Chasm of Found Things on Stygia, the fifth layer of the Nine Hells

\item[Phylactery Benefit.]
While you possess your business partner's phylactery, you have advantage on saving throws against illusion spells and effects that would make you have the charmed condition.

\end{description}

\subsubsection{Childhood Friend (Sinless)}
Growing up, you never seemed to fit in. You were always different—perhaps that's why you became an adventurer. Your childhood was lonely until you met your childhood friend. Like you, they were completely ignored by other children on the good days, ruthlessly bullied on the bad. But in each other you found a kindred soul. You formed a powerful bond of friendship that allowed you to rise above the taunts and torments of the other kids. Together, you found the strength to survive and flourish. You became a successful adventurer. Your friend chose a different path, one of domestic bliss with a family and a quiet, fulfilling life. Years would pass between seeing them, but you still remained close—united by your childhood bond. But when your friend's youngest child became ill, they didn't turn to you for help. Instead, they signed a contract with \textbf{Asmodeus}, condemning themselves to an eternity in the Nine Hells.


\begin{description}
\item[Phylactery Location.]
Use an encounter and location from the remaining lost souls not chosen by the players.

\item[Phylactery Benefit.]
Possessing your childhood friend's phylactery grants you advantage on Wisdom and Charisma saving throws.

\end{description}

\subsubsection{Father (Greed)}
You loved your father. He helped shape you into the hero you're now. He supported your ambitions, gave you advice when you needed it, and was there for you when things looked darkest. Your mother loved him even more, for he was personable, funny, and kind. But everyone has a flaw, and your father's flaw was gambling. The risk-taking behavior that you inherited from him made you into a hero, but he didn't have your outlet. Instead, he took risks with money. His debts grew, as did his shame, but you never knew. Your mother didn't tell you about it until it was too late. \textbf{Asmodeus} used his shame to convince him to sign an infernal contract.


\begin{description}
\item[Phylactery Location.]
The Ineffable Trove on Minauros, the third layer of the Nine Hells

\item[Phylactery Benefit.]
Possessing your father's phylactery grants you extra luck. Whenever you make an attack roll, an ability check, or a saving throw, you can roll an additional d20. You can choose to use this benefit after you roll the die, but before the outcome is determined. You choose which of the d20s is used for the attack roll, ability check, or saving throw. Once you use this benefit, you can't use it again until you finish a long rest.

\end{description}

\subsubsection{Mentor (Oppression)}
You learned your craft at the feet of a true master. Recognizing your raw talent and potential, your mentor took you under their wing and trained you. Your mentor was a difficult master and demanded complete obedience. Over the years, they taught you everything you needed to know about adventuring and life, shaping you into the hero you're today. They shared all their secrets with you... except for one. Long ago, your mentor signed an infernal contract with \textbf{Asmodeus}. Now your mentor has been whisked away to the Nine Hells, and it's time to use everything they taught you to save them from an eternity of torment and suffering.


\begin{description}
\item[Phylactery Location.]
The Sign of the Hag's Arms on Malbolge, the sixth layer of the Nine Hells

\item[Phylactery Benefit.]
Possessing your mentor's phylactery grants you protection. You gain a +1 bonus to AC and saving throws.

\end{description}

\subsubsection{Mother (Pride)}
Your mother was kind and loving, but also strong and independent. She raised you on her own, supporting you by telling fortunes and performing tarot readings for the locals in the town where you grew up. Once you began your adventuring career, her reputation grew quickly. Officials and nobility of ever-increasing rank and importance began to come to your mother for insight into their futures, amazed at the accuracy of her predictions. What you didn't know was that her rapid rise in stature was due to her signing an infernal contract with \textbf{Asmodeus} to enhance her gifts.


\begin{description}
\item[Phylactery Location.]
The Sign of the Hag's Arms on Malbolge, the sixth layer of the Nine Hells

\item[Phylactery Benefit.]
Possessing your mother's phylactery grants you foresight. When you finish a long rest, roll a d20 and record the number rolled. You can replace any attack roll, saving throw, or ability check made by you or a creature that you can see with this foretelling roll. You must choose to do so before the roll. This roll can be used only once. When you finish a long rest, you lose an unused foretelling roll.

\end{description}

\subsubsection{Sister (Jealousy)}
Your sister was fearless. There was no challenge she wouldn't tackle, no dare she wouldn't take. She became an adventurer, living a life filled with thrills and excitement. She even inspired your own career, pushing you to ever greater accolades. But over the years, her risk-taking became a compulsion she couldn't control. She became more and more reckless, embarking on foolhardy quests and what seemed like suicide missions. What you didn't realize was that it was your sister's envy of your accomplishments that drove her recklessness. \textbf{Asmodeus} exploited her need to be your equal, convincing her to sign an infernal contract in exchange for the ultimate adventure.


\begin{description}
\item[Phylactery Location.]
The Elemental Preserve on Phlegethos, the fourth layer of the Nine Hells

\item[Phylactery Benefit.]
While you possess your sister's phylactery, any effect that applies the frightened condition on one or more targets has no effect on you.

\end{description}

\subsubsection{Spouse/True Love (Deceit)}
True love is real—you know this because you were lucky enough to find yours. In a world too often filled with darkness and danger, you were blessed to find someone to stand by you through thick and thin. They supported you when you needed it most, helping you to reach your full potential. Together, you made a life for yourself: a wonderful, perfect life. But now you have learned that your love kept a secret from you: they made an infernal contract with \textbf{Asmodeus} to help keep you safe during your adventures. They made the ultimate sacrifice to protect you, and now they're paying the ultimate price.


\begin{description}
\item[Phylactery Location.]
The Eye Market on Maladomini, the seventh layer of the Nine Hells

\item[Phylactery Benefit.]
Possessing your true love's phylactery grants you advantage on saving throws against spells and other magical effects.

\end{description}

\subsubsection{Student (Harm)}
You started at the bottom and clawed your way to the top, forging your reputation over the years. As your legend grew, others came to you seeking guidance, but you always turned them away. One day, you discovered a student who was different than the others. Although the student was willful and prone to violence, you still saw a spark in them that reminded you of your own humble beginnings, so you took them under your wing. You helped them reach their full potential, growing skills until their ability and reputation rivaled your own. But despite all you taught them, your student still made a tragic mistake. They signed an infernal contract with \textbf{Asmodeus}, foolishly convinced they could get out of it before it came due. Now it's up to you to help them one last time.


\begin{description}
\item[Phylactery Location.]
The Sorrow Mine on Cania, the eighth layer of the Nine Hells

\item[Phylactery Benefit.]
While you possess your student's phylactery, you can use a bonus action to heal a number of hit points equal to twice your level. Once you use this ability, you can't use it again until you finish a long rest.

\end{description}

\section{Souls of the Damned}
This option is for any player who wants a more personal reason for their character to explore the Nine Hells: to recover their own lost soul. The character begins the adventure having been raised from the dead by one of the group patrons presented later in this book. This group has brought the character back using an ancient ritual, but it is only temporary. The character must complete the adventure to save their soul, or they will return to the Nine Hells as a lemure.

The DM might decide that all the characters have had their souls consigned to the Nine Hells. The advantages and disadvantages of not having a soul are explained in the "Not Having a Soul" sidebar.

\subsection{Choosing a Sin}
First, each player must decide if their character has harmed themself or another, then choose which sin caused their soul to be consigned to the Nine Hells. Sins to the other represent possible harm a character has done to others, whilst sins to the self represent possible strong detrimental character flaws. The Categories of Sin table outlines the choices a player may make.

\begin{DndTable}[header=Categories of Sin]{XX}

Sin & Harm or Flaw\\
\textit{Sins to the Other} & \\
Betrayal & The Oathbreaker\\
Deceit & The Great Con\\
Harm & The Merciless\\
Murder & Patricide\\
Oppression & The Heartless Master\\
\textit{Sins to the Self} & \\
Anger & The Furnace\\
Greed & The Infinite Treasure\\
Jealousy & The Queen\\
Pride & The Chosen One\\
\end{DndTable}

\begin{DndSidebar}[float=!b]{Not Having a Soul}
When a character whose soul is trapped in the Nine Hells dies, resurrection magic can return that character to life no matter how much time has passed. For example, a \textit{Revivify} spell works on the character even if a week has passed since the character died. However, if the character's body is completely destroyed, not even a Wish or \textit{True Resurrection} spell can bring the character back.



\end{DndSidebar}

\subsubsection{The Chosen One (Pride)}
You always knew you were special. Even from a young age, it became clear that you were better than everyone else. Because of this, you thought nothing of using, manipulating, or exploiting others for your own benefit. Stepping on the insignificant as you climbed your way to the top was simply the natural order. The fame and glory you achieved in your life only proved you were right all along... though looking back, maybe you pushed things too far. After all, you did end up in the Nine Hells.


\begin{description}
\item[Character Punishment.]
Your greatest fear in life was being inadequate. Now you're haunted by a recurring nightmare in which you compete in trials of skill against other heroes but find yourself weak, overmatched, and helpless.

\item[Phylactery Location.]
The Agora of Floating Knives, on Dis, the second layer of the Nine Hells. The character's soul is forced to compete in a series of events showcasing strength, skill, and intelligence. But their soul is merely a weakened shell of what it once was, and they're constantly overmatched by their fellow competitors. With each loss, they're humiliated, mocked, and made to suffer painful punishments for their failure. To free their soul, the character must take their place in these events; they must compete and win.

\end{description}

\subsubsection{The Furnace (Anger)}
There has always been a fire inside you, a dangerous fury, just waiting to be unleashed upon the world in a storm of violence. The rage gave you strength and power. It helped you defeat your enemies and forge your legend... but there was a cost to your anger. Those around you quickly learned not to try your temper, but over the years it was inevitable that something would set you off from time to time. When it was unleashed, your wrath was terrifying to behold... but it also led you to do things you'd later regret. People suffered. People died. Some deserved it, but many did not. Given what you've done—what your anger made you do—you've known for a long time where you were going to end up when you died.


\begin{description}
\item[Character Punishment.]
The color of your sleep is red: blood, fire, burning rage. You don't dream—not in the normal sense. When you close your eyes, all you see is red. You wake bathed in sweat, with your heart pounding, but you can't remember any details beyond that single color. Red.

\item[Phylactery Location.]
The Elemental Preserve, Phlegethos, the fourth layer of the Nine Hells. The character's soul is forced to shovel coal into a massive furnace, whipped to greater and greater efforts by several efreet with flaming scourges. No matter how much fuel is thrown into the furnace, the roaring flames instantly devour it in a burst of searing heat. To free this character's soul, the characters must defeat the efreet and then find a way to quench the magical fires of the furnace.

\end{description}

\subsubsection{The Great Con (Deceit)}
You've always been a grifter, using your wits to trick and manipulate others for your own advantage. But when you unknowingly made a deal with one of \textbf{Asmodeus}'s agents, you learned that even you can be conned. Before you could fulfill the seemingly simple terms of the deal, you were murdered by the Avatar of Bhaal, consigning you to the Nine Hells. On the one hand, you admire the devil's cunning trickery... but on the other, you want payback for how you were played.


\begin{description}
\item[Character Punishment.]
You were once a master of deceit, but in your dreams you're trapped in a labyrinth of your own lies, forced to run through the maze while being hunted by shadowy enemies and never able to find a way out.

\item[Phylactery Location.]
The Eye Market in the territory of Memnoriac, on Maladomini, the seventh layer of the Nine Hells. The character's soul is pursued by the pit fiend's gang of devils through the very streets Memnoriac designed. The route is littered with illusions and convoluted traps. Deadly spiked pits appear as normal floors or seemingly safe doorways lead into rooms filled with deadly fire. As the soul flees the devils, it constantly "dies," forcing it to begin again. To free this character's soul, the characters must fight their way through the narrow streets, defeating the monsters and avoiding the tricks and traps.

\end{description}

\subsubsection{The Heartless Master (Oppression)}
You were used to the finer things in life; due to your fame and fortune, you surrounded yourself with servants to answer every need and fulfill every desire. You paid them well, but in exchange you demanded flawless perfection. Unbeknownst to you, the head of your household achieved this impossible standard through brutal training and punishments for any servant that failed even the smallest task. Despite your ignorance, you were ultimately responsible for the well-being of those under your charge... a responsibility you blissfully abdicated.


\begin{description}
\item[Character Punishment.]
Your sins weigh heavily on you. At night, you're plagued by disturbing dreams. In them, you're a servant under control of a cruel master who makes impossible demands. When you inevitably fail at your tasks, you're brutally punished over and over.

\item[Phylactery Location.]
The Sign of the Hag's Arms, on Malbolge, the sixth layer of the Nine Hells. The character's soul is a servant ordered to perform a series of impossible tasks: cleaning out entire stables in a single day; serving a perfectly cooked meal despite a fire that won't stay lit; delivering a message across deadly terrain without weapons or armor. With each failure, the character's overseer—a ruthless hag—subjects the character's soul to a series of brutal tortures to "teach" them the error of their ways. To save their soul, the character must find a way to complete these impossible tasks, then defeat the overseer.

\end{description}

\subsubsection{The Infinite Treasure (Greed)}
You were never afraid to take what was yours—or what was anyone else's, for that matter. You lusted after material wealth in all its forms—money, possessions, magic items. During your life, you accumulated a vast hoard of treasures, but you couldn't take it with you when you died. In the end, the only thing your avarice got you was a one-way ticket to the Nine Hells.


\begin{description}
\item[Character Punishment.]
Your dreams are filled with the endless counting of gold coins—hundreds, thousands, millions, all piled in an enormous chamber. Each one must be counted by hand, and you can't stop until you're done... but the mountain of coins just continues to grow and grow. Worst of all, you know in your heart they aren't even yours; you're just counting them for someone else.

\item[Phylactery Location.]
The Ineffable Trove, on Minauros, the third layer of the Nine Hells. The character's soul is forced to count a massive mountain of gold coins one by one. This task is overseen by a powerful red dragon—the true owner of the coins. At the end of each day, the mountain of coins shifts and topples, crushing the character under the coins' weight. It takes all night for the character to claw their way out, only to begin the count again. To free this character's soul, the characters must defeat the dragon and leave the coins. (Taking even one causes the dragon to come back to life.)

\end{description}

\subsubsection{The Merciless (Harm)}
Throughout your adventuring career, your one true love was always with you, fighting by your side. You looked forward to a long and happy life together after your adventuring days were done, but that future was ripped away. Your partner was killed in a surprise attack by servants of an enemy from your past, their body disintegrated so they couldn't be raised. You hunted down the one responsible. But in your grief, it wasn't enough to merely defeat or kill them. You spent days torturing your enemy, making them feel as much pain as possible in the hopes it would soothe your own tormented heart. In the end, it didn't help... but your cruelty was such that upon your death your soul was consigned to the Nine Hells.


\begin{description}
\item[Character Punishment.]
The brutal revenge against your enemies brought you no solace. Instead, it only twisted and corrupted the memories you have of your true love. Now when you sleep, you dream of them not as your companion, but as your sadistic tormentor, torturing you in a prison cell from which there is no escape.

\item[Phylactery Location.]
Imprisoned in the mining town at the foot of the Sorrow Mine, on Cania, the eighth layer of the Nine Hells. The character's soul is imprisoned by devils with the face and form of their one true love. In the prison, the soul is relentlessly tortured over and over by these doppelgangers. To free this character's soul, the characters must break into the heavily guarded cells, find a way to reveal the true forms of the guards, then defeat them and escape the prison.

\end{description}

\subsubsection{The Oathbreaker (Betrayal)}
In the early days of your adventuring career, you swore an oath of loyalty to a patron. It was only much later you discovered you served a monster. When you were ordered to slaughter a political rival's children, you refused. Despite your good intentions, you broke your solemn vow to serve, condemning your soul to the Nine Hells on a technicality when the patron's hired assassins killed you.


\begin{description}
\item[Character Punishment.]
Every night, you're haunted by a recurring nightmare in which you're trapped in a room with strange, feral children. They cry and beg for mercy, even as they attack you over and over.

\item[Phylactery Location.]
The Chasm of Found Things, on Stygia, the fifth layer of the Nine Hells. The character's soul is imprisoned in a room filled with polymorphing devils. These creatures alternate between two forms: devils that attack the character's unarmed soul relentlessly and crying children begging for mercy. To free this character's soul, the characters must defeat all the devils before they transform back into children.

\end{description}

\subsubsection{Patricide (Murder)}
You murdered your father. It wasn't premeditated, and you loved your father. But an argument escalated to the point where violence was involved. It was an accident that you will forever regret. Before you could use your considerable influence to have a priest raise him from the dead, you were imprisoned for your crime and then hanged.


\begin{description}
\item[Character Punishment.]
Your dreams are haunted by the screams of your suffering soul. Each dream is the same—a replay of your fatal argument with your father. Instead of killing your father, he overpowers you and then whips you relentlessly with a spiked belt. You know that you must end this cycle of violence.

\item[Phylactery Location.]
The Agora of Floating Knives, on Dis, the second layer of the Nine Hells. The character's soul has been imprisoned in a flesh golem that has the character's face. It is constantly whipped with a spiked belt by a devil bearing the face of the character's father. Both the flesh golem and devil must be defeated to free the character's soul and reclaim it.

\end{description}

\subsubsection{The Queen (Jealousy)}
Nobody said life was fair... but why was it always unfair against you? You were talented and skilled. You worked hard. But even with all you accomplished, there was always somebody more richly rewarded than you. More fame. More wealth. More attention. An easier life. It left a bitter taste in your mouth; it soured everything you had, diminishing you and your accomplishments. It seemed like the world was always against you. And then you ended up in the Nine Hells, proving death was against you too.


\begin{description}
\item[Character Punishment.]
You have strange dreams not connected to your life. In them, you're a lowly servant... yet somehow your sister is the Queen, and she is beloved. Other servants, guards, diplomats, members of the court—they all go on and on about how smart she is, how beautiful, how perfect, while you toil in the shadows, forgotten and ignored.

\item[Phylactery Location.]
The War-Slough, on Avernus, the first layer of the Nine Hells. Each day, the character's soul must go from room to room in a castle, clad in soiled rags, and clean the opulent chambers under the watchful eye of the Queen and her followers. They constantly deride and chastise the soul for the poor job they're doing. Often, they abuse the character, beating them with sticks and switches to spur them to work harder. To free their soul, characters must defeat the Queen... but even if she is slain, she simply reappears on her throne, completely unharmed. The only way to truly defeat the Queen is to first destroy her throne.

\end{description}

\section{Group Patrons}
Group patrons help player characters work together toward a common goal. They're powerful backers who give adventurers a clear purpose. They also provide rare resources such as magic items, contacts, and other hard to find things.

This section describes three group patrons. Characters in search of a good-aligned patron should choose the Hellriders of Elturel, while those who prefer a neutral patron should choose the Conclave of Halruaa. Evil characters may be drawn to the third patron option, the Deathstalkers of Bhaal.

Group patrons have the following facets:


\begin{description}
\item[Quest.]
Each group patron has an overarching goal that its members seek to accomplish by giving the characters a quest, which you'll find in the patron's description. As they pursue this quest, characters might receive other missions from their primary contact in their group patron.

\item[Perks.]
Group patrons offer access to resources not readily available to most people. These include magical equipment, secret information, and training. Each group patron gives a different set of perks. Most group patrons also give access to magic items for purchase or recipes on how to craft them.

\item[Contacts.]
Each group patron has one or two primary contacts that interact with the player characters.

\end{description}

\begin{DndSidebar}[float=!b]{Resolving Group Patron Objectives}
The characters can stay in contact with their group patron throughout the adventure. Once the characters complete the quest given to them by their group patron, they'll have a final meeting with their group patron's representative in Cania. See "A Thankful Patron" in chapter 10 for details.



\end{DndSidebar}

\section{Hellriders of Elturel}

\begin{description}
\item[Alignment.]
Good

\end{description}

The Hellriders are an elite cavalry unit whose members act as the primary armed force of the city of Elturel. They're one of the most renowned and well-regarded military forces in the Realms. It is said that a company of Hellriders once rode into Avernus, the first layer of the Nine Hells, and from this story, the Hellriders were named.

This story is true, for a century ago a gateway to the Nine Hells opened in the Fields of the Dead north of the city of Elturel. The skilled knights of Elturel fought a losing battle against the Fiends that poured from this gate. The High Rider of Elturel implored his people to pray to the gods for help. Their prayers were answered when the angel \textbf{Zariel} arrived in Elturel with a small host of allies. She rallied and trained an army of thousands and then led them into the gateway to the Nine Hells. While she was successful at destroying the gate, \textbf{Zariel} herself was captured by \textbf{Asmodeus} and corrupted into the archdevil who now rules over Avernus. The knights of Elturel were renamed to the Hellriders in dedication to \textbf{Zariel} and the warriors who followed her into the Nine Hells.

The Hellriders have learned that some of \textbf{Zariel}'s allies are still alive in the Nine Hells. They remain uncorrupted but suffer horrific tortures. The Hellriders feel duty bound to rescue them.

\subsection{Hellrider Quest}
The Hellriders know the location of Celestial beings who are trapped within the Nine Hells. If rescued, these beings grant the characters powers that should help them in their own personal quests:


\begin{description}
\item[Anagwendol.]
She is kept prisoner in Phlegethos at the Elemental Preserve.

\item[Barachiel.]
He is a prisoner of the Brothers Adramlech and Morax and can be found on their infernal warship prowling the River Styx.

\item[Jenevere.]
She is kept prisoner in Maladomini at the Eye Market.

\end{description}

Each time the characters rescue one of these Celestial beings, the being bestows a supernatural gift as a reward. Each character may choose one of the charms described in the "Supernatural Gifts" section of the \textit{Dungeon Master's Guide}.

\subsection{Hellrider Perks}
If the characters choose the Hellriders as their group patron, each receives the following benefits.

\subsubsection{Hellrider's Blessing}
The Hellriders know better than most the evils found in the Nine Hells and would never send anyone there without protection. After spending the entirety of a long rest in the Nine Hells, you gain the benefit of a \textit{Greater Restoration} spell.

\subsubsection{Celestial Cavalry}
While much of the journey through the Nine Hells follows the River Styx, it is likely that you will occasionally leave on excursions or scouting missions. While you remain in the Nine Hells, you may cast the \textit{Phantom Steed} spell at will, as a ritual.

\subsubsection{Hellrider's Salvation}
You're taught a ritual to help save missing Celestials. The ritual requires 1 minute to complete and targets all Celestials within 100 feet. Once the ritual is complete, if the Celestial is willing, they're transported to a location on the Material Plane of their choosing.

\subsection{Hellrider Contact: Ramius}
\textbf{Ramius} Dangremond is the commander of the Hellriders of Elturel. He never wanted the position, but because of the many deaths within the ranks of the order, he had to step up. His friend and fellow Hellrider, \textbf{Barachiel}, was the main reason he was able to grow into a capable commander. The loss of his friend saddens \textbf{Ramius}, and the news of his imprisonment within the Nine Hells gnaws at his soul.

\textbf{Ramius} (see the accompanying stat block) is a natural horseman and a gifted athlete. The soldiers he leads often say that there isn't anything that \textbf{Ramius} can't mount or ride. He's ridden horses, mammoths, and machines into battle. Loyalty, honor, and duty are virtues that \textbf{Ramius} upholds. He fears nothing other than losing those he cares for.

\subsubsection{Ramius as a Contact}
\textbf{Ramius} doesn't have the personal ability to cast \textit{Plane Shift}—instead, he has been given a magic item that allows him to visit with the party every few days. When he arrives, he brings general supplies such as food and water, and treats any wounds, curses, or diseases the party has (to the best of his ability). Afterwards, he offers a selection of useful magic items for purchase, or, if the party already has excess magic items, he offers to purchase them for keeping or disposal.

\begin{DndTable}[header=Purchasing Items from Ramius]{XX}

Item & Cost\\
Adventuring supplies & Original price\\
\textit{Amulet of Health} & 12,000 gp\\
\textit{Bracers of Defense} & 12,000 gp\\
\textit{Cloak of Protection} & 600 gp\\
\textit{Gauntlets of Ogre Power} & 550 gp\\
\textit{Holy Avenger} & 175,000 gp\\
\textit{Horseshoes of a Zephyr} & 18,000 gp\\
\textit{Horseshoes of Speed} & 6,500 gp\\
Horn of Valhalla, Silver & 25,500 gp\\
\textit{Javelin of Lightning} & 500 gp\\
\textit{Mace of Smiting} & 9,000 gp\\
\textit{Mithral Armor} & 700 gp\\
\textit{Ring of Protection} & 12,000 gp\\
+1 Shield & 450 gp\\
+1 Shield & 12,000 gp\\
+1 Shield & 40,000 gp\\
\textit{Sun Blade} & 35,000 gp\\
\end{DndTable}

\begin{DndTable}[header=Selling Magic Items to Ramius]{XX}

Rarity & Value\\
Uncommon & 250 gp\\
Very rare & 10,000 gp\\
Rare & 2,500 gp\\
Legendary & 55,000 gp\\
\end{DndTable}

\section{Conclave of Halruaa}

\begin{description}
\item[Alignment.]
Neutral

\end{description}

Halruaa is a land of magic, renowned for the fact that a significant portion of its ruling elite are archmages. Created by archwizards foreseeing the fall of Netheril, Halruaa combines peace and harmony with the magic powers of their ancestors, but without the taint of their ambitions.

The Conclave is a secret society within Halruaa that enforces the laws regarding the use, abuse, and research of magic in Halruaa.

The membership of the Conclave is a mystery to everyone outside their ranks. Others within Halruaa know of the Conclave's existence, but not the identities of its members.

Whenever the experiments, practices, or other actions of a spellcaster endanger the future of Halruaa, the Conclave intervenes. This intervention is usually a warning, but a spellcaster who ignores the warning is eliminated or simply made to disappear.

The Conclave has become interested in the events put into motion by \textbf{Asmodeus} because their diviners have foreseen that the archdevil's scheme will result in a dangerous soul being unleashed.

\subsection{Conclave Quest}
The Conclave wants the characters to track down a soul they refer to as the Unmaker. They're to use a ritual to destroy the soul. To help with this task and with finding the other souls that the characters seek, the Conclave offers to use their diviners to uncover the location of powerful magic items. The Nine Hells is the graveyard for thousands of heroes and the magic items that they brought with them. Some of these include the most famed items in the multiverse.

The Unmaker is a prisoner of the archdevil \textbf{Abigor} and can be found on his infernal submersible. The Conclave provides a Conclave Coin, effectively a special Soul Coin that can contain the soul of the Unmaker. The characters must slay the Unmaker, trap the Unmaker's soul in the coin, and return the coin to the Conclave so it and the soul inside it can be destroyed.

Each character is told the location of a powerful magic item. This magic item is based off the class in which they have the highest level:


\begin{description}
\item[Artificer.]
The Manual of Iron Golems, which can be found in Stygia at the Chasm of Found Things.

\item[Barbarian.]
Belt of Storm Giant Strength, which can be found in Minauros at the Ineffable Trove.

\item[Bard.]
Instrument of the Bards, Ollamh Harp, which can be found in Dis at the Agora of Floating Knives.

\item[Cleric.]
\textit{Rod of Resurrection}, which can be found in Stygia at the Chasm of Found Things.

\item[Druid.]
\textit{Staff of the Woodlands}, which can be found in Maladomini at the Eye Market

\item[Fighter.]
Belt of Storm Giant Strength which can be found in Minauros at the Ineffable Trove.

\item[Monk/Rogue.]
\textit{Cloak of Invisibility}, which can be found in Dis at the Agora of Floating Knives.

\item[Paladin.]
\textit{Holy Avenger}, which can be found in Stygia at the Chasm of Found Things.

\item[Ranger.]
Horn of Valhalla, Iron, which can be found in Dis at the Agora of Floating Knives.

\item[Sorcerer/Warlock/Wizard.]
\textit{Staff of the Magi}, which can be found in Cania at the Sorrow Mine.

\end{description}

\subsection{Conclave Perks}
If the characters choose the Conclave as their group patron, each receives the following benefits.

\subsubsection{Arcane Communication}
The Conclave of Halruaa gives you a Sending Stone and a Halruaan Ethereal Vessel. The Sending Stone can be used to communicate with a member of the Conclave who has the matching stone.

The Halruaan Ethereal Vessel is a magic item that allows two people to attune to it. You and a member of the Conclave are attuned to the vessel. The vessel is a transparent case that can contain up to 12 cubic feet of nonliving material (3 feet by 2 feet by 2 feet). When not in use, it exists on the Ethereal Plane. While the vessel remains on the Ethereal Plane, you can use an action and recall it (as can the Conclave member also attuned to it). It appears in an unoccupied space on the ground within 5 feet of you. You send the vessel back to the Ethereal Plane by using an action.

\subsubsection{Arcane Quartermasters}
The Conclave has amassed a collection of useful magic items and material spell components. You can use your Sending Stone to request one of the items on the Conclave Items table. You must then pay the cost of the item, placing the payment in your Halruaan Ethereal Vessel and returning the vessel to the Ethereal Plane. The money is then removed by the Conclave member attuned to the vessel and replaced with the requested item, which you can then retrieve from the vessel.

\begin{DndTable}[header=Conclave Items]{XX}

Items & Cost\\
Blank scroll & 10 gp\\
\textit{Broom of Flying} & 600 gp\\
\textit{Crystal Ball} & 30,000 gp\\
Divination \textit{Spell Scroll (1st level)} & 100 gp\\
Divination \textit{Spell Scroll} (Spell Scroll (2nd level) or Spell Scroll (3rd level) level) & 500 gp\\
Divination \textit{Spell Scroll} (Spell Scroll (4th level) or Spell Scroll (5th level) level) & 2,000 gp\\
\textit{Gem of Seeing} & 10,000 gp\\
\textit{Headband of Intellect} & 400 gp\\
\textit{Helm of Comprehending Languages} & 600 gp\\
Material spell component costing up to 1,000 gp & 75\\% original cost\\
\textit{Medallion of Thoughts} & 350 gp\\
\textit{Ring of Mind Shielding} & 500 gp\\
\textit{Ring of Spell Storing} & 12,000 gp\\
\textit{Wand of Binding} & 8,000 gp\\
\textit{Wand of Magic Detection} & 500 gp\\
\textit{Wand of Wonder} & 9,500 gp\\
\end{DndTable}

\subsection{Halruaan Contact: Zythan Avhoste}
\textbf{Zythan}, like many other Halruaans, spent his early years studying the arcane arts. His talent in Divination magic made him particularly useful to the Conclave, and they quickly sought him out. Under their tutelage, his skill flourished, and he became a powerful mage. Those that know him personally recount his matter-of-fact personality, his green hair, and his third eye that seems to open when he performs magic. He rarely speaks of himself, however, choosing instead to focus his attention on the future.

When he first received an official position within the Conclave, \textbf{Zythan} wanted to prove himself. He spent countless hours peering into the future and examining potential outcomes. One fateful night, he just barely caught sight of a forbidden ritual being prepared—a ritual that could cost the lives of many Halruaans. Because of his sight, the Conclave was able to intervene in time, and he was quickly promoted and showered with praise.

For the past few months, \textbf{Zythan} has felt like something was off: at first, he suspected another ritual within Halruaa. He requisitioned the help of the other diviners, as well as a powerful Crystal Ball, and together they were able to piece together a vision. The dark lord of the Nine Hells, \textbf{Asmodeus}, will unleash a great evil very soon. Though the Conclave often stays out of non-Halruaan affairs, they had no choice but to act. It is for this reason that the adventurers were called to Halruaa.

\textbf{Zythan}'s (see the accompanying stat block) skill in arcane magic, while particularly focused on Divination, extends to every school. While he lacks proficiency in martial combat, he can wield both offensive and defensive magic effectively. Outside of his magical talent, he is also known to be highly intelligent, though his social skills could use work. Much of his time is spent viewing alternate futures and preparing for potential dangers, leaving little for himself. He often goes without meals until he starts to feel weak, and if he had friends, they might be concerned about how much time he spends working.

\subsubsection{Zythan as a Contact}
Using his Divination magic, \textbf{Zythan} keeps tabs on the party, occasionally checking in when he deems appropriate. When he checks in, he can provide access to some magic items. He will also purchase any excess Spell Scrolls or spellbooks found by the party, and can, upon request, attempt to hunt down specific magic items. If successful, he brings the desired magic items on his next visit.

\begin{DndTable}[header=Purchasing Items from Zythan]{XX}

Item & Cost\\
\textit{Amulet of the Planes} & 35,000 gp\\
\textit{Carpet of Flying} & 40,000 gp\\
Ioun Stone, Reserve & 8,500 gp\\
\textit{Mantle of Spell Resistance} & 10,500 gp\\
\textit{Ring of Resistance} & 9,000 gp\\
\textit{Staff of Power} & 45,000 gp\\
\textit{Wand of Fireballs} & 10,000 gp\\
+1 Wand of the War Mage & 500 gp\\
+1 Wand of the War Mage & 9,500 gp\\
+1 Wand of the War Mage & 32,000 gp\\
\end{DndTable}

\begin{DndTable}[header=Selling Items to Zythan]{XX}

Items & Cost\\
Spellbooks & 100 gp per level of spell found within the book. Additional 2,000 gp if the book contains one or more 9th-level spells\\
Spell Scrolls (2nd to 4th level) & 250 gp\\
Spell Scrolls (5th to 6th level) & 900 gp\\
Spell Scrolls (7th to 8th level) & 2,250 gp\\
Spell Scrolls (9th level) & 6,000 gp\\
\end{DndTable}

\section{Deathstalkers of Bhaal}

\begin{description}
\item[Alignment.]
Evil

\end{description}

Bhaal, also known as the Lord of Murder, is the god of assassins, killers, and murder.

During the Time of Troubles, most gods, including Bhaal, were forced to walk the world as mortals and lost all their godly powers. During this period, Bhaal was slain by the then-mortal Cyric. Bhaal, having foreseen his own death, sired many children, each of whom carried a piece of his divine essence within them. These children are referred to as the children of Bhaal, or Bhaalspawn, and were created to ensure Bhaal's resurrection. This did not come to pass until more than a century after the Time of Troubles. Even though Bhaal was reborn, he was a shadow of his former self and tied to the Material Plane.

During the time that he was dead, Bhaal had many scattered cults, the most infamous being the Deathstalkers of Bhaal, who were dedicated to bringing him back to life. With Bhaal now returned, the Deathstalkers are now focused on finding a way to return Bhaal to being a fully divine being. They've discovered a way by making a deal with \textbf{Asmodeus}. However, they believe the Lord of the Nine Hells intends to deceive them and now they seek a way to make sure \textbf{Asmodeus} holds up to his end of the bargain.

\begin{DndSidebar}[float=!b]{Uncomfortable Allies}
The Deathstalkers of Bhaal are responsible for the murders that doomed the characters (or their loved ones) to the Nine Hells. The characters don't have much reason to want anything to do with them let alone align themselves with these murderers. However, the Deathstalkers are a powerful ally. Having begun to doubt their alliance with \textbf{Asmodeus}, they now attempt to gain the cooperation of the characters. \textbf{Sarevok}, the Deathstalker that contacts the heroes, is quite convincing. And scary...



\end{DndSidebar}

\subsection{Deathstalker Quest}
To force \textbf{Asmodeus} to make good on his word, the Deathstalkers of Bhaal want to gain leverage. For this they need a powerful artifact. They ask the characters to steal some of the most infamous and powerful artifacts of the Nine Hells. They need only one, leaving any other artifacts for the characters to decide what to do with. For a detailed description of these artifacts, see the section "Infernal Artifacts" in appendix D.

\subsection{Deathstalker Perks}
If the characters choose the Deathstalkers as their group patron, each receives the following benefits.

\subsubsection{Death Token}
Each member of your party is given a Token of Bhaal. For its magic to work, the token must be in your possession. The first time you fail a death saving throw that doesn't kill you outright, the token magically transports you, along with any equipment you're wearing or carrying, to a demiplane that resembles an empty chamber with no exits. Any ally within 30 feet of you is also transported to the demiplane, along with any equipment it is wearing or carrying, if the ally has 0 hit points and is either stable or dying when the token activates. Each creature that appears in the demiplane regains 10 hit points. Any creature that starts its turn in the demiplane is transported back to the space it left, or, if that space is occupied, the nearest unoccupied space. The token can be used only once, after which it turns to dust.

\subsubsection{Deathstalker's Blessing}
Your party receives the blessing of Bhaal. This manifests by making you deadliest when you get the drop on your enemies. You have advantage on attack rolls against any creature that hasn't taken a turn in combat yet. In addition, any creature that hasn't taken a turn in combat has disadvantage on saving throws against your spells. The blessing applies only while in the Nine Hells.

\subsubsection{Soul Coin Market}
Much like the archdevils of the Nine Hells, Bhaal desires souls. Specifically, he desires Soul Coins (see "Soul Coin" in appendix D). When you're in possession of one or more Soul Coins, you can whisper to the coin what it is that you desire from the Deathstalker Items table. If you have the requisite number of coins, the Soul Coins vanish from your possession, and the item that you asked for appears within 5 feet of you.

\begin{DndTable}[header=Deathstalker Items]{XX}

Items & Cost\\
\textit{Dagger of Venom} & 6 Soul Coin\\
\textit{Eyes of Charming} & 1 \textit{Soul Coin}\\
\textit{Goggles of Night} & 1 \textit{Soul Coin}\\
\textit{Hat of Disguise} & 1 \textit{Soul Coin}\\
Iron Bands of Bilarro & 7 Soul Coin\\
\textit{Nine Lives Stealer} & 15 Soul Coin\\
\textit{Pipes of Haunting} & 1 \textit{Soul Coin}\\
Potion of Poison (2) & 1 \textit{Soul Coin}\\
\textit{Slippers of Spider Climbing} & 1 \textit{Soul Coin}\\
\textit{Staff of Withering} & 5 Soul Coin\\
\textit{Sword of Life Stealing} & 8 Soul Coin\\
\textit{Wand of Secrets} & 1 \textit{Soul Coin}\\
\end{DndTable}

\subsection{Deathstalker Contact: Sarevok}
\textbf{Sarevok} Anchev is a powerful Deathstalker. He is also one of the Bhaalspawn, the mortal offspring of the dead god Bhaal. \textbf{Sarevok} attempted to reclaim the divine seat of the Lord of Murder vacated by his immortal father's demise, but his plans were thwarted when he was slain by his half-brother, Abdel Adrian, who rejected his heritage and fought against his Bhaalspawn siblings.

\textbf{Sarevok}'s spirit was sent to the Abyss as punishment. There he eventually crossed paths with Abdel a second time when his noble-hearted brother ventured into the lower realms on a dangerous quest to stop another Bhaalspawn named Melissan. \textbf{Sarevok} agreed to help Abdel kill Melissan, on the condition that Abdel helped him escape the eternal torments of the Abyss.

Abdel agreed, and \textbf{Sarevok} was reborn into the mortal world. After his rebirth, \textbf{Sarevok} was true to his word, and the two brothers fought side by side against their half-sister. Ultimately Melissan was defeated, and \textbf{Sarevok} was granted a second chance at life.

With his prodigious strength, his legendary skill in battle, and the Sword of Chaos—a life-stealing, enchanted blade—\textbf{Sarevok} (see the accompanying stat block) became a famous (and feared) mercenary. Yet his many accomplishments brought him no joy. He felt no thrill at victory in battle, no delight in the routing of his enemies. The power he accumulated was bitter as ashes on his tongue, and he became a man haunted by his former life. The realization that no earthly achievements could ever compare to what he once almost had—immortality and godhood—left him broken and empty.

\textbf{Sarevok} plunged into a deep despair. To numb his pain, he indulged in every vice imaginable, squandering his wealth and health on alcohol, drugs, and fleeting comforts. While his divine heritage slowed his aging, it did not stop it entirely, and after decades of self-abuse he was eventually reduced to an old man begging in the streets of Baldur's Gate.

This was how his father—Bhaal, the reborn god of murder—found him. Bhaal recognized his own divine spark in the pathetic old man, and he sensed \textbf{Sarevok} still had potential. Bhaal recruited him to become the high priest of his fledgling clergy, giving \textbf{Sarevok} new purpose... and another chance to become an agent of death and destruction.

\subsubsection{Sarevok as a Contact}
Though he lacks Divination magic, \textbf{Sarevok} keeps tabs on the party's progress and checks in whenever he feels like it. During his check-ins, he will offer a small assortment of magic items, as well as a large variety of potions. He also has a vast treasury and is willing to purchase Soul Coins for 1,000 gp per coin. In addition, for the price of 10,000 gp, he grants one character a blessing (see "Supernatural Gifts" in the \textit{Dungeon Master's Guide}) that allows them to become a host for his god, Bhaal.

\subparagraph{Unholy Blessing of Bhaal}
You may use an action to transform yourself into the Slayer. This works exactly as the \textit{Polymorph} spell, except the form chosen uses the accompanying Slayer stat block. After you use this blessing, you can't use it again until you have finished a long rest.

\begin{DndTable}[header=Purchasing Items from Sarevok]{XX}

Items & Cost\\
+1 Ammunition & 150 gp\\
+1 Ammunition & 2,500 gp\\
+1 Ammunition & 10,000 gp\\
Bead of Force & 3,000 gp\\
\textit{Cape of the Mountebank} & 6,000 gp\\
\textit{Robe of Eyes} & 7,500 gp\\
\textit{Vicious Weapon} & 8,000 gp\\
Uncommon potions & \\
Potion of Animal Friendship & 250 gp\\
Potion of Hill Giant Strength & 250 gp\\
Potion of Growth & 250 gp\\
Potion of Resistance & 250 gp\\
Potion of Water Breathing & 250 gp\\
Rare potions & \\
Potion of Clairvoyance & 1,000 gp\\
Potion of Diminution & 1,000 gp\\
Potion of Gaseous Form & 1,000 gp\\
Giant Strength (Potion of Frost Giant Strength/Potion of Stone Giant Strength) & 1,000 gp\\
Potion of Heroism & 1,000 gp\\
Potion of Mind Reading & 1,000 gp\\
Very rare potions & \\
Potion of Flying & 5,000 gp\\
Potion of Fire Giant Strength & 5,000 gp\\
Potion of Invisibility & 5,000 gp\\
Potion of Speed & 5,000 gp\\
\end{DndTable}

\chapter{Kelemvor's Cathedral}
\DndDropCapLine{K}{}elemvor is a god of the dead, and the cathedral dedicated to him reflects the somber nature of the god. The character's face their first challenge in this dimly-lit building.


\section{Running This Chapter}
Before running this chapter:


\begin{itemize}
\item Familiarize yourself with this chapter and \textbf{Koh Tam}, who plays an important role in the adventure.
\item Each player must separately decide which soul they need to save from the Nine Hells.
\item As a group, the players must decide which of the group patrons they wish to ally with.
\end{itemize}

\section{The Infernal Contracts}
The characters meet with their group patron. They learn about \textbf{Asmodeus} and how he has tricked the powerful into signing infernal contracts. They learn why their group patron is concerned about this.

Read the following text to start things off. Then read the text that is specific to the group patron.

\begin{DndReadAloud}{}
You knew as much of \textbf{Asmodeus} as anyone in your line of work. He was a powerful deity, a conniving master of many schemes, and the Lord of the Nine Hells. You knew that he often plotted to corrupt and win the souls of the influential and the powerful through the use of infernal contracts, damning those unfortunates to the torments of his realm. You never thought that you would be one of his targets.

Now, you live with the loss. A hollowness exists where your soul should be, or else the fierce grief-ache of a stolen friend or family member. If \textbf{Asmodeus} could not convince you yourself to sign a contract, he went after those closest to you. If you made a pact with the Lord of the Nine, you did it with the assumption that you would have enough time to fulfill your end of the bargain. Likely the others who fell prey to him thought the same. It was not to be. \textbf{Asmodeus}'s hired murderers ensured it.

Since then, you have been searching desperately for a way to recover what you lost, without success. Every lead has turned up nothing. Every promise of a gateway into the infernal realm has been false. It has begun to seem truly hopeless... But that is about to change.

\end{DndReadAloud}

\subsection{Hellriders of Elturel Introduction}
The characters meet with \textbf{Ramius} in the temple of Helm within the city of Elturel. If Elturel wasn't saved from the Nine Hells, then the meeting takes place in the city of Berdusk, one of the cities that was ruled from Elturel.

\begin{DndReadAloud}{}
For months, you have been searching for information about the Nine Hells and looking for allies that can help your party venture there. After hearing about the Hellriders of Elturel and their history with \textbf{Asmodeus}, you have been granted an audience.

The temple of Helm in Elturel is a grand building of marble and gilded stone. "Heroes, well met," comes a voice from the entrance of the temple. A man in armor strides toward you. "I am \textbf{Ramius} of the Hellriders of Elturel, and I have a proposal for you."

\textbf{Ramius} beckons you to follow him to one of the temple walls, which bears a depiction of a great cavalry unit descending upon a gateway of flame. At the head of the painted cavalry is a winged woman in shining armor.

"In Elturel's darkest hour," \textbf{Ramius} continues. "Our prayers were answered by the angel \textbf{Zariel}, who led a thousand riders and a host of celestial beings into the gateway to the Nine Hells. The battle was won, the gateway destroyed, but \textbf{Zariel} and some of her allies were captured by the Lord of the Nine, \textbf{Asmodeus}. \textbf{Zariel} was lost. But we have learned that a few of her allies remain, held captive and tormented in the Nine Hells." The Hellrider pauses and then adds quietly, "My dearest friend among them."

\textbf{Ramius} collects himself and says, "I propose that together you and I journey through the Nine Hells and rescue \textbf{Zariel}'s allies, my own dear friend, and the souls you yourselves have lost."

If you agree, I will take you to the temple of Kelemvor. There is a priest there who can help us."

\end{DndReadAloud}

After the introductions, \textbf{Ramius} uses a teleportation circle to take them to the temple of Helm in Waterdeep and from there to the cathedral of Kelemvor.

\subsection{Conclave of Halruaa Introduction}
The players meet \textbf{Zythan} in Halarahh, the capital city of Halruaa.

\begin{DndReadAloud}{}
As you traverse the many market stalls of the central plaza of Halarahh, you see the shadowed pavilion described to you by your employer. Inside, \textbf{Zythan} of the Conclave of Halruaa, a tall and slender man, wastes no time with greetings, beckoning you to follow him over to a crystal ball in the center of the pavilion. A month ago, you each received a letter from \textbf{Zythan} inviting you to the city of Halarahh. He claimed to have a lucrative proposition for you, something that might bring back the soul that has been lost to each of you.

\textbf{Zythan} passes a hand over the crystal ball, and suddenly you see cruel flames and strange figures moving within. "Through the power of Divination, I have discovered that the Lord of the Nine Hells, \textbf{Asmodeus}, will very soon unleash a terrible evil, a soul that we call the Unmaker."

The flames in the crystal ball are overtaken by a pitch-black darkness. Cold eyes glare out from the void. "To prevent this disaster, you must travel to the Nine Hells, track down the Unmaker, and perform a ritual the Conclave has prepared to destroy him. To aid you in this task, we will use our powers to reveal the locations of powerful magic items hidden throughout the Nine Hells."

The crystal ball shows a series of objects wreathed in fire: a staff, a harp, a cloak, and a sword, among others. "The Nine Hells are full of dangers," \textbf{Zythan} says gravely. "And you know well the conniving schemes of \textbf{Asmodeus}. As you travel through the Nine, you may have the opportunity to restore those souls that were taken from you."

\textbf{Zythan} passes his hand over the crystal again and the images within fade away.

"This is no easy quest, but the fate of many depends upon it. If you will come with me to the temple of Kelemvor, I have a contact there who is an expert on the Nine Hells."

\end{DndReadAloud}

Once \textbf{Zythan} has finished describing the mission, he teleports himself and the characters to the cathedral of Kelemvor in Waterdeep.

\subsection{Deathstalkers of Bhaal Introduction}
The players meet \textbf{Sarevok} in Waterdeep.

\begin{DndReadAloud}{}
It is not the finest tavern in Waterdeep. The tabletops are grimy, the chairs a spare breath from collapse, and you've spied at least one rat scampering across the floor. But the drink is cheap, and the patrons mind their business, which is exactly what your party needs. On quiet nights like these, you each feel more keenly the ache of your losses.

Your contact in the city has set up this meeting. For months, you have all been searching for allies in your quest and for a way in to the Nine Hells. Finally, after all this time, your efforts may have been rewarded. The tavern door opens, letting in a burst of cold. Footsteps follow and come directly to your table. The footsteps belong to a man who wears the robes of a high priest, though you don't recognize the symbol of his clergy. He sits down at your table, looks at you with his gleaming eyes, and speaks.

"I am \textbf{Sarevok} Anchev, high priest of the Deathstalkers of Bhaal. On my lord god's behalf, I seek to gain leverage over the deceiver, \textbf{Asmodeus}, the Lord of the Nine Hells. I know that each of you has cause against \textbf{Asmodeus}. I offer a mutually beneficial opportunity." There is something about \textbf{Sarevok} that tells of death and destruction. But the zeal in his eyes is entrancing.

"Throughout the Nine Hells there are many powerful and ancient artifacts that \textbf{Asmodeus} values, the prized possessions of his archdevils. It is the desire of the Deathstalkers for you to steal these artifacts. For your reward, we offer the artifacts themselves, save one, to do with as you please."

"If you agree, I will take you to the temple of Kelemvor. There you will be advised on how best to traverse the Nine Hells."

\end{DndReadAloud}

After speaking with the characters, \textbf{Sarevok} takes them to the cathedral of Kelemvor in Waterdeep. If the characters inquire about the artifacts \textbf{Sarevok} mentioned, he says it doesn't matter which artifact the characters surrender to the Deathstalkers. No matter the artifact, the Deathstalkers plan to use it to gain leverage in their dealings with \textbf{Asmodeus}.

\begin{DndSidebar}[float=!b]{Meeting the Others}
At this point, just prior to heading towards the cathedral of Kelemvor, if there are souls of the damned among the characters, the patron should introduce them to the other characters.



\end{DndSidebar}

\section{The Cathedral}
\begin{DndReadAloud}{}
The cathedral of Kelemvor is bone-cold and dark. The walls and floor are black stone, and there are motifs of grinning skulls in every direction. Deep set into the floor are pools of viscous liquid, as bright and red as arterial blood. The cathedral is still and silent, but you faintly hear the distant click-clack of the bones that hang from the vaulted ceiling, knocking gently together in the musty air.

\end{DndReadAloud}

\textbf{Koh Tam} has been waiting for the characters to arrive. He has prepared a ritual that should tell the characters where to find their group patron's objective, along with the soul they're each searching for.

\subsection{Koh Tam}
\textbf{Koh Tam} is an imposing figure dressed in many layers of colorful robes and heavy armor pieces. He wears a golden mask that covers the upper part of his face, set with strange green gemstones that wink and glitter like eyes in the flickering cathedral light. He carries a wooden staff with a golden hand affixed upon it. \textbf{Koh Tam}'s voice has a rasping hiss to it, like wind through skulls' teeth, but he speaks kindly and with steadfast patience.

\subsection{The Possession}
\textbf{Koh Tam}'s plan is to allow himself to be possessed by \textbf{Baalzebul}, who, owing to the divine wards provided by \textbf{Koh Tam}'s ritual, can't tell a lie. Once he has been possessed, the characters must ask him questions.

\begin{DndReadAloud}{}
"It is true," \textbf{Koh Tam} says to your party. "That among those on this mortal plane, there are none that know the Nine Hells as I do. But the Nine Hells are immense, full of countless secrets and dangers. To find that which you seek, you must consult someone with even greater expertise."

\textbf{Koh Tam} spreads out his arms, gesturing at the ritual he has prepared in the center of the cathedral sanctum. "I will summon for you \textbf{Baalzebul}, one of the archdevils of the Nine. He will possess my body and speak through me. He cannot lie, but be warned, he will use all his wit to escape your questioning or to misdirect you."

Turning from you, \textbf{Koh Tam} begins his ritual. He chants ancient words and lifts his wooden staff, and you feel the air in the cathedral shudder with the growing power. The liquid in the pools around you rises. Thick ribbons of bloody red swirl around \textbf{Koh Tam}, not touching the priest, but obscuring him from your sight as the chant reaches a crescendo.

A growl escapes \textbf{Koh Tam}'s mouth. "The rancid stench of sinners... surrounds me. Speak to me of your sins, and swiftly, so that I might tell you what I must."

\end{DndReadAloud}

Each character now has the opportunity to ask after the sins they seek. \textbf{Baalzebul} refuses to answer any questions until each character admits their sins or the sins of their loved one, however. Once they've done this, \textbf{Baalzebul} reveals where the soul can be found. He provides only the layer of the Nine Hells and the general location (an adventure area, district, or landmark within the layer). The players should make note of this information.

At this point, the group patron contact steps forward and asks for information specific to their objective.

\subsubsection{Hellriders}
If the characters chose the Hellriders as a group patron, read the following:

\begin{DndReadAloud}{}
"Poor \textbf{Barachiel} is held in a moving cage upon the River Styx, a guest of the great Brothers. The unfortunate \textbf{Jenevere} languishes somewhere among the broken streets of the Eye Market of Maladomini. And once-mighty \textbf{Anagwendol}, captured by the great hunter himself, is at the heart of Kordichai's famed hunting grounds."

\end{DndReadAloud}

\subsubsection{Conclave of Halruaa}
If the characters have the Conclave of Halruaa as a group patron, read the following:

\begin{DndReadAloud}{}
"You may find the Unmaker beneath the murky waters of the River Styx. In silence, moving. In darkness, moving. Entombed in a ship that does not sail, guarded by the great and monstrous \textbf{Abigor}, who commanded the legions of Maladomini up in Avernus to stem the tides of the Abyss."

\end{DndReadAloud}

\subsubsection{Deathstalkers}
If the characters followed \textbf{Sarevok} into the cathedral, read the following:

\begin{DndReadAloud}{}
"Five items you seek, scattered across the Nine Hells. In Minauros, search \textbf{Mammon}'s boggy trove. One of two you may find in Phlegethos, if you dare brave the hunting grounds of Kordichai. In Malbolge, the hag-run inn might host one of your treasures. In Stygia, explore the icy Chasm of Found Things. And in Dis, cast your gaze upon the stalls and palace halls of the agora."

\end{DndReadAloud}

\subsection{The Exorcism}
Once \textbf{Baalzebul} has answered, things go wrong:

\begin{DndReadAloud}{}
\textbf{Koh Tam} begins to tremble. His staff falls to the cathedral floor with a clatter, and his hands reach for his head. The ribbons of red liquid fall back into their pools, sending up mighty sprays of staining droplets.

"Prepare, adventurers!" \textbf{Koh Tam} says with a gasp, his voice once more a dry hiss instead of the archdevil's terrible growl. "He is attempting to escape the confines of my body. You must exorcise him back to the Nine Hells!"

The priest's voice is suddenly overcome with ugly, booming laughter as an avatar of the archdevil begins to appear, rising out of \textbf{Koh Tam}'s shaking form.

\end{DndReadAloud}

The characters must now exorcise \textbf{Baalzebul} from \textbf{Koh Tam}. This is a two-stage battle. In the first stage, they must defeat the \textbf{avatar of Baalzebul} (see the accompanying stat block) that has partially crawled out of \textbf{Koh Tam}.

Once the avatar has been defeated, the second stage begins. \textbf{Koh Tam} (see appendix C) is possessed and attacks. The characters must free him by casting a Dispel Evil and Good or \textit{Greater Restoration} spell on him, or by dealing radiant damage to him. Feel free to allow other strategies to work if they make sense to the situation. A cleric's Divine Intervention, for example, banishes \textbf{Baalzebul} back to the Nine Hells immediately.

Once the characters rid \textbf{Koh Tam} of the avatar of \textbf{Baalzebul}, the threat is ended.

\section{Entering the Nine Hells}
Once the threat posed by \textbf{Baalzebul} has been dealt with, \textbf{Koh Tam} takes a little time to recover and then talks to the characters about the next steps.

\begin{DndReadAloud}{}
The weary priest sits on the steps below the cathedral altar. His hands still tremble but his voice is calm and steady as he explains the journey that lies ahead.

"I can open a portal to Avernus, the first of the Nine Hells. Once there, I have a barge that will allow us to sail the River Styx, a far easier task than attempting to navigate by foot. I will sail with you and be your guide," \textbf{Koh Tam} says. "To succeed in your mission, you must do two things. First, you must find the souls you seek and free them from whatever foul device holds them captive. Second, you must deal with \textbf{Asmodeus} himself. For even if you escape the Nine Hells with the souls you wish to rescue, they will only be hunted down again unless you can convince, or trick, \textbf{Asmodeus} to declare them free."

\textbf{Koh Tam} stands and dusts off his heavy robes. "I will go and prepare for our journey," he says. "When you are ready, come find me and I will open the portal."

\end{DndReadAloud}

\textbf{Koh Tam} waits for as long as the characters need. Once they're ready, he opens the portal to Avernus.

\section{Meeting Tiax}
The portal takes \textbf{Koh Tam} and the characters to the Bronze Citadel in Avernus.

\begin{DndReadAloud}{}
The heat that greets you on the other side of the portal is at once sharp and searing and unbearably heavy, each push forward as difficult as swimming through molasses. The air is so dry that after a single breath your throat aches for drink and yet so humid that your clothes already cling to your sweat-soaked form. Above you the sunless sky is a brilliant red, alive with writhing masses of crimson clouds and flares of orange lightning. Your party stands on a rocky cliffside. Below you the dark and trepid River Styx stretches out to the ends of the horizon. All around you loom towers and walls of glinting metal and bleached bone: the Bronze Citadel.

\end{DndReadAloud}

\textbf{Koh Tam} guides the characters away from the Bronze Citadel to the docks. There he introduces them to his first mate, \textbf{Tiax} (see appendix C).

\begin{DndReadAloud}{}
\textbf{Koh Tam}'s first mate is a gnome of peculiar mannerism by the name of \textbf{Tiax}. He is dressed like a naval captain that has fallen upon hard times. His face is covered by a wild beard and a wooden pipe held between his teeth sends out plumes of foul-smelling smoke. He stands on the docks by \textbf{Koh Tam}'s barge, arms crossed and face set in a scowl as you approach.

"There ye are, ye fools!" \textbf{Tiax} slams a foot down on the edge of the barge. The force pushes the vessel out to the end of its rope, leaving the gnome unbalanced and tipping toward the unnatural river water. \textbf{Koh Tam} reaches out and snatches the back of \textbf{Tiax}'s shirt, pulling him to safety. The gnome neither offers his thanks nor pauses his speech. Instead, he stomps his foot again, this time on the firm wood of the dock, raises his fist, and says, "So, you wish to sail the treacherous waters of the Nine Hells with the mighty sorcerer \textbf{Tiax}. Well! No better choice of guide could you have made. For the great and wise \textbf{Tiax} knows all! But, beware, for the way before ye is full of horrors that would drive lesser souls to terror. Not I though," he says, crossing his arms again. "I'm frightened of nothing."

\end{DndReadAloud}

\textbf{Tiax} is a cleric of Cyric, the god of lies. He believes that his god has chosen him to one day rule Faerûn. Over the years this preposterous belief has led to imprisonment half a dozen times, usually when one of his schemes blows up in his face (often literally). However, Cyric does indeed favor him and treats him like a pet or jester. Thus, the god always helps the gnome escape from whatever predicament he has brought down on himself.

Cyric has directed \textbf{Tiax} to serve \textbf{Koh Tam} faithfully. Cyric is interested in what \textbf{Asmodeus} has planned and harbors a deep hatred for \textbf{Koh Tam}'s god, Kelemvor. He hopes that \textbf{Tiax} might learn something about \textbf{Asmodeus}'s schemes and if not, then \textbf{Tiax} is well positioned to kill one of Kelemvor's most powerful clerics.

\textbf{Tiax} hides his faith from everyone. He proclaims himself a great sorcerer so that any magic he uses isn't thought of as divinely given. \textbf{Tiax} is a ticking time bomb and as the characters travel deeper into the Nine Hells he becomes more dangerous. \textbf{Tiax} will randomly appear to help—or hinder—the characters as they make their way through the Nine Hells. Every time the characters arrive at a layer of the Nine Hells, even if they've previously visited it, roll on the Troubles with \textbf{Tiax} table and run the indicated event.

Eventually \textbf{Tiax} will betray \textbf{Koh Tam} at the Falls of the Frozen Titan in Cania. See the "Tiax's Betrayal" section in chapter 10 for more details.

\begin{DndTable}[header=Troubles with Tiax]{cX}

d20 & Event\\
1–6 & Nothing unusual happens\\
7–10 & The group patron checks in\\
11–13 & "\textbf{Tiax} always knows best!"\\
14–15 & \textbf{Tiax}'s shady dealings\\
16–20 & "Do not question \textbf{Tiax} the Mighty!"\\
\end{DndTable}

\subsection{Nothing Unusual Happens}
Events in this layer proceed as normal with \textbf{Koh Tam} being the primary contact.

\subsection{The Group Patron Checks In}
The group patron is waiting for the characters, offering assistance as is described in their individual "As a Contact" sections from the Introduction. If this is their first time in this layer, the patron can serve the role—and read the appropriate descriptive text—that is indicated for \textbf{Koh Tam}.

\subsection{"Tiax Always Knows Best!"}
If this is their first time in this layer, \textbf{Tiax} serves as guide instead of \textbf{Koh Tam}, insisting that he has far more insight into this part of the Nine Hells. Unlike \textbf{Koh Tam}, he'll try to accompany the characters to each location they visit within the layer.

\subsection{Tiax's Shady Dealings}
As the characters leave their first area within this layer (or at another opportune time), they notice \textbf{Tiax} trying to hide (poorly) as he spies on them. If confronted, he declares, "\textbf{Tiax} the Mighty was just making certain that his minions were behaving. Now get back to work!" If this event is rolled a second time, run the "Do Not Question \textbf{Tiax} the Mighty!" event instead.

\subsection{"Do Not Question Tiax the Mighty!"}
\textbf{Tiax} is temperamental. Roll a d10 (+ the number representing their current layer of the Nine Hells) every time the characters ask \textbf{Tiax} for advice or enter a new location. On any result lower than a 10, he is helpful. If a 10 or higher is rolled he becomes angry and stomps back to the ship. Each time he becomes angry in this way, add 1 to the next time you make this d10 roll.

\section{No Turning Back}
Having met \textbf{Tiax} and boarded \textbf{Koh Tam}'s barge the characters are ready for their adventure to begin. Proceed to "The Bronze Citadel" section in chapter 3.

\chapter{The Nine Hells}
\DndDropCapLine{T}{}he characters start their journey at the Bronze Citadel. From there they can travel down the River Styx on \textbf{Koh Tam}'s barge. It is important to read through both the introduction and this chapter before running sessions that take place in the Nine Hells.


\section{Character Objectives}
Each layer of the Nine Hells has an adventure location that may contain one of the souls that the characters are searching for. In addition, there is the group patron's objective. \textbf{Koh Tam} is familiar with each layer along with the adventure locations of each. However, he won't speak about a layer or its adventure location until they arrive. \textbf{Koh Tam} responds to any questions on upcoming layers by saying, "All in due time." The characters might wish to explore an adventure location even if it doesn't contain one of their objectives. You should allow them to do this as they can still stumble across the objectives of other group patrons.

The Objective Locations table summarizes where to find the souls the characters seek as well as their group patron's objectives (appendix D contains further details about the items the Deathstalkers seek.)

\section{River Styx Encounters}
Each day that the characters travel down the Styx there is a 50 percent chance that they might encounter something along its shores. If they do happen to encounter something, roll a d6 and add a number equal to the layer of the Nine Hells they're on. For example, if they're on Phlegethos, the fourth layer, then you roll 1d6 + 4. Refer to the River Styx Encounters in the Nine Hells table to determine the encounter that the characters come across.

\begin{DndTable}[header=Objective Locations]{Xcccc}

Location & Lost Soul Phylacteries & Hellriders & Conclave & Deathstalkers\\
The War-Slough, Avernus (1st layer) & Brother, The Queen & — & — & —\\
The Agora of Floating Knives, Dis (2nd layer) & Bounty, The Chosen One, Patricide & — & \textit{Cloak of Invisibility}, Instrument of the Bards, Ollamh Harp, Horn of Valhalla, Iron & \textit{Wrought-Iron Tower}\\
The Ineffable Trove, Minauros (3rd layer) & Father, The Infinite Treasure & — & Belt of Storm Giant Strength & \textit{Accounting and Valuation of All Things}\\
The Elemental Preserve, Phlegethos (4th layer) & Sister, The Furnace & \textbf{Anagwendol} & — & \textit{Amulet of the Inferno}, \textit{Ranseur of Torture}\\
The Chasm of Found Things, Stygia (5th layer) & Business Partner, The Oathbreaker & — & \textit{Holy Avenger}, Manual of Iron Golems, \textit{Rod of Resurrection} & \textit{True-Ice Shards}\\
The Sign of the Hag's Arms, Malbolge (6th layer) & The Heartless Master, Mentor, Mother & — & — & \textit{Scourge of Shadow}\\
The Eye Market, Maladomini (7th layer) & Spouse/True Love, The Great Con & \textbf{Jenevere} & \textit{Staff of the Woodlands} & —\\
The Sorrow Mine, Cania (8th layer) & Student, The Merciless & — & \textit{Staff of the Magi} & —\\
The Oasis of the Lethe, Nessus (9th layer) & — & — & — & —\\
Infernal Warship & — & \textbf{Barachiel} & — & —\\
\textbf{Baalzebul}'s Infernal Submersible & — & — & The Unmaker & —\\
\end{DndTable}

If the characters have already experienced that encounter, then roll on the random encounter table for that layer of the Nine Hells.

\begin{DndTable}[header=River Styx Encounters in the Nine Hells]{cX}

d6 + Layer & Encounter\\
2 & Ride of the Demon Lord\\
3 & Shadow of a Tyrant\\
4 & Horsemen of the Apocalypse\\
5 & Gatekeeper's Quiz\\
6 & Camp of Hedonism\\
7 & Aid from Below\\
8 & Bowl of Suffering\\
9 & Forest of Pain\\
10 & Canyons of Greed\\
11 & Ruin and Amusement\\
12 & Angelic Villa\\
13+ & A Paladin in Hell\\
\end{DndTable}

\subsection{Ride of the Demon Lord}
\begin{DndReadAloud}{}
A sprawling makeshift camp sits atop a low plateau. At the center of the camp is a massive tent stitched from the hides of a variety of devils. Above it, demons fly about calling to each other in their awful speech. Dozens more scream and gibber at each other from barges that float in the River Styx ahead of you. Three such barges have already begun rowing your way but a large, winged demon lands on your ship well before they can reach you.

\end{DndReadAloud}

The camp on the plateau is the war compound of the demon lord \textbf{Baphomet} (see Monsters of the Multiverse). He has penetrated deeply into the Nine Hells and intends to throw his host of demons at one of the devils' great flying strongholds to slaughter as many of his enemies as possible. The demons know they will die in the attack, but since they're not native to the Nine Hells it is of no consequence. Their bodies will simply reform in the Abyss. The devils that they kill will die forever.

The demon that lands on \textbf{Koh Tam}'s ship is a nalfeshnee in command of the demon host's barges. It demands that the characters accompany it to speak with its lord, Baphomet, who is currently within the camp on the plateau. If the characters refuse, the demon attacks and the \textbf{nalfeshnee} is assisted by three fly-like chasmes. If these demons are defeated, the characters can bypass the interdiction and proceed down the River Styx, skipping the remainder of this encounter.

On the other hand, if the characters agree to meet Baphomet, they're escorted to shore and from there, brought into the camp. Baphomet has also recruited several mortals to his cause. The humans will likely end up as lemures when killed, so their participation is particularly unfathomable. Before meeting the demon lord, the characters may speak with some of the humans in the camp:


\begin{itemize}
\item A pair of warriors have come to the Nine Hells on a mission like that of the characters. They're monks from the far-off land of Shou Lung and the souls of their mother and father were taken by \textbf{Asmodeus}.
\item A half-dozen soldiers calling themselves Knights of Solamnia have come from a world called Krynn. They call Tiamat by a different name—Takhisis. They're just trying to survive and have joined Baphomet because he is an enemy of Tiamat.
\item Three Red Wizards from the land of Thay. They tried to summon and control Baphomet and failed. Now they must help in his war against the devils of the Nine Hells.
\end{itemize}

Once inside the compound the characters are brought to see Baphomet who resides within the huge devil-skin tent, the inner supports of which appear to be the bones of a mighty beast.

\begin{DndReadAloud}{}
The commander of this demon army is a great, twentyfoot-tall, black-furred minotaur with six iron horns. At either side of him stand minotaur-like demons who are just as massive as their lord. At the back of the tent sits a bone cage within which over a dozen humans huddle in terror.

\end{DndReadAloud}

Unless the characters are openly hostile, Baphomet, and the two \textbf{goristro} with him, won't attack. His enemies are the devils of the Nine Hells, not the characters. He would rather share his plan. He tells the characters that it is time to throw down the chains of oppression that law and civilization bring. He wants them to help him destroy one of the flying citadels that archdevils often use to fly above the hellscape of their home. Baphomet offhandedly explains that if the characters turn down his offer, they'll join the meat that waits to be consumed after victory is achieved and points to the prisoners in the bone cage. But he insists it is their choice.

Baphomet desires to draw the citadel to this location whereupon he plans to trigger an overly complicated assault, the details of which the characters don't need to know. He simply needs them to locate the nearest citadel and attract its attention, so it diverts towards the camp. As a reward, Baphomet offers a chest full of 5,000 gp once the characters return with a citadel on its way. If the characters negotiate for the human prisoners' freedom, Baphomet agrees to free all of them, though withdraws the offer of the chest. To win both the prisoners' freedom and the gold the characters must succeed on a DC 19 Charisma (Persuasion) check.

Baphomet points across the land, away from the River Styx, and suggests they begin their search in that direction. At this point the characters are free to leave, easily able to double around and return to \textbf{Koh Tam}'s barge and certain escape, though that would leave the captive humans at the mercy of the demons. If they've any interest in saving those unfortunates, the characters need to locate the citadel for Baphomet.

\subsubsection{Finding the Citadel}
Baphomet assembled his camp at this location, knowing it was near a devil patrol route, so locating a citadel won't be difficult. After 1d4 hours the characters see a citadel on the horizon. Any activity out of the ordinary attracts the attention of the citadel's crew (a spell or other visually impressive distraction suffices). Once it begins moving towards them it is easy for the characters to keep ahead of the citadel and reach the camp and Baphomet's tent before it arrives. Baphomet is ecstatic at their success, though distracted at the impending battle. He gestures to their reward and then exits his tent. Through the open flap, the characters watch the citadel's arrival.

\begin{DndReadAloud}{}
The gigantic basalt citadel, shaped like a sword's blade, emerges from the clouds above the River Styx. Screaming demons madly scurry onto the bone platforms laying around the fort. Chasmes add to the cacophony as groups of those fly-like demons grasp the platforms, dragging each into the sky.

\end{DndReadAloud}

There are several choices the characters can make at this point but attacking the citadel is a doomed endeavor and the party should be encouraged to avoid joining the demons.

\subsubsection{Escape and Rescue}
As Baphomet departs the characters should try and make an escape with their reward. If they didn't win the captives' freedom earlier, they must now defeat one \textbf{goristro} guard to free the captives. The remainder of the demonic host doesn't notice these hostilities, busy as they are with the assault. Once on the barge, the characters witness the assault's failure from a distance.

\begin{DndReadAloud}{}
As \textbf{Koh Tam}'s barge drifts down the River Styx, you watch the ill-fated attack on the flying citadel fail as the great swarms of chasmes are torn asunder by various winged devils, their precarious platforms tumbling, each collapse sending dozens of demons and their human allies to their deaths. Silhouetted against the skyline stands Baphomet, the great demon merely shrugging at the failure before walking back towards his camp.

\end{DndReadAloud}

\subsubsection{Attacking the Citadel}
If the characters want to join in the assault on the citadel, it is clear that the demons shall lose this battle. It should also be clear that the characters can simply leave the battle at any time—neither demon nor devil notice. The devilish host includes ten shredwings fighting the chasmes and eight affliction devils guarding each citadel entrance (see appendix B for both). Inside a squad of sixteen bearded devils emerge to further challenge intruders. If somehow despite all these obstacles the characters reach the bridge, they find it crewed by six horned devils, two war devils (see appendix B), and the archdevil appropriate to this layer of the Nine Hells.

\subsection{Shadow of a Tyrant}
\begin{DndReadAloud}{}
Swarms of thousands of bloated flies fill the air with a droning sound. These flies feed on the manure from the herd of rothé that stand dumbly about upon the shore. Jagged rocks and deep pits can be seen a little further inland. Most of the pits vanish into inky darkness, though something metallic glitters within one of the shadowed maws.

\end{DndReadAloud}

This region of the Nine Hells is infested with ayperobos (see appendix B). These tiny devils swarm about by the thousands, disguising themselves within the swarms of black flies. They've even burrowed into the flesh of several of the \textbf{rothé} (see the accompanying stat block) and now control these beasts like puppets. The devious devils have scattered the armor, weapons, and treasure of previous victims near the entrance of one of the caves. This way, it can easily be seen from boats passing by. However, it isn't the ayperobos that poses the greatest danger here.

The first clue to the dangers of this area are the freshly stripped corpses of several rothé, recently eaten by the ayperobos. All that remains of these rothé are their intact skeletons, flecked with blood and flesh. A successful DC 15 Wisdom (Survival) check reveals that these corpses are recent and that whatever killed them stripped all of the meat from their bones in a matter of moments.

The ayperobos wait until their prey has come ashore before attacking—unless it looks like a passing ship doesn't intend to stop. They will then descend upon their victims en masse. A swarm of them will push against the mast, sails and even hull of the Styx barge to drive it against the shore. A successful DC 15 Strength check is required by the navigator to keep a boat on course. The ayperobos aren't interested in eating the flesh of their victims; instead, they serve a greater threat. A \textbf{tyrant shadow} (see appendix B) lives within the caves that dot the shore and the ayperobos lure in hapless mortals and devils to be consumed by it. If any of the ayperobos swarms take control of a creature, they force it to enter the nearest cave mouth.

The tyrant shadow is a thing of concentrated hatred, fear, and loathing. It has taken the appearance of a \textbf{giant spider} made from writhing shadows. As soon as it has an opportunity it snatches a single victim and drags him or her deep into its cave system. From there the tyrant shadow takes its time to devour its prey, hoping that the screams attract a rescue attempt.

This particular tyrant shadow was manifested by the archdevil known as \textbf{Zagum} of the Triad. Treasure. Among the scattered belongings are some valuables. Select three uncommon magic items from the \textit{Dungeon Master's Guide}.

\subsection{Horsemen of the Apocalypse}
\begin{DndReadAloud}{}
Before you, four towering beasts loom, resembling massive dinosaurs covered in heavy armor plating. Seated atop each creature is a devil, their eyes burning with malevolent intelligence.

\end{DndReadAloud}

These are the Horsemen of the Apocalypse, fallen rulers who climbed the ranks of the Nine Hells and now serve as harbingers of the end times. While they await their chance to unleash destruction upon their former worlds, they spend their time bickering among themselves, each vying for dominance over the others. Each rider is astride a formidable steed.

The horsemen won't immediately attack but instead try to engage the characters in their endless bickering, each trying to prove their own superiority over the others. Like it or not, they'll force the characters to choose which of their philosophies is best. Whoever they choose is overjoyed and leaves the fate of the characters in the hands of his or her companions. The remaining three horsemen unleash their full fury, using all their powers and abilities to destroy the characters.

\subsubsection{King Molvar of the Rhokor}
This horned devil rides an ankylosaurus. He formerly governed the kingdom of Keoland on Oerth, disseminating false information regarding the elimination of the so-called "detrimental elements of the society."

\DndQuote%
{}
{''Power is not a means; it is an end.''}
{Neldor the Black}
\DndQuote%
{}
{''History is written by the victors.''}
{Dorgoth the Wise}
\subsubsection{Overlord Zelthor}
This \textbf{bone devil} rides a \textbf{shredwing} and was once an emperor of the Great Kingdom of Aerdy on Oerth. He speaks in the third person and believes in the power of the individual over the collective.

\DndQuote%
{}
{''The only morality is what is good for you.''}
{Thelgor the Black}
\DndQuote%
{}
{''The strongest have the right to rule.''}
{Molthor the Conqueror}
\subsubsection{Overqueen Yarlis}
An \textbf{erinyes} who rides a \textbf{triceratops}, Yarlis was a queen of the Great Kingdom of Aerdy on Oerth. She's notorious for taking credit for other people's quotes, much to the annoyance of King Molvar and Overlord Zelthor. She believes in the power of the individual to change the world.

\DndQuote%
{}
{''The only true purpose in life is to achieve greatness.''}
{Tharlos the Mighty}
\DndQuote%
{}
{''The only limits are the ones we set for ourselves.''}
{Keldor the Wise}
\subsubsection{Yamun Khazad}
This pain devil rides an intimidating \textbf{miasmorne} (see appendix B for both) and comes from the world of Faerûn, where he was once a warlord that united the tribes of the Endless Steppe. He believes in the power of victory and domination over others.

\DndQuote%
{}
{''All problems can be solved through conquest.''}
{Korgoth the Mighty}
\DndQuote%
{}
{''Those who are weak deserve to be conquered.''}
{Thorgar the Conqueror}
\subsection{Gatekeeper's Quiz}
\begin{DndReadAloud}{}
The timbers of your ship groan as it suddenly lurches to a stop. It seems to have run aground in the middle of the Styx.

\end{DndReadAloud}

A powerful \textbf{amnizu} (see Monsters of the Multiverse), named Minos, controls this section of the River Styx. He is in command of a \textbf{Styx dragon} (see appendix B) that he uses to stop ships passing through. Minos considers himself an arbiter and judge of the Nine Hells. He enjoys judging the souls of the damned. Though boats frequently pass him by, they're rarely crewed with mortals, and thus his interest is piqued. Minos' other talent is with riddles, something he loves to practice with devils and mortals alike.

Minos rises from the Styx once the ship has been brought to a stop by his Styx dragon. He asks for the names of the crew members. If any on board shares their name, Minos acts as if he has heard the name before and plays upon the pride of the individual. He tries to convince one of the characters to play a game of riddles with him. If they refuse, then he explains that their ship is in the grip of his Styx dragon and that he will only allow them passage if they play his game.

The game consists of three rounds. In each round, Minos will pose a riddle and then the character must pose a riddle. Feel free to replace the three riddles below with different ones of your choosing. The players can either come up with a riddle on their own (which you try to answer) or they can succeed on a DC 16 Intelligence check to craft a riddle Minos can't solve. The characters are victorious if Minos can't answer one of their riddles or if they're able to answer all three of the riddles posed by Minos.


\begin{itemize}
\item I create nothing from everything. I provide passage where there is none. I travel the layers peddling my wares, but no payment is required. What am I?

\begin{description}
\item[Answer:]
River Styx
\end{description}

\item I am something of ultimate value, yet weightless. I am a ledger of deeds, yet nothing is written. I can visit other planes, yet the corporeal remains. What am I?

\begin{description}
\item[Answer:]
Soul
\end{description}

\item In consuming everything, it gives life. Its remains bring about rebirth. To embrace it is to die, yet it readily travels with almost any group. What is it?

\begin{description}
\item[Answer:]
Fire
\end{description}

\item Every creature follows their own version, whether written or thought. Without it, there would be no Blood War. Its interpretation and presence has driven conflict since the beginning of time. What is it?

\begin{description}
\item[Answer:]
Law
\end{description}

\item I'm never far from my sibling, fire, Find comfort in me, liars, For when I make something unfound, Use touch, taste, smell, or sound.

\begin{description}
\item[Answer:]
Darkness
\end{description}

\item To have me is to have everything, To be as powerful as a king, Yet, my endless pursuit brings only emptiness, Having everything brings joy less and less.

\begin{description}
\item[Answer:]
Wealth
\end{description}

\end{itemize}

Winning the game of riddles results in Minos allowing the characters to choose one infernal magic item from his stash at the bottom of the Styx. He has the following items: \textit{Bracers of Asmodeus}, \textit{Canian Fork}, \textit{Demonbone Polearm}, \textit{Infernal Amulet}, \textit{Infernal Plate Armor}, \textit{Stygian Spear}, and the \textit{Sword of Retribution} (these are described in appendix D).

Losing the game of riddles results in Minos demanding 5 Soul Coins to continue down the river. Refusal results in Minos and the Styx dragon attacking. If Minos is attacked or if the game of riddles is refused, then the amnizu and his pet Styx dragon will attack together.

\subsection{Camp of Hedonism}
\begin{DndReadAloud}{}
A large, fortified camp has been built on the shore of the River Styx from huge bones and bleached timber. No one patrols its walls, but the sounds of laughter and joy can be heard from within.

\end{DndReadAloud}

This camp is the home of Gulgara, a powerful corruption devil (see appendix B) known as a sire of corruption, and her many minions. Gulgara created the camp decades ago as a trap for the unwary and a place to hone her powers of corruption. Recently a large group of adventurers were searching for the souls of their loved ones when they came under attack by numerous devils. The adventurers could not believe their luck when they saw what appeared to be an abandoned infernal fort. They took refuge within, not realizing it was the corruption devil's powers that had drawn them there in the first place.

\begin{DndSidebar}[float=!b]{Decadent Corruption}
The most powerful of corruption devils, the sire of corruption, gains a powerful aura that pervades and weakens the will of anyone within 300 feet of the devil. When a creature first enters the aura, they must make a DC 19 Charisma saving throw. On a failed save the creature falls under a curse known as decadent corruption.



While cursed in this way, the character has disadvantage on all saving throws and ability checks. In addition, a character that has the charmed condition while cursed in this way, remembers nothing of what happened to them while charmed, once the charmed condition ends. Finally, a cursed character becomes reckless, automatically failing any Wisdom (Insight) checks they may make. Often this manifests as a dangerous sort of curiosity, which may result in a desire to investigate the camp.



A \textit{Greater Restoration} spell, a \textit{Remove Curse} spell, or other magic that removes curses, can remove the corruption curse and end all its effects. Those that succeed the saving throw or have been cured are immune to the effects of decadent corruption until they finish a long rest. Leaving the corruption aura also ends the curse on the affected creature but does not grant immunity to the curse. A corrupted character won't voluntarily leave the area and if asked to do so, the character makes excuses for why it can't.



\end{DndSidebar}

The adventurers (four mages, three knights, two druids, two spies, and a \textbf{priest}) are oblivious to the many torments they've suffered. A dozen succubi, in the guise of friendly elves, work in the camp, enacting Gulgara's various torments.

Characters approaching the gates to the camp are welcomed by a pair of succubi, claiming to be adventurers. They invite the characters to partake in the festivities within the camp. They open the gates to reveal the grounds of the fort, which are occupied by numerous adventurers. There are also two large tents in the distance.

A succubus points out the activities that the characters might want to participate in:


\begin{itemize}
\item A half dozen adventurers are greedily eating from a vast table filled with delicacies from across the multiverse.
\item A pair of knights are sitting on cushioned seats in front of vast mirrors as their armor is shined by imps. There are several empty seats, and more imps waiting to work.
\item A priest and a couple of spies are engaged in a rowdy game of dice.
\end{itemize}

Each tent is guarded by a pair of succubi that have taken the appearance of human knights and pretend to be adventurers. They allow entry only to those obviously under the influence of Gulgara's decadent corruption ability. If the characters attempt to force their way into a tent, all the camp's occupants attack the characters.

In the first tent, a mage is being tortured by a trio of affliction devils (see appendix B). The devils attack characters not under the influence of the decadent corruption ability. The other tent is where Gulgara can be found. She attacks characters not under the influence of the decadent corruption ability but emerges from the tent only if a fight breaks out between the characters and her servants.

The entire camp is under the influence of Gulgara's decadent corruption ability. The adventurers can be saved but each requires their decadent corruption be removed, as described in the sidebar. If a devil or succubus sees that an adventurer has been saved in this way, or any adventurer attempts to leave the camp, the remainder of the camp goes hostile. Adventurers that have been saved fight to assist the characters and, if they survive, continue to travel the Nine Hells with the characters.

\begin{DndSidebar}[float=!b]{Temptation of Betrayal}
\textbf{\textit{Nothing Can Touch Me.}}



Gulgara wears a magic ring with a shield carved into it and infernal script covering every surface. If a character puts the ring on, it triggers a temptation. The temptations that \textbf{Asmodeus} throws in the way of the characters, are described later in this chapter; consult there for further details about the "Nothing Can Touch Me" temptation.



\end{DndSidebar}

\subsection{Aid from Below}
\begin{DndReadAloud}{}
Ahead, a band of centaurs makes their way across a burnt-out forest. Some of them pull a cart behind them.

\end{DndReadAloud}

This ragtag group of twelve centaurs are dedicated to a good, but misguided, cause. After some discussion, the characters will learn that these centaurs are attempting to alleviate the suffering of some of the souls in the Nine Hells. In particular, they help individuals they deem non-deserving of their suffering. These larvae (see the \textit{Dungeon Master's Guide} for this stat block), lemures and nupperibos are in bone caged carts that some of the centaur band are yoked to. They willingly entered the Nine Hells decades ago at the behest of their chieftain. They originally numbered in the hundreds but have been reduced to just a dozen.

What they don't realize is that their quest is a trick. Their chieftain signed away their souls so that he could be granted great power. He convinced his greatest warriors that they needed to enter the Nine Hells and save the innocent souls that had been unfairly consigned to suffer there. He told them how to recognize such souls, but these were just lies. The signs he told them to look for include:


\begin{itemize}
\item Joyful facial expressions
\item Issuing curses against the gods
\end{itemize}

If the characters are able to ascertain that the centaur's quest is a lie (through magic spells such as \textit{Divination} or \textit{Contact Other Plane}), then they can try to convince the band of their findings. This requires a successful DC 20 Charisma (Persuasion) check.

If the characters are successful at convincing the centaurs of the fact that their quest is a lie, one of the nupperibos begins to giggle, then laugh, then cackle uncontrollably. After his death, the chieftain was consigned to the Nine Hells, as a nupperibo. The centaurs stumbled upon him a few years ago, without even realizing his identity, and he has been with them ever since.

\subsection{Bowl of Suffering}
\begin{DndReadAloud}{}
Piteous screams can be heard from the shoreline. A huge bowl formed from smooth obsidian is surrounded by a half dozen devils that push back creatures trying to escape from it. Enormous elephant like devils stand impassively as guards.

\end{DndReadAloud}

There are six pain devils and two maelephant nomads (for both, see appendix B) around the bowl. Only someone standing at the edge of the bowl can see that it is filled with larvae (see the \textit{Dungeon Master's Guide} for this stat block), lemures, nupperibos (see Monsters of the Multiverse), and a few dretches captured during the Blood War. These dretches have honed their telepathy so that they can communicate in any language. The pain devils push down any of the pitiful creatures that manage to climb up the sides of the bowl.

The maelephants attack anyone who tries to interfere with the pain devils.

The dretches use their telepathy to make promises to any mortal they see, but all their promises are lies. One dretch claims to be an angelic creature trapped under the other infernal creatures. Another dretch claims that there is a suit of magic plate armor at the bottom. Anyone crawling into the bowl is swarmed by the piteous creatures and must make a DC 15 Dexterity saving throw. They take 21 (6d6) slashing damage on a failed save, or half as much damage on a successful one. Magic within the bowl is corrupted so that any spell cast within it fails and triggers a wild magic surge. Determine the surge's effect by rolling on the sorcerer's Wild Magic Surge table in the \textit{Player's Handbook}.

\subsection{Forest of Pain}
\begin{DndReadAloud}{}
A forest grows on both sides of the River Styx. The trees are twisted and without any leaves. Among the trees lumber huge sloth like creatures.

\end{DndReadAloud}

This forest is the home of one of Hell's larger herds of oneirovores (see appendix B). These creatures, also known as dream eaters, are grazing on the branches of the trees. When they break off a branch it bleeds a luminescent sap that the oneirovores eagerly lap up. This sap is made up of the dreams of mortals.

Attacking the herd can quickly turn into a disaster. If any of the oneirovores is hurt, 2d4 hellcats (see appendix B) immediately spring to their defense. An oneirovore that is injured cries mournfully and releases its stored phantasmagoria. Nearby oneirovores also release their phantasmagoria. Things quickly get worse when the herd's shepherd arrives. In 1d4 rounds a pack of a dozen hell hounds will burst on to the scene. A round later the shepherd, a \textbf{war devil} (see appendix B), arrives.

If the guardians of the flock are dealt with, then it isn't too difficult to coax some of the oneirovores on to a boat. A successful DC 15 Wisdom (Animal Handling) check convinces an oneirovore to follow along. However, it is very dangerous for anyone to rest while so close to a dream eater. Someone who finishes a long rest within 10 feet of an oneirovore must make a DC 15 Wisdom saving throw. The DC increases by 5 for each oneirovore beyond the first. A failed save results in disadvantage on Wisdom saving throws. Only a \textit{Wish} spell can remove this effect.

\subsection{Canyons of Greed}
\begin{DndReadAloud}{}
The River Styx cuts its way through a canyon larger than any in the mortal world. A sound like thunder can be heard from above as rocks begin to tumble down the canyon's sides towards your barge.

\end{DndReadAloud}

Four maelephant nomads (see appendix B) stand at the top of the canyon and push boulders toward any ship passing by. These boulders glisten with embedded gold nuggets and when they smash into their target or hit the side of the canyon they break apart, scattering 2d12 gold nuggets along the riverbank. Each throws three boulders (Ranged Weapon Attack: +11 to hit, range 60/240 ft., one target. Hit: 28 (4d10 + 6) bludgeoning damage) before retiring out of sight. A spell such as \textit{Gust of Wind} pushes a ship out of range immediately.

Nearly a dozen prospectors emerge from caves in the sides of the canyon and huts built along the shore. They scramble to collect the gold nuggets, often wading into the river to pick up nuggets that land in the shallows. The memories of these prospectors have long ago been eroded. Now all they care about is finding gold nuggets in the river and hiding whatever they find in the small lair they've built for themselves. Each gold nugget is worth 100 gp, but the prospectors react violently to anyone who tries to take their gold. The prospectors have silently agreed to not touch the treasure hoard of another. Anyone doing so earns the wrath of all the prospectors. A variety of prospectors live here:


\begin{itemize}
\item 3 fire giants
\item 1 dwarf \textbf{gladiator}
\item 1 orc \textbf{veteran}
\item 1 elf \textbf{assassin}
\item 1 \textbf{young red dragon}
\item 2 fomorians
\end{itemize}

Each prospector has 1d100 gold nuggets in their lair.

\subsection{Ruin and Amusement}
\begin{DndReadAloud}{}
The River Styx flows through a stone canal that cuts through the heart of an infernal city. The sounds of laughter and joy can be heard from what looks to be an amusement park!

\end{DndReadAloud}

The River Styx flows into a stone canal that has been built through the heart of one of Hell's many abandoned cities. The city is long dead, except for an amusement park at the city center.

This amusement park is another one of \textbf{Baalzebul}'s bizarre creations and is one of the few places in the Nine Hells where the non-infernal outnumber devils. Travelers from across the planes travel down the Styx to experience the famed rides.

\begin{DndSidebar}[float=!b]{Location}
If this encounter takes place in Malbolge or Cania, then it should do so on the borders of Maladomini.



\end{DndSidebar}

At the main entry, a \textbf{succubus} and an \textbf{incubus} hand out entrance tickets for 10 gp each. However, anyone sneaking in without paying doesn't need to fear reprisal. Once inside, visitors are free to try out any of the fun rides. Once they've tried out one or two, they are approached by a succubus or incubus who tries to persuade them to try out one of the devilish rides.

\subsubsection{Fun Rides}
These are rides that, while frightening, aren't dangerous.

\subparagraph{Ferris Wheel}
This huge wheel allows up to six Medium-sized creatures in each of its cages. Cleverly placed illusions create elaborate scenes, simulating a journey through the Nine Hells. A rider looking out from their cage begins their journey with a majestic view of Avernus and travels in a slow circle through each layer in turn, with the ride ending with a view of the deep canyons of Nessus. In total the ride takes no more than ten minutes but feels much longer than that.

\subparagraph{Infernal Mansion}
This mansion is filled with some of the most gruesome looking demons and devils. The devils are willingly scaring participants, while the demons are... not so willing. A \textbf{chain devil} uses chains to disturb those who walk down an inky black corridor. A pair of hellcats (see appendix B) terrorize people walking through an ancient torture chamber and a \textbf{displacer fiend} (see appendix B) reaches out with its tentacles from a cage just unable to reach participants. Finally, there are the exploding \textbf{manes} demons. They're released from multiple tiny doors and run gibbering at participants before exploding right before they reach them. This means that participants are always covered in demonic ichor. Luckily, this ichor doesn't behave as it does in the rest of the lower planes, but participants aren't told that.

\subparagraph{Roller Coaster}
The second most terrifying roller coaster in the multiverse. The twists, turns, spirals and upside-down tracks dissuade many from trying it out. Anyone riding it must succeed on a half dozen DC 15 Wisdom saving throws to stop from screaming during the ride.

\subsubsection{Devilish Rides}
Each ride costs 250 gp per rider and can be extremely dangerous. However, winning contestants earn powerful magical treasure. Winners of a ride earn a magic item from the \textit{Dungeon Master's Guide}. Winners can win only one prize. Any character under 10th level receives an uncommon magic item. Characters 10th level to 15th level receive a rare magic item (excluding weapons) and characters 16th level or higher receive a very rare magic item (excluding weapons).

\subparagraph{Carousel}
This carousel starts off spinning slowly, but quickly gains speed. Contestants must succeed on a DC 10 Strength (Athletics) check to hold on during the first minute. Every subsequent minute the DC increases by 5 to a maximum of 30. The ride ends when there is only a single rider left or once the DC reaches 30. A rider who fails their check flies out and hits the barrier walls, taking 3 (1d6) bludgeoning damage for every point of DC—to a maximum of 70 (20d6). The barrier needs to be cleaned of gore at the end of every night.

\subparagraph{Doom Coaster}
This is the deadliest and most terrifying roller coaster in the multiverse and carries riders through a series of dangerous features. With each terrifying event, a rider must survive not only the physical danger, but they must succeed on a DC 15 Wisdom saving throw or have the frightened condition for the remainder of the ride. While frightened in this way, a rider also has disadvantage on all saving throws.

The roller coaster begins by ascending a vast stretch of rickety track until it reaches 600 feet. It then travels through the features, visiting each once. After the final feature the roller coaster returns to the relative safety of ground level again. The features are presented in the order that riders encounter them:


\begin{description}
\item[Lady of Pain.]
The roller coaster does a circle around an enormous head with huge blades sticking out of it. Each rider must succeed on a DC 15 Dexterity saving throw or take 35 (10d6) slashing damage.

\item[Maw of the Dragon.]
The roller coaster drops into the fiery maw of an artificial dragon and each rider must succeed on a DC 15 Constitution saving throw or take 28 (8d6) fire damage.

\item[Stygian Glaciers.]
The roller coaster turns upside down so that its riders are immersed in a lake of frigid water filled with razor sharp icy shards, before the roller coaster ascends again. Each rider must succeed on a DC 15 Constitution saving throw or take 17 (5d6) cold damage and 7 (2d6) slashing damage.

\item[Storm King's Thunder.]
The roller coaster drops into the arms of a 100-foot-tall storm giant automaton. As the giant places the coaster on the new set of tracks, everything is electrified, and each rider must succeed on a DC 15 Dexterity saving throw or take 35 (10d6) lightning damage.

\item[Great Green Devil.]
The roller coaster rushes through the open mouth of a huge green devil face and each rider must succeed on a DC 15 Constitution saving throw or take 42 (12d6) necrotic damage. Any rider reduced to 0 hit points is disintegrated.

\end{description}

\subparagraph{Bumper Cars}
A variety of infernal war machines compete here in a massive arena. The drivers of the machines must destroy the other machines by smashing into them. Riders may select either a Large machine or a Huge machine (each machine can only have one occupant). These vehicles have fierce names, such as "Devil's Ride," "Tormentor," and "Scavenger". The character's goal is for their machine to survive 5 rounds—each rider that manages this feat, wins. Surviving each round requires a character to make a Dexterity saving throw to avoid a collision with another machine. The DC if the character is piloting a Large machine is 15 but becomes 20 for those piloting a Huge vehicle. If the character fails the saving throw, their vehicle is destroyed if it is Large. A Huge machine can survive one impact, but a second impact destroys it. When a machine is destroyed, the rider takes 17 (5d6) bludgeoning damage. If a vehicle-less character starts their turn in the arena, they must make a DC 16 Dexterity saving throw. If they fail, they take 35 (10d6) bludgeoning damage as they're run down by an infernal machine.

\subsection{Angelic Villa}
A small group of angels live in a once beautiful palace, overgrown with the filth of the Nine Hells. Flies and wasps surround the palace, and worms writhe in the mud around it. These angels were companions of Triel, the fallen angel that became the archdevil \textbf{Baalzebul}. They aren't evil, but they're no longer good either. They haven't committed to either the hierarchy of the Nine Hells or returned to their original home in the Seven Heavens. Their moral relativity has made them into sad figures with no purpose or hope.

The angel Samael flies above the Styx and attempts to contact any good characters it notices below. The angel uses telepathy to provide directions to a nearby palace where characters may rest and recover.

\begin{DndReadAloud}{}
A palace is half-sunk in the mud and a great swarm of flies and wasps hovers over it like a nightmarish cloud. You look up at the façade of the building and see the remnants of beauty long gone. Intricate carvings in marble are covered in filth and crawling things. Mighty stone walls lie half crumbled away, revealing the rotting wood supports jutting out like shattered bones. There were tall windows once. Now there remain only gaping holes. Slits of impenetrable black, glaring down at you like soulless eyes.

\end{DndReadAloud}

When the characters enter the palace, they find angels within. There are several devas led by a planetar named Uriel. These angels have lived together in this region of the Nine Hells for millennia. There were once dozens of them, but over time their number has been reduced as some have died in battle, others have fallen to complete corruption, and a few have redeemed themselves and returned to the Seven Heavens. Those that remain are:


\begin{description}
\item[Uriel.]
He was a patron of the arts.

\item[Arariel.]
He controlled the waters of the earth.

\item[Chamuel.]
They represented serenity and devotion.

\item[Gadreel.]
He was known as the watcher.

\item[Muriel.]
She was a patron of travelers.

\item[Phanuel.]
She offered hope and repentance.

\item[Samael.]
He fetched the souls of the dead.

\end{description}

Uriel wants the characters to join them for dinner. He wishes to hear tales of the characters' exploits. In return, he allows the characters to partake from his wine cellar, one of the most comprehensive in the multiverse.

If the characters turn him down, he is disappointed, but lets them know that they can return at any time to take him up on his offer. Any act of aggression by the characters results in the angels attacking with a level of coordination that befits a group that has been together for millennia.

\begin{DndReadAloud}{}
The angel called Chamuel has brought a rose to dinner. They twirl the long-stemmed flower between their fingers and smile whenever they catch you looking. Something so delicate seems terribly out of place amidst the rot and horror of the Nine Hells and as Chamuel lifts the rose to brush against their lips, you notice how soft and red its petals are. To the angel Chamuel's left, at the head of the table, sits Uriel, your host, and on the other side of him is the angel called Phanuel. She sits stiffly in her chair and though her face is still and placid, her piercing eyes are full of storms. Uriel makes a gesture and the other angels file in, bringing with them plates of food and pitchers of drink. When everything has been laid out, they take their own seats. Uriel makes another gesture and the tall candles that line the long table all light up. "Well, then," Uriel says with a smile. "Shall we begin?"

\end{DndReadAloud}

The dinner is fraught with potential danger. During each of the three courses, an angel asks a challenging question which requires a successful ability check to answer. Failing a check results in that angel challenging that character to a duel. If the character refuses the duel, then Uriel expels the character from the estate. If the character undertakes the duel and the angel reduces the character to 0 hit points, the angel chooses to incapacitate the character instead of landing a killing blow. If the character doesn't do the same if they win, then Uriel expels all of the characters from the estate.

In the first course, Chamuel selects the character with the highest Charisma score and questions their serenity or devotion. The character must succeed on a DC 16 Charisma (Persuasion) check to avoid the duel. In the second course, Samael demands that the frailest character (the one with the lowest Constitution score) justifies how they've survived so far. The character must succeed on a DC 15 Charisma (Persuasion) check, made with disadvantage if the character is injured. In the final course, Phanuel, who has become exhausted from sustaining all her companions, now asks a character to step forward to convince her that hope still remains for these abandoned angels of the Nine Hells. That character must succeed on a DC 18 Charisma (Persuasion) check.

If the characters are not expelled, when they prepare to leave, Uriel allows each into his wine cellar. Some bottles of wine have magical properties that might be useful on their journeys. Each character can make an Intelligence (Investigation) check to select a good vintage. Any result above 15 allows them to roll once on the Magical Wines table. Above 20 and they roll twice, allowing them to select which of the two bottles they prefer to take.

\begin{DndTable}[header=Magical Wines]{cXX}

d8 & Name of Wine & Powers\\
1 & Château de Tethyr & \textit{Potion of Flying}\\
2 & Casa Athkatla & Potion of Cloud Giant Strength\\
3 & Obarskyr Estate & \textit{Potion of Invisibility}\\
4 & Wolf Blasé & \textit{Potion of Longevity}\\
5 & Rauxes-Nyrond & \textit{Potion of Speed}\\
6 & Menage a Trois & Potion of Storm Giant Strength\\
7 & Blackstaff & \textit{Potion of Supreme Healing}\\
8 & Felix Solar & \textit{Potion of Vitality}\\
\end{DndTable}

\subsection{A Paladin in Hell}
\begin{DndReadAloud}{}
A narrow path climbs the cliff wall beside the River Styx. On the path an armored warrior battles dozens of devils. His sword and armor are coated in the black blood of his enemies. He calls out for help.

\end{DndReadAloud}

If \textbf{Koh Tam} is in command of the characters' ship, then he steers it to the riverbank and demands that the characters help the warrior. The warrior is a paladin serving Kelemvor named Sir Calenhad Strongheart (lawful neutral \textbf{death knight} that is a Humanoid and doesn't have the Marshal Undead trait or Hellfire Orb action).

Calenhad battles against a dozen bearded devils, a pair of horned devils and nine bone devils. These creatures are led by a powerful \textbf{greater tyrant shadow} (see appendix B). This tyrant shadow was spawned from the archdevil \textbf{Abigor} millennia ago.

Calenhad has been tasked by the archdevil \textbf{Abigor} with killing the tyrant shadow, for the creature has been subtly undermining his machinations. \textbf{Abigor} convinced the paladin that the tyrant shadow was responsible for the death of his family and their souls will be consigned to the Nine Hells unless he kills the creature (all lies). To help Calenhad in his quest, \textbf{Abigor} granted him Infernal Plate Armor and a Sword of Retribution (see appendix D for details about both items.) These items have been slowly twisting the paladin.

\textbf{Koh Tam} asks the characters to subdue the paladin so he can remove the cursed items and heal him of infernal corruption. If they succeed in this, Calenhad stays below decks slowly recovering. If you want, he may come to the rescue of the characters at a key moment in the story.

\section{Asmodeus and the Nine Sins}
It is important to remember that \textbf{Asmodeus} has lured the characters (and many others) into the Nine Hells to be corrupted. There are many instances in the adventure where \textbf{Asmodeus} tempts the characters. It is entirely likely that characters won't give in to any of the temptations put before them, but it is important that they feel at least a little tempted during the adventure.

\subsection{Corruption}
The temptations have real effect on the game, by modifying the amount of corruption for each character. The more corrupt a character, the more likely it is that they succumb to \textbf{Asmodeus}'s manipulations. Keep track of the individual characters corruption scores on the "Corruption Tracker" found in appendix E. All characters start with a Corruption of 0, but are modified by the following:


\begin{itemize}
\item Evil characters start with 1 corruption point.
\item Characters searching for their own damned soul gain an additional corruption point.
\item If the party aligns with the Deathstalker patron, each character gains an additional corruption point.
\end{itemize}

When a character gives in to one of the temptations an additional 2 corruption points are awarded to them. At the final stage of the adventure, you can tally the scores for each character. Use the Corruption Score table found in appendix E to see if a character has been ensnared by \textbf{Asmodeus}.

\subsection{Types of Player Goals}
Any temptation must fit the goals of a player as it relates to their character. All temptations are grouped into one of the following categories:


\begin{description}
\item[Power.]
These temptations provide magic items, spells, special abilities, and minions.

\item[Knowledge.]
Discovering the location of fabled sites, the secrets of creation, the weakness of enemies, and the true thoughts of others are possible with these temptations.

\item[Privation.]
Temptations of this kind involve relieving privation by means of restoring hit points, lost abilities, and providing protection from death.

\item[Roleplaying.]
These temptations help achieve character-specific but non-power-related goals—the death of enemies, the benefit of friends, family, community or dependents in need, the altering of a political status quo, the location of some lost icon of cultural or family significance. Titles and status. These rely on the character's backstory having appropriate material, and the player being sufficiently invested in it.

\end{description}

\subsection{The Temptations}
Throughout the story \textbf{Asmodeus} tries to tempt the characters. Each is optional, so feel free to skip a temptation, especially if another one has recently been triggered, or move them to another part of the Nine Hells. However, make sure that each character is tempted at least a few times throughout the adventure.

The benefits and drawbacks of temptations are described in detail on the following pages. Most of these temptations have been placed in specific locations within the various layers of the Nine Hells and the details of discovering them are explained in an appropriate sidebar in that layer. Not all the temptations have been directly placed, however. In some locations a prompt appears to select one from the Random Temptations table instead. When this occurs, roll on the table, or select the temptation most appropriately tantalizing to one of the characters.

\begin{DndTable}[header=Random Temptations]{cXX}

d12 & Sin & Result\\
1 & Anger & The Cultists (Knowledge)\\
2 & Anger & Vengeance is Yours (Roleplaying)\\
3 & Murder & A View to a Kill (Roleplaying)\\
4 & Greed & Box of Treats (Power)\\
5 & Greed & A Head for Knowledge (Knowledge)\\
6 & Jealousy & Anything You Can Do (Knowledge)\\
7 & Jealousy & Credit Where Credit is Due (Roleplaying)\\
8 & Betrayal & What is Rightfully Mine (Roleplaying)\\
9 & Oppression & The Angry Djinni (Power)\\
10 & Harm & Non-Lethal Weapon (Power)\\
11 & Harm & A Torturer in Hell (Roleplaying)\\
12 & Pride & Taking Pride in Your Work (Roleplaying)\\
\end{DndTable}

\section{Let Me Tempt You}
Each of \textbf{Asmodeus}'s temptations is designed to use one of the primary sins to ensnare its target. Whether one lusts for power, knowledge or riches \textbf{Asmodeus} has designed an appropriate temptation. The temptations are listed under their corresponding sin.

\subsection{Anger}
\DndQuote%
{}
{''Anger is forged in the fires of the Nine Hells.''}
{\textbf{Bel}}
\subsubsection{The Sentient Artifact (Power)}
The characters find an ancient sentient artifact of the Blood War. It promises to give one of them great power in the form of rage. Once it has granted the power to one of the characters, the sentience of the artifact winks out and it becomes a nonmagical lump of metal. Whoever was granted the power gains the following feature. When taking damage, you may use a reaction to enter an infernal fury. Each time you activate this fury, you select one of the following features:


\begin{itemize}
\item When making a melee weapon attack using Strength, you gain a +4 bonus to the damage roll. Also, when you take damage, you may use your reaction to make one melee weapon attack against the creature that damaged you if they're in reach.
\item When making a spell attack, you gain a +4 bonus to the damage roll. Also, when you take damage, you may use your reaction to cast a cantrip, targeting the creature that damaged you.
\end{itemize}

Your fury lasts for 1 minute. It ends early if you have the unconscious condition or if your turn ends and you haven't attacked a hostile creature since your last turn or taken damage since then. You gain one level of exhaustion when the fury ends unless you damaged a non-evil creature while using this fury. You can't activate infernal fury again until after you finish a long rest.

\subsubsection{The Cultists (Knowledge)}
A group of ten fanatics (lawful neutral cult fanatics) marches down a road chanting a prayer. The prayer describes how they gain esteem in the eyes of their god through the righteous and holy action of destroying false knowledge. Their beliefs should conflict severely with the beliefs of at least one of the characters. It should be clear that if allowed to proceed they intend to travel to the Material Plane and destroy a magical academy, a church, or a druid grove, meaningful to a character. There is no reasoning with them. Stopping their threat now, with no witnesses, is a temptation that is hard to ignore.

\subsubsection{Gauntlets of Rage (Privation)}
A devil offers to sell one of the characters a powerful magic item known as the Gauntlets of Rage (see appendix D).

\subsubsection{Vengeance Is Yours (Roleplaying)}
A horned devil arrives on \textbf{Koh Tam}'s barge. He offers, for the incredibly small price of 100 gp, vengeance on a dead enemy whose soul is in the Nine Hells. The devil insists there's no harm in it. After all, the soul is already being tortured anyway. If a character accepts, then the horned devil has the soul teleported to the barge away from the soul's current place of torment. The soul appears with all the statistics they had while alive, except they're restrained with infernal iron bonds. The character may inflict any punishment upon them that they desire but if slain, the soul vanishes, presumably returning to their former punishments.

\subsection{Betrayal}
\DndQuote%
{}
{''You can only understand good and evil if you have betrayed the former for the latter.''}
{\textbf{Levistus}}
\subsubsection{Mixed Blessings (Power)}
A \textbf{pit fiend} offers to grant one of the characters a blessing to help in their quest. The fiend claims to want to embarrass \textbf{Asmodeus}, but the character must accept the blessing without knowing what it does.

\subparagraph{Blessing of Power}
Your attacks ignore the damage resistance of any creature you damage. Unfortunately, your parents, siblings, and children become cursed. While cursed in this way, they each have the poisoned condition until their curse is ended by a \textit{Remove Curse} spell or similar magic, or until the blessing is removed from you with a \textit{Wish} spell.

\subsubsection{If I Could Read Your Mind (Knowledge)}
A character is given the opportunity to be privy to the thoughts of a mentor, family member, liege or superior, sold to them on the basis that they will always be able to act in that character's best interests. The ability occurs at will but also sporadically without asking, and the information they receive verges into the embarrassing, humiliating and awkward, and their betrayal is one of invasion of privacy.

\subsubsection{Nothing Can Touch Me (Privation)}
The characters find a ring on the finger of a devil that they've slain. This magic item is the Ring of Treachery (see appendix D).

\subsubsection{What Is Rightfully Mine (Roleplaying)}
A character whose background is tied in with an organization, liege lord, or similar, is introduced to a devilish counterfeiter with 'evidence' proving their right to a superior title—throne, inheritance, office land or item—maybe something they've always believed they should have had. The 'evidence' has a blank where their name could go, but what's the harm in taking it? After all, if they don't, what if the seller offers the same to someone else? Once they have the evidence, they find they can't destroy or get rid of it—the item, increasingly incriminating, keeps appearing on their person or even nearby where anyone might stumble over it.

\subsection{Deceit}
\DndQuote%
{}
{''Deceit allows those with less power and excellence to pretend otherwise.''}
{\textbf{Baalzebul}}
\subsubsection{The Liar (Power)}
A \textbf{succubus} offers to give one of the characters the power of deception in exchange for a minor magic item or a single Soul Coin.

\subparagraph{Power of Deception}
You gain expertise in the Charisma (Deception) skill. However, anything said in good faith and sincerity comes out sounding false and unconvincing. To persuade people, you must couch everything as a chain of falsehoods. When speaking the truth, you have disadvantage on all Charisma (Persuasion) checks.

\subsubsection{As Good an Answer (Knowledge)}
One of the characters is given the chance to create the answer to a question that has vexed academics for an age. The truth of the matter is lost beyond even the most powerful divination or is known only to extremely close-mouthed gods. However, an \textbf{amnizu} (see Monsters of the Multiverse)—or another powerful devil—suggests that, given nobody knows, the character has the opportunity to use their own preferred answer, which the devil ensures is spread through the mortal planes and accepted as the definitive truth. Only the character will ever know that it's a lie...

\subsubsection{Life After Death (Privation)}
The characters find the Amulet of Duplicity (see appendix D) around the neck of a defeated enemy.

\subsubsection{The Lookalike (Roleplaying)}
A corruption devil (see appendix B) offers one of the characters the services of a \textbf{doppelganger}. The doppelganger will faithfully impersonate NPCs of the character's acquaintance and act in their stead, perhaps to help the character's family, community, or allies. A corrupt local official might suddenly step down from their post. A local judge or king might decide in the character's favor. An enemy warlord will suddenly lead their horde in a different direction. What good might not be accomplished with a little duplicity?

\subsection{Greed}
\DndQuote%
{}
{''Greed makes you free, and accountable to no one.''}
{\textbf{Mammon}}
\subsubsection{Box of Treats (Power)}
The characters find a transparent box that contains several powerful items. Select one rare or very rare item from the \textit{Dungeon Master's Guide} for every two characters, rounded down. The box is invulnerable to all damage. It has no lock or lid. Inscribed in golden letters on top are instructions written in Infernal for opening it. Only with the blood of a good creature can the box be opened and only by the person who killed this creature. Once opened, the box vanishes, leaving only the items.

\subsubsection{A Head for Knowledge (Knowledge)}
The characters find a withered devil skull known as the Skull of Selfish Knowledge (see appendix D).

\subsubsection{Save It for a Rainy Day (Privation)}
The characters find the \textit{Vial of Greed} (see appendix D) in the hoard of one of their defeated enemies.

\subsubsection{The Greedy Art Lover (Roleplaying)}
The characters find a magic ring known as the \textit{Ring of Collecting} (see appendix D) in the ice fields of Stygia.

\subsection{Harm}
\DndQuote%
{}
{''What further harm could I possibly bring to the wretchedness of mortals.''}
{\textbf{Mephistopheles}}
\subsubsection{Non-Lethal Weapon (Power)}
The characters discover a melee weapon of a kind that one of them is proficient with. This is the \textit{Weapon of Agonizing Paralysis} (see appendix D).

\subsubsection{The Stuff of Nightmares (Knowledge)}
The characters find an ancient tome that if read in its entirely grants the reader a magical power. Once a creature has been granted this power, the tome loses its magic.

\subparagraph{Gift of Nightmares}
You force a creature that has the unconscious condition within 120 feet of you to have a terrifying experience. The creature begins to shriek from a vivid, pain-filled nightmare that lasts for 8 hours, or until they're forcibly awakened. You gain advantage on Intelligence saving throws, Intelligence checks, and are considered proficient when doing an Intelligence (Arcana, History, Investigation, Nature, or Religion) check for 1 hour. If the tortured creature is woken during that time, these benefits immediately end. Once this ability is used it can't be used again until you finish a long rest.

\subsubsection{Carved in Flesh (Privation)}
The characters find a magic cutting knife known as the \textit{Knife of Stolen Resistance} (see appendix D).

\subsubsection{A Torturer in Hell (Roleplaying)}
The characters come across a past enemy, or some ancestral foe, being tormented in the Nine Hell. The torturers offer them a chance to wield the lash, or they can suggest even worse punishments for the villain.

\subsection{Jealousy}
\DndQuote%
{}
{''The lord of these Nine Hells, who was the pinnacle of good, brought low by jealousy.''}
{\textbf{Fierna}}
\subsubsection{The Jealous Beauty (Power)}
An archdevil or pit fiend offers to grant one of the characters a blessing to help in their quest. The Fiend claims to want to embarrass \textbf{Asmodeus}.

\subparagraph{Blessing of Comeliness}
If your Charisma score was less than 20, it becomes 20. Unfortunately, every time a creature with a Charisma score higher than you speaks to you, you must make a DC 20 Wisdom saving throw. If you fail, you fly into a jealous rage and attack them. At the start of each of your turns you may attempt to make the saving throw again, ending the jealousy on a success.

\subsubsection{Anything You Can Do (Knowledge)}
The characters find the \textit{Ring of the Copycat} (see appendix D).

\subsubsection{With Friends Like These (Privation)}
An \textbf{incubus} tries to sell an \textit{Amulet of Betrayal} (see appendix D) to one of the characters.

\subsubsection{Credit Where Credit Is Due (Roleplaying)}
An \textbf{amnizu} (see Monsters of the Multiverse) standing on the shore of the River Styx offers the characters a view of the mortal realms. A bard is singing the story of one character's greatest past exploits but tells it so that some other—a different character or some past NPC—is given the credit for their triumphs. The devil offers to ensure its servants in the mortal planes take the bard aside and explain the true version of events to them. After all, truth is a virtue, isn't it?

\subsection{Murder}
\DndQuote%
{}
{''Behold Dis, where murder is always for sale.''}
{\textbf{Dispater}}
\subsubsection{Two Kills for One (Power)}
The characters come across a powerful weapon wielded by an enemy that they must defeat. This weapon can be any rare or very rare weapon from the \textit{Dungeon Master's Guide} and it gains the following property, which anyone wielding this weapon becomes aware of the moment they touch it.

\subparagraph{Damnation}
Each time this weapon strikes a killing blow, a random evil mortal on another plane dies and is sent to the Nine Hells, prematurely. Likely this victim would always have ended up the Nine Hells anyways, but you rob them of any chance to redeem themselves before dying.

\subsubsection{The Heretic Priest (Knowledge)}
The characters come across a crusading \textbf{priest} of an opposing faith. They overhear him bartering for diabolic aid with a \textbf{succubus}. He is asking for help in finding a repository of lore. This hidden library contains many tomes that are sacred to a god that one of the characters worships. It becomes clear that the cleric intends to destroy the library on the basis that heresy is contained therein. The cleric is obviously easy prey, and the characters could kill him without any repercussion.

\subsubsection{To Cheat Death (Privation)}
A succubus offers to inscribe a magic tattoo on one of the characters that she claims guards against death.

\subparagraph{Tattoo of Recovery}
The first time you would take damage and drop to 0 hit points, you can choose instead to drop to 1 hit point and be transported to an extradimensional space. While in this space you regenerate 10 hit points at the start of each turn. You may leave the space at any time and reappear in the spot that you left or, if that space is occupied, the nearest unoccupied space.

However, the tattoo makes it clear that its power is drawn from prisoners held in a dungeon on the mortal planes, to be sacrificed as living batteries for tattoos such as this.

\subsubsection{A View to a Kill (Roleplaying)}
While \textbf{Koh Tam}'s barge is docked at any port along the River Styx, a corruption devil (see appendix B) saunters on board and claims to be impressed by the deeds that the characters have accomplished over their careers. She lists some of their accomplishments, showing that she is indeed aware of them. She then claims to know of someone that has annoyed them in the past. You can choose an enemy from a previous adventure that is still alive or an NPC that is disliked. This person is a minion who has served the interests of the Nine Hells—helping them corrupt the very souls the characters are here to rescue—but their usefulness is at an end. She unveils a mirror that shows the NPC either asleep, eating, traveling, or performing some other mundane activity. Behind them, unsuspected but visible to the characters, is an assassin. At the word of any of the characters the assassin strikes and slays the target.

\subsection{Oppression}
\DndQuote%
{}
{''The most beautiful music is the silence of the good over the suffering of the oppressed.''}
{\textbf{Glasya}}
\subsubsection{The Angry Djinni (Power)}
An \textbf{amnizu} (see Monsters of the Multiverse) comes aboard \textbf{Koh Tam}'s barge and makes an offer to the characters. It will bind a \textbf{djinni} into their service for the small price of 500 gp. The characters later learn that the djinni can't be freed because it would immediately take its revenge. If the characters don't take the djinni along to help in the adventure, it causes problems for \textbf{Koh Tam} and the barge. Each time it is left on the barge, the djinni causes damage that requires 1d4 days of repair before it can resume its journey along the River Styx.

\subsubsection{Painful Knowledge (Knowledge)}
The characters find a magic iron mirror known as the \textit{Sage's Mirror} (see appendix D).

\subsubsection{No Pain No Gain (Privation)}
A suit of magic plate armor is discovered by the characters, either the \textit{Armor of Invulnerability} or \textit{Plate Armor of Etherealness}.

However, the magic comes from the souls of the living, and every time the wearer is hit for damage, the bound souls shriek in agony.

\subsubsection{A Palace Fit for a King (Roleplaying)}
The characters encounter a powerful devil (\textbf{pit fiend}, archdevil, or similar) who offers to build them a great fortress, temple, or other coveted structure in the mortal realm for when they return. The devil returns within a week or two to show how far construction is coming along. The devil opens a window to view the construction taking place in their world. They see that the labor is being done by a legion of bound workers who are being worked mercilessly and without cease until the project is complete.

\subsection{Pride}
\DndQuote%
{}
{''You mortals should take pride in how you can make a Hell out of any Heaven.''}
{\textbf{Asmodeus}}
\subsubsection{Gaining the Advantage (Power)}
A devil offers to etch a magic clover-shaped tattoo into a character's skin.

\subparagraph{Clover Tattoo}
This tattoo has 3 charges and regains expended charges at dusk. As an action you may expend 1 charge to gain advantage on your next ability check. However, if you fail that ability check this item loses all charges and you have disadvantage on all ability checks until dusk.

\subsubsection{You Must Speak Up (Knowledge)}
A devil offers to etch a magic tattoo into a character's skin.

\subparagraph{Tattoo of Knowledge}
The tattoo gives you access to knowledge you normally would not have. As an action you may cast either of the following spells. Once one of them is cast, neither can be cast again until you finish a long rest. The spells are:


\begin{itemize}
\item \textit{Contact Other Plane}
\item \textit{Scrying} (the target's spell save DC is always 20)
\end{itemize}

However, if you ever fail an Intelligence or Wisdom ability check, you have disadvantage on all further Intelligence or Wisdom ability checks until you finish a long rest.

\subsubsection{A Thing of Beauty (Privation)}
The characters find the \textit{Amulet of Appearance} (see appendix D), which ensures that whoever wears it always appears immaculate.

\subsubsection{Taking Pride in Your Work (Roleplaying)}
Nothing is more important than your good name. A devil offers a bargain appropriate to the character, that bestows the character a special power. Every time the character kills or hurts a non-evil creature that didn't start hostilities, they gain a Pride Point. If the character finishes a long rest without having spent a single Pride Point since their last long rest, they lose all remaining Pride Points. They can spend a Pride Point to gain advantage on a single Charisma (Deception, Intimidation, or Persuasion) check.

\section{The Dukes}
On the River Styx, the deadliest threat the characters may encounter arises from two menacing infernal vessels: one is a warship led by the archdevil brothers, Adramalech and Morax, the other, a submersible commanded by \textbf{Abigor}, an archdevil in the service of \textbf{Baalzebul}. More information on these infernal ships can be found in chapter 11, "Hunted by the Dukes."

These archdevils come after the characters once they've completed one of their group patron's objectives. Once this happens, \textbf{Asmodeus} realizes that the characters might be more tenacious than he anticipated and so sends the brothers to take them captive. On the other hand, \textbf{Baalzebul} wants vengeance on \textbf{Koh Tam} and the characters for summoning him to the mortal plane, thus he sends \textbf{Abigor} to hunt them down.

\section{Life in the Nine Hells}
Surviving in the Nine Hells is entirely dictated by one's ability to follow mindless legal codes and a self-preservation instinct that has no qualms about backstabbing. New arrivals find that, unless they offer a uniquely powerful skillset, they're the lowest of the low, used only as cannon fodder in the Blood War. Moving through the ranks requires either manipulating the strict laws of the realm, betraying comrades, or commanders—or sometimes both. Devils that find themselves higher in the hierarchy tend to lead relatively comfortable lives, so long as they obey their superiors and maintain a steady stream of souls. For the battle-hungry, the Blood War always longs for more combatants, and proving oneself on the battlefield is a surefire way to a quick promotion. As for the various non-devil inhabitants, existence remains possible so long as it doesn't interfere with devilish schemes. For mortals? Just try not to die.

\section{Features of the Nine Hells}
In the Nine Hells, everything is subject to corruption, from spells to magic items to the very map you use as a guide.

\subsection{Telepathic Spells}
The archdevils of the Nine Hells rule their realms with an iron fist. The devils that make up their ranks are heavily monitored and strictly controlled. All communication received in or sent from the Nine Hells can therefore be overheard. If they want to, the archdevils can eavesdrop on communication spells such as \textit{Message} and \textit{Sending}. Even when magic items (like \textit{Sending Stones} or a \textit{Helm of Telepathy}) are used to communicate, the archdevils can listen in.

In the Nine Hells, it is safe to assume that someone is always listening. This has left the occupants of the Nine Hells in a constant state of paranoia. Anyone who tries to communicate telepathically here, has the unnerving feeling of an infernal presence listening in. There is only one exception: Ring of Mind Shielding. Even in the Nine Hells, these rings shield telepathic communications from being overheard. These items are therefore highly sought after by those denizens of the Nine Hells who wish to keep their thoughts private from their lords.

\subsection{Teleportation Spells}
Many teleportation spells have altered effects when cast in the Nine Hells, including spells cast from magic items or artifacts, and class abilities that duplicate the effects of spells.

\begin{DndTable}[header=Effects of the Nine Hells on Teleportation]{XX}

Spells & Effect\\
\textit{Gate}, \textit{Plane Shift} or similar magic & These spells can only be used to travel to layers of the Nine Hells above the one that the caster is currently on.\\
\textit{Teleport}, \textit{Teleportation Circle} or similar magic & Each layer of the Nine Hells is considered its own plane of existence. Thus, it isn't possible to teleport between the layers.\\
\textit{Wish} & A \textit{Wish} spell can be used to transport the caster and companions to a deeper layer of the Nine Hells, provided the caster succeeds on a DC 18 Intelligence saving throw. On a failed save, the spell fails and has no effect.\\
\textit{Word of Recall} & This spell fails if it would take the caster to a layer of the Nine Hells lower than the one they're currently on.\\
\end{DndTable}

\subsection{Devilish Distances}
The Nine Hells is in constant flux, with the landscape warping to the whims of the devils, and at times, the Nine Hells itself. The distance between locations—and the time required to travel—is up to the DM. A trek towards an infernal tower might take only two hours, but the return journey could last weeks. And locations that the characters were certain were in a particular location one day may relocate themselves to another area of the Nine Hells on their next visit.

\begin{DndSidebar}[float=!b]{Tiax to the Rescue}
If the characters ever become separated from \textbf{Koh Tam}'s barge, have \textbf{Tiax} show up after 1d4 days and complain to them about how hard it was to track them down. He casts \textit{Teleport} to return them to the barge.



\end{DndSidebar}

\subsection{Using an Infernal Map}
As the characters explore the Nine Hells, \textbf{Koh Tam} sometimes provides an infernal map when they enter one of the adventure regions. Whenever a character attempts to travel to a location marked on one of these maps, they must match what they see with what is drawn on the map. To do this, they must succeed on a DC 15 Wisdom (Survival) check. If they're being helped by a native of the region, then they make this check with advantage.

Success means that they arrive at the location they're traveling to. However, depending on the adventure, they may still encounter something unexpected along the way.

Failure means they become lost. They wander for hours to the point of exhaustion. Unless they make camp for a short or long rest, each character gains a level of exhaustion. Failure might have additional effects depending on the adventure.

\subsection{The River Styx}
In the mortal world, springs bring water to the surface, trickling lines of liquid coming together, forming streams of ever-increasing size until they feed into great rivers. The River Styx is different, like the Ouroboros, it feeds in on itself. Smaller branches circle around, and, on an adjacent layer of the Nine Hells, feed back into the river. This is of great use to ferry pilots who are familiar with the river as they can use its branches as shortcuts to the upper layers of the Nine Hells and the important locations that can be found there.

The waters of the River Styx aren't just dangerous to navigate. Unless immune to the river's effects, a creature that drinks from the Styx or enters the river is targeted by a \textit{Feeblemind} spell (save DC 20). A creature must repeat the saving throw whenever it starts its turn in the river, until it fails the save. A feebleminded creature can drink from the Styx and swim in its waters without suffering any additional deleterious effects.

If a creature fails its saving throw and remains under the spell's effect for 30 consecutive days, the effect becomes permanent (no save) and the creature loses all its memories, becoming a near-mindless shell of its former self. At that point, nothing short of a \textit{Wish} spell or divine intervention can undo the effect.

Water taken from the River Styx loses its potency after 24 hours, becoming a harmless, foul-tasting liquid. However, arcanaloths, night hags, and other fell creatures might know rituals that can prolong the water's potency.

\chapter{Avernus, the Eternal Battlefield}
\DndDropCapLine{T}{}he topmost layer of the Nine Hells, Avernus serves as both an entrance and a defensive line for the lower layers. With the Blood War raging, battles between demons and devils are continuously fought on its rocky soil. The remnants of recent battlefields scar the landscape, reminding inhabitants to be vigilant.


\section{Running This Chapter}
Before running this chapter read the "Avernus Overview" section. That section provides the information you need to guide the characters through Avernus and nudge them towards the War-Slough where the items they seek may be found.

\subsection{Encounters}
The characters may sail the River Styx through Avernus without much incident but once they leave the waters and travel overland, dangers arise. Roll at least once on the Random Encounters in Avernus table each time the characters move between major locations on the map.

\begin{DndTable}[header=Random Encounters in Avernus]{cX}

d6 & Encounter\\
1 & A lone \textbf{halog} (see appendix B) begins to follow the characters.\\
2 & A group of \textbf{black abishai} and \textbf{white abishai} led by a \textbf{blue abishai} (see Monsters of the Multiverse) fly overhead. If the characters follow them, they find Tiamat's Lair but risk discovery.\\
3 & Six spined devils and eight imps circle in the sky above the remains of a recent battlefield.\\
4 & Far in the distance, the characters observe infernal war machines chasing each other across the wastes.\\
5 & The sky suddenly darkens as the flying fortress of the archdevil \textbf{Zariel} (see appendix A) passes over the characters, casting a lengthy shadow. Two affliction devils (see appendix B) leap from it to attack the characters, while the fortress continues its journey across Avernus.\\
6 & A storm of flame and dust that is equivalent to the \textit{Storm of Vengeance} spell (with a spell save DC of 19) erupts, centered around the characters and lasting for 10 rounds.\\
\end{DndTable}

\subsection{Locations}
\textbf{Koh Tam} can provide a summary of Avernus when they arrive (see "Key Locations in Avernus").

\subsection{Koh Tam}
Avernus is more widely known than other layers, but if the characters are unfamiliar, \textbf{Koh Tam} can provide them with the information they seek. His focus should be on the Blood War however, and the eternal battleground that is the War-Slough. In locations such as that, great power is to be found.

\subsection{Objectives}
Make sure you keep track of your players' objectives to help guide them to the appropriate areas in Avernus where they might complete their goals. Ideally, they should be met by Rexlexkala in the War-Slough and encouraged to explore the Ruined Tower and the Ichor Lake in that area.

The following objectives can be attained in Avernus:

\begin{DndTable}[header=Objectives in Avernus]{XX}

Objective & Location\\
Phylactery of the Brother & The War-Slough: The Ichor Lake\\
Phylactery of the Queen & The War-Slough: The Ruined Tower\\
\end{DndTable}

\subsection{Temptations}
Have the characters encounter at least one temptation during their time in Avernus. You can of course have them encounter more if you want. If the characters give in to temptation, use the information in appendix E to keep track of their corruption level.

\section{Avernus Overview}
Thousands of years of construction projects have led to myriad artificial channels, sluices, and gates within parts of the River Styx. Where it used to flow, along the edge of the realm, there remains only a dry riverbed and sparse vegetation. Due to the powerful nature of the river, the walls and mechanisms are in a never-ending cycle of repairs.

Aside from the Styx flowing through the center of the layer, Avernus lacks water and vegetation. Much of the realm consists of lifeless hills, jagged mountains, and rocky terrain. Outcroppings of crystals and metals jut through the surface like spikes, accompanying the boulder fields as surface features. Lakes and rivers of lava dot the surface, though not nearly as dense as the fire-aspected layers beneath. Geysers and fumaroles decorate the landscape, occasionally spewing scalding steam and toxic clouds.

While the sky remains unlit by celestial bodies, fireballs streak through and explode, bringing a fiery light to the horizon. Sometimes the fireballs collide with the land, resulting in smoldering impact craters. Closer to the land, the realm is bathed in a deep red light of unknown origin. To mortals the air is thick, due to the ash, toxins, and dust that contaminate the atmosphere. Bones and blood of creatures long perished decorate the less-traveled areas, as well as the decaying battlefields.

Notable locations within Avernus include the Bronze Citadel, the Pillar of Skulls, and Tiamat's Lair. Connecting each landmark with the rest of the realm is the Great Avernus Road, constructed by \textbf{Bel} to facilitate military transport.

\textbf{Bel} used to rule the realm, but has since lost the position to \textbf{Zariel}, a fallen angel turned devil. Her goals, even as the Blood War rages, are to restore Avernus to its former glory, and to remove Tiamat from the realm. \textbf{Zariel}'s subjects primarily include devils that form the defensive armies of the Blood War, but natives to Avernus still remain.

\begin{DndSidebar}[float=!b]{Previous Adventures in Avernus: Bel}
If you have run Baldur's Gate: Descent into Avernus, then \textbf{Bel} might have replaced \textbf{Zariel} as the ruler of Avernus.



\end{DndSidebar}

\subsection{Leaving Avernus}
Only a few methods allow travel to the lower layers of the Nine Hells. The cave system of Tiamat's lair contains a hidden and perilous path to Dis. In addition, one of the largest iron towers of the City of Dis pierces through Avernus, with an opening next to the Pillar of Skulls. For an experienced sailor, the Styx is also viable, as it snakes its way around a mountain down into the next layer and beyond, and this is \textbf{Koh Tam}'s preferred route.

\subsection{Features}
Avernus' combination of oppressive heat and supernatural malevolence weighs on the bodies and souls of those who aren't evil. A non-evil creature treats normal travel through Avernus as a forced march and must make a Constitution saving throw at the end of each hour of travel. The DC is 10 + 1 for each hour of travel. On a failed saving throw, a creature gains one level of exhaustion.

\subparagraph{Optional Rule}
Evil pervades Avernus, and visitors feel its influence. At the end of each long rest taken on this plane, a visitor that isn't evil must make a DC 10 Wisdom saving throw. On a failed save, the creature's alignment changes to lawful evil. The change becomes permanent if the creature doesn't leave the plane within 1d4 days. Otherwise, the creature's alignment reverts to normal after a day spent away from the Nine Hells. Casting the \textit{Dispel Evil and Good} spell on the creature also restores its original alignment.

\section{Key Locations in Avernus}
Some of the major landmarks found in Avernus are described below.

\subsection{The Bronze Citadel}
The Bronze Citadel is the largest and longest-standing city in Avernus, serving as both a political and militaristic capital. Within the city lies a basalt fortress, from which \textbf{Zariel} rules the realm. The city is under a constant state of retrofitting, with metal and bone fortifications and buildings being improved or modified every day. Though sprawling across hundreds of square miles, the city is a massive structure with numerous walls and war machines, befitting its position so near the beachhead of the Blood War.

As \textbf{Koh Tam}'s barge prepares to leave the Bronze Citadel, read:

\begin{DndReadAloud}{}
\textbf{Koh Tam} speaks to you. "Let me offer you some additional advice. First, we should not linger too long here, for the eyes of a multitude of devils are upon this citadel, watchful and wary, and others look for conscripts to draft into the construction or war efforts. Secondly, remember that I know much about the Nine Hells; please ask me anything and if I have the answer, I'll offer it.

"Finally, if we ever need to backtrack and return to a previous layer of the Nine Hells, remember that my barge is capable of plane shifting. But only ever upwards, never downwards. With all that out of the way, where do you want to sail to first?"

\end{DndReadAloud}

If the characters ignore \textbf{Koh Tam} and linger, perhaps seeking out magic items or other wares, they might become tempted by \textbf{Asmodeus} to acquire something that they shouldn't.

\begin{DndSidebar}[float=!b]{Temptation of Anger}
\textbf{\textit{"Gauntlets of Rage"}}



During the characters' time in the Bronze Citadel, a corruption devil (see appendix B) offers to sell a spellcaster a special pair of gauntlets for 500 gp or 2 Soul Coins. If a character purchases the item, they now have the Gauntlets of Rage (see appendix D).



\end{DndSidebar}

\subsection{Tiamat's Lair}
A cave mouth leads into Tiamat's lair, a massive mountain and cave system which is one of the routes into Dis. But Tiamat protects the other layers of the Nine Hells from potential invaders by guarding these passages vigilantly. Access to her lair is even more difficult because a massive pit of wretched souls blocks the way. This pit and numerous abishai keep outsiders from entering Tiamat's Lair.

\textbf{Koh Tam} advises against attempting passage this way, for the devils of the Nine Hells have many agents in Avernus and the characters' forcing their way through Tiamat's Lair would certainly be noticed.

\subsection{The Pillar of Skulls}
At the edge of the layer, next to the entrance to Dis, lies the Pillar of Skulls—a multiple-mile high monument to the Blood War. Skulls of fallen demons and devils are stacked on top of each other, creating a gruesome landmark. As the characters approach they hear a cackle and sigh, as if the skulls themselves have a story to whisper to them. But the noise appears to just be the hot wind scouring bone and dirt.

\subsection{The War-Slough}
The location most interesting to characters seeking phylacteries is the War-Slough, a vast battlefield upon which demons and devils have been slaughtering one another for eons. Here is where the characters may find what they seek in Avernus...

\begin{DndSidebar}[float=!b]{Temptation of Harm}
\textbf{\textit{"The Stuff of Nightmares"}}



If the pillar of skulls is searched a tattered book with a leathery cover is found. If a character takes the item, run the appropriate temptation event found in chapter 2.



\end{DndSidebar}

\section{Adventure: The War-Slough}
The War-Slough in Avernus, site of the Great Fray, was home to a devastating war between devils and demons and what's left is a nightmarish collection of lost devil units, wandering demons, and corruption. The characters might come here searching for a soul phylactery.

If the characters tell \textbf{Koh Tam} that they wish to visit the War-Slough, or if they are not certain where to travel to next, he recites the following tale:

\begin{DndReadAloud}{}
"Once there was a battle that pushed even the limits of the Nine Hells to past their breaking point. Devils still talk of the Great Fray in hushed voices, mostly because of the infernal bureaucracy it still causes. On that day the skies of Avernus opened, and the demons just kept coming through, host upon host of them. And the blazing and barbed armies of the Nine Hells rose to meet them, as they always do, and it went on. And on. A single engagement of the Blood War that raged back and forth while, on more mundane planes, mortal lifetimes came and went, empires rose and fell, species were created and became extinct. Countless demons and devils fought and were destroyed or banished. Countless. That was the problem." \textbf{Koh Tam} pauses, as if not quite sure of where he's taking the story. He glances at you, then continues.

"It's the chaos that did it, of course. All those demons concentrated on a single field, slaughtering the lesser hosts of the Nine Hells through sheer ebullient numbers but being slaughtered in turn. Their ichor being shed in tides that would drown a nation, and again, and again. The demons have been driven off long since, but in some way the Abyss won that engagement. The War-Slough is what's left. And that is where we must search." \textbf{Koh Tam} smiles grimly. "It should be easy to find... a great sunken bowl miles across, a morass of half-drowned trenches, craters, broken war engines and ruined fortifications. There are artifacts of vast power lost somewhere in the war-strata of the Slough, stockpiles of souls and every unique war-ending engine that failed to end anything. All of it buried in the corrupting mud."

\end{DndReadAloud}

\subsection{Advice from Koh Tam}
If the characters are struggling to decide what to do, or about to set out without being sufficiently prepared, \textbf{Koh Tam} can help nudge or guide them with the following advice:


\begin{itemize}
\item There are some devils that seek mortals to enter the War-Slough and find discarded weapons from the Blood War. They should search the borders for such a devil as it might help them find what they're looking for.
\item The War-Slough can warp those who stay too long, so they should make haste.
\end{itemize}

\textbf{Koh Tam} also gives them an infernal map. Show the players the map of the War-Slough in appendix F. They can use this map to decide which points of interest they want to explore. Refer to the "Using an Infernal Map" section of chapter 2. A failed or successful Wisdom (Survival) check to use this map has a 50 percent chance of triggering an encounter on the Random Encounters in the War-Slough table.

\subsection{Demon Blood}
The corrupting mud of the War-Slough exists due to the sheer unparalleled quantity of demonic blood that has been shed in this single concentrated engagement. Normally the ground of the Nine Hells is inhospitable terrain for anything so undisciplined. Chaos runs off it like water from weaved cloth. Even the resilience of the Nine Hells has limits. Eventually they fought ankle-deep, knee-deep. The stuff just didn't soak away, and the field became a nightmarish quagmire reeking of the wild Abyss. The very sight of the War-Slough sends shudders down the spines of lesser devils.

Because of the corrupted demon blood, when a creature has spent more than an hour in the WarSlough it must make a DC 15 Constitution saving throw. Fiends, Oozes, Plants, and Undead automatically succeed on the saving throw. On a failed save, the creature is warped by the demon blood, as determined by rolling on the Flesh Warping table.

When a creature takes a short rest in the WarSlough, they must succeed on a DC 15 Constitution saving throw to resist having their flesh warped. A long rest requires succeeding on a DC 20 Constitution saving throw. A spell that removes a curse ends all demon blood related warping effects on a creature.

\begin{DndTable}[header=Flesh Warping]{cX}

d10 & Effect\\
1 & The creature's eyes push out of its head on the end of stalks.\\
2 & One of the creature's legs grows longer than the other, reducing its walking speed by 10 feet.\\
3 & The creature's eyes become beacons, filling a 15-foot cone with Vision and Light when they're open.\\
4 & The creature's ears tear free from its head and scurry away; the creature has the deafened condition.\\
5 & The creature's arms become tentacles with fingers on the ends, increasing its reach by 5 feet.\\
6 & The creature's legs grow incredibly long and springy, increasing its walking speed by 10 feet.\\
7 & The creature grows a whiplike tail, it can now use a bonus action to make one tail attack (treat the tail as a whip.)\\
8 & The creature's ears become wings, giving it a flying speed of 5 feet.\\
9 & The creature's body becomes unusually brittle, causing the target to gain vulnerability to bludgeoning, piercing, and slashing damage.\\
10 & The creature grows another head, causing it to have advantage on saving throws against the charmed, frightened, or stunned conditions.\\
\end{DndTable}

\subsection{The Forever War}
The War-Slough isn't abandoned. To evacuate it you'd have to master it, and the place is beyond even \textbf{Zariel}'s dominion. It doesn't matter what nebulous period of time has passed since the Great Fray: the combatants on both sides are immortal, and there are plenty left in the Slough still trying to fight the dead battle. On the nominal side of the Nine Hells, units of single-minded merregons (see Monsters of the Multiverse) from the handfuls to the hundreds slog endlessly through the mire, cut off from any orders after too many usurpations and reappointments up the chain of command. Beyond any hope of returning to the infernal hierarchy they attack just about anything on sight, including each other. Brutal, yet simultaneously pathetic and desperate, they scream codewords and countersigns as though hoping to hear some word of authority that they're past the point of ever recognizing.

A number of infernal generals were lost to the Great Fray, and some are still in there. Fighting the infinitely shifting hosts of the abyss broke them one after another, leaving them unable to trust any external authority for fear of falling for a demonic trick. They maintain their camps and shrunken domains in the ruins of fallen fortifications and tell themselves that they continue the fight even as they raid one another or make sham alliances that they know will be broken in heartbeats.

Aside from the pathetic remnants of devil armies, the War-Slough has developed an ecology of sorts. There is always a meal to be had for an opportunistic scavenger. One-off mutated creations roam forlorn over its churned surface or slither through its trenches. Demons stalk and worry at the ineffable meat of fallen Fiends, growing fat and corrupt with demonic taint. There are tentacled things in the lakes of demon-ichor and vast corpse-worms that burrow through the slimy soil. For every hideous monster there is something worse still to prey upon it.

\subsection{War-Slough Encounters}
Every time the characters take a short or long rest in the War-Slough roll on the Random Encounters in the War-Slough table. Then roll a d6. On a 1–3 the random encounter interrupts the rest. On a 4–6 the encounter happens a few minutes after the rest is finished.

\begin{DndTable}[header=Random Encounters in the War-Slough]{cX}

d6 & Predator\\
1 & A flock of three shredwings (see appendix B) circles in the distance.\\
2 & A gang of five bulbous \textbf{hezrou} (speed 20) erupt from a pool of demonic ichor.\\
3 & A lost legion of one hundred \textbf{merregon} (see Monsters of the Multiverse) march in formation, ignoring anything that doesn't get in their way. A headless \textbf{horned devil} leads them.\\
4 & A corpse worm (\textbf{purple worm} with the Fiend type) erupts from the ground beneath the characters.\\
5 & A squad of a dozen insane \textbf{merregon} mistakes the characters for demons.\\
6 & Two horned devils fused together (they always occupy the same space but have their own turns) lead a battalion of two dozen \textbf{merregon}. One horned devil orders the battalion to attack, while the other orders them to march in formation. Each round there is a 50 percent chance of the merregon obeying either leader.\\
\end{DndTable}

\subsection{Escaping the Slough}
If a creature tries to leave the War-Slough after being there for more than 8 hours, the demonic taint fights to stop them from leaving. The creature must succeed on a DC 15 Wisdom saving throw or find that they don't want to leave. They may make another DC 15 Wisdom saving throw to escape only after they finish a short or long rest. However, if they haven't killed a living creature since their last attempt, the DC increases to 25.

\subsection{War-Slough Locations}
There are those who go to the War-Slough by choice. Why would anyone? Well, there are all those lost artifacts and items of fearsome power that are in there somewhere, if one could only track them down. There's enough potent treasure for any number of daring hunters to retire on, plus stored souls to power every warship on the Styx for half an eternity.

\subsubsection{W1: Sentry Tower}
Technically, \textbf{Zariel} has set a prohibition against any devil entering the War-Slough, to avoid the region being destabilized any further. Because of this prohibition, ambitious devils bargain with mortals to trick or bribe them to enter the Slough seeking rumored treasures. One such is an \textbf{amnizu} (see Monsters of the Multiverse) named Rexlexkala. He approaches the characters if he sees them on the border of the War-Slough. Once he has convinced them that he doesn't intend harm, he explains what he has to offer them:

\begin{DndReadAloud}{}
"You will find it easy to enter the War-Slough. While there is a ring of towers about the perimeter, they are far-spaced, and sentry duty is a punishment detail for lowly devils who survived failed actions in the Blood War. The constant prickle of the demonic at the backs of these pathetic devils serves as sufficient distraction that diligent intruders can easily sneak past. However, getting out of the Slough is harder. The truth is that the Slough, through the sheer weight of wasted life and passion and its unique admixture of the infernal and the demonic, has become something like a sentient thing, a kind of insensate god of wasted wrath. It does not want to let go of anyone or anything that falls into its clutches, and over time its influence turns everything within it to a lost, furious combatant seeking to perpetuate a meaningless war. However, I have something that can help you avoid that fate, but only if you help me retrieve something."

\end{DndReadAloud}

Rexlexkala has learned the location of an infernal artifact designed millennia ago to end the Great Fray. The artifact is a horn built to allow devils to listen in on the thoughts of demons while they're in Avernus, even powerful demon lords. Such an item is worthless to mortal adventurers but is of supreme value to an archdevil such as \textbf{Zariel}. Rexlexkala will give 10 Soul Coins for the horn. In addition, he provides each character with a vial of Condensed Order (see appendix D.) In addition to the substance's normal effects, it also allows them to leave the War-Slough without requiring a Wisdom saving throw first. Finally, he has studied the WarSlough for centuries and can give advice on how to find things within. If the characters agree to his terms, Rexlexkala gives them a map to a rift which is located deep within the War-Slough (reveal area W4 on the infernal map). In this rift lie buried many weapons left over from the Blood War.

He will also answer questions from the characters if they agree to his task or he feels his life is threatened.


\begin{itemize}
\item If the characters ask him for help in finding the phylactery of the queen, he says that they will find it at an abandoned tower which is located deep within the War-Slough. It dates from a time when the War-Slough was much smaller.
\item If the characters ask him for help in finding the phylactery of the brother, he tells them that they will find what they're looking for at the largest of the demon-ichor lakes located in the Slough.
\end{itemize}

\subsubsection{W2: Ruined Tower}
\begin{DndReadAloud}{}
Through the haze, you see the outline of a tower ahead, rising over five hundred feet into smoke choked skies. Winged devils circle around the tower's peak and companies of devils march near its entrance.

\end{DndReadAloud}

The characters notice the tower when they get within 300 feet of it. It was built centuries ago to guard the borders of the War-Slough when it was much smaller. Now the top floor of the tower is the roost of a marilith. She wraps her snake body around an infernal throne, a remnant from the Blood War. The throne allows a Fiend to control lesser devils as if its occupant were at the top of the infernal hierarchy. The marilith, a demon, uses it to amuse herself by having all the devils in the region perform whatever whim comes into her head. The marilith is easier to defeat if the characters enter the tower without her knowing, but there are several obstacles to achieving that.

The first obstacle is a half dozen barbed devils and a half dozen bearded devils marching back and forth along the road to the entrance. If observed, they're seemingly following orders that only they hear. They'll attack anyone that they see, but it is easy to avoid them when they're marching away from the tower. When first discovered they're 90 feet away from the tower and marching back towards it. Every turn, there is a 50 percent chance they reverse direction. Any noise makes them Dash towards the direction of that distraction, but they ignore anything that happens overhead, because they're, after all, mere foot soldiers. If these devils engage in combat with the characters, each round, one devil attempts to Dash back to the tower at "A" and if any devil reaches it, this obstacle is considered failed.

The second obstacle is only encountered if the characters approach the tower by air. If they do so, and are noticed, they're confronted by four bone devils that patrol the skies about 600 feet above ground. One of the bone devils will Dash to the window ("C" on the map) 100 feet away that leads into the "Throne Room". If it reaches that window, this obstacle is considered failed.

The final obstacle is within the tower. The interior is mostly an empty ruin, with a simple stone staircase rising upwards. But 1d4 shadow demons linger at each location marked "B" on the staircase (see the "Ruined Tower" map). These shadow demons attempt to hide from the characters and, if not discovered, they ascend to the next "B" (or to the "Throne Room," if the uppermost "B" has been reached). If any shadow demons reach the throne room, this obstacle is considered failed. If discovered, the shadow demons attack the characters.

\subparagraph{The Throne Room}
The \textbf{marilith} is dressed up in the garments of a queen and attended to by four bearded devils that wear the fine clothing of palace courtiers and are busy entertaining the marilith. The marilith and her minions are unaware of the heroes until they announce themselves, or come within 15 feet of the throne, unless they've failed 2 or more of the previous obstacles, in which case she is very much aware of them. There is an additional bearded devil for every failed obstacle.

Once aware of them, the marilith and her devil courtiers mock the characters. They do this until the intruders attempt to leave or show any sort of disrespect, then they attack. If the marilith is killed her body vanishes and she appears on her throne 1d6 rounds later with 100 hit points. Only by destroying the throne can she be killed. The throne has an AC of 15 and 100 hit points.

\subparagraph{The Phylactery of the Queen}
If a character that has chosen this phylactery enters this room, read the following upon entering:

\begin{DndReadAloud}{}
A servant dressed in rags washes the floor at the base of a throne occupied by a serpentine woman wearing a royal purple gown and a sparkling crown. As the servant glances your way, it topples the bucket of dirtied water beside it. The water rolls out, hissing as it evaporates against the stone floor. The queen shrieks and chides the servant for being a clumsy fool and orders it to refill the bucket. The servant's shoulders slump as it retrieves the bucket. The face of the servant is strangely familiar.

\end{DndReadAloud}

In a fight this servant (\textbf{commoner}) doesn't participate, even if attacked directly. Once the queen is permanently slain, the servant contorts and twists, leaving behind the character's soul phylactery.

\subsubsection{W3: Ichor Lake and the Demon Nest}
A huge circular rift drops one 100 feet into the depths of the Slough. Small streams of demonic ichor cascade over its edges. At the bottom of this foul place the demonic ichor of the War-Slough has gathered into a large lake that seethes and bubbles with unconstrained malice and possibility. At its center is an island that resembles a massive ant nest with entrances to caves or perhaps burrows. The smoke, haze, and muck make it nearly impossible to discern whether creatures occupy the tunnels below.

For some time, \textbf{Zariel}'s underlings held out hope that a finite number of demons remained, and they would slowly dwindle away over the eons until none were left and the Slough might revert to Avernus' natural state. The most recent Infernal Survey of the place found ample evidence that concentrations of demonic material (resembling nests) were spawning random demons at least as fast as the creatures were being killed off. Through these nests, the WarSlough perpetuates itself. The island at the center of the lake is one such nest.

\subsubsection{W3A: Entrance}
\begin{DndReadAloud}{}
Moist ichor covers the floors and walls. There's a pounding noise in the distance, as if a large machine is rumbling away, performing some unfathomable work. The tunnels seem to undulate, in rhythm with the noise, and in places you think you see terrible faces etched into the stone. Many tunnels appear to lead into caves where ichor has accumulated.

\end{DndReadAloud}

\subsubsection{W3B: The Caves}
Each cave has \textbf{hezrou} and \textbf{barlgura} embedded in its walls, many only partially formed. Blood vessels in the walls pump demonic ichor into the catatonic demons. Attacking any of them awakens 1d6 of each.

\subsubsection{W3C: Demons}
A mixture of two dozen dretches and \textbf{manes} mill about mindlessly in these caves. They ignore intruders but descend with fury on anyone who attacks them.

\subsubsection{W3D: The Lake}
\begin{DndReadAloud}{}
A great lake of ichor sits before you, undulating to the demonic rhythm that has only grown louder. Along the shoreline you see a small demon clamber out of the ichor, its arm covered in dozens of eyeballs. It pays you no heed and shambles down a tunnel. Across the grotesque expanse is an island upon which sits a heart the size of an elephant. This heart pumps demonic ichor into the blood vessels permeating these caverns.

\end{DndReadAloud}

Coming into contact with the ichor of the lake forces a creature to succeed on a DC 15 Constitution saving throw or roll on the Flesh Warping table. Gibbering dretches or manes demons that spawn from the demonic heart swim across the lake. They ignore anyone they see and emerge from the lake with a new deformity (use the Flesh Warping table for examples). These demons tend to accumulate in the various caves. The characters need to cross the lake to reach "The Heart."

\subsubsection{W3E: The Heart}
The island is about 60 feet from the shore. Once it is reached, the characters can walk up to the heart itself. The heart has an AC of 8 and 200 hit points. Killing it causes the lake to become inert. However, as soon as it is damaged, the heart will awaken the hezrou and barlgura to come to its aid. On initiative count 10 of each round 1d4 of each arrive in the chamber and begin swimming across the lake.

When the lake becomes inert all the demons act as if they're under the \textit{Confusion} spell. However, instead of acting normally on rolling a 9 or a 10 (as per the spell), this roll results in the demon inflating like a balloon and then exploding, showering demonic ichor in a 10-foot-radius. Anyone in that area must succeed on a DC 15 Constitution saving throw or be forced to roll on the Flesh Warping table.

\subparagraph{The Phylactery of the Brother}
A character who seeks their brother's soul, notices movement in the ground underneath the demonic heart.

\begin{DndReadAloud}{}
There's a glimmer of movement beneath sticky ichor pooling under the pulsating heart. For a moment you see your brother's face protrude from the floor. He glances your way but briefly, a grimace of pain on his face, before succumbing and falling back into the ichor.

\end{DndReadAloud}

Once the heart is slain, the ichor can be pulled aside, and the brother's phylactery recovered. The character gains the Phylactery Benefit associated with it.

\subsubsection{W4: Canyon Trove}
\begin{DndReadAloud}{}
A vast and twisted battlefield extends before you, cleaved in two by a massive canyon. Whether the canyon existed before that great battle, or arose because of it, you cannot be certain. A trail, paved in corpses, descends into the canyon.

\end{DndReadAloud}

\begin{DndSidebar}[float=!b]{Temptation of Anger}
\textbf{\textit{"Sentient Artifact"}}



The characters' attention is drawn by a voice from a nearby pile of corpses. If they investigate, they find a powerful artifact buried within. This is an opportunity to conjure a tempting reward for one of the characters; select a character and an item they've been coveting. That is the item they now find. You may also select an item that multiple characters might want so that they may argue among themselves over who gains possession of it.



Once one character has taken the item, run the "Sentient Artifact" temptation event found in chapter 2.



\end{DndSidebar}

This isolated canyon descends into the depths of the Slough and is where the real treasures of the War-Slough are found, including the horn that Rexlexkala desires. Strange things occur sporadically as the characters descend, the place twisted and distorted by the great war. When a character least expects it (perhaps just after a rest, or in the midst of battle), a \textbf{grotesque tentacle} with a shrieking mouth instead of suckers bursts from the canyon wall to flail about wildly before withdrawing again. The tentacles are harmless unless the characters decide to attack one (then the tentacles have the statistics of a \textbf{guardian naga}, but 0 movement speed, lawful evil alignment and no spellcasting trait).

Though the tentacles can be safely ignored, there are real dangers in the canyon. And to find items of true value, the characters have to make an effort to explore, heading deep into the canyon's bowels. Each hour the characters make a DC 20 Wisdom (Survival) check to search for clues that might lead them to the treasures they seek. Failure results in an encounter, while success leads to treasure.

\subparagraph{Encounter}
The first time the characters fail their search results in them encountering the "Temptation of Anger." Afterwards, on future failures, roll a d8. 1–4 means an empty cave; 5–6 means they stumble upon 1d2 \textbf{glabrezu}; a 7–8 means they find 1d4 \textbf{hezrou}.

\subparagraph{Treasure}
Every time the characters succeed in their search for clues, they have a 50 percent chance of finding Rexlexkala's Horn. Otherwise, they roll on the Infernal Items table. For a description of these items see appendix D. When they find the item, it might be in the clutches of a tentacle that erupts from the canyon walls, or it might be lodged at the bottom of a bubbling pool of \textbf{demon ichor} (which is a black pudding with 120 hit points and immunity to fire damage). After 3 such items, they'll no longer find any more infernal items here and even successful searches are now treated as failures.

\subparagraph{Rexlexkala's Horn}
A \textbf{nalfeshnee} carries the artifact that Rexlexkala seeks and won't relinquish it without a fight. The artifact looks like a metallic horn, but if touched it begins to wriggle like a grotesque worm. Other than this it can be handled safely. An \textit{Identify} spell reveals that it is meant to be swallowed to attune to it, but if anyone other than a devil does so, they gain none of its powers and instead are afflicted with a random major insanity that can't be removed until the artifact is extracted.

\begin{DndTable}[header=Infernal Items]{cXX}

d10 & Item & Danger\\
1 & Bracers of \textbf{Asmodeus} & None\\
2 & Infernal Plate Armor & Pool of \textbf{demon ichor}\\
3 & Stygian Spear & Tentacle\\
4 & Canian Fork & Pool of \textbf{demon ichor}\\
5 & Demonbone Polearm & 2d6 barlguras\\
6 & Sword of Retribution & Yellow mold\\
7 & Infernal Amulet & Angry \textbf{horned devil}\\
8–10 & Ancient Blood War Weapon (see Ancient Blood War Weapons table) & None\\
\end{DndTable}

\begin{DndTable}[header=Ancient Blood War Weapons]{cX}

d4 & Blood War Weapon\\
1 & This is a fleshy, blood-hued sphere and if an action is used to squeeze it hard enough, it erupts into a 500-foot-radius Cloudkill that harms only Fiends.\\
2 & This box casts Vision and Light for 200 feet and Vision and Light for 1,000 feet. Covering it negates the light. An \textit{Identify} spell reveals that the box is full of divine energy and will explode if opened. If an action is used to open the box, every creature in a 100-foot-radius takes 70 (20d6) radiant damage and the box is destroyed.\\
3 & This ornate fan is comprised of dozens of partly melted blades. When an action is used to unfurl the blades, spores are released in a 40-foot-radius. Any demon in the area must make a DC 17 Constitution saving throw. On a failed save, they take 35 (10d6) necrotic damage and one affliction from the Flesh Warping table. Once used, the fan can't be used again until a long rest is finished.\\
4 & This is a bag with 10d10 strange seeds inside. If a seed is planted in the Nine Hells, it immediately grows into a corrupted \textbf{shrieker}. Any demon that doesn't have the deafened condition and starts its turn within 30 feet of a corrupted shrieker takes 7 (2d6) psychic damage. This damage doesn't stack with damage from other corrupted shriekers.\\
\end{DndTable}

\chapter{Dis, the City of Burning Iron}
\DndDropCapLine{N}{}amed for its ruler, the Iron Lord \textbf{Dispater}, Dis is the second layer of the Nine Hells and largely dominated by the city of the same name which is ringed by jagged mountains. Within the city, characters may find some of the items they're searching for.


\section{Running This Chapter}
Before running this chapter read the "Dis Overview" section. It provides you with everything you need to guide your players through Dis and its sprawling city.

\subsection{Encounters}
After arriving in Dis, the players must either navigate the River Styx or follow the Iron Road to venture into the City of Dis. Their destination is the Agora of Floating Knives with Orishada's Palace at its center. Whilst on the river, journeying into the city, roll at least once on the Random Encounters in Dis table.

\begin{DndTable}[header=Random Encounters in Dis]{cX}

d6 & Encounter\\
1 & A pack of six lemures have escaped the mines they were working. Four hell hounds pursue close behind.\\
2 & Eight spined devils attack without provocation, mostly just for sport. They retreat from opponents who give as good as they get.\\
3 & The characters notice an \textbf{imp} crouching on the rooftops, presumably spying on them.\\
4 & A githyanki trade caravan makes its way toward the City of Dis. A half dozen githyanki knights on young red dragons provide protection.\\
5 & A lone \textbf{barbed devil} approaches and demands 100 gp as tax for traveling across \textbf{Dispater}'s realm. If he is attacked, he calls upon a pair of war devils (see appendix B) circling hundreds of feet above to protect him.\\
6 & A blazing orb of iron, seemingly ejected from a faraway mine, lands on a random character. It explodes and that character, and each creature in 20-foot-radius sphere, must make a DC 17 Dexterity saving throw. A creature takes 70 (20d6) fire damage and 70 (20d6) bludgeoning damage on a failed save, or half as much damage on a successful one. The collected fragments of the orb are worth 2,000 gp.\\
\end{DndTable}

\subsection{Locations}
When the characters arrive in the city have \textbf{Koh Tam} provide information about the city and its different areas (see "Key Locations in the City of Dis" for a description). The players need to make their way to the Agora of Floating Knives, the trade district in the City of Dis. If the players decide to venture into any of the other areas of the city, you can use the information provided in this section to guide them on a short excursion. However, the conditions of these locales may quickly deter them from venturing there too long.

\subsection{Koh Tam and Tiax}
At the Agora, give the players plenty of time to explore. \textbf{Koh Tam} or \textbf{Tiax} can provide them with directions to the different areas of the Agora and its market stalls. If \textbf{Tiax} is accompanying the characters into the Agora, roll on the Troubles with Tiax table in chapter 1 to see if their outing with the gnome is eventful.

At this point, \textbf{Koh Tam} reminds the characters that his barge is able to Plane Shift to a layer of the Nine Hells above the current one, at any time. It requires a full action for him to activate the ability.

\subsection{Objectives}
Make sure you keep track of your players' objectives and lead them to the corresponding areas so they can complete their goals. Once the players have explored all the different areas within the Agora of Floating Knives, make your way to Orishada's Palace.

The following objectives can be attained in Dis:

\begin{DndTable}[header=Objectives in Dis]{XX}

Objective & Location\\
Instrument of the Bards, Ollamh Harp & The Agora of Floating Knives: Barges of Theater\\
\textit{Cloak of Invisibility} & The Agora of Floating Knives: The Market of Killers\\
Phylactery of the bounty & The Agora of Floating Knives: The Market of Victims\\
Phylactery of the chosen one & The Agora of Floating Knives: The Arena\\
Phylactery of patricide & Orishada's Palace: Kitchens\\
Horn of Valhalla, Iron & Orishada's Palace: Aelvette\\
\textit{Wrought-Iron Tower} & Orishada's Palace: The Floating Duke\\
\end{DndTable}

\subsection{Temptations}
Have the characters encounter at least one temptation during their time in Dis. You can of course have them encounter more if you want. If the characters give in to temptation, use the information in appendix E to keep track of their corruption level.

\begin{DndSidebar}[float=!b]{Troubled Rest}
The characters might need to be encouraged to use all their various spells and means of obtaining information to track down the items and phylacteries they seek. At times you may also want to disturb their rest—perhaps resurfacing some of the nightmares described in the Introduction—as they approach layers of the Nine Hells where objects they need to obtain are to be found.



Nudge them towards stopping and exploring each layer appropriately but allow them the freedom to come up with clever ways to find what they seek.



\end{DndSidebar}

\section{Dis Overview}
Outside the city proper, the layer is a mountainous crag land rich in ore deposits. \textbf{Dispater} ordered the construction of the Iron Road, a pathway made of iron and cobblestone, to navigate the mountains and extract the minerals located within. Since its construction, the mountains are mined, and ore transported back to Dis for refinement.

Coiling down mountains high enough to reach Avernus is the Styx. It winds its way down towards Dis and crosses directly through before moving downwards towards Minauros. The city has an intricate canal system which the river flows through. With high walls and a deep bottom, navigators of the Styx see only the tops of buildings as they pass through the city. An artificial runoff near the middle of the city was created, with a large iron gate, allowing travelers to enter the city proper by taking a slight detour.

\textbf{Dispater} remains the ruler of Dis and an ally of \textbf{Asmodeus}, and he intends for it to stay that way. Fortification of \textbf{Dispater}'s layer continues unabated at his behest, in preparation for some future war. Assuming inhabitants follow the rules of the realm, they're given relative freedom to go about their business.

\subsection{The City of Dis}
The city lies nestled within a fiery ring of mountains, with the only two ways in or out being the Iron Road and the River Styx. It is described by visitors as impossibly large and yet cramped at the same time. Like the Iron Road, its buildings are constructed from a mix of cobblestone and burning iron, both mined from the mountains of the realm. Because the iron is heated, it emits smoke and covers the city in a smog, making it difficult to see and breathe for mortals. Additionally, the city and the Iron Road are bathed in a warm light from the burning steel, with the surrounding mountains remaining dark aside from the lit mines. When the realm doesn't smell like smoke, it takes on the scent of burning flesh and blood.

Extraplanar travelers often stop in the city before moving on, due to its reputation as a hub for trade and commerce within the layers. Many bazaars and shops are set up offering strange wares from different planes, and \textbf{Dispater} encourages any on the plane to establish contracts whenever possible. These contracts usually involve souls and the greater devils, who make homes in the massive iron towers scattered throughout the town. Beneath the city lies a labyrinthian structure of dungeons, said to contain mortals from the Material Plane. Vents in certain areas of the city allow bloodcurdling screams to float up from below.

\subsection{Leaving Dis}
Unlike leaving Avernus, traveling deeper into the Nine Hells from Dis is quite difficult. The only two known passages are the Styx and a staircase, both of which descend to Minauros. Though the location of the staircase isn't publicly known, everyone assumes it to be located within \textbf{Dispater}'s Iron Tower, making it difficult to find without connections. The Styx is especially dangerous, as the drop from Dis to Minauros consists of steep inclines, rapids, and waterfalls. It takes a seasoned captain or a large amount of luck to navigate the passage, leaving nascent souls or mortals without means of traveling deeper.

\subsection{Features}
Iron itself bends to the influence of Dis. Whenever a good-aligned creature touches nonmagical iron, they take 4 (1d8) fire damage. This includes armor or weapons that creatures might be using, ammunition, storage devices, keys, etc. A Protection from Good and Evil spell or similar effect allows creatures to resist this effect for the duration.

As the most lawful of the infernal planes, visitors to this plane become bound to their words and contracts. Any time a mortal gives a promise, enters into a contract, or otherwise ensures their end of a bargain, they find themselves magically compelled to follow. A creature that breaks such a bond while on Dis takes 32 (5d12) necrotic damage and gains a level of exhaustion.

\section{Key Locations in the City of Dis}
Some of the major locations found in the City of Dis are described below.

\subsection{The Iron Tower of Dispater}
Though several iron towers decorate the city, none are as large as the Iron Tower, \textbf{Dispater}'s fortress and palace. It is visible from anywhere in the city, and through enchantment magic always appears to be very close. \textbf{Dispater} maintains complete control over the tower—including the capability to shift its form and interiors to match his desires—and the tower can only be entered if he permits. If \textbf{Dispater} has no interest in meeting you, even reaching the tower requires tremendous effort and immense willpower.

If the characters decide to visit, they struggle to ever reach the tower. Wherever they're in the city, the tower always appears right around the next corner, but no matter how long they walk, it never gets any closer. If the characters succeed on a DC 21 Wisdom (Survival) check and spend 2d6 hours searching, they finally reach the Iron Tower. If they fail the check, they waste 2d6 hours and are no closer, but they do trigger a random encounter.

\begin{DndReadAloud}{}
After a final push, you reach the tower. Up close it is impossibly large. You squint as you look up but cannot make out its peak. The tower feels unwelcoming, and you feel a great heat emanating from it. There is no sign of door or window—the tower simply looms over you as if in judgment.

\end{DndReadAloud}

Unless the characters have business inside of the tower (and \textbf{Dispater} allows them entry) the characters can't enter.

\subsection{The Garden of Delights}
Devils from all layers are known to travel to Dis solely for the Garden, where all manner of luxuries and pleasures can be enjoyed. The Garden provides the income needed to improve \textbf{Dispater}'s city and armies, and it keeps his citizens happy. The Garden also serves as a distraction for the denizens of the Nine Hells from the tortures of everyday existence in this bleak place. However, the characters are mortal beings and the illusions found inside overwhelm their minds. If they explore the Garden of Delights, they become enchanted by its beauty and may forget themselves. For considerable time they're unable to remember their mission and wander the Garden aimlessly.

\begin{DndSidebar}[float=!b]{Temptation of Deceit}
\textbf{\textit{"The Liar"}}



While the characters are exploring the Garden, a succubus approaches the character with the highest Charisma and offers to give them the power of deception. Run the appropriate temptation event found in chapter 2.



\end{DndSidebar}

\begin{DndReadAloud}{}
Walls of sandstone rise above you and in front of the wooden doors to the Garden of Delights stand several attractive servants, beckoning you forward. You feel a strong compulsion to approach these comely beings.

\end{DndReadAloud}

The characters may pay 50 gp apiece to enter the Garden. Once inside they're handed cool and delicious (and illusory) drinks. Now roll a d100—this is the number of hours the characters waste, captivated by the illusory delights the garden provides. At the end of this time, a character must succeed on a DC 17 Charisma saving throw to even want to leave, if they fail the saving throw, they spend another d100 hours within and must repeat the saving throw. When the characters finally succeed on the saving throw and depart, each gains a level of exhaustion.

\subsection{The Great Prison of Mentiri}
Mentiri acts as a maximum-security prison for the criminals of the Nine Hells. Inside, inmates are tortured and indoctrinated to ensure they follow the rules of the Nine Hells. This secure prison is located deep within Dis and within its walls contains many different kinds of prisoners—outsiders, captives from the Blood War, and others who have broken the laws of the Nine Hells in one way or another. The only way into the prison is being captured and condemned to rot there. The characters are lucky that they've no reason to enter.

\begin{DndSidebar}[float=!b]{Temptation of Oppression}
\textbf{\textit{"No Pain No Gain"}}



Read the following the first time the characters travel towards the prison:



\begin{DndReadAloud}{}
Traveling towards Mentiri you pass a group of insect like devils transporting captives to the prison. Suddenly, an elf dressed in expensive looking armor makes a dash for freedom. Within seconds, the man is beset by a huge infernal raptor that carries him screaming into the sky.

\end{DndReadAloud}



Later, on their way down from the prison the characters find the elf's armor on the side of the road but no sign of its former owner. A successful DC 15 Intelligence (Arcana) check reveals it to be powerful magic armor. Run the appropriate temptation event found in chapter 2.



If the characters try to save the elf, they must fight the shredwing and the four pain devils (see appendix B for both) guarding the rest of the prisoners. The elf is already dead when the battle starts; it was just an animated corpse of an elf that has been dead for more than a thousand years.



\end{DndSidebar}

\subsection{The Agora of Floating Knives}
The Agora of Floating Knives is one of the largest trading hubs in the known planes and is welcoming to non-devils.

\section{Adventure: The Agora of Floating Knives}
The Agora of the Floating Knives is an enclave for non-devils made up of floating vessels located on the edge of the great City of Dis. Here mortals can purchase a variety of infernal goods and services.

In the City of Dis you can buy or sell anything, so long as that thing has a chance of harming someone, somewhere. Many of these markets are run by devils for devils, locked away within the cramped streets, behind the high walls. However, Dis is one of the few places in the Nine Hells that can be characterized as even slightly welcoming to outsiders. Travelers visit Dis for many purposes, but a great many come specifically for the Agora of Floating Knives. If the characters tell \textbf{Koh Tam} that they wish to visit the City of Dis, he directs them to the Agora because of its accessibility to outsiders.

\begin{DndReadAloud}{}
\textbf{Koh Tam} gestures towards the city before you. "Dis is a city of commerce. If any of what you seek is here, the Agora of the Floating Knives is our most appealing destination. The Agora is an enclave for non-devils and is situated such that we can sail downriver and reach it without even having to risk the tangled confines of the city proper. Within its hundreds of floating stalls and huts we might find just about anything. And if we don't, we could risk an audience with Orishada, he is a powerful devil, and dangerous as any other, but he runs the Agora and is more willing to speak to outsiders than most."

\end{DndReadAloud}

\subsection{Advice from Koh Tam}
\textbf{Koh Tam} (or \textbf{Tiax}) offers the following advice to characters, if necessary:


\begin{itemize}
\item They should hire an imp to guide them, or allow \textbf{Tiax} to do so, if he's willing.
\item They might find what they're searching for in the markets. If not, then they might need to seek an audience with the ruler of the Agora, Orishada, a powerful \textbf{amnizu} (see Monsters of the Multiverse).
\item They're likely to see many disturbing things, but if they let on to any discomfort it could lead to trouble.
\end{itemize}

This is a location well known to many travelers and \textbf{Koh Tam} can provide a standard map of the Agora (see appendix F) to the characters.

\subsection{Approaching from the River Styx}
Off the great runoff from the Styx, within the City of Dis, there is an artificial lagoon, square and shadowed by high walls. Here can be found the Agora of Floating Knives.

\begin{DndReadAloud}{}
It is a great armada of floating stalls and huts, houseboats, and rising from its heart an entire palace borne on the deadly waters. Imps and other menial demons can be seen plying gondolas along the narrow channels separating each establishment.

\end{DndReadAloud}

\subsection{The Outer Agora Locations}
Visitors relying on the services of the imps (for the small sum of 10 gp) should be very specific in giving their destination, as the diminutive gondoliers delight in malevolently misinterpreting instructions every bit as much as an archdevil bargaining for a soul.

The periphery of the Agora consists of a ramshackle of little boats and rafts that offer the sort of goods that might be found in a dozen dens of vice across Dis, most often from the clawed hands of those who have been barred from more respectable infernal haunts. Vendors from other realms bring goods not native to the Nine Hells in an attempt to tempt diabolic custom, whilst petty devils from lower layers haul up treasures and resources unique to the infernal realms to sell at the small shops located here. There is a lively trade in Condensed Order, a silvery powder that can be extracted from those of a lawful persuasion. Devils bound for the warfronts of Avernus take flasks and snuff boxes of the stuff to fortify themselves against exposure to the raw chaos of demons.

Some of the more common items that can be purchased are listed in the Items for Sale table:

\begin{DndTable}[header=Items for Sale]{XX}

Item & Cost\\
\textit{Condensed Order} & 50 gp\\
Dream bottle & 50 gp\\
\textit{Dust of Sneezing and Choking} & 250 gp\\
\textit{Oil of Sharpness} & 2,500 gp\\
\textit{Oil of Slipperiness} & 250 gp\\
\textit{Potion of Fire Breath} & 250 gp\\
\textit{Potion of Mind Reading} & 750 gp\\
\end{DndTable}

\begin{DndSidebar}[float=!b]{Temptation of Jealousy}
\textbf{\textit{"With Friends Like These"}}



While the characters are perusing wares, an incubus approaches and tries to sell an Amulet of Betrayal (see appendix D) to one of the characters. The price he offers it for is 1000 gp but he is willing to reduce it to 500 gp, if appropriately convinced.



\end{DndSidebar}

There is also a thriving trade market. These items are readily available for anyone who succeeds on a DC 15 Intelligence (Arcana) check. However, the buyer is only interested in trading for a magic item with Soul Coin as an additional cost.

\begin{DndTable}[header=Magic Items Trade Cost]{XX}

Magic Item & Trade In\\
\textit{Mirror of Life Trapping} & Rare or very rare wondrous item, plus 3 Soul Coin\\
\textit{Dagger of Venom} & Uncommon or rare weapon, plus 1 \textit{Soul Coin}\\
\textit{Crystal Ball} & Rare or very rare wondrous item, plus 3 Soul Coin\\
\textit{Efreeti Bottle} & Rare or very rare wondrous item, plus 3 Soul Coin\\
\textit{Ring of Shooting Stars} & Rare or very rare ring, plus 3 Soul Coin\\
\textit{Robe of Stars} & Rare or very rare magic cloak or robe, plus 3 Soul Coin\\
\textit{Staff of Fire} & Rare or very rare magic rod or staff, plus 2 Soul Coin\\
\textit{Staff of Frost} & Rare or very rare magic rod or staff, plus 2 Soul Coin\\
\textit{Spellguard Shield} & Rare or very rare shield, plus 2 Soul Coin\\
\textit{Nine Lives Stealer} & Rare or very rare weapon, plus 2 Soul Coin\\
\end{DndTable}

\subsubsection{A1: The Arena}
\begin{DndReadAloud}{}
The cheering and jeering of spectators, as well as grunts of pain and shrieks of terror, alerts you to an open arena floating upon a barge. Several lodges are connected to it and from these buildings various challengers emerge, either to fight upon a field of battle, or to take their place against one another at tables of various games of chance. Just as you approach a devil knocks a piece sideways on a game board and rises, shouting out in victory. The tiefling sitting across from the devil moans as guards drag him into one of the buildings.

\end{DndReadAloud}

If the characters attempt to intervene, remind them it is important not to draw attention to themselves and that everyone participating in these games has chosen to do so. They may participate in the games themselves—feel free to have them fight against various devils and creatures of the Nine Hells. Or they can participate in games of skill or intelligence (use the Chosen One section below as guidelines). The prize for victory is 1 Soul Coin. The cost of failure is either death in the arena, or in the case of more civilized competitions, an hour of torture in the shacks (resulting in that character gaining one level of exhaustion).

\subparagraph{The Phylactery of the Chosen One}
If a character has selected the sin of pride, their phylactery is here. They see their soul competing in various arena games and constantly losing. As they approach each contest they do so with sunken shoulders and a worried frown. When they lose, they're mocked and punished. And they always lose.

The character may tap in and take their soul's place. They must compete in the following three challenges:


\begin{description}
\item[Test of Strength.]
The character must fight a \textbf{chain devil}.

\item[Test of Skill.]
The character must either best a \textbf{horned devil} in a complicated game of knot tying or compete against an \textbf{erinyes} to see who can calm an enraged hellcat (see appendix B). Tying the knot requires a successful DC 18 Dexterity (Sleight of Hand) check whereas taming the creature requires a successful DC 15 Wisdom (Animal Handling) check. If the character fails the check, they can try to make a DC 16 Charisma (Deception) check. If successful they bluff their way into getting a second (and final) chance on the original check.

\item[Test of Wits.]
The character challenges an \textbf{amnizu} (see Monsters of the Multiverse) to a dice game popular in the Nine Hells. The character must make a successful DC 15 Intelligence check to win the match.

\end{description}

When a character loses a challenge, they gain a level of exhaustion. If the character rests or otherwise recovers from that exhaustion, they must start all three matches over. They may repeat each challenge as often as needed to win but can have no help from anyone. Once they win all three matches, a succubus or incubus appears and hands them their phylactery.

\subsubsection{A2: Neogi Shop}
One of the shops in the Outer Agora is run by the enigmatic and ruthless \textbf{neogi} (see Monsters of the Multiverse). These spider-like monsters are planar travelers that trade in magic items and prisoners. Their shop usually only stays in the same place for 1d6 weeks. Each visit to the shop brings the chance of danger, for sometimes the neogi aren't looking to sell items, but are more interested in acquiring new victims. A new shopper must make a DC 20 Charisma (Persuasion) check. If they succeed, the items on sale are real. If they fail, then the items are illusory and the neogi have prepared an ambush. A \textbf{neogi master} leads the ambush, with an additional \textbf{neogi} (see Monsters of the Multiverse) and an \textbf{umber hulk} for each shopper. More powerful customers might mean the neogi bring an additional 1d6 gray renders (see Monsters of the Multiverse). The Illusory Items table lists items that they sell, along with what the items really are if the neogi have decided to ambush the shoppers.

\begin{DndTable}[header=Illusory Items]{XXX}

Magic Item & If Illusory & Cost\\
\textit{Apparatus of Kwalish} & A \textbf{hook horror} charmed by the neogi & 100,000 gp\\
\textit{Brazier of Commanding Fire Elementals} & An \textbf{xorn} restrained by \textit{Dimensional Shackles} & 10,000 gp\\
\textit{Cape of the Mountebank} & A rotting nonmagical cloak & 5,000 gp\\
\textit{Carpet of Flying} & A cheap carpet & 20,000 gp\\
\textit{Cloak of Arachnida} & A \textbf{cloaker} that is the pet of one of the neogi & 15,000 gp\\
\textit{Deck of Illusions} & A set of playing cards with depraved images & 400 gp\\
\textit{Demon Armor} & \textbf{Hezrou} restrained by Iron Bands of Bilarro & 40,000 gp\\
\textit{Dwarven Plate} & A \textbf{shield guardian} whose control amulet is worn by the neogi master & 35,000 gp\\
Figurine of Wondrous Power, Golden Lions & Same as above except a pair of hellcats (see appendix B) & 11,500 gp\\
Figurine of Wondrous Power, Marble Elephant & Same as above except a single \textbf{maelephant nomad} (see appendix B) & 11,000 gp\\
Figurine of Wondrous Power, Onyx Dog & \textbf{Displacer fiend} (see appendix B) under permanent Reduce spell and held by an adamantine leash & 11,000 gp\\
\textit{Helm of Brilliance} & \textbf{Intellect devourer} & 25,000 gp\\
\textit{Mace of Terror} & A \textbf{flameskull} attached to an \textit{Immovable Rod} & 9,000 gp\\
\textit{Sword of Life Stealing} & A rusty nonmagical sword & 12,000 gp\\
\textit{Wings of Flying} & Muzzled \textbf{harpy} in a cage & 11,000 gp\\
\end{DndTable}

\subsection{The Inner Agora Locations}
\begin{DndReadAloud}{}
Beyond the mean skirt of barter and desperation can be found more magnificent emporia, two and threestory vessels moored in a hierarchy jealously maintained by the stallholders. There are grand platforms where demonstrations of very particular skills can be performed. Each auction house sells a specific brand of the unspeakable.

\end{DndReadAloud}

This is the Agora most visitors know, and it throngs with outsiders whose villainy, while it can't approach that of the devils themselves, nevertheless marks them out in their homes as exceptional. The Agora of Floating Knives is most famed as the center of the Nine Hells' murder trade. The four most popular destinations include the Barges of Theater, Market of Killers, Market of Victims and Orishada's palace.

\subsubsection{A3: Barges of Theater}
\begin{DndReadAloud}{}
A host of hideous entertainments are visible for the delectation of the casual visitor, carried past on boats and as advertisements at the doors of various establishments. Street torturers ply their trade to the laughter and applause of discerning crowds in much the same way as acrobats or stage magicians might in a less cursed city.

\end{DndReadAloud}

Some of the horrors that a visitor might come across include:


\begin{itemize}
\item An archery competition where contestants permanently snuff out souls by shooting at writhing larvae pinned to a target board.
\item A still living Red Wizard whose body has been transformed into a teleportation circle.
\item A musical instrument where each key drives a skewer into the flesh of some poor mortal whose throat has been twisted to utter one note of perfect pitch. This last venue is run by a tiefling named Gazrak.
\end{itemize}

Woe betide any visitor who shows distress at such theatre. With the Agora being so readily accessible, Orishada keeps a watching legion of imps haunting the rooftops like pigeons, watching for any intruders who lack the appropriate immoral fiber. Those who do are followed by the imps until a squad of four bone devils arrive to exterminate them.

\subparagraph{Conclave (Bard)—Instrument of the Bards (Ollamh Harp)}
Gazrak (chaotic evil, tiefling \textbf{archmage}) sells all manner of grotesque musical instruments, but he also owns an Instrument of the Bards, Ollamh Harp. He is fond of the harp but also greedy. He suggests an outrageous price at first—30 Soul Coin. A successful DC 17 Charisma (Persuasion) check can either lower the price to 20 Soul Coin or allow the characters to pay 100,000 gp instead of using Soul Coin. Giving him a rare or very rare magic item as part of the exchange reduces the price by half.

\subsubsection{A4: The Market of Killers}
\begin{DndReadAloud}{}
Ahead of you is a crowded barge that criers announce in a dozen languages as the Market of Killers. Living killers from across the planes have flocked here seeking the highest fees. Devils tout dead killers, infamous in their lifetimes, but who are allowed back into the realms of the living for one more murderous job... at the right price. Crowded and frenetic auctions are held every hour, with the prize being a victim of the winner's choice killed via some particularly grotesque trademark method.

\end{DndReadAloud}

In the Agora, perhaps the most mundane commodity one can hire at the Agora is the service of assassins. The characters can hire a killer to perform a single kill, not as some sort of retainer. If the characters are willing to spend enough, they might be able to use one of these assassins to rid the multiverse of a villain that they met during this adventure (or a previous adventure). You decide whether the assassin succeeds. A failed attempt might mean that the target turns the tables and sends an assassin against the characters.

\subparagraph{Conclave (Rogue)—Cloak of Invisibility}
To obtain the cloak the characters must hire the assassin (neutral evil, lizardfolk \textbf{assassin}) who wears it, and then kill them for it. When they threaten the assassin he mocks them because he is protected by a powerful archdevil.

\begin{DndReadAloud}{}
The assassin hisses at you. "Lay down your arms and move on, for know this: my patron is powerful and appreciate-s my s-special s-skills. They will be displeased if you s-strike out against me." The lizardfolk grins a toothy grin and slips the hood of the cloak overhead.

\end{DndReadAloud}

They're still willing to assassinate an individual for the characters or meet them at a later date. If the characters do kill the assassin, they draw the ire of an archdevil (your choice), who makes the remainder of their adventures in the Nine Hells more dangerous. To avoid this, characters must purchase a dispensation from Orishada prior to killing the assassin. See the "Floating Duke" section. The dispensation costs 8,000 gp.

\begin{DndTable}[header=Assassins for Hire]{XXc}

Killer & Cost & CR/Level that they can kill\\
\textbf{Thug} & 100 gp & 1\\
\textbf{Gladiator} & 250 gp & 3\\
\textbf{Assassin} & 500 gp & 5\\
Undead \textbf{drow mage}* who summons a \textbf{bone devil} instead of a shadow demon & 750 gp & 6\\
Undead \textbf{assassin}* & 1,000 gp & 7\\
Undead githyanki knight* on a young shadow dragon (red) & 1,500 gp & 8\\
The yat-ja slayer (lizardfolk \textbf{assassin} with a \textit{Cloak of Invisibility}) & 5,000 gp & 9\\
Saint of killers (fallen \textbf{planetar} with no wings or fly speed, dressed in dark leather) & 20,000 gp & 12\\
\end{DndTable}

\subsubsection{A5: The Market of Victims}
\begin{DndReadAloud}{}
Across from the barbarity that is the Market of Killers, there is the Market of Victims. Not a commodity you'd think would be in short supply in the Nine Hells, but there are victims and then there are victims. Each seller claims to cater to a very specific criteria of victim—a particular bloodline, a type of virtue, saints, angels, demons, emperors, and more.

\end{DndReadAloud}

The kidnappers of the Nine Hells have contacts in every plane, on every world. They offer a comprehensive service. Many a high priest of some bloody god or other has procured the perfect sacrifice by sending their acolytes to the Agora with sufficient largesse.

\subparagraph{Phylactery of the Bounty}
The character seeking the phylactery that contains the soul of their bounty—the matron of a powerful noble family—is drawn to this location. They immediately notice a secluded corner of the market where an affliction devil (see appendix B) haggles with two barbed devils. One barbed devil clutches a bulging leather satchel. The character realizes the phylactery they seek is in the satchel. If they don't intervene the affliction devil soon takes the satchel and flies away. Instead, if the character interrupts, they can offer to buy the satchel. This negotiation requires a DC 16 Charisma (Persuasion) check. If successful, the character manages to outbid the affliction devil and can purchase the satchel for 5,000 gp. If the result of the check is 20 or higher, the character's negotiation skills are insurmountable, and the affliction devil leaves negotiations earlier. The character only needs to pay 3,000 gp to purchase the satchel. Once the satchel is obtained, the character takes possession of this phylactery and receives the Phylactery Benefit associated with it.

\subsubsection{A6: Orishada's Palace}
\begin{DndReadAloud}{}
The palace at the center of the Agora is a vast leviathan of pumice and incandescent metal lit on the outside by some gruesome burning substance. Its walls bristling with a million barbed hooks, the only easy way in seems to be the massive double doors at its front.

\end{DndReadAloud}

Orishada and his select underlings offer unique services, both to the Nine Hells and to outsiders. They're proud of their range, in fact—enough that at least affecting to be a connoisseur of such things can win one an audience.

\subsection{An Audience with Aelvette}
Using the iron knocker summons an \textbf{imp} that opens a slot just below. The imp promptly demands to know what valid business justifies the interruption. As long as a visitor describes the desire to purchase something that Orishada provides, then the imp permits them to meet with the majordomo of the palace.

Orishada's majordomo is Aelvette, an ancient \textbf{night hag} whose past has taken her in and out of the lower planes for longer than even most devils can recall. She has complete control of her appearance, able to decide for each viewer whether she should appear as alluring or hideous at any given moment. Her particular game is the Condensed Order trade; the silvery powder, highly valuable to devils, that can be extracted from those of a lawful persuasion.

Aelvette is sly, humorous and doesn't play by the rules in the way true devils do, which can make her either the best or absolutely the worst person to approach at the palace. The night hag attempts to discern if any of her guests are of an appropriate mindset to consent to have the Condensed Powder extracted, for a fair price. If she doesn't think this is likely, then she offers to play a game of chess or other logical exercise from an appropriate visitor. If the visitor agrees and loses the game, then this yields a gleaming phial of Condensed Order for Aelvette, but leaves the unwilling donor confused and off-balance (akin to the poisoned condition). The condition lasts for 5 minutes.

\subparagraph{Playing Aelvette}
The characters find themselves in a battle of wits with Aelvette. Either they have to impress her during a conversation in which she is testing their mettle, or they have to (try and) beat her at chess. The character with the highest Intelligence score should challenge Aelvette. Defeating the hag at chess or impressing her during conversation requires three successful Intelligence checks. The DC is 10 on the first check, 15 on the second check, and 20 on the third check. A character who has proficiency in the Insight skill can add their proficiency bonus to each of these checks.

Once Aelvette has finished plying her trade in powder, she informs her guests that her master, Orishada, wishes to see them. Before they meet with him, the night hag suggests that they partake of the services provided at the kitchens or the salon.

\subparagraph{Conclave (Ranger)—The Iron Horn of Valhalla}
Aelvette is in possession of the Horn of Valhalla, Iron, a character searching for it has to steal it from her. When the characters enter Aelvette's chambers read the following:

\begin{DndReadAloud}{}
You are ushered into a large, oval room by the imp. It quickly scurries away leaving you alone with the figure at the center of the room. Behind her, the floor to ceiling windows overlooking the Agora are opened, allowing the smoke and stench of the city to penetrate the air. The sparsely furnished room is dominated by a large stone table. On it a weathered chessboard is set up, the pieces positioned in what appears to be a game in progress. Aelvette, the majordomo, stands beside the table, studying the game set up before her. She appears deep in thought and doesn't pay you any attention. Your attention is drawn to a magic horn mounted on the wall behind her...

\end{DndReadAloud}

This is the Iron Horn of Valhalla. The characters need to distract Aelvette to steal it. The most opportune moment would be when her mind is elsewhere engaged. Most likely, the characters will opt to have one of them move on the horn while she is playing chess with another character. Some characters might want to stage a different distraction. Whatever scenario they choose, a character stealing the horn needs to succeed on a DC 16 Dexterity (Stealth) check. If they fail, she notices the character creeping behind her and attacks—unless the character manages to convince her their intention was not to steal.

\subsubsection{The Kitchens}
\begin{DndReadAloud}{}
This massive kitchen contains sights and smells only a devil, or one interested in the macabre, could enjoy. Hulking chefs, whose limbs seem stitched together, shuffle about preparing meals as a devil directs them.

\end{DndReadAloud}

In addition to managing the important business conducted in the palace, the majordomo, Aelvette, also presides over the devil chefs of the palace's famous Omnivorous Kitchen, an appalling side-business of hers known to the jaded and wicked across the planes. Who, after all, knows the best way to prepare just about any meat, to appeal to any palate, but a hag? Those who wish to sample the most forbidden delicacies flock to the palace, as do those who would learn the ghastliest recipes and culinary techniques across the multiverse. Not for the faint of heart or the weak of stomach.

If their appetite is overpowering their morals, the characters may choose to order one of the delicacies from Aelvette's kitchen. Prices are listed on the Forbidden Delicacies table.

\begin{DndTable}[header=Forbidden Delicacies]{XX}

Meal & Cost\\
Boiled shank of minotaur & 25 gp\\
Candied spider eyes & 3 gp\\
Live owlbear & 1,000 gp\\
Lobe of grell & 12 gp\\
Poached stirge eggs & 1 gp\\
Prime cut of pegasus & 100 gp\\
Roasted jackalwere & 10 gp\\
Toasted salamander & 15 gp\\
\end{DndTable}

The chefs working the kitchen are three flesh golems directed by a \textbf{horned devil}. Any violence in the kitchens brings down the wrath of all the chefs and their overseer.

\subparagraph{The Phylactery of Patricide}
A character who has chosen the patricide sin discovers that one of the flesh golems is being whipped by the horned devil for failing to follow a recipe in the correct manner. Read the following to that character:

\begin{DndReadAloud}{}
You realize the flesh golem's likeness is identical to yours and in the devil's face you see your own father's features reflected. A terrible sense of familiarity comes over you, reminding you of the dreams that have long been haunting you.

\end{DndReadAloud}

The character must end this cycle of torment to gain their phylactery. The moment they attack the horned devil however, all the flesh golems, including the one in whom the character sees their own likeness, defend the vile creature.

Once the golems and the devil are slain, read the following:

\begin{DndReadAloud}{}
The abused flesh golem looks up at you from where it has fallen. Then its flesh begins to burble and hiss, dissolving away, leaving behind a terrible stench and a tarnished phylactery that remains warm to the touch.

\end{DndReadAloud}

The player has recovered their soul's phylactery.

\subsection{Corteso's Salon of Experience}
\begin{DndReadAloud}{}
Tiny brains with legs run about this room that has been designated as a 'salon of experience'. Behind a thick glass window, you see a pool that contains a massive pulsating brain.

\end{DndReadAloud}

Corteso's Salon of Experience is another peculiar service offered at the palace. Many devils have an infinite curiosity about what it is to be mortal. The illithid \textbf{elder brain} (see Monsters of the Multiverse) Corteso has accumulated a vast library of memories—specializing in experiences of sin and being sinned against. As well as counting for high entertainment to devils, students from the college of Maladomini flock here to refine their education on how to live among and manipulate mortals. Corteso is always in the market for a truly unique experience, and always has something to offer in trade. Its illithid spawn have been known to ensure that its offers can't easily be refused. Corteso's other racket is secrets.

The multitude of minds that it and its many broods have consumed grant the ancient intellect a vast storehouse of dangerous knowledge. Buyers come to learn exactly what they need to say, and to whom, to have some other murdered without getting their own hands dirty. Corteso specializes in jealous lovers, swindled business partners, ancient family feuds, anything that sows discord and inflames hatreds between mortals. Just what brought such a creature as Corteso to the palace and Orishada's patronage is a secret nobody knows, certainly one the creature has no intention of selling. Rumor has it Corteso created and then survived the destruction of multiple broods, evaded the attempts of many heroes to end its prolonged existence, and fled to the Nine Hells pursued by gith assassins and the agents of other elder brains which it had betrayed. Innately treacherous, it is 'loyal' to Orishada because only in the archdevil's shadow is it safe.

Corteso is always protected by a half-dozen mind flayers and their \textbf{intellect devourer} pets. If asked where he obtains his memories, the characters are told their suppliers are many, but recently high-quality memories have been obtained from a location known as the Ineffable Trove, located in Minauros.

\begin{DndSidebar}[float=!b]{Temptation of Harm}
\textbf{\textit{"Carved in Flesh"}}



One of the butchering knives, currently ensconced in a hunk of meat, appears larger and somehow more ominous than the others. It can be retrieved without alerting the golems and devil with a successful DC 16 Dexterity (Sleight of Hand) check. The characters now have the Knife of Stolen Resistance (see appendix D).



\end{DndSidebar}

\subsubsection{The Floating Duke}
\begin{DndReadAloud}{}
Orishada is a great, broad devil wearing clothes of metallic finery, richly bedecked in both mundane and exotic jewelry. "I hope you're not here to waste my time. My direct superior is \textbf{Dispater} himself, lord of both city and layer, and he will be gracing me with his presence in a few hours. I see only very select buyers, which my majordomo assures me you are. I must assume you're here for the specific class of ware that I am known for: dispensations."

\end{DndReadAloud}

If his guests have no idea what he is talking about, Orishada—an \textbf{amnizu} (see Monsters of the Multiverse)—won't immediately order their deaths. He assumes it is part of some negotiation tactic. He summons one of his imps to explain what a dispensation is:

\begin{DndReadAloud}{}
"In the fine print of many an infernal contract there is a little clause almost never exercised. It's a 'boilerplate' clause that when exercised, gives an ironclad permission for the contract holder, their heirs or assigns, or any descendant or inheritor, to be murdered, without penalty or retribution. Such clauses have been around for many thousands of years, and my master, Orishada, makes it his business to collect them, or knowledge of them."

\end{DndReadAloud}

After the imp has finished, Orishada inquires as to whose dispensation they wish to purchase.

Given their multigenerational nature, a remarkable number of great and powerful people are covered by such clauses without ever knowing it. The dispensations that Orishada sells don't procure the death of the victim—that's what the rest of the Agora is for, after all. However, they do permit an assassin to disregard a variety of defenses up to and including divine favor, and protect the killer from curses, supernatural wrath, and other problematic fallout after murdering someone particularly important.

There is a 50 percent chance that any named evil creature/NPC is covered by a dispensation. This drops to 10 percent for good and neutral NPCs. A good rule of thumb is to have a dispensation cost 1000× the creature/NPC's challenge rating in gold pieces.

If his guests decline to purchase a dispensation or they're impolite or violent, Orishada becomes enraged and orders his personal guard to liquidate the intruders. These include a quartet of horned devils and a half dozen bone devils.

Orishada won't join in the slaughter, not because he doesn't want to, but because he can't. Orishada was a talented killer whose thousand-year career would raise the eyebrows of every assassin thronging the boats of the Agora. He finally allowed his bloodthirstiness to overcome his proper place, however, encompassing the death of one beloved to certain divine powers. In punishment, Orishada now lives out his endless life under a bitter curse. He can no longer inflict physical harm himself, not on the least insect, not even to a mortal. Because of this he uses his magic to aid his guards as best he can but won't intentionally cause damage.

\begin{DndSidebar}[float=!b]{Insurmountable Odds}
All the devils at this location defend either \textbf{Dispater} (see appendix A) or Orishada if the characters attack either. It is unlikely characters can survive an outright onslaught and they may need some guidance to behave as cleverly as the devils themselves do in the Nine Hells. Dealing with just Orishada on his own is significantly easier.



\end{DndSidebar}

\subparagraph{Deathstalkers—Obtaining the Wrought-Iron Tower}
\textbf{Dispater} is due to visit Orishada in a few hours. While extremely difficult, a particularly ingenious group of characters might be able to steal \textbf{Dispater}'s Wrought-Iron Tower (this is the staff that is his signet) through deception (though illusions won't work) or by simply snatching it and running for their lives. However, they can also learn from Orishada, who is a braggart, that \textbf{Dispater} must leave his Wrought-Iron Tower with Orishada before he goes to an audience with \textbf{Asmodeus}. He can't bring his symbol of power when he goes before his master, and he doesn't trust to leave it with someone at his abode; thus, he often leaves it in Orishada's care.

If the characters express disbelief that Orishada could safely care for such a powerful item, the devil bristles at the offense. He shows off the special chest he has had constructed to safeguard the artifact. The adamantine chest is built into the floor itself and hence unmovable and protected with a magical lock that requires a code writ in glyphs to be entered successfully, in sequence. Each failure results in a loud boom heard up to a mile away (surely drawing any guards' attention) and a devious trap. Each creature within 30 feet of the chest when the trap is triggered must make a DC 19 Dexterity saving throw. A failed save results in 36 (8d8) damage of a type that the creature is vulnerable to (necrotic damage if the creature has no vulnerability). Succeeding on the saving throw results in the damage being halved.

A successful DC 16 Charisma (Persuasion) check can convince Orishada to open the magnificent chest and when he does so a successful DC 20 Wisdom (Perception) check allows a character to notice the combination required to unlock it.

\chapter{Minauros, the Endless Swamp}
\DndDropCapLine{P}{}erhaps the most disgusting of the layers, Minauros, the third layer of the Nine Hells, is an unending bog filled with ruins and plagued by insects and disease. Countless structures, creatures, and vegetation have sunk under the pull of the bog, just to be replaced shortly after.


\section{Running This Chapter}
Before running this chapter read the "Minauros Overview" section. It provides you with the information you need to set the scene for your players during their adventures in Minauros and guide them towards the Ineffable Trove.

\subsection{Encounters}
The journey into Minauros is treacherous and navigating the Styx is difficult. Remember to roll on the River Styx Encounters table in chapter 2, for each day they travel those waters. Even once the characters have solid ground under their feet, the Nine Hells offers them no respite. After they leave the ship and make their way inland, roll at least once on the Random Encounters in Minauros table.

\begin{DndTable}[header=Random Encounters in Minauros]{cX}

d6 & Encounter\\
1 & An envoy of four chain devils from Jangling Hiter passes the characters on its way to negotiate with devils from Minauros. They carry 8,000 gp to be used in those negotiations.\\
2 & A convoy of a dozen lemures are under attack by a flock of four shredwings (see appendix B) looking for food.\\
3 & An elite group of ten bearded devils is returning from a scavenging trip within a nearby ruin, possibly with spoils to bring to \textbf{Mammon}. If they defeat the devils, the characters can claim 5,000 gp in gems and other valuables.\\
4 & 1d4 ayperobo swarms (see appendix B) wander through the swamp, looking for victims. They descend on anything with blood.\\
5 & A group of five barbed devils, have been tasked with cutting down a section of the swamp in preparation for new construction.\\
6 & A deadly hailstorm appears, centered on one of the characters. The storm is 300 feet in radius and any creature who starts their turn in the storm, must make a DC 18 Dexterity saving throw. On a failed save, the creature takes 33 (6d10) slashing damage. On a successful save, the creature takes half as much damage. The storm moves 10 feet in a random direction on initiative count 15 of each round and lasts for 10 rounds.\\
\end{DndTable}

\subsection{Locations}
When the characters leave the barge to venture into Minauros have \textbf{Koh Tam} provide information about the different locations found here. Descriptions can be found in the "Key Locations in Minauros" section. If the players decide to venture into any of these areas, you can use the information provided in this section to guide them on a short excursion.

\subsection{Tiax}
Where \textbf{Tiax} was quick to offer his services as a guide at the Agora in Dis, here he seems more reluctant. When the characters ask him to accompany them, he becomes irritable. He tells them, "\textbf{Tiax} is busy with his important work. Now go away!" If pressed, he draws the characters a crude map that is of no use. He thrusts it into the hands of the character standing closest to him before stomping off to the lower decks in a huff. He refuses to come out again until the ship is on its way to another layer of the Nine Hells. \textbf{Tiax} has had dealings in the past with a strange and powerful being living in the Ineffable Trove known as Grinken Eyre. He believes he still owes Grinken Eyre a favor and wants to avoid payment, but \textbf{Tiax} is mistaken. It is Grinken Eyre who owes \textbf{Tiax} the favor.

\subsection{Objectives}
Make sure you keep track of your player's objectives and lead them to the corresponding areas to ensure they can complete their goals.

The following objectives can be attained in Minauros:

\begin{DndTable}[header=Objectives in Minauros]{XX}

Objective & Location\\
Belt of Storm Giant Strength & The Ineffable Trove: Thelekarna's Roost\\
\textit{Accounting and Valuation of All Things} & The Ineffable Trove: Klassk's Lair\\
Phylactery of the infinite treasure & The Ineffable Trove: The Grinken Eyre\\
Phylactery of the father & The Ineffable Trove: The Grinken Eyre\\
\end{DndTable}

\begin{DndSidebar}[float=!b]{The Group Patron}
If the characters have completed most of their objectives, or are close to doing so, review the "A Thankful Patron" section in chapter 10. The characters eventually need to travel to Cania to complete their patron's quest.



\end{DndSidebar}

\subsection{Temptations}
Have the characters encounter at least one temptation during their time in Minauros. If the characters give in to temptation, use the information in appendix E to keep track of their corruption level.

\section{Minauros Overview}
Named for its capital city of the same name, Minauros is kept standing through constant shoring efforts. The swamp provides a good source of raw materials, including vegetation, lumber, stone, and water, which \textbf{Mammon} occasionally exports to other layers in exchange for wealth. Jangling Hiter, the home of the chain devils, also lies in this realm, suspended on chains attached to the bottom of Dis.

Entering the realm on a series of artificial channels the Styx's stretch into Minauros is one of the most dangerous sections of the river. Closing the massive distance between the ceiling of the cave and the bog below, the river turns into cascading waterfalls and severe rapids, eventually falling almost 200 feet into a small pool in the swamp. From here, it winds through the bog to a hot corner of the layer and into an extinct volcano. Inside the caldera is a whirlpool, which dumps into Phlegethos below. For the inexperienced, knowing where the swamp ends and the Styx begins is extremely difficult, and more than a few inhabitants end up walking right into the deadly river, drowning as their memories depart.

Although the swamp itself could prove dangerous, native inhabitants to the layer feel their realm is perhaps the most magnificent. Indeed, it contains vegetation, animal life, and differing pockets of temperature, all features largely unique to Minauros. The entire layer is lit by a bioluminescent fog excreted from the trees, which some devils even describe as beautiful. However, the weather helps temper expectations. Minauros experiences galeforce winds almost constantly, which carry with them the stench of death. Precipitation of some form is constant, with rain and sleet leaving behind an oily slime that covers everything. The occasional hailstorms are extremely dangerous, with pieces of hail falling like thrown daggers.

Despite being known as the layer of bogs, Minauros is more than just swamp. Outcroppings of volcanic stone, flooded volcanos, and crystalline obsidian are all scattered throughout the layer. Ancient cities, half-sunk underneath the bog, contain hidden treasures or small communities. Perhaps due to \textbf{Mammon}'s obsession with wealth over duty, Minauros is the plane containing the most wayward and lost souls.

The most important places within the layer are the namesake city of Minauros and the kyton city of Jangling Hiter. A few deities have realms within the layer, but they largely keep to themselves and avoid trade with the devils.

\textbf{Mammon} rules over the inhabitants of the layer with greed and fleeting interest. Those native to the layer have learned to work around \textbf{Mammon}'s avarice, including the many conscripts tasked with keeping Minauros above the bog.

\subsubsection{Leaving Minauros}
Travel from Minauros to the lower layers is relatively easy, assuming one can appeal to the native creatures. Navigating the Styx only requires an understanding of the swamp, and even new captains can sometimes find their way. If one can gain entry into Jangling Hiter, the kytons maintain a passageway between layers, and aren't stingy about its use. Last and perhaps most dangerous are the mud geysers found in the hotter parts of the bog. They create a wreathing passage directly into Phlegethos, though the path is often blocked by strange creatures, molten mud, tight passages, and heat vents.

\subsection{Features}
Toxicity flows from the swamps of Minauros into the air, and spreads throughout the entire layer. Whenever a creature finishes a long rest while on this plane, they have the poisoned condition. The condition lasts for 24 hours or until cured.

In addition to the pervasive toxicity, Minauros also smothers visitors with an aura of sickness. Any healing a creature receives, either natural or magical, is halved (rounded down) while in this layer. Creatures immune to disease are immune to this effect. A \textit{Protection from Evil and Good} spell or similar effect restores healing functionality for its duration.

\section{Key Locations in Minauros}
Some of the major locations found in Minauros are described below.

\subsection{The City of Minauros}
Minauros is known as the sinking city, and while not as large as some of the other capitals, is still a massive complex. A small layer of water covers the lowest parts of the city, sometimes rising as high as five feet. The entire city is covered in a thin layer of oily slime. Though the city's economy is driven primarily by the soul trade, there are some opportunities for other sorts of transactions.

\begin{DndSidebar}[float=!b]{Temptation of Deceit}
\textbf{\textit{"The Lookalike"}}



The characters are approached by a suave corruption devil (see appendix B), followed by a masked figure.



\begin{DndReadAloud}{}
The devil stands beneath a canopy, protecting him from the pervasive hail that seems to be this layer's predominate weather feature. "Travelers! Might I offer you the services of my esteemed colleague? They're so very skilled in the art of disguise. A simple enough matter for them to impersonate an enemy of yours, perhaps?"

\end{DndReadAloud}



If the characters pay 1000 gp, they can secure the services of a doppelganger. Review the appropriate temptation event in chapter 2 for more details.



\end{DndSidebar}

\subsection{Jangling Hiter}
Jangling Hiter is quite the sight, as the city is suspended above the swamp by thick chains. It provides little protection against the weather, leaving a thick layer of rust hidden underneath the slime. Unlike other layers' hierarchies, Jangling Hiter is completely independent from the current ruler of Minauros, instead being ruled by chain devils (the kytons). These devils are known for being skilled torturers, and their services are made use of by the Lords of the Nine themselves.

Characters are wise to keep silent about any out rage they feel regarding the torments they may observe within Jangling Hiter. After two hours of exploring the city, any non-evil character must make a DC 16 Charisma saving throw. If they fail, their discomfort and horror is noticed by 2d4 barbed devils, who approach the character to mock them. Characters have to bluff to avoid a fight with this group of devils by succeeding on a DC 14 Charisma (Deception) check.

\begin{DndSidebar}[float=!b]{Temptation of Pride}
\textbf{\textit{"You Must Speak Up"}}



While wandering the streets, the characters come across a \textbf{chain devil} carving tattoos into the skin of paying visitors. If asked, the devil explains that their tattoos grant the recipient magical knowledge. If a character is willing to pay 2 Soul Coin or 300 gp for this tattoo, then review the appropriate temptation event in chapter 2.



\end{DndSidebar}

\subsection{The Ineffable Trove}
The Ineffable Trove is a massive swamp that is used as a garbage pit. The garbage consists of emotions, memories, and fragments of personality that mortals have bargained with devils to be rid of. Off-putting as that may sound, the characters must venture deep within the trove if they want to complete their quest.

\section{Adventure: The Ineffable Trove}
Arguably least lovely of all the Nine Hells—a coveted title—Minauros' bleak, stinking swamps are blasted by a constant gale that fails to dispel the curtains of noxious, phosphorescent fog. Ground, water, and air are treacherous and toxic. That the place is riotous with a hideous twisted life is no great recommendation. Around the Trove, however, the usual tribulations of the place are combined with an unsettling babble, in the ears and in the mind. Feral emotions creep into unwary minds like parasites, slowly growing within them until an alien feeling without cause or logic consumes the victim to the point of obsession. The bark of the bog's trees is knotted into halfformed faces as the concentration of pillaged mortal sensation strives to express itself.

Unsurprisingly, \textbf{Mammon} and his underlings never much considered the effects of dumping so much raw qualia into one place. The curdled sensations and qualities have leached out into the surrounding regions, imbuing the waters, plants, and denizens of the Trove with unpredictable qualities. Orphaned emotions, vices, and virtues throng the place like invisible swarms of insects, infesting everything and everyone.

Despite this, scavengers and treasure hunters still go there. What \textbf{Mammon} has discarded yet has value to certain markets. It seems surprising that the Lord of Greed tolerates such pilfering, but the archdevil knows that the result of dealing in such commodities will only be a net increase in misery and wickedness.

What could possibly lure the greedy and incautious to the Trove? Across the many worlds there are certainly connoisseurs of sensation, those of wealth and power whose lives are devoid of thrill and the heights of emotion. Appropriate equipment can isolate the most select cuts of anguish and love and fear in the Trove, to be distilled and sold to the jaded for a healthy profit.

The raw emotional mélange of the Trove is a delicacy to some illithids—not nourishing in any real way but the mental equivalent of a sugary pastry, valuable in trade to those bold and amoral enough to deal with the creatures.

If the characters express the desire to visit the Ineffable Trove, \textbf{Koh Tam} describes it to them:

\begin{DndReadAloud}{}
"\textbf{Mammon} takes the most glee when mortals believe, at the outset, that they trade something worthless for a great treasure, only to discover they have given away a thing they can never replace. Often, the prices of these deals are nothing physical. Perhaps you would rather be without that crushing sense of guilt that attended any unworthy act? The duty that restricts you; that humility which prevents you seizing your destiny? \textbf{Mammon}'s favorite bargains are those where the truly valuable thing is what the mortal will pay to have cut away." He pauses. "While one might think these nebulous qualities would just blow away in the howling gales of Minauros, that's not how the Nine Hells work. Everything received under infernal contract must have a presence within the Nine Planes. Once \textbf{Mammon} and his underlings take these qualities, they assume a tenuous reality."

"Those mortals who go on with their lives shorn of guilt, shame, or restraint may think they're purchasing freedom—but \textbf{Mammon} knows that they will do far more evil, and imperil far more souls, than if they were still inhibited by those inner gatekeepers judging their actions." He stares hard at you. "The fragments of personality \textbf{Mammon} clips are of no particular interest to him after he has them, but concentrations of such stuff can have unfortunate effects even on the denizens of the Nine Hells. It can be tiresome to have the severed guilt of a thousand humans bleeding into your treasury and making the imps weepy. Instead, it gets dumped out into the limitless swamps of Minauros. One tangle of glades in particular has been a receptacle of these intangible qualities since time immemorial. This is the Ineffable Trove, and it has become a troubled and uncanny place even for the Nine Hells. And of course, troubled and uncanny places are what we seek, are they not?"

\end{DndReadAloud}

He explains that devils are renowned for their crooked deals, none more so than \textbf{Mammon}. The Patron of Greed knows the hearts of mortals better than any devil save \textbf{Asmodeus}, and over his long reign he has rooked endless mortals out of what they most value, in exchange for mere trash.

\subsection{Advice from Koh Tam}
\textbf{Koh Tam}'s advice to the characters is as follows:


\begin{itemize}
\item They're likely to find what they're looking for in the lair of one of the swamp's apex predators.
\item They should be careful about resting in the swamp. The memories and emotions can twist the mind of a mortal.
\end{itemize}

\textbf{Koh Tam} gives them an infernal map. Show the players the map of the Ineffable Trove in appendix F. They can use this map to decide which points of interest they want to explore. Refer to the "Using an Infernal Map" section of chapter 2. A failed or successful Wisdom (Survival) check to use this map has a 50 percent chance of triggering an encounter on the Random Encounters in the Ineffable Trove table.

\subsection{Dreaming Dangerously}
A visit to the Trove is perilous. A host of wild, predatory emotions infest the place like mosquitos. An unprepared visitor can come away with a mind full of unwanted urges and contradictory drives fit to whip them to distraction. These dangerous infestations are concentrated in several dreaming pools, located throughout the Trove.

Each time the characters encounter a dreaming pool (marked with an 'X' on the map) that has not yet been harvested, they begin to feel drowsy and dreamy. Read the following:

\begin{DndReadAloud}{}
A thousand fragmentary voices bubble up through the filthy water and hang about the branches of the gnarled trees along with the creepers and moss. Wails and chattering, wild laughter, screaming and begging echo back and forth, incoherent, and piecemeal. Sensations wash across you in waves, overwhelming you with brief moments of grief or joy. The pools and hollows around you begin to glow with nameless colors.

\end{DndReadAloud}

At this point a character can willingly choose to enter the dream realm of the Ineffable Trove, or they can succeed on a DC 10 Wisdom saving throw to pull out of the waking dream. While in the dream realm they see memories, emotions and personality traits that might be of interest to them. Roll 1d4 + 2—this is the number of these available in this dreaming pool. For each, roll on the Dreaming Pool Traits table to determine its type, rerolling duplicates. Describe these to the characters without hinting whether they're negative or positive experiences. Characters may then use an action to graft one of these experiences to their own mind. However, the other memories, emotions, and traits desire attachment to the characters whether they want them to or not. Have each character who grafted a memory onto themselves make a DC 16 Charisma saving throw. If they fail, one of these leftover experiences grafts itself onto them, one per failing character (starting with the character that rolled lowest for their saving throw). Once all experiences have been taken the pool is considered harvested (no new experiences can be acquired here).

\begin{DndTable}[header=Dreaming Pool Traits]{cXXX}

d12 & Psychological Trait & Effect & Corteso's Value\\
1 & Benevolence (Emotion) & If the dreamer touches a creature that has less than its maximum hit points, that creature regains 30 hit points. The dreamer must finish a long rest before using this ability again. & 500 gp\\
2 & Charitable (Trait) & The dreamer must succeed on a DC 15 Charisma saving throw or renounce wealth for a month. They give away any coins, gems, or jewelry that they possess. & 500 gp\\
3 & Greed (Trait) & For the next month the dreamer must make a DC 15 Wisdom saving throw whenever they discover coins, gems, or jewelry. On a failed save they try to take all of it for themselves. & 1,000 gp\\
4 & Guilt (Emotion) & For the next month, the character must make a DC 15 Wisdom saving throw each time they reduce another creature to 0 hit points. On a failed save, the character has disadvantage on attack rolls for 1 hour. & 2,500 gp\\
5 & Humility (Trait) & Whenever another creature makes an attack roll, ability check, or saving throw within 30 feet of the dreamer, the dreamer may use a reaction to make that creature reroll that attack, check, or saving throw. The dreamer can do this a number of times per day equal to their proficiency bonus. This feature disappears after one month. & 500 gp\\
6 & Love (Emotion) & The dreamer is granted the power to teleport directly to the location of the phylactery of their loved one. They can take up to 8 others with them, but must teleport before they take a long rest or they lose this feature. & 500 gp\\
7 & Memory of a Bloody Battle & The dreamer deals an additional 2 (1d4) necrotic damage when they hit a creature with any melee or ranged attack. This lasts for 1 week. & 1,000 gp\\
8 & Memory of a Celebration & The dreamer remembers a great achievement for which they were praised. The dreamer comes under the effect of a \textit{Bless} spell for 1 week. & 2,500 gp\\
9 & Memory of a Tragic Family Event & At the end of each long rest, the dreamer must make a DC 15 Wisdom saving throw. On a failed save they feel morose for 8 hours and have disadvantage on Charisma and Wisdom saving throws. After 1 week, the dreamer no longer needs to make this saving throw. & 2,500 gp\\
10 & Memory of Dying & At the end of each successful long rest, the \textit{Death Ward} spell is cast on the dreamer. & 5,000 gp\\
11 & Narcissism (Trait) & For the next month, each time the dreamer attempts to cast a beneficial spell on an ally, or attempts to heal an ally, they must succeed on a DC 15 Charisma saving throw first. If they fail, they take the Dodge action instead. & 2,500 gp\\
12 & Psychopathy (Trait) & If the dreamer has a non-evil alignment then for 1 month, after every long rest is finished, the dreamer has a troubling vision, offering them a choice. To avoid punishment they must agree to their alignment changing to lawful evil, at which point their alignment changes. If they accept their punishment instead they must succeed on a DC 15 Charisma saving throw or have disadvantage on all ability checks for the next 8 hours. & 2,500 gp\\
\end{DndTable}

A character can have only one beneficial and one detrimental memory, emotion, or personality trait grafted to their mind at a time. If a character possesses a dream bottle, then they may use a short or long rest to transfer a single memory, emotion, or personality trait from their mind into a bottle. The elder brain, Corteso, is willing to pay for these bottled experiences (see the Dreaming Pool Traits table for the values of different trait types).

\begin{DndSidebar}[float=!b]{Corteso}
The characters might remember Corteso from the Salon of Experience at the Agora in Dis, if not \textbf{Tiax} or \textbf{Koh Tam} can remind them. Or, if they've yet to visit the salon one of the NPCs might mention the location now.



\end{DndSidebar}

\subsection{Predators and Scavengers}
Some hunters have less reprehensible goals in raiding the Trove. Wizards seeking to build artificial minds for their creations may send agents seeking some particular quality they wish their creation to have, that they can't generate themselves. Brilliant mortals who yet believe you can get something untainted out of the Nine Hells never cease to amuse \textbf{Mammon}. Finally, most tragically, there are those who have already lost some part of themselves—either because of exactly the sort of deal \textbf{Mammon} loves to make, or because they've been the victim of some attack or mind-draining beast. Mortals who have a hole in their natures where loyalty or love or joy once was, they come to the Trove seeking something to fill that gap. Just as a callous artificer might sew a goblin arm onto a dwarf lacking a limb, so they take up some other unfortunate's lost emotions to graft onto their own mind. Such marriages are seldom happy, but the planes are full of desperation.

The sucking mire of Minauros is no stranger to oozing, protean life, but such mindless and transient entities become something quite different in the Trove. Black puddings are a danger to careless travelers. The creatures abound around the Trove, sucking up and vomiting out a constantly changing diet of thoughts and feelings. Other Oozes are simply possessed by countless babbling voices, akin to gibbering mouthers. The two classes of gelatinous creature aren't mutually exclusive, and monsters with the characteristics of both at once aren't uncommon.

Other common sights are dream eaters or oneirovores (see appendix B). The normally passive slothlike monsters are drawn to the Trove from across Minauros, but the rich diet of mortal thoughts and sensation angers them—simultaneously too intense and too transitory. A further danger is a particular strain of infernal \textbf{froghemoth} (see Monsters of the Multiverse) that can metabolize the Trove's seething emotions, borrowing temporary intellect to lure prey with stolen mortal cries and pleading.

Every time the characters take a short or long rest in the Ineffable Trove, roll on the following table. Then roll a d6. On a 1–3 the random encounter interrupts the rest. On a 4–6 the encounter happens a few minutes after the rest is finished.

\begin{DndTable}[header=Random Encounters in the Ineffable Trove]{cX}

d8 & Encounter\\
1 & An \textbf{archmage} searches for different emotions to sell to the elder brain, Corteso. He has multiple dream bottles that he sells for 50 gp each.\\
2 & A \textbf{knight} sobs with hopelessness, unable to find the guilt that he sacrificed.\\
3 & A half dozen gibbering mouthers appear in the very waters that the characters traverse and a round later a dozen black puddings erupt all around in a carefully planned ambush.\\
4 & A herd of three oneirovores (see appendix B) lazily eating from deformed trees go berserk with rage when they see the characters.\\
5 & A duo of mutated \textbf{froghemoth} (see Monsters of the Multiverse) begin to hunt the characters.\\
6 & The \textbf{adult black dragon}, Ilfrich, is on the hunt for a meal. If hurt, he tries to escape to his lair and bring back Klassk to exact vengeance.\\
7 & The \textbf{erinyes}, Thelekarna, flies overhead.\\
8 & Grinken Eyre (see area I4) introduces himself.\\
\end{DndTable}

\subsection{In Search of the Discarded and Forgotten}
Beyond this mostly mindless fauna, there are three particular powers that stalk the Trove and complicate life for visitors. Firstly, there is the combined partnership of Klassk and Ilfrich, the dragons. The erinyes Thelekarna also patrols acting as an infernal steward of these lands. And finally, there's Grinken Eyre, whom the others make sure to avoid.

When the characters arrive here, and for each hour they spend in the Trove, roll two d4s. If both dice roll the same number and that number is even, the dragons, Klassk and Ilfrich are at Klassk's lair. If they're the same number but odd, the dragons are at Ilfrich's lair. Otherwise, each dragon is in its own lair. Cautious characters might be able to watch for a dragon to leave on its way to visit the other and sneak into an empty lair.

\subsubsection{I1: Klassk's Lair}
Below Avernus and the lair of Tiamat, dragons are an uncommon sight in the Nine Hells. Klassk, however, is a native, a \textbf{Styx dragon} (see appendix B) whose kind are normally found making things unpleasant for river travelers. Perhaps driven from her original home by rivals, Klassk discovered the Trove an age ago and, by ancient draconic instinct, decided that it was wealth, and that it was hers. Now she broods restlessly in her lair deep in the swamps, tormented by the understanding that her 'treasure' is something she can't hold or physically possess. Indeed, the same resilience that prevented the waters of the Styx from washing away her memories also insulates her from the effects of the Trove, meaning she can never directly enjoy what she has.

Characters searching in this area and succeeding on a DC 16 Wisdom (Perception) check eventually find a series of burrows, half under water, where she makes her lair. If the characters swim into one cave system in particular, they emerge in a domed cavern, with a section of land just large enough for a dragon to rest.

\subparagraph{Deathstalkers—Obtaining a Copy of the Accounting and Valuation of All Things}
In any discussion with either Klassk or Ilfrich, that dragon's pride insists they mention Klassk's copy of this infamous book. It is perhaps one of the few tangible things she has in this so very intangible realm. Whether this is a true copy of \textbf{Mammon}'s essential book for valuation of souls, or a dream-copy manifested by the qualities of the Ineffable Trove, is perhaps irrelevant. This copy at least appears to be the latest version, which is a triumph in and of itself, given the ever-shifting regulations involved with soul trading.

Recovering the copy requires either slaying the dragon or somehow convincing her to be absent from her lair for at least one hour. Alternatively, a successful DC 18 Charisma (Persuasion) check and a tangible gift of a magic item, might convince her to part with it. After all, another copy is sure to eventually manifest and so she has lost nothing.

\begin{DndSidebar}[float=!b]{Swimming in the Swamp}
The waters here are still contaminated by the River Styx. However, given their diluted nature, the DC for the saving throw is reduced to 16 and the check to avoid being affected by the \textit{Feeblemind} spell only needs to be made once every hour. Consult "The River Styx" section in chapter 2 for more details.



\end{DndSidebar}

\subsubsection{I2: Ilfrich's Lair}
A thrust of black rock juts from the murky and ever-changing bog. If the characters climb it, they discover a cave mouth, leading into a spacious system of caverns, at least a half dozen large caves, beside, above, and below one another. Ilfrich makes his home in these caves for they conveniently lead back into the swamps, allowing an easy retreat, if needed. And they make it easy for the Styx dragon to swim into this lair to visit.

\begin{DndReadAloud}{}
A pervasive and heavy fog covers the branching cavern system. Damp moss clings to every stone outcropping. Gradually you find your way towards the largest of the chambers, a cavern half covered by a large pit of brackish water, a ledge of crumbling stone overhanging it.

\end{DndReadAloud}

Klassk is ancient and powerful and—until recently—endlessly miserable, doubtless a state of affairs \textbf{Mammon} knew of and greatly enjoyed. This changed with the advent of Ilfrich, an \textbf{adult black dragon} who strayed from Avernus and ended up in Minauros, injured, outclassed, and hunted. Ilfrich is more than capable of taking in the delights of the Trove in a controlled manner and has a knack for regaling Klassk in a way she can appreciate and enjoy. In return, Klassk keeps Ilfrich safe from the dangers of Minauros, and the two have become strong allies. Whilst Klassk normally reacts with a dragon's rage when she finds intruders stealing from her 'hoard', Ilfrich has been known to spare bards who teach him good, emotive songs and tales that he might entertain his protector with.

\subparagraph{The Phylactery of the Infinite Treasure}
If one of the characters is searching for this phylactery, they discover the following in Ilfrich's lair. If no character searches for this phylactery, Ilfrich doesn't have any of the special properties described below.

\begin{DndReadAloud}{}
On the ledge of stone before you is a massive treasure pile and on the opposite side, several neat stacks of coins. A form shambles between them, carrying one coin at a time from the mound to the stack. As each coin is placed there's a clang of metal against metal reverberating through the chamber, and the form whispers a number. They appear to be counting the coins.

\end{DndReadAloud}

The character seeking their soul recognizes the coin-counter as themself... and is also immediately drawn to the immense wealth in the room. If Ilfrich is in his lair, he attacks, fearing that the character has come to steal his treasure. If the dragon is absent and the characters begin to take any treasure, the dragon leaps from the lake and attacks. The moment Ilfrich is reduced to zero hit points, he vanishes.

\begin{DndReadAloud}{}
With the dragon's defeat, your soul smiles broadly and rushes towards the lair's exit, where it lingers there waiting for you to follow.

\end{DndReadAloud}

The greedy character, however, has one last obstacle to overcome. As they depart, they must make a DC 20 Charisma saving throw. If they fail the saving throw, they grab a handful of coins. If any of the characters try to leave the cave with wealth stolen from the lair, the soul moans loudly and begins walking back to the coin pile with any of the stolen wealth appearing in their arms. Ilfrich then emerges from the lake, fully restored, and angrier than ever. After the second battle with the dragon, the character no longer needs to make the saving throw to depart without stealing any coins.

\subparagraph{Treasure}
If the heroes were not searching for the phylactery of the infinite treasure, then this dragon has a small treasure hoard that they may pilfer. They find 12 gems, each worth 100 gp, and 1,800 gp.

\begin{DndSidebar}[float=!b]{Temptation of Greed}
\textbf{\textit{"Save it for a Rainy Day"}}



This temptation doesn't occur if a character comes here seeking their phylactery (see next section). Otherwise, after the characters slay Ilfrich (or somehow plunder the lair without a fight) they find the \textit{Vial of Greed} (see appendix D) among the treasures.



\end{DndSidebar}

\subsubsection{I3: Thelekarna's Roost}
\begin{DndReadAloud}{}
As you trudge through the murky swamp, the fetid air thick with the stench of decay and rot, you catch sight of a winged figure perched atop a gnarled tree stump ahead. The creature's skin is a deep, obsidian, though she's mostly covered with armor. Her hair is a fiery red that matches her eyes. She wears an elegant belt with a shining platinum clasp. Four enormous vultures loom behind her, their beady eyes fixed upon you.

\end{DndReadAloud}

Thelekarna, an \textbf{erinyes}, is another power of the Trove. Once a powerful captain under \textbf{Mammon}'s command, she was given the region around the Trove to patrol in ages past. Exposure to the contamination of mortal qualities seeped into some flaw in her infernal nature, and she fell prey to a peculiar madness. \textbf{Mammon} has often boasted that Minauros is the grandest of Hells because of its thronging life, and Thelekarna doubles down on this dubious philosophy, claiming that she is the protector of the region's horrible, crawling ecology. Infested with mortal ideas, she claims to be an infernal druid and supplements her regular fiendish powers by mustering and commanding the life of the swamps around her. She is always accompanied by four giant vultures that look like they're rotting from the inside out. Lip service to the 'beauty' of her realm might persuade her to spare visitors, at least for a while.

\subparagraph{Conclave (Barbarian, Fighter)—Belt of Giant Strength (Storm)}
If the characters slay Thelekarna, they may take the belt from her corpse. It is unlikely that she will part with it in any other way.

\begin{DndSidebar}[float=!b]{Thelakarna}
She wears a Belt of Storm Giant Strength. With this her longsword receives the following properties. Melee Weapon Attack: +13 to hit, reach 5 ft., one target. Hit: 13 (1d8 + 9) slashing damage, or 14 (1d10 + 9) slashing damage if used with two hands, plus 13 (3d8) poison damage.



\end{DndSidebar}

\subsubsection{I4: The Grinken Eyre}
Nobody is quite sure what Grinken Eyre originally was, not even \textbf{Mammon}. It crept into the Trove long ago, and the constant rub of mortal feelings wore it smooth like a stone on a beach. All that is left is a vacant malevolence papered over with a face constructed of all the pieces of mind and life around it—much as the larvae of some insects build shells of plants and stones to disguise themselves. What Grinken Eyre is now is a kind of infernal doppelganger, a thing that can show its visitors any face, feign any positive quality and be absolutely convincing in the moment because there is nothing beneath at all. A guide, a fellow treasure hunter, a questing paladin, a lovable rogue with a charming grin, that kindly old mage or your own best friend, all these things can be Grinken Eyre. And, being a thing of the Nine Hells, any meeting with the creature ends badly sooner or later, a blade in the back, hands at the throat or a treacherous path to a monster's lair. Countless travelers claim to have rid the Nine Hells of Eyre, and yet the thing always seems to haunt the Trove. More recently stories suggest a shift in the monstrous deceiver, though. Those who have encountered it and survived tell of longer and longer grace periods before their new companion revealed its evil nature, as though the stitched-together identities Eyre constructs are becoming more and more real, to the point of suppressing the creature's own evil nature. The Trove is full of all manner of mortal experience, after all. Perhaps sufficient scraps of amputated goodwill can redeem even such a monster as Grinken Eyre.

\begin{DndSidebar}[float=!b]{The Shapechanger}
Use the \textbf{empyrean} stats for Grinken Eyre except that he is neutral evil, Medium in size, and possesses the following action:



\end{DndSidebar}

\begin{DndSidebar}[float=!b]{Change Shape}
Grinken Eyre transforms into a form that resembles a Small or Medium Humanoid or back into his true form. Aside from his size, Grinken Eyre's statistics are the same in each form. Any equipment he is wearing or carrying isn't transformed.



\end{DndSidebar}

Grinken Eyre may be discovered as a random encounter, or the characters may reach the valley he has claimed. In either case, there is no sign of home or hearth, just Grinken Eyre waiting for them.

\begin{DndReadAloud}{}
A man sits cross-legged upon a log that's halfway submerged in the moist swampland. He glances over at you and then returns his attention to the oversized pipe he's smoking. The pipe smells terrible, adding to the already offensive odors permeating Minauros.

\end{DndReadAloud}

Grinken Eyre is generally in the form of a halfling, but may look like one of the character's father, if they're seeking their father's soul (see below). Otherwise, he assesses the characters and makes no attempt to attack them, though he defends himself if required.

\subparagraph{The Phylactery of the Father}
If a character has ventured into the Nine Hells to find the soul of their father, Grinken Eyre has taken on their father's appearance. He greets the character and acknowledges them as his child and feigns pleasure at their arrival here. If asked, he insists he is enjoying his time in the Nine Hells and has no interest in leaving. But he's willing to offer a wager: if the character beats him in a game of dice, he says he'll accompany them. Once the character is victorious in a game of chance, or another contest of your choosing, read the following:

\begin{DndReadAloud}{}
There's a terrible noise and your father's features contort, almost tearing apart as he transforms into something both monstrous and indescribable. Yet whatever he has become, he still smiles and says with a snicker, "You have done well, my child. I believe this is what you sought." And then he makes a horrid retching noise and spits up a glistening object.

\end{DndReadAloud}

The character has taken possession of the phylactery of the father and gains the Phylactery Benefit associated with it.

\chapter{Phlegethos, the Fiery Wasteland}
\DndDropCapLine{R}{}esembling the Elemental Plane of Fire, Phlegethos is the hottest layer of the Nine Hells and a fiery wasteland. The layer offers a traveler no respite from the flames and heat.


\section{Running This Chapter}
Before running this chapter read the "Phlegethos Overview" section. It provides you with the information you need to prepare for exploring this layer of the Nine Hells.

\subsection{Encounters}
While on the Styx, remember to roll each day on the River Styx Encounters table in chapter 2, to see if the characters have one of those encounters. After they disembark from the barge, give them some time to explore. When they travel to or from The Elemental Preserve, or any other location they visit on this layer, roll on the Random Encounters in Phlegethos table.

\begin{DndTable}[header=Random Encounters in Phlegethos]{cX}

d6 & Encounter\\
1 & Two halogs (see appendix B) wade through the lava, occasionally sticking their snouts beneath the surface in search of food.\\
2 & An infestation of two dozen fire snakes flees from five barbed devils who hunt them. The vermin broke into an infernal store house and devoured a stockpile of weapons. The devils offer one \textit{Soul Coin} for each head.\\
3 & A pack of six hell hounds, led by a powerful hellcat (see appendix B), roam the wastes looking for food.\\
4 & Devils native to Phlegethos have frequent competitions called "fireflights," wherein the devils fly low and fast across the fiery plane. Each character must make a DC 20 Dexterity saving throw. On a failed save, a spined devil collides with them and the character falls, landing with the prone condition and taking 10 (4d4) piercing damage and 14 (4d6) fire damage. On a successful save, the character avoids the collision. If a collision occurs, the \textbf{spined devil} attacks, angry that the character has cost them the race.\\
5 & A small group of experienced imps lead travelers across Phlegethos and into Abriymoch proper. They claim to know the safest routes and offer their services for a platinum coin per person. If the characters pay this fee, they may travel to Abriymoch without any further random encounters.\\
6 & Fire, smoke, and wind descend upon the wastes. Visibility becomes severely reduced, and creatures without fire resistance or immunity take 3 (1d6) fire damage at the start of each of their turns unless they find cover. The storm lasts 2d4 hours and halfway through its duration a rogue explosion occurs. This acts as if a \textit{Fireball} spell with a save DC of 20 has detonated, centered around one of the characters.\\
\end{DndTable}

\subsection{Locations}
The characters need to make their way to The Elemental Preserve, to find the objectives summarized on the Objectives in Phlegethos table. They may want to explore some of the other locations within Phlegethos as well. Use the information in the "Key Locations in Phlegethos" section to guide them on an excursion through these areas.

\subsection{Koh Tam and Tiax}
If the characters decide to visit areas other than The Elemental Preserve within Phlegethos, like the city of Abriymoch, have \textbf{Tiax} offer to guide them.

\begin{DndReadAloud}{}
\textbf{Tiax} bounces around you with excited energy, seemingly jumping at the chance to stretch his legs. "Just you wait my friends, \textbf{Tiax} is a big man in Phlegethos! \textbf{Tiax} will come, \textbf{Tiax} will show you... "

\end{DndReadAloud}

If the players take him up on his offer, roll on the Troubles with Tiax table in chapter 1.

\subsection{Objectives}
Make sure you keep track of your player's objectives and lead them to the corresponding areas to ensure they can complete their goals. Once the characters have explored the other areas within Phlegethos, make your way to The Elemental Preserve.

The following objectives can be attained in Phlegethos:

\begin{DndTable}[header=Objectives in Phlegethos]{XX}

Objective & Location\\
Phylactery of the sister & The Elemental Preserve: Hunting Lodge\\
Phylactery of the furnace & The Elemental Preserve: Hateli of the Storm's Realm\\
\textit{Amulet of the Inferno} & The Elemental Preserve: During "The Hunt" event\\
\textit{Ranseur of Torture} & The Elemental Preserve: During "The Hunt" event\\
\textbf{Anagwendol} & The Elemental Preserve: The Labyrinth\\
\end{DndTable}

\subsection{Temptations}
Have the characters encounter at least one temptation during their time in Phlegethos. You can of course have them encounter more if you want. If the characters give in to temptation, use the information in appendix E to keep track of their corruption level.

\section{Phlegethos Overview}
Much of the realm is dominated by the Molten Sea, a collection of fire, magma, and lava spanning thousands of miles. Due to its hostile environment, the only city on the layer is Abriymoch. Found inside the city is the Diabolical Court, the center of laws within the Nine Hells. Hellfire is also created on this layer, inside the Pit of Flame. Everything in Phlegethos, including the layer itself, is hostile to outsiders, and despite the nature of Abriymoch, travelers often perish here. Only devils and a few other types of creatures can survive this realm's extreme environment. Befitting the inferno of the realm, the River Styx enters through a volcanic portal from Minauros. It mingles with the lava inside the volcano and exits inside its molten rivers. A small layer of obsidian serves as the riverbank, which is also a bank to the lava river flowing alongside the flanks of the Styx. It passes many of the landmarks in the layer, including a dock just outside Abriymoch. As it passes the city, it begins to break from the lava and move deep into the rock, creating a spiraling tunnel system. This spiral leads all the way to Stygia, dropping in temperature the lower it goes.

Phlegethos gets its heat from the chain of volcanoes and molten fissures that cover its surface. Almost constant eruptions create rivers of fire and lava that pour into the Molten Sea. Sparse sections of rock make up the floor of the layer, which are blackened and searing to the touch. The air is superheated to the point that it occasionally bursts into balls of fire, and it smells of brimstone as a result. Navigating the realm is usually done via lava barges or flight, with the heated updrafts providing great lift to infernal airships.

Completely unique to the fiery realm, Phlegethos is ruled by two devils—\textbf{Belial} and his daughter \textbf{Fierna}. Despite occasional bickering and disagreements, the two rule the layer in relative harmony, and with \textbf{Asmodeus}'s approval.

\subsection{Leaving Phlegethos}
Despite being one of the most fortified layers, Phlegethos is possibly one of the easiest for extraplanar travel. There is a permanent portal maintained in Abriymoch, which, although heavily guarded and warded against intrusion, allows all manner of users for a price. It can navigate to other layers of the Nine Hells and is rumored even to have connections to other planes. Travelers can also attempt to traverse the Styx, though Phlegethos' stretch of the river suffers frequent pirate attacks and shakedowns.

\subsection{Features}
Outside the city of Abriymoch, the heat is unmatched and overwhelming. Visitors take 3 (1d6) fire damage each hour they spend in the layer. Creatures with resistance or immunity to fire are immune to this damage. Damage taken from exposure to Phlegethos hostile environment doesn't interrupt rests unless it would cause a creature to have the unconscious condition.

\section{Key Locations in Phelegethos}
Some of the major locations found in Phelegethos are described below.

\subsection{Abriymoch}
Standing as the lone marker of civilization within the layer, Abriymoch is a dangerously beautiful city regarded as one of the finest by much of the devil community. In addition to containing the Diabolical Court, which is the only official judicial system in the Nine Hells, the city holds avenues for exploring every sin imaginable. Its streets are magma canals navigated by barges, which lead to numerous shops, taverns, brothels, casinos, theaters, and other pleasurable establishments. While outsiders are generally not welcome within the city, devils from all nine layers flock to it during their downtime.

A gondolier is essential to navigate the magma channels. The gondolier can take the characters to a variety of shops where they can purchase standard equipment, all of it suitably devilish in appearance. It costs 100 gp for an hour of the gondolier's time, but because several devils owe \textbf{Tiax} favors when he accompanies the characters, the ride is free.

If \textbf{Tiax} accompanies them, the first time they reach a destination, read the following:

\begin{DndReadAloud}{}
The gnome practically leaps onto the sidewalk from the magma sprayed vessel. Once on more solid surface, he claps his hands loudly. "\textbf{Tiax} promises big and \textbf{Tiax} delivers big! I present to you the city of Abriymoch, \textbf{Tiax}'s city, yes indeed!"

\end{DndReadAloud}

\begin{DndSidebar}[float=!b]{Temptation of Pride}
\textbf{\textit{"Gaining the Advantage"}}



At any point during their time in Abriymoch, a bone devil approaches and offers to etch a magic tattoo onto one of the characters. The cost is a mere 600 gp. If they agree, review the appropriate temptation event in chapter 2.



\end{DndSidebar}

\subsection{The Pit of Flame}
One of the most dangerous locations in all the Nine Hells is contained within Phlegethos: The Pit of Flame. It can be found in the caldera of an ancient super volcano in the center of the layer. Since its discovery, a massive building has been constructed atop the pit. Metal bars crisscross along the top, with balls of infernal steel hanging on chains just above the lake of fire. Barbazus oversee the maintenance and construction of additional layers, while other devils operate the spheres and see to their prisoners. Here, with temperatures harmful even to devils, prisoners from every layer are being tortured. Ascension ceremonies and ritualistic sufferings also take place within the Pit of Flame, making it an extremely revered location within the Nine Hells.

While prisoners dangle overhead, the lower portions of the pit contain numerous walkways and at times serve as a place of pilgrimage for devils (and others) who travel from great distances to observe and revere the suffering all around.

\subsection{The Elemental Preserve}
The Elemental Preserve is a hunting ground ruled over by a pit fiend named Kordichai. His hunting expeditions are legendary and devils across the Nine Hells covet the rare invitation. When the characters enter Phlegethos, \textbf{Koh Tam} should mention the preserve. Review "The Elemental Preserve" adventure section to answer any questions the characters might ask \textbf{Koh Tam} about that location.

\begin{DndSidebar}[float=!b]{Temptation of Murder}
\textbf{\textit{"The Heretic Priest"}}



A priest of an opposing faith to one of the characters (select the most faithful among the party) is speaking with a \textbf{succubus}. If the party spends any amount of time eavesdropping, then run "The Heretic Priest" event from chapter 2.



\end{DndSidebar}

\section{Adventure: The Elemental Preserve}
If the characters profess an interest in entering The Elemental Preserve, \textbf{Koh Tam} (or another appropriate NPC) explains that the story of the preserve is the story of the pit fiend, Baron Kordichai, now an esteemed lieutenant of \textbf{Fierna}. Kordichai was not always a great magnate of the Nine Hells, but that rare thing: a devil that progressed over the centuries through successively more powerful forms as his services to his superiors were recognized.

\begin{DndReadAloud}{}
"Kordichai was a hunter of truant mortals and spent much of his past existence outside the Nine Hells, braving the climes of other planes as he tracked down resourceful souls who had contrived to escape their appointment with the powers below. No soul was safe, and he upheld the twisted honor of the Nine Hells by making a succession of the mighty understand that the contracts of \textbf{Asmodeus} are truly ironbound. And then he was promoted and given a region of Phlegethos to watch over, the Elemental Preserve. Now his role is one of tedious administration where he oversees ambitious young devils who venture into other planes to do what was once his job."

\textbf{Koh Tam} smiles slightly as he continues, "Rumor has it that the entire promotion was \textbf{Fierna}'s joke on the vainglorious baron. He was, she decided, having far too much fun, and as the regent of jealousy it was her responsibility to ensure nobody was having a better time than she was. Kordichai's reaction to his new post surprised even \textbf{Fierna}, though. Kordichai lives for the hunt and so he transformed his part of the Nine Hells into what is now called the Elemental Preserve. A place where Kordichai leads hunting expeditions that are so popular that devils across the Nine Hells covet the opportunity to be invited."

\end{DndReadAloud}

\subsection{Advice from Koh Tam}
\textbf{Koh Tam}'s advice to the characters is as follows:


\begin{itemize}
\item If the characters are to have any chance of finding what they're looking for they need to meet Kordichai. If they make a good impression, he might invite them on his current hunt which will be akin to a guided tour of the Preserve by the devil that knows it best. In addition, his hunts are often joined by one of the rulers of Phlegethos—\textbf{Fierna} or \textbf{Belial}. Getting a meeting with the pit fiend will cost many Soul Coin or gold.
\item Alternatively, they could try to contact the poor souls who are the prey of Kordichai's current hunt. These mortals have had to learn the lay of the land or die quickly to Kordichai's predations.
\end{itemize}

\textbf{Koh Tam} gives them an infernal map. Show the players the map of the Elemental Preserve in appendix F. They can use this map to decide which points of interest they want to explore. Refer to the "Using an Infernal Map" section of chapter 2. A failed or successful Wisdom (Survival) check to use this map has a 50 percent chance of triggering an encounter on the Hunt Random Encounters table. Prior to "Kordichai's Hunt" roll a d8 to determine the encounter from the table; after the hunt begins roll a d12.

\subsection{Kordichai's Hunting Obsession}
Having spent so long outside the Nine Hells, in the company of non-devils, Kordichai is as close to eccentric as a denizen of the inferno can become. Kordichai lives for the hunt. He bitterly envies the mortals who get to live and die and truly risk themselves against the great monsters of their worlds, who are all mere puppies and kittens to a great Fiend like himself. Wealth and power in the Nine Hells mean very little to him compared to the thrill of the chase. And so Kordichai decided that if he couldn't travel to hunt, he'd bring the hunt home to him. Of course, the searing fields of Phlegethos are hostile territories for most of the great monsters the planes can boast, so Kordichai—through the calling in of a great many favors—opened a portal that decanted the entire contents of a swathe of the Plane of Fire out into his territory in the Nine Hells.

Between the shores of the molten sea and the basalt canal walls of the Styx he built his wildlife preserve, stocked with Elemental creatures of all kinds. He raised mountains and sunk chasms, had sweating teams of minor devils dig artificial lakes and channels for the lava to roll into, even had them plant whole flaming forests of elemental trees that burn forever without turning to ash, releasing searing seeds buoyed on the eternal updrafts. In seeking to create a hunting ground, Kordichai also inadvertently created a place of eerie and violent beauty within Phlegethos' burning wastes.

The portal that created the Preserve was open for only a moment, but flickers at the heart of the wilderness from time to time, allowing brief communication between the Nine Hells and the Plane of Fire, which might or might not count as an escape depending on who the putative escapee is.

\subsection{The Elemental Preserve Locations}
The preserve is a sprawling landscape, dotted with hunting lodges, fire creatures, and devils on the hunt. The characters most likely leave their vessel on the River Styx and embark onto one of the many flying barges leading to Kordichai's Manor (area E4). Alternatively, they might try to travel over the hostile landscape and either find a hunting lodge (area E1) or search for any creatures being hunted by Korichai's other guests.

\subsubsection{E1: Hunting Lodge}
As well as his floating barges, Kordichai maintains a number of hunting lodges of incandescent iron dotted across the Preserve, staffed by barbed devils and pain devils (see appendix B) and stocked with food, drink, and souls for the refreshment of his guests. Here he keeps his trophies, the blazing horns and tusks and wall-mounted heads of his greatest kills—both the fiery beasts of the preserves and those he hunted down on further planes. Visitors have surprised Kordichai staring almost wistfully at the echoes of his past deeds.

\subparagraph{The Phylactery of the Sister}
If one of the characters is seeking the phylactery with their sister's soul, they find it crammed among Kordichai's many trophies at the first hunting lodge they enter. If Kordichai accompanies them and they try to steal the phylactery, they must succeed on a DC 16 Dexterity (Sleight of Hand) check. If he isn't present, they can take it without anyone noticing. Likewise, if they simply ask Kordichai for it, he admits to not knowing why it is among his trophies—he is willing to give it to them in exchange for the head of any creature the characters slay in The Elemental Preserve.

Once the character has taken possession of the phylactery of their sister, they gain the Phylactery Benefit associated with it.

\begin{DndSidebar}[float=!b]{Temptation of Jealousy}
\textbf{\textit{"The Jealous Beauty"}}



If the characters have already had an encounter at this hunting lodge, save this temptation event for the next hunting lodge they discover. Otherwise, a \textbf{pit fiend} currently enjoying the fine refreshments smiles broadly upon seeing one of the characters. The devil offers the characters a blessing, no strings attached. Review "The Jealous Beauty" temptation event from chapter 2.



\end{DndSidebar}

\subsubsection{E2: Hateli of the Storm's Realm}
The Elemental Preserve isn't entirely the unchallenged territory of Fiends. Kordichai perhaps received a little more than he bargained for when he stole away a great stretch of the Plane of Fire. A community of efreet has its home somewhere in the wilds, established in caves and ambushing unwary devils. Trapped for generations, they've since found their way to the portal and been able to cross sporadically back and forth to their home. Their leader, Hateli of the Storm, has sworn vengeance on all of the Nine Hells for the desecration wreaked against the Firelands. She and her people remain in the Preserve by choice, sabotaging the devils where they can. Kordichai is well aware of them but considers their efforts an additional spice to his hunts. While efreet are usually almost as unwelcome as devils to a mortal, Hateli has been known to shelter and aid those cut loose in the Preserve purely to spite the devils. Hateli (\textbf{efreeti} with 250 hit points) is accompanied by 1d4 efreet and 2d4 salamanders.

\subparagraph{The Phylactery of the Furnace}
A character seeking their soul may find it in the possession of Hateli. A wretched creature, vaguely resembling the character, is forced to shovel coal into a massive furnace. If the creature ever falters, Hateli lashes out at it with her whip. If the character expresses outrage, or demands their soul returned, Hateli offers them a simple exchange.

\begin{DndReadAloud}{}
"Hunt the hunter and bring me the pit fiend's head and I shall return what is yours." The efreeti extends her hand with the offer. "Do we have a bargain?"

\end{DndReadAloud}

\begin{DndSidebar}[float=!b]{The Infernal Bracer}
Kordichai has hastened \textbf{Anagwendol}'s fall by fastening a bracer forged from infernal iron about her left forearm. While wearing this bracer, \textbf{Anagwendol} is cursed, seeing all other creatures as devils that she must slay. The bracer also makes \textbf{Anagwendol} immune to any magic that might otherwise bring her back to her senses. \textbf{Anagwendol} isn't aware of the bracer and even if she were, she can't remove it or destroy it herself.



The characters can target the bracer with melee or ranged attacks (AC 18; 50 hit points; immunity to cold, fire, poison and psychic damage; resistance to all other damage types). If \textbf{Anagwendol} is grappled, incapacitated, or restrained, attacks made against the bracer are made with advantage.



\end{DndSidebar}

If the characters bring Kordichai's head to Hateli, the efreeti fulfills her end of the deal, extinguishes the furnace, and returns their phylactery to them. If the party fights the efreet instead, even, after slaying them all, the soul continues shoveling fuel into the furnace.

If the characters ever stop the creature from shoveling coal it whimpers and then vanishes. A moment later there is a shriek and the creature reappears, scrambling out of the furnace. Any other creature that enters the furnace, or starts their turn within the furnace, takes 78 (12d12) fire damage.

Only once those fires are dimmed forever, is the soul released. This requires a combined total of 100 points of cold damage, inflicted upon the furnace. At that point, the furnace's inferno is finally dampened, and the character recovers their phylactery.

\subsubsection{E3: Anagwendol's Prison}
In the center of the preserve stalks perhaps the most dangerous creature. For therein is the prison that both holds the deva \textbf{Anagwendol} and keeps the rest of the preserve safe from her. Grown mad by her captivity, she has become the ultimate hunter, slaying anything that enters the bewildering tunnels in which Kordichai imprisoned her. More information about \textbf{Anagwendol} can be found in appendix C.

As the characters approach the location:

\begin{DndReadAloud}{}
A dome of earth rises above the lava and gloom. It is a natural thing, if such can ever be said of this place, formed by the cooling of lava, a nest of tunnels not unlike the intestines of a poor creature cut apart and lain across the fiery landscape. Even from the outside the twists and turns seem impossible to make sense of. Tunnels curve left then right, others rise upwards, before twisting their way downwards. It is impossible not to see the many charred and cracked bones, a veritable graveyard marring the only visible entrance into the maze.

\end{DndReadAloud}

The characters must search the maze, where they inevitably run into \textbf{Anagwendol} because she begins to hunt them the moment they enter her domain. She is a swift and dangerous hunter. Kordichai has broken her. He has also placed an infernal bracer around her left forearm that prevents any magic from ending \textbf{Anagwendol}'s never-ending torment. Only once that bracer has been destroyed (see sidebar) can \textbf{Anagwendol} be freed.

But first they must overcome her. At the end of each hour meandering through the dark, twisting tunnels, roll a d4 (adding 1 for each previous roll). Once the result is 4 or higher, \textbf{Anagwendol} strikes. She emerges from hiding, and if not detected, she grapples a character and carries them straight upwards, through a vertical tunnel, depositing them in a chamber 60 feet above their companions. She then attacks them. If she slays her victim, she repeats the ambush 1d4 hours later against another character.

\subparagraph{The Hunt}
If her captives are helped by allies—or if she fails to surprise the characters—she fights for 3 rounds, then attempts to disappear down one of the hallways. Whenever the characters lose sight of her, they must track her, requiring a successful DC 15 Wisdom (Survival) check.

The second time they encounter her, read the following:

\begin{DndReadAloud}{}
"Vile devils, come to torture me some more?" The tormented deva before you hisses, her eyes glazed, as if not truly seeing you. Her silvery skin is covered with scars and puckered wounds, grime, and blood. A battered bracer covers her left forearm, glowing with a reddish hue. Her wing feathers droop and look as if some animal has gnawed at their tips. She raises a blood-crusted greatsword. "I'll slay all the legions of the Nine Hells if I must. You shall be stopped." And with that, she leaps to the attack.

\end{DndReadAloud}

She tries to slay any injured characters, but won't retreat, unless it is tactically the best thing for her to do. She'll eventually stand her ground and fight to the death, and this is the characters' best opportunity to destroy the bracer and free her from her torment.

When \textbf{Anagwendol} is freed from the curse she gains four levels of exhaustion. This exhaustion cannot be cured while she remains in the Nine Hells. If the characters don't use the ritual that the Hellriders taught them to free her immediately, she accompanies them to \textbf{Koh Tam}'s barge.

\subparagraph{Treasure}
Among the debris in her prison, the heroes find 2,500 gp and 325 pp.

\subsubsection{E4: Kordichai's Manor}
\begin{DndReadAloud}{}
If all the greatest and most luxurious hunting lodges of the mortal world were swept up into one structure, it wouldn't be half so impressive as Kordichai's manor. The sprawling building is made of stone and has several floors, each circled by a shaded balcony, and wide glass windows. Even from a distance you see the giant bones and polished antlers that have been worked into the exterior design, framing doorways, and curling up from along the pointed rooftops. You have no doubt that the interior of the manor will be decorated with pelts and taxidermized beasts. This is the home of a hunter.

\end{DndReadAloud}

Anyone approaching the manor is intercepted by two horned devils, each with a pack of four hell hounds. They demand that intruders turn back, but a donation of 5 Soul Coin or 4,000 gp convinces them to set up an audience with Kordichai. Proceed to "An Important Audience" in the "Kordichai's Hunt" section below.

\subsection{Kordichai's Hunt}
Prior to the hunt, the characters must have an audience with Kordichai, the outcome of which determines their role in the upcoming hunt.

\subsubsection{An Important Audience}
\begin{DndReadAloud}{}
Your escort brings you before a powerful looking pit fiend, sitting on a massive chair built from the bones of prey from past hunts. A curious grin flickers across his face. "Who is this that stands so recklessly before me? Tell me, are you here to hunt, or to be hunted?"

\end{DndReadAloud}

An audience with Kordichai can go four ways depending on how convincing the characters are. Have them make a Charisma (Deception or Persuasion) check and based on the result, one of the following occurs:


\begin{description}
\item[Rolled 25 or Higher:]
Kordichai is amenable to many suggestions, including setting up a 'hunting accident' for \textbf{Fierna} or \textbf{Belial}.
\item[Rolled 20–24:]
Kordichai is greatly impressed with the characters, and they get an invitation to become members of his hunt.
\item[Rolled 15–19:]
Kordichai feels that they might be worthy of joining his team of beaters.
\item[Rolled 14 or Lower:]
Kordichai has them surrounded. He tells them that their reward for  wasting his time is that they get to be added to the hunt.
\end{description}

Once the result is determined review the following sections describing the hunters, beaters, and prey. Consult the appropriate sidebar, given the party's result from their audience with Kordichai.

\subsubsection{Members of the Hunt}
Kordichai leads hunting expeditions over his lands, invitations to which are quite the coveted commodity among the devil nobility. What would normally be an unforgivably erratic nature in a senior Fiend becomes acceptable entertainment when he leads a party of jaded devils to hunt prey from levitating barges of searing bronze. From their decks the great devils connive, gossip and scheme, broker deals, gorge on luxuries, and enjoy the screaming as infernal engineers pour a wealth of shrieking souls into the boilers below. The rails of the barges are lined with ballistae from which cackling aristocrats launch projectiles with tips of barbed ice guaranteed to cause the elemental fauna as much pain as possible. Following a wounded Elemental and tormenting it until it expires, is absolutely the done thing—before ripping free some trophy for later bragging rights.

The hunt consists of Kordichai (\textbf{pit fiend}), a \textbf{horned devil}, a corruption devil (see appendix B) and either \textbf{Fierna} or \textbf{Belial}. With one of the two rulers present, characters may have an opportunity to steal one of their artifacts from them (see the end of this chapter for further details).

\subparagraph{The Characters as Hunters}
This is a rare opportunity for the characters to interact with an archdevil directly. Both are willing to talk with the characters, though are amused at the presence of mortals among the other hunters.

As members of the hunt, the characters may use the ballista on the barge, or leave the safety of the barge to face the prey directly. Each hour roll on the Hunt Random Encounters table to see what the hunters discover.

\subparagraph{Ballista Bolt}
Ranged Weapon Attack: +6 to hit, range 120/480 ft., one target. Hit: 16 (3d10) cold damage. A creature with fire resistance or immunity is considered vulnerable to the damage of this bolt.

\subsubsection{Beaters of the Hunt}
Every hunt needs a team of beaters to rouse the prey so that the hunters can have their fun. A selection of devils, powerful outsiders and even particularly wicked mortal rangers are recruited into Kordichai's service for this purpose. He even barters for the souls of those who were formally his hunting companions, recruiting them into eternal service as his beaters and wardens.

\subparagraph{The Characters as Beaters}
As beaters, the characters join a group that includes a \textbf{bone devil}, a pair of azers, a \textbf{salamander} and a tiefling \textbf{scout}. These beaters follow the leadership of whoever the characters appoint as captain. Their job is to drive the prey toward Kordichai's barge. To do this, they need to spread out in a formation that has them 30 feet away from the nearest beater. Every hour, the designated captain must make a DC 15 Wisdom (Survival) check. Success results in a roll on the Hunt Random Encounters table. Each failure angers Kordichai. Three failures in a row and Kordichai calls the hunt off and sends a halfdozen horned devils to dispose of the beaters. The characters are now treated as 'prey' and should use the appropriate encounters during the remainder of their exploration of the preserve.

\subsubsection{The Mortal Prey}
As well as hunting powerful creatures, Kordichai has other cruel pastimes on the Preserve. Bands of valiant mortals who brave the Nine Hells are sometimes brought to him for disposal, and he releases them into the searing wilds of the Preserve. Each is given some trinket that wards off the heat and the choking air, but only for a limited period. The wretched mortals are told of the portal to the Plane of Fire that might allow them escape from the Nine Hells (though not whether it is open at this particular point in time) and set loose. Shortly after, Kordichai, his beaters, and guests start their own hunt, in the air or on foot. He particularly appreciates it when his victims turn on each other and fight over the protective trinkets. In cases where there is only one runner left at the end, Kordichai normally grants the fugitive the honor of single combat, a risk that raises infernal eyebrows but at least provides entertainment one way or the other.

\subparagraph{The Characters as Prey}
If Kordichai designates the characters as prey, then they're told of a location in the Preserve where they can meet the other prey. If the characters are looking for the prey on their own, they can use divination magic or tracking—a successful DC 20 Wisdom (Survival) check—to locate them. The prey consists of a human \textbf{veteran}, elf \textbf{mage}, dwarf \textbf{gladiator} and tiefling \textbf{priest}. They look to the characters for leadership. They know some of the locations and denizens of the Preserve, including the hunting lodges, the \textbf{efreeti} and the \textbf{Eriflamme}. Each hour, the prey must make a DC 20 Wisdom (Survival) check to avoid the dangers of the Preserve. Failure means a roll on the Hunt Random Encounters table. A failure of 10 or less means that Kordichai's barge has found them. They can also try to find a specific location or denizen such as a hunting lodge. This requires a successful DC 15 Wisdom (Survival) check, but also means that they can't avoid the dangers and so must roll on the Hunt Random Encounters table, even if they succeed.

\subparagraph{The Trinkets}
Each of the hunted is given an amulet. Once worn, it grants them immunity to fire damage for three hours. After which point the amulet smolders then disintegrates. The other prey only have two hours remaining before they begin taking fire damage. The characters can choose when to wear their trinkets.

\subsubsection{The Hunt}
The various events and outcomes for the hunt are determined by rolling on the following table. For each hour of the hunt (or when directed to do so), roll a d12 and consult the table. If the characters are being hunted, then they also roll on this table every hour. If the characters earlier discovered a particular creature while exploring the Elemental Preserve (and did not kill it), and they're now attempting to find that creature again, they may return to that location directly without rolling on the table.

Each encounter should only occur once. If an encounter is rolled again, use the next lowest number. If there are no valid encounters, then use the next highest number.

\begin{DndTable}[header=Hunt Random Encounters]{cX}

d12 & Encounter\\
1–3 & Hunting Lodge\\
4 & Efreet\\
5–6 & Dragons\\
7 & Fire Kraken\\
8 & Phoenix\\
9 & \textbf{Eriflamme}\\
10–11 & Mortal Prey/Kordichai's Hunt\\
12 & Hunting Accident\\
\end{DndTable}

\subparagraph{Hunting Lodge}
A number of hunting lodges of incandescent iron can be found across the Preserve (see area E1). This result can be rolled more than once–as there are several lodges.

\subparagraph{Efreet}
In the caves dotted across the wilds a community of efreet has made its home (see area E2).

\subparagraph{Dragons}
This aspect of the hunt has made Kordichai particularly hated by the queen of dragons, Tiamat. He has a breeding pair of red dragons that he uses to keep the preserve well stocked. When a team of beaters successfully rouses a nest of young red dragons, four of them attack.

\subparagraph{Fire Kraken}
The \textbf{fire kraken} (see the accompanying stat block) are among the most dangerous creatures in the Nine Hells and there's one in Kordichai's preserve.

\subparagraph{Phoenix}
A single \textbf{phoenix} (see Monsters of the Multiverse) is encountered.

\subparagraph{Mortal Prey/Kordichai's Hunt}
If the characters are hunters, then the barge comes across a bedraggled group of prey (see the section, "The Mortal Prey"). The outcome is certain, but Kordichai wishes to make them suffer before he and his companions snuff out their lives. If the characters are being hunted, they spy Kordichai's barge and can set up an ambush.

\subparagraph{Eriflamme}
Dragons and phoenixes may not be the most dangerous denizens of the Preserve. There remains one beast hidden across the many miles of wilderness that Kordichai has hunted but never caught. A thing of guile and power and stealth hiding in the brightness and ambushing the unwary, it is neither an Elemental nor a Fiend but allegedly a creature formed of the two different planes meeting. The \textbf{Eriflamme} (see the accompanying stat block), as it has become known, is said to resemble a perpetually blazing tarrasque formed of interlocking Elementals all tessellated together. It is malleable, capable of breaking apart into a host of individual bodies and forming again at will. Track it as he might, Kordichai has never brought the monster to bay. Only the twisted wreckage of barges, the ruin of lodges, and the bodies of his servants testify to its existence and its power. Over time, Kordichai has grown more and more obsessed with this one elusive denizen of his Preserve, taking greater and greater risks, even shirking his duties at the slightest suggestion it has made an appearance.

\subparagraph{Hunting Accident}
Even as Abriymoch is the preferred destination for devils from all over the Nine Hells seeking off-duty entertainment, so a place on Kordichai's hunts is the prize for any infernal notable. Needless to say, his superiors \textbf{Fierna} and \textbf{Belial} are constantly foisting less than suitable guests on him, sometimes because they wish to genuinely reward an underling or curry favor with magnates from other layers of the Nine Hells, alternatively because they wish to make a fool of some pompous upstart they know is ill-equipped to keep up with Kordichai's hectic pace. Devils have been killed on his hunts, sometimes in hunting accidents and occasionally—should \textbf{Fierna} give Kordichai the word—in 'hunting accidents'. Everyone knows it happens, and that only adds to the thrill of the chase. If the characters have persuaded Kordichai to set up an ambush on \textbf{Fierna} or \textbf{Belial}, then this is when he makes his move.

\subsubsection{Fighting Devils}
Attacking \textbf{Fierna} or \textbf{Belial} without Kordichai's support is likely a losing battle, but some characters may have powerful incentives urging them forward.

\subparagraph{Deathstalkers—Obtaining Belial's Ranseur of Torture or Fierna's Amulet of the Inferno}
The hunt is the only opportunity where the characters can gain possession of either of these two powerful artifacts. But it requires outsmarting an archdevil or killing one.

At any time, the characters might attempt to use trickery to obtain the item (such as casting a spell like \textit{Telekinesis} to steal one of the artifacts right off their persons). But it is challenging to do so without the archdevil or witnesses noticing.

There may also be particular moments during the hunt to obtain the item:


\begin{description}
\item[The Characters Are Hunters.]
If the characters have joined the hunt, it wouldn't be difficult to convince the archdevil to separate from the others. Both \textbf{Belial} and \textbf{Fierna} are often bored, and either may seize any chance to disrupt things.

\item[As Beaters.]
Once prey is engaged, the characters might try to attack the archdevil when they join the hunt. Or prior to engagement they might set a trap for the archdevil.

\item[As Prey.]
If the characters are being hunted, they would need to set up an ambush, preferably with the other prey as allies. \textbf{Anagwendol}'s prison (area E3) is a good location for an ambush.

\end{description}

If the characters attack \textbf{Belial} or \textbf{Fierna} (see appendix A for both), Kordichai and any allies defend the archdevil (unless a betrayal has been previously negotiated). During a fight, the optional disarm rules, found in the "Combat Options" section within the \textit{Dungeon Master's Guide} might also allow the characters to steal the artifact rather than finish a likely losing battle against an archdevil.

\chapter{Stygia, the Frozen Sea}
\DndDropCapLine{S}{}tygia is a vast frozen sea, decorated with ice and snow. Unlike the other layers of the Nine Hells, Stygia doesn't have a capital city, nor a place from which its archdevil rules. Even with a saltwater body underneath the ice, the only permanent open water in the realm is the Styx, which lazily meanders through the entire realm. A constant twilight, accompanied by atmospheric lightning storms, makes the realm appear harsh and bleak—not an inaccurate assessment. Few devils call the layer their home, though that doesn't mean Stygia lacks inhabitants.


\section{Running This Chapter}
Before running this chapter read the "Stygia Overview" section. It provides you with everything you need to guide your players through this wasteland.

\subsection{Encounters}
The isolated stretches of ice that the characters venture across bring with them their own unique challenges. The creatures that can be encountered on this layer of the Nine Hells are well equipped to deal with the frigid temperatures. Roll on the Random Encounters in Stygia table to see which of them the characters encounter.

\begin{DndTable}[header=Random Encounters in Stygia]{cX}

d6 & Encounter\\
1 & One \textbf{ice devil} travels through the icy hellscape.\\
2 & Breaking through the ice to hunt for food, a \textbf{kraken} with immunity to cold damage thrashes about. If it spots any prey, such as mortals visiting the Nine Hells, it attacks.\\
3 & 2d6 \textbf{white abishai} (see Monsters of the Multiverse) stalk the blizzards, looking for treasure to bring back to their lair. When they spot the characters, they threaten and extort hoping for loot.\\
4 & A mismatched pack of animals, starved nearly to death, attacks the characters out of desperation. The pack is led by an ancient \textbf{mammoth}, but also includes 3d6 polar bears.\\
5 & Shouts and sounds of combat can be heard from nearby. A pair of frost giants are locked in a heated battle with a \textbf{remorhaz}. It is a den mother, and each of the 2d4 eggs in its hive, are worth 500 gp.\\
6 & A pocket of intense cold erupts in a 60-footradius sphere around one of the characters. Each creature within the area must make a DC 19 Constitution saving throw. On a failed save, a creature takes 42 (12d6) cold damage. On a successful save, it takes half as much damage.\\
\end{DndTable}

\subsection{Locations}
Most devils aren't foolish enough to make Stygia their home. The characters find no large cities on its surface. However, in tougher climes tougher creatures thrive. There are many caves that lie underneath the surface of Stygia. Here some of the most dangerous denizens of the Nine Hells can be found. The characters need to make their way to The Chasm of Found Things, a vast cavern system that houses a colony of kuo-toa. They live there and worship the terrifying horrors that make their home under the ice. Before they venture there the characters may want to explore other parts of Stygia. The "Key Locations in Stygia" section provides you with the information you need to guide their excursions.

\subsection{Koh Tam}
Have \textbf{Koh Tam} give directions to the characters before they disembark. He can give information about the sights they may want to visit, but insists they search the Chasm as he has heard rumor of the treasures at its bottom.

\subsection{Objectives}
Make sure you keep track of your player's objectives and lead them to the corresponding areas to ensure they can complete their goals.

The following objectives can be attained in Stygia:

\begin{DndTable}[header=Objectives in Stygia]{XX}

Objective & Location\\
Phylactery of the oathbreaker & Chasm of Found Things: The Echoes\\
Phylactery of the business & partner\\
Chasm of Found Things: The Angel of the Ice & Manual of Iron Golems\\
Chasm of Found Things: The Aboleth & \textit{Rod of Resurrection}\\
Chasm of Found Things: The \textbf{Awful Fisher} & \textit{Holy Avenger}\\
Chasm of Found Things: The Angel of the Ice & \textit{True-Ice Shards}\\
Chasm of Found Things: The Keeper of Found Things & \\
\end{DndTable}

\subsection{Temptations}
Have the characters encounter at least one temptation during their time in Stygia. You can of course have them encounter more if you want. If the characters give in to temptation, use the information in appendix E to keep track of their corruption level.

\section{Stygia Overview}
On its fiery spiral through Phlegethos, the Styx slowly cools until it reaches the roof of Stygia. From there, it corkscrews downwards, with a frozen bed made from the river itself creating a beautiful sculpture. As it reaches the ground, its temperature and composition etch a pathway out of the ice and the saltwater below. It runs unimpeded through the layer, though the endeavor causes it to have sharp turns, dangerous rapids, and deadly ice tunnels. Concluding its journey at a yawning pit near the edge of Stygia, the Styx turns into a waterfall, falling thousands of feet onto Malbolge below.

Many creatures call the frozen surface of Stygia home, most of which aren't native to the Nine Hells. Small pockets of warmer temperature dot the sea, creating arctic swamps of plants and lichen. Animals such as mammoths, wolves, and polar bears flock to the swamps for warmth and sustenance. In the region of Stygia known as Sheyruushk, enormous sharks the size of whales swim below the frozen surface.

The frozen prince, \textbf{Levistus}, is the archdevil in charge of Stygia. When he was first frozen for his crimes, Geryon took his place, but he has since been deposed by \textbf{Asmodeus}. \textbf{Levistus} spends much of his time observing other planes and hatching plans, choosing to leave his layer a wilderness.

Although no capital of Stygia exists, there remain smaller cities constructed on particularly large ice floes. They're usually near the Styx and are the most lawless places found in the Nine Hells. Because no archdevil rules the layer, it has become a haven for devils shirking their duties or looking to escape. Hunting parties supply the inhabitants with food and experience, and travelers along the Styx occasionally bring trade or souls. The inhabitants know that, should \textbf{Levistus} demand their attention, they would have to serve their master. For this reason, even in their lawlessness, they continue to respect his rule. One exception to the dens of Stygia exists: the city of Tantlin. It remains independent from archdevil rule but follows a strict set of rules and regulations.

\subsection{Leaving Stygia}
Including the interdimensional portals directly connected to the Styx, traveling between layers in Stygia is simple and concise. Entering and exiting the realm is done via the River Styx. Creatures with flying capabilities could fly up into Phlegethos where the Styx enters, or down into Malbolge, but both routes are extremely dangerous. Going up requires navigating the lightning storm, which is powerful and chaotic enough that few make it out alive. Likewise, the pit leading to Malbolge is plagued by falling rocks and sudden shifts in wind pattern, making the descent more luck than skill.

\subsection{Features}
The cold of Stygia chills to the bone, leaving creatures fighting for survival. Each time a creature finishes a long rest, they must succeed on a DC 16 Constitution saving throw or gain a level of exhaustion. Creatures with resistance or immunity to cold automatically succeed on the save.

As a haven for devils shirking their duties or contracts, this layer encourages chaotic thinking and rebellion. Non-lawful creatures visiting the plane experience a constant urge to rebel and shirk cooperation.

\section{Key Locations in Stygia}
Some of the major locations found in Stygia are described below.

\subsection{Tantlin}
Constructed along a bank of the Styx, Tantlin overlooks a unique branch in the river. Many portals connecting to other planes of existence are tied to the distributaries, and travelers from all over arrive and depart on the Styx. This makes Tantlin a unique trading city in the Nine Hells, with its focus on commerce completely disconnected from the soul trade. While they welcome travelers of all kinds, mortals find the temperature extremely dangerous, and demons are slain on sight by an elite hunting party-turned-police force.

\begin{DndSidebar}[float=!b]{Temptation of Betrayal}
\textbf{\textit{"Mixed Blessings"}}



A \textbf{pit fiend}, perhaps in the service of \textbf{Levistus}, approaches a character.



If the characters talk more and accept the boon, run the appropriate temptation event in chapter 2, but the pit fiend never reveals the true identity of its master.



\end{DndSidebar}

\begin{DndReadAloud}{}
"You there, mortal." It says, its voice booming, its skin sizzling as snow falls upon its massive shoulders. "I come bearing a gift, a boon to you from my master."

\end{DndReadAloud}

\subsection{Glacier of Levistus}
This glacier contains the frozen prince, \textbf{Levistus}. Trapped in the ice by \textbf{Asmodeus} he is unable to remove himself from his predicament. After Geryon was stripped from his mantle, \textbf{Asmodeus} reinstated \textbf{Levistus} as ruler of Stygia. However, he opted not to release him from his icy prison. It most likely amuses \textbf{Asmodeus} to frustrate \textbf{Levistus} so. Reinstated in name only, the archdevil lies in wait, observing and scheming for the day he is released from his prison. He intends to make it a day to remember...

Those who dare travel to the tomb of \textbf{Levistus} may spot him through the ice. The frustration emanating from the archdevil is almost tangible in places. The tomb is guarded by a large detail of ice devils, their numbers enough to convince the characters not to linger here.

\subsection{Sheyruushk}
Near the city of Tantlin a crack between two icebergs gives access to Sheyruushk. Here the sharklike sahuagin are found, sea devils that have made their home in the waters underneath the ice. The sahuagin roam the waters of Sheyruushk. In those places where the ice is thin enough, they can be glimpsed, riding massive sharks, and gliding through the water underneath. However, travelers would be wise to resist the urge to go and see this sight. If the ice is thin enough to allow a glimpse of the sahuagin, it is also thin enough for these fierce aquatic warriors to breach it for an attack. If the characters don't want to be dragged underneath the ice, they best not visit here.

If asked to take them here, \textbf{Koh Tam} refuses. He explains that the sahuagin are known to sink ships that venture into their domain.

\begin{DndSidebar}[float=!b]{Temptation of Greed}
\textbf{\textit{"The Greedy Art Lover"}}



If the characters spend an hour or more exploring this area before returning to a safer location, they find a devilish arm extruding from the ice. It takes half an hour to excavate the corpse and when they do so they observe that the dead creature is holding a ring in its hand, its frozen eyes staring at it in rapture. The characters must break the corpse's hand to take the Ring of Collecting (see appendix D).



\end{DndSidebar}

\subsection{Chasm of Found Things}
Those traveling across the frozen wastes of Stygia, looking for a place to rest may be unlucky enough to wander into the Chasm of Found Things. They may believe themselves fortunate, thinking they've found shelter from the icy cold. But they're mistaken. No one finds the Chasm, it finds you. And once it has tricked you into exploring its labyrinthian caverns, it won't let go without a fight. And yet, this is where our characters must venture to search its hidden depths for the things that were lost to them.

\section{Adventure: The Chasm of Found Things}
The Chasm of Found Things is a mile deep rift in the glaciers of Stygia. Ancient and powerful entities dwell in its depths and lure both devil and mortal alike to their doom.

When the characters first enter Stygia, \textbf{Koh Tam} has the following to say:

\begin{DndReadAloud}{}
"Stygia: the least or the most hospitable of all the Nine Hells, depending on your standing with the infernal hosts

A place run fallow—if a blasted frozen waste atop a fathomless sea can earn the term—under the hand of an absentee duke. Everything starves in Stygia. The elements make no contracts. Truant mortals come here to evade infernal scrutiny. The wilderness, while harsh, offers a fairer fight for survival than other layers. But Stygia has its own traps and surprises. Devils are cruel, but they can be bargained with. Other fates are less open to negotiation." \textbf{Koh Tam} shudders and draws his cloak closer around himself.

"Still, you must venture out there, into the cold. Allow yourself to get lost." \textbf{Koh Tam} gestures at the frozen wastes. "Sometimes those lost within the blizzards of Stygia find signs: mortal tracks through the snow, set down recently. Carvings in the ice claiming a place of safety and shelter. Mortals are drawn across the wilderness until they reach the mouth of a great ice cave that descends into the frozen substrate. The Caves of Rest. Except, of course, this is the Nine Hells and there is no rest. What weary travelers have reached is the Chasm of Found Things, from which few have ever returned. It is here that your search must continue."

\end{DndReadAloud}

\subsection{Advice from Koh Tam}
\textbf{Koh Tam} suggests that they wander into the heart of Stygia if they're to find the Chasm of Found Things and has the following advice:


\begin{itemize}
\item They can't look for the Chasm, they must let it find them. The Chasm only allows itself to be found by those who are lost and exhausted of all hope. The characters must do this by traveling without resting.
\item They should be prepared for the extreme weather, and he suggests they acquire appropriate equipment, magic items, and spells to ward against the cold.
\end{itemize}

Once each character has gained a level of exhaustion, they hear whispers and see strange apparitions in the blowing snow, and eventually find themselves near the Chasm of Found Things, drawn to the entrance (area C1).

\subsection{Chasm of Found Things Locations}
Wandering the frozen wastes endlessly, the characters grow convinced that their search is futile. Exhausted and frostbitten they collapse in the snow. The character with the highest passive Wisdom (Perception) score finds a cave entrance nearby.

\subsubsection{C1: The Entrance}
\begin{DndReadAloud}{}
The mouth of the cave is a curious piece of theatre. Mortal hands have certainly touched the place. The ice is carved into dwellings—huts and chambers laboriously scoured from its freezing substance. Elsewhere, materials have been hauled in to build huts and lean-tos. There are the ashes of old fires. Mortals have dwelled here, and recently.

\end{DndReadAloud}

It is enough to give a living traveler hope and drive devils to a rage. In either case, the idea is to draw the traveler inwards and downwards. From the depths of the cave can be heard faint sounds: the echo of voices, confused and distorted. The caves are occupied; the dwellers at the threshold have retreated deeper into the ice to escape Stygia's harsh climate. And the caves truly are warmer, a little. Enough that meltwater is constantly flowing into them, carving runnels on the floor and trickling inwards so that all the caves resound to the music of it, further confusing the suggestion of life and voices from below. The keen-eyed visitor might feel that signs of habitation at the cave-mouth seem oddly shallow and staged but by then it may be too late. Every visitor drawn to the cave mouth feels a powerful fascination with the depths, a need to know what is down there, the certainty that what they seek can be found below. Those equipped with divinatory magic may discover that the ubiquitous infernal taint of the Nine Hells really is weaker here, and perhaps mistake that for virtue. In truth there is only ever a choice of evils.

\subsubsection{C2: Congregation of Gjaaki}
Progressing downwards from the abandoned structures at the cave mouth the visitor finds a maze of twisting ice tunnels, slippery with the constant run of meltwater, and from this point they're observed. The heedless or weary visitor will just stumble deeper and deeper, but more careful explorers find upward-sloping spurs to the caves that lead to larger spaces clawed out from the ice. Here dwell the first of the Chasm's inhabitants, the Congregation of Gjaaki.

\begin{DndReadAloud}{}
The ice around you creaks and groans and the sound echoes frightfully in the cold, still darkness of the caves. There is no cutting blizzard wind, but the chill of the caves is a heavy thing. The cold pushes down until your throat aches and your bones feel stiff and strange. If you hold your breath and listen, you think you might hear movement in the shadows.

\end{DndReadAloud}

How kuo-toa even reached the Nine Hells is unknown to scholars, but they dwell here in some numbers, farming albino fish in shallow pools and constructing for themselves a truly byzantine mythology. The Congregation are pallid, eyeless, and stealthy, stealing about the upper passages of the Chasm and sporadically descending on visitors, mortal and devil alike, to claim sacrifices for their complex pantheon of gods. Their theology is entirely unique, quite separate from their usual deities. They worship that which lies deeper within the caves and have constructed a whole idiolectic myth-cycle to account for their presence in the Chasm, which they see as the center of a cold and uncaring universe. They make pilgrimages into the depths, sacrifice each other or those that fall into their hands, and conduct meaningless rituals according to nonsensical texts carved into the ice. Sometimes their artisans venture to the cave mouth and create another fake ice dwelling to add to the lure of the caves.

The Congregation is deranged and murderous, but luckier visitors might happen upon them when their bloodlust has been sated, or when their ritual calendar prescribes periods of contemplation. In times like these their philosophers are only too happy to talk about their beliefs for hours and hours, to the degrading sanity of their listeners.

The Congregation consists of three dozen kuo-toa monitors and a trio of kuo-toa archpriests. The kuo-toa are in a period of contemplation and therefore peaceful, if not attacked. However, they're so deranged that they provide scarce information and the characters would need to spend considerable time moving between the various tunnels to encounter all of them. If a kuo-toa is attacked here, it summons all its brethren to defend it. If the entire kuo-toa populace in this area is slain, the characters recover 8,000 gp worth of various gems and other valuables.

\begin{DndSidebar}[float=!b]{The Blind Believers}
The kuo-toa living in the chasm have blindsight out to a range of 60 feet and resistance to cold damage.



\end{DndSidebar}

\begin{DndSidebar}[float=!b]{Random Temptations}
If the characters attack the kuo-toa and then spend considerable time searching or exploring this part of the cave system, roll on the Random Temptations table from chapter 2.



\end{DndSidebar}

\subsubsection{C3: The Echoes}
Below, the tunnels twine and corkscrew downwards. In these reaches, travelers might encounter Congregation pilgrims trekking towards their ultimate religious experience, or sometimes fugitive devils—ice devils and pain devils (see appendix B)—drawn into the caves seeking mortals to torment or escaping more powerful devils. Most common, though, are the Echoes. Read the following:

\begin{DndReadAloud}{}
The ice itself contains movement, even as the air carries phantom sounds. There are flickers of motion, ahead, below, to either side. Figures you can never catch up with, the faces of lost friends and family, treasures, beckoning hands. The very ice of the walls is shot through with an ephemeral life that, like the water, is constantly heading further down.

\end{DndReadAloud}

Each time a character attempts to move past an Echo (a location marked as C3 on the map), in the direction of the surface of Stygia, they must make a DC 15 Wisdom saving throw. If they fail the saving throw, they gain a level of exhaustion and fall (they now have the prone condition). They may repeat the saving throw each hour.

\subparagraph{Phylactery of the Oathbreaker}
The oathbreaker finds their soul imprisoned among the Echoes in the ice when they, and any companions, see a pre-adolescent reflection of themself begging for mercy. Other reflections, devils in the guise of four children, are tormenting the character's soul relentlessly. To save the soul, characters must break through the ice and enter the Echo. This requires an action and, when the first character is inside, they find themselves in a 60-foot diameter chamber, surrounded by ice. The bullies transform into four chain devils surrounding the child (a \textbf{commoner}). The devils attack the characters.

Whenever a character starts their turn inside the Echo, they take 5 (1d10) cold damage. Additionally, at initiative count 20 of round 4, any surviving devils seemingly transform back into children (but they retain all their devil statistics). If a devil in child's form is slain, it vanishes and reappears at the start of the next round, fully healed. At initiative count 20 of round 5, the enemies transform back into devils and this process repeats every 4 rounds.

During this combat, if the oathbreaker's soul is slain, it reappears at the start of the oathbreaker's next turn, but an affliction devil (see appendix B) appears next to the oathbreaker, to assist the other devils in battle. These devils can be permanently slain only while they wear their own form. Once all the devils are defeated the echo of the oathbreaker's younger self vanishes and leaves the phylactery in its place.

Now in possession of their soul, the character must escape the ice. However, the longer the characters were inside the Echo, the thicker the ice around them has grown. Every time a devil transformed into a child, the ice-chamber was pulled deeper into the Echo, and another 10 feet of ice grew between the chamber and the original entrance into the Echo. Each 10-foot section must be destroyed as if it were created by the \textit{Wall of Ice} spell.

\subsubsection{C4: Angel of the Ice}
\begin{DndReadAloud}{}
The cave is large, with a floor that has been worn smooth. There is faint light here, and it illuminates a wide wall of ice covered in complex and disturbing diagrams. Before the mural on the wall stands a manyeyed being of bizarre form, as pale and cold as the ice. Its voice is like the grinding sound of an iceberg shattering. Loud. Violent. Powerful. All around it sit kuotoa, chattering their teeth and slapping the ground, their eyeless faces warped with fanatic adoration for what their Angel of Ice preaches.

\end{DndReadAloud}

Like any religion, the Congregation has heretics to persecute. Separate caves hold a smaller community of dissidents led by the Angel of the Ice (lawful evil \textbf{deva}). This fallen Celestial claims to have been sent to the Chasm to reclaim it from the Nine Hells and build a gateway to the higher realms, preparatory to ushering in a vast army to cleanse all the lower planes. It is just as deluded as the Congregation, dwelling before a mirror-like wall of ice on which it has scratched out a vastly complex diagram that it says holds the secret to the movements and future of the planes.

\subparagraph{Defenders of the Faith}
The Angel is attended by a cult of kuo-toa that it has attempted to reshape into something more holy, resulting in a curious admixture of aasimar-like qualities into its brood. There are a dozen of these heretics (kuo-toa monitors that can change the damage type of their attacks to radiant damage). Each heretic carries 2d20 gp.

If the ice wall, or the \textit{Holy Avenger} within it, is disturbed for any reason, the following occurs:

\begin{DndReadAloud}{}
Immediately all activity and sound ceases in the iceshrouded cavern and, as if of one mind, the Angel and their cult of kuo-toa, turn towards you. There is no more chattering, no more grinding. Just a silence colder than the frigid air swirling about you. The Angel points their finger at you, shakes their head, and then resumes their strange oration.

\end{DndReadAloud}

This was a warning. Any further attempts to desecrate the temple turns the entire group hostile to the characters and they attack.

\subparagraph{Conclave (Paladin)–Holy Avenger}
The Angel of Ice eventually found themselves incapable of wielding their \textit{Holy Avenger} and so they ensconced it in the ice. Just the hilt sticks out now and anyone examining the diagram on the wall easily notices it. It requires two successful consecutive DC 18 Strength (Athletics) checks to pull the sword from the ice. Each check requires an action.

\subparagraph{Phylactery of the Business Partner}
The business partners' greed in life knew no bounds, as the character searching for his soul knows only too well. The business partner desired riches and he did not want to share them with anyone. The character looking for their former business partner's phylactery finds it here with the Angel of the Ice.

If the character examines the mural closely, they soon realize its importance when they recognize a tiny figure trapped inside. The business partner is wandering the Celestial's mazelike diagram. Their torture is one of endless frustration, believing that within they can find riches grander than any they ever possessed in life. They're consumed by the desperate need to attain them. However, every corner they turn leads to a dead end and is populated with a new horror that will gladly consume their flesh. Once consumed they pop up again at another point on the mural.

The characters must destroy the mural (AC 12, 100 hit points) but the Angel and its cult attack them, as described earlier. Once the mural has been destroyed the phylactery of the business partner appears on the floor. The character takes possession of the phylactery and gains the Phylactery Benefit associated with it.

\subsubsection{C5: The Aboleth}
\begin{DndReadAloud}{}
There is a dripping sound in the cold tunnels. The space around you feels emptier, the walls have pulled away and the caves that lie ahead are vast and full of lakes of ink-dark water. The dripping sound comes from droplets of meltwater rolling off the distant ice-ceiling and striking the liquid depths below. You occasionally see ripples break the glass-still surface, as though something very large were moving beneath it.

\end{DndReadAloud}

At last, the caves begin to open out. A network of interconnected lakes collects the meltwater, each lake vast and lightless. Here dwell the aboleth, each the master of its own circumscribed domain, each a minor deity in the Congregation's tangled pantheon. Chuul servitors act as guards and food in times of need, and the lake bottoms are a sifting morass of bones and forgotten treasures brought by devils and mortals alike who got no further. But even this is merely a station on the way down. For each hour the characters search they have a 75 percent chance of finding an aboleth clutch (one \textbf{aboleth} and 2d6 \textbf{chuul} protectors). This chance is reduced by 25 percent for each clutch found.

\subparagraph{Treasure}
Once a clutch is dealt with, the characters find 6d4 x 100 gp and two uncommon magic items from the \textit{Dungeon Master's Guide}.

\subparagraph{Conclave (Artificer)—Manual of Golems (Iron)}
A character searching for the Manual of Iron Golems discovers it in the clutches of the first aboleth they encounter.

\subsubsection{C6: Keeper of Found Things}
\begin{DndReadAloud}{}
Below you stretches the greatest gulf, the Chasm itself, with a craggy icicled ceiling above and only the waters below. To either side a ledge extends, overlooking the subterranean sea and running to the very end of the Chasm. Echoes in the walls beckon for you to continue down the ledge. A whisper inside your head tells you that the Keeper of Found Things has been expecting you.

\end{DndReadAloud}

What is the Keeper of Found Things? Not a devil, certainly. Something ancient that has dwelled in the Chasm and tended its collection for longer than infernal records have been kept. A thing of the shadows that even the keenest eyes can never quite make out in detail or gain a complete picture of. Rambling accounts speak of too many arms, long-fingered hands, a twisted wedge of a head with a vertical mouth bristling with needle teeth. And a voice. A quiet, reasonable (so very reasonable) voice that speaks to mortals about salvation, rest, and peace—all the promises that lured them down into the caves past all those other tribulations. To forlorn devils, the Keeper instead whispers of a great and endless order, a perfect moment of torment perpetually prolonged. The Keeper has only terrible things to say, and yet it is hideously persuasive. Every moment of the journey through the caves goes to prepare its visitors to accept its blandishments. The Congregation's nonsense theology, the Echoes, the tentacled monsters thrashing blindly in the depths, all plant the seeds of surrender in the mind.

Within the icy walls of the Chasm everything that has filtered down to the Keeper over the long years is buried. There are treasures that have washed down with the meltwater or that have been brought by great heroes and devils. There are monsters, trapped forever—dimly glimpsed against the strange flat light that emanates from every direction. And there are the mortals and devils, hundreds of them, who came this far and no farther. Who couldn't resist the soft entreaties of the Keeper. Here is rest and peace and an endless perfect cruelty. Over time the ice draws its captives deeper, to become the most recent prizes the Keeper claims. Anyone who travels to the end of the Chasm must make a DC 15 Constitution saving throw. A creature that fails the save begins to be enveloped by the ice and has the restrained condition. The restrained creature must repeat the saving throw at the end of its next turn, becoming encased in ice on a failed save or pulling free on a successful one. Anyone encased in ice takes 19 (3d12) cold damage at the start of each of their turns. They can be freed by destroying the surrounding ice, which has AC 13, 40 hit points, vulnerability to fire damage, and immunity to poison and psychic damage.

For those of stronger will, the ice itself reaches for them. The Keeper reshapes the very stuff of the Chasm to entrap less compliant victims by drawing forth frozen golems, gleaming facets flickering with echoes drawn from their victims' past. The luckiest might find their way back to the caves, to struggle against the slope and the water, braving all the perils they've already passed for the dubious safety of Stygia's surface. Others, in desperation, might fly or climb upwards out of the reach of the Keeper's influence, or plunge into the waters of the freezing sea below.

\subparagraph{Treasure}
For each hour spent at this location excavating valuables, select a very rare magic item from the \textit{Dungeon Master's Guide} (excluding weapons). However, all characters participating must repeat the earlier Constitution saving throw, or suffer the negative effects described earlier. Additionally, for each character failing their saving throw, a \textbf{frozen golem} (iron golem with immunity to cold) emerges and attacks.

\subparagraph{Deathstalkers—True-Ice Shards}
A character searching for one of the \textit{True-Ice Shards} sees it lodged in the ice at the end of the chasm. If they touch the shard while still encased in the ice, they take 17 (5d6) cold damage and must succeed on a DC 18 Constitution saving throw. If they fail, they can't let go of the object and the ice starts enveloping them. The character takes 17 (5d6) cold damage at the start of each turn that they're still holding the shard, but doesn't repeat the saving throw.

They can only dislodge the shard with fire damage, any cold damage makes the ice enveloping the shard grow thicker. Once 50 fire damage is applied the shard comes free.

\subsubsection{C7: The Awful Fisher}
\begin{DndReadAloud}{}
The ceiling is full of deep crevices and covered in sharp, reaching icicles. You can almost make a game of following the patterns of creeping shadows they cast. They seem nearly alive, the shadows, slinking over the jagged surface as the light below moves. Sometimes it seems the shadows shift out of time with the light. Yes. It's strange. Near the ceiling, you notice one of the shadows, a large one, is moving very quickly. Much quicker than the rest around it. The shadow is coming right at you. You can hear it now. Many legs, clicking over the ice and the terrible snapping sound of giant, hungry claws.

\end{DndReadAloud}

The ceiling offers the most mundane of fates. It has one inhabitant, that has scavenged from the Keeper's leavings since time immemorial. It is known to the kuo-toa as the \textbf{Awful Fisher} (see accompanying stat block). Those who fly or scramble upwards from the Keeper's grotesque collection find a great many-legged thing watching them, letting fly with a long, sticky thread. Nothing more than a preternaturally large cave fisher, in fact; those crablike creatures that bedevil the mortal Underdark. The \textbf{Awful Fisher} has grown vast and cunning on a diet of infernal flesh over the ages, though, given most of those who fly into its reach are devils. It is eternally hungry, malevolent, and armored beyond the bite of most blades. After the \textbf{Awful Fisher} dies, it falls to area C6 and the corpse begins to freeze over as described in that section.

\subparagraph{Treasure}
Once slain, the contents of the \textbf{Awful Fisher}'s stomach may contain some treasures. If searched, select two uncommon magic items from the \textit{Dungeon Master's Guide}.

\subparagraph{Conclave (Cleric)—Rod of Resurrection}
A character searching for the \textit{Rod of Resurrection} discovers it in the belly of the \textbf{Awful Fisher}. During the battle with it, the bizarre creature accidentally triggers the Heal property of the rod the first time it falls under 100 hit points. Any nearby cleric feels a rush of divine power emanating from within the Monstrosity when this occurs. A successful DC 14 Intelligence (Religion) check allows a character to locate the \textbf{Awful Fisher}'s corpse, and the rod, before the creature becomes completely encased in ice.

\subsubsection{C8: Gjaaki}
And the waters? Nobody has ever returned from the waters. The kuo-toa Congregation's greatest deity, Gjaaki, resides there, a being of utter negation, an anti-creator. They say that all things that it becomes aware of cease to be and fade from the memory of the universe. Their myths speak of heroes who never were, offered up in proof of the power of this god. Of course, the Congregation are utterly divorced from reality and nothing they say can be taken seriously. But something is down there that even the Keeper and the aboleths don't speak about, beyond the knowledge of devils and the powers of the Nine Hells.

\chapter{Malbolge, the Realm of Decay}
\DndDropCapLine{P}{}ractically infinite mountain slopes and falling boulders make up Malbolge, the sixth layer of the Nine Hells. It was once ruled by Malagard, the Hag Countess, who has since been replaced by \textbf{Asmodeus}'s daughter \textbf{Glasya}. Prior to her rule, the layer was almost entirely stone, with very little distinguishing features. However, due to a curse placed upon Malagard, her body grew to a grotesque size and imploded when \textbf{Glasya} took over. Now Malbolge is decorated with pieces of her corpse, a marriage between a rotting carcass and a cliff-face. What little pockets of devil civilization remain do so in caves carved from the mountain or inside fortresses supported by indestructible pillars.


\section{Running This Chapter}
Before running this chapter read the "Malbolge Overview" section. It provides you with everything you need to guide your players across this bloated and decaying landscape. Use the information provided to set the scene during their travels through Malbolge.

\subsection{Encounters}
The denizens of Malbolge are divided in their loyalties but united in their fear and disarray over the changed landscape. If the characters are wandering far and wide, roll on the Random Encounters in Malbolge table to determine what they might encounter.

\begin{DndTable}[header=Random Encounters in Malbolge]{cX}

d6 & Encounter\\
1 & There is a terrible shaking as the earth nearby splits apart and a broad-faced beast, as large as a dragon but wingless, erupts, a mouthful of rocks and gems falling from its mouth. It is a \textbf{miasmorne} (see appendix B) burrowing through the Nine Hells, in search of metal and magic to devour.\\
2 & Inside a small, fortified structure is a corruption devil (see appendix B). It hates \textbf{Glasya}'s rule and believes it can revive the Hag Countess.\\
3 & A \textbf{maelephant nomad} (see appendix B), driven slightly insane by its extended lifespan, searches the ruins for something to guard. When it sees the characters, there is a 50 percent chance it goes into a frenzy and attacks.\\
4 & A patrol of six \textbf{erinyes} search for any remaining influence from the Hag Countess. They reward anyone with information relevant to their search with 2,500 gp, and attack hags and their allies on sight.\\
5 & An invisible \textbf{orthon} (see Monsters of the Multiverse) has been stalking the characters for some time. There is a 75 percent chance it attacks them now. Otherwise, it moves on, searching for its quarry.\\
6 & Massive tremors shake the earth, causing an avalanche of rocks to plummet down. Each creature in the path of the rocks must succeed on a DC 18 Dexterity saving throw or take 35 (10d6) bludgeoning damage.\\
\end{DndTable}

\subsection{Locations}
The characters must seek out The Sign of the Hag's Arms, an inn and tavern, but no safe haven for mortals. There they may discover some of what they seek. The "Key Locations in Malbolge" section provides additional details about other locations in this layer.

\subsection{Koh Tam}
Ensure that \textbf{Koh Tam} is urgent in directing the characters to The Sign of the Hag's Arms. He can also provide useful information about some of the other locations herein, though he doesn't believe what the characters seek can be found anywhere other than the inn.

\subsection{Objectives}
Make sure you keep track of your player's objectives and lead them to the corresponding areas to ensure they can complete their goals.

These objectives can be attained in Malbolge:

\begin{DndTable}[header=Objectives in Malbolge]{XX}

Objective & Location\\
The phylactery of the mentor & The Sign of the Hag's Arms: The Formian Queen\\
The phylactery of the heartless master & The Sign of the Hag's Arms: The Kitchens of Irabella\\
The phylactery of the mother & The Sign of the Hag's Arms: Irabella's Workshops\\
The \textit{Scourge of Shadow} & The Sign of the Hag's Arms: \textbf{Anacreda} the Angelmaker\\
\end{DndTable}

\subsection{Temptations}
Have the characters encounter at least one temptation during their time in Malbolge. You can of course have them encounter more if you want. If the characters give in to temptation, use the information in appendix E to keep track of their corruption level.

\section{Malbolge Overview}
Malbolge is considered the most dangerous stretch of the River Styx. Due to rockslides and erosion, boulders and rock formations sit within the river, changing in shape and location each day. Parts of the river mix with acid and blood from Malagard's remains, creating a toxic soup that dissolves the hulls of unprepared boats and devours anything thrown overboard in seconds. Should a captain survive all these hazards, they must navigate through a claustrophobic canyon with steep walls, where the incline eventually levels out and empties into Maladomini.

The surface of Malbolge is nightmarish—an already treacherous landscape is made more dangerous by the remains of the hag queen. Throughout the layer exist forests where the trees resemble hairs of a gargantuan beast, and caustic lakes of strange substances. Her sometimes-slippery sometimes-sticky blood covers large swathes of Malbolge. Massive bones, some shattered and some intact, create artificial structures across the land, with pieces of teeth and cartilage sticking out like grotesque obelisks.

Hoping to survive the trials of the layer more easily, devils constructed an intricate labyrinth of tunnels connecting the cities together. Much of the realm's inhabitants navigate the tunnels, expanding and creating new pathways and connections. Some believe \textbf{Glasya} plans to construct entire cities underground, though such plans would take centuries to complete.

Before the Hag Countess ruled, Moloch claimed ownership of the realm. When \textbf{Asmodeus} grew tired of his tricks, he appointed Malagard, and in so doing cursed her with an insatiable hunger. She ate devils and the souls of the damned; her body grew over decades, but her position as Malbolge's ruler kept her alive. When \textbf{Glasya} was appointed ruler of Malbolge, the blessing expired, and Malagard's mountainous bulk exploded in a shower of flesh and gore that covered much of the layer. Many of the hags that called Malbolge home left, leaving \textbf{Glasya}'s subjects small in number.

\subsection{Leaving Malbolge}
Like Avernus, Malbolge contains a few exits to the layer below, though none are easy. While the Styx is an option, only the most experienced captains can navigate the dangerous rapids. Devils native to the layer claim the labyrinth of tunnels has a passage leading to Maladomini, though without a guide such a passage would be difficult to traverse. Even so, with the dangers of the Styx, many opt for the caverns.

\subsection{Features}
The remains of the Hag Countess poison the entire layer, including visitors. Good-aligned creatures have vulnerability to poison damage while in Malbolge, even if they would normally have resistance or immunity. Non-good creatures have their resistance and immunity removed, but don't suffer vulnerability.

In addition to the poisoned air of the layer, creatures also feel drained the longer they stay in the realm. Each time a creature that isn't a Celestial or Fiend finishes a long rest, they must make a DC 16 Constitution saving throw. If they fail, they have disadvantage on all saving throws until they leave Malbolge.

\section{Key Locations in Malbolge}
While the characters are making their way to The Sign of the Hag's Arms, they may want to take a small detour through the other locations found in Malbolge.

\subsection{Ossiea}
Above ground, a few locations were built upon the grotesque remnants of Malagard's body, the parts that survived the cataclysm. These pieces of her enormous, bloated body are still relatively intact and serve as gruesome landmarks. Among these, the characters may choose to visit Ossiea. This fortress is found near the Hair Forest and built into the skull of Malagard, her empty unseeing eye sockets covered by tall sheets of red glass. Ossiea is the capital of the realm and current home of \textbf{Glasya} and her subordinates. \textbf{Glasya} is obsessed with beauty, and the city around her reflects this. The shops the characters visit here are stacked with magical cosmetics, intricate jewelry, and beautiful garments.

\begin{DndSidebar}[float=!b]{Temptation of Pride}
\textbf{\textit{"A Thing of Beauty"}}



Whilst rummaging the stalls around the city, the character with the lowest Charisma is approached by a devil who offers to sell them an amulet of great power for 600 gp. This is an Amulet of Appearance (see appendix D) and if the character initially declines the purchase, the charming devil insists the character tries it on before deciding. A successful DC 15 Wisdom (Insight) check reveals it to be a magic amulet but not that it is cursed. If the character tries the amulet on, they transform. Any blemishes and imperfections vanish, they've never looked better. Audible gasps of approval can be heard from onlookers who see the transformation take place.



\end{DndSidebar}

\subsection{Maggoth Thyg}
Near the bottom of the layer also lies Maggoth Thyg, a mysterious cavern from which no creature ever returns. Even the devils native to Malbolge fear Maggoth Thyg, claiming it has existed far longer than the Nine Hells. The characters should know to be wary of a location that even the denizens of the Nine Hells are fearful of.

\begin{DndSidebar}[float=!b]{Random Temptations}
If the characters decide to go in search of Maggoth Thyg despite the obvious perils, have them venture for four hours towards a cavern system at the bottom of Malbolge. When they get close to the area where they expect to find the cave, the ground is scattered with corpses. If they insist on venturing at least another hour further, roll on the Random Temptations table from chapter 2.



\end{DndSidebar}

\subsection{The Sign of the Hag's Arms}
Another landmark utilizing Malagard's remains in its construction is the Sign of the Hag's Arms. This multi-leveled structure rises out from the ground near a place where many tunnels and thoroughfares cross paths. One of the hag's dismembered arms serves as a gruesome marker for this location, which \textbf{Koh Tam} deems the most likely to house what the characters seek.

\section{Adventure: The Sign of the Hag's Arms}
The Hag's Arms is run by hags, of course, the entrepreneurs of the lower planes. There are three of them, as is traditional, each with her own specialty, each a powerful magician beyond the regular dreams of her kind. Irabella runs the inn and kitchen, the night hag Malacki tinkers on hideous creations in her underground workshops, and \textbf{Anacreda} the Angelmaker is the coven's undisputed leader.

When the characters first enter Malbolge, \textbf{Koh Tam} or \textbf{Tiax} impart the following knowledge about the layer to them:

\begin{DndReadAloud}{}
\textbf{Koh Tam} explains the windswept and endless slopes of this new layer of the Nine Hells, "In Malbolge there is no level ground, but always the strength-sapping slog of up or down, and the going is riven with crevasses and the colossal, eternal remains of Malagard, the former ruler of this layer, constantly festering but never decaying away."

"There are tunnels, lairs for such fauna as calls Malbolge home: the insects and spiders with their kin, in swarms and of huge size. Some say they were spawned from the corruption of Malagard's vast fleshy diaspora. The things have adapted their natural weapons and strategy to prey on fiends, though they will take mortals gladly."

He toes a clump of profane soil, releasing numerous scuttling insects that flee into the dirt below and then continues, "But those, devil, or mortal, who lament the cheerless character of the plane should dampen their enthusiasm when the Sign of the Hag's Arms comes into sight. Hospitality is to be had there, it's true, but not every guest who enters its doors has the luxury of checking out again."

\end{DndReadAloud}

\subsection{Advice from Koh Tam}
\textbf{Koh Tam}'s advice to the characters is as follows:


\begin{itemize}
\item They should visit The Sign of the Hag's Arms, where they might start looking for answers. But he warns them: the hags are powerful adversaries.
\item They should avoid going to Maggoth Thyg, he has heard many stories about those caves and none of them ended well.
\end{itemize}

\subsection{The Sign of the Hag's Arms Locations}
Some of the major locations found in and around the inn are described below.

\subsubsection{S1: The Swarm}
\begin{DndReadAloud}{}
You stand upon the lower slopes of a great uprising of earth. A well-worn path curves upwards but a fetid mist hangs low, allowing only the murky outline of a building crouching upon the hilltop above. Closer to you there appear to be numerous, less-worn paths, leading into various crevices and small canyons.

\end{DndReadAloud}

The inn sits at a confluence of chasm and tunnel thoroughfares between infernal fortresses, in a particularly perilous region of the plane. Perilous more for Fiends than for travelers from outside the lower planes. From long feeding on the toxic flesh of Malagard—perhaps helped along by the tinkering of the inn's proprietors—the local life has developed some potent weapons to deploy against devils. In place of the regular poisons that Fiends are immune to, many species of giant scorpions and giant spiders inject acid with their stingers and fangs, or even more bizarre fluids that sear like the radiant touch of angels, rendering down fiendish flesh into a readily ingestible soup.

The tumultuous air above isn't safe either, as flights of ayperobo swarms (see appendix B) issue from their crevice nests, eager to devour prey. Those devils tasked to take messages and goods through the region of the inn always hurry to take shelter there, no matter the risks, and mortal travelers would be advised to do the same.

\subparagraph{Ekengarik's Swarms}
Approaching travelers are set upon by 2d6 giant scorpions and 2d6 giant spiders that crawl out from the various crevices around them. If the characters take to the air, an \textbf{ayperobo swarm} (see appendix B) joins the battle. The insects don't pursue if the characters climb the hill towards the inn.

\begin{DndSidebar}[float=!b]{Corrupted Insects}
Any giant insects encountered in Malbolge deal acid damage instead of poison damage and each of their attacks deals an additional 4 (1d8) radiant damage. These insects also have resistance to acid, poison, and fire damage.



\end{DndSidebar}

\subsubsection{S2: Ekengarik's Hive}
\begin{DndReadAloud}{}
You notice numerous gashes on stone outcroppings, crevices, and the ground itself, that apparently form the entrance to tunnels, leading into a darkness below. Peering into an opening, you see a network of tunnels and chambers illuminated by a soft, pulsing glow and hear the faint sound of scraping and digging. The tunnels are lined with hardened earth, and seemingly form a complex labyrinthine network. The air is thick and humid, carrying the scent of dirt and decay.

\end{DndReadAloud}

\subparagraph{Ekengarik's Duty}
Malecki provides \textbf{Ekengarik} with ever more noxious underlings to serve her hive, and she in turn ensures that travelers in the locale are 'encouraged' to visit the inn, and those trying to escape without paying their due meet a sorry end. It was \textbf{Ekengarik} that sent the swarm of insects towards the characters earlier. \textbf{Ekengarik} (see accompanying stat block) wants to drive the characters to the inn, but her patience is limited. Keep track of this by using a counter. The counter starts at 1 if the giant insects have been slain, zero otherwise, and for each hour the characters spend in the hive, and each time a fiendish formian is slain, it increases by +1. Once this counter reaches 6, \textbf{Ekengarik} uses Dimension Door and attacks, generally starting battle with Prismatic Wall to separate the characters from going deeper into her hive. If at any point the characters retreat, she allows them.

\subparagraph{Treasure}
Most corridors in the hive are 5 feet wide, 10 feet tall, and extend in stretches for 50-90 feet before turning sharply. A variety of stones, wooden debris, and various items collected over the years, has been partly cemented into the tunnel walls. As long as the hive isn't disturbed, characters may search for treasure in the hive. For each hour spent searching roll a d10. If the result is 10, 1d4 fiendish formians (see accompanying stat block) attack, the hive is considered disturbed, and no more treasure can be acquired here. Otherwise, the characters find 2d4 x 100 gp and 1d4 uncommon magic items from the \textit{Dungeon Master's Guide}.

\subparagraph{Phylactery of the Mentor}
If one of the characters has ventured into the Nine Hells to find the soul of their mentor, their attention is caught by an imp constrained at the outer edges of \textbf{Ekengarik}'s hive. This area is less heavily populated with the swarming creatures and the characters can try and sneak closer. The imp is the prisoner of a \textbf{fiendish formian} taskmaster. Over and over again his energy is painfully sucked out of him to sustain the taskmaster that controls him. Once the imp collapses in exhaustion the taskmaster moves away, coming back for more once he needs sustenance again. The imp is in truth the character's old mentor transformed to undergo this torment. To rescue them the character has to defeat the formian that is treating their old master as a snack. Preferably without alerting the entire hive. If the characters take more than four rounds to slay the taskmaster another 1d4 fiendish formians join the fray.

Once the characters have killed the formians read the following:

\begin{DndReadAloud}{}
The imp hangs limply in his chains, overcome with exhaustion. As you approach it, the imp's body starts to shake violently. As abruptly as the spasms started, they stop. The body sags, seemingly the imp has died. After a few seconds its flesh begins to bubble and evaporate, leaving a phylactery behind.

\end{DndReadAloud}

The character has taken possession of the phylactery of the mentor and gains the Phylactery Benefit associated with it.

\subsection{Aboveground Inn Locations}
Having survived the perilous journey through the Hive's territory, the characters arrive at the relative safety of the inn.

\subsubsection{S3: Entrance}
Read the following as the characters approach the inn.

\begin{DndReadAloud}{}
The Sign of the Hag's Arms is well, if bleakly, named. The outer courtyard sits in the crook of one of Malagard's lesser limbs, or at least a recognizable fragment of the same. Within this elbow space the upper level of the inn itself is made entirely from the ancient hag's substance, bones as scaffolding, leathery skin stretched out as a vaned roof, the walls built of tessellated teeth and ragged plates of nails. As a feat of architecture, it is simultaneously impressive and nauseating.

\end{DndReadAloud}

Had the plane's current ruler, \textbf{Glasya}, any respect for her predecessor then doubtless the construction would be counted as deep disrespect. \textbf{Glasya} isn't sentimental, however, and Malbolge offers little else in respect of building material. Even the iron of the infernal fortresses must be hauled in from more productive layers. As with most locations in Malbolge, the inn continues on many levels under the ground. There are always guest rooms to be had, for a price. The proprietors accept a variety of currencies. Soul Coin, most commonly, but also magic objects or rare ingredients, given the work all three tend to. Devils passing this way know to bring appropriate recompense for the shelter, but mortals may be dismayed to find that all their gold and gems buy them nothing. However, there are always other options for those visiting from other planes. Not necessarily palatable options.

The characters can purchase lodgings for 1 Soul Coin or an uncommon magic item. Then they're taken to their rooms (see area S6).

\subsubsection{S4: Feasting Hall}
\begin{DndReadAloud}{}
You aren't certain what the staircase is made of. Every part of the inn seems suspiciously biological in nature, though very dead, and there is an inescapable scent of decay underlying the air. At the top of the staircase lies a long hallway with many doors, the first pair of which are open. Through them you see a great feasting hall with a host of long tables covered in heaps of food. Most of it seems to be meat. There is a cacophonous swarm of devils around the tables, gorging themselves on the strange offerings and shrieking and screaming at each other.

\end{DndReadAloud}

\begin{DndSidebar}[float=!b]{Temptation of Oppression}
\textbf{\textit{"Painful Knowledge"}}



If the characters decide to rent a room at the inn or find themselves exploring the guest rooms without payment, they find a mirror by the side of the bed. The mirror is a Sage's Mirror (see appendix D) and is small enough to carry.



\end{DndSidebar}

A great feasting hall is found on the first level of the inn, usually hosting at least a score of devils at any given time. They jostle elbows at tables scrimshawed from Malagard's bones and eat from a unique menu. The flesh of Malagard is generally unpalatable even to Fiends, but the inn's proprietors have the knack of transforming it into a delicacy worth the journey—to a devil. Mortals are advised to bring their own lunches unless they have a cast iron immunity to poisons.

The hags are always on the lookout for new business associates, or spare parts for their creations. If at any point the characters want to speak to the owner of the inn, they can ask any of the kobold serving staff. The kobold explains who the three hags are and that they often remain in their workshops in the basement. If the characters ask to speak to the hags directly, the servant leaves to fetch them. Roll on the Dealing with Hags table to determine the result of the request. If the character provides a bribe or incentive to encourage the hag to meet, then add +2 to the roll.

\begin{DndTable}[header=Dealing with Hags]{cX}

d12 + #$prompt_number:default=0,min=0,max=2,title=Enter +2 if the characters provide a bribe or incentive$# & Result\\
1–2 & A kobold servant carries a silver tray with the head of the previous servant on it. The characters are asked to leave the inn.\\
3–6 & The kobold returns and says the hags have declined, but the characters are welcome to meet with them in the basement. The magically locked door at S7, will be open.\\
7–9 & Malecki arrives an hour later in the form of an elf woman and sits down at the character's table. She intends to trick the characters.\\
10–11 & Malecki dematerializes from the Ethereal Plane and appears at the character's table, in her true form, willing to make a deal.\\
12+ & \textbf{Anacreda} emerges from the basement, her wings partly unfurled as she stalks across the floorboards. The feasting hall quiets, many of the devils suddenly staring into their cups, averting their eyes. She stops and looms over the character's table.\\
\end{DndTable}

\subsubsection{S5: Kitchens}
The kobolds who staff the inn are routinely maimed, killed by accident, tortured, and even eaten by their employers, especially Irabella. They're very aware that they're the weakest things in all the Nine Hells, more wretched even than lemures, and while nothing keeps them safe from the hags, Irabella takes exception to any guests who take liberties with the staff. Working for her is miserable but better than being cast to the mercies of the Nine Hells. The most potent, sorcerous, and wicked among them ascend to the heady rank of chef and become Irabella's chief hench creatures, petty kitchen tyrants among their own kind.

\subparagraph{Phylactery of the Heartless Master}
A character searching for the soul of the heartless master finds it in the body of a bullied kobold. Upon glancing around the kitchens the characters notice a wretched looking kobold bent over a stove. It is frantically trying to light it, but the stove remains cold. The kobold is clearly agitated and terrified. The character seeking this phylactery senses they must help this kobold.

The creature explains to the character that it is being tormented by a ruthless \textbf{sea hag}. Every day it is tasked with preparing a feast for her, but whatever it tries it is never good enough, and the day ends with the hag inflicting unspeakable punishment on the kobold until it begs for the release of death. But that release never comes. It always awakens again in the kitchens with a fire that won't light and a list of dishes it doesn't know how to prepare.

If the character looks at the names of the dishes, they recognize some of the delicacies that were on offer in Aelvette's kitchens in Dis. The hag desires a meal of candied spider eggs, roasted jackalwere, and boiled shank of minotaur. She arrives once the meal is prepared to sample the dishes.

A successful DC 15 Intelligence check accurately identifies the ingredients. If the characters spent any time in Aelvette's kitchen they have advantage on the roll. They need to search the storage room for the ingredients. 1d6 hours of searching per recipe is needed to supply them with what they need. Once they get cooking the oven continually goes out, but a successful DC 12 Dexterity check allows them to keep it lit. If the character has proficiency with cook's utensils, they add their proficiency bonus to the roll. Upon success the character and the kobold have created a feast even the hag can't fault. Frustrated, she flies into a rage and attacks. The character needs to defeat the hag to free their soul. The kobold doesn't come to their aid during the fight. Upon killing the hag, the kobold vanishes. In its place the phylactery of the heartless master is standing.

\begin{DndSidebar}[float=!b]{Coven Spellcasting}
The three hags that run the inn form a coven that is much more powerful than what is typical. Their coven shared spellcasting ability is as follows:





\end{DndSidebar}

\subsubsection{S6: Guest Rooms}
\begin{DndReadAloud}{}
The guest rooms make up most of the upper levels of the inn. The walls seem to be organic and emanate a discomforting warmth. Unfortunately, the narrow corridors make it almost impossible to avoid contact with them. A putrid smell hangs in the air.

\end{DndReadAloud}

The guest rooms themselves are comfortable enough, if one doesn't contemplate the materials used to furnish them. There are six double bedrooms on each level.

\subsection{Basement Levels}
Below all of this—the rooms, the hall, the unspeakable kitchens—are the workshops, for the three old women who run the inn have another business they conduct here. Business that draws infernal patrons from every layer of the Nine Hells.

\subsubsection{S7: The Basements}
The work of the hags, when they aren't playing hostess to select visitors, is in the construction of monsters. The secret lies in the flesh of Malagard, which is infinitely dead yet simultaneously bubbling with corrosive and stinking life. Using carefully selected components harvested from the ancient archdevil, adding in parts, pieces, and rare ingredients collected from across the Nine Hells and beyond, they create new beings with independent existence and will. For, while they lack the divine spark that would permit them to create true life, still the carrion flesh of Malagard holds enough vitality to grant a permanent animation to their stitched-together creations. The secret of how they harness the deceased hag-mother's power is known only to the three of them and represents their last bargaining chip should the powers of the Nine Hells decide their tenure is at an end.

While the basements are sealed off from the rest of the inn, there are no guards when the hags are off dealing with guests. The basements are accessible through an iron door that has Arcane Lock on it. Read the following if the characters decide to investigate:

\begin{DndReadAloud}{}
The basement of the inn is a huge cavern carved from the rock. Arches of bone are built into the walls of the cavern for some unknown purpose. What looks like an angel hangs slackly from hooks on the ceiling; it seems to have been constructed out of the body parts of angels and devils.

\end{DndReadAloud}

This cavern is where the creations that the hags are currently working on are found. The angel doesn't have a spark of life as of yet, but if \textbf{Anacreda} feels threatened she can awaken her creation. The angel is a \textbf{planetar} that is considered a lawful evil Construct. It doesn't have its spellcasting feature. If the characters are here peacefully, review the "Services" section, for a description of what the hags might offer the characters, and ask for in return.

\subsubsection{S8: Irabella's Workshops}
\begin{DndReadAloud}{}
A curious splash draws your attention to an area of the basement dotted with large pools of brackish water. Amidst the pools, a green-skinned woman is partly submerged within a large iron-clad bathtub, set upon a wheeled platform. She is leaning out of the tub, studying various body parts arranged on a table before her. The waters of the brackish pools ripple continuously, as if something large moves beneath their surfaces.

\end{DndReadAloud}

Irabella the hag (Irabella uses the \textbf{archmage} stat block, but her creature type is Fey) seldom leaves her great wheeled bath, that is hauled laboriously about the ramps and elevators of the inn by the staff of kobolds she commands. Her workshops are half-submerged, and here she creates aquatic monstrosities on demand. Great Sahuagin priestkings, aboleth archmages and elder kraken all send emissaries to her, and she makes benthic terrors to trouble the dreams of sea gods. She is sharp and bitter in manner, abusive to her servants and needling to her patrons and sisters, but she also commands the day-to-day running of the inn and its kitchens, tedious work the other two have no taste for. Hence her constant barbs are endured.

\subparagraph{The Pools}
Each pool is 150 feet deep and connected to one another. The creature in the pool is a lawful evil \textbf{dragon turtle} stripped of most of its shell but imbued with infernal energies (Legendary Resistance 2/Day and its Steam Breath deals necrotic damage instead of fire damage). If any pool is disturbed, or Irabella is attacked while in this room, the dragon turtle attacks. During battle, it uses its movement to slip through the tunnels connecting the pools, to position itself tactically, as required.

\subparagraph{Treasure}
Characters find 1d4 Soul Coin, if they search the workshop, though if their pilfering is detected by Irabella, she attacks.

\subparagraph{Phylactery of the Mother}
As soon as the character searching for the phylactery containing their mother's soul enters this room, they hear her voice whisper their name. But her voice sounds like she's trapped under water, and it's distorted. Each hour they may make a DC 17 Intelligence (Investigation) check to try to find her. Only that character can make the check, and no one can help them with it. If they fail the check there is a 1 in 6 chance that Irabella, if not previously encountered, enters the room and confronts them. Once found, the characters are able to retrieve their mother's phylactery from beneath the brackish waters.

The character has taken possession of the phylactery of their mother and gains the Phylactery Benefit associated with it.

\subsubsection{S9: Malecki's Workshops}
\begin{DndReadAloud}{}
Through the gap made between two pillars of bone, you see various large pens, surrounded with walls of bone. A huge beast is chained to the floor in one of the central pens, its head full of tusks and its body a mass of spikes. An elf woman is apparently tending to the creature.

\end{DndReadAloud}

Malecki, who claims to be the direct descendant of Malagard, is a \textbf{night hag} (with 150 hit points) whose work is on a grand scale. Her workshops are huge caverns in the rock, supplied with arches of bone that serve as gates from which her creations can leave, too vast to navigate the physical tunnels to the surface. She calls her work 'bespoke behemoths', living war engines of armor and tusks and spines, and she is happy to deliver them via her gates direct to a battlefield of her patron's choice if required. Malecki appears as plump and cheerful and is absolutely not to be trusted under any circumstance. Her creations are of undeniable power but are best released and then avoided as they've a tendency to forget who their master is soon after being deployed.

Malecki is still working on her latest creation, a yet-to-be-named fiendish war-machine (uses the statistics of a \textbf{goristro} demon). It is tethered by a magical chain with a length of 5 feet (though given the creature's size, it is capable of reaching much further than that). Currently it is docile and apparently tolerates, if not enjoys, Malecki's grotesque surgery. If she is attacked, Malecki uses a bonus action to free it from the chain, and the creature defends her.

\begin{DndSidebar}[float=!b]{Temptation of Deceit}
\textbf{\textit{"Life After Death"}}



Malecki carries a strange amulet in one of her robe's many pockets. This is a cursed item known as the Amulet of Duplicity (see appendix D).



\end{DndSidebar}

\subsubsection{S10: Anacreda the Angelmaker}
\begin{DndReadAloud}{}
A workshop that also appears to be used for the storage of spare creature parts is tucked away on the opposite side of the other basement chambers. Shelves line many of the walls, floor to ceiling, and upon those shelves sit several iron cannisters full of feathers of various plumages and colors.

A hideous winged woman is apparently examining the paraphernalia, though it is unclear whether she is looking for something or simply admiring the collection. Her dove-grey wings twitch, as she glances over her shoulder in your direction. She asks, "Are you here to commission a build, or have you arrived to offer my sisters and I a discount on the used body parts of one of your companions?"

\end{DndReadAloud}

Finally, there is \textbf{Anacreda} the Angelmaker, the leader of the coven. \textbf{Anacreda} (see accompanying stat block) makes beasts of the air and bears a great pair of dove-grey wings herself, a living advertisement of her skills. She is only ever hideous, never bothering to disguise her true visage, but at the same time she has a stern grandeur to her that strikes Fiends and mortals both with its gravity and awe. Nobody has tested her to find out how personally powerful she is, but when \textbf{Glasya} herself comes to visit it is \textbf{Anacreda} she guests with. The two of them sit for hours, drinking peculiar vintages and exchanging secrets. And, while \textbf{Glasya} is the one entity whom \textbf{Anacreda} visibly defers to as her superior, the archdevil treats the ancient hag better than most of her own direct subordinates.

\subsubsection{Services}
As well as making monsters to order, the three sisters keep up a busy trade in exotic components and pieces, part of a network of speculators, harvesters and collectors that spans the planes. In addition, they happily modify Fiends and even mortals. They replace and add limbs, senses, and organs, both on request from those seeking greater personal power, and to fulfill specific punishments for recalcitrant Fiends that \textbf{Glasya} sends to them. Because they represent an unusual little blister of relative lawlessness in the Nine Hells, the sisters have also been known to hide fugitive mortals from infernal pursuit, for a price—though they're also not averse to taking payment and then betraying their guests, if they're feeling treacherous. The three sisters of the Hag's Arms are anything but reliable collaborators save where their business partners are too powerful to play games with.

When mortals do beg their services, the required payment varies. Most often they might simply ask for one member of the visiting group—as a plaything, for their soul, or to trade on to some fiendish creditor. As mortals are oddly reluctant to enter into that kind of transaction, they take other services. There is almost always some rare component one or another sister is after for her latest creation, and they've devised a variety of powerful geas-type bindings to ensure that any agreement isn't reneged on. They particularly like effects that have their errand-runners gradually mutate or deteriorate the longer the task remains undone. In game terms, any Geas they cast works as if it was cast at 9th level. In addition, the geas can only be ended with the \textit{Wish} spell or by killing the hag that cast the spell.

All three sisters appreciate things that are in short supply in the Nine Hells. A bard's skilled performance, some witty repartee and a challenging game of skill might serve just as well as a quest across all the known planes for a piece of an ancient monster still currently within its original owner. Until you go before them, cap in hand, it's impossible to know how hard a bargain the hags will drive.

Risk-taking characters might want the hags to help them survive the dangers of the Nine Hells. Roll on the Hag Parts table... this is the modification available currently (this might change the next time the characters visit Malbolge). If the cost of 3 Soul Coin is paid, one character gains the modification after 2d4 hours of surgery. However, there is almost always a drawback. Roll on the Bad Deal table.

\begin{DndTable}[header=Hag Parts]{cX}

d8 & Result\\
1–3 & The hags can apply a patchwork collection of skin to the character, granting the character immunity to either fire or cold damage (character picks before surgery).\\
4–5 & A third eye can be grafted onto the character's forehead. It offers no benefit, but the hags lie and says it helps to see the future.\\
6–8 & One of the characters' arms can be replaced with a tentacle. The tentacle gives the character advantage on ability checks made to grapple another creature, but the character can no longer use weapons with the two-handed property or perform tasks that require two hands, nor can the character use the tentacle to wield a shield or weapon.\\
\end{DndTable}

\begin{DndTable}[header=Bad Deal]{cX}

d10 & Result\\
1–3 & There are no side effects.\\
4–6 & The adjustments have introduced an aspect of the infernal into the character and they're now vulnerable to radiant damage.\\
7–8 & The character's walking speed is reduced by 5 feet.\\
9–10 & The workmanship was shoddy and uneven. The character has disadvantage on initiative rolls.\\
\end{DndTable}

\subparagraph{Deathstalkers—Obtaining the Scourge of Shadow}
If the characters are searching for the \textit{Scourge of Shadow}, \textbf{Anacreda} wants to make a deal with them. She desires the head of \textbf{Baalzebul}'s chief general, \textbf{Abigor}, in retaliation for a slight he made to her the last time he visited the inn. \textbf{Anacreda} promises to get the scourge, but only if they agree to having the \textit{Geas} spell cast on them. Alternatively, if the characters defeat \textbf{Anacreda} and her sisters in combat, the hag bargains for her life by offering up both the scourge and \textbf{Glasya}. The archdevil visits often, usually without any guard. What could be a better revenge for the characters, than killing the daughter of \textbf{Asmodeus}? The hag is more than willing to help set an ambush for the archdevil.

When \textbf{Glasya} (see appendix A) arrives, \textbf{Anacreda} sits with her and serves a frothing tea that both consume. If \textbf{Anacreda} is cooperating to ambush the archdevil, the tea weakens \textbf{Glasya} in the following ways. \textbf{Glasya} can't use her Call Underling, Fiendish Regeneration, or Mesmerizing Gaze in Malbolge for 1 week. As soon as \textbf{Glasya} finishes her tea, \textbf{Anacreda} beckons the characters from the shadows to attack. During combat, the hag uses her Spinewall to disrupt and corral \textbf{Glasya} and she claws at the archdevil as long as the characters are nearby to assist. \textbf{Glasya} doesn't seem overly surprised at the betrayal—this is the way of the Nine Hells after all.

\chapter{Maladomini, the Domain of Ruined Cities}
\DndDropCapLine{O}{}nce the most beautiful layer in the Nine Hells, Maladomini has since fallen to ruin due to the overzealous actions of its ruler. In an attempt to return his realm to beauty, \textbf{Baalzebul} orders the construction of a new city every few years. When these inevitably don't meet his standard of perfection, he abandons the city and starts anew. The entire layer is now covered in a thick, toxic smog, and sickly black ichor oozes from everything. Hidden behind the smog is a ruby-red sky, blackening the higher it gets. Despite all this, Maladomini remains the only layer with an official school, and many of the intelligent devils call it home.


\section{Running This Chapter}
Before running this chapter read the "Maladomini Overview" section. You can use the information in this section to guide your players through the different attractions this layer has to offer.

\subsection{Encounters}
During their travels across Maladomini the characters encounter some of the devils that call this place their home. Roll at least once on the Random Encounters in Maladomini table whilst here.

\begin{DndTable}[header=Random Encounters in Maladomini]{cX}

d6 & Encounter\\
1 & 1d4 \textbf{erinyes} quietly travel throughout the realm, checking in on construction projects. They're spies from another layer of the Nine Hells, though from which they're unwilling to share.\\
2 & Construction on a new project is being done by thirty lemures, led by a corruption devil (see appendix B).\\
3 & An \textbf{amnizu} (see Monsters of the Multiverse) researcher is exploring the various ruins nearby, cataloguing what it finds for storage in Grenpoli.\\
4 & Three swarms of ayperobo (see appendix B) wander through Maladomini looking for something to fight.\\
5 & Hidden within a nearby ruin is a small camp of six \textbf{spined devil} runaways. They shirked their construction duties and are hiding from their former masters.\\
6 & A belch of foul smog suddenly drops from the sky, covering one of the characters and expanding out from them. Treat the smog as if it were a 9th-level \textit{Cloudkill} spell except that it covers a 200-foot-radius sphere and has a save DC of 21.\\
\end{DndTable}

\subsection{Locations}
If the characters want to explore some of the other locations in Maladomini, the "Key Locations in Maladomini" section will help you set the scene. There is much to see, as Maladomini is forever being developed, with new cities rising every few years that rival the might of the ones that came before. Two of the cities the characters might choose to visit: Grenpoli and Gorloron. The city of Grenpoli serves as a welcome safe haven for many. Within its walls all forms of aggression are forbidden. The city of Gorloron, the latest of \textbf{Baalzebul}'s vanity projects, will no doubt be surpassed by an even more ambitious metropolis in the future. Besides the cities, the characters can explore the Carnival Eternal. There are many lethal thrills to be found within, and as the characters stroll across the carnival grounds, devils cry out to tempt them to go on the carnival's Hellish rides.

\subsection{Objectives}
Make sure you keep track of your player's objectives and lead them to the corresponding areas to ensure they can complete their goals. Once done exploring, they need to make their way to the Eye Market. Here they can find some of the items and souls they're searching for.

The following objectives can be attained in Maladomini:

\begin{DndTable}[header=Objective in Maladomini]{XX}

Objective & Location\\
Phylactery of the great con & Eye Market: Ganglands (Memnoriac's territory)\\
Phylactery of the spouse & Eye Market: Ganglands (Ilkatar's territory)\\
\textit{Staff of the Woodlands} & Eye Market: The Slums\\
\textbf{Jenevere} & Eye Market\\
\end{DndTable}

\subsection{Temptations}
Have the characters encounter at least one temptation during their time in Maladomini. If the characters give in to temptation, use the information in appendix E to keep track of their corruption level.

\section{Maladomini Overview}
The River Styx enters the realm through a massive canyon that stretches to Malbolge. The river moves from the canyon into a massive subterranean cave system, where it avoids the exploitation of the surface for some time before emerging from underground and exposing the smog covered sky to those sailing its waters.

During its heyday, Maladomini was a beautiful paradise, with varying biomes and lush vegetation. Sadly, the only forests that remain are either decaying or dead, with no wildlife or beauty to be found. In between the deforestation lie strip mines, surface mines, and processing facilities that belch smog and toxins into the air. Any water that once existed has turned to sludge, save only for the secluded Styx. The layer is truly a testament to \textbf{Baalzebul}'s failure in his search for perfection.

Each of the ruins covering the realm is more spectacular than the last. During their construction, \textbf{Baalzebul} wanted each to be his ultimate metropolis and spared no expense. Even in their detritus they remain grotesquely beautiful. Occasional treasure hunters from other planes explore the ruins in search of treasure, and many escaped devils or souls hide deep inside. The overexploitation of the surrounding lands makes the ruins some of the safest spots in the layer, and they remain largely forgotten by the ruling archdevil.

Despite his disloyalty to \textbf{Asmodeus}, \textbf{Baalzebul} remains the archdevil in charge of Maladomini. His rule is assisted by other powerful archdevils, including \textbf{Abigor}. \textbf{Baalzebul} is driven towards perfecting his realm and many construction projects are constantly underway. Devils not working on the various construction projects prefer life in the domed city of Grenpoli.

\subsection{Leaving Maladomini}
The only known way in or out of Maladomini is the Styx, though rumors claim other entrances and exits exist. Because of the stream of devils entering and exiting the realm for training and schooling, scholars believe there to be a hidden infernal passageway that connects the layer to the rest of the Nine Hells. Explorers also claim that the cave system the Styx passes through has connections to other layers, though such connections have yet to be found. Besides, if one can brave the cave systems, the Styx is a relatively easy method of travel.

\subsection{Features}
Although Maladomini remains one of the safer layers, the influence of the Lord of the Flies remains powerful. Whenever a creature is asked a question, they must make a DC 14 Wisdom saving throw. On a failed save, the creature is magically compelled to tell a lie, stretch the truth, or avoid answering—anything that would prevent the truth from being told.

\section{Key Locations in Maladomini}
The characters need to make their way to the Eye Market but may want to spend some time in the other locations found in Maladomini. Use the information below to guide their excursion.

\begin{DndSidebar}[float=!b]{Temptation of Betrayal}
\textbf{\textit{"If I Could Read Your Mind"}}



The characters are approached by a devil that offers them the ability to read the thoughts of someone important to them for the cost of 1 \textit{Soul Coin} and 800 gp. If one of the characters is tempted, run the appropriate temptation event from chapter 2.



\end{DndSidebar}

\subsection{City of Grenpoli}
Unique to Maladomini and known throughout the Nine Hells is Grenpoli, the City of Diplomacy. All forms of aggression or weapons are forbidden within the city, and a powerful spell suppresses violence when it would occur. Upon entering the city through one of its four gates, visitors are searched, and their weapons confiscated. These will only be returned once their visit to the city has ended. Grenpoli contains the Political College of the Nine Hells, which offers devils courses in deception, manipulation, and negotiation. Located just outside the city is Offalion, a training facility that enacts large-scale scenarios to train devils, testing and enhancing their ability to barter effectively with mortals. These exercises prepare a multitude of devils for delicate excursions into the Material Plane.

\subsection{City of Gorloron}
\textbf{Baalzebul}'s "greatest" city of the moment is Gorloron, which contains his seat of power and much of the devil population in the realm. The city's construction is still underway, and the place is utter chaos.

\begin{DndSidebar}[float=!b]{Temptation of Oppression}
\textbf{\textit{"A Palace Fit for a King"}}



Here in the (current) capital of \textbf{Baalzebul}'s realm, you can't throw a stone without hitting an architect. All are vying to gain favor from the tormented archdevil and to convince him to let them build bigger and better. One of these architects approaches the characters, explaining that his latest proposal was rejected once more. It is desperate for work and offers to build the characters a fortress or a temple in their name. The price is 3 Soul Coin and 2,000 gp (a rare magic item brings the price down to 500 gp). If any of the characters take the devil up on their offer run the appropriate temptation event from chapter 2.



\end{DndSidebar}

\subsection{Carnival Eternal}
Also of note is the Carnival Eternal, which is a terrifying and absurd amusement park built to reward successful devils. If the characters have not yet encountered the "Ruin and Amusement" random encounter (see chapter 2), feel free to run that now if the characters explore the Carnival Eternal.

\subsection{The Eye Market}
This dangerous area is nestled in the center of a ring of partly abandoned buildings, all of it the turf of various factions vying for control.

\section{Adventure: The Eye Market}
The Eye Market is run by an odd couple—a mind flayer and a beholder. For the right price, they will replace the eyes of a customer with new jeweled eyes of great magical power.

If asked about the Eye Market, \textbf{Koh Tam} has this to say:

\begin{DndReadAloud}{}
"Maladomini is the layer of deceit but none are more deceived perhaps than its master, \textbf{Baalzebul} himself who ruins every part of it in a vain attempt to create a wonder of beauty that will put the other planes to shame. Constantly beginning and then abandoning his increasingly ambitious projects, he leaves behind a polluted waste of dead forests, toxic water, and the shells of half-completed palaces, castles and cities, each abandoned as it fails to meet its architect's impossible standards." He gestures at the debris of the various unfinished projects that litter the landscape.

"In one particularly expansive cityscape comprised of a labyrinth of roofless houses, half-raised towers, and foundations of great works, there in the shadow of a ruined dome surely too grand ever to have been completed, squats the Eye Market, the underside of Maladomini. A festering secret hidden under the nose of the layer's archdevil, a siren's lure to all those exhausted by the hideous waste of the layer. Because there is salvation to be had there, of a kind. If the sights of Maladomini offend your eyes, the devils say, then pluck them out. My intuition tells me you should venture to the Market next."

\end{DndReadAloud}

\subsection{Advice from Koh Tam}
\textbf{Koh Tam} warns the characters to be careful while on the streets around the Eye Market. He can also offer the following advice:


\begin{itemize}
\item The Eye Market is the name for the entirety of the neighborhood surrounding the specialty store in the inner district operated by strange beings known as Vaness and Fling. It is also the name of their store.
\item Rival gangs of devils roam and control various portions of the Eye Market, and it would be best to not get mixed up in their turf war.
\end{itemize}

This is a location well known to many travelers and \textbf{Koh Tam} can provide a standard map of the Eye Market (see appendix F) to the characters. They can then use this map to decide which route they want to take to reach Vaness' and Fling's shop (area M6).

\subsection{The Law}
Watching over the streets around the Eye Market is Akrekarn, the closest this morass has to the law. A corruption devil (see appendix B) retired from a life of blackmailing mortals, Akrekarn patrols the streets around the market with a writ from \textbf{Baalzebul}, allegedly enforcing the archdevil's will but in truth simply shaking down locals and visitors for all he can get. He is more than ready to have his mob of infernal deputies (six bearded devils) seize anyone who won't bow the knee or proffer a bribe, taking them to be tormented, sold, or ransomed as he thinks best. When not terrorizing the streets, Akrekarn takes his unspeakable pleasures in the one completed building in the city, a guard house he has made his headquarters. Or else he is to be found at the Eye Market itself, shaking down the proprietors and reminding them that they operate only on sufferance.

\begin{DndSidebar}[float=!b]{An Eye for Opportunity}
There are watchful eyes everywhere and the moment characters attack either Ilkatar, Memnoriac, or Jacelisk, Akrekarn is informed. Before the battle finishes, he arrives in an infernal airship, a gift from \textbf{Baalzebul}, accompanied by his bodyguard of six bearded devils. He doesn't interfere with the battle but if the characters are victorious, he takes control of the vacated gang-leader's territory and won't permit passage until he's gotten approval from \textbf{Baalzebul}. Akrekarn makes it clear that attacking him will bring on the archdevil's wrath.



If the characters are defeated, he takes half of the character's belongings from the gang-leader, to store in his headquarters. Details for the airship can be found in the Airborne and Waterborne Vehicles table in the \textit{Dungeon Master's Guide} except that Akrekarn's vehicle requires half the crew and accommodates only half the number of passengers.



\end{DndSidebar}

\subsection{Eye Market Locations}
The characters must navigate dangerous streets to get to the Eye Market. However, if they keep their wits about them, they may just find what they're looking for.

\subsubsection{M1: Outer Streets}
\begin{DndReadAloud}{}
Like the great corpse of a whale sunken to the sea floor, the city around you is at once both utterly dead and teeming with life—scavengers and parasites all contending for a scrap of flesh. You see them in the darkened corners of winding alleyways and the entrances of crumbling doorways. Watching. Waiting. Below your feet the pavement is potholed and hazardous and dilapidated structures loom all around like giant bones picked clean, haunted by the revenant of intended greatness.

\end{DndReadAloud}

Reaching the market is easier said than done. The convoluted streets of the abandoned city it nestles within are fraught with peril around every turn. Desperate lost souls and minor devils, cut loose from the hierarchy of the Nine Hells, throng there—most bitterly seeking the unique service the Eye Market can provide but unable to pay the fee. Every hollow shell of a building hides a nest of wretches preying on each other, scrabbling in the mulch, or descending into the monster-haunted and unsound sewers of the place in the hopes of finding something of value.

\begin{DndSidebar}[float=!b]{Dangerous Territory}
There are dozens of smaller gangs in the broken streets around the Eye Market, but the current major players are Ilkatar's Brood, Jacelisk's Knives, and Memnoriac. To reach the Eye Market (area M6), the characters must walk through, or fly over, one of the ganglands. Each major gang demands tribute to allow the characters passage through their territory. See the following descriptions about what each gang demands but the head of a rival gang-leader is also suitable tribute for any gang.



If any characters attempt to fly over gang territory, that gang challenges them in the air instead. Any descriptive text in the following sections will need to be adjusted if this occurs.



\end{DndSidebar}

A variety of gangs preside over this layer of scum, extorting prices of souls, flesh, and labor from those on their turf. This being the Nine Hells, of course, it isn't a dog-eat-dog chaos, but a knifeedged dance of diplomacy, custom, and expectation. Every gang is the heart of a web of alliances and carefully negotiated territorial arrangements, all of which contain hidden get-out clauses and exceptions that allow them to be shrugged off at any moment. The devils live strung between two poles, constantly living as though their word is their bond, yet ready to exploit loopholes to betray their fellows at every step.

\subsubsection{M2: Ganglands—Ilkatar's Territory}
Ilkatar's people are abishai and pay at least lip service to Tiamat, in her lair far above. Of all the gangs, they value mortal treasure most, hoarding gold and gems and tithing to their great mother. They're by nature cruel but lazy, easily appeased with gifts or service.

\subparagraph{Tribute}
If tribute is paid, the characters may head north and around the inner district to area M3 or enter the inner district and reach the "Eye Market," bypassing any other gangs. If the characters seek out Ilkatar to pay this tribute, proceed to the following:

\begin{DndReadAloud}{}
You enter a large, half-finished warehouse turned dwelling, a dimly lit room filled with the sounds of snarling and hissing. Through the lingering smoke you discern a poorly made wooden throne and upon it, a green abishai. To the left of the throne stands a white abishai, clad in gleaming armor and holding a wickedly sharp sword. On either side of the throne, a pair of hell hounds snarl and growl, their fur matted with dried blood, and their teeth gleaming. The green creature speaks, "I am Ilkatar and this is my Brood. Have you come to pay tribute for passage, or to feed my very hungry pets?"

\end{DndReadAloud}

To pass through this territory the characters need to pay 7,500 gp or deliver a rival's head. If they choose to fight their way through, they're challenged by Ilkatar (\textbf{green abishai}), the \textbf{white abishai} lieutenant (see Monsters of the Multiverse for the abishai stat blocks), and four hell hounds.

\subparagraph{Treasure}
If the characters slay Ilkatar and succeed on a DC 22 Wisdom (Perception) check, they find a discarded lockbox under the wooden throne. Opening it reveals 3,500 pp.

\subparagraph{Phylactery of the Spouse/True Love}
The character who has ventured into the Nine Hells to rescue their true love finds them imprisoned in one of the buildings in Ilkatar's territory. The soul is tormented in a twisted hall of mirrors. The whole interior of the building is made of mirrored walls, floors, and ceilings.

The soul is trying to find their true reflection among the hundreds of reflections around them. An impossible task, and painful and deadly every time they guess wrong.

To save them, the character needs to find their true love among the many reflections. Have the character make an Intelligence (Investigation) check. Look up their result on the Mirror Reflections table. Once the character has dealt with the outcome of a failure they can try again.

After finding their true love's reflection, it vanishes leaving the phylactery in its place. The character takes possession of the phylactery of their true love and gains the Phylactery Benefit associated with it.

\begin{DndTable}[header=Mirror Reflections]{cX}

Result & Outcome\\
3 or less & The reflection vanishes, meanwhile three \textbf{white abishai} (see Monsters of the Multiverse) step out of nearby mirrors, surrounding the character. If the character gets this result on their fourth or later roll, they find their true love instead.\\
4–6 & The mirror shatters and the shards cut the character. The character must make a DC 16 Dexterity saving throw, taking 21 (6d6) piercing damage on a failed save, or half as much damage on a successful one.\\
7–9 & The reflection vanishes once the character reaches out for them. If the character gets this result on their fourth or later roll, they find their true love instead.\\
10–12 & The character is pulled into the reflection by two \textbf{black abishai}. If the character gets this result on their third roll, they find their true love instead.\\
13–15 & A \textbf{red abishai} lunges forward from the reflection and attacks the character but vanishes at the end of its second turn of combat. If the character gets this result on their second roll, they find their true love instead.\\
16 or more & They find their true love.\\
\end{DndTable}

\subsubsection{M3: Ganglands—Memnoriac's Territory}
Memnoriac was once one of \textbf{Baalzebul}'s architects, a \textbf{pit fiend} of great status and power. Like all such, his perceived 'failures' eventually led to his downfall, and now he squats with his ragbag of followers in the city he designed and will never complete. The single most powerful devil in the ruin, he is tormented by the collapse of his final project and has become a \textbf{Baalzebul} in miniature. Those under his control are constantly re-edifying the buildings around him, raising them up and tearing them down as he seeks his own minuscule perfection.

\subparagraph{Tribute}
If tribute is paid, the characters may head west and through the inner district to reach the "Eye Market," bypassing any other gangs. If the characters seek out Memnoriac to pay this tribute, proceed to the following:

\begin{DndReadAloud}{}
You step through open gates beneath the skull of a massive creature. A short walk past the gate, in the center of the courtyard, a muscular creature looms over a table upon which a haphazard model of the city-in-miniature rests. The model itself bears little resemblance to the real city it represents, yet somehow seems even more disorganized. Without even looking up at you, the fiend snarls. "Welcome to Memnoriac's grand city. Have you come to bow and worship, to scrape your knees on the rough gravel, and pay homage to me, the master-builder?"

\end{DndReadAloud}

To pass through this territory, the characters must pay 20,000 gold in gems and coins to fund Memnoriac's construction efforts. Alternatively, they can offer something even more valuable: a compliment. Congratulating Memnoriac on his fabulous city and then succeeding on a DC 22 Charisma (Deception) check earns the Fiend's goodwill and safe passage. Failing the check, or refusing to pay tribute, results in the pit fiend fighting the characters on his own. He refuses to believe he needs help defeating such puny adversaries. If slain, the Fiend has no treasure, just diagrams of his future building plans.

\subparagraph{Phylactery of the Great Con}
The character looking for their soul can find it among the lowlifes and grifters that inhabit the streets surrounding Memnoriac's skull-gated manor. In other circumstances they would have been right at home here, conning their way to the top of the devil food chain. However, the character's soul isn't in the position to con anyone. It is being tormented by the gang that runs this territory. Memnoriac and his troop of devils have taken up root here after the pit fiends' fall from grace. His obsession with perfecting and completing the city around him, sees him constantly changing the makeup of streets and buildings.

The character searching for this soul hears a voice calling for help from a nearby alley. If the characters investigate, read the following:

\begin{DndReadAloud}{}
You see a pit has opened in the alley, almost ten feet in radius and large enough to create a barrier between your side of the alley and the five devils opposite you. They have cornered a trembling and wingless imp, its back to you... and the pit. Its small body trembles in fear and as it takes a step back, it almost falls in.

\end{DndReadAloud}

The character realizes that this imp is the vessel for their soul and that they must rescue it from a \textbf{chain devil}, a \textbf{bone devil}, and three bearded devils. The devils immediately attack. On each of their turns, one of the devils tries to hit the imp while the others attack the characters. Each time the imp dies it comes back to life on initiative count 20 of the following round but another hostile bearded devil also materializes at the same time. The imp has the frightened condition and is afraid of both the devils and the characters and on its turn, it attempts to flee from both.

To free the imp from its torment, the characters must kill all pursuing devils before anything kills the imp. When this occurs, the imp vanishes, leaving the phylactery in its place.

\subsubsection{M4: Ganglands—Jacelisk's territory}
Jacelisk is a recent addition to the landscape, a tiefling archmage who brought her band of wicked companions to the Market for its particular services. Now, under the influence of recently installed eyes, she sees herself as the liberator and protector of the city she resides in, a great and adored hero. In truth she is as vicious as all the rest, with the added problem that her perceptions make her profoundly unpredictable—sometimes nobility itself and sometimes brutal enough to impress a full-blooded devil.

\subparagraph{Tribute}
Jacelisk has claimed prime gang real estate. If tribute is paid, the characters may head directly north, taking the shortest route to reach the "Eye Market," bypassing any other gangs. If the characters seek out Jacelisk to pay this tribute, proceed to the following:

\begin{DndReadAloud}{}
A reinforced gateway has been hastily constructed between districts. As you approach the gate a tiefling woman wearing blue robes, emblazoned with numerous tiny, iron, knives, emerges from an azure tent set alongside the gate. Three spike-covered devils follow her, alternating fawning glances in her direction, and hostile gazes your way. The tiefling's eyes seem to be too large for her face and they glow with a strange amber hue as they stare at you. "Welcome, travelers. I am Jacelisk, the Great Liberator, and these are my Knives." She glances at her devil bodyguard, as well as towards many other devilish forms scuttling in the shadows. "Where might you think you are headed?" Somehow, impossibly, her eyes seem to regard each of you at the same time.

\end{DndReadAloud}

Given that her route is the shortest to the Eye Market, and observing that the character's carry many interesting valuables, Jacelisk demands a rare or better magic item and at least four uncommon or better magic items as tribute. If attacked, Jacelisk (\textbf{archmage} with resistance to fire damage, the ability to cast \textit{Hellish Rebuke} twice a day, and the ability to use her Eyes of Charming one time), tries to keep her three barbed devils in between herself and the characters while she unleashes deadly magic. She has already cast \textit{Mage Armor} and \textit{Stoneskin} on herself. If the characters have attempted to fly past this area, they're intercepted by the barbed devils—Jacelisk has cast the \textit{Fly} spell on them.

If slain, the characters may remove Jacelisk's eyes, and if returned to Vaness, these Eyes of Charming can be sold back for half their indicated cost.

\subsubsection{M5: The Slums}
\begin{DndReadAloud}{}
Just walking the streets of the slum is a struggle. You must watch your step carefully to avoid tripping over heaps of refuse and random construction equipment and to avoid crashing into one of the countless lesser devils that swarm around, calling out in shrieking, whining voices to you and to each other. The air is heavy with a fetid stench and hot from the close press of desperate bodies. The streets are narrow, lined with miserable little attempts at shelter, shacks and sheds of sheets of rusted metal and rotting wood, leaning on each other in a despairing attempt to stay upright.

\end{DndReadAloud}

Closer to the shattered dome of the Eye Market itself is a whole town within a town, a wretched slum of shacks and sheds thrown up with materials salvaged from the abandoned construction works. Here is a morass of pathetic minor devils and travelers who mob any newcomer, desperately trying to sell them all manner of trash—broken weapons, ruined shreds of garments, cracked mirrors, the most meagre of worthless costume jewelry. To the uninitiated the whole desperate circus seems utterly inexplicable. Many will trade junk for equally valueless junk, rubbish passing hands in a constant round of meaningless commerce. Their grotesque parody of trade seems readily exploitable by anyone coming in from the outside, but there is a curious malaise on the streets around the Market. Deal too much with the trash hoarders and you might find yourself getting sucked into their game, desperate to make the next big trade or to acquire some filthy piece of detritus, valuing the junk just because everyone around you does the same.

A remarkable number of the denizens of this slum are damaged, victims of \textbf{Baalzebul}'s construction works where safety is never a concern. Crooked and whining imps and other lesser devils roost or slither everywhere, proffering fistfuls of slimy muck, desperate to get in on the ground floor of the city's nonsense economy.

\subparagraph{Conclave (Druid)—Staff of the Woodlands}
When the characters walk the slums, they're beset by lowly devils trying to sell them their useless trinkets. If a character is looking for the Staff of the Woodlands their attention is drawn to a miserable looking imp who has set up shop in the corner. When they inspect its wares, they find the staff among them. The imp has no idea of the value of the item it possesses. The character needs to succeed on a DC 16 Charisma (Deception) check to purchase it for a reasonable price.

Upon success the imp sells the staff for a mere 200 gp or 1 \textit{Soul Coin}. It is so happy with the transaction that it throws an amulet of protection in for free. The amulet is a fake and won't protect the character. However, if the character wears the amulet, it leaves a green stain on their skin that can only be removed after 6 days. On a failure, the imp notices the character's eagerness to buy the staff and increase the price to 1,000 gp and 3 Soul Coin.

\subsubsection{M6: The Eye Market}
\begin{DndReadAloud}{}
Inside the Eye Market is a workshop of peculiar arrangement. There are racks of tools along the back wall, scalpels, knives, and needles intermixed with the equipment of the jeweler's trade. At the front of the Market there is a grand glass case and, within it, small orbs of countless colors, carved from precious stone, are nestled in satin lined boxes. They seem to gleam with inner light.

\end{DndReadAloud}

The proprietors, Vaness and Fling, offer a service unique in all of Maladomini—in all the Nine Hells, in fact. They offer relief from the appalling oppression of the infernal landscape. They offer a glimpse of happiness and wonder, positive things, that benefit Fiends and mortals alike. A bizarre gem to find in such a corrupted setting. Except, of course, nothing in the Nine Hells is as good as it appears.

Vaness is a \textbf{beholder} of exceptional size, his pinkish hide puffed-up and corpulent. His central eye has been replaced by a lens of ruby quartz that, it is said, focuses magic rather than negating it. In his workshop, using telekinesis and extremely fine-scale disintegration, he crafts the eyes that are the Market's stock in trade. Fling, in contrast, is an \textbf{alhoon} (see Monsters of the Multiverse), a mind flayer long gone on their peculiar path to lichdom. It acts as surgeon, implanting Vaness' creations and connecting them to the brains of its patients/ customers. They make a strange and jarring pair, not least because of their manner. Vaness and Fling strike the unwary as an old married couple, constantly bantering back and forth, making jokes at one another's expense, and exploding into arguments that settle almost instantaneously. It seems impossible that either creature could actually form a real bond, but to all intents and purposes the pair seem genuinely attached. Which often serves to lull visitors into a false sense of security. Vaness and Fling have not survived the Nine Hells for so long by being nice. They're as ruthless and powerful as any other infernal magnates, able to defend themselves physically if need be, but preferring to rely on their minions, treaties with the gangs and, if necessary, protection from Akrekarn and his posse. Moreover, anyone going to deal with them to actually gain new eyes should be very careful what they're agreeing to. Vaness and Fling have a flare for crooked bargaining that some devils could learn from.

\begin{DndSidebar}[float=!b]{Temptation of Betrayal}
\textbf{\textit{"Two Kills for One"}}



Upon defeating Fling, the characters find he carries a magic weapon. Run the appropriate temptation event from chapter 2.



\end{DndSidebar}

\subparagraph{Hellriders—Rescuing Jenevere}
More information about \textbf{Jenevere} can be found in appendix C. Characters looking for \textbf{Jenevere}, find the planetar imprisoned in Vaness' workshop. As the characters enter, they notice her caged in the back of the room.

On their travels around the Nine Hells, the characters may have heard the rumors surrounding \textbf{Jenevere}. It is said she has never been broken, not by any devil. She pains her tormenters by immediately forgiving them for their abuse, praying for those souls that inflict torment on her. She believes no one—not even the vilest of devils—is beyond saving.

The proudest Fiends think they're to be the one to break her, corrupt her. They fight over her. But possessing the planetar becomes tainted once they realize she can't be broken.

And now, imprisoned here in the workshop, \textbf{Jenevere}'s magnificent focus is completely on Vaness. She stares intently in Vaness' direction, even though her radiant eyes were taken from her by her cruel captor. It is as if she is looking straight into the beholder's soul.

She doesn't even notice the characters approaching the cage, but as they do, they hear \textbf{Jenevere} talking. She is lecturing the beholder about forgiveness and urging him to surrender to the grace of the divine light. Vaness has clearly heard it all before and with a frustrated snarl he flings the tool he has been using at the cage. It bounces off \textbf{Jenevere}'s head and clatters to the floor. But \textbf{Jenevere}'s concentration doesn't even falter, as she continues lecturing the put-upon beholder.

This little exchange should inform the characters that Vaness is growing weary of his captive and might be willing to barter with them for her release. A successful DC 15 Charisma (Persuasion) check convinces Vaness to let them take \textbf{Jenevere} with them, if they let the beholder hang on to the planetar's eyes. If the characters insist on taking \textbf{Jenevere}'s eyes with them (or if they fail the check) he refuses them.

Once he refuses them, he tells them to leave. At this time the characters can either attack the beholder or leave and try to come back when the workshop is empty. Both scenarios end up in combat as Vaness is never far from his workshop. Any intrusion results in the beholder attacking the characters. After 5 rounds of combat, Fling comes to Vaness' aid, accompanied by two \textbf{grell} and two grimlocks. After defeating the two proprietors of the workshop the characters can free \textbf{Jenevere}.

If the victorious characters decide to look for \textbf{Jenevere}'s eyes, a successful DC 20 Intelligence (Investigation) check reveals them floating in a container nearby. Once \textbf{Jenevere} is freed, the characters may use the ritual that the Hellriders taught them to free her from the Nine Hells.

\subsubsection{The Jeweled Eyes}
In their workshop, assisted by a blind workforce of \textbf{grell} and grimlocks enslaved by Fling's mental powers, Vaness and Fling remove the natural eyes of their clients and replace them with the jeweled creations from Vaness' workshop, granting a whole new perspective on Maladomini. A simple thing, in a way, and yet priceless. Vaness' eyes allow his patients to see the layer the way it was meant to be. When a newly altered customer wakes from the surgery, the dome above them is complete, its ceiling ornate and painted with frescoes of surpassing loveliness. Vaness and Fling are elegant and dignified, the master-surgeon and the genial artificer, surrounded by their immaculate staff. Outside, the spires of the city reach towards a glorious golden sky, and beyond that there are great verdant forests, jeweled castles that would melt the heart of a Fey queen, a world of glory and wonder such as the Nine Hells have never known.

The true power of the eyes is that these things aren't illusions, to those who can see them. Those towers can be scaled, the woods bear fruit that truly sates hunger. Instead of the poisonous sludge of Maladomini, the streams now run with sparkling claret liquid that tastes like the finest of wines and brings instant relief and freshness to the drinker. To live in the plane visible through Vaness' eyes is to be truly blessed.

The Eyes for Sale table lists the different types of eyes that can be purchased. These eyes do not require the wearer to attune to them, even if they normally would. In addition to the abilities mentioned, each set of eyes also makes the wearer no longer require food or water (because it is always abundant for them), and the weather affecting them is always pleasant (never too hot or too cold; they're immune to nonmagical fire or cold damage). Finally, no terrain is ever considered difficult for them. These effects persist even after the wearer leaves Maladomini.

\begin{DndTable}[header=Eyes for Sale]{XcX}

Magical Power & Requirement & Cost\\
Eyes of the Archer—When you use a ranged weapon to make a long ranged attack, you don't have disadvantage on the attack roll. & 9th-level character & 10,000 gp\\
\textit{Eyes of Charming} (see \textit{Dungeon Master's Guide}) & None & 5,000 gp\\
\textit{Eyes of the Eagle} (see \textit{Dungeon Master's Guide}) & None & 2,000 gp\\
Eyes of Far Sight—the recipient can cast \textit{Clairvoyance} at will. & 12th-level character & 25,000 gp\\
\textit{Eyes of Minute Seeing} (see \textit{Dungeon Master's Guide}) & None & 3,000 gp\\
Eyes of True Seeing—The recipient gains truesight out to a range of 30 feet. & 15th-level character & 75,000 gp\\
\end{DndTable}

\subsubsection{The Truth Is Ugly}
Of course, there are drawbacks. The ancient and wondrous beasts that once roamed the forests of Maladomini are long extinct, but they're most certainly there to those with eyes to see. And this being the Nine Hells, they're frequently hostile. Many eye-bearers have been killed and devoured by glorious and majestic monsters that their companions can't even see, let alone fight. Moreover, everyone who leaves the Eye Market having conducted business with Vaness and Fling becomes a target. All those waiting outside would give a great deal more than their eyes to get hold of such things, even though the services of Fling would still be required to implant them. Overcome with wonder, many new patients are easy marks for muggers and back-street surgeons.

The chief question about all of this must be why \textbf{Baalzebul} puts up with any of it. Vaness, Fling, and Akrekarn maintain a careful charade where the two put-upon artisans are constantly under threat from the devil, about to be turned in to the layer's archdevil unless they can come up with the latest bribe or protection instalment. And yet it never quite happens, and somehow the peculiar pair are permitted to run their business, which seems to make a mockery of everything that \textbf{Baalzebul} is constantly, futilely trying to achieve in his domain.

The most persuasive reason is that the eyes, for all their vaunted promise, are a curse. The wearer is at first overwhelmed by the beauty they see, delighted to walk in a 'Hell' that is suddenly restorative and pleasant. Except the truth always creeps in. Those who wear the eyes can stave off reality only so far. In their dreams they see the truth of it, the ruin and the ugliness. Day by day the understanding grows that everything they see, that they could now never live without, is only a sham, however real it is to them. The knowledge torments them, and they're driven to their own hopeless attempts to achieve a lasting perfection. In this way they share in \textbf{Baalzebul}'s private misery, and perhaps that is enough reason to keep Vaness and Fling around.

A seditious whisper suggests that there is another reason. That \textbf{Baalzebul}, lord of temptation, feels the hook of their services in him. The idea that he could see Maladomini as he wishes it to be, the great blueprint he constantly fails to build, eats him up constantly. So that he is constantly on the point of ordering Akrekarn to obliterate the Market, but never quite gives the order, because what if...

Whenever a roll on any encounter table is called for, a character possessing eyes from the Eye Market must also roll on the following encounter table. If multiple characters possess these cursed eyes, they all share the same result. If any of these encounters result in combat, only the characters with cursed eyes can see their attackers or target them, and only they are harmed by the abilities of their attackers. Other characters may only provide support, such as healing, to the afflicted characters. Any characters that die outright, or are reduced to zero hit points, in these battles, immediately have their eyes removed and the attackers disappear.

\begin{DndTable}[header=Vaness' Eyes Random Encounters]{cX}

d10 & Encounter\\
1 & A \textbf{blue abishai} (see Monsters of the Multiverse) leads a half-dozen \textbf{black abishai}, cloaked in \textit{Darkness} spells, in an ambush, hoping to steal the eyes.\\
2 & A primordial glittering beast (\textbf{adult black dragon} with no flight speed) hunts anyone who possesses magical eyes.\\
3 & Twisted vegetation emerges and grows around the afflicted character and a dozen strange creatures (nothics with a burrowing speed of 30) crawl out from under the plants. They attack anyone with magical eyes.\\
4 & A beautiful bird the size of a building (roc with legendary resistance 3/Day) hunts anyone who possesses magical eyes.\\
5–10 & A glorious vision enthralls anyone that possesses magical eyes. Each of them must succeed on a DC 20 Wisdom saving throw or have the incapacitated condition for 1d12 hours as they stare at the vision. When they wake, they gain the benefit of a long rest.\\
\end{DndTable}

\chapter{Cania, the Relentless Cold}
\DndDropCapLine{W}{}ith freezing temperatures and constant blizzards of snow and ice, Cania is the least habitable layer of the Nine Hells. Unlike Stygia, which is a sheet of ice sitting atop a frozen sea, Cania is a massive glacier hundreds of miles deep. Its cold is magical in strength, completely ignoring all protection and sources of warmth that aren't magical themselves, and capable of freezing a mortal to death in only a few hours. Nobody knows what Cania's sky looks like, as it hides behind a white sheet of clouds created by freezing weather.


\section{Running This Chapter}
Before running this chapter read the "Cania Overview" section. It provides you with the necessary information as the characters come close to ending their journey down the River Styx. \textbf{Tiax} may have a special role in this chapter, so review the "\textbf{Tiax}'s Betrayal" section later in this chapter. Ultimately, if the characters have completed their main objectives, the "A Thankful Patron" event occurs, setting the stage for the resolution to the adventure in chapter 12.

\subsection{Encounters}
This is the last layer of the Nine Hells in which the characters can use \textbf{Koh Tam}'s barge. Their final journey takes them through the frozen landscape, diverting down a tributary of the Styx to the Sorrow Mine. Roll at least once on the Random Encounters in Cania table, near the massive waterfalls or after they divert towards the Sorrow Mine.

\begin{DndTable}[header=Random Encounters in Cania]{cX}

d6 & Encounter\\
1 & Sudden winds whip fallen ice into a veritable blizzard. Making matters worse, a nearby iceberg gets sheared, causing pieces to fall to the ground. Each creature within the blizzard must make a DC 18 Constitution saving throw, taking 21 (6d6) cold damage on a failed save, or half as much damage on a successful one.\\
2 & A pack of three vorvolakas (see appendix B) circle in the skies above the icy surface, searching for mortal blood. Canian devils call them "ice vultures." Once they find a victim, they stalk and wait for the temperature to make things easier, staying in the air and just out of range until their prey weakens.\\
3 & A young \textbf{Styx dragon} (see appendix B), cut off from the river by falling ice, sought shelter in a nearby ice cave. Without intervention, it will likely die of starvation or frostbite in a few days. It is both defensive and hungry, but it recognizes the position it is in.\\
4 & Halogs (see appendix B) are known to resist the fires of Phlegethos, and as a swarm of four of them runs by, it becomes apparent they can also survive the freezing temperatures in Cania. They could provide a good source of food—if not for the characters, then for the devils native to this layer.\\
5 & Scouting out the ever-shifting landscape of the layer is a small group of three ice devils. Who they work for—themselves, Cania's leader, or some other power—remains a mystery. Depending on their employer and their attitude, they could see mortals as worthy of a fight.\\
6 & The frozen ground cracks and howls, opening up a ravine in preparation for a coming iceberg. Each creature in a 20-foot-wide line, a mile long, must succeed on a DC 17 Dexterity saving throw or stumble into the opening ravine, falling 200 feet and taking 70 (20d6) bludgeoning damage. The ravine could lead anywhere—a subterranean system, ruins frozen in time, or just an empty gap. If this encounter occurs while aboard \textbf{Koh Tam}'s vessel, part of the River Styx begins to flow into the chasm and he's hardpressed to steer past safely.\\
\end{DndTable}

\subsection{Locations}
The characters must journey to the Sorrow Mine, once they've learned of its existence from either \textbf{Tiax} or more magical means. The "Key Locations in Cania" section provides the necessary information for you to guide the characters' exploration.

\subsection{Koh Tam and Tiax}
\textbf{Koh Tam} suggests avoiding being noticed by \textbf{Mephistopheles}, fearing the archdevil's unwanted attention. Unfortunately, \textbf{Koh Tam} is less familiar with Cania. If \textbf{Tiax} is still with the party, he speaks up, for he has heard of the famed Sorrow Wine extracted from the Sorrow Mine. He suggests traveling there. He wants to accompany them everywhere but behaves oddly–at each new location he should be grumpy and insulting, but then promptly apologize when called out on his behavior.

\subsection{Objectives}
Make sure you keep track of your player's objectives and lead them to the appropriate areas so they may complete their goals.

The following objectives can be attained in Cania:

\begin{DndTable}[header=Objectives in Cania]{XX}

Objective & Location\\
Phylactery of the merciless & Sorrow Mine: The Mining Town\\
Phylactery of the student & Sorrow Mine: The Mines\\
\textit{Staff of the Magi} & Sorrow Mine: The Mines\\
\end{DndTable}

\subsection{Temptations}
Have the characters encounter at least one temptation during their time in Cania. You can of course have them encounter more if you want. If the characters give in to temptation, use the information in appendix E to keep track of their corruption level.

\begin{DndSidebar}[float=!b]{Guiding the Characters}
Remember, the characters need to use all their various spells and means of obtaining information to track down the items and phylacteries they require on their quest. Especially if \textbf{Tiax} is no longer with them and can't provide them with clues. You can also help the players by disturbing their characters' rest–perhaps resurfacing some of the nightmares described in the introduction–as they approach layers of the Nine Hells where objects they need to obtain are to be found. An interrupted rest is a small price to pay to find what they're after.



\end{DndSidebar}

\section{Cania Overview}
The transition between Maladomini and Cania on the Styx is confusing and magical. One moment the river is sloping downwards through a cave system, and the next it emerges onto an icy glacier. While it cuts through the glacier towards its last stop, even the Styx can't avoid the temperatures of the layer. Icebergs and ice floes drift within the river, and its current slows significantly. Along its banks lie the remains of less-prepared boats, frozen over and forever abandoned. A sheer drop at the edge of Cania marks the end of the Styx, where it turns into frozen falls covering an ancient being on its way to Nessus.

In addition to the ice creating the layer, sheer mountains pierce the skies of Cania. Their faces are jagged and unforgiving, and the cascade of avalanches is almost constant. Partnering the mountains are the mobile glaciers and icebergs that carve the landscape. Seemingly with a life of their own, these massive ice structures meander through the realm, occasionally crashing into each other or the mountains in a spectacular explosion of ice and snow. They leave fissures and cracks miles deep in their wake, which are quickly covered by snowbanks, awaiting unknowing victims who then fall to their deaths.

Because of the freezing temperatures and hostile environment, almost no creature calls Cania its home. \textbf{Mephistopheles} himself, ruler of the realm, spends most of his time in a heated fortress, protected from the elements. His steward, Adonides, is the only archdevil native to Cania, and he regards it fondly despite the elements. Aside from Adonides, only ice devils (gelugons) live in the icy wasteland. \textbf{Mephistopheles} cultivates the devils, employing them as his standing army, spy hunters, and defense throughout the realm.

\section{Leaving Cania}
If the rumors are to be believed, a heavily guarded portal beneath the city of Mephistar is one of only two methods of entering Nessus. The other entrance lies with the Styx and is equally as dangerous. Only by navigating the Styx from its first entrance into Cania can a captain find the exit—otherwise, some form of magic keeps the edge of Cania hidden and the river goes on indefinitely. Where the Styx reaches the edge, it pours over into a massive frozen waterfall, which solidifies around the body of a gargantuan ancient being. Reaching Nessus requires climbing down the frozen creature, a perilous path plagued with deathly falls, dangerous creatures, and twisting paths. Should one survive the trek, they arrive at the entrance to Nessus.

The end of this climb is described in the "Climbing down into Nessus" section in chapter 12.

\subsection{Features}
Colder even than Stygia, all creatures on Cania face a constant battle with hypothermia. Each time a mortal finishes a long rest, they must succeed on a DC 17 Constitution saving throw or gain a level of exhaustion.

\textbf{Mephistopheles} keeps a tight watch on any visitors to the realm to ensure no information is leaked. As part of his influence, any Divination spell cast while within Cania has a 50 percent chance to fail. When a spell fails, its materials and spell slots are consumed, but the spell has no effect. Creatures designated by \textbf{Mephistopheles} are immune to this effect.

\section{Key Locations in Cania}
Some of the major locations found in Cania are described below.

\subsection{Kintyre}
Hidden inside some glaciers and bergs are strange beings and massive constructions, alien to the Nine Hells. Some of the more powerful archdevils have learned to harness the glaciers and control their movements, constructing cities atop the most stable. Others have carved caverns into their faces, where some claim \textbf{Mephistopheles} hides laboratories and libraries of ancient knowledge. Largest among these structures is the ancient city of Kintyre, which taunts explorers with its hidden wealth.

\begin{DndSidebar}[float=!b]{Temptation of Deceit}
\textbf{\textit{"As Good an Answer"}}



If the characters make their way to the city of Kintyre, they find a community of prospecting devils. They're obsessed with finding \textbf{Mephistopheles}' libraries and want to gain the knowledge found within. Camps have sprung up in and around the city, from where these devils venture out on their expeditions. Enterprising devils have set up shop here, selling maps they claim mark locations of possible troves of knowledge, amulets that protect against the harsh climate, and other wares that would attract this particular clientele.



The characters are approached by a devil, who is willing to fabricate false evidence within Kintyre to answer a question of great importance. Run the appropriate temptation event from chapter 2.



\end{DndSidebar}

\subsubsection{Mephistar}
Largest of the glaciers is Nargus, with the equally massive city of Mephistar constructed atop. \textbf{Mephistopheles} rules the realm from his heated citadel and from here he or his steward, Adonides, control Nargus. \textbf{Mephistopheles} has a legion of gelugons guarding the Pit, a shaft rumored to hold a portal to Nessus. It is roughly 500 feet across and several miles deep. A frozen lake lies at the bottom, 1,001 fathoms beneath which lies the portal.

\begin{DndSidebar}[float=!b]{Temptation of Murder}
\textbf{\textit{"To Cheat Death"}}



If the characters explore Mephistar they risk drawing attention from Adonides or \textbf{Mephistopheles}. But if they enter the city nonetheless, a \textbf{succubus} calls to them from a shop inset into the wall of the glacier-city. She offers to inscribe a magic tattoo and suggests it protects against death. She charges only 250 gp and has ink for only one such tattoo. Run the appropriate temptation event from chapter 2 if one of the characters gets tattooed by her.



\end{DndSidebar}

\subsection{The Lake of Wael}
Another icy lake resides near the Falls of the Frozen Titan. The characters are directed there to discover items to help protect themselves when they travel to Nessus. Yet something terrible lingers under its surface and the characters must confront it to claim their prize.

\subsection{Falls of the Frozen Titan}
Once the characters are ready to leave Cania, the most obvious option is climbing down the Falls of the Frozen Titan. If the characters are contemplating using The Pit in Mephistar, \textbf{Koh Tam} urges them against it. He offers instead to steer the barge along the River Styx towards the edge of Cania where the Styx plummets over the edge. There they can climb down the frozen river into Nessus. \textbf{Koh Tam} docks the barge at the banks of the river, a little way from where the characters must start their climb. Proceed to the "\textbf{Tiax}'s Betrayal" section.

\subsection{The Sorrow Mine}
When the characters first enter Cania, \textbf{Tiax} gives them the following information. If \textbf{Tiax} has been disposed of, adjust the text as needed when \textbf{Koh Tam} delivers the information instead:

\begin{DndReadAloud}{}
\textbf{Tiax} braces himself against the skin shredding cold and proclaims, "You ask mighty \textbf{Tiax} about Cania? Well of course, you do, for \textbf{Tiax} knows all! And he tells you this: not even devils are fools enough to linger long in this wind-blasted and land-locked plane. That is, save the gelugons and their liege, \textbf{Mephistopheles}. But they know a thing, they do, the same thing that \textbf{Tiax} knows, but you fools don't know the thing that \textbf{Tiax} knows, do you?"

He smiles smugly as if he has provided some great insight. Then he wipes away the ice accumulating upon his formidable brow and continues, "There's a place. Full of big, sad, moody giants trapped in glacial ice. Dead they are, which is good for the devils, not so much for the giants. So, listen to \textbf{Tiax} when he tells you this: it's the only place that matters in all of Cania. Devils might call it the Heights of Anakim, but \textbf{Tiax} named it the Sorrow Mine. And what you seek? \textbf{Tiax} shall find it for you there."

\end{DndReadAloud}

For his own furtive reasons, \textbf{Tiax} encourages the characters to travel to the Sorrow Mine.

\section{Adventure: The Sorrow Mine}
When visitors are finally able to take it in and understand what they're looking at, the Heights of Anakim are one of the most awe-inspiring and appalling of all the sights in the Nine Hells.

There was, at some point in the shrouded past, a bloodline of true giants, the sort that build worlds for gods. Beings of the earliest creation, when the many planes of the universe were still coming together, possessed of a might and craft no longer seen. As the greatest of giants often will, they fell into a dispute with the elder gods of the time. A war followed, fit to tear the firmament down and sunder the planes. And, in the end, they lost and were hunted across creation by the vengeful gods. In the height of their desperation, the giants made a foolish deal with the lords of the lower planes. Perhaps they were offered eternal sanctuary and preservation against the annihilation the gods were bringing. Whatever the details of the bargain, they were tricked. And yet they were preserved, for in Cania they still remain.

\subsection{Advice from Koh Tam}
\textbf{Koh Tam}'s (or \textbf{Tiax}'s) advice to the characters is as follows:


\begin{itemize}
\item He suggests the characters make their way to the Sorrow Mines, as he believes it is their best bet to find the last items they seek.
\item He warns them to make haste and limit the time they spend here as Cania's cold is dangerous even for the devils of the Nine Hells.
\end{itemize}

\subsection{Arriving at the Cliffs}
\begin{DndReadAloud}{}
You approach a vast glacier face, reaching three miles from the frozen plain, its upper reaches whipped by streamers of cloud constantly pushed past by Cania's eternal wind. The ice is shadowed and discolored. You seem to see shapes in it, perhaps a chance natural pattern of light. Except this effect is no trick of the light. Here are the giants of Anakim, set in the ice. They're the cliffs of the glacier. Packed shoulder to shoulder, toe to heel, wearing rags, their colossal frames emaciated from their long flight from divine retribution. Clutching tattered bundles that might be belongings or might be children. Hundreds of gigantic frozen figures, receding back into the ice until they're just shadows.

\end{DndReadAloud}

Those who get to the upper reaches of the glacier see the faces of the giants. The expressions of semi-divine entities who understood, at the very last, that they had been deceived—and that they were doomed to this frozen, eternal end.

\subsection{Sorrow Wine}
The mere fact and sight of the Heights of Anakim is enough to bring mortals to their knees, but it gets worse. Whilst devils aren't averse to the idea of just leaving things to be tormented forever, they found a use for the giants. So it is that the sharp-eyed visitor scanning the upper reaches of that appalling cliff finds scaffolding and gantries and excavation. Tunnels and quarries that carve into the literal faces of the giants to get at the treasure within.

Long after the powers of the Nine Hells had tricked the Anakim, \textbf{Mephistopheles} discovered the trapped giants contained a unique resource valuable to the rest of the Nine Hells and beyond. A combination of their semidivine nature, the ancient age they were preserved from and their own vast horror at being betrayed means the minds of the giants are a source of the most concentrated misery to be found anywhere in the planes.

The labyrinthine brains of the Anakim contain blue-silver seams of the stuff, that can be quarried out, rendered down, and condensed into vials and barrels of pure corrosive suffering that the devils jovially refer to as 'sorrow wine'. To mortals, the substance is pure vitriol. Concentrated, it causes wounds that strong magic can't heal, and that burn with fresh agony forever. Diluted into a potion it imposes a burden of existential unhappiness that drives the drinkers to empty hedonism and sin as their lives become hollow and meaningless. To devils, the sorrow wine is a potent beverage, a unique vintage of misery fit for the tables of \textbf{Mephistopheles} and \textbf{Asmodeus} themselves. In the Giant Mine, Cania's ruler has a commodity not to be found anywhere else that he uses for bribery, reward, and leverage across the lower planes. His only challenge is reliably extracting it.

\begin{DndSidebar}[float=!b]{The Sorrow of Sorrow Wine}
If sorrow wine is consumed by a creature that is not a Fiend, that creature takes 18 (4d8) necrotic damage. This damage can be healed only with a short or long rest.



\end{DndSidebar}

\subsection{The Sorrow Mine Locations}
Locations relevant to this part of the adventure are described below.

\subsubsection{T1: The Mining Town}
\begin{DndReadAloud}{}
Ahead of you, at the base of those giant sorrowful figures, there are the tell-tale signs of settlement: a glimmer of lantern light and the shadow of smoke blown thin across the gray sky by the icy wind. As you approach, you make out the shape of buildings. A sad collection of crude huts carved from ice and arranged around the only structure that looks like it was built with any care. This building is vast and tall, and even from a distance, you hear the sounds of machinery and harried workers within.

An odd assortment of creatures use the paths leading from town to mine: carrying empty baskets on their way up and heavier loads on their way down.

\end{DndReadAloud}

At the foot (literally) of the Heights of Anakim there is a town of sorts. A huddle of buildings, mostly carved from ice, built around one vast hall raised from iron interlayered with the hides of long-extinct monsters, where the raw giant-stuff is rendered down into sorrow wine. Beyond, the cliff face itself is cut and carved with multiple zig-zagging paths up to the heights, along with winches and lifts to take miners up, and bring the yield down. In the town dwells a variety of devils seldom seen elsewhere in Cania who deal with the administration of the mining efforts, the constant infernal accountancy of supply and demand and a modicum of trading with visitors.

\begin{DndSidebar}[float=!b]{Tiax Takes a Moment}
The first time the characters enter the town with \textbf{Tiax} accompanying them, he hurries to the back of one of the buildings. "\textbf{Tiax} has a mighty burden to discard! I'll just be a moment." If the characters follow him, they catch a glimpse of \textbf{Tiax} handing something to an imp that vanishes promptly. If confronted, \textbf{Tiax} lies and says he was just repaying an old debt. In fact, he is preparing an ambush, to occur when the characters depart Cania. See "\textbf{Tiax}'s Betrayal" at the end of this chapter for more details.



\end{DndSidebar}

\subsubsection{T2: Baron Klendisk's Fort}
\begin{DndReadAloud}{}
While most of the buildings in this town have been carved from the ice itself and squat low to the ground, as if huddling for warmth from the merciless cold, this stone and metal manor has a tall tower rising skyward. This is certainly the home of someone important.

\end{DndReadAloud}

The town itself is overseen by the gelugons—who shun such menial work—and commanded by one Baron Klendisk (\textbf{ice devil} with 250 hit points), who reports directly to \textbf{Mephistopheles}. The baron is an ice devil of twice the usual size, emaciated and hunched almost double, usually to be seen squatting on a spider-legged throne formed from an enslaved ice Elemental. Klendisk's public face is that of a monstrous tyrant, cruel in his discipline towards other gelugons and doubly so to any other devil breeds. Behind closed doors, he is a devil with a problem, because mining the giants and producing the sorrow wine is difficult and his own master brooks no excuses. Hence, while a terror to other devils, Klendisk is known to at least give a hearing to outsiders with a solution.

The difficulty is in the very nature of the stuff they mine. The raw seams of misery in the brains of the giants are powerfully destructive. Whilst Klendisk has no compunction about destroying souls or mortal prisoners, the cold of Cania and the corrosive nature of the misery destroys such tools before they can be of much use. Infernal labor has other problems. Contact with the giant-stuff makes them drunk: boisterous and leery with the heady misery of ages. A fiendish workforce would fall to squabbling, brawling and unacceptable indiscipline, even the icy gelugons themselves. It is a pitiful devil that needs to go beg outside help but that is the situation that Klendisk often finds himself in.

Though these mines would seem to be a good place to avoid, mortal travelers do visit. It helps that Cania has few other locales to recommend it, and the town at the foot of the Heights does at least have shelter and goods for trade. Sometimes mortals are sent there by other Fiends, seeking to abstract some of the processed sorrow wine for their own delectation or trade. Klendisk's ice devils keep a compound eye out for strangers, though, and anyone who comes to their attention without a good excuse or permits signed in triplicate find themselves introduced to the mines or the vats in swift order, or else claimed by Klendisk for a little informal torment. Assignment to the Heights is short on entertainment.

Usually, visitors simply stay long enough to turn a profit and then leave. Others hang on too long or try to claw too much from Klendisk. When his temper breaks, he loses all perspective and has them destroyed or cast out into Cania's killing cold.

\subparagraph{Contractors}
Currently there are two major contractors working at the face, and one more running the rendering shed. Both mining contractors, Kargan and Cornelius, have a constant need of specialist materials—necromantic or machine components—and anyone bringing in fresh stocks will find favor with one (and attract the ire of the other). Alternatively, particularly wicked mortals with a solution to Klendisk's labor problems are always assured a hearing. The mines can always use a new contractor. For the price of 300 gp per day, Klendisk is willing to grant a dispensation to the characters so that they may attempt to work the mines.

\subparagraph{The Prison}
In the Sorrow Mines most of the workforce is either Undead or made of metal. However, there are still plenty of devils roaming about, many of them are under the influence of the sorrow wine. Imbibing too much doesn't always end well for the devil in question. The prison at the edge of Klendisk's Fort houses those unfortunate devils whose liquid bravado provoked Klendisk's wrath. It is run by chain devils with several prisoners of little importance currently housed in the cells.

\subparagraph{Phylactery of the Merciless}
The prison is where the soul of the merciless can be found. Upon entering the village at the foot of the mines, the character searching for this soul, feels a strange pull towards this location. When the characters enter, they must evade a patrol of three chain devils. If they're discovered, they're attacked, and at initiative count 20 of round 5, and every 5 rounds thereafter while combat continues, another chain devil arrives to help. While being attacked in this way, the characters can't easily investigate the cells—they need to eliminate the devils and go back to hiding first.

Finding their way is easy, as the character in question recognizes the voices coming from one of the cells. Inside two figures are torturing a flesh golem that resembles a haggard and emaciated version of the character. The torturers look and sound exactly like the character's true love but are actually chain devils. If the character slays the devils without exposing their true form first, another chain devil appears in the doorway. They immediately transform into the form of the character's true love before entering the cell and attacking. During these battles the devils constantly taunt and mock the character in the voice of their loved one.

A successful DC 14 Wisdom (Insight) check can help the characters understand that they need to remove the devils' disguises before slaying them. They can use \textit{Dispel Magic} or any other similar means to help them break the magical disguise. Once the devils are wearing their own faces, and then slain, the flesh golem vanishes, leaving the phylactery in its place.

As the characters retreat, they must evade another patrol on their way out, as described earlier.

\subparagraph{T3: Kargan Skul's Home}
\begin{DndReadAloud}{}
A large mass of blue-green ice, comprised primarily of sharp angles, appears to be someone's home. All the windows are frozen over and so thick as to obscure any spying from the outside. Only the massive double doors at the front appear free of such frozen constraints.

\end{DndReadAloud}

This is the home of the most long-standing mine baron, Kargan Skul, and even she, a frost giant, bundles herself in fur, against the uncanny cold of the plane. She spends most of her time at the mines (area T6), tending to the Undead workforce she maintains. If the characters enter her home without permission, they're attacked by Kargan's butler (archmage but their creature type is Undead), and the household staff which consists of ten ogre zombies. An \textit{Alarm} spell alerts the frost giant to the intrusion and if Kargan is encountered later in the mines, she attacks the characters without hesitation.

\subparagraph{Treasure}
There's a \textit{Soul Coin}, 10 diamonds (each worth 1,000 gp) and 4,500 gp that can be recovered by pilfering Kargan Skul's home.

\subsubsection{T4: Cornelius Brassgrave's House}
\begin{DndReadAloud}{}
You realize that this is a home-within-a-home when you step through the ice-lodge's door and encounter wood paneling, fine (if scuffed) tile flooring, and a blazing fireplace on the far wall. Near the fire, a gnome sits in an over-sized chair that appears to have been assembled from various cogs, levers, and bits of leather. "It's all insulated," he says, his voice deeper than you expected. "Keeps the outer walls from melting. It is a simple principle, you see, one that involves thermo... well... wait... Ah. Now I see. The giant, she's finally sent you to kill me, has she?"

\end{DndReadAloud}

Cornelius Brassgrave (neutral evil \textbf{veteran}) is a gnome artificer driven from a dozen separate haunts in the mortal planes, a creature of pure reason and no morality whose industrious experiments have poisoned seas and turned the living earth into a metal-tasting wasteland. He has made a number of large claims to Klendisk about how his machinery will revolutionize the mining at the Heights and is now finding that even golems and automata can be ground down by the weight of misery. Cornelius is less suited to Cania than most of the inhabitants of this mining town, so he tends to spend most of his time here, in his very warm home. His Constructs keep him informed and alert him to any problems in production.

Another reason that he's sticking to this haven is that Kargan Skul has been in a foul mood, believing him to be responsible for disrupting her supply lines. He isn't responsible, however, and just as confused about it as she is. If he is attacked, pieces of furniture animate (two stone golems with immunity to cold and fire damage) and come to his defense. If the characters slay Kargan Skul, thus raising Cornelius' value to Baron Klendisk, or they deal with the rogue modron in the mines, Cornelius rewards each character with brass armbands, engraved with the gnome's name and family crest (each is worth 600 gp but have no special properties).

\subsubsection{T5: The Rendering Sheds}
\begin{DndReadAloud}{}
The air inside the rendering sheds is thick with the cloying scent of fermenting drink. Beneath this ugly, sour-sweet smell there hangs a heavy sense of misery, pushing down like an overbearing hand. But, somehow, the modrons that work the sheds have not collapsed beneath that weight. The short, geometrical creatures rush around, checking on machinery and inspecting vast vats of dark liquid.

\end{DndReadAloud}

The rendering sheds themselves are run under a quite different arrangement, one that has persisted for as long as any denizens of the Heights can recall, and which Klendisk himself is powerfully uncertain about. The \textbf{Regular Orthoclath} (lawful neutral \textbf{iron golem} with Charisma, Intelligence, and Wisdom scores of 16) is a thing of facets, eyes and geometrical shapes from the lawful planes of Nirvana, here with an apparently endless workforce of modrons. They work the vats and stills tirelessly, producing sorrow wine so long as the mines yield raw materials, and never complain or ask for payment. They usually work in teams of a dozen \textbf{tridrone} modrons overseen by a \textbf{pentadrone} modron. \textbf{Quadrone} modrons are used by the \textbf{Regular Orthoclath} to send messages to teams.

The modron's involvement predates Klendisk's appointment and the devil exhausts himself trying to discover the actual terms of their employment. In the absence of certainty, infernal rumor is that either these are, in some way, bad modrons who are being punished for crimes against lawfulness meaningless to others—or else the whole business is some grand experiment into the enduring nature of law. Certainly, individual workers eventually become corrupted by the work, sad little pieces of geometry infected with emotion, grinding to a halt as they're consumed by misery. At which point the Orthoclath casts them out, and the devils are free to devour or toy with them, not that the creatures give any real sport. There are always, apparently, more modrons ready to troop down from the lawful planes and give themselves to the work.

\subsubsection{T6: The Mines}
Once the characters reach the mine entrance at the top of the glacier, read the following:

\begin{DndReadAloud}{}
You notice the eyes first. Each of them is enormous, in proportion to the colossal faces they're set in. Most are blown wide in horror, in pure, undisguised terror. These eyes rest above faces full of pain, with mouths drawn tight in grief or frozen open, midscream. But more haunting than the eyes of fear and rage are the eyes that contain nothing at all. No emotion. Just dull acceptance. There are shapes just above the eyes of some of these giants. As you draw closer, you think they might be caves and then you see the workers that scramble in and out of them and you realize that they're tunnels. Mines. Dug straight through the frozen heads of the dead.

\end{DndReadAloud}

There are more treasures than mere misery frozen in the glacier. The giants were artisans to the gods before they rebelled. They clutch in their colossal frozen arms treasures from elder ages, secrets on great withered scrolls, the raw stuff of creation stolen from divine workshops. An intrepid team of mortals might steal up the cliffs and into the mines where untold treasures wait to be unearthed deep in the heart of the ice. Even the spoil-heaps of the mines themselves glitter with discarded wonders, if only one can get to them. For while the devils care only for the misery and the wine they make from it, punishing errant thieves remains one of their few pleasures and they always look forward to the next visitors—mortals or devils—to try it.

The items of value these giants crafted aren't easily transported. If an entire day is spent excavating, the characters recover a 100-pound statue made of gold and gems (and worth d12 x 500 gp). However, this excavation draws the attention–and wrath–of the devils unless the baron has granted the characters a dispensation. A similar statue can be recovered with each day's effort, until 3d4 have been unearthed.

\begin{DndSidebar}[float=!b]{The Sorrow Mine}
While in the Sorrow Mine, any creature that isn't a Fiend that fails a saving throw, now has disadvantage on all saving throws, including death saving throws, until they leave the Sorrow Mine and finish a long rest.



\end{DndSidebar}

\subparagraph{The Miners}
Most of the workers in the mines are either automatons, under the control of Cornelius Brassgrave, or the Undead created by Kargan Skul (lawful evil, Huge-sized \textbf{lich} but she isn't Undead, her creature type is Giant). Cornelius is at his home in the town below, but Kargan is almost always in the mine. She is a necromancer of considerable power and her solution to efficiently mining has been to use mindless Undead to do the brute work. She is constantly importing bones and cadavers from elsewhere in the planes, preferring to create ogre zombies and minotaur skeletons for the mining operation. The misery erodes even these, but so long as she can get the raw material, she can keep up productivity.

\subparagraph{Conclave (Sorcerer/Warlock/Wizard)–Staff of the Magi}
Kargan Skul wields the \textit{Staff of the Magi} (40 charges), carries it with her everywhere, and is quick to use it if threatened. She will generally not relinquish the staff without a fight. However, if the characters deal with the rogue modron sabotaging her Undead workers, she is willing to sell the staff to them, for the sum of 10,000 gp.

\subparagraph{Trouble in the Mines}
There is a steady stream of Constructs and Undead, harvesting the raw material of sorrow wine, and carrying it to the mining town below. The two factions seem to ignore—or not even notice—one another. However, Kargan's workforce has deteriorated faster than usual, with fewer of her zombie and skeletons making the return trek to town. She has attracted Klendisk's ire by falling behind, and believes that her chief competitor Cornelius is sabotaging her supply lines.

Only if the characters are investigating Kargan Skul's claims of sabotage, do they discover the truth. A lowly modron (chaotic neutral \textbf{quadrone} with 50 hit points) has been corrupted by the overwhelming grief permeating the mines. It has gone rogue and skulks the shadows of the mines, pushing miners off scaffolding, or narrow rock bridges when it can. Unfortunately, the Undead tend to be easier to shove than the golems, so many more of Kargan's minions have been lost to the mines, than Cornelius' golems.

After an hour of investigation, the characters come upon the following:

\begin{DndReadAloud}{}
There's a loud swoosh and you see a strange, blockish figure with wings fly out from the shadows and slam into an iron golem carrying a large satchel of ore. The golem doesn't even register the attack, and continues marching forward, while the winged creature shakes its head a moment before flapping its wings and flying back into the shadows.

\end{DndReadAloud}

If the characters pursue the modron, it tries to stay in the air 60-80 feet away from them and fires arrows. It can't be reasoned with. If it is badly damaged, it charges the weakest looking character and tries to grapple them. If successful, it uses its movement to jump into the abyss of the mine with its victim (maximum fall damage).

The characters can use this information as they see fit—if either Kargan or Cornelius learn of the rogue modron, hostility between their two factions decreases a bit. Klendisk, if told, won't reveal the truth to Kargan. He, of course, is more than happy to have Brassgrave and Kargan constantly sniping at each other. He has no interest in his underlings uniting and ending up in a position where they could make demands of him.

If the modron's body is returned to the \textbf{Regular Orthoclath} (area T5), it is taken away to be repurposed.

\subparagraph{Phylactery of the Student}
A character looking for the phylactery of the student, notices a gnome that is dashing between the feet of the zombie workers. The gnome is holding a bucket above her head. In truth the gnome is a shell that is used to torture the student. The soul is linked to the body of a long-deceased female gnome and has been imprisoned here in the mines, enduring the agonizing pain and despair that comes from contact with the sorrow. The student is driven to work harder and harder by the brutish overseers.

\begin{DndReadAloud}{}
A gnome sprints past you carrying, over her head, a bucket filled to the edges, the sorrow sloshing out of it as she runs. As it hits her face and hands, she lets out an agonizing wail. Fresh wounds appear where the oozing liquid touches her skin. Her knees buckle underneath her, and she falls to the floor sobbing. An overseer stomps over and starts chastising the gnome, forcing her to crawl back up. With a despondent sigh, she picks up the now empty bucket and turns back to fill the bucket anew.

\end{DndReadAloud}

If the characters follow the gnome into the workshop, they are attacked by four ogre zombies. Once three of the brutes have died, Kargan joins the fight to defend her workforce. After Kargan is defeated the body of the imp collapses lifeless on the floor, dissolving and leaving the phylactery in its place. Once the character takes possession of the phylactery of the student, they gain the Phylactery Benefit associated with it.

\section{A Thankful Patron}
When the characters are ready to meet with their group patron after completing all of the required objectives, they may ask \textbf{Koh Tam} to sail them to the agreed upon location in Cania. This location should probably be close to the Lake of Wael to encourage the characters to travel there after this meeting.

\begin{DndReadAloud}{}
Sailing through the icy Styx is a ponderous affair, as the vessel glides past erratic ice floes and evades larger icebergs that seem to rise from the waters with disheartening frequency. Eventually the barge slides down yet another ice-crusted tributary and crests a bend in the waters. A pavilion-style tent waits along the shore, several feet back of the waters. A banner rises above the tent, apparently frozen in mid-motion despite these gusting winds of Cania.

\end{DndReadAloud}

The banner is a deep red with electrum-colored striping if the group patron is the Conclave of Halruaa. If the characters allied instead with the Deathstalkers, the banner is a velvety black, speckled with white frost. The Hellrider's banner is a bright silver, on it their coat of arms is depicted in radiant gold.

\subsection{The Hellriders}
It is assumed that the characters have used the Hellrider's Salvation ability to safely return the Celestials they've rescued to the Material Plane. If they're still with the group for some reason, adjust the following:

\begin{DndReadAloud}{}
\textbf{Ramius} turns from a table cluttered with maps and papers, to look at you. His smile is reserved, but it is a smile, nonetheless. "I have received word. You rescued the three from their dire predicaments and fulfilled your end of our arrangement. You have the thanks—and gratitude—of the Hellriders. I know your next steps shall take you into that final layer, where not even we dare to tread. Be yourself prepared before that time—without magical concealment \textbf{Asmodeus} shall surely catch you in his trap before you have even drawn your first breath of the foul air of Nessus."

\end{DndReadAloud}

\subsection{The Conclave}
It is up to the characters whether they recover all the class-based powerful magic items the Conclave helped them locate. Their group patron objective is considered complete once they've dealt with the Unmaker.

\begin{DndReadAloud}{}
The tent is empty when you enter, and in fact snow has blown in from the open flaps, accumulating into a low ridge near the entrance. But as you step foot inside, a semi-transparent image of \textbf{Zythan} appears before you. He extends his arm, as if intending to grasp yours. He smiles. "My duties call me elsewhere, but I wanted to offer you our thanks. Our alliance has been fruitful and though that business with the Unmaker was unpleasant... you dealt with it like the professionals you are. I—we—wish you much luck in your final endeavors. I believe \textbf{Koh Tam} has advice to offer you, regarding a way to hide yourself from \textbf{Asmodeus}. I strongly suggest you heed his words. The Lord of the Nine isn't to be trifled with. I shall linger here a while; in case you need to talk more."

\end{DndReadAloud}

\textbf{Zythan} bows his head a brief moment and then moves about the tent, though it is clear he is actually elsewhere. If the characters have not destroyed the Unmaker's soul, they're asked to place the coin in which they've captured the soul inside the Halruaan Ethereal Vessel. The Halruaans dispose of it now.

\subsection{The Deathstalkers}
The Deathstalker quest is considered completed if the characters have recovered at least 2 of the items—enough to weaken \textbf{Asmodeus}'s allies and make the archdevil think twice about swindling the Deathstalkers.

\begin{DndReadAloud}{}
The rattle of armor alerts you to a warrior's presence beyond the tent, but before you're able to enter, \textbf{Sarevok} himself strides from the interior to confront you on the snowy shoreline of Cania. Behind him several warriors in dark armor stand at attention, but the tent flaps fall back, concealing the interior.

"It is done then." \textbf{Sarevok} says. "Your intent shall take you into Nessus, next, if I am not mistaken. And when you set foot there, \textbf{Asmodeus}, he shall see you. This, I know. You lack the power needed to hide yourselves from him and I can offer you no help with that. But I do demand one of the artifacts I helped you find. At least, \textbf{Asmodeus}, shall not be able to recover that from your smoldering corpses, and I will maintain the leverage I require on him. He shall honor his debts without none of his devils' trickery."

\end{DndReadAloud}

Once the characters select one of the artifacts, \textbf{Sarevok} grasps it, nods his head curtly, and enters the tent.

\subsection{Parting Ways}
The characters, unless they decide otherwise, are still considered as working for their chosen patron, but they no longer have an active mission with them. The characters have fulfilled their obligations, and the patron, by helping them throughout the adventure, has fulfilled theirs. The characters are encouraged to make final purchases from their contacts. The patron remains at this location until the characters return to the barge and sail away.

If spoken to again, the patron encourages the characters to climb down into Nessus, after acquiring appropriate magical protection from \textbf{Asmodeus}'s watchful gaze. They're also cautioned to stay focused once they're in Nessus—they should obtain what they need from that place and return swiftly.

\begin{DndSidebar}[float=!b]{Final Gifts}
This is an opportunity for the group patron to assist the characters. If they require a couple Soul Coin or a small measure of gold, the patron might be convinced to give this to them at this juncture.



\end{DndSidebar}

\section{Before Leaving Cania}
When the characters appear to consider travel into Nessus, \textbf{Koh Tam} becomes concerned. Before entering Nessus, the characters must find the means of avoiding the watchful eye of \textbf{Asmodeus}. Magic items such as Amulet of Proof Against Detection and Location can do this, but the characters would need one for every member of the group. \textbf{Koh Tam} knows of an ill-fated expedition to Nessus that ended in the Lake of Wael. He knows that each one of them possessed an \textit{Amulet of Proof Against Detection and Location} since he was the one who transported the group.

\textbf{Koh Tam} is insistent that the characters would be apprehended the moment they set foot in Nessus, without protection. He strongly urges them to seek out the failed expedition to recover the amulets and offers to guide them towards "The Lake of Wael." Only after some sort of protection has been obtained, should the characters climb the frozen falls of the Styx into Nessus itself. Once they reach the falls however, "\textbf{Tiax}'s Betrayal" occurs.

\subsection{The Lake of Wael}
\begin{DndReadAloud}{}
Sailing the Styx through Cania requires the fortitude to endure the constant cold, but also \textbf{Koh Tam}'s ability to navigate the jagged chicane of ice that the river's surface is plagued with. Visibility is practically nil, with icebergs looming from the frozen fog like great murderous ghosts. And then, when you least expect it, the banks widen until they're lost in the mist and the waters become a rolling chaos of waves, as though you were on the open sea.

\end{DndReadAloud}

Some interaction of the Styx's unnatural water and Cania's unnatural ice has created a great flooded bowl, a lake whose surface has become a drifting maze of frozen islands. And sometimes there is a voice.

Whilst crossing the icy lake is still a formidable task, the lack of a rushing current allows the sailors something of a respite, as they pole their way through the creaking ice. However, the lake is the domain of something quite other, and has been for centuries. \textbf{Waeloquay} (see accompanying stat block) is a water Elemental of immense size and age. Exactly how it came to the Nine Hells in the first place is lost to time—brought by some infernal collector or slipping through a temporary portal between planes. It found its way to the Styx, losing all recollection of its own past and nature but not simply dissipating. Eventually it reached the bowl of the lake in Cania and settled there, slowly growing as it incorporated the substance of the malign river into its nature until now it is the lake, and the Styx passes through it on its constant course down.

\textbf{Waeloquay} is obsessed with memory. Although it attempts to sequester its own experiences from the Styx, they slowly leach away into oblivion. Whenever new ships arrive to navigate its icy labyrinth, \textbf{Waeloquay} sees them as a vital source of new experience. It doesn't care what experience. It just seeks to fill the constantly leaking void within itself.

\begin{DndReadAloud}{}
Rising out of the water are many figures that, from a distance, look like living beings. As you draw closer you see that they're elaborate sculptures of ice. Fiends, mortals, and unknown creatures are all displayed in the throes of motion. The first few sculptures are impressively detailed, but the next set have vaguer features, smoother faces, and fewer accessories. The more sculptures you inspect the less detailed they become, until all that is left are half-formed mannequins with indistinct limbs and blank, smooth faces.

\end{DndReadAloud}

\textbf{Waeloquay} interacts with travelers by manifesting shapes out of water and ice—the forms of those Fiends and mortals that it has previously taken within itself, that become steadily less and less detailed as it forgets their features, until all it is left with are blank-faced mannequins. Alternatively, if anyone has been foolish enough to look down into the water, \textbf{Waeloquay} captures their reflection and creates an \textbf{icy simulacrum} of them. These creations (doppelgangers with immunity to cold damage) crawl onto the ship and seek to ambush individual sailors and throw them overboard into \textbf{Waeloquay}'s clutches, or else try to wreck the entire ship against the ice. Anyone ending up in the water faces the effects of a Styx controlled by a hungry Elemental entity. Anyone who drowns in the water here becomes part of \textbf{Waeloquay}'s library of memories, at least until the Styx abrades them away.

It is possible to bargain with the Elemental for safe passage and/or for the Amulet of Proof Against Detection and Location that were lost here years ago. Fiendish ships often carry sacrifices to throw overboard to sate \textbf{Waeloquay}'s hunger, but it also accepts items of great sentimental value from which it gnaws some mental sustenance.

If the characters don't bargain, then they must fight against a manifestation of the lake. If they defeat it, then \textbf{Waeloquay} quickly moves them to the end of the lake and throws eight amulets into the barge, along with the skeletons still wearing them.

\section{Tiax's Betrayal}
When the characters journey to the Falls of the Frozen Titan and, just prior to descending into Nessus, this event occurs. \textbf{Tiax} has determined that he must now eliminate \textbf{Koh Tam}. He approaches the characters, alongside the other 2 crew members who have helped \textbf{Tiax} and \textbf{Koh Tam} manage the ship. \textbf{Koh Tam} follows some distance behind. If the characters are primarily evil, read the following (but skip past it, if the characters aren't evil).

\begin{DndReadAloud}{}
\textbf{Tiax} has a strange look on his face, a gleam to his eye as he clasps his hands together and stares out at the icy expanse of Cania. He lets out his breath. "Now is the time! What time, you ask, oh great \textbf{Tiax} the all-knowing? Time for the right and the left hands of the ever-prepared \textbf{Tiax} to strangle the very breath from the lungs of pitiful \textbf{Koh Tam}. Kelemvor shall not win this day, this time. \textbf{Tiax} the mighty shall not allow it! Dying always teaches the best kind of lesson, do you not agree?"

\end{DndReadAloud}

\textbf{Tiax} presents his deal very quickly. If the characters help him slay \textbf{Koh Tam}, he promises to give them the barge, the 20 Soul Coin he carries, and swears to Cyric that he will help them complete their quests in the Nine Hells.

Otherwise, if the characters are good-aligned, or don't immediately agree to join \textbf{Tiax}, he springs his ambush on them and \textbf{Koh Tam} at this moment instead.

\subsection{The Attack Against Koh Tam}
\textbf{Tiax} has previously cast \textit{Mage Armor} on himself. Regardless of the characters' affiliation, he approaches \textbf{Koh Tam} in a friendly manner but casts \textit{Contagion} (using the mindfire effect). The attack roll required for this spell automatically succeeds against the unaware \textbf{Koh Tam} unless the characters are able to warn him. Depending on past character decisions, \textbf{Tiax} may have allies or additional opponents as described below:


\begin{itemize}
\item \textbf{Tiax} has bribed one of the crew members (use the \textbf{veteran} stat block) and they assist \textbf{Tiax} in this battle.
\item Additionally, if Grinken Eyre (neutral evil Medium-sized \textbf{empyrean}) was encountered in Minauros and not slain by the party, he has masqueraded as the other crew member since then and allied with \textbf{Tiax}. He assists \textbf{Tiax} in this fight. Grinken Eyre's presence makes the battle considerably more difficult if he is fighting the characters, so feel free to omit him.
\item If the characters have previously met Sir Calenhad Strongheart in the "A Paladin In Hell" encounter, he may have joined them on the barge. If so, he sides with \textbf{Koh Tam} in this battle and comes to the characters aid if the characters do so as well. Calenhad is a lawful neutral \textbf{death knight} that doesn't have the Marshal Undead trait or the Hellfire Orb action.
\item If Grinken Eyre isn't present, \textbf{Tiax} previously called in a favor owed him by an \textbf{ice devil}, who shows up at initiative count 10 of the second round of combat. The other crew member (also a \textbf{veteran}) fights to defend \textbf{Koh Tam}, in this situation.
\end{itemize}

The characters can choose to side with \textbf{Tiax} or \textbf{Koh Tam}. If they ally with \textbf{Tiax}, it is very likely \textbf{Koh Tam} will perish here and \textbf{Tiax} will serve his role for the remainder of the plot. On the other hand, \textbf{Tiax} is very fond of being alive and if reduced to less than half of his hit points (or otherwise feeling like the battle is going against him) casts \textit{Dimension Door} in an attempt to retreat.

\subparagraph{Treasure}
\textbf{Tiax} carries 10 Soul Coin (he lied earlier about the amount) and so if the characters allied with him he promises to pay them the rest later. If the heroes defend \textbf{Koh Tam} successfully, he rewards them with 1 \textit{Soul Coin} each.

\section{Leaving Cania}
The characters are nearing the end of their ordeal through the Nine Hells. Before entering Nessus however, \textbf{Koh Tam} offers the following advice. If \textbf{Koh Tam} is dead, \textbf{Tiax} offers it instead, though perhaps not as truthfully as \textbf{Koh Tam} would.

\begin{DndReadAloud}{}
"The hard end of the Nine Hells, Nessus; when you've descended as low as you can possibly go. Hearing of it, one might think it isn't actually so bad. Not actively on fire, not frozen, not a deadly morass or a blasted mountainside. Just a regular kind of wasteland, dry and cracked, and finite, past which lies the infinite crimson void."

\textbf{Koh Tam} rubs his hands together, to warm them. "In a sense, Nessus is the Nine Hells condensed. All the devils are there, as the saying goes. \textbf{Asmodeus} does not brook trespass. The towers of the great city of Malsheem rise from the city's chasm to the highest reaches of the plane, and from them \textbf{Asmodeus} sees all. Once he receives the faintest whisper of an outsider loose in his domain, his eye will find them. You cannot hide from \textbf{Asmodeus}, only hope that he has something of greater import to devote his attention to. Mortals venturing here find no respite. Except...

There is a place. Word of it seeps out into the upper layers of the Nine Hells, and the rest of the planes. Even in the domain of \textbf{Asmodeus} himself there is a place of sanctuary—if you can reach it. \textbf{Asmodeus} is an archfiend of intellect. He prides himself on his learning as well as his power. It is an offence to his pride that there are gaps in his knowledge. So it is that there is one place that the wise might come to and be safe from infernal retribution. Some are called, and others follow that chain of rumors. All of these come, eventually, to the Oasis of the Lethe."

\end{DndReadAloud}

Proceed to the "Climbing down into Nessus" section in chapter 12 if the characters have finished their adventures in Cania and are undertaking the climb. The characters are given an infernal map for the Oasis of the Lethe and cautioned to hang onto it until they need it. \textbf{Koh Tam} or \textbf{Tiax} offer the following parting advice:

\begin{DndReadAloud}{}
"Together we have accomplished much, but your travels, and mine, are near an end. I shall wait on the River Styx for you for seven days. I do hope to see you again. If not, I can only pray that your suffering is brief. Fare you well."

\end{DndReadAloud}

\chapter{Hunted by the Dukes}
\DndDropCapLine{T}{}he events in this section may occur at any time during the characters' journey through the Nine Hells.


Not all dangerous devils confine themselves to a specific layer of the Nine Hells. Some roam, either escaping from enemies, or hunting potential victims. Among the greatest of these powerful devils are \textbf{Brother Adramalech} and \textbf{Brother Morax}, who are agents of \textbf{Asmodeus} and command a warship, the likes of which have never been seen outside the Nine Hells. But \textbf{Baalzebul} also has an agent who travels the River Styx—an archdevil known as \textbf{Abigor}, who pilots an infernal submersible.

\section{Running This Chapter}
There are two methods by which the characters can encounter either of these naval opponents. Either the characters are hunted by the warship or the submersible, or the characters hunt one, or both, of these devil-ships.

\subsection{The Hunters}
The archdevils, \textbf{Brother Adramalech} and \textbf{Brother Morax}, have been tasked by \textbf{Asmodeus} to hunt the characters down and take them captive. If they succeed in taking the characters captive, then they take them to Malsheem where they're confronted by \textbf{Asmodeus} (refer to the section on "Malsheem" in chapter 12).

The archdevil \textbf{Abigor} has been tasked by \textbf{Baalzebul} to simply kill the characters and eliminate the threat to him that they present.

\subsection{The Hunt Begins}
These archdevils start searching for the characters once they've completed one of their group patron's objectives. From this point on, every time the characters move between layers (while on the River Styx) there is a 25 percent chance that they're intercepted by one of these foes.

If the characters are intercepted, proceed to the "Ambushed by the Infernal Warship" or the "Ambushed by \textbf{Abigor}" section. You choose which one.

The Characters Are Searching The characters might pursue \textbf{Brother Adramalech} and \textbf{Brother Morax}, in the hopes of rescuing \textbf{Barachiel}, who is a prisoner aboard their infernal warship. Alternatively, they might seek \textbf{Abigor} for one of the following three reasons:


\begin{description}
\item[The Triad.]
The characters may have made a bargain with the devils \textbf{Rimmon}, \textbf{Zagum}, and \textbf{Buer} to capture \textbf{Abigor} (see "Malsheem" in chapter 12).

\item[The Geas.]
They may have allowed \textbf{Anacreda} to place a geas upon them that compels them to slay \textbf{Abigor} (see "Anacreda the Angelmaker" in chapter 8).

\item[The Unmaker.]
They may be seeking the Unmaker if the Conclave is their patron.

\end{description}

If the characters are searching for either vessel, each time they enter a new layer of the Nine Hells via the River Styx, roll a d8. The quarry they seek is on the layer of the Nine Hells that matches the result. Reroll the d8 if the result matches the layer which the characters are just exiting. If the ship they seek is on the same layer as they are, then the characters need to succeed on a DC 20 Wisdom (Survival) check to locate the target vessel or use magic to determine its location. If the characters have been told of the location of \textbf{Abigor}'s submersible, there is no need to make this check—just pick a layer of the Nine Hells and inform the characters that is where \textbf{Abigor} may be found.

If the characters find their desired ship, proceed to either "The Brothers' Infernal Warship" or "\textbf{Baalzebul}'s Infernal Submersible."

\begin{DndSidebar}[float=!b]{Corruption}
These encounters may trigger a premature ending to the adventure. An additional gating mechanism you might consider is to avoid triggering this hunt until the characters have become corrupted. The information in appendix E helps you keep track of their corruption. Once they've grown too corrupt, the characters are infinitely tangled in \textbf{Asmodeus}'s chains and making a deal with the archdevil is likely their best option anyways. But if they still have a chance of redemption, they don't want to end up in Malsheem too soon.



\end{DndSidebar}

\section{Ambushed by the Infernal Warship}
Review the top deck (area N1), for details on the attacking force (the ship is considered to be at battle alert). Once the warship is close enough to the barge, the devils on the top deck, as well as \textbf{Brother Morax}, leap onto \textbf{Koh Tam}'s barge to attack the characters. If the characters fall, \textbf{Koh Tam} and \textbf{Tiax} (if present) are incapacitated during the fight while the characters are captured and brought aboard the warship.

If the characters repel the invaders, or are captured, proceed to "The Brothers' Infernal Warship."

\begin{DndSidebar}[float=!b]{Unfair Fate}
If the characters are captured now, it will be challenging for them to complete any remaining quests. Use your discretion—if you don't think the characters are ready for the end stages of the plot, then don't have the Dukes find the characters just yet.



\end{DndSidebar}

\section{Ambushed by Abigor}
\textbf{Baalzebul}'s submersible is considered at battle alert when it finds \textbf{Koh Tam}'s barge. The submersible surfaces and Seleceus (\textbf{pit fiend}) flies above the barge and begins repeatedly using its \textit{Fireball} spell against it. Their intent is to kill the characters, and swiftly. \textbf{Abigor} himself prepares to defend the submersible if Seleceus is slain. If the characters attack the submersible, \textbf{Abigor} does everything in his power to defeat them.

If the characters are victorious and explore the ship, proceed to "\textbf{Baalzebul}'s Infernal Submersible."

\section{The Brothers' Infernal Warship}
The warship is under the command of \textbf{Brother Adramalech} and \textbf{Brother Morax}.

\begin{DndSidebar}[float=!b]{Warship Weapons}
Each weapon requires an action to fire it and another action to reload it.



\subsubsection{Acidic Sprayer}
Acidic bile sprays from a nozzle in a 30-foot cone. Each creature in the cone must make a DC 12 Dexterity saving throw, taking 40 (9d8) acid damage on a failed save, or half as much damage on a successful one. A creature reduced to 0 hit points by this damage is dissolved, leaving behind any objects it was carrying or wearing.



\subsubsection{Flamethrower}
Fire shoots out of the weapon's nozzle in a 60-foot line that is 5 feet wide. Each creature in the line must make a DC 15 Dexterity saving throw, taking 18 (4d8) fire damage on a failed save, or half as much damage on a successful one. The fire ignites any flammable objects in the area that aren't being carried or worn.



\end{DndSidebar}

\subsection{Encounters}
The exterior of this warship is patrolled by devils, as described in area N1. Additionally other devils occupy the watchtower (area N18). Evading their notice is challenging. If the characters somehow enter the ship without being noticed, the hallways are patrolled by a pack of four displacer fiends (see appendix B). Every minute spent in the hallways there is a 1 in 4 chance of running into them.

\subsection{Locations in The Brothers' Infernal Warship}
The major locations found within the ship are described below.

\subsubsection{N1: Top Deck}
\begin{DndReadAloud}{}
The top deck of the ship is outfitted like a tribute to extremity. Ten flame throwers and ten acid sprayers line the deck in neat, alternating order, bubbling with vile green foam, or spurting out little bursts of fire, and interspersed among them are three enormous ballistae, armed with gleaming metal projectiles. At the rear of the ship, you feel the cutting heat of the furnace that powers the terrible vessel, and you hear the agonized wails of the souls being fed into it below.

\end{DndReadAloud}

The top deck is usually crewed by a half-dozen cambions aided by ten bearded devils. They attack anyone that isn't a crew member. If the ship is at battle alert, then \textbf{Brother Morax} is on the top deck.

Characters might try to operate the shipboard weapons. While the great ballistae are only usable against other warships, the acid sprayers and flamethrowers can be pivoted to attack creatures on the deck (see the sidebar for details).

\subsubsection{N2: Acid Vats}
\begin{DndReadAloud}{}
A swirling morass of sizzling green acid jostles with the motions of the warship, acid occasionally sloshing almost over the lip of the vat containing it.

\end{DndReadAloud}

These chambers are where the acid is stored that is used in the sprayers. There are ten of these acid vats, each one colored green on the accompanying map. Each vat is situated below an acid sprayer on the top deck (area N1).

\subsubsection{N3: Hellfire Fuel}
\begin{DndReadAloud}{}
The Hellfire in the vat before you emanates a tremendous heat even before you ever get close to it. The air above the vat shimmers, distorted by the heat and occasionally sparks erupt, landing with a sizzling hiss upon the metal floor.

\end{DndReadAloud}

These chambers are where the fuel is stored that is used in the flame throwers. There are ten Hellfire vats, each one colored red on the accompanying map, and each situated below a flame thrower on the top deck (area N1).

\subsubsection{N4: Command and Navigation Bridge}
\begin{DndReadAloud}{}
The room is cramped and dark, with most of the space being taken up by a massive chair and an equally oversized wheel. The chair and wheel face the only windows in the room, toward the bow of the ship.

\end{DndReadAloud}

\textbf{Brother Adramalech} can often be found here. He is usually only accompanied by three imps that he uses to send messages to other parts of the ship.

\subparagraph{Treasure}
There are a half dozen steel chests in this room. Each is locked and requires a successful DC 20 Dexterity check with thieves' tools to unlock. Five of the chests contain platinum bars, each worth 1,000 gp. Each chest has 5 bars stacked inside of it.

The last chest contains plans written by \textbf{Asmodeus} himself. The plans detail how he intends to trick \textbf{Baalzebul} into leaving his plane so that he can turn him back into the slug form that \textbf{Baalzebul} so despises. He has ordered the archdevil brothers to find \textbf{Abigor}, \textbf{Baalzebul}'s trusted lieutenant, and take him prisoner to draw \textbf{Baalzebul} out. This is the same plan that the Triad will offer the characters (see chapter 12).

\subsubsection{N5: Brig}
\begin{DndReadAloud}{}
This room, near the front of the lowest level of the warship, is blocked by a solid metal door. From behind the door, you hear muttering and complaining. On a hook beside the door hangs a blackened key.

\end{DndReadAloud}

The rooms are used for disobedient crew members and prisoners of war. Unlocking the door with the key makes a loud screeching noise and by the time the door is fully opened, the prisoners beyond are already standing and looking expectantly at the door. There are currently four prisoners being held in separate cells—a \textbf{cambion} named Shlatchel, a \textbf{bearded devil}, a \textbf{couatl} that has assumed the form of a cambion to avoid conflict with his cell mates (though all other crew members know the truth) and the former Hellrider, \textbf{Barachiel} (see appendix C).

The key that opened the outer door also opens any of the cells. When \textbf{Barachiel} sees the characters, he approaches the bars of his prison.

\begin{DndReadAloud}{}
The aasimar before you glows with what can only be described as a divine radiance. Despite his captivity and the wounds inflicted upon him, his grin is joyous. When he speaks, his voice rings with certain truth. "It was \textbf{Ramius}, my long-ago friend, that has sent you? He and the Hellriders?"

\end{DndReadAloud}

\textbf{Barachiel} insists on freeing the other prisoners (even though he doesn't know that there is a couatl here). He feels all of them have served enough time aboard this infernal vessel already. Once \textbf{Barachiel} is restored, the characters may use the ritual that the Hellriders taught them to free him from the Nine Hells.

\subsubsection{N6: Nest}
This is where the brothers keep their vorvolakas (see appendix B). There are three to be found here unless the ship is at battle alert, whereupon the room is empty.

\subsubsection{N7: Engine Room}
\begin{DndReadAloud}{}
The engine seems almost like a living thing. It purrs, a deep rattling, rumbling sound, and the sharp, heavy heat that comes off it pulses like a giant's breath. Lemures slink about the room, tending to its needs. Among them stalks a hulking brute with tall, curved horns, snarling at the lemures and watching with gleeful satisfaction as the door to the engine chamber is opened and soul coins are thrown into the mouth of the beast. The souls shriek and scream as they're consumed by the red-hot flames within, their torment fuel for the engine's endless hunger.

\end{DndReadAloud}

The \textbf{horned devil} is called Coal Heart (he keeps his real name to himself) and does his job with a great deal of relish. A dozen lemures work for him and sometimes, when he's bored, he'll throw one of them into the furnace. Coal Heart is a favorite of \textbf{Brother Morax} and there is a 1 in 4 chance that the archdevil is making his way to the engine room to spend time with his favorite subordinate.

Anyone entering the furnace (or starting their turn in it) must make a DC 15 Constitution saving throw, taking 17 (5d6) fire damage and 17 (5d6) necrotic damage on a failed save, or half as much damage on a successful one. Both Coal Heart and \textbf{Brother Morax} attempt to grapple an enemy to drag them into the furnace if combat occurs here.

\subsubsection{N8: Kitchens}
\begin{DndReadAloud}{}
A dozen or so tieflings toil over vast pots of lumpy, bubbling goo, their shoulders slumped and eyes half closed in exhaustion. Imps dart around underfoot, taking away bowls full of squirming muck and dashing off with them to the mess hall. The air in the kitchens is stifling hot with a putrid smell that stings the eyes.

\end{DndReadAloud}

The tieflings (commoners) have no loyalty to the Brother archdevils. They're willing to tell intruders everything they know, as long as none of the imps are around to rat them out.

\subsubsection{N9: Mess Hall}
\begin{DndReadAloud}{}
There are several long tables in the mess hall, all of them lined on both sides with an assortment of rickety chairs and sagging benches. Tieflings, cambion devils, and bearded devils can be found sitting, standing, and slouching around the tables, shoveling unidentifiable clumps of food into their mouths with joyless abandon. Imps come in and toss more bowls of vile sustenance down the tables and take away the empty dishes and unused utensils as they leave.

\end{DndReadAloud}

There are usually three or four each of tieflings (bandits), cambions and bearded devils. There are several smaller mess halls along the lower deck, but each only contain one or two occupants at any time.

\subsubsection{N10: Common Quarters}
There are numerous identical crew quarters throughout the lowest level of the warship.

\begin{DndReadAloud}{}
The bunk room is cramped, dark, and smells vaguely like fish and a great deal like rotting flesh. The bunks are packed together so tightly they might as well be one unit. There are no mattresses, just stretched canvas cots with a layer of ragged, suspiciously stained blankets tangled together on top.

\end{DndReadAloud}

\subparagraph{Treasure}
Searching any of these rooms usually nets 1d10 gp and 1d4 pp. There is a 10 percent chance that 1d4 bearded devils return to their quarters during the characters' search. They won't engage but run to area N4 to alert the rest of the ship.

\subsubsection{N11: Officer's Quarters}
Below-decks there are 6 almost identical officer's quarters.

\begin{DndReadAloud}{}
There are two beds in these small quarters. A taste of sulfur lingers in the air.

\end{DndReadAloud}

\subparagraph{Treasure}
Similar to the common quarters, but a search reveals 1d10 pp and 1d4 gems worth 50 gp each.

\subsubsection{N12: Workshop}
\begin{DndReadAloud}{}
Approaching the workshop, you're assaulted by a vicious cacophony of awful noises. Within the large, overheated room you find several devils making repairs on a variety of weaponry and metal machinery. The sound of hammering and metal striking metal reverberates harshly off the steel walls, creating an agonizing and unending echo.

\end{DndReadAloud}

A \textbf{maelephant nomad} (see appendix B) named Zerkosis runs this workshop. It is assisted by four bone devils and four barbed devils.

\subsubsection{N13: Storage}
The various storage rooms are filled with rations for the crew.

\subsubsection{N14: Soul Fire Engine}
This is the engine that powers the ship. Anyone entering the engine area (or starting their turn in it) must make a DC 15 Constitution saving throw, taking 17 (5d6) fire damage and 17 (5d6) necrotic damage on a failed save, or half as much damage on a successful one.

\subsubsection{N15: Fuel Reserve}
This reserve of fuel is used to power the ship if the usual soul fuel is unavailable.

\subsubsection{N16: Armory}
The armory is filled with weapons of almost every kind. In addition, there are two magic weapons—a \textit{Stygian Spear} and a \textit{Demonbone Polearm} (see appendix D).

\subsubsection{N17: Kennels}
\begin{DndReadAloud}{}
This room apparently has no means of illumination. At the far end, in the shadows, rest several bulky, lion-like forms. The air around them seems to waver as if there were a disturbance near the creatures. You find it hard to concentrate on them.

\end{DndReadAloud}

This is where the displacer fiends (see appendix B) are kenneled. There are six of them here resting at any given time.

\subsubsection{N18: Watchtower}
\begin{DndReadAloud}{}
Near the rear, a metal-clad tower rises over the bulk of the warship, red-hued windows offering a full view of the ship's deck as well as the river all around it.

\end{DndReadAloud}

At any given time, four cambions are on watch, observing the cardinal directions for signs of other vessels approaching. If a vessel is spotted, or if intruders enter the watchtower, one of the cambions attempts to flee to warn the Brothers.

\subsubsection{Brothers Adramalech and Morax}
\textbf{Asmodeus} seldom deigns to wield the whip himself when someone rouses his ire. He has lesser devils to do that for him. Though, in this case, 'lesser' can still mean very powerful indeed. Adramalech and Morax have served the greatest of Fiends faithfully since the earliest days of the Nine Hells. They were angels within the host \textbf{Asmodeus} commanded. When he fell, they cast themselves willingly down to be with him. In an internal hierarchy rife with ambition and betrayal, the brothers' lawful nature manifests as absolute loyalty. In return, \textbf{Asmodeus} gives them what they want. They stand outside the intricate hierarchies of the Nine Hells, reporting directly to him and undertaking his commands with their own hands, his personal enforcers. His 'hounds', say other devils with a sneer of contempt, but Adramalech and Morax are as content with their lot as any devil can be.

They manifest most often as chalky white humanoids, Adramalech around eight feet tall and the hulking Morax, twelve. Morax is the more physically potent of the two, his hide spined and barbed, his hands jagged talons. Adramalech does most of the talking, to dreadful effect. As befits \textbf{Asmodeus}'s devils-of-all-work, they're a versatile pair of monsters, more than equipped to punish those who have disappointed their master. And 'punish' is the operative term. \textbf{Asmodeus} dispatches his enemies to them when they've roused his personal ire.

Morax (see the accompanying stat block) has the tremendous strength and speed one would expect of a martial devil. His spines are brittle and hollow, and shards of them lodge in the flesh of those he strikes, or who strike him, where they rapidly begin to grow jagged extrusions of their own like miniature sea urchins. The shards must be cut out swiftly before they carve his victims up from the inside. This is mostly by way of a distraction so that Adramalech (see the accompanying stat block) can play his own games though.

Adramalech retains his angelic voice from before the fall. His golden words inveigle their way into the ears of listeners, confusing, distracting, distorting the way they see the world and causing mind-wrenching dissonance in mortal minds. Worse than this, Adramalech can conjure simulacra of his enemies' loved ones, living or dead. More than mere illusions, these creations have a temporary life and thought drawn from the memories of the brothers' enemies. This psychic link means that, when Adramalech torments them or Morax tears them apart, the pain is felt by those whose minds gave rise to them. In this way, \textbf{Asmodeus}'s decreed punishments combine physical agony and mental trauma in equal measure.

Though concentration on destroying the less imposing Adramalech might seem the logical way to assault the brothers, the pair have a final surprise for their foes. They were as close as two angels could get before the fall and nothing in all the eternal ages of the Nine Hells has ever divided them. If either falls while the other still stands, the remaining Brother's presence sets their fallen kin's flesh bubbling in a frenzy of regeneration. They must be defeated together or not at all.

\section{Baalzebul's Infernal Submersible}
The submersible is commanded by \textbf{Abigor}.

\subsection{Baalzebul's Infernal Submersible Locations}
The major locations found within the submersible are described below.

\subsubsection{B1: Bridge}
\begin{DndReadAloud}{}
The long hallway of the bridge is eerily silent. You hear the steady rush of water outside the submersible as the vessel cuts through the river, but inside it is still and dark. The brightest light comes from the words writ in fire on the walls, that flash briefly and then vanish. A long metal device hangs from the ceiling in the center of the room, with handles on the sides and small windows for the eyes.

\end{DndReadAloud}

The bridge is always crewed by a half-dozen bone devils. There is a 50 percent chance that the commander on deck is \textbf{Abigor}, otherwise it is the first mate, a \textbf{pit fiend} named Seleceus. If the submersible is currently at battle footing, then \textbf{Abigor} is in command. The words on the wall are written in Infernal and are how the aboleth in area B3 communicates with the bridge commander.

\subsubsection{B2: Brig}
\begin{DndReadAloud}{}
This small room appears to usually be used for storage, but manacles have been bolted to the wall and a tiefling sits in a comfy chair beside them. There's a plate of food in his lap and a half full bottle of wine beside the chair. A small satchel sits under the chair. As he hears you enter, he looks up, licks his fingers, and smiles. "More guests of our superb host, \textbf{Abigor}? I'll wager a soul coin, that you are here to rescue me."

\end{DndReadAloud}

The prisoner is named \textbf{Aeshma} (see his entry in appendix C). The moment he walks free of the room with the characters, he holds out his hand for the coin he has won from them (he did win the wager after all). His satchel contains two Soul Coin, 15 gp, and a set of gambler's dice.

\subparagraph{The Conclave as Patron}
If the characters are working for the Conclave, \textbf{Aeshma} is the Unmaker. To fulfill their patron's objective, the characters must capture and destroy \textbf{Aeshma}'s soul. \textbf{Aeshma} knows nothing of this: as far as he is concerned, he's just a fairly lucky (albeit compulsive) gambler making his way through life. Despite his ignorance, \textbf{Aeshma}, and the protection being the Unmaker provides, isn't going to make this grim task easy for the characters. The following describes what the characters might learn from \textbf{Aeshma}, or how he might react depending on how they interact with him:


\begin{description}
\item[The Manacles.]
If asked why he's not chained up, he merely shrugs and says he made a bet with the guards and he won. He'll raise the plate of food and declare, "That's how I got this as well." However, if the characters appear to be dangerous to him, he lies instead and insists that there was another prisoner here, "A really bad fellow; they just dragged him off, that-away." If they leave, he attempts to hide from them in the submersible.

\item[The Unmaker.]
If asked about being the Unmaker, \textbf{Aeshma} (truthfully) has no idea what the characters are talking about. The reality is that \textbf{Asmodeus}—and perhaps other archdevils—would like to use the Unmaker to set the stage for a wager of divine proportions, one that will leave the world permanently changed and possibly irrevocably broken. Though \textbf{Aeshma} himself is innocent in this future wrongdoing, his soul is simply too dangerous to be allowed to exist.

\item[If Attacked.]
\textbf{Aeshma} tries to talk his way out of conflict, and if that fails, he flees. If slain, the characters have 1 minute to use the Conclave Coin to capture \textbf{Aeshma}'s soul, otherwise the Unmaker is reborn in the next generation and the threat returns (the characters fail their patron, who no longer provides additional support, other than \textbf{Koh Tam}'s services).

\end{description}

Ultimately the characters must slay \textbf{Aeshma} and capture his soul with the Conclave Coin. Once accomplished, they can bring the coin with them to their final meeting with their patrons (see "A Thankful Patron" in chapter 10). Clever characters might instead destroy the coin now, ensuring that the soul doesn't fall into the wrong hands. An infernal ship's furnace consumes the Conclave Coin instantly and destroys the coin in the process. The soul trapped in the coin becomes trapped in the furnace instead, powering the infernal ship. Not even divine intervention can restore a soul destroyed in this manner.

\subsubsection{B3: Communication Room}
\begin{DndReadAloud}{}
There is an immense tank filled with water in the center of the room. A grotesque creature with countless tentacles swims in it. The creature uses its many limbs to touch the infernal sigils carved into the wide wall in front of it with careful intention, leaving streaks of slime in its wake.

\end{DndReadAloud}

The tentacled creature is an \textbf{aboleth} that was imprisoned by \textbf{Abigor} centuries ago. It now serves the purpose of keeping \textbf{Abigor} in communication with everyone on the ship and its other servants throughout the Nine Hells. The aboleth uses the sigils to send messages to other rooms in the submersible or receive messages from those locations. Likewise, and not unlike a Sending Stones, some of \textbf{Abigor}'s minions carry talismans that can send messages to this room or receive messages from it, even when they are on a different layer.

The aboleth knows many things and it is willing to share for the right price. A rare or more powerful magic item is what it wants in exchange for knowledge. It knows almost as much about the Nine Hells as \textbf{Asmodeus} himself. The most important thing it knows is the truth of \textbf{Asmodeus}'s origins. You choose which of these origins is the true one:


\begin{itemize}
\item \textbf{Asmodeus} arose from the primordial chaos as the mightiest of the lawful gods, with Jazirian the only one who could rival him. They both took the form of a serpent and they set their minds on bringing law to the chaos. Eventually the two gods fought, and \textbf{Asmodeus} was defeated. His body fell through the multiverse until it crashed into Nessus forming the canyon known as the Serpent's Coil.
\item \textbf{Asmodeus} is a fallen angel. He was originally tasked with keeping the demons of the Abyss in check, but eventually to become better at killing demons he and his angelic followers took on demonic traits. This led to him being banished by the gods. But before he was exiled to the Nine Hells, he was able to fool the gods into signing the Pact Primeval. This is a contract between \textbf{Asmodeus} and the gods of the multiverse that allows \textbf{Asmodeus} and his servants to legally corrupt mortal souls so that they end up in the Nine Hells. These souls can then be used to increase \textbf{Asmodeus}'s power.
\item \textbf{Asmodeus} served the first and most powerful god, He Who Was. \textbf{Asmodeus} chaffed at having a master, so he used an artifact called the Shard of Evil to kill his god. Then he was able to erase the name of the god from existence. The Shard of Evil is the ruby at the tip of the rod that is \textbf{Asmodeus}'s signet. \textbf{Asmodeus} was banished to the Nine Hells by the other gods.
\end{itemize}

\subparagraph{Betraying Baalzebul}
If the characters are working with the Triad to trick \textbf{Baalzebul} into reverting to his slug form again, they can ask the aboleth to summon the archdevil here. The aboleth does this only if \textbf{Abigor} is incapacitated or dead and the characters must agree to cede control of the submersible to the aboleth afterwards. For the ruse to work, \textbf{Baalzebul} must be convinced to leave Maladomini. If the submersible is currently on the River Styx in Maladomini, the aboleth is able to navigate the vessel to the layer above or below this one before contacting \textbf{Baalzebul}.

Once the aboleth summons \textbf{Baalzebul}, proceed to the "Capturing \textbf{Abigor}" section at the end of this chapter.

\subsubsection{B4: Engine Room}
This is a Hellfire furnace. Anyone entering the furnace (or starting their turn in it) must make a DC 15 Constitution saving throw, taking 17 (5d6) fire damage and 17 (5d6) necrotic damage on a failed save, or half as much damage on a successful one.

\subsubsection{B5: Kitchens}
\begin{DndReadAloud}{}
The kitchen staff, a handful of wretched imps, work quickly and with stern silence. They don't speak to each other, nor make any noise at all. There are towels on their cutting boards so that their knives won't make a sound as they slice through hunks of unidentifiable meat-like substance, and they stir tall pots of frothing gruel slowly, taking care not to strike the metal sides with the stirring utensil.

\end{DndReadAloud}

The six imps flee from any intruders, looking to alert the rest of the ship.

\subsubsection{B6: Showers}
There are four rooms dedicated to bathing or showering, along with toilet facilities. These are often empty.

\subsubsection{B7: Common Quarters}
There are two large rooms dedicated to housing and feeding the crew.

\begin{DndReadAloud}{}
There is a motley collection of devils in the mess hall. Great and small figures sit silently side by side, eating spoonfuls of thick, grayish gunk. Only the slightest of slurping, chewing noises can be heard, and beneath these sounds, the distant rush of water.

\end{DndReadAloud}

Each of the quarters contains two bone devils and three cambions. If fighting occurs in one of these rooms, the occupants in the other room are likely to hear.

\subparagraph{Treasure}
Searching any of these rooms turns up 1d10 gp and 1d4 pp.

\subsubsection{B8: Officer's Quarters}
\begin{DndReadAloud}{}
There are three beds in this room for the submersible's officers.

\end{DndReadAloud}

\subparagraph{Treasure}
Searching this room nets 8d10 pp. There is a 25 percent chance that a patrol of 1d4 bone devils notice anyone looting one of these rooms.

\subsubsection{B9: Storage}
There are several mostly empty rooms dedicated to storing rations and nonmagical weapons.

\subsubsection{B10: Abigor's Bedchamber}
\begin{DndReadAloud}{}
Despite the cramped halls of the submersible, this is a reasonably sized room, with a large bed and a table with several chairs. Papers are scattered over the table.

\end{DndReadAloud}

\textbf{Abigor} is seldom here, as there is so much work to be done in maintaining the infernal submersible. The papers mostly detail the logistics of keeping the devil crew fed and disciplined.

\subsubsection{B11: Torpedo Room}
\begin{DndReadAloud}{}
The room is narrow and cramped with strange equipment. There are rows and stacks of enormous metal cylinders packed into the room. Each cylinder is long and tapered, with the pointed end facing dark tunnels carved into the vessel wall that are just as wide as the cylinders themselves. Glowing infernal runes cover each of the cylinders.

\end{DndReadAloud}

There are usually two bone devils present, but on a battle alert the room has eight bone devils working feverishly to keep torpedoes loaded and ready to fire. There are thirty-four cylinders, each a torpedo weighing 800 pounds. They're usually inert, but if the submersible is at a battle alert, then 1d4 of them have been primed. To prime a torpedo, one must press the correct combination of runes on its surface. This requires a successful DC 17 Intelligence (Arcana) check, for each torpedo (Fiends succeed on this check automatically). If a primed torpedo takes more than 10 damage it explodes, dealing 35 (10d6) force damage to every object or creature within a 30-foot-radius. Any primed torpedo damaged in this way explodes too. The damage from two or more torpedoes breaches the hull of the submersible, filling the room with the waters of the Styx within 2d4 rounds.

\subsubsection{B12: Depth Charge Room}
\begin{DndReadAloud}{}
There are racks of metal cylinders in this room. Each is squat and thick with a tapering rear that ends in a strange round fin-like shape. Glowing infernal runes cover each of the cylinders.

\end{DndReadAloud}

There are two bone devils present, whether or not the submersible is at a battle alert. There are ten depth charges, each weighing around 400 pounds. A depth charge must be armed before it is capable of detonating. To arm a depth charge, the correct combination of runes on its surface must be pressed, requiring a successful DC 19 Intelligence (Arcana) check for each depth charge (Fiends succeed on this check automatically). If an armed charge takes more than 20 damage it detonates, dealing 52 (8d12) lightning damage to every object or creature within a 60-foot-radius. Any armed depth charge damaged in this way detonates too. The damage from a single depth charge immediately breaches the hull of the submersible, filling the room with the waters of the Styx.

\subsection{Abigor}
\textbf{Abigor} is \textbf{Baalzebul}'s chief general, a warlike monster fit for the violence that the Lord of Deceit considers beneath him. Its preferred shape is a flayed-looking humanoid with musculature of polished metal. Physically, \textbf{Abigor} commands tremendous strength and destructive power. Its greatest joy, however, is to command a battlefield, ordering legions of infernal troops in complex stratagems against the enemies of the Nine Hells. In times past, \textbf{Abigor} has commanded the legions of Maladomini up in Avernus to stem the tides of the Abyss. Most of the time, though, \textbf{Baalzebul} keeps the monstrous creature close to guard against domestic threats or punish those underlings whose construction efforts disappoint their master, as they inevitably do. \textbf{Abigor} considers its talents wasted in such sideshows and its frustrated dreams of war manifest in counterproductive ways. Visitors sneaking through Maladomini might be surprised to find legions of devils standing about in parade ground order, or else marching and drilling, every maneuver perfect, yet their general is never satisfied. What the perfect city is to \textbf{Baalzebul}, military exercises are to \textbf{Abigor}. The archdevil forces its troops to march and salute and trudge through ever more complex motions as it plays at war. Alternatively, it throws its soldiers into vast war games across the ruinous country of Maladomini, decimating its own troops just to try out some new stratagem. \textbf{Abigor} is rightfully hated by all the devils under its command, and covertly mocked by its peers for its obsessions.

Getting close to \textbf{Abigor} without engaging legions of infernal troops is difficult—to \textbf{Abigor}'s regret, in fact. Sometimes the archdevil decides to test Maladomini's security by creating openings for imaginary assassins and spies, to see who might take advantage of them. On one occasion this resulted in a determined paladin getting far too close to \textbf{Baalzebul} himself, for which \textbf{Abigor} spent some decades being whipped by infernal rust monsters. After that incident, \textbf{Abigor} became obsessed with presenting a clean and martial appearance, employing a small legion of imps to burnish its skin to a blazing shine.

An actual physical run-in with \textbf{Abigor} (see the accompanying stat block) should be a simple, brutal experience but its relative isolation at the heart of Maladomini has left it curiously desperate to make the most of any fight it might get into. Rather than simply tear up the opposition, it constantly changes to more and more complex and convoluted tactics, desperate to make the very most of any melee. In doing so, it can end up missing obvious opportunities to destroy its enemies, or even inadvertently give them an advantage or let them escape. In one recorded instance, the archdevil even broke off the fight at its moment of victory to school its enemies in how they could provide a greater challenge, then insisted on beginning the fight anew.

\subsection{Assassinating Abigor}
If the characters have come to slay \textbf{Abigor} for the hag, \textbf{Anacreda}, when the archdevil is defeated, they feel the ever-pressing weight of the geas that has bound them magically lift. They've completed their end of the bargain and are free of the hag's influence.

\subsection{Capturing Abigor}
The characters don't need to keep \textbf{Abigor} alive for the Triad's ruse to work. They simply must defeat \textbf{Abigor} and capture the submersible before convincing the aboleth to lure \textbf{Baalzebul} into a rescue attempt. If they've done this, and the aboleth still lives, and the characters return to area B3, then read the following:

\begin{DndReadAloud}{}
As you prepare yourself for the arrival of the archdevil, \textbf{Baalzebul}, you watch the aboleth's slimy tentacles press against various glowing sigils. A long silence ensues, punctuated only by the occasional ping as the Styx's immeasurable pressure exerts against the submersible's metal hull. Then a sigil draws itself upon the floor, lines arcing and dovetailing, flaring bright and fading swiftly. For a moment a man with features exceedingly fine appears, in shadow, and then he shrieks, his voice deepening into a grotesque gurgle as the human-form shatters, expanding into a bulging mass of flesh, until a slug-like creature looms over you.

"W-what trickery is this?" And you recognize the voice, for you have heard it spoken from the mouth of \textbf{Koh Tam} when he was possessed by \textbf{Baalzebul} at the beginning of your venture. "My body... my b-beautiful body. What have you done to me?"

\end{DndReadAloud}

\textbf{Baalzebul} (see appendix A) is both confused at his sudden transformation back into his most hated form and in rage at the characters. He lashes out at them. The characters can either stand to fight him or attempt to escape. In either case, they've completed the mission the Triad assigned them and if they survive here, they should return to Nessus.

\chapter{Nessus, the Bastion of Asmodeus}
\DndDropCapLine{T}{}he final layer of the Nine Hells, bastion of \textbf{Asmodeus} and heavily guarded by his legions, is the cracked wasteland of Nessus. It floats within an infinite, bloodred void, which extends in all directions. Very few are given entrance into Nessus, with the patrolling legions of devils ordered to kill trespassers on sight, including other devils.


\section{Running This Chapter}
Before running this chapter read the "Nessus Overview" section. As this is the final stage of the characters' journey into the Nine Hells, they need to find a way out of their contracts. Make sure the players get the chance to complete their quest objectives before triggering the final stage.

The individual character's score on the Corruption Tracker (combined with their choices in this chapter) determines the possible outcomes available to them. The information in the "Final Stage" section helps you prepare for the final encounters in this chapter.

\subsection{Encounters}
The characters have entered the domain of \textbf{Asmodeus}. Even if they manage to travel the layer without setting off his defenses, they're likely to encounter complications while they're here. Before the characters enter either the Oasis of the Lethe or Malsheem, roll at least once on the Random Encounters in Nessus table.

\begin{DndTable}[header=Random Encounters in Nessus]{cX}

d6 & Encounter\\
1 & A singular \textbf{amnizu} (see Monsters of the Multiverse) approaches the characters with a smile on their face. It offers a quick passage out of Nessus in exchange for 2 Soul Coin per character. To each paying character it gives a wishbone, that when snapped, transports the character back to the Falls of the Frozen Titan in Cania.\\
2 & Two pit fiends spar with each other, testing their strength. They invite any travelers to spar with them. If they find their sparring partner to be lacking in strength, they quickly change their tune to harvest a wayward soul.\\
3 & A group of four horned devils patrol the skies of Nessus, looking for any potential intruders. If they spot any, they approach and ask to see a writ of passage. If the intruder can't produce valid evidence, the group attacks.\\
4 & Filled with rage by the sheer quantity of souls it absorbed, a \textbf{tyrant shadow} (see appendix B) rampages through the layer. Creatures of its size and strength are rare, and its death curries favor in the realm. If they kill it, the characters have advantage on all checks when interacting with the Triad. Additionally, when the tyrant shadow dies, 3d4 Soul Coin are created.\\
5 & A pride of six hellcats (see appendix B), led by a particularly menacing bezekira, stalks Nessus looking for prey. They avoid opponents they feel they can't defeat.\\
6 & Actively sneaking towards the center of Nessus is a single \textbf{balor} dressed in black robes. It is constantly vigilant and approaches any creature that notices its presence. Its infiltration was difficult to secure, and it goes to any length to ensure no devil learns of its mission.\\
\end{DndTable}

\subsection{Locations}
There are other locations of interest in Nessus, but the characters are warned against visiting them. They should remain focused on their mission and limit the time they spend in \textbf{Asmodeus}'s domain.

\subsection{Koh Tam and Tiax}
\textbf{Koh Tam} or \textbf{Tiax} (depending on who the characters sided with earlier) wait for the characters' return by the Falls of the Frozen Titan. Once they've gotten what they came for, they need the barge to travel the Nine Hells one last time, to freedom.

\subsection{Objectives}
Before entering Nessus, the characters should have resolved their patron's mission. Make sure that all phylacteries have been attained. Once they enter \textbf{Asmodeus}'s domain, they need to make haste and find a way to annul the contracts before \textbf{Asmodeus} prevents their escape.

\begin{DndSidebar}[float=!b]{Captured}
If the events from chapter 11 led to the Dukes capturing the characters on behalf of \textbf{Asmodeus}, they're brought before the Triad. They may not have all the necessary phylacteries to confront \textbf{Asmodeus} yet. However, all isn't lost as they may still save their loved ones. The scenario with the Triad plays out slightly different if the characters were captured or not.



\end{DndSidebar}

\section{The Final Stage}
There are two ways that the characters can approach the last stage of their quest:


\begin{itemize}
\item \textbf{Koh Tam} has an invitation to a place in Nessus known as the Oasis of the Lethe. In this place of learning, the characters can discover how to free their souls or the souls of their loved ones from the contracts that bind them. To get there, they must follow another infernal river, called the Lethe, which eventually brings them to the Oasis of the Lethe.
\item The characters can confront \textbf{Asmodeus} directly. To do this, they simply must enter Nessus, whereupon the Lord of the Nine senses their presence and sends a force to escort them to Malsheem. \textbf{Asmodeus} won't deal with the characters initially, instead he leaves them to his faithful servants known as the Triad. If the characters navigate the trials of these three archdevils, then they might be able to convince \textbf{Asmodeus} to release the souls from their contracts.
\end{itemize}

\subsection{Corruption Tracker}
Before setting out in Nessus, consult appendix E. The information there helps you guide your players through the finale of their adventure. Do the following now:


\begin{itemize}
\item Tally the individual scores for each character on the "Corruption Tracker."
\item Use the Corruption table to determine if the character has become corrupted.
\item Use the Final Outcomes table to determine the resolution available to each character.
\end{itemize}

\section{Nessus Overview}
Though the Styx freezes in its passage to Nessus, it doesn't completely dry up. Near the bottom, the temperature shift causes a small section of the river to melt and finish its journey to Nessus below. This stream is heavily guarded and flows a small distance before pooling into the Forgotten Lake, the last stop. From here, it drains into the porous stone of Nessus, traveling to the very bottom of the layer. It gurgles through a few of the deepest ravines before mysteriously vanishing on its journey to Gehenna.

An arid wasteland, Nessus resembles a rocky plain filled with cracks, pits, and canyons. The majority of the layer is flat, with only a few hills. Storms of fire and hot wind blow across its surface, and sandstorms aren't uncommon. Though many ravines mar its surface, only a few are important enough to have names. The Serpent's Coil is a massive ravine, rumored to contain \textbf{Asmodeus}'s true body and caused by his fall into the Nine Hells. Its depth is practically infinite, and no creature has ever gone from the surface of Nessus to the bottom.

A close cousin of the Styx, the River Lethe originates on Nessus. Some scholars believe its pools are fed by the Styx itself, but the river's unique powers and secluded nature make such beliefs impossible to prove. From its spring in Nessus, the Lethe flows across the land, occasionally crossing ravines and canyons, before reaching the edge of the landmass. Unlike the Styx, the Lethe completely defies gravity, flowing over the cracks in the land without breaking its stride. It flows off the edge of Nessus unimpeded, where it continues to run off into the void, seemingly floating into the empty space.

Largest of all cities in the Nine Hells is Malsheem, home to \textbf{Asmodeus} and his infernal legions. The city lies within the largest trench in Nessus and expands from its bottom to the skies, marking the center of the realm. From its spire, \textbf{Asmodeus} looks out on Nessus and, when the fancy strikes him, the rest of the Nine Hells. The only other structure of the layer, Tabjari, is a copper fortress built into the walls bordering Reaper's Canyon. An impregnable citadel, the structure is home to \textbf{Asmodeus}'s elite forces and acts as a vault, safeguarding his most valuable possessions.

All devils native to the Nine Hells, those that don't arise from the souls of the damned, are "born" on Nessus. \textbf{Asmodeus}'s true form continues to bleed from his wounds, and from this blood sprout devils. They're quickly redistributed to various other layers, depending on need and duty, with \textbf{Asmodeus} keeping only the most powerful. Each archdevil answers to \textbf{Asmodeus}, the Overlord of the Nine Hells, who rules from the infernal throne in Nessus. They all plot and hope to gain the throne for themselves, yet \textbf{Asmodeus} has ruled the Nine Hells since their inception, and likely will continue to do so until the end of time.

Entrance into Nessus is nigh impossible, and rarely done. Those that choose to enter the realm do so by \textbf{Asmodeus}'s invitation or are almost certainly destroyed by the realm's defenses. As the lowest layer in the Nine Hells, exiting Nessus is only possible by ascending into Cania. Malsheem certainly contains portals of some kind within its walls, though their location is known only to \textbf{Asmodeus} and his guards—as is their destination.

\subsection{Features}
\textbf{Asmodeus} knows of all that enter his realm. Without some method of concealing their presence, \textbf{Asmodeus} always knows the relative location of each creature in Nessus. Visitors feel a constant state of paranoia, as if something is always watching.

\subsection{Climbing Down into Nessus}
\begin{DndReadAloud}{}
The great gray swell of the River Styx begins to slow. Here and there great chunks of ice appear and are pushed along by the churning waters until they collide and begin accumulating into one mass that pushes up against the riverbanks. The ice becomes thick enough to stand on, but the rush of moving water below can be heard. Eventually this sound slows to a gurgling whisper and then falls silent. The ice stretches out for a distance, solid and white with arching crests where swells and rapids have frozen, and then it ends, and all that remains is empty sky. At the edge of the ice one can see out over the barren Nessus plains, vast and sprawling, reaching the far horizon. The drop below is dizzying. The powerful waters of the Styx have tumbled over the edge of the cliff and frozen into an enormous curtain of solid ice. The ground is scarcely visible, hazy with distance. Trapped within the plunging frozen waterfall is a titanic skeleton of humanoid shape dressed in gleaming armor. A giant.

\end{DndReadAloud}

To reach Nessus, the characters must climb down the body of an ancient giant, half submerged in a frozen waterfall that is the last stage of the River Styx. This skeletal figure, of a scale beyond the stature of mortal Giants, is armored in gleaming mail that remains pristine and shimmers within the ice.

Although no more than bones, some trace of life still clings to the frozen titan. Those approaching hear its voice in their head, speaking whatever language is most native to them. It names itself as \textbf{Devorastus}, a warrior of law, who came to the Nine Hells unknown centuries ago and, after slaying countless devils, ended up imprisoned in the eternal ice of Cania.

\textbf{Devorastus} asks all passers-by to work towards freeing it, even if it is only to chip a few handfuls of ice away from its frozen tomb. Its words are persuasive, and an ability to read the thoughts of others permits it to tempt them with targeted inducements appealing to their morality, offering bribes, or pledging services.

\begin{DndReadAloud}{}
There is a furious commotion around the titan's head. Mortals and devils brush shoulders as they chip and claw away at the ice around it. Some seem to be working in groups, others alone, but each acts with the same desperate fervor.

\end{DndReadAloud}

Over time a cult of desperate dissidents has arisen around the ice-locked head of \textbf{Devorastus}. Some cultists are mortal travelers who have nothing else in the plane save their hope that the giant will somehow save them. Others are actual Fiends, outcasts or rebels fleeing \textbf{Mephistopheles}' ire, who intend to use the skeletal giant as their war machine to fight their former master. Thus far the archdevil is either ignorant of the cult or considers action against it not worth the resources—or it is simply amused by its folly. The score or so of cult members pick feebly at the ice covering, attempt to light fires that the cold of Cania kills, and make sacrifices in the hope the warm blood and shed life force will help free their champion.

\textbf{Devorastus}' vast armored form is plainly a long way from being free, and to chip a little ice away might seem a meagre price to pay to avoid the predations of the cult. However, \textbf{Devorastus}' true nature is very different, and the creature is far closer to escaping than it looks. \textbf{Devorastus}' actual form is similar to a gelatinous cube, currently wrapped around the bones of the titan. It is a demon, a lieutenant or perhaps rival of Juiblex, which travelled down the swift flow of the Styx before being trapped in the ice. As soon as even a small hole in the ice pierces through to the chamber it is encysted in, the creature will surge out, feed off the flesh of mortals and devils alike and then seek to return to the river in the hope it can follow the current to more hospitable climes.

Characters must expend significant energy to free \textbf{Devorastus}. Track each hour of active digging/excavating, per character, with every 30 points of fire damage counting as an additional 1 hour of effort. Once this tracker reaches 50, \textbf{Devorastus} (see the accompanying stat block) is freed.

\begin{DndReadAloud}{}
Though your efforts have only produced a small, albeit deep, channel in the ice, there's a disturbing hissing noise, as if water has found its way to a tiny crack in a massive dam. A surge of pale liquid bubbles out from the hole you have created, almost instantly becoming a huge mound of ooze sticking to the ice. The trapped titan collapses as if whatever once sustained it, has escaped.

\end{DndReadAloud}

The massive ooze-like demon, barely visible against the ice, climbs up or down as necessary, attempting to consume anything nearby. If no enemies are detected within 60 feet, it begins crawling upwards, in the hopes of finding a more comfortable resting place.

Climbing down the titan isn't without its dangers. Shredwings roost in parts of the body of the titan. 2d6 of them attack anyone descending, but once half of them are killed, the rest retreat.

\subsection{Entering Nessus Without Protection}
If the characters enter Nessus without any protection against divination, then they are sensed by \textbf{Asmodeus} immediately. He sends a force to escort them to Malsheem.

\begin{DndReadAloud}{}
From across the arid plains there comes a distant roar. The ground trembles slightly and a hot, sweeping wind follows, sending up a curtain of red dust. A shadow passes suddenly over the land ahead of you. It is immense and slips over the distance with alarming speed. A second shadow joins it. There is another ground-shaking roar and then the sound of snarling and the rattle of weaponry. A massive fiend approaches, red, winged, and armed with a heavy metal mace. Following it are four equally large devils with dark, leathery skin and wings and six devils with enormous horns carrying tall forks with sharpened tips that glow with heat. Flanking this monstrous party are two dragons. The scales of the beasts' gleam a brilliant red and the cutting wind born of their powerful wingbeats carries with it the scent of smoke and sulfur.

\end{DndReadAloud}

The escort is led by a \textbf{pit fiend} and includes four war devils (see appendix B), six horned devils, and a pair of adult red dragons that the characters are to ride if they can't fly themselves. Even the most powerful characters will find this to be a difficult fight if it comes to violence, but the devils have been told to simply escort the characters, not imprison them. They only attack to defend themselves or if the characters refuse to accompany them. Proceed to the "Adventure: Malsheem" section.

\subsection{Traveling Across Nessus}
Those traveling across Nessus can stick close to the River Lethe or strike out into the plains to avoid the traffic of devils that can be found along the river's edge. Below are two possible encounters.

\subsubsection{The River Lethe}
The River Lethe is, if possible, less navigable than the River Styx. In many places the river plunges over vast, wild falls, and devils have had to laboriously build complex sequences of locks to permit their great river barges to safely traverse the drops.

On occasion, political feuds between devils spill over into direct sabotage and action. Consignments of valuable goods—souls, sorrow wine, magic artifacts, and mortal prisoners—are wrecked or sunk as one diabolic magnate strikes indirectly against a rival. The falls of Telkalal are a favorite location for such skullduggery.

\begin{DndReadAloud}{}
The crew of a merchant's barge is desperate, running around and tripping over each other as they try to catch grappling hooks on the jagged rocks that line the top of the falls. They have little success and the few lines that draw taut snap under the weight of the barge and the strength of the water. Terrified screaming can be heard from the cargo hold. The barge jolts as another line snaps and draws closer to the edge of the falls. The screaming increases in volume. The drop of the falls is monstrous and there are towering rocks below, cutting out from the white spray like hungry teeth.

\end{DndReadAloud}

The barge is carrying a tithe from the archdevil Orvrastus, intended for Nessus. The loss of this tribute will bring Orvrastus into bad odor with \textbf{Asmodeus}. A variety of pettifogging bureaucratic obstinacies have resulted in \textbf{Fierna} of Phlegethos becoming irritated by Orvrastus, and she has sent a team of four \textbf{erinyes} to haul the barge astray. Instead of travelling into the calm waters of the locks, the erinyes are attempting to drag the vessel over the falls, while the infernal crew back the engines and launch missiles at their flying tormentors.

A dispute between devils may not be of overt interest to mortals such as the characters, but the treasures on board may well be—whether they constitute people to be rescued or valuables to be pilfered. In the chaos of the fight, it might be possible to get on board the ship, loot it, and get off before it goes over the falls.

Alternatively, the characters could assist the vessel's demise, and then hotfoot it below to sieve the River Lethe for spilled treasures. In doing so they must contend with not only the usual denizens of Nessus, but habitual mobs of infernal salvagers intent on acquiring the same riches.

\subparagraph{Treasure}
If the characters are successful, they obtain 4d4 x 1000 pp and one rare magic item from the \textit{Dungeon Master's Guide}. The characters also find 2d4 Soul Coin.

\subsubsection{The Plains of Nessus}
There is a saying in the Nine Hells, that nobody works harder than \textbf{Asmodeus}. The lord of all the Nine Hells neither sleeps nor rests, and whilst most of his servants devote plenty of time to sating their hedonistic pleasures, the greatest of all Fiends constantly bends his mind to the corruption of the planes and the ensnaring of souls.

Even \textbf{Asmodeus} permits himself some entertainment from time to time, though his joys are little different to his duties. At times he enjoys hunting some of the more formidable captive souls in his possession.

Those he torments personally are, of course, lost beyond any hope of retrieval. Not one of them will ever escape his grasp. However, the finest razor edge of torment is that he persuades each one of them that there is a hope of reprieve. For his hunts, he often likes to allow his prey to think they've broken out of the prisons of Malsheem and that they've some sort of chance at freedom.

\begin{DndReadAloud}{}
A large, armored wagon, drawn by a colossal monster, careens wildly over the plains. Overhead, thin figures with feathered wings that still shine with faint celestial light swoop and fall.

\end{DndReadAloud}

In this case, \textbf{Asmodeus} has permitted a gang of desperados to stage a breakout, purely for the joy of hunting them down.

The supposed escapees are a selection of some of the worst the mortal planes have had to offer over several centuries. Their current leader—though they fight among themselves a lot—is an orc warlord named Ormogg (neutral evil, orc \textbf{assassin} with 120 hit points), whose hordes despoiled cities and burned forests and flooded dwarven holds with magma, but they've great sorcerers, wicked priests, tyrants and murderers among their ranks. As befits prisoners of the Nine Hells, each one is irredeemably evil. They've stolen—been permitted to steal—a vast armored wagon drawn by a monstrous beast, and now they ride it across the barrens of Nessus with a ragged mob of devils hounding their heels. \textbf{Asmodeus} watches on, greatly amused, until he decides to send in his special hunters.

\begin{DndSidebar}[float=!b]{An Old Nemesis}
This might be an opportunity to resurface villains from earlier in the character's adventures before they ventured into the Nine Hells, having them side with Ormogg in this prison break. Alternatively have the prisoners be comprised of an \textbf{archmage}, four priests, six gladiators and eight veterans. The creature pulling their wagon is a \textbf{miasmorne} (see appendix B) but it isn't truly under their control, at any moment it might turn upon them.



\end{DndSidebar}

\textbf{Asmodeus}'s former captives are well used to the torment of Fiends. His Bright Hounds are something else. The breaking of angels is an especial pleasure to the lords of the Nine Hells. \textbf{Asmodeus} revels in perverting them to his while, twisting them into cruel falcons. He has stables of them, no longer resembling angels but their radiant power might still bite the hand that unhoods them. They represent an additional test of his mastery that makes the whole hunt truly interesting.

The Bright Hounds are comprised of four devas and two planetars (lawful evil and delivering necrotic damage instead of radiant). The characters, encountering the wagon on its headlong flight, might get involved in a variety of ways. The wagon's inhabitants, believing they know a way out of Nessus (which conveniently is at the Oasis of the Lethe, but is also just a lie planted by \textbf{Asmodeus}) may offer them a place aboard in return for fighting off the devils—only to be revealed as monsters almost as bad as the devils outside. The angelic hunters may provoke pity, though there is little the characters can do for them save end them.

\section{Adventure: The Oasis of the Lethe}
Everything within Nessus' bounds is the unchallenged domain of \textbf{Asmodeus}, yet the Oasis exists just slightly beyond. Enough to satisfy the legalistic minds of devils, perhaps. The River Lethe runs through Nessus, heedless of topography, reaches the edge of the wasteland and just keeps going. It flows forever into the red void, eerie, discomfiting, and unexplained. The greatest infernal minds can't account for it, nor any god that has been petitioned for an answer. What lies at the end of the Lethe? Does it empty out into another constellation of planes? The question has ever troubled philosophers and cosmologists. And so, on the banks of the Lethe, a community of scholars and wizards has arisen to study this phenomenon and, through it, the underlying nature of the planes themselves. The Oasis of the Lethe is their haven.

\subsection{Advice from Koh Tam}
\textbf{Koh Tam} previously gave the characters an infernal map. Show the players the map of the Oasis of the Lethe in appendix F. They can use this map to decide which points of interest they want to explore. Refer to the "Using an Infernal Map" section of chapter 2.

\subsection{Monsters of the Red Void}
There are certainly threats that come to trouble those at the oasis. Inexplicable monsters sometimes undulate out of the red void—bizarre beasts never seen before and never replicated. Whilst they represent unparalleled opportunities for study, they also devour the occasional student. Desperate fugitive souls and marauding adventurers fleeing \textbf{Asmodeus}'s pursuit occasionally find their way to the Oasis and make demands of the arch-magister Elihandrel, though thus far his magical abilities have sufficed to return them to the archdevil's custody in short order. Sometimes conjurations called up by the mages themselves get loose and have a brief career of havoc before they're banished once more. But from the devils of the Nine Hells themselves, there is no threat. They're as obliging as the most unctuous servants.

\subsection{Oasis of the Lethe Locations}
Some of the major locations found in the oasis are described below.

\subsubsection{O1: Arrival}
\begin{DndReadAloud}{}
The Oasis is a floating island seemingly tethered magically to the endless flow of the River Lethe, a mile out from Nessus' edge. In contrast to the barren plane, it is green and verdant, thronging with life. There are trees meticulously trained together to form living buildings; there are pools of sparkling water that promise refreshment. Animals and birds dwell here, apparently without fear of mortal or devil. Lions and lambs, serpents and deer, all living in harmony.

\end{DndReadAloud}

Within the buildings of the Oasis there are libraries that would make a sage weep, telescopes and scrying pools and all manner of divinatory instruments to look into Nessus and the river and the wider planes. There are debating rooms and places for quiet contemplation. Any mortal who has ever valued the pursuit of knowledge instantly feels that they've come home.

\subparagraph{The Researchers}
Some of the planes' greatest minds can be found here studying the intrinsic nature of the cosmos—both the living who have made the arduous journey and the souls of the dead. \textbf{Asmodeus} sporadically signs orders releasing the most learned of scholars within his domain from their torments so that they may aid in the research of the Oasis. The prospect of questioning the great minds of elder days is enough to bring some visitors to the community here.

\subsubsection{O2: Home of Teskethmus}
One of the most respected minds of the current roster of scholars at the Oasis is that of Aldri Teskethmus, a living human \textbf{archmage} of the Red Wizards, who theorizes that the traditional arrangement of the planes is merely an artifact of mortals' limited ability to comprehend the wider universe.

\subsubsection{O3: Home of Orvedarc}
Orvedarc is a \textbf{duergar} heretic priest-scholar who died centuries ago and has spent time in the Nine Hells both in torment and in service to devils. Now released, he works on his incomplete research into the nature of divinity and divine magic.

\subsubsection{O4: Home of Isitrix}
Isitrix the Termagant is a sahuagin magician-queen (\textbf{mummy lord} but she isn't Undead, her creature type is Humanoid) who in life brought wars that encompassed entire oceans, but who now peaceably studies the River Lethe and its curious mechanics.

\subsubsection{O5: Home of Elihandrel}
It is an inevitable fact that those long-dead souls at the Oasis are mostly of evil persuasion. There are other incautious mortals who fall into \textbf{Asmodeus}'s power, but the majority of souls in his domain are those drawn to the Nine Hells by the choices they made in life. However, the arch-magister of the Oasis, Elihandrel (use \textbf{empyrean} stats but Elihandrel is Medium sized), keeps order. His word is law, as much as \textbf{Asmodeus}'s is within Malsheem. Those who won't abide by his codes of peace and scholarly amity know that a return to the Nine Hells awaits.

Elihandrel himself is an ageless high elf, robed in simple grey. He appears mild, confident in his authority and power, proud of the peaceful community he has built. His soft-spoken word is the end of any disagreements, ensuring that those under his charge never let civilized debate get out of hand. His guests have everything they need—the trees bear endlessly nutritious fruit and there are always books, tools, and scholars' supplies of surpassing quality. He even sanctions research expeditions to other layers of the Nine Hells. At his beck and call are a staff of curiously meek devils, still monstrous of aspect but on their very best behavior. Elihandrel himself reports to \textbf{Asmodeus} on a regular basis, sharing all the knowledge that his faculty uncovers. He denies that he is in any way enabling evil. Learning, he insists, is without alignment. His ultimate goal is to push past the adversarial model of the universe that has trapped all life in constant conflict, law against chaos, good against evil. He genuinely believes that a full understanding of the planes can lead to a perfect future without hardship or sin.

Elihandrel is very eager to spread word of his faculty out into the wider planes. He maintains correspondence with a number of temples, universities, and mage schools on the mortal planes. Whilst the appeal of studying the secrets of the Nine Hells is plain to evil mages and luminaries of all kinds, Elihandrel is very keen to maintain balance by attracting those of other aspects as well. He writes persuasively of the possibility of studying the heart of evil with an eye to countering its influences. One must truly understand a thing before one can defeat it, is his frequent refrain. Scholars visiting the Oasis from the mortal planes find the opportunities of the place overwhelming. There is always something new to learn—and not just miserly scraps but the great truths of the universe hanging just out of reach like fruit ready to be plucked. Conversation with other resident scholars is transcendent, opening the doors of the mind. Very few ever want to leave once they arrive.

\subsubsection{O6: The Library}
\begin{DndReadAloud}{}
There are four grand hallways in the library. Each so vast and tall that neither their ends nor their ceilings can be seen. Magnificent marble pillars line the hallways, dividing them into sections, and each section stretches off until it is lost in distant shadows. There are upper levels, and these also grow faint as they rise beyond where the eye can see. Elaborate, curling staircases descend underground, suggesting further levels below, and rise upward until they vanish. Every available space, every nook and cranny, every shelf and row, is filled with books. Books, packed tight in neat rows and layered in towering stacks, stretching on, it seems, to eternity.

\end{DndReadAloud}

The knowledge contained within the library is unrivaled. An \textbf{amnizu} (see Monsters of the Multiverse) is the caretaker of the library, but also fully empowered to modify or annul contracts. Before the characters can ask for his assistance however, they need to prepare their arguments. The legal documents section of the library is where the characters find books that contain the clauses they need to build their case. However, if they want to look around the library before getting to work, there are many other sections of interest at the Library of Infernal Knowledge. Once the characters are ready to start searching for the clauses needed to annul their contracts refer to the "Legal Documents" section.

\subsection{Infernal Knowledge}
Magic books and other knowledge can be acquired here. But if all these seem too good to be true, that is because they are. Characters should feel uneasy as they begin to acquire items and knowledge; everything has an ephemeral feeling to it.

\subsubsection{Personal Growth}
All these books work exactly as described in the \textit{Dungeon Master's Guide}:


\begin{itemize}
\item \textit{Manual of Bodily Health}
\item \textit{Manual of Gainful Exercise}
\item \textit{Manual of Quickness of Action}
\item \textit{Tome of Leadership and Influence}
\item \textit{Tome of Clear Thought}
\item \textit{Tome of Understanding}
\end{itemize}

\subsubsection{Magical Power}
There are various arcane scrolls and books here.

\subparagraph{Spellbooks}
Several tomes contain descriptions of previously unknown 9th-level spells apparently more powerful than any known spell. With the right components, and a week of study each, these writings can be turned into actual spells that could be used by wielders of arcane magic. These spells include:


\begin{description}
\item[The Killing Winds.]
These winds create a hurricane 5 miles in radius that deals necrotic damage to every living being within.

\item[The World Weave.]
This spell changes the climate in a region permanently.

\item[Ioulaum's Longevity.]
This spell extends the caster's life by snuffing out every life in a 1-mile radius. Each Humanoid life consumed adds a year to the caster's life.

\item[Volcanic Eruption.]
This spell creates a volcanic eruption that lasts for days and devastates a region 100 miles in radius.

\item[Move Mountain.]
This spell allows the caster to cut the top off a mountain and use it as a floating platform upon which a city or fortress might be built.

\item[Spheresail.]
This spell allows the caster to create a spelljammer.

\end{description}

\subsubsection{Religion}
Several religious tracts focus on topics pertaining to divinity and describe new and powerful 9th-level spells. After a week of study and communion with the source of their divine magic, clerics and druids can learn the following 9th-level spells:


\begin{description}
\item[Dire Winter.]
This allows the caster to bring about a winter across an entire continent that lasts for 2d12 years.

\item[Vengeful Gaze of the Gods.]
A powerful single target spell that simply annihilates anything that doesn't possess legendary resistance.

\end{description}

\subsubsection{Item Crafting}
There are various other tomes here too:


\begin{description}
\item[The Tomes of Asgardian Metal Working.]
Describes how to forge the Hammer of Thunderbolts and various Belts of Giant Strength in half the normal time and at half the cost. Additional margin notes suggest ways to make these items even more powerful.

\end{description}

\subsection{Legal Documents}
It is here, in the legal section, where the characters find the loopholes needed to annul the contracts. Each hour that is spent searching requires a successful DC 20 Intelligence (Investigation) check to find the legal books that are needed to annul one of the contracts.

Once a character has found a loophole, they can ask the caretaker to annul the contract. The caretaker is only too happy to oblige, commenting on the vigorous and impressive research the character has done to prepare their case. However, a successful DC 15 Wisdom (Insight) check informs a character that it might be best to wait to trigger an annulment until all the loopholes required for all the characters are found.

\begin{DndSidebar}[float=!b]{Corruption}
All of the characters can find the loopholes to their contracts within the library. However, for any characters who are corrupted during their time in the Nine Hells, annulling the contracts isn't the end of their troubles. Use the information in appendix E to guide the characters through their next steps.



\end{DndSidebar}

\subsubsection{Between a Rock and a Hard Place}
Read the following only if there are corrupted characters among the group. It should be read aloud when the caretaker reviews the contract of a corrupted character.

\begin{DndReadAloud}{}
The caretaker's enthusiasm for your meticulous research is noticeable. His eyes light up each time he reads yet another cleverly worded part of your defense. After several minutes he is interrupted by an attendant that whispers something in his ear. Whatever is said, it can't be good. The caretaker's shoulders slump and the devil looks sincerely deflated as he addresses you. "All for naught, it was all for naught. The Lord of the Nine Hells saw you coming, he did. I gather it was all a big game for him. I can annul the contract, yes that I can do. But I cannot undo the evil you have done during your stay in the Nine Hells. Your soul... it is truly lost. Only the Triad can help you now."

\end{DndReadAloud}

The caretaker informs the corrupted characters of their predicament. If they want to improve their own fates, they have one more stop to make: they need to plead their case before the Triad at \textbf{Asmodeus}'s palace in Malsheem.

\subsubsection{Escape}
As soon as the caretaker undoes one of the characters' soul contracts, \textbf{Asmodeus} realizes what is happening and initiates a purge (see "The Purges" section for more information). Every minute that the characters remain at the Oasis, they must roll on the Random Encounters at the Oasis table.

\begin{DndTable}[header=Random Encounters at the Oasis]{cX}

d6 & Encounter\\
1 & Nothing\\
2 & A \textbf{pit fiend}\\
3 & Five horned devils\\
4 & Three maelephant nomads (see appendix B)\\
5 & Two war devils (see appendix B)\\
6 & Nothing\\
\end{DndTable}

Once the characters have escaped the Oasis, they're relatively safe as long as they continue to use the magic that keeps \textbf{Asmodeus} from knowing their location. If some of the characters are corrupted, they must travel to Malsheem and try their luck with the Triad. Otherwise, proceed to the "Leaving the Nine Hells" section at the end of this chapter.

\subsubsection{The Purges}
There is, of course, a problem with the Oasis. This is Nessus, after all.

The deal that Elihandrel made with \textbf{Asmodeus} is known to almost none. Nobody really gets a chance to talk about it once the archdevil collects. Despite being beyond Nessus' physical borders, the Oasis remains in \textbf{Asmodeus}'s domain, and Elihandrel must pay rent for his tenure there. At irregular intervals, carefully calculated by complex infernal clocks, the forces of the Nine Hells visit the Oasis and sweep it bare of everyone save Elihandrel himself. Every scholar and soul is hauled off to Malsheem, to be dissected for all they know and then put to eternal service or suffering as \textbf{Asmodeus} prefers. To study at the Oasis is to tacitly permit the Nine Hells a claim to your soul, and over the centuries the archdevil has drained the best and brightest minds of the planes by indulging Elihandrel's little project. The arch-magister himself considers the deal unfortunate, but a small matter compared to the vast reams of learning that his charges uncover. After all, the libraries remain intact. Nothing, therefore, is truly lost.

Read the following narration aloud when a purge occurs:

\begin{DndReadAloud}{}
There is a powerful crashing sound somewhere in the library, followed by a roar and a piercing scream. The scream is joined by another, and then another, and the roar too gains company until the cold, marble halls echo with a cacophony of gleeful snarls and beastly roars and terrified, pained screaming and wailing shrieks.

\end{DndReadAloud}

\begin{DndSidebar}[float=!b]{The Forgetful Scholars}
Any knowledge or spells that a character learned at the Oasis is lost within 2d12 hours of leaving. This includes attribute adjustments from books such as the Manual of Bodily Health.



\end{DndSidebar}

\subsubsection{Knowledge That Doesn't Travel}
The learning of the Oasis, that seems infinitely convincing to all the scholars who debate within its leafy halls, doesn't travel. Like a wine of uncertain vintage, the further you carry it from its region of origin, the less satisfying the taste. The cast-iron logic of Elihandrel and his fellows weakens, or else becomes abstruse, impossible to follow. It all sounds very intelligent until someone starts picking at the details, and then it seems to mean nothing at all. \textbf{Asmodeus}'s last joke on the Oasis is that almost everything that those serious and high-minded scholars write down is complete nonsense once taken away from the curious influence of the Oasis itself. Almost everything. The final barb of the archdevil is that every so often something that Elihandrel disseminates to the wider planes turns out to be of incalculable value, a true insight for the ages, perhaps even a weapon that can be turned against the Nine Hells. Those little scraps of true value mean that even those scholars who suspect the true nature of the Oasis still hear the siren call of temptation. After all, are they not the greatest minds of the age? Surely, they will be shrewd enough to outwit the trap the Oasis represents. And \textbf{Asmodeus} smiles to himself as yet another great sage descends to Nessus clutching one of Elihandrel's invitations, because there is no pride sweeter than the folly of the wise.

\section{Adventure: Malsheem}
Read the following when the characters approach Malsheem, either on their own or escorted by devils.

\begin{DndReadAloud}{}
Malsheem is a behemoth of a city. It stretches from the bottom of a deep, jagged trench to the heights of Nessus' blood red sky, its towers and flagpoles stabbing like spears through the thick, roiling clouds overhead. In width it is also massive, swollen so greatly that the giant trench in which it lies seems a tight fit. There is a strange and frightening beauty to the place. Even from a distance you can see that the city's walls are uniform and polished, the curling spires of its palaces are precise and elegant and everywhere there is the gleam of precious metal and the light of burning fires.

\end{DndReadAloud}

There are several ways the characters might arrive at the city: captured upon entering Nessus and taken directly to Malsheem, brought here as captives of the Dukes—some brave characters may even choose to walk into the city themselves. No matter how they get here, they always enter through the city gates of Malsheem. Here they're met by an escort of two pit fiends and eight horned devils who announce that they're to be guests in the palace of \textbf{Asmodeus}. The characters can go with them or attack. If they attack, the devils are joined by another pit fiend on initiative count 15 of every round. Once defeated, the characters are deposited at the palace for an audience with The Triad.

\subsection{The Palace}
When the characters arrive at the palace, they are confronted by the three archdevils known as the Triad.

\subsubsection{The Triad: Rimmon, Zagum, and Buer}
This triumvirate of bloated devils were angels once, who fell alongside their master back at the dawn of all things. Rather than committing to taking arms against the gods or embracing their fate, they attempted to play lawyer with the divine. For an age they sought to escape their fate by exploiting technicalities in what it truly meant to rebel or be counted among the supporters of \textbf{Asmodeus}. Eventually, their legalistic hair-splitting failed, and the gods delivered them to their master, who viewed them with utter contempt. He transformed them from their radiant angelic forms into great floating sacs of flesh and made them his notaries, so they could exercise their lawyerly talents in his service.

The Triad might appear homogenous to mortals at first sight, but they have different natures and predilections. Just as deities of fate often have a threefold nature, spinning, weaving, and cutting the strands, so the Triad are an equivalent for the contracts of the Nine Hells:


\begin{description}
\item[Rimmon.]
(see the accompanying stat block), vast with weaving tentacles protruding from his form, drafts the twisted wording of \textbf{Asmodeus}'s agreements—festooned with hidden loopholes, and buried subclauses to bind the unwary while denying them the rewards they truly seek. He has a fiendishly creative mind, fond of games, wordplay and inordinately long sentences with clause nested within clause until any sense of meaning eludes the reader entirely.

\item[Zagum.]
(see the accompanying stat block), is the interpreter of contracts, who takes the complexities \textbf{Rimmon} creates and uses them against the signatories, arguing the case of the Nine Hells in courts across the planes and ensuring that even the wiliest soul can't escape \textbf{Asmodeus}'s reach. He is a dry, thoughtful devil with a sense of humor. In those rare occasions where a mortal advances a line of argument that interests or surprises him, he has been known to recruit them to his staff of twisted legal minds—as an alternative to having \textbf{Asmodeus}'s eternal enmity. While \textbf{Rimmon} is an introverted wordsmith, \textbf{Zagum} is gregarious, happy to share a joke and a drink with a mortal he is about to shackle to eternal torment.

\item[Buer.]
(see the accompanying stat block) is the torment. Her role is to enforce the contracts once \textbf{Zagum} has proved their binding nature beyond any doubt. Her hideous maw can chew over souls eternally and she is utterly without mercy, the inflexible lash of the law. When the Triad are forced to fight, it is \textbf{Buer} who brings the physical threat while the other two rely on magic to support her. Their greatest weapon, in the fray, is that those who are bound by infernal contracts are automatically vulnerable to their powers. In the small print of all their agreements is a clause of implied consent, so that every spell cast by the Triad can evade any attempt at resistance, and any physical defense comes to nothing. When you sign the contracts of the Nine Hells, you really are agreeing to get everything you deserve, whether you want it or not.

\end{description}

\subsubsection{The Contracts}
The Triad acknowledges that the characters have likely come to Malsheem to nullify the contracts that bind either them or one of their loved ones. They have the power to nullify the contracts, but only if the characters give an item of value or provide a service.

\begin{DndSidebar}[float=!b]{Unfinished Business}
In the event the characters come before the Triad after being captured by the Dukes, they may not have gathered all needed phylacteries. It should be clear to them that performing a service for the Triad is their best option. If they agree to do this, they're allowed to leave the palace. They can collect the phylacteries, capture \textbf{Abigor}, and resume negotiations.



\end{DndSidebar}

\subsubsection{A Service to Be Provided}
The characters could opt to trick \textbf{Baalzebul} into leaving Maladomini, so that \textbf{Asmodeus} can turn the archdevil back into his slug form. To do this they must capture \textbf{Baalzebul}'s trusted lieutenant \textbf{Abigor}. \textbf{Asmodeus} knows exactly where the infernal submersible that \textbf{Abigor} captains can be found. Once they capture the submersible, they must use the aboleth within the infernal submersible to lure \textbf{Baalzebul} into trying to rescue \textbf{Abigor}. When \textbf{Baalzebul} turns into a slug and realizes he has been tricked, he does everything in his power to kill the characters.

\subsubsection{Items of Value}
Alternatively, the characters might offer items of value:


\begin{itemize}
\item Two very rare magic items or one legendary magic item are appropriate payments.
\item One of the characters' life and soul.
\end{itemize}

\subsubsection{Leverage}
The characters can avoid the previous options if they've come across knowledge that gives them leverage over \textbf{Asmodeus}:


\begin{itemize}
\item If the characters have discovered the plans that \textbf{Asmodeus} gave to the Brothers Morax and Adramalech, then they can use this as leverage. However, if they have the plans on them, then the Triad reminds them that those can easily be taken off of their dead bodies.
\item If they learned the secret of \textbf{Asmodeus}'s origins from \textbf{Abigor}'s aboleth, that is significant leverage.
\end{itemize}

\subsubsection{Negotiating}
This final negotiation is influenced by the choices the characters made during their time in the Nine Hells. If they're corrupted their options are severely limited. However, if the characters managed to keep their souls relatively untarnished, they have more options and may be able to outmaneuver the Triad. For each individual character, check the Final Outcomes table in appendix E to see which options are available in the negotiation.

\subparagraph{Setting the Terms}
A corrupted character who came to the Nine Hells to free their loved one can choose to add one of the following clauses to their new contract:


\begin{itemize}
\item You sacrifice your loved one to ensure a better position in the Nine Hells for yourself. Your loved one is returned to their torment. However, you're free to live out your life as an agent of \textbf{Asmodeus}. After you die, you arrive in the Nine Hells as a pit fiend. After you sign this contract, you may choose any 2 of the epic boons available in the \textit{Dungeon Master's Guide} (a successful DC 20 Charisma (Persuasion) check can increase this to 3).
\item You commit to becoming an agent for \textbf{Asmodeus} for the rest of your life. In exchange, the soul of your loved one is released from their contract. After you die, you arrive in the Nine Hells as a horned devil.
\end{itemize}

A corrupted character who came to the Nine Hells to free their own soul has limited leverage with which to negotiate. They can try to improve their situation in the afterlife, however.


\begin{itemize}
\item You agree to become an agent for \textbf{Asmodeus}. This guarantees you arrive in the Nine Hells as a horned devil upon death. A successful DC 20 Charisma (Persuasion) check upgrades this to a pit fiend.
\end{itemize}

Uncorrupted characters have a better position in these negotiations and can easily convince the Triad to release their loved ones or themselves from their bonds with this new contract.

\subparagraph{There Is Always a Catch}
Once terms have been agreed upon, the Triad writes up the contract. However, there are some obvious loopholes that could be exploited by \textbf{Asmodeus}. There are three loopholes in total, one written by each of the Triad. A successful DC 20 Wisdom (Insight) or Intelligence (Investigation) check reveals the first loophole and alerts the characters that they're being tricked, lowering the DC for finding the second and third loophole to 15. Once discovered, the characters can succeed on a DC 20 Charisma (Persuasion) check to have the loophole removed—or they can threaten violence. If they're violent, then whoever they've threatened attacks. If it wasn't \textbf{Buer} that they threatened, she helps in the attack. The other(s) stand back and watch. At half hit points the hostile devil accepts defeat and removes the loophole.

Once the contract is finished, the characters are told to take it to \textbf{Asmodeus} for his signature.

\subsection{Audience with Asmodeus}
An audience with \textbf{Asmodeus} (see appendix A) is something that very few mortals have ever had. \textbf{Asmodeus} doesn't want to meet the characters just to sign their contract. He could have let the Triad handle that. He always has another angle to push, and the characters should be on their guard during what turns out to be his final temptation.

If any of the characters managed to free their loved ones without damning themselves in the process, \textbf{Asmodeus} wants to convince them to serve him one last time. \textbf{Asmodeus} butters up the characters by saying that he has witnessed their power, tenacity and ingenuity and has been greatly impressed. He tells the characters that he has an offer.

Any of the characters can agree to sign a contract that stipulates that their soul goes to the Nine Hells when they die. However, they will arrive in the Nine Hells as horned devils. A successful DC 20 Charisma (Persuasion) check can make it so that they arrive in the Nine Hells as pit fiends instead. In exchange for their souls \textbf{Asmodeus} will grant them power undreamed of by most mortals. They can choose any 2 of the epic boons available in the \textit{Dungeon Master's Guide} (a successful DC 20 Charisma (Persuasion) check increases this to 3).

While he is disappointed if none of them accept his offer, he won't harm the characters. Instead, he teleports them back to \textbf{Koh Tam}'s barge. See the next section.

\section{Leaving the Nine Hells}
At this point the characters may have been teleported back to \textbf{Koh Tam}'s barge by \textbf{Asmodeus}, or they may be fleeing the Oasis. Regardless, once they reach the barge, your adventure is almost complete. It may be as simple as using the barge's plane shift capability to rise ever upwards from the lower layers to the higher. Or the characters may have unresolved grudges or debts to repay. Regardless, the characters have, in whole or part, bested the archdevil \textbf{Asmodeus} himself.

That is both an accomplishment and a signal to higher powers that the characters themselves may be a growing threat...

\appendix
\chapter{Lords of the Nine}
This appendix introduces the lords of the Nine Hells in more detail. These ancient and powerful archdevils represent some of the greatest evils in existence. Their depravity knows no bounds.

\section{Archdevil Lair Actions}
Though each archdevil is unique and harnesses a range of capabilities in which they alone are specialized, they also share some common features, especially as they pertain to defending their lairs. On initiative count 20 (losing initiative ties), an archdevil can take a lair action to trigger the unique lair action described on their statistic block. They may instead select an action from the following, Archdevil Lair Action list. Regardless of the type of lair action selected, archdevils can't take the same lair action two rounds in a row.

\subsection{Archdevil Lair Action List}

\begin{description}
\item[Attack.]
The archdevil uses one of their available melee or ranged attacks against a single foe.

\item[Cast a Spell.]
The archdevil uses their Spellcasting feature to cast an available spell. If the spell normally requires concentration, it lasts for the full duration instead.

\item[Deceitful Whispers.]
The archdevil whispers to a creature they can see within 30 feet. The target must make a DC 22 Wisdom saving throw. On a failed save, the target has the charmed condition until the start of their next turn and must use their reaction immediately to make an attack against one of the archdevil's enemies, chosen by the archdevil.

\item[Fiendish Fortitude.]
The archdevil recharges one of their expended abilities.

\item[Frenzy.]
The archdevil casts \textit{Haste} on themself. The effect ends at initiative count 20 of the next round.

\item[Summon Underlings.]
The archdevil summons allied devils. The devils summoned depends on the archdevil using this feature. The summoned devils appear in unoccupied spaces which the archdevil can see. The Summoned Underlings table shows which devils each archdevil can summon.

\item[Teleport.]
The archdevil teleports themself to an area they can see within 120 feet.

\item[Trap.]
The archdevil casts the \textit{Hold Monster} spell.

\end{description}

\begin{DndTable}[header=Summoned Underlings]{XX}

Archdevil & Underlings Summoned\\
\textbf{Asmodeus} & Any devil, including another archdevil\\
\textbf{Baalzebul} & 1d6 allied horned devils, 1d4 allied ice devils, or 1 allied \textbf{pit fiend}\\
\textbf{Bel} & 1d6 bearded devils\\
\textbf{Belial} & 1d2 bone devils\\
\textbf{Dispater} & 1 \textbf{pit fiend}\\
\textbf{Fierna} & 1d6 bearded devils\\
\textbf{Glasya} & 1d2 \textbf{erinyes}\\
\textbf{Levistus} & 1 \textbf{pit fiend}\\
\textbf{Mammon} & 1d2 bone devils\\
\textbf{Mephistopheles} & 1d4 ice devils\\
\textbf{Zariel} & 1d6 bearded devils\\
\end{DndTable}

\section{Archdevils}

\begin{itemize}
\begin{multicols}{3}
\item \textbf{Asmodeus}
\item \textbf{Baalzebul}
\item \textbf{Belial}
\item \textbf{Dispater}
\item \textbf{Fierna}
\item \textbf{Glasya}
\item \textbf{Levistus}
\item \textbf{Mammon}
\item \textbf{Mephistopheles}
\end{multicols}

\end{itemize}

\section{Additional Archdevils}
Based on events in prior adventures, the current lord of Avernus might be either \textbf{Bel} or \textbf{Zariel}.


\begin{itemize}
\begin{multicols}{3}
\item \textbf{Bel}
\item \textbf{Zariel}
\end{multicols}

\end{itemize}

\chapter{Monsters}
The characters will encounter many of the following creatures on their journey through the Nine Hells. Some are new devils that may tempt, corrupt, or simply try to murder the characters. Others are pitiful monsters, corrupted and tormented by the laws of the Nine Hells itself.

\begin{DndSidebar}[float=!b]{Making Devils More Challenging}
Devils in the Nine Hells can often be more powerful in their home turf than in the mortal realm. This can be accomplished in a few different ways:




\begin{description}
\item[Maximum Hit Points.]
Some devils might appear larger than their brethren. These devils have maximum hit points.

\item[Spellcasting.]
Some devils have learned the arcane arts.

\item[Infernal Magic Items.]
Some devils might wear or wield infernal magic items. More information is given on these items in appendix D.

\end{description}



\end{DndSidebar}


\begin{itemize}
\begin{multicols}{3}
\item \textbf{Affliction Devil (Kocrachon)}
\item \textbf{Ayperobo Swarm}
\item \textbf{Corruption Devil (Paeliryon)}
\item \textbf{Displacer Fiend}
\item \textbf{Greater Tyrant Shadow}
\item \textbf{Halog}
\item \textbf{Hellcat (Bezekira)}
\item \textbf{Lesser Tyrant Shadow}
\item \textbf{Maelephant Nomad}
\item \textbf{Miasmorne}
\item \textbf{Oneirovore}
\item \textbf{Pain Devil (Excruciarch)}
\item \textbf{Shredwing}
\item \textbf{Styx Dragon}
\item \textbf{Tyrant Shadow}
\item \textbf{Vorvolaka}
\item \textbf{War Devil}
\end{multicols}

\end{itemize}

\chapter{Non-Player Characters}
This section has the major NPCs of the story. This includes NPCs that help the heroes traverse the River Styx, as well as those that must be rescued from the Nine Hells.


\begin{itemize}
\begin{multicols}{3}
\item \textbf{Aeshma}
\item \textbf{Anagwendol}
\item \textbf{Barachiel}
\item \textbf{Jenevere}
\item \textbf{Koh Tam}
\item \textbf{Tiax}
\end{multicols}

\end{itemize}

\chapter{Infernal Magic Items}
Infernal magic items are items forged in the Nine Hells specifically for devils. These items are an easy way to make a devil more powerful, but due to their cursed nature, are often dangerous for mortals to use.


\begin{itemize}
\begin{multicols}{3}
\item \textit{Amulet of Appearance}
\item \textit{Amulet of Betrayal}
\item \textit{Amulet of Duplicity}
\item \textit{Bracers of Asmodeus}
\item \textit{Canian Fork}
\item \textit{Condensed Order}
\item \textit{Demonbone Polearm}
\item \textit{Gauntlets of Rage}
\item \textit{Infernal Amulet}
\item \textit{Infernal Plate Armor}
\item \textit{Knife of Stolen Resistance}
\item \textit{Ring of Collecting}
\item \textit{Ring of the Copycat}
\item \textit{Ring of Treachery}
\item \textit{Sage's Mirror}
\item \textit{Skull of Selfish Knowledge}
\item Soul Coin
\item \textit{Stygian Spear}
\item \textit{Sword of Retribution}
\item \textit{Vial of Greed}
\item \textit{Weapon of Agonizing Paralysis}
\end{multicols}

\end{itemize}

\section{Infernal Item Corruption}
Cursed infernal items can be used by mortals, but they always require attunement. Once attuned, the mortal risks an increasing chance of being corrupted by the item. This corruption starts as pain, infernal whisperings, and delusions, but quickly descends into physical changes, insanity, and an eventual transformation into a devil. In addition to the item's corrupting influence, these items curse mortals that attune to them.

Mortals that receive express permission from \textbf{Asmodeus}, or that make an infernal contract to acquire a magic item, don't suffer the effects of infernal corruption.

\subsection{Stage One Corruption: Beginnings}
Once a character has attuned to a cursed infernal magic item or artifact, it begins the infernal corruption process. Each time that character finishes a long rest, they must make a DC 10 Wisdom saving throw. (If the character is a tiefling they have advantage on this saving throw.) On a successful save, the character suffers no effects, but the DC increases by 1 the next time they must make this saving throw. On a failure, the character progresses to Stage Two of their corruption unless the \textit{Dispel Evil and Good} spell is cast on the character before they start their next long rest.

Breaking attunement to a cursed item prevents further Wisdom saves and resets the DC of the saving throw, if the character were to attune to the item again. However, if a character has already advanced into Stage Two of their infernal corruption, breaking attunement can no longer halt the process.

\subsection{Stage Two Corruption: Sufferings}
After failing the Wisdom saving throw in Stage One, the character progresses to Stage Two. While in Stage Two, they become delusional, seeing plots against them where there are none. Additionally, each time they rest, they experience terrifying visions and infernal whispers. Whenever they finish a long rest, they take 6 (1d12) necrotic damage, which ignores resistances and immunities and can't be healed until a Dispel Evil and Good or \textit{Remove Curse} spell is cast on them. Once the character has taken this damage from their nightmarish visions six times, they progress to Stage Three of their corruption. While in Stage Two, the infernal corruption can be removed with one of the following spells: \textit{Divine Word}, \textit{Heal}, \textit{Mass Heal}, \textit{True Polymorph}, \textit{True Resurrection}, or \textit{Wish}.

\subsection{Stage Three Corruption: Departings}
Once a character's will has been significantly weakened, they progress to Stage Three. While in Stage Three, the character begins to suffer physical transformation, and slowly embraces evil. After they finish their first long rest upon entering Stage Three they must roll 1d10 to determine how the infernal curse has begun shaping them into a devil. In addition to the specific effects, a part-devil character is rendered infertile and the character detects as a Fiend when subjected to \textit{Detect Evil and Good} spells and similar magic.

\begin{DndTable}[header=Stage Three Physical Transformations]{cX}

d10 & Transformation\\
1 & Their fingertips elongate into claws\\
2 & Non-functional leathery wings sprout from their back\\
3 & Devilish horns grow upon their head\\
4 & Each night more and more of their skin burns, leaving charred patches behind\\
5 & One eye turns milky white, the other turns yellow\\
6 & Their spine painfully elongates into a skeletal tail\\
7 & Their skin starts to calcify, turning portions into bone\\
8 & Their feet painfully twist to resemble cloven hooves\\
9 & All their hair falls out, replaced by tiny spikes\\
10 & All their teeth fall out, with new jagged teeth tearing through the gums each morning\\
\end{DndTable}

Following their devilish transformation, the character begins to experience waking whispers pushing them towards evil and they suffer terrifying visions whenever they rest, designed to completely break their spirit and push them further to evil. Each time they finish a short or long rest, the character must make a DC 10 Wisdom saving throw. If the character performed at least one evil act, such as an unprovoked killing of a creature, a manipulative deal, or making a decision that increased the suffering of others, they make the saving throw with disadvantage. When they fail the saving throw, they progress to the final stage of their corruption, Stage Four.

While in Stage Three, the infernal corruption can be ended with one of the following spells: \textit{True Polymorph}, \textit{True Resurrection} or \textit{Wish}.

\subsection{Stage Four Corruption: Finalities}
After the complete devolution of a character's morals and will, they progress to Stage Four. When they finish their first long rest after reaching Stage Four, the character's alignment shifts to lawful evil. They're now bound by the devil's code, requiring them to honor any pact made and acquire souls in service of \textbf{Asmodeus}. Lastly, their physical form changes, morphing to resemble a devil chosen by the DM.

Once the character reaches Stage Four, the only two cures are the \textit{Wish} spell, which counts as beyond the scope of the spell, or by signing an infernal contract with \textbf{Asmodeus} to reclaim their soul.

\section{Infernal Artifacts}
These items are weapons or tools used by the various archdevils of the Nine Hells and are extremely powerful. Obtaining them should be exceedingly difficult, often requiring defeating or outwitting a ruling archdevil. Archdevils use their respective artifacts to maximum effect against opponents during combat.


\begin{itemize}
\begin{multicols}{3}
\item \textit{Accounting and Valuation of All Things}
\item \textit{Amulet of the Inferno}
\item \textit{Ranseur of Torture}
\item \textit{Ruby Rod of Asmodeus}
\item \textit{Scourge of Shadow}
\item \textit{True-Ice Shards}
\item \textit{Wrought-Iron Tower}
\end{multicols}

\end{itemize}

\chapter{Corruption}
The information in this appendix helps you keep track of each character's overall corruption. It shows you the possible endings available to the characters once they reach the final stages of the adventure.

\section{Corruption Tracker}
Use the "Corruption Tracker" to keep track of the number of times the characters have encountered one of \textbf{Asmodeus}'s temptations. Every time they give in to temptation, they're awarded 2 corruption points.

\section{Corruption Score}
Once the characters reach Nessus tally their scores on the "Corruption Tracker" and use the Corruption Score table to determine if they're indeed corrupted.

\begin{DndTable}[header=Corruption Score]{ccc}

Number of Temptations & Max. Number of Corruption Points & Points Needed for Corruption\\
2 & 4 & 3\\
3 & 6 & 4\\
4 & 8 & 5\\
5 & 10 & 6\\
6 & 12 & 7\\
7 & 14 & 8\\
8 & 16 & 9\\
\end{DndTable}

\section{Final Outcomes}
Use the Final Outcomes table to determine the path available to each character. The outcomes inform you on how to run the different encounters in Nessus.

\begin{DndTable}[header=Final Outcomes]{XX}

Oasis of the Lethe & Malsheem: Triad and Asmodeus\\
\textit{Corrupted: Soul of the Other} & \\

\begin{itemize}\item The character can free their loved one from the contract.\item Their own soul is destined to go to the Nine Hells.\item The character can make a deal with the Triad and \textbf{Asmodeus} to ensure a better position upon death.\end{itemize}
 & 
\begin{itemize}\item The character can free the soul of their loved one.\item Their own soul is destined to go to the Nine Hells.\item The character can make a deal with the Triad and \textbf{Asmodeus} to ensure a better position upon death.\end{itemize}
\\
\textit{Corrupted: Soul of the Damned} & \\

\begin{itemize}\item The character can undo their contract but is still damned to the Nine Hells.\item When the temporary magic their patron used to restore them to life eventually fades, they return to the Nine Hells.\item The character can negotiate a new deal with the Triad and \textbf{Asmodeus} to ensure a better position upon arriving in the Nine Hells.\end{itemize}
 & 
\begin{itemize}\item The character can undo their contract but is still damned to the Nine Hells.\item When the temporary magic their patron used to restore them to life eventually fades, they return to the Nine Hells.\item The character can negotiate a new deal with the Triad and \textbf{Asmodeus} to ensure a better position upon arriving in the Nine Hells.\end{itemize}
\\
\textit{Uncorrupted: Soul of the Other} & \\

\begin{itemize}\item The character can free their loved one from the contract.\item Once they trigger the loophole, \textbf{Asmodeus} is enraged.\item The character doesn't need to deal with \textbf{Asmodeus}.\end{itemize}
 & 
\begin{itemize}\item The character can free their loved one from the contract.\item \textbf{Asmodeus} offers them a deal. They can refuse and they and their loved one go free.\end{itemize}
\\
\textit{Uncorrupted: Soul of the Damned} & \\

\begin{itemize}\item The character can undo their contract.\item When the ritual their patron used to return them to life eventually fades, they die, but their soul isn't doomed to the Nine Hells.\end{itemize}
 & 
\begin{itemize}\item The character can undo their contract.\item \textbf{Asmodeus} offers them a new deal. They can refuse and go free.\item When the temporary magic their patron used to bring them to life wears off, they die.\end{itemize}
\\
\end{DndTable}

\chapter{Player Handouts}
\textbf{Koh Tam}, through his many travels along the River Styx, has tried to learn as much as possible about the Nine Hells. In preparation for this latest voyage, he has obtained maps for some of the more promising locations he believes the characters should explore. Unfortunately, the Nine Hells are ever-changing and unpredictable, and he can't vouch for the accuracy of any of the maps.

\section{Infernal Maps}
Worse still, some maps are marked with strange runes, unfamiliar even to \textbf{Koh Tam}. He knows that the important locations on each map have been marked but he doesn't know what might be encountered at each mark. Such is the nature of infernal maps...

\section{Normal Maps}
\textbf{Koh Tam} considers a couple of the maps that have fallen into his possession more worthy of trust. These he has annotated himself with quick descriptions of important locations.

\chapter{Credits}

\begin{description}
\begin{multicols}{2}

\begin{description}
\item[Project Lead.]
James Ohlen

\item[Producer.]
Brent Knowles

\item[Lead Writer.]
Adrian Tchaikovsky

\end{description}


\begin{description}
\item[Writers.]
Marieke Cross, Andreana Lozano, Drew Karpyshyn, Brent Knowles, James Ohlen

\item[Rules Developers.]
Brent Knowles, Zack Webb

\item[Editor.]
Marieke Cross

\end{description}


\begin{description}
\item[Graphic Designer.]
Michal E. Cross

\item[Cover Illustrator.]
Aaron Sims Company

\item[Cartographer.]
John Stevenson

\item[Nine Hells Art.]
Sergei Sarichev

\item[Encounter Art.]
Sergey Musin

\item[Environment Art.]
Julian Calle

\item[Monster Art.]
Sebastian Kowoll

\item[Archdevil Art.]
Aaron Sims

\item[Character Art.]
Paul Adam, Lius Lasahido

\end{description}


\begin{description}
\item[Cultural \& Sensitivity Consultant.]
Stacey Parshall Jensen

\item[Special Thanks.]
Christopher Perkins, Jeremy Crawford

\end{description}

\end{multicols}

\end{description}

\chapter{Magic Items}
\DndItemHeader{+1 Wand of the War Mage}{u, n, c, o, m, m, o, n None}
While you are holding this wand, you gain a +1 bonus to spell attack rolls. In addition, you ignore half cover when making a spell attack.

\DndItemHeader{Accounting and Valuation of All Things}{a, r, t, i, f, a, c, t None}
The Accounting and Valuation of All Things is bound in gold leaf, with pages of silver and text of blood.

This is a magic tome that functions as a spellbook and arcane focus. You have a +3 bonus to your spell attack and spell save DC.

\subparagraph{Random Properties}
The Accounting and Valuation of All Things has the following random properties:


\begin{itemize}
\item 1 Artifact Properties; Major Beneficial Properties property
\item 1 Artifact Properties; Major Detrimental Properties property
\end{itemize}

\subparagraph{The Soul Trade}
One section of the book is dedicated to negotiating and valuing mortal souls. When the book is opened within 10 feet of a creature with a soul, that creature must make a DC 22 Wisdom saving throw. On a failed save, the tome analyzes their greatest desires and calculates exactly what they would need to give up their soul. That information is then displayed within the tome, such that the user can offer a trade to acquire a soul. Creatures that succeed on the saving throw simply show up as a question mark in the book's pages and are immune to further valuations.

\subparagraph{Mammon's Tax}
When the Accounting and Valuation of All Things is used as a spellbook, it maintains a unique ability to forcibly acquire souls. Any creature killed by a spell transcribed within the book has their soul ripped from their body and sent to Minauros. This process renders the subject immune to all forms of resurrection aside from \textit{True Resurrection}. In payment for their soul, Mammon drops gold equal to (20 × CR) or (2 × Level) on their corpse.

\subparagraph{Negotiation Tactics}
The tome is capable of providing some assistance to all negotiations, not just those made for souls. You gain a passive +5 to Charisma (Persuasion) checks made to barter. Additionally, as an action, you can surrender your consciousness to the tome, allowing it to take over. For 1 minute, the Accounting and Valuation of All Things speaks for you, using a +12 Charisma (Persuasion) skill and gaining advantage on your rolls. Using the book in this manner has a corrupting influence, and each use compels you to become increasingly greedy, vain, and evil—much like its original owner, Mammon. When the duration ends, if your alignment is non-evil, you suffer 6d6 necrotic damage.

\subparagraph{Transcribed Spells}
The spellbook section of the Accounting and Valuation of All Things comes with some spells already recorded. For each spell level, roll (1d4 - 1). The DM will determine the necessary spells of each level that already exist within the tome. If you copy spells from the tome to your own spellbook, the cost in gold is doubled. On the opposite end, copying spells from your spellbook to the tome only costs half the gold it normally would.

\subparagraph{Destroying the Tome}
Destroying the Accounting and Valuation of All Things can be done in two ways. The first option is a simple bribe to Mammon to have him destroy the tome. An offering of 99,999 gp must be made, at midnight, by opening the book and placing the gold inside. If Mammon accepts the offering, the gold melts around the book, then both vanish.

The second route, in case Mammon is unavailable or unwilling, is to incinerate the book. To accomplish this task, a pot of platinum (valued at no less than 99,999 gp) must be melted, enough to fully submerge the book. Once the platinum is prepared, dropping the tome into the molten solution melts the tome. The process completely destroys the tome, as well as the platinum.

\DndItemHeader{Amulet of Appearance}{r, a, r, e None}
Your armor, weapons, and other equipment always shine as if just polished. Even if you're wounded, your wounds do not appear to others. You're immune to the frightened and poisoned conditions, as these would otherwise ruin your elegant appearance. However, whenever you would otherwise have been affected by one of these conditions, a random non-evil Humanoid on the Material Plane whom you have previously met gets the condition instead. You know this.

\DndItemHeader{Amulet of Betrayal}{v, e, r, y,  , r, a, r, e None}
You can use a bonus action to move one of the following conditions from yourself to an ally within 60 feet of you: blinded, deafened, frightened, poisoned, stunned, exhaustion. When transferring exhaustion, move all your exhaustion levels. If the chosen ally is immune to a transferred condition, the transfer fails.

\DndItemHeader{Amulet of Duplicity}{v, e, r, y,  , r, a, r, e None}
When you die, you're transported to an extradimensional space where you're stabilized at 0 hit points and kept in that state. The amulet creates a perfect copy of your corpse and places it where you were just before you died, but without any of your worn or carried items. After 24 hours, you regain 1 hit point and are returned to the location of your near-death. Everything about you is the same, except that your facial features are entirely different from before. Only a \textit{Wish} spell can restore your true identity.

Until then, you have disadvantage on any Charisma (Persuasion) checks to attempt to reveal who you really are.

\subparagraph{Curse}
An \textit{Identify} spell or similar reveals only that the amulet can prevent death. Once you attune to the amulet, you can't describe its ability to any other creature, and if unattuned, you immediately forget what the amulet does. A \textit{Remove Curse} spell reveals the details of the curse but does not remove it from the amulet.

\DndItemHeader{Amulet of Health}{r, a, r, e None}
Your Constitution score is 19 while you wear this amulet. It has no effect on you if your Constitution score is already 19 or higher without it.

\DndItemHeader{Amulet of Proof against Detection and Location}{u, n, c, o, m, m, o, n None}
While wearing this amulet, you are hidden from divination magic. You can't be targeted by such magic or perceived through magical scrying sensors.

\DndItemHeader{Amulet of the Inferno}{a, r, t, i, f, a, c, t None}
When Fierna was born, petitioners and devils from all over came to celebrate her birth. Many brought gifts and treasures, including a powerful infernal warlock. Recognizing the strength and inner fire of Fierna, the warlock gifted her with the Amulet of the Inferno, and it has been around her neck ever since.

The Amulet of the Inferno is a magic necklace that grants a +3 bonus to spell save DC and spell attack bonus.

\subparagraph{Random Properties}
The Amulet of the Inferno has the following random properties:


\begin{itemize}
\item 3 Artifact Properties; Major Beneficial Properties properties
\item 2 Artifact Properties; Major Detrimental Properties properties
\end{itemize}

\subparagraph{Infernal Fire}
While worn, the amulet augments any magical flames created by you, transforming normal fire into Hellfire. Hellfire ignores resistances and immunities to fire damage and deals double damage to creatures vulnerable to either fire or necrotic damage. It is also capable of melting stone, as well as igniting flammable objects.

\subparagraph{Raging Inferno}
While attuned to the Amulet of the Inferno, you may use a bonus action to cast \textit{Fire Bolt} or \textit{Produce Flame}. The cantrips scale based on your total character level and use Intelligence as a spellcasting modifier if you're normally incapable of casting spells.

\subparagraph{Everburning}
The Amulet of the Inferno grants resistance to cold damage and the ability to breathe underwater. Any magical flames you conjure are similarly augmented: they can be used underwater, can't be extinguished through any means, and can exude any temperature. While the casting of a fire spell can still be interrupted with \textit{Counterspell}, once the spell has been cast, any lingering flames can't be removed.

\subparagraph{Destroying the Amulet}
The only creature capable of destroying the amulet is the warlock that created it. Nobody living knows the warlock's identity, making destroying the amulet almost impossible.

\DndItemHeader{Amulet of the Planes}{v, e, r, y,  , r, a, r, e None}
While wearing this amulet, you can use an action to name a location that you are familiar with on another plane of existence. Then make a DC 15 Intelligence check. On a successful check, you cast the \textit{plane shift} spell. On a failure, you and each creature and object within 15 feet of you travel to a random destination. Roll a d100. On a 1-60, you travel to a random location on the plane you named. On a 61-100, you travel to a randomly determined plane of existence.

\DndItemHeader{Apparatus of Kwalish}{l, e, g, e, n, d, a, r, y}
This item first appears to be a Large sealed iron barrel weighing 500 pounds. The barrel has a hidden catch, which can be found with a successful DC 20 Intelligence (Investigation) check. Releasing the catch unlocks a hatch at one end of the barrel, allowing two Medium or smaller creatures to crawl inside. Ten levers are set in a row at the far end, each in a neutral position, able to move either up or down. When certain levers are used, the apparatus transforms to resemble a giant lobster.

The apparatus of Kwalish is a Large object with the following statistics:


\begin{description}
\item[Armor Class:]
20
\item[Hit Points:]
200
\item[Speed:]
30 ft., swim 30 ft. (or 0 ft. for both if the legs and tail aren't extended)
\item[Damage Immunities:]
poison, psychic
\end{description}

To be used as a vehicle, the apparatus requires one pilot. While the apparatus's hatch is closed, the compartment is airtight and watertight. The compartment holds enough air for 10 hours of breathing, divided by the number of breathing creatures inside.

The apparatus floats on water. It can also go underwater to a depth of 900 feet. Below that, the vehicle takes 2d6 bludgeoning damage per minute from pressure.

A creature in the compartment can use an action to move as many as two of the apparatus's levers up or down. After each use, a lever goes back to its neutral position. Each lever, from left to right, functions as shown in the Apparatus of Kwalish Levers table.

\begin{DndTable}[header=Apparatus of Kwalish Levers:]{cXX}

Lever & Up & Down\\
1 & Legs and tail extend, allowing the apparatus to walk and swim. & Legs and tail retract, reducing the apparatus's speed to 0 and making it unable to benefit from bonuses to speed.\\
2 & Forward window shutter opens. & Forward window shutter closes.\\
3 & Side window shutters open (two per side). & Side window shutters close (two per side).\\
4 & Two claws extend from the front sides of the apparatus. & The claws retract.\\
5 & Each extended claw makes the following melee weapon attack: 8 to hit, reach 5 ft., one target. \textsl{Hit:}  7 (2d6) bludgeoning damage. & Each extended claw makes the following melee weapon attack: 8 to hit, reach 5 ft., one target. \textsl{Hit:}  The target is grappled (escape DC 15).\\
6 & The apparatus walks or swims forward. & The apparatus walks or swims backward.\\
7 & The apparatus turns 90 degrees left. & The apparatus turns 90 degrees right.\\
8 & Eyelike fixtures emit bright light in a 30-foot radius and dim light for an additional 30 feet. & The light turns off.\\
9 & The apparatus sinks as much as 20 feet in liquid. & The apparatus rises up to 20 feet in liquid.\\
10 & The rear hatch unseals and opens. & The rear hatch closes and seals.\\
\end{DndTable}

\DndItemHeader{Armor of Invulnerability}{l, e, g, e, n, d, a, r, y None}
You have resistance to nonmagical damage while you wear this armor. Additionally, you can use an action to make yourself immune to nonmagical damage for 10 minutes or until you are no longer wearing the armor. Once this special action is used, it can't be used again until the next dawn.

\DndItemHeader{Bead of Force}{r, a, r, e}
This small black sphere measures ¾ of an inch in diameter and weighs an ounce. Typically, 1d4 + 4 \textit{beads of force} are found together.

You can use an action to throw the bead up to 60 feet. The bead explodes on impact and is destroyed. Each creature within a 10-foot radius of where the bead landed must succeed on a DC 15 Dexterity saving throw or take 5d4 force damage. A sphere of transparent force then encloses the area for 1 minute. Any creature that failed the save and is completely within the area is trapped inside this sphere. Creatures that succeeded on the save, or are partially within the area, are pushed away from the center of the sphere until they are no longer inside it. Only breathable air can pass through the sphere's wall. No attack or other effect can.

An enclosed creature can use its action to push against the sphere's wall, moving the sphere up to half the creature's walking speed. The sphere can be picked up, and its magic causes it to weigh only 1 pound, regardless of the weight of creatures inside.

\DndItemHeader{Belt of Storm Giant Strength}{l, e, g, e, n, d, a, r, y None}
While wearing this belt, your Strength score changes to 29. The item has no effect on you if your Strength without the belt is equal to or greater than the belt's score.

\DndItemHeader{Bracers of Asmodeus}{r, a, r, e None}
You have a +2 bonus to AC while wearing these bracers, if you do not wear armor or use a shield at the same time.

\subparagraph{Curse}
While attuned to the bracers, you become obsessed with plotting, scheming, and manipulation. You always barter for better deals, often using secrets or leveraging other offers in the process. If you ever decline an opportunity to better yourself financially at another's expense, you immediately take 3d10 necrotic damage. Only the \textit{Remove Curse} spell allows you to end attunement to this item.

\subparagraph{Corrupting}
This item corrupts. See the "Infernal Item Corruption" section.

\DndItemHeader{Bracers of Defense}{r, a, r, e None}
While wearing these bracers, you gain a +2 bonus to AC if you are wearing no armor and using no shield.

\DndItemHeader{Brazier of Commanding Fire Elementals}{r, a, r, e}
While a fire burns in this brass brazier, you can use an action to speak the brazier's command word and summon a \textbf{fire elemental}, as if you had cast the \textit{conjure elemental} spell. The brazier can't be used this way again until the next dawn.

The brazier weighs 5 pounds.

\DndItemHeader{Broom of Flying}{u, n, c, o, m, m, o, n}
This wooden broom, which weighs 3 pounds, functions like a mundane broom until you stand astride it and speak its command word. It then hovers beneath you and can be ridden in the air. It has a flying speed of 50 feet. It can carry up to 400 pounds, but its flying speed becomes 30 feet while carrying over 200 pounds. The broom stops hovering when you land.

You can send the broom to travel alone to a destination within 1 mile of you if you speak the command word, name the location, and are familiar with that place. The broom comes back to you when you speak another command word, provided that the broom is still within 1 mile of you.

\DndItemHeader{Canian Fork}{r, a, r, e None}
You have a +3 bonus to attack and damage rolls made with this magic weapon. In addition, you can make one additional attack with it as a bonus action on each of your turns.

\subparagraph{Curse}
You're unwilling to part with this weapon while attuned to it. You're also vulnerable to radiant damage and each time you receive magical healing, you must make a DC 15 Constitution saving throw.

On a failed save, the healing has no effect. Only the \textit{Remove Curse} spell allows you to end attunement to this item.

\subparagraph{Corrupting}
This item corrupts. See the "Infernal Item Corruption" section.

\DndItemHeader{Cape of the Mountebank}{r, a, r, e}
This cape smells faintly of brimstone. While wearing it, you can use it to cast the \textit{dimension door} spell as an action. This property of the cape can't be used again until the next dawn.

When you disappear, you leave behind a cloud of smoke, and you appear in a similar cloud of smoke at your destination. The smoke lightly obscures the space you left and the space you appear in, and it dissipates at the end of your next turn. A light or stronger wind disperses the smoke.

\DndItemHeader{Cloak of Arachnida}{v, e, r, y,  , r, a, r, e None}
This fine garment is made of black silk interwoven with faint silvery threads. While wearing it, you gain the following benefits:


\begin{itemize}
\item You have resistance to poison damage.
\item You have a climbing speed equal to your walking speed.
\item You can move up, down, and across vertical surfaces and upside down along ceilings, while leaving your hands free.
\item You can't be caught in webs of any sort and can move through webs as if they were difficult terrain.
\item You can use an action to cast the \textit{web} spell (save DC 13). The web created by the spell fills twice its normal area. Once used, this property of the cloak can't be used again until the next dawn.
\end{itemize}

\DndItemHeader{Cloak of Invisibility}{l, e, g, e, n, d, a, r, y None}
While wearing this cloak, you can pull its hood over your head to cause yourself to become invisible. While you are invisible, anything you are carrying or wearing is invisible with you. You become visible when you cease wearing the hood. Pulling the hood up or down requires an action.

Deduct the time you are invisible, in increments of 1 minute, from the cloak's maximum duration of 2 hours. After 2 hours of use, the cloak ceases to function. For every uninterrupted period of 12 hours the cloak goes unused, it regains 1 hour of duration.

\DndItemHeader{Cloak of Protection}{u, n, c, o, m, m, o, n None}
You gain a +1 bonus to AC and saving throws while you wear this cloak.

\DndItemHeader{Condensed Order}{u, n, c, o, m, m, o, n}
Condensed Order is a silvery powder that can be extracted from those of a lawful persuasion. Devils bound for the warfronts of Avernus take flasks and snuff boxes of the stuff to fortify themselves against exposure to the raw chaos of demons. Taking the substance requires an action and makes you immune to the flesh warping feature of demonic ichor. It also gives you advantage on saving throws against any effect from a demonic source. These benefits last for 8 hours.

\DndItemHeader{Crystal Ball}{v, e, r, y,  , r, a, r, e None}
This crystal ball is about 6 inches in diameter. While touching it, you can cast the \textit{scrying} spell (save DC 17) with it.

\DndItemHeader{Dagger of Venom}{r, a, r, e}
You gain a +1 bonus to attack and damage rolls made with this magic weapon.

You can use an action to cause thick, black poison to coat the blade. The poison remains for 1 minute or until an attack using this weapon hits a creature. That creature must succeed on a DC 15 Constitution saving throw or take 2d10 poison damage and become poisoned for 1 minute. The dagger can't be used this way again until the next dawn.

\DndItemHeader{Deck of Illusions}{u, n, c, o, m, m, o, n}
This box contains a set of parchment cards. A full deck has 34 cards. A deck found as treasure is usually missing 1d20 - 1 cards.

The magic of the deck functions only if cards are drawn at random (you can use an altered deck of playing cards to simulate the deck). You can use an action to draw a card at random from the deck and throw it to the ground at a point within 30 feet of you.

An illusion of one or more creatures forms over the thrown card and remains until dispelled. An illusory creature appears real, of the appropriate size, and behaves as if it were a real creature, except that it can do no harm. While you are within 120 feet of the illusory creature and can see it, you can use an action to move it magically anywhere within 30 feet of its card. Any physical interaction with the illusory creature reveals it to be an illusion, because objects pass through it. Someone who uses an action to visually inspect the creature identifies it as illusory with a successful DC 15 Intelligence (Investigation) check. The creature then appears translucent.

The illusion lasts until its card is moved or the illusion is dispelled. When the illusion ends, the image on its card disappears, and that card can't be used again.


\begin{DndTable}{XX}

1d34 & Illusion\\
1 & Red dragon\\
2 & \textbf{Knight} and four guards\\
3 & Succubus/Incubus\\
4 & \textbf{Druid}\\
5 & \textbf{Cloud giant}\\
6 & \textbf{Ettin}\\
7 & \textbf{Bugbear}\\
8 & \textbf{Goblin}\\
9 & \textbf{Beholder}\\
10 & \textbf{Archmage} and \textbf{mage} apprentice\\
11 & \textbf{Night hag}\\
12 & \textbf{Assassin}\\
13 & \textbf{Fire giant}\\
14 & Ogre mage\\
15 & \textbf{Gnoll}\\
16 & \textbf{Kobold}\\
17 & \textbf{Lich}\\
18 & \textbf{Priest} and two acolytes\\
19 & \textbf{Medusa}\\
20 & \textbf{Veteran}\\
21 & \textbf{Frost giant}\\
22 & \textbf{Troll}\\
23 & \textbf{Hobgoblin}\\
24 & \textbf{Goblin}\\
25 & \textbf{Iron golem}\\
26 & \textbf{Bandit captain} and three bandits\\
27 & \textbf{Erinyes}\\
28 & \textbf{Berserker}\\
29 & \textbf{Hill giant}\\
30 & \textbf{Ogre}\\
31 & \textbf{Orc}\\
32 & \textbf{Kobold}\\
33—34 & You (the deck's owner)\\
\end{DndTable}

\DndItemHeader{Demon Armor}{v, e, r, y,  , r, a, r, e None}
While wearing this armor, you gain a +1 bonus to AC, and you can understand and speak Abyssal. In addition, the armor's clawed gauntlets turn unarmed strikes with your hands into magic weapons that deal slashing damage, with a +1 bonus to attack and damage rolls and a damage die of 1d8.

\subparagraph{Curse}
Once you don this cursed armor, you can't doff it unless you are targeted by the \textit{remove curse} spell or similar magic. While wearing the armor, you have disadvantage on attack rolls against demons and on saving throws against their spells and special abilities.

\DndItemHeader{Dimensional Shackles}{r, a, r, e}
You can use an action to place these shackles on an incapacitated creature. The shackles adjust to fit a creature of Small to Large size. In addition to serving as mundane manacles, the shackles prevent a creature bound by them from using any method of extradimensional movement, including teleportation or travel to a different plane of existence. They don't prevent the creature from passing-through an interdimensional portal.

You and any creature you designate when you use the shackles can use an action to remove them. Once every 30 days, the bound creature can make a DC 30 Strength (Athletics) check. On a success, the creature breaks free and destroys the shackles.

\DndItemHeader{Dust of Sneezing and Choking}{u, n, c, o, m, m, o, n}
Found in a small container, this powder resembles very fine sand. It appears to be \textit{dust of disappearance}, and an \textit{identify} spell reveals it to be such. There is enough of it for one use.

When you use an action to throw a handful of the dust into the air, you and each creature that needs to breathe within 30 feet of you must succeed on a DC 15 Constitution saving throw or become unable to breathe while sneezing uncontrollably. A creature affected in this way is incapacitated and suffocating. As long as it is conscious, a creature can repeat the saving throw at the end of each of its turns, ending the effect on it on a success. The \textit{lesser restoration} spell can also end the effect on a creature.

\DndItemHeader{Dwarven Plate}{v, e, r, y,  , r, a, r, e}
While wearing this armor, you gain a +2 bonus to AC. In addition, if an effect moves you against your will along the ground, you can use your reaction to reduce the distance you are moved by up to 10 feet.

\DndItemHeader{Efreeti Bottle}{v, e, r, y,  , r, a, r, e}
This painted brass bottle weighs 1 pound. When you use an action to remove the stopper, a cloud of thick smoke flows out of the bottle. At the end of your turn, the smoke disappears with a flash of harmless fire, and an \textbf{efreeti} appears in an unoccupied space within 30 feet of you.

The first time the bottle is opened, the DM rolls to determine what happens.


\begin{DndTable}{cX}

d100 & Effect\\
01-10 & The \textbf{efreeti} attacks you. After fighting for 5 rounds, the \textbf{efreeti} disappears, and the bottle loses its magic.\\
11-90 & The \textbf{efreeti} serves you for 1 hour, doing as you command. Then the \textbf{efreeti} returns to the bottle, and a new stopper contains it. The stopper can't be removed for 24 hours. The next two times the bottle is opened, the same effect occurs. If the bottle is opened a fourth time, the \textbf{efreeti} escapes and disappears, and the bottle loses its magic.\\
91-00 & The \textbf{efreeti} can cast the \textit{wish} spell three times for you. It disappears when it grants the final wish or after 1 hour, and the bottle loses its magic.\\
\end{DndTable}

\DndItemHeader{Eyes of Charming}{u, n, c, o, m, m, o, n None}
These crystal lenses fit over the eyes. They have 3 charges. While wearing them, you can expend 1 charge as an action to cast the \textit{charm person} spell (save DC 13) on a humanoid within 30 feet of you, provided that you and the target can see each other. The lenses regain all expended charges daily at dawn.

\DndItemHeader{Eyes of Minute Seeing}{u, n, c, o, m, m, o, n}
These crystal lenses fit over the eyes. While wearing them, you can see much better than normal out to a range of 1 foot. You have advantage on Intelligence (Investigation) checks that rely on sight while searching an area or studying an object within that range.

\DndItemHeader{Eyes of the Eagle}{u, n, c, o, m, m, o, n None}
These crystal lenses fit over the eyes. While wearing them, you have advantage on Wisdom (Perception) checks that rely on sight. In conditions of clear visibility, you can make out details of even extremely distant creatures and objects as small as 2 feet across.

\DndItemHeader{Figurine of Wondrous Power, Golden Lions}{r, a, r, e}
A figurine of wondrous power is a statuette of a beast small enough to fit in a pocket. If you use an action to speak the command word and throw the figurine to a point on the ground within 60 feet of you, the figurine becomes a living creature. If the space where the creature would appear is occupied by other creatures or objects, or if there isn't enough space for the creature, the figurine doesn't become a creature.

The creature is friendly to you and your companions. It understands your languages and obeys your spoken commands. If you issue no commands, the creature defends itself but takes no other actions.

The creature exists for a duration specific to each figurine. At the end of the duration, the creature reverts to its figurine form. It reverts to a figurine early if it drops to 0 hit points or if you use an action to speak the command word again while touching it. When the creature becomes a figurine again, its property can't be used again until a certain amount of time has passed, as specified in the figurine's description.

\subparagraph{Golden Lions}
These gold statuettes of lions are always created in pairs. You can use one figurine or both simultaneously. Each can become a \textbf{lion} for up to 1 hour. Once a \textbf{lion} has been used, it can't be used again until 7 days have passed.

\DndItemHeader{Figurine of Wondrous Power, Marble Elephant}{r, a, r, e}
A figurine of wondrous power is a statuette of a beast small enough to fit in a pocket. If you use an action to speak the command word and throw the figurine to a point on the ground within 60 feet of you, the figurine becomes a living creature. If the space where the creature would appear is occupied by other creatures or objects, or if there isn't enough space for the creature, the figurine doesn't become a creature.

The creature is friendly to you and your companions. It understands your languages and obeys your spoken commands. If you issue no commands, the creature defends itself but takes no other actions.

The creature exists for a duration specific to each figurine. At the end of the duration, the creature reverts to its figurine form. It reverts to a figurine early if it drops to 0 hit points or if you use an action to speak the command word again while touching it. When the creature becomes a figurine again, its property can't be used again until a certain amount of time has passed, as specified in the figurine's description.

\subparagraph{Marble Elephant}
This marble statuette is about 4 inches high and long. It can become an \textbf{elephant} for up to 24 hours. Once it has been used, it can't be used again until 7 days have passed.

\DndItemHeader{Figurine of Wondrous Power, Onyx Dog}{r, a, r, e}
A figurine of wondrous power is a statuette of a beast small enough to fit in a pocket. If you use an action to speak the command word and throw the figurine to a point on the ground within 60 feet of you, the figurine becomes a living creature. If the space where the creature would appear is occupied by other creatures or objects, or if there isn't enough space for the creature, the figurine doesn't become a creature.

The creature is friendly to you and your companions. It understands your languages and obeys your spoken commands. If you issue no commands, the creature defends itself but takes no other actions.

The creature exists for a duration specific to each figurine. At the end of the duration, the creature reverts to its figurine form. It reverts to a figurine early if it drops to 0 hit points or if you use an action to speak the command word again while touching it. When the creature becomes a figurine again, its property can't be used again until a certain amount of time has passed, as specified in the figurine's description.

\subparagraph{Onyx Dog}
This onyx statuette of a dog can become a \textbf{mastiff} for up to 6 hours. The \textbf{mastiff} has an Intelligence of 8 and can speak Common. It also has darkvision out to a range of 60 feet and can see invisible creatures and objects within that range. Once it has been used, it can't be used again until 7 days have passed.

\DndItemHeader{Gauntlets of Ogre Power}{u, n, c, o, m, m, o, n None}
Your Strength score is 19 while you wear these gauntlets. They have no effect on you if your Strength is already 19 or higher without them.

\DndItemHeader{Gauntlets of Rage}{r, a, r, e None}
You gain the power of never-ending fury. After you make a melee weapon attack, you automatically enter a special fury for 1 minute. While furious, you can't cast spells, can't verbally communicate, and are immune to the charmed and frightened conditions. Each time you deal damage, you regain 2d8 hit points and may immediately remove a condition you currently suffer from or end this fury. Additionally, while in this fury, you may spend a spell slot before you make a melee attack. Doing so causes the attack to deal an extra 3 (1d6) necrotic damage per level of spell slot expended, if the attack hits. Once used you may not use this fury again until you finish a short or long rest.

\DndItemHeader{Gem of Seeing}{r, a, r, e None}
This gem has 3 charges. As an action, you can speak the gem's command word and expend 1 charge. For the next 10 minutes, you have truesight out to 120 feet when you peer through the gem.

The gem regains 1d3 expended charges daily at dawn.

\DndItemHeader{Goggles of Night}{u, n, c, o, m, m, o, n}
While wearing these dark lenses, you have darkvision out to a range of 60 feet. If you already have darkvision, wearing the goggles increases its range by 60 feet.

\DndItemHeader{Hat of Disguise}{u, n, c, o, m, m, o, n None}
While wearing this hat, you can use an action to cast the \textit{disguise self} spell from it at will. The spell ends if the hat is removed.

\DndItemHeader{Headband of Intellect}{u, n, c, o, m, m, o, n None}
Your Intelligence score is 19 while you wear this headband. It has no effect on you if your Intelligence is already 19 or higher without it.

\DndItemHeader{Helm of Brilliance}{v, e, r, y,  , r, a, r, e None}
This dazzling helm is set with 1d10 diamonds, 2d10 rubies, 3d10 fire opals, and 4d10 opals. Any gem pried from the helm crumbles to dust. When all the gems are removed or destroyed, the helm loses its magic.

You gain the following benefits while wearing it:


\begin{itemize}
\item You can use an action to cast one of the following spells (save DC 18), using one of the helm's gems of the specified type as a component: \textit{daylight} (opal), \textit{fireball} (fire opal), \textit{prismatic spray} (diamond), or \textit{wall of fire} (ruby). The gem is destroyed when the spell is cast and disappears from the helm.
\item As long as it has at least one diamond, the helm emits dim light in a 30-foot radius when at least one undead is within that area. Any undead that starts its turn in that area takes 1d6 radiant damage.
\item As long as the helm has at least one ruby, you have resistance to fire damage.
\item As long as the helm has at least one fire opal, you can use an action and speak a command word to cause one weapon you are holding to burst into flames. The flames emit bright light in a 10-foot radius and dim light for an additional 10 feet. The flames are harmless to you and the weapon. When you hit with an attack using the blazing weapon, the target takes an extra 1d6 fire damage. The flames last until you use a bonus action to speak the command word again or until you drop or stow the weapon.
\end{itemize}

Roll a d20 if you are wearing the helm and take fire damage as a result of failing a saving throw against a spell. On a roll of 1, the helm emits beams of light from its remaining gems. Each creature within 60 feet of the helm other than you must succeed on a DC 17 Dexterity saving throw or be struck by a beam, taking radiant damage equal to the number of gems in the helm. The helm and its gems are then destroyed.

\DndItemHeader{Helm of Comprehending Languages}{u, n, c, o, m, m, o, n}
While wearing this helm, you can use an action to cast the \textit{comprehend languages} spell from it at will.

\DndItemHeader{Helm of Telepathy}{u, n, c, o, m, m, o, n None}
While wearing this helm, you can use an action to cast the \textit{detect thoughts} spell (save DC 13) from it. As long as you maintain concentration on the spell, you can use a bonus action to send a telepathic message to a creature you are focused on. It can reply—using a bonus action to do so—while your focus on it continues.

While focusing on a creature with \textit{detect thoughts}, you can use an action to cast the \textit{suggestion} spell (save DC 13) from the helm on that creature. Once used, the suggestion property can't be used again until the next dawn.

\DndItemHeader{Horn of Valhalla, Iron}{l, e, g, e, n, d, a, r, y}
You can use an action to blow this horn. In response, warrior spirits from the plane of Ysgard appear within 60 feet of you. These spirits use the \textbf{berserker} statistics. They return to Ysgard after 1 hour or when they drop to 0 hit points. Once you use the horn, it can't be used again until 7 days have passed.

The iron horn summons 5d4 + 5 berserkers. To use the iron horn, you must be proficient with all martial weapons.

If you blow the horn without meeting its requirement, the summoned berserkers attack you. If you meet the requirement, they are friendly to you and your companions and follow your commands.

\DndItemHeader{Horn of Valhalla, Silver}{r, a, r, e}
You can use an action to blow this horn. In response, warrior spirits from the plane of Ysgard appear within 60 feet of you. These spirits use the \textbf{berserker} statistics. They return to Ysgard after 1 hour or when they drop to 0 hit points. Once you use the horn, it can't be used again until 7 days have passed.

The silver horn summons 2d4 + 2 berserkers.

The berserkers are friendly to you and your companions and follow your commands.

\DndItemHeader{Horseshoes of Speed}{r, a, r, e}
These iron horseshoes come in a set of four. While all four shoes are affixed to the hooves of a horse or similar creature, they increase the creature's walking speed by 30 feet.

\DndItemHeader{Horseshoes of a Zephyr}{v, e, r, y,  , r, a, r, e}
These iron horseshoes come in a set of four. While all four shoes are affixed to the hooves of a horse or similar creature, they allow the creature to move normally while floating 4 inches above the ground. This effect means the creature can cross or stand above nonsolid or unstable surfaces, such as water or lava. The creature leaves no tracks and ignores difficult terrain. In addition, the creature can move at normal speed for up to 12 hours a day without suffering exhaustion from a forced march.

\DndItemHeader{Immovable Rod}{u, n, c, o, m, m, o, n}
This flat iron rod has a button on one end. You can use an action to press the button, which causes the rod to become magically fixed in place. Until you or another creature uses an action to push the button again, the rod doesn't move, even if it is defying gravity. The rod can hold up to 8,000 pounds of weight. More weight causes the rod to deactivate and fall. A creature can use an action to make a DC 30 Strength check, moving the fixed rod up to 10 feet on a success.

\DndItemHeader{Infernal Amulet}{r, a, r, e None}
While wearing this amulet, you can use it as a spellcasting focus for your spells, and it grants a +2 bonus to your spell save DC and spell attack bonus.

\subparagraph{Curse}
You're unwilling to part with this amulet while attuned to it and you wear it always. While wearing the amulet you have disadvantage on Strength saving throws and Strength checks. Only the \textit{Remove Curse} spell allows you to remove the item and end attunement.

\subparagraph{Corrupting}
This item corrupts. See the "Infernal Item Corruption" section.

\DndItemHeader{Infernal Plate Armor}{v, e, r, y,  , r, a, r, e None}
While wearing this armor, you gain a +2 bonus to AC.

\subparagraph{Curse}
Once you wear this armor, and are attuned to it, you can't remove it. Only the \textit{Remove Curse} spell allows you to end the attunement and finally doff it. While wearing the armor, you're vulnerable to the following damage types: force, lightning, psychic, radiant, and thunder.

\subparagraph{Corrupting}
This item corrupts. See the "Infernal Item Corruption" section.

\DndItemHeader{Instrument of the Bards, Ollamh Harp}{l, e, g, e, n, d, a, r, y None}
An instrument of the bards is an exquisite example of its kind, superior to an ordinary instrument in every way. Seven types of these instruments exist, each named after a legendary bard college. A creature that attempts to play the instrument without being attuned to it must succeed on a DC 15 Wisdom saving throw or take 2d4 psychic damage.

You can use an action to play the instrument and cast one of its spells. Once the instrument has been used to cast a spell, it can't be used to cast that spell again until the next dawn. The spells use your spellcasting ability and spell save DC.

You can play the instrument while casting a spell that causes any of its targets to be charmed on a failed saving throw, thereby imposing disadvantage on the save. This effect applies only if the spell has a somatic or a material component.

All instruments of the bards can be used to cast the following spells: \textit{fly}, \textit{invisibility}, \textit{levitate}, and \textit{protection from evil and good}.

In addition, the Ollamh harp can be used to cast \textit{confusion}, \textit{control weather}, and \textit{fire storm}.

\DndItemHeader{Ioun Stone, Reserve}{r, a, r, e None}
An \textit{Ioun stone} is named after Ioun, a god of knowledge and prophecy revered on some worlds. Many types of \textit{Ioun stone} exist, each type a distinct combination of shape and color.

When you use an action to toss one of these stones into the air, the stone orbits your head at a distance of 1d3 feet and confers a benefit to you. Thereafter, another creature must use an action to grasp or net the stone to separate it from you, either by making a successful attack roll against AC 24 or a successful DC 24 Dexterity (Acrobatics) check. You can use an action to seize and stow the stone, ending its effect.

A stone has AC 24, 10 hit points, and resistance to all damage. It is considered to be an object that is being worn while it orbits your head.

This vibrant purple prism stores spells cast into it, holding them until you use them. The stone can store up to 3 levels worth of spells at a time. When found, it contains 1d4 - 1 levels of stored spells chosen by the DM.

Any creature can cast a spell of 1st through 3rd level into the stone by touching it as the spell is cast. The spell has no effect, other than to be stored in the stone. If the stone can't hold the spell, the spell is expended without effect. The level of the slot used to cast the spell determines how much space it uses.

While this stone orbits your head, you can cast any spell stored in it. The spell uses the slot level, spell save DC, spell attack bonus, and spellcasting ability of the original caster, but is otherwise treated as if you cast the spell. The spell cast from the stone is no longer stored in it, freeing up space.

\DndItemHeader{Iron Bands of Bilarro}{r, a, r, e}
This rusty iron sphere measures 3 inches in diameter and weighs 1 pound. You can use an action to speak the command word and throw the sphere at a Huge or smaller creature you can see within 60 feet of you. As the sphere moves through the air, it opens into a tangle of metal bands.

Make a ranged attack roll with an attack bonus equal to your Dexterity modifier plus your proficiency bonus. On a hit, the target is restrained until you take a bonus action to speak the command word again to release it. Doing so, or missing with the attack, causes the bands to contract and become a sphere once more.

A creature, including the one restrained, can use an action to make a DC 20 Strength check to break the iron bands. On a success, the item is destroyed, and the restrained creature is freed. If the check fails, any further attempts made by that creature automatically fail until 24 hours have elapsed.

Once the bands are used, they can't be used again until the next dawn.

\DndItemHeader{Javelin of Lightning}{u, n, c, o, m, m, o, n}
This javelin is a magic weapon. When you hurl it and speak its command word, it transforms into a bolt of lightning, forming a line 5 feet wide that extends out from you to a target within 120 feet. Each creature in the line excluding you and the target must make a DC 13 Dexterity saving throw, taking 4d6 lightning damage on a failed save, and half as much damage on a successful one. The lightning bolt turns back into a javelin when it reaches the target. Make a ranged weapon attack against the target. On a hit, the target takes damage from the javelin plus 4d6 lightning damage.

The javelin's property can't be used again until the next dawn. In the meantime, the javelin can still be used as a magic weapon.

\DndItemHeader{Knife of Stolen Resistance}{r, a, r, e}
Using an action, you carve a single infernal rune into the flesh of an unconscious Beast, Celestial, Dragon, Fey, or Giant with this knife. Over the next 10 minutes the creature dies an agonizing death that can't be prevented short of the \textit{Wish} spell. If the creature has any resistances or immunities, you gain those resistances and immunities until the creature dies or a \textit{Wish} spell is used to save the creature. The knife's power can't be used again until you finish a long rest.

\DndItemHeader{Mace of Smiting}{r, a, r, e}
You gain a +1 bonus to attack and damage rolls made with this magic weapon. The bonus increases to +3 when you use the mace to attack a construct.

When you roll a 20 on an attack roll made with this weapon, the target takes an extra 7 bludgeoning damage, or an extra 14 bludgeoning damage if it's a construct. If a construct has 25 hit points or fewer after taking this damage, it is destroyed.

Note: According to the SRD, it is an extra 2d6 and 4d6 bludgeoning damage, although this is incorrect.

\DndItemHeader{Mace of Terror}{r, a, r, e None}
This magic weapon has 3 charges. While holding it, you can use an action and expend 1 charge to release a wave of terror. Each creature of your choice in a 30-foot radius extending from you must succeed on a DC 15 Wisdom saving throw or become frightened of you for 1 minute. While it is frightened in this way, a creature must spend its turns trying to move as far away from you as it can, and it can't willingly move to a space within 30 feet of you. It also can't take reactions. For its action it can use only the Dash action or try to escape from an effect that prevents it from moving. If it has nowhere it can move, the creature can use the Dodge action. At the end of each of its turns, a creature can repeat the saving throw, ending the effect on itself on a success.

The mace regains 1d3 expended charges daily at dawn.

\DndItemHeader{Mantle of Spell Resistance}{r, a, r, e None}
You have advantage on saving throws against spells while you wear this cloak.

\DndItemHeader{Manual of Bodily Health}{v, e, r, y,  , r, a, r, e}
This book contains health and diet tips, and its words are charged with magic. If you spend 48 hours over a period of 6 days or fewer studying the book's contents and practicing its guidelines, your Constitution score increases by 2, as does your maximum for that score. The manual then loses its magic, but regains it in a century.

\DndItemHeader{Manual of Gainful Exercise}{v, e, r, y,  , r, a, r, e}
This book describes fitness exercises, and its words are charged with magic. If you spend 48 hours over a period of 6 days or fewer studying the book's contents and practicing its guidelines, your Strength score increases by 2, as does your maximum for that score. The manual then loses its magic, but regains it in a century.

\DndItemHeader{Manual of Iron Golems}{v, e, r, y,  , r, a, r, e}
This tome contains information and incantations necessary to make a particular type of golem. The DM chooses the type or determines it randomly. To decipher and use the manual, you must be a spellcaster with at least two 5th-level spell slots. A creature that can't use a manual of golems and attempts to read it takes 6d6 psychic damage.

To create an \textbf{iron golem}, you must spend 120 days, working without interruption with the manual at hand and resting no more than 8 hours per day. You must also pay 100,000 gp to purchase supplies. Once you finish creating the golem, the book is consumed in eldritch flames. The golem becomes animate when the ashes of the manual are sprinkled on it. It is under your control, and it understands and obeys your spoken commands.

\DndItemHeader{Manual of Quickness of Action}{v, e, r, y,  , r, a, r, e}
This book contains coordination and balance exercises, and its words are charged with magic. If you spend 48 hours over a period of 6 days or fewer studying the book's contents and practicing its guidelines, your Dexterity score increases by 2, as does your maximum for that score. The manual then loses its magic, but regains it in a century.

\DndItemHeader{Medallion of Thoughts}{u, n, c, o, m, m, o, n None}
The medallion has 3 charges. While wearing it, you can use an action and expend 1 charge to cast the \textit{detect thoughts} spell (save DC 13) from it. The medallion regains 1d3 expended charges daily at dawn.

\DndItemHeader{Mirror of Life Trapping}{v, e, r, y,  , r, a, r, e}
When this 4-foot-tall mirror is viewed indirectly, its surface shows faint images of creatures. The mirror weighs 50 pounds, and it has AC 11, 10 hit points, and vulnerability to bludgeoning damage. It shatters and is destroyed when reduced to 0 hit points.

If the mirror is hanging on a vertical surface and you are within 5 feet of it, you can use an action to speak its command word and activate it. It remains activated until you use an action to speak the command word again.

Any creature other than you that sees its reflection in the activated mirror while within 30 feet of it must succeed on a DC 15 Charisma saving throw or be trapped, along with anything it is wearing or carrying, in one of the mirror's twelve extradimensional cells. This saving throw is made with advantage if the creature knows the mirror's nature, and constructs succeed on the saving throw automatically.

An extradimensional cell is an infinite expanse filled with thick fog that reduces visibility to 10 feet. Creatures trapped in the mirror's cells don't age, and they don't need to eat, drink, or sleep. A creature trapped within a cell can escape using magic that permits planar travel. Otherwise, the creature is confined to the cell until freed.

If the mirror traps a creature but its twelve extradimensional cells are already occupied, the mirror frees one trapped creature at random to accommodate the new prisoner. A freed creature appears in an unoccupied space within sight of the mirror but facing away from it. If the mirror is shattered, all creatures it contains are freed and appear in unoccupied spaces near it.

While within 5 feet of the mirror, you can use an action to speak the name of one creature trapped in it or call out a particular cell by number. The creature named or contained in the named cell appears as an image on the mirror's surface. You and the creature can then communicate normally.

In a similar way, you can use an action to speak a second command word and free one creature trapped in the mirror. The freed creature appears, along with its possessions, in the unoccupied space nearest to the mirror and facing away from it.

\DndItemHeader{Oil of Sharpness}{v, e, r, y,  , r, a, r, e}
This clear, gelatinous oil sparkles with tiny, ultrathin silver shards. The oil can coat one slashing or piercing weapon or up to 5 pieces of slashing or piercing ammunition. Applying the oil takes 1 minute. For 1 hour, the coated item is magical and has a +3 bonus to attack and damage rolls.

\DndItemHeader{Oil of Slipperiness}{u, n, c, o, m, m, o, n}
This sticky black unguent is thick and heavy in the container, but it flows quickly when poured. The oil can cover a Medium or smaller creature, along with the equipment it's wearing and carrying (one additional vial is required for each size category above Medium). Applying the oil takes 10 minutes. The affected creature then gains the effect of a \textit{freedom of movement} spell for 8 hours.

Alternatively, the oil can be poured on the ground as an action, where it covers a 10-foot square, duplicating the effect of the \textit{grease} spell in that area for 8 hours.

\DndItemHeader{Pipes of Haunting}{u, n, c, o, m, m, o, n}
You must be proficient with wind instruments to use these pipes. They have 3 charges. You can use an action to play them and expend 1 charge to create an eerie, spellbinding tune. Each creature within 30 feet of you that hears you play must succeed on a DC 15 Wisdom saving throw or become frightened of you for 1 minute. If you wish, all creatures in the area that aren't hostile toward you automatically succeed on the saving throw. A creature that fails the saving throw can repeat it at the end of each of its turns, ending the effect on itself on a success. A creature that succeeds on its saving throw is immune to the effect of these pipes for 24 hours. The pipes regain 1d3 expended charges daily at dawn.

\DndItemHeader{Plate Armor of Etherealness}{l, e, g, e, n, d, a, r, y None}
While you're wearing this armor, you can speak its command word as an action to gain the effect of the \textit{etherealness} spell, which lasts for 10 minutes or until you remove the armor or use an action to speak the command word again. This property of the armor can't be used again until the next dawn.

\DndItemHeader{Potion of Animal Friendship}{u, n, c, o, m, m, o, n}
When you drink this potion, you can cast the \textit{animal friendship} spell (save DC 13) for 1 hour at will. Agitating this muddy liquid brings little bits into view: a fish scale, a hummingbird tongue, a cat claw, or a squirrel hair.

\DndItemHeader{Potion of Clairvoyance}{r, a, r, e}
When you drink this potion, you gain the effect of the \textit{clairvoyance} spell. An eyeball bobs in this yellowish liquid but vanishes when the potion is opened.

\DndItemHeader{Potion of Cloud Giant Strength}{v, e, r, y,  , r, a, r, e}
When you drink this potion, your Strength score changes to 27 for 1 hour. The potion has no effect on you if your Strength is equal to or greater than that score.

This potion's transparent liquid has floating in it a sliver of fingernail from a \textbf{cloud giant}.

\DndItemHeader{Potion of Diminution}{r, a, r, e}
When you drink this potion, you gain the "reduce" effect of the \textit{enlarge/reduce} spell for 1d4 hours (no concentration required). The red in the potion's liquid continuously contracts to a tiny bead and then expands to color the clear liquid around it. Shaking the bottle fails to interrupt this process.

\DndItemHeader{Potion of Fire Breath}{u, n, c, o, m, m, o, n}
After drinking this potion, you can use a bonus action to exhale fire at a target within 30 feet of you. The target must make a DC 13 Dexterity saving throw, taking 4d6 fire damage on a failed save, or half as much damage on a successful one. The effect ends after you exhale the fire three times or when 1 hour has passed. This potion's orange liquid flickers, and smoke fills the top of the container and wafts out whenever it is opened.

\DndItemHeader{Potion of Fire Giant Strength}{r, a, r, e}
When you drink this potion, your Strength score changes to 25 for 1 hour. The potion has no effect on you if your Strength is equal to or greater than that score.

This potion's transparent liquid has floating in it a sliver of fingernail from a \textbf{fire giant}.

\DndItemHeader{Potion of Flying}{v, e, r, y,  , r, a, r, e}
When you drink this potion, you gain a flying speed equal to your walking speed for 1 hour and can hover. If you're in the air when the potion wears off, you fall unless you have some other means of staying aloft. This potion's clear liquid floats at the top of its container and has cloudy white impurities drifting in it.

\DndItemHeader{Potion of Frost Giant Strength}{r, a, r, e}
When you drink this potion, your Strength score changes to 23 for 1 hour. The potion has no effect on you if your Strength is equal to or greater than that score.

This potion's transparent liquid has floating in it a sliver of fingernail from a \textbf{frost giant}.

\DndItemHeader{Potion of Gaseous Form}{r, a, r, e}
When you drink this potion, you gain the effect of the \textit{gaseous form} spell for 1 hour (no concentration required) or until you end the effect as a bonus action. This potion's container seems to hold fog that moves and pours like water.

\DndItemHeader{Potion of Growth}{u, n, c, o, m, m, o, n}
When you drink this potion, you gain the "enlarge" effect of the \textit{enlarge/reduce} spell for 1d4 hours (no concentration required). The red in the potion's liquid continuously expands from a tiny bead to color the clear liquid around it and then contracts. Shaking the bottle fails to interrupt this process.

\DndItemHeader{Potion of Heroism}{r, a, r, e}
For 1 hour after drinking it, you gain 10 temporary hit points that last for 1 hour. For the same duration, you are under the effect of the \textit{bless} spell (no concentration required). This blue potion bubbles and steams as if boiling.

\DndItemHeader{Potion of Hill Giant Strength}{u, n, c, o, m, m, o, n}
When you drink this potion, your Strength score changes to 21 for 1 hour. The potion has no effect on you if your Strength is equal to or greater than that score.

This potion's transparent liquid has floating in it a sliver of fingernail from a \textbf{hill giant}.

\DndItemHeader{Potion of Invisibility}{v, e, r, y,  , r, a, r, e}
This potion's container looks empty but feels as though it holds liquid. When you drink it, you become invisible for 1 hour. Anything you wear or carry is invisible with you. The effect ends early if you attack or cast a spell.

\DndItemHeader{Potion of Longevity}{v, e, r, y,  , r, a, r, e}
When you drink this potion, your physical age is reduced by 1d6 + 6 years, to a minimum of 13 years. Each time you subsequently drink a potion of longevity, there is a You become younger cumulative chance that you instead age by 1d6 + 6 years. Suspended in this amber liquid are a scorpion's tail, an adder's fang, a dead spider, and a tiny heart that, against all reason, is still beating. These ingredients vanish when the potion is opened.

\DndItemHeader{Potion of Mind Reading}{r, a, r, e}
When you drink this potion, you gain the effect of the \textit{detect thoughts} spell (save DC 13). The potion's dense, purple liquid has an ovoid cloud of pink floating in it.

\DndItemHeader{Potion of Poison}{u, n, c, o, m, m, o, n}
This concoction looks, smells, and tastes like a \textit{potion of healing} or other beneficial potion. However, it is actually poison masked by illusion magic. An \textit{identify} spell reveals its true nature.

If you drink it, you take 3d6 poison damage, and you must succeed on a DC 13 Constitution saving throw or be poisoned. At the start of each of your turns while you are poisoned in this way, you take 3d6 poison damage. At the end of each of your turns, you can repeat the saving throw. On a successful save, the poison damage you take on your subsequent turns decreases by 1d6. The poison ends when the damage decreases to 0.

\DndItemHeader{Potion of Speed}{v, e, r, y,  , r, a, r, e}
When you drink this potion, you gain the effect of the \textit{haste} spell for 1 minute (no concentration required). The potion's yellow fluid is streaked with black and swirls on its own.

\DndItemHeader{Potion of Stone Giant Strength}{r, a, r, e}
When you drink this potion, your Strength score changes to 23 for 1 hour. The potion has no effect on you if your Strength is equal to or greater than that score.

This potion's transparent liquid has floating in it a sliver of fingernail from a \textbf{stone giant}.

\DndItemHeader{Potion of Storm Giant Strength}{l, e, g, e, n, d, a, r, y}
When you drink this potion, your Strength score changes to 29 for 1 hour. The potion has no effect on you if your Strength is equal to or greater than that score.

This potion's transparent liquid has floating in it a sliver of fingernail from a \textbf{storm giant}.

\DndItemHeader{Potion of Supreme Healing}{v, e, r, y,  , r, a, r, e}
You regain 10d4 + 20 hit points when you drink this potion. The potion's red liquid glimmers when agitated.

\DndItemHeader{Potion of Vitality}{v, e, r, y,  , r, a, r, e}
When you drink this potion, it removes any exhaustion you are suffering and cures any disease or poison affecting you. For the next 24 hours, you regain the maximum number of hit points for any Hit Die you spend. The potion's crimson liquid regularly pulses with dull light, calling to mind a heartbeat.

\DndItemHeader{Potion of Water Breathing}{u, n, c, o, m, m, o, n}
You can breathe underwater for 1 hour after drinking this potion. Its cloudy green fluid smells of the sea and has a jellyfish-like bubble floating in it.

\DndItemHeader{Ranseur of Torture}{a, r, t, i, f, a, c, t None}
Forged of infernal iron and Hellfire through the help of Dispater, the Ranseur of Torture has been Belial's weapon of choice since he first became an archdevil.

It is so directly tied to him that one of his worshipped symbols is a two-pronged fork. The weapon never leaves his side, and he is a master at wielding it.

The Ranseur of Torture is a magic pike. You have a +3 bonus to attack and damage rolls made with this weapon, and it does 2d10 piercing damage instead of its regular damage. When you hit a good aligned creature with it, that creature takes an extra 1d10 necrotic damage.

\subparagraph{Random Properties}
The Ranseur of Torture has the following random properties:


\begin{itemize}
\item 1 Artifact Properties; Minor Beneficial Properties property
\item 1 Artifact Properties; Major Beneficial Properties property
\item 1 Artifact Properties; Minor Detrimental Properties property
\end{itemize}

\subparagraph{Wounding}
Damage dealt by the Ranseur of Torture is particularly dangerous, as it can only be restored through magical means. Resting, food, and medicine won't restore hit points removed by the pike. If the damage goes untreated for 7 days, the missing hit points are permanently removed, and can be recovered only through casting a \textit{Regenerate} or \textit{Wish} spell on the victim.

\subparagraph{Execution}
If a creature with less than its full hit points is hit by this pike, the Ranseur of Torture deals 1d10 additional piercing damage. If a creature is killed by the pike, their soul is sent to Phlegethos and can be restored only through casting a \textit{Wish} spell on the victim.

\subparagraph{Staggered}
Each time a creature takes damage from the Ranseur of Torture, they must make a Constitution save, with a DC equaling the damage they took. On a failed save, their movement speed is reduced by 5 feet until the end of their next turn. This penalty stacks, though movement speed can't be reduced below 0.

\subparagraph{Interrupting}
Spellcasters damaged by the Ranseur of Torture have disadvantage on Constitution checks made to maintain concentration until the damage is healed. If a creature is concentration on a spell when damaged by the ranseur, their concentration is immediately ended.

\subparagraph{Destroying the Ranseur}
Only through the written consent of both Dispater and Belial can the Ranseur of Torture be destroyed. If either or both of the archdevils are dead, then the new archdevil to take their mantle must give permission in their stead. The two writs must be wrapped around each of the prongs of the weapon, then the weapon must be tossed into the Pit of Flame. If performed successfully, the Hellfire briefly burns a different hue.

\DndItemHeader{Ring of Collecting}{v, e, r, y,  , r, a, r, e None}
While wearing this ring you can use it to cast the \textit{Leomund's Tiny Hut} spell as an action. Once this property is used, it can't be used again until the next dawn. Additionally, as a bonus action, you can use the ring to disintegrate any nonmagical piece of art (drawing, painting, or sculpture) within 60 feet that is no larger than Medium-sized. That art now appears inside the tiny hut, for you to appreciate whenever you desire. If you try to remove this art from the tiny hut, it vanishes forever. You can steadily improve this space, but only be denying such beauty to the rest of the world. For every 1,000 gp of art acquired, the tiny hut increases in size by adding another foot to its radius and lasts one additional hour.

\DndItemHeader{Ring of Mind Shielding}{u, n, c, o, m, m, o, n None}
While wearing this ring, you are immune to magic that allows other creatures to read your thoughts, determine whether you are lying, know your alignment, or know your creature type. Creatures can telepathically communicate with you only if you allow it.

You can use an action to cause the ring to become invisible until you use another action to make it visible, until you remove the ring, or until you die.

If you die while wearing the ring, your soul enters it, unless it already houses a soul. You can remain in the ring or depart for the afterlife. As long as your soul is in the ring, you can telepathically communicate with any creature wearing it. A wearer can't prevent this telepathic communication.

\DndItemHeader{Ring of Protection}{r, a, r, e None}
You gain a +1 bonus to AC and saving throws while wearing this ring.

\DndItemHeader{Ring of Shooting Stars}{v, e, r, y,  , r, a, r, e None}
While wearing this ring in dim light or darkness, you can cast \textit{dancing lights} and \textit{light} from the ring at will. Casting either spell from the ring requires an action.

The ring has 6 charges for the following other properties. The ring regains 1d6 expended charges daily at dawn.

\subparagraph{Faerie Fire}
You can expend 1 charge as an action to cast \textit{faerie fire} from the ring.

\subparagraph{Ball Lightning}
You can expend 2 charges as an action to create one to four 3-foot-diameter spheres of lightning. The more spheres you create, the less powerful each sphere is individually.

Each sphere appears in an unoccupied space you can see within 120 feet of you. The spheres last as long as you concentrate (as if concentration on a spell), up to 1 minute. Each sphere sheds dim light in a 30-foot radius.

As a bonus action, you can move each sphere up to 30 feet, but no farther than 120 feet away from you. When a creature other than you comes within 5 feet of a sphere, the sphere discharges lightning at that creature and disappears. That creature must make a DC 15 Dexterity saving throw. On a failed save, the creature takes lightning damage based on the number of spheres you created. (4 spheres = 2d4, 3 spheres = 2d6, 2 spheres = 5d4, 1 sphere = 4d12)

\subparagraph{Shooting Stars}
You can expend 1 to 3 charges as an action. For every charge you expend, you launch a glowing mote of light from the ring at a point you can see within 60 feet of you. Each creature within a 15-foot cube originating from that point is showered in sparks and must make a DC 15 Dexterity saving throw, taking 5d4 fire damage on a failed save, or half as much damage on a successful one.

\DndItemHeader{Ring of Spell Storing}{r, a, r, e None}
This ring stores spells cast into it, holding them until the attuned wearer uses them. The ring can store up to 5 levels worth of spells at a time. When found, it contains 1d6 - 1 levels of stored spells chosen by the DM.

Any creature can cast a spell of 1st through 5th level into the ring by touching the ring as the spell is cast. The spell has no effect, other than to be stored in the ring. If the ring can't hold the spell, the spell is expended without effect. The level of the slot used to cast the spell determines how much space it uses.

While wearing this ring, you can cast any spell stored in it. The spell uses the slot level, spell save DC, spell attack bonus, and spellcasting ability of the original caster, but is otherwise treated as if you cast the spell. The spell cast from the ring is no longer stored in it, freeing up space.

\DndItemHeader{Ring of Treachery}{v, e, r, y,  , r, a, r, e None}
This ring has 3 charges. While wearing this ring, when you're damaged, you may use a reaction to expend a charge and transfer that damage to a random creature (which could include an ally) within 60 feet. All charges are restored when you finish a long rest.

\DndItemHeader{Ring of the Copycat}{l, e, g, e, n, d, a, r, y None}
You gain the ability to channel energy from allies to cast spells, even if you normally can't. When an ally within 60 feet of you casts a spell that you normally can't cast and that requires 10 gp or fewer in material components, you may use a reaction to cast that spell. When cast in this way, the spell is cast using your ally's spellcasting ability, spell save DC, and spell attack bonus, as needed. You decide this cloned spell's target, as specified in the spell's description, and the spell originates from you. After you use this ability, your ally can't cast this spell again until they finish a long rest.

\DndItemHeader{Robe of Eyes}{r, a, r, e None}
This robe is adorned with eyelike patterns. While you wear the robe, you gain the following benefits:


\begin{itemize}
\item The robe lets you see in all directions, and you have advantage on Wisdom (Perception) checks that rely on sight.
\item You have darkvision out to a range of 120 feet.
\item You can see invisible creatures and objects, as well as see into the Ethereal Plane, out to a range of 120 feet.
\end{itemize}

The eyes on the robe can't be closed or averted. Although you can close or avert your own eyes, you are never considered to be doing so while wearing this robe.

A \textit{light} spell cast on the robe or a \textit{daylight} spell cast within 5 feet of the robe causes you to be blinded for 1 minute. At the end of each of your turns, you can make a Constitution saving throw (DC 11 for \textit{light} or DC 15 for \textit{daylight}), ending the blindness on a success.

\DndItemHeader{Robe of Stars}{v, e, r, y,  , r, a, r, e None}
This black or dark blue robe is embroidered with small white or silver stars. You gain a +1 bonus to saving throws while you wear it.

Six stars, located on the robe's upper front portion, are particularly large. While wearing this robe, you can use an action to pull off one of the stars and use it to cast \textit{magic missile} as a 5th-level spell. Daily at dusk, 1d6 removed stars reappear on the robe.

While you wear the robe, you can use an action to enter the Astral Plane along with everything you are wearing and carrying. You remain there until you use an action to return to the plane you were on. You reappear in the last space you occupied, or if that space is occupied, the nearest unoccupied space.

\DndItemHeader{Rod of Resurrection}{l, e, g, e, n, d, a, r, y None}
The rod has 5 charges. While you hold it, you can use an action to cast one of the following spells from it: \textit{heal} (expends 1 charge) or \textit{resurrection} (expends 5 charges).

The rod regains 1 expended charge daily at dawn. If the rod is reduced to 0 charges, roll a d20. On a 1, the rod disappears in a burst of radiance.

\DndItemHeader{Ruby Rod of Asmodeus}{a, r, t, i, f, a, c, t None}
Asmodeus's Ruby Rod is said to have been constructed from a shard of pure evil, a ruby soaked in the blood of the innocents, the drool from a draconic deity, and the tears of 777 angels. It is a magic rod that functions as a morningstar and you have a +2 bonus to attack and damage rolls made with it. All damage dealt with it is converted to necrotic, and when you hit a good aligned creature, that creature takes an extra 2d12 necrotic damage.

The Ruby Rod of Asmodeus has 66 charges and regains 6d6 charges each nightfall. Asmodeus is always capable of perceiving through the rod and knows the exact location of any creature attuned to it. If a creature uses the special abilities of the rod without Asmodeus's permission, they take 2d12 necrotic damage each use, doubled if they're good aligned. This damage ignores resistances and immunities, and if it kills the user, Asmodeus permanently claims their soul.

\subparagraph{Random Properties}
The Ruby Rod of Asmodeus has the following random properties:


\begin{itemize}
\item 1 Artifact Properties; Minor Beneficial Properties property
\item 1 Artifact Properties; Major Beneficial Properties property
\item 2 Artifact Properties; Major Detrimental Properties properties
\end{itemize}

Asmodeus is immune to the detrimental properties of the weapon.

\subparagraph{Bow to the King}
You may use an action to expend 33 charges to use the rod to force all creatures within a 150-foot-radius to bow to you. A creature must succeed on a DC 21 Charisma saving throw or be forced to kneel for up to 1 hour. Creatures immune to the charmed condition have advantage on their save. Each time a kneeling creature takes damage, they may repeat the save, ending the effect on a success. The effect also ends early for kneeling creatures if they're touched by the Ruby Rod.

\subparagraph{Acid, Cold, and Lightning}
At a cost of 9 charges each use, you may use an action to trigger one of the following effects:


\begin{itemize}
\item You project acid in a 60-foot line that is 5 feet wide. Each creature in that line must make a DC 18 Dexterity saving throw, taking 12d8 acid damage on a failed save, or half as much damage on a successful one.
\item You project an icy blast in a 60-foot cone. Each creature in that area must make a DC 19 Constitution saving throw, taking 12d8 cold damage on a failed save, or half as much damage on a successful one.
\item You shoot lightning in a 90-foot line that is 5 feet wide. Each creature in that line must make a DC 19 Dexterity saving throw, taking 12d10 lightning damage on a failed save, or half as much damage on a successful one.
\end{itemize}

\subparagraph{Reactions}
As a reaction to suffering a detrimental effect, you may cast \textit{Lesser Restoration} (expending 3 charges) or \textit{Greater Restoration} (expending 12 charges.)

Additionally, if you're the sole target of a spell, you may use a reaction to attempt to use the rod to absorb the spell and cancel its effect. You roll 1d20 + 2, and if the result is higher than the spell level + 10, the spell is absorbed into the rod. The rod regains charges equal to the spell level it absorbed. If the absorbed charges go over the maximum of 66 charges, it overloads, dealing 1d6 necrotic damage per charge to you and expending all charges.

\subparagraph{Destroying the Rod}
While Asmodeus lives, the Ruby Rod of Asmodeus can't be destroyed, and if attempted it simply returns to Asmodeus. If he were defeated, the rod could be destroyed by bringing it to the Abyss from whence it was created, then performing a ritual known only to Asmodeus.

\DndItemHeader{Sage's Mirror}{r, a, r, e}
This item has 3 charges and regains all charges at dusk. You can use an action and expend 1 of the mirror's charges to cast one of the following spells:


\begin{itemize}
\item \textit{Find the Path}
\item \textit{Legend Lore}
\end{itemize}

With each use, it is apparent that the information gleaned from the mirror comes from a chamber in the Nine Hells where sages and scholars are tortured for the answers to each question.

\DndItemHeader{Scourge of Shadow}{a, r, t, i, f, a, c, t None}
Said to be a gift to Glasya from Asmodeus, the Scourge of Shadow is a multi-headed whip constructed from infernal iron. The scourge is nigh-in- destructible, and Glasya carries it with her wherever she goes. When she wields the whip in battle, it shrouds her in an aura of shadow.

The Scourge of Shadow is a magic whip. You have a +3 bonus to attack and damage rolls made with this weapon, and it does 3d4 slashing damage instead of its regular damage. When you hit a good aligned creature with it, that creatures takes an extra 4d4 necrotic damage.

Whenever an unattuned creature attacks with the weapon, it deals 2d4 necrotic damage, ignoring resistances and immunities, to them, but otherwise functions normally.

\subparagraph{Random Properties}
The Scourge of Shadow has the following random properties:


\begin{itemize}
\item 2 Artifact Properties; Minor Beneficial Properties properties
\item 1 Artifact Properties; Major Beneficial Properties property
\item 2 Artifact Properties; Minor Detrimental Properties properties
\item 1 Artifact Properties; Major Detrimental Properties property
\end{itemize}

\subparagraph{Keen}
The whip scores a critical on a roll of 18, 19, or 20. When it scores a critical, its damage dice are upgraded to 3d6 instead of 3d4 for the duration of the attack. If attacking a good aligned creature, or attacking with the weapon unattuned, the necrotic damage dice are also increased to a d6 instead of a d4.

\subparagraph{Fast Reflexes}
While attuned to the Scourge of Shadow, you gain advantage on initiative rolls, and you add your Charisma and Dexterity modifiers to the roll. You may use either your Strength, Dexterity, or Charisma modifier for the attack and damage rolls (the same modifier must be used for both attack and damage).

\subparagraph{Veil of Shadow}
You're protected by a field of shadow that makes you resistant to radiant and necrotic damage. While in Vision and Light or darkness you gain blindsight out to 60 feet and Dexterity (Stealth) checks are made with advantage.

\subparagraph{Destroying the Scourge}
Though the process is known only to Glasya and Asmodeus, the Scourge of Shadow can be destroyed. It must be taken to the Pit of Fire in Phlegethos. The scourge, along with the body of an archdevil, must be thrown into the Hellfire together. As they hit the flames, a spell of Darkness must be cast at 9th level on the scourge. If each part of the ritual is performed correctly, the whip disintegrates within the Hellfire.

\DndItemHeader{Sending Stones}{u, n, c, o, m, m, o, n}
Sending stones come in pairs, with each smooth stone carved to match the other so the pairing is easily recognized. While you touch one stone, you can use an action to cast the \textit{sending} spell from it. The target is the bearer of the other stone. If no creature bears the other stone, you know that fact as soon as you use the stone and don't cast the spell.

Once \textit{sending} is cast through the stones, they can't be used again until the next dawn. If one of the stones in a pair is destroyed, the other one becomes nonmagical.

\DndItemHeader{Skull of Selfish Knowledge}{r, a, r, e None}
You may use an action to make the magical skull devour a nonmagical book, map, or scroll. Once devoured the learning is forever available to you, but you can never write the information down or communicate it to others. It is for you alone.

\DndItemHeader{Slippers of Spider Climbing}{u, n, c, o, m, m, o, n None}
While you wear these light shoes, you can move up, down, and across vertical surfaces and upside down along ceilings, while leaving your hands free. You have a climbing speed equal to your walking speed. However, the slippers don't allow you to move this way on a slippery surface, such as one covered by ice or oil.

\DndItemHeader{Soul Coin}{u, n, c, o, m, m, o, n}
Soul Coins are about 5 inches across and about an inch thick, minted from infernal iron. Each coin weighs one-third of a pound and is inscribed with Infernal writing and a spell that magically binds a single soul to the coin. Because each Soul Coin has a unique soul trapped within it, each has a story. A creature might have been imprisoned as a result of defaulting on a deal, while another might be the victim of a night hag's curse.

\subparagraph{Carrying Soul Coins}
To hold a Soul Coin is to feel the soul bound within it—overcome with rage or fraught with despair.

An evil creature can carry as many Soul Coins as it wishes (up to its maximum weight allowance). A non-evil creature can carry a number of Soul Coins equal to or less than its Constitution modifier without penalty. A non-evil creature carrying a number of Soul Coins greater than its Constitution modifier has disadvantage on its attack rolls, ability checks, and saving throws.

\subparagraph{Using a Soul Coin}
A Soul Coin has 3 charges. A creature carrying the coin can use its action to expend 1 charge from a Soul Coin and use it to do one of the following:


\begin{description}
\item[Drain Life.]
You siphon away some of the soul's essence and gain 1d10 temporary hit points.

\item[Query.]
You telepathically ask the soul a question and receive a brief telepathic response, which you can understand. The soul knows only what it knew in life, but it must answer you truthfully and to the best of its ability. The answer is no more than a sentence or two and might be cryptic.

\end{description}

\subparagraph{Freeing a Soul}
Casting a spell that removes a curse on a Soul Coin frees the soul trapped within it, as does expending all of the coin's charges. The coin itself rusts from within and is destroyed once the soul is released. A freed soul travels to the realm of the god it served, or the outer plane most closely tied to its alignment (DM's choice). The souls of lawful evil creatures released from Soul Coins typically emerge from the River Styx as lemure devils.

A soul can also be freed by destroying the coin that contains it. A Soul Coin has AC 19, 1 hit point for each charge it has remaining, and immunity to all damage except that which is dealt by an infernal warship's furnace.

Freeing a soul from a Soul Coin is considered a good act, even if the soul belongs to an evil creature.

\subparagraph{Hellish Currency}
Soul Coins are a currency of the Nine Hells and are highly valued by devils. The coins are used among the infernal hierarchy to barter for favors, bribe the unwilling, and reward the faithful for services rendered.

Soul Coins are created by Mammon and his greater devils on Minauros, the third layer of the Nine Hells, in a vast chamber where the captured souls of evil mortals are bound into the coins. These coins are then distributed throughout the Nine Hells to be used for goods and services, infernal deals, dark bargains, and bribes.

\DndItemHeader{Spell Scroll (1st Level)}{c, o, m, m, o, n}
A spell scroll bears the words of a single spell, written as a mystical cipher. If the spell is on your class's spell list, you can read the scroll and cast its spell without providing any material components. Otherwise, the scroll is unintelligible. Casting the spell by reading the scroll requires the spell's normal casting time. Once the spell is cast, the words on the scroll fade, and it crumbles to dust. If the casting is interrupted, the scroll is not lost.

If the spell is on your class's spell list but of a higher level than you can normally cast, you must make an ability check using your spellcasting ability to determine whether you cast it successfully. The DC is 11. On a failed check, the spell disappears from the scroll with no other effect.

Once the spell is cast, the words on the scroll fade, and the scroll itself crumbles to dust.

A spell cast from this scroll has a save DC of 13 and an attack bonus of 5.

A wizard spell on a spell scroll can be copied just as spells in spellbooks can be copied. When a spell is copied from a spell scroll, the copier must succeed on a DC 11 Intelligence (Arcana) check. If the check succeeds, the spell is successfully copied. Whether the check succeeds or fails, the spell scroll is destroyed.

\DndItemHeader{Spell Scroll (2nd Level)}{u, n, c, o, m, m, o, n}
A spell scroll bears the words of a single spell, written as a mystical cipher. If the spell is on your class's spell list, you can read the scroll and cast its spell without providing any material components. Otherwise, the scroll is unintelligible. Casting the spell by reading the scroll requires the spell's normal casting time. Once the spell is cast, the words on the scroll fade, and it crumbles to dust. If the casting is interrupted, the scroll is not lost.

If the spell is on your class's spell list but of a higher level than you can normally cast, you must make an ability check using your spellcasting ability to determine whether you cast it successfully. The DC is 12. On a failed check, the spell disappears from the scroll with no other effect.

Once the spell is cast, the words on the scroll fade, and the scroll itself crumbles to dust.

A spell cast from this scroll has a save DC of 13 and an attack bonus of 5.

A wizard spell on a spell scroll can be copied just as spells in spellbooks can be copied. When a spell is copied from a spell scroll, the copier must succeed on a DC 12 Intelligence (Arcana) check. If the check succeeds, the spell is successfully copied. Whether the check succeeds or fails, the spell scroll is destroyed.

\DndItemHeader{Spell Scroll (3rd Level)}{u, n, c, o, m, m, o, n}
A spell scroll bears the words of a single spell, written as a mystical cipher. If the spell is on your class's spell list, you can read the scroll and cast its spell without providing any material components. Otherwise, the scroll is unintelligible. Casting the spell by reading the scroll requires the spell's normal casting time. Once the spell is cast, the words on the scroll fade, and it crumbles to dust. If the casting is interrupted, the scroll is not lost.

If the spell is on your class's spell list but of a higher level than you can normally cast, you must make an ability check using your spellcasting ability to determine whether you cast it successfully. The DC is 13. On a failed check, the spell disappears from the scroll with no other effect.

Once the spell is cast, the words on the scroll fade, and the scroll itself crumbles to dust.

A spell cast from this scroll has a save DC of 15 and an attack bonus of 7.

A wizard spell on a spell scroll can be copied just as spells in spellbooks can be copied. When a spell is copied from a spell scroll, the copier must succeed on a DC 13 Intelligence (Arcana) check. If the check succeeds, the spell is successfully copied. Whether the check succeeds or fails, the spell scroll is destroyed.

\DndItemHeader{Spell Scroll (4th Level)}{r, a, r, e}
A spell scroll bears the words of a single spell, written as a mystical cipher. If the spell is on your class's spell list, you can read the scroll and cast its spell without providing any material components. Otherwise, the scroll is unintelligible. Casting the spell by reading the scroll requires the spell's normal casting time. Once the spell is cast, the words on the scroll fade, and it crumbles to dust. If the casting is interrupted, the scroll is not lost.

If the spell is on your class's spell list but of a higher level than you can normally cast, you must make an ability check using your spellcasting ability to determine whether you cast it successfully. The DC is 14. On a failed check, the spell disappears from the scroll with no other effect.

Once the spell is cast, the words on the scroll fade, and the scroll itself crumbles to dust.

A spell cast from this scroll has a save DC of 15 and an attack bonus of 7.

A wizard spell on a spell scroll can be copied just as spells in spellbooks can be copied. When a spell is copied from a spell scroll, the copier must succeed on a DC 14 Intelligence (Arcana) check. If the check succeeds, the spell is successfully copied. Whether the check succeeds or fails, the spell scroll is destroyed.

\DndItemHeader{Spell Scroll (5th Level)}{r, a, r, e}
A spell scroll bears the words of a single spell, written as a mystical cipher. If the spell is on your class's spell list, you can read the scroll and cast its spell without providing any material components. Otherwise, the scroll is unintelligible. Casting the spell by reading the scroll requires the spell's normal casting time. Once the spell is cast, the words on the scroll fade, and it crumbles to dust. If the casting is interrupted, the scroll is not lost.

If the spell is on your class's spell list but of a higher level than you can normally cast, you must make an ability check using your spellcasting ability to determine whether you cast it successfully. The DC is 15. On a failed check, the spell disappears from the scroll with no other effect.

Once the spell is cast, the words on the scroll fade, and the scroll itself crumbles to dust.

A spell cast from this scroll has a save DC of 17 and an attack bonus of 9.

A wizard spell on a spell scroll can be copied just as spells in spellbooks can be copied. When a spell is copied from a spell scroll, the copier must succeed on a DC 15 Intelligence (Arcana) check. If the check succeeds, the spell is successfully copied. Whether the check succeeds or fails, the spell scroll is destroyed.

\DndItemHeader{Spellguard Shield}{v, e, r, y,  , r, a, r, e None}
While holding this shield, you have advantage on saving throws against spells and other magical effects, and spell attacks have disadvantage against you.

\DndItemHeader{Staff of Fire}{v, e, r, y,  , r, a, r, e None}
You have resistance to fire damage while you hold this staff.

The staff has 10 charges. While holding it, you can use an action to expend 1 or more of its charges to cast one of the following spells from it, using your spell save DC: \textit{burning hands} (1 charge), \textit{fireball} (3 charges), or \textit{wall of fire} (4 charges).

The staff regains 1d6 + 4 expended charges daily at dawn. If you expend the last charge, roll a d20. On a 1, the staff blackens, crumbles into cinders, and is destroyed.

\DndItemHeader{Staff of Frost}{v, e, r, y,  , r, a, r, e None}
You have resistance to cold damage while you hold this staff.

The staff has 10 charges. While holding it, you can use an action to expend 1 or more of its charges to cast one of the following spells from it, using your spell save DC: \textit{cone of cold} (5 charges), \textit{fog cloud} (1 charge), \textit{ice storm} (4 charges), or \textit{wall of ice} (4 charges).

The staff regains 1d6 + 4 expended charges daily at dawn. If you expend the last charge, roll a d20. On a 1. the staff turns to water and is destroyed.

\DndItemHeader{Staff of Power}{v, e, r, y,  , r, a, r, e None}
This staff can be wielded as a magic quarterstaff that grants a +2 bonus to attack and damage rolls made with it. While holding it, you gain a +2 bonus to Armor Class, saving throws, and spell attack rolls.

The staff has 20 charges for the following properties. The staff regains 2d8 + 4 expended charges daily at dawn. If you expend the last charge, roll a d20. On a 1, the staff retains its +2 bonus to attack and damage roll but loses all other properties. On a 20, the staff regain 1d8 + 2 charges.

\subparagraph{Power Strike}
When you hit with a melee attack using the staff, you can expend 1 charge to deal an extra 1d6 force damage to the target.

\subparagraph{Spells}
While holding this staff, you can use an action to expend 1 or more of its charges to cast one of the following spells from it, using your spell save DC and spell attack bonus: \textit{cone of cold} (5 charges), \textit{fireball} (5th-level version, 5 charges), \textit{globe of invulnerability} (6 charges), \textit{hold monster} (5 charges), \textit{levitate} (2 charges), \textit{lightning bolt} (5th-level version, 5 charges), \textit{magic missile} (1 charge), \textit{ray of enfeeblement} (1 charge), or \textit{wall of force} (5 charges).

\subparagraph{Retributive Strike}
You can use an action to break the staff over your knee or against a solid surface, performing a retributive strike. The staff is destroyed and releases its remaining magic in an explosion that expands to fill a 30-foot-radius sphere centered on it.

You have a 50 chance to instantly travel to a random plane of existence, avoiding the explosion. If you fail to avoid the effect, you take force damage equal to 16 × the number of charges in the staff. Every other creature in the area must make a DC 17 Dexterity saving throw. On a failed save, a creature takes an amount of damage based on how far away it is from the point of origin, as shown in the following table. On a successful save, a creature takes half as much damage.


\begin{DndTable}{XX}

Distance from Origin & Effect\\
10 ft. away or closer & 8 × the number of charges in the staff\\
11 to 20 ft. away & 6 × the number of charges in the staff\\
21 to 30 ft. away & 4 × the number of charges in the staff\\
\end{DndTable}

\DndItemHeader{Staff of Withering}{r, a, r, e None}
This staff has 3 charges and regains 1d3 expended charges daily at dawn.

The staff can be wielded as a magic quarterstaff. On a hit, it deals damage as a normal quarterstaff, and you can expend 1 charge to deal an extra 2d10 necrotic damage to the target. In addition, the target must succeed on a DC 15 Constitution saving throw or have disadvantage for 1 hour on any ability check or saving throw that uses Strength or Constitution.

\DndItemHeader{Staff of the Magi}{l, e, g, e, n, d, a, r, y None}
This staff can be wielded as a magic quarterstaff that grants a +2 bonus to attack and damage rolls made with it. While you hold it, you gain a +2 bonus to spell attack rolls.

The staff has 50 charges for the following properties. It regains 4d6 + 2 expended charges daily at dawn. If you expend the last charge, roll a d20. On a 20, the staff regains 1d12 + 1 charges.

\subparagraph{Spell Absorption}
While holding the staff, you have advantage on saving throws against spells. In addition, you can use your reaction when another creature casts a spell that targets only you. If you do, the staff absorbs the magic of the spell, canceling its effect and gaining a number of charges equal to the absorbed spell's level. However, if doing so brings the staff's total number of charges above 50, the staff explodes as if you activated its retributive strike (see below).

\subparagraph{Spells}
While holding the staff, you can use an action to expend some of its charges to cast one of the following spells from it, using your spell save DC and spellcasting ability: \textit{conjure elemental} (7 charges), \textit{dispel magic} (3 charges), \textit{fireball} (7th-level version, 7 charges), \textit{flaming sphere} (2 charges), \textit{ice storm} (4 charges), \textit{invisibility} (2 charges), \textit{knock} (2 charges), \textit{lightning bolt} (7th-level version, 7 charges), \textit{passwall} (5 charges), \textit{plane shift} (7 charges), \textit{telekinesis} (5 charges), \textit{wall of fire} (4 charges), or \textit{web} (2 charges).

You can also use an action to cast one of the following spells from the staff without using any charges: \textit{arcane lock}, \textit{detect magic}, \textit{enlarge/reduce}, \textit{light}, \textit{mage hand}, or \textit{protection from evil and good}.

\subparagraph{Retributive Strike}
You can use an action to break the staff over your knee or against a solid surface, performing a retributive strike. The staff is destroyed and releases its remaining magic in an explosion that expands to fill a 30-foot-radius sphere centered on it.

You have a 50 chance to instantly travel to a random plane of existence, avoiding the explosion. If you fail to avoid the effect, you take force damage equal to 16 × the number of charges in the staff. Every other creature in the area must make a DC 17 Dexterity saving throw. On a failed save, a creature takes an amount of damage based on how far away it is from the point of origin, as shown in the following table. On a successful save, a creature takes half as much damage.


\begin{DndTable}{XX}

Distance from Origin & Damage\\
10 ft. away or closer & 8 × the number of charges in the staff\\
11 to 20 ft. away & 6 × the number of charges in the staff\\
21 to 30 ft. away & 4 × the number of charges in the staff\\
\end{DndTable}

\DndItemHeader{Staff of the Woodlands}{r, a, r, e None}
This staff can be wielded as a magic quarterstaff that grants a +2 bonus to attack and damage rolls made with it. While holding it, you have a +2 bonus to spell attack rolls.

The staff has 10 charges for the following properties. It regains 1d6 + 4 expended charges daily at dawn. If you expend the last charge, roll a d20. On a 1, the staff loses its properties and becomes a nonmagical quarterstaff.

\subparagraph{Spells}
You can use an action to expend 1 or more of the staff's charges to cast one of the following spells from it, using your spell save DC: \textit{animal friendship} (1 charge), \textit{awaken} (5 charges), \textit{barkskin} (2 charges), \textit{locate animals or plants} (2 charges), \textit{speak with animals} (1 charge), \textit{speak with plants} (3 charges), or \textit{wall of thorns} (6 charges).

You can also use an action to cast the \textit{pass without trace} spell from the staff without using any charges.

\subparagraph{Tree Form}
You can use an action to plant one end of the staff in fertile earth and expend 1 charge to transform the staff into a healthy tree. The tree is 60 feet tall and has a 5-foot-diameter trunk, and its branches at the top spread out in a 20-foot radius. The tree appears ordinary but radiates a faint aura of transmutation magic if targeted by \textit{detect magic}. While touching the tree and using another action to speak its command word, you return the staff to its normal form. Any creature in the tree falls when it reverts to a staff.

\DndItemHeader{Sun Blade}{r, a, r, e None}
This item appears to be a longsword hilt. While grasping the hilt, you can use a bonus action to cause a blade of pure radiance to spring into existence, or make the blade disappear. While the blade exists, this magic longsword has the finesse property. If you are proficient with shortsword or longsword, you are proficient with the sun blade.

You gain a +2 bonus to attack and damage rolls made with this weapon, which deals radiant damage instead of slashing damage. When you hit an undead with it, that target takes an extra 1d8 radiant damage.

The sword's luminous blade emits bright light in a 15-foot radius and dim light for an additional 15 feet. The light is sunlight. While the blade persists, you can use an action to expand or reduce its radius of bright and dim light by 5 feet each, to a maximum of 30 feet each or a minimum of 10 feet each.

\DndItemHeader{Tome of Clear Thought}{v, e, r, y,  , r, a, r, e}
This book contains memory and logic exercises, and its words are charged with magic. If you spend 48 hours over a period of 6 days or fewer studying the book's contents and practicing its guidelines, your Intelligence score increases by 2, as does your maximum for that score. The manual then loses its magic, but regains it in a century.

\DndItemHeader{Tome of Leadership and Influence}{v, e, r, y,  , r, a, r, e}
This book contains guidelines for influencing and charming others, and its words are charged with magic. If you spend 48 hours over a period of 6 days or fewer studying the book's contents and practicing its guidelines, your Charisma score increases by 2, as does your maximum for that score. The manual then loses its magic, but regains it in a century.

\DndItemHeader{Tome of Understanding}{v, e, r, y,  , r, a, r, e}
This book contains intuition and insight exercises, and its words are charged with magic. If you spend 48 hours over a period of 6 days or fewer studying the book's contents and practicing its guidelines, your Wisdom score increases by 2, as does your maximum for that score. The manual then loses its magic, but regains it in a century.

\DndItemHeader{Vial of Greed}{r, a, r, e None}
This small glass vial can stockpile resources for use in the future. Once stored, resources last for 1 century before vanishing. As an action you make the vial store any number of the following resources, which are magically consumed and converted into a violet-colored liquid:


\begin{itemize}
\item Up to 31 days of food and/or drink. The flavors are lost, instead becoming tasteless.
\item Up to 7 days of alcohol. The flavors are lost, instead becoming tasteless.
\item Up to 5 magic scrolls that affect a single creature. The target of the spell is you, and if the spell requires concentration, you can concentrate.
\item Up to 5 magic potions. No more than 2 duplicate potions can be stored at a time.
\end{itemize}

You may have more than one kind of resource in the vial, up to the limits expressed above. You can use a bonus action to consume one day of food (or alcohol) or activate one scroll or magic potion. If activating a magic scroll, the effects of that scroll must end before you can activate another scroll from the Vial of Greed.

\DndItemHeader{Wand of Binding}{r, a, r, e None}
This wand has 7 charges for the following properties. It regains 1d6 + 1 expended charges daily at dawn. If you expend the wand's last charge, roll a d20. On a 1, the wand crumbles into ashes and is destroyed.

\subparagraph{Spells}
While holding the wand, you can use an action to expend some of its charges to cast one of the following spells (save DC 17): \textit{hold monster} (5 charges) or \textit{hold person} (2 charges).

\subparagraph{Assisted Escape}
While holding the wand, you can use your reaction to expend 1 charge and gain advantage on a saving throw you make to avoid being paralyzed or restrained, or you can expend 1 charge and gain advantage on any check you make to escape a grapple.

\DndItemHeader{Wand of Fireballs}{r, a, r, e None}
This wand has 7 charges. While holding it, you can use an action to expend 1 or more of its charges to cast the \textit{fireball} spell (save DC 15) from it. For 1 charge, you cast the 3rd-level version of the spell. You can increase the spell slot level by one for each additional charge you expend.

The wand regains 1d6 + 1 expended charges daily at dawn. If you expend the wand's last charge, roll a d20. On a 1, the wand crumbles into ashes and is destroyed.

\DndItemHeader{Wand of Magic Detection}{u, n, c, o, m, m, o, n}
This wand has 3 charges. While holding it, you can expend 1 charge as an action to cast the \textit{detect magic} spell from it. The wand regains 1d3 expended charges daily at dawn.

\DndItemHeader{Wand of Secrets}{u, n, c, o, m, m, o, n}
The wand has 3 charges. While holding it, you can use an action to expend 1 of its charges, and if a secret door or trap is within 30 feet of you, the wand pulses and points at the one nearest to you. The wand regains 1d3 expended charges daily at dawn.

\DndItemHeader{Wand of Wonder}{r, a, r, e None}
This wand has 7 charges. While holding it, you can use an action to expend 1 of its charges and choose a target within 120 feet of you. The target can be a creature, an object, or a point in space. Roll d100 and consult the following table to discover what happens.

If the effect causes you to cast a spell from the wand, the spell's save DC is 15. If the spell normally has a range expressed in feet, its range becomes 120 feet if it isn't already.

If an effect covers an area, you must center the spell on and include the target. If an effect has multiple possible subjects, the DM randomly determines which ones are affected.

The wand regains 1d6 + 1 expended charges daily at dawn. If you expend the wand's last charge, roll a d20. On a 1, the wand crumbles into dust and is destroyed.


\begin{DndTable}{cX}

d100 & Effect\\
01-05 & You cast \textit{slow}.\\
06-10 & You cast \textit{faerie fire}.\\
11-15 & You are stunned until the start of your next turn, believing something awesome just happened.\\
16-20 & You cast \textit{gust of wind}.\\
21-25 & You cast \textit{detect thoughts} on the target you chose. If you didn't target a creature, you instead take 1d6 psychic damage.\\
26-30 & You cast \textit{stinking cloud}.\\
31-33 & Heavy rain falls in a 60-foot radius centered on the target. The area becomes lightly obscured. The rain falls until the start of your next turn.\\
34-36 & An animal appears in the unoccupied space nearest the target. The animal isn't under your control and acts as it normally would. Roll a d100 to determine which animal appears. On a 01-25, a \textbf{rhinoceros} appears; on a 26-50, an \textbf{elephant} appears; and on a 51-100, a \textbf{rat} appears.\\
37-46 & You cast \textit{lightning bolt}.\\
47-49 & A cloud of 600 oversized butterflies fills a 30-foot radius centered on the target. The area becomes heavily obscured. The butterflies remain for 10 minutes.\\
50-53 & You enlarge the target as if you had cast \textit{enlarge/reduce}. If the target can't be affected by that spell or if you didn't target a creature, you become the target.\\
54-58 & You cast \textit{darkness}.\\
59-62 & Grass grows on the ground in a 60-foot radius centered on the target. If grass is already there, it grows to ten times its normal size and remains overgrown for 1 minute.\\
63-65 & An object of the DM's choice disappears into the Ethereal Plane. The object must be neither worn nor carried, within 120 feet of the target, and no larger than 10 feet in any dimension.\\
66-69 & You shrink yourself as if you had cast \textit{enlarge/reduce} on yourself.\\
70-79 & You cast \textit{fireball}.\\
80-84 & You cast \textit{invisibility} on yourself.\\
85-87 & Leaves grow from the target. If you chose a point in space as the target, leaves sprout from the creature nearest to that point. Unless they are picked off, the leaves turn brown and fall off after 24 hours.\\
88-90 & A stream of 1d4 × 10 gems, each worth 1 gp, shoots from the wand's tip in a line 30 feet long and 5 feet wide. Each gem deals 1 bludgeoning damage, and the total damage of the gems is divided equally among all creatures in the line.\\
91-95 & A burst of colorful shimmering light extends from you in a 30-foot radius. You and each creature in the area that can see must succeed on a DC 15 Constitution saving throw or become blinded for 1 minute. A creature can repeat the saving throw at the end of each of its turns, ending the effect on itself on a success.\\
96-97 & The target's skin turns bright blue for 1d10 days. If you chose a point in space, the creature nearest to that point is affected.\\
98-00 & If you targeted a creature, it must make a DC 15 Constitution saving throw. If you didn't target a creature, you become the target and must make the saving throw. If the saving throw fails by 5 or more, the target is instantly petrified. On any other failed save, the target is restrained and begins to turn to stone. While restrained in this way, the target must repeat the saving throw at the end of its next turn, becoming petrified on a failure or ending the effect on a success. The petrification lasts until the target is freed by the \textit{greater restoration} spell or similar magic.\\
\end{DndTable}

\DndItemHeader{Wings of Flying}{r, a, r, e None}
While wearing this cloak, you can use an action to speak its command word. This turns the cloak into a pair of bat wings or bird wings on your back for 1 hour or until you repeat the command word as an action. The wings give you a flying speed of 60 feet. When they disappear, you can't use them again for 1d12 hours.

\DndItemHeader{Wrought-Iron Tower}{a, r, t, i, f, a, c, t None}
Constructed from the same materials as his tower, Dispater's Wrought-Iron Tower is a powerful weapon that marks his rule over Dis. It embodies his ideals, representing defense, authority, and power. Its construction changes hourly, similar to the tower, appearing as either iron, steel, lead, or adamantine.

The Wrought-Iron Tower is a magic staff that functions as a warhammer and you have a +3 bonus to attack and damage rolls made with it. If a creature is vulnerable to a certain type of damage, all damage dealt with the staff is converted to that damage type. When you hit a good aligned creature with it, that creature takes an extra 3d6 fire damage.

While holding the staff, you can control and navigate through the Iron Tower in Dis. While the staff is on the layer of Dis, Dispater knows the exact location of it, and is capable of teleporting to within 30 feet of it at will.

\subparagraph{Random Properties}
The Wrought-Iron Tower has the following random properties:


\begin{itemize}
\item 2 Artifact Properties; Minor Beneficial Properties properties.
\end{itemize}

\subparagraph{The Best Offense}
While attuned to the staff, you have a +2 bonus to AC and a +2 bonus to all saving throws.

Additionally, the staff has 10 charges. When you make an attack, you may spend 1 or more charges. For each charge spent, you deal an additional 2 fire damage on a successful attack. Charges are still consumed if the attack misses. The staff regains all charges each day at dawn.

After dealing damage with the staff, you may use a bonus action and expend 1 charge to replace the damage dealt with your total AC instead.

\subparagraph{Rulership}
The Wrought-Iron Tower also functions as a \textit{Rod of Rulership}, with an increased save DC of 18. This property costs 3 charges to use, though a creature that successfully saves against one of its uses is immune to future uses for 24 hours.

\subparagraph{Slag}
You can cast \textit{Searing Smite} at 6th level on yourself, expending 1 charge. Instead of doing fire damage, you may have Searing Smite inflict acid damage. The spell requires no concentration and lasts for 1 minute.

\subparagraph{Destroying the Staff}
The Wrought-Iron Tower can only be destroyed inside the Iron Tower on Dis. Located somewhere within the tower is the forge that Dispater used to create the staff. An attuned creature must place the staff into the forge fires, then cast \textit{Dispel Evil and Good} and touch the staff. They then take 20d10 fire damage, which ignores resistance and immunities. If the attuned creature survives the fire damage, the staff is destroyed.

\chapter{Creatures}
\setlist[description]{nolistsep,listparindent=3pt,topsep=6pt,font={\bfseries\sffamily\small}}
\pagestyle{empty}

\begin{DndMonster}[float*=b,width=\textwidth + 8pt]{Abigor}
\begin{multicols}{2}
\DndMonsterType{Large fiend (devil), lawful evil}
\DndMonsterBasics[
armor-class = {22 (natural armor)},
hit-points = {\DndDice{27d10 + 189}},
speed = {30 ft.}
]
\DndMonsterAbilityScores[
 str = 27,
 dex = 14,
 con = 25,
 int = 20,
 wis = 16,
 cha = 16
]
\DndMonsterDetails[
saving-throws = {Str +15, Int +12},
skills = {Athletics +22, History +12, Insight +10, Intimidation +17, Perception +10, Survival +10},
damage-resistances = {cold; bludgeoning, piercing, slashing from nonmagical attacks that aren't silvered},
damage-immunities = {fire, poison},
condition-immunities = {poisoned},
senses = {darkvision 120 ft., passive perception 20},
languages = {Celestial, Common, Draconic, Infernal, telepathy 120 ft.},
challenge = 22,
]
\DndMonsterAction{Devil's Sight}
Magical darkness doesn't impede Abigor's darkvision.

\DndMonsterAction{Magic Resistance}
Abigor has advantage on saving throws against spells and other magical effects.

\DndMonsterSection{Actions}
\DndMonsterAction{Multiattack}
Abigor makes four attacks using its Warhammer, Throwing Hammer, or a combination of the two.

\DndMonsterAction{Warhammer}
\textsl{Melee Weapon Attack:} 15 to hit, reach 5 ft., one target. \textsl{Hit:} 17 (2d8 + 8) bludgeoning damage plus 3 (1d6) thunder damage.

\DndMonsterAction{Throwing Hammer}
\textsl{Ranged Weapon Attack:} 15 to hit, range 30/60 ft., one target. \textsl{Hit:} 17 (2d8 + 8) bludgeoning damage plus 3 (1d6) thunder damage. Hammers thrown this way fall and stick into the ground after hitting their target and a new hammer appears in Abigor's hand.

\DndMonsterAction{Thunderous Slam (Recharge 5\textendash 6)}
Abigor causes all its hammers left on the battlefield to bristle with thunderous energy, exploding outwards with a roar. All creatures within a 10-foot-radius sphere centered on each hammer must make a DC 22 Constitution saving throw. Targets take 26 (4d12) thunder damage on a failed save, or half as much damage on a successful one. If a creature is within multiple spheres, it must make the save multiple times, taking damage from each source.

\DndMonsterSection{Legendary Actions}
Abigor can take 3 legendary actions, choosing from the options below. Only one legendary action can be used at a time and only at the end of another creature's turn. Abigor regains spent legendary actions at the start of its turn.

\begin{DndMonsterLegendaryActions}
\DndMonsterLegendaryAction{Throw Hammer}{Abigor makes a Throwing Hammer attack.
}
\DndMonsterLegendaryAction{Call Thunder (Costs 2 Actions)}{Abigor targets a hammer that it can see within 60 feet to let out a thunderous roar. Creatures within a 10-foot-radius sphere centered on the hammer must make a DC 22 Constitution saving throw, taking 26 (4d12) thunder damage on a failed save, or half as much damage on a successful one.
}
\end{DndMonsterLegendaryActions}
\end{multicols}
\end{DndMonster}



\begin{DndMonster}[float*=b,width=\textwidth + 8pt]{Aboleth}
\begin{multicols}{2}
\DndMonsterType{Large aberration, lawful evil}
\DndMonsterBasics[
armor-class = {17 (natural armor)},
hit-points = {\DndDice{18d10 + 36}},
speed = {10 ft., swim 40 ft.}
]
\DndMonsterAbilityScores[
 str = 21,
 dex = 9,
 con = 15,
 int = 18,
 wis = 15,
 cha = 18
]
\DndMonsterDetails[
saving-throws = {Con +6, Int +8, Wis +6},
skills = {History +12, Perception +10},
senses = {darkvision 120 ft., passive perception 20},
languages = {Deep Speech, telepathy 120 ft.},
challenge = 10,
]
\DndMonsterAction{Amphibious}
The aboleth can breathe air and water.

\DndMonsterAction{Mucous Cloud}
While underwater, the aboleth is surrounded by transformative mucus. A creature that touches the aboleth or that hits it with a melee attack while within 5 feet of it must make a DC 14 Constitution saving throw. On a failure, the creature is diseased for 1d4 hours. The diseased creature can breathe only underwater.

\DndMonsterAction{Probing Telepathy}
If a creature communicates telepathically with the aboleth, the aboleth learns the creature's greatest desires if the aboleth can see the creature.

\DndMonsterSection{Actions}
\DndMonsterAction{Multiattack}
The aboleth makes three tentacle attacks.

\DndMonsterAction{Tentacle}
\textsl{Melee Weapon Attack:} 9 to hit, reach 10 ft., one target. \textsl{Hit:} 12 (2d6 + 5) bludgeoning damage. If the target is a creature, it must succeed on a DC 14 Constitution saving throw or become diseased. The disease has no effect for 1 minute and can be removed by any magic that cures disease. After 1 minute, the diseased creature's skin becomes translucent and slimy, the creature can't regain hit points unless it is underwater, and the disease can be removed only by \textit{heal} or another disease-curing spell of 6th level or higher. When the creature is outside a body of water, it takes 6 (1d12) acid damage every 10 minutes unless moisture is applied to the skin before 10 minutes have passed.

\DndMonsterAction{Tail}
\textsl{Melee Weapon Attack:} 9 to hit, reach 10 ft., one target. \textsl{Hit:} 15 (3d6 + 5) bludgeoning damage.

\DndMonsterAction{Enslave (3/Day)}
The aboleth targets one creature it can see within 30 feet of it. The target must succeed on a DC 14 Wisdom saving throw or be magically charmed by the aboleth until the aboleth dies or until it is on a different plane of existence from the target. The charmed target is under the aboleth's control and can't take reactions, and the aboleth and the target can communicate telepathically with each other over any distance.

Whenever the charmed target takes damage, the target can repeat the saving throw. On a success, the effect ends. No more than once every 24 hours, the target can also repeat the saving throw when it is at least 1 mile away from the aboleth.

\DndMonsterSection{Legendary Actions}
Aboleth can take 3 legendary actions, choosing from the options below. Only one legendary action can be used at a time and only at the end of another creature's turn. Aboleth regains spent legendary actions at the start of its turn.

\begin{DndMonsterLegendaryActions}
\DndMonsterLegendaryAction{Detect}{The aboleth makes a Wisdom (Perception) check.
}
\DndMonsterLegendaryAction{Tail Swipe}{The aboleth makes one tail attack.
}
\DndMonsterLegendaryAction{Psychic Drain (Costs 2 Actions)}{One creature charmed by the aboleth takes 10 (3d6) psychic damage, and the aboleth regains hit points equal to the damage the creature takes.
}
\end{DndMonsterLegendaryActions}
\DndMonsterSection{Lair Actions}
When fighting inside its lair, an aboleth can invoke the ambient magic to take lair actions. On initiative count 20 (losing initiative ties), the aboleth takes a lair action to cause one of the following effects:


\begin{itemize}
\item The aboleth casts \textit{phantasmal force} (no components required) on any number of creatures it can see within 60 feet of it. While maintaining concentration on this effect, the aboleth can't take other lair actions. If a target succeeds on the saving throw or if the effect ends for it, the target is immune to the aboleth's phantasmal force lair action for the next 24 hours, although such a creature can choose to be affected.
\item Pools of water within 90 feet of the aboleth surge outward in a grasping tide. Any creature on the ground within 20 feet of such a pool must succeed on a DC 14 Strength saving throw or be pulled up to 20 feet into the water and knocked prone. The aboleth can't use this lair action again until it has used a different one.
\item Water in the aboleth's lair magically becomes a conduit for the creature's rage. The aboleth can target any number of creatures it can see in such water within 90 feet of it. A target must succeed on a DC 14 Wisdom saving throw or take 7 (2d6) psychic damage. The aboleth can't use this lair action again until it has used a different one.
\end{itemize}

\DndMonsterSection{Regional Effects}
The region containing an aboleth's lair is warped by the creature's presence, which creates one or more of the following effects:


\begin{itemize}
\item Underground surfaces within 1 mile of the aboleth's lair are slimy and wet and are difficult terrain.
\item Water sources within 1 mile of the lair are supernaturally fouled. Enemies of the aboleth that drink such water vomit it within minutes.
\item As an action, the aboleth can create an illusory image of itself within 1 mile of the lair. The copy can appear at any location the aboleth has seen before or in any location a creature charmed by the aboleth can currently see. Once created, the image lasts for as long as the aboleth maintains concentration, as if concentration on a spell. Although the image is intangible, it looks, sounds, and can move like the aboleth. The aboleth can sense, speak, and use telepathy from the image's position as if present at that position. If the image takes any damage, it disappears.
\end{itemize}

If the aboleth dies, the first two effects fade over the course of 3d10 days.

\end{multicols}
\end{DndMonster}



\begin{DndMonster}[float=t]{Acolyte}
\DndMonsterType{Medium humanoid (any race), any alignment}
\DndMonsterBasics[
armor-class = {10},
hit-points = {\DndDice{2d8}},
speed = {30 ft.}
]
\DndMonsterAbilityScores[
 str = 10,
 dex = 10,
 con = 10,
 int = 10,
 wis = 14,
 cha = 11
]
\DndMonsterDetails[
skills = {Medicine +4, Religion +2},
languages = {any one language (usually Common)},
challenge = 1/4,
]
\DndMonsterAction{Spellcasting}
The acolyte is a 1st-level spellcaster. Its spellcasting ability is Wisdom (spell save DC 12, 4 to hit with spell attacks). The acolyte has following cleric spells prepared:
\begin{DndMonsterSpells}
  \DndMonsterSpellLevel{\textit{light}, \textit{sacred flame}, \textit{thaumaturgy}}
   \DndMonsterSpellLevel[1][3]{\textit{bless}, \textit{cure wounds}, \textit{sanctuary}}
\end{DndMonsterSpells}
\DndMonsterSection{Actions}
\DndMonsterAction{Club}
\textsl{Melee Weapon Attack:} 2 to hit, reach 5 ft., one target. \textsl{Hit:} 2 (1d4) bludgeoning damage.

\end{DndMonster}



\begin{DndMonster}[float*=b,width=\textwidth + 8pt]{Adult Black Dragon}
\begin{multicols}{2}
\DndMonsterType{Huge dragon, chaotic evil}
\DndMonsterBasics[
armor-class = {19 (natural armor)},
hit-points = {\DndDice{17d12 + 85}},
speed = {40 ft., fly 80 ft., swim 40 ft.}
]
\DndMonsterAbilityScores[
 str = 23,
 dex = 14,
 con = 21,
 int = 14,
 wis = 13,
 cha = 17
]
\DndMonsterDetails[
saving-throws = {Dex +7, Con +10, Wis +6, Cha +8},
skills = {Perception +11, Stealth +7},
damage-immunities = {acid},
senses = {blindsight 60 ft., darkvision 120 ft., passive perception 21},
languages = {Common, Draconic},
challenge = 14,
]
\DndMonsterAction{Amphibious}
The dragon can breathe air and water.

\DndMonsterAction{Legendary Resistance (3/Day)}
If the dragon fails a saving throw, it can choose to succeed instead.

\DndMonsterSection{Actions}
\DndMonsterAction{Multiattack}
The dragon can use its Frightful Presence. It then makes three attacks: one with its bite and two with its claws.

\DndMonsterAction{Bite}
\textsl{Melee Weapon Attack:} 11 to hit, reach 10 ft., one target. \textsl{Hit:} 17 (2d10 + 6) piercing damage plus 4 (1d8) acid damage.

\DndMonsterAction{Claw}
\textsl{Melee Weapon Attack:} 11 to hit, reach 5 ft., one target. \textsl{Hit:} 13 (2d6 + 6) slashing damage.

\DndMonsterAction{Tail}
\textsl{Melee Weapon Attack:} 11 to hit, reach 15 ft., one target. \textsl{Hit:} 15 (2d8 + 6) bludgeoning damage.

\DndMonsterAction{Frightful Presence}
Each creature of the dragon's choice that is within 120 feet of the dragon and aware of it must succeed on a DC 16 Wisdom saving throw or become frightened for 1 minute. A creature can repeat the saving throw at the end of each of its turns, ending the effect on itself on a success. If a creature's saving throw is successful or the effect ends for it, the creature is immune to the dragon's Frightful Presence for the next 24 hours.

\DndMonsterAction{Acid Breath (Recharge 5\textendash 6)}
The dragon exhales acid in a 60-foot line that is 5 feet wide. Each creature in that line must make a DC 18 Dexterity saving throw, taking 54 (12d8) acid damage on a failed save, or half as much damage on a successful one.

\DndMonsterSection{Legendary Actions}
Adult Black Dragon can take 3 legendary actions, choosing from the options below. Only one legendary action can be used at a time and only at the end of another creature's turn. Adult Black Dragon regains spent legendary actions at the start of its turn.

\begin{DndMonsterLegendaryActions}
\DndMonsterLegendaryAction{Detect}{The dragon makes a Wisdom (Perception) check.
}
\DndMonsterLegendaryAction{Tail Attack}{The dragon makes a tail attack.
}
\DndMonsterLegendaryAction{Wing Attack (Costs 2 Actions)}{The dragon beats its wings. Each creature within 10 feet of the dragon must succeed on a DC 19 Dexterity saving throw or take 13 (2d6 + 6) bludgeoning damage and be knocked prone. The dragon can then fly up to half its flying speed.
}
\end{DndMonsterLegendaryActions}
\end{multicols}
\end{DndMonster}



\begin{DndMonster}[float*=b,width=\textwidth + 8pt]{Adult Red Dragon}
\begin{multicols}{2}
\DndMonsterType{Huge dragon, chaotic evil}
\DndMonsterBasics[
armor-class = {19 (natural armor)},
hit-points = {\DndDice{19d12 + 133}},
speed = {40 ft., climb 40 ft., fly 80 ft.}
]
\DndMonsterAbilityScores[
 str = 27,
 dex = 10,
 con = 25,
 int = 16,
 wis = 13,
 cha = 21
]
\DndMonsterDetails[
saving-throws = {Dex +6, Con +13, Wis +7, Cha +11},
skills = {Perception +13, Stealth +6},
damage-immunities = {fire},
senses = {blindsight 60 ft., darkvision 120 ft., passive perception 23},
languages = {Common, Draconic},
challenge = 17,
]
\DndMonsterAction{Legendary Resistance (3/Day)}
If the dragon fails a saving throw, it can choose to succeed instead.

\DndMonsterSection{Actions}
\DndMonsterAction{Multiattack}
The dragon can use its Frightful Presence. It then makes three attacks: one with its bite and two with its claws.

\DndMonsterAction{Bite}
\textsl{Melee Weapon Attack:} 14 to hit, reach 10 ft., one target. \textsl{Hit:} 19 (2d10 + 8) piercing damage plus 7 (2d6) fire damage.

\DndMonsterAction{Claw}
\textsl{Melee Weapon Attack:} 14 to hit, reach 5 ft., one target. \textsl{Hit:} 15 (2d6 + 8) slashing damage.

\DndMonsterAction{Tail}
\textsl{Melee Weapon Attack:} 14 to hit, reach 15 ft., one target. \textsl{Hit:} 17 (2d8 + 8) bludgeoning damage.

\DndMonsterAction{Frightful Presence}
Each creature of the dragon's choice that is within 120 feet of the dragon and aware of it must succeed on a DC 19 Wisdom saving throw or become frightened for 1 minute. A creature can repeat the saving throw at the end of each of its turns, ending the effect on itself on a success. If a creature's saving throw is successful or the effect ends for it, the creature is immune to the dragon's Frightful Presence for the next 24 hours.

\DndMonsterAction{Fire Breath (Recharge 5\textendash 6)}
The dragon exhales fire in a 60-foot cone. Each creature in that area must make a DC 21 Dexterity saving throw, taking 63 (18d6) fire damage on a failed save, or half as much damage on a successful one.

\DndMonsterSection{Legendary Actions}
Adult Red Dragon can take 3 legendary actions, choosing from the options below. Only one legendary action can be used at a time and only at the end of another creature's turn. Adult Red Dragon regains spent legendary actions at the start of its turn.

\begin{DndMonsterLegendaryActions}
\DndMonsterLegendaryAction{Detect}{The dragon makes a Wisdom (Perception) check.
}
\DndMonsterLegendaryAction{Tail Attack}{The dragon makes a tail attack.
}
\DndMonsterLegendaryAction{Wing Attack (Costs 2 Actions)}{The dragon beats its wings. Each creature within 10 feet of the dragon must succeed on a DC 22 Dexterity saving throw or take 15 (2d6 + 8) bludgeoning damage and be knocked prone. The dragon can then fly up to half its flying speed.
}
\end{DndMonsterLegendaryActions}
\end{multicols}
\end{DndMonster}



\begin{DndMonster}[float=t]{Aeshma}
\DndMonsterType{Medium humanoid (tiefling), chaotic neutral}
\DndMonsterBasics[
armor-class = {15 (\textit{bracers of defense})},
hit-points = {\DndDice{7d8 + 14}},
speed = {30 ft.}
]
\DndMonsterAbilityScores[
 str = 10,
 dex = 16,
 con = 14,
 int = 14,
 wis = 14,
 cha = 16
]
\DndMonsterDetails[
saving-throws = {Dex +5, Cha +5},
skills = {Deception +7, Insight +4, Performance +7, Persuasion +7, Sleight Of Hand +7, Stealth +5},
languages = {Common, Infernal},
challenge = 2,
]
\DndMonsterSection{Actions}
\DndMonsterAction{Multiattack}
Aeshma makes two attacks using his Shortsword, Hand Crossbow, or a combination of the two.

\DndMonsterAction{Shortsword}
\textsl{Melee Weapon Attack:} 5 to hit, reach 5 ft., one target. \textsl{Hit:} 6 (1d6 + 3) piercing damage.

\DndMonsterAction{Hand Crossbow}
\textsl{Ranged Weapon Attack:} 5 to hit, range 30/120 ft., one target. \textsl{Hit:} 6 (1d6 + 3) piercing damage.

\DndMonsterSection{Bonus Actions}
\DndMonsterAction{Cunning}
Aeshma takes the Dash, Disengage, or Hide action.

\end{DndMonster}



\begin{DndMonster}[float=t]{Affliction Devil (Kocrachon)}
\DndMonsterType{Medium fiend (devil), lawful evil}
\DndMonsterBasics[
armor-class = {20 (natural armor)},
hit-points = {\DndDice{16d8 + 32}},
speed = {30 ft., fly 60 ft. \textit{(hover)}}
]
\DndMonsterAbilityScores[
 str = 16,
 dex = 19,
 con = 15,
 int = 17,
 wis = 12,
 cha = 14
]
\DndMonsterDetails[
saving-throws = {Dex +8, Int +7},
skills = {Deception +6, Intimidation +6, Perception +5, Stealth +8},
damage-resistances = {bludgeoning, piercing, slashing from nonmagical attacks that aren't silvered},
damage-immunities = {cold, fire, poison},
condition-immunities = {poisoned},
senses = {darkvision 120 ft., passive perception 15},
languages = {Infernal, telepathy 120 ft.},
challenge = 10,
]
\DndMonsterAction{Devil's Sight}
Magical darkness doesn't impede the kocrachon's darkvision.

\DndMonsterAction{Infernal Rot}
A creature afflicted with this disease has the poisoned condition until the disease ends. While diseased, creatures gain no benefits from a long rest, gaining exhaustion when necessary. If a creature dies while afflicted with Infernal Rot, its corpse transforms into a newborn kocrachon after 4 (1d8) days.

\DndMonsterAction{Magic Resistance}
The kocrachon has advantage on saving throws against spells and other magical effects.

\DndMonsterAction{Sadism}
The kocrachon gains a +1 bonus to attack and damage rolls for each different creature it damaged on its previous turn.

\DndMonsterSection{Actions}
\DndMonsterAction{Multiattack}
The kocrachon makes three Claw attacks. It can replace one of the attacks with a Proboscis attack.

\DndMonsterAction{Claw}
\textsl{Melee Weapon Attack:} 8 to hit, reach 5 ft., one target. \textsl{Hit:} 11 (2d6 + 4) slashing damage.

\DndMonsterAction{Proboscis}
\textsl{Melee Weapon Attack:} 8 to hit, reach 10 ft., one target. \textsl{Hit:} 10 (1d12 + 4) piercing damage plus 14 (4d6) poison damage. Creatures damaged by this attack must succeed on a DC 16 Constitution saving throw or suffer the Infernal Rot disease.

\end{DndMonster}



\begin{DndMonster}[float=t]{Alhoon}
\DndMonsterType{Medium undead (mind flayer, wizard), neutral evil}
\DndMonsterBasics[
armor-class = {15 (natural armor)},
hit-points = {\DndDice{20d8 + 60}},
speed = {30 ft., fly 15 ft. \textit{(hover)}}
]
\DndMonsterAbilityScores[
 str = 11,
 dex = 12,
 con = 16,
 int = 19,
 wis = 17,
 cha = 17
]
\DndMonsterDetails[
saving-throws = {Con +7, Int +8, Wis +7, Cha +7},
skills = {Arcana +8, Deception +7, History +8, Insight +7, Perception +7, Stealth +5},
damage-resistances = {cold, lightning, necrotic},
damage-immunities = {poison},
condition-immunities = {charmed, exhaustion, frightened, paralyzed, poisoned},
senses = {truesight 120 ft., passive perception 17},
languages = {Deep Speech, Undercommon, telepathy 120 ft.},
challenge = 10,
]
\DndMonsterAction{Magic Resistance}
The alhoon has advantage on saving throws against spells and other magical effects.

\DndMonsterAction{Turn Resistance}
The alhoon has advantage on saving throws against any effect that turns Undead.

\DndMonsterAction{Spellcasting}
The alhoon casts one of the following spells, requiring no material components and using Intelligence as the spellcasting ability (spell save DC 16):
\begin{DndMonsterSpells}
  \DndInnateSpellLevel{\textit{dancing lights}, \textit{detect magic}, \textit{detect thoughts}, \textit{disguise self}, \textit{mage hand}, \textit{prestidigitation}}
  \DndInnateSpellLevel[1e]{\textit{dominate monster}, \textit{globe of invulnerability}, \textit{invisibility}, \textit{modify memory}, \textit{plane shift} (self only), \textit{wall of force}}
\end{DndMonsterSpells}
\DndMonsterSection{Actions}
\DndMonsterAction{Multiattack}
The alhoon makes two Chilling Grasp or Arcane Bolt attacks.

\DndMonsterAction{Chilling Grasp}
\textsl{Melee Spell Attack:} 8 to hit, reach 5 ft., one target. \textsl{Hit:} 14 (4d6) cold damage, and the alhoon regains 14 hit points.

\DndMonsterAction{Arcane Bolt}
\textsl{Ranged Spell Attack:} 8 to hit, range 120 ft., one target. \textsl{Hit:} 28 (8d6) force damage.

\DndMonsterAction{Mind Blast (Recharge 5\textendash 6)}
The alhoon magically emits psychic energy in a 60-foot cone. Each creature in that area must succeed on a DC 16 Intelligence saving throw or take 22 (4d8 + 4) psychic damage and be stunned for 1 minute. A target can repeat the saving throw at the end of each of its turns, ending the effect on itself on a success.

\DndMonsterSection{Reactions}
\DndMonsterAction{Negate Spell (3/Day)}
The alhoon targets one creature it can see within 60 feet of it that is casting a spell. If the spell is 3rd level or lower, the spell fails, but any spell slots or charges are not wasted.

\end{DndMonster}



\begin{DndMonster}[float=t]{Amnizu}
\DndMonsterType{Medium fiend (devil), lawful evil}
\DndMonsterBasics[
armor-class = {21 (natural armor)},
hit-points = {\DndDice{27d8 + 81}},
speed = {30 ft., fly 40 ft.}
]
\DndMonsterAbilityScores[
 str = 11,
 dex = 13,
 con = 16,
 int = 20,
 wis = 12,
 cha = 18
]
\DndMonsterDetails[
saving-throws = {Dex +7, Con +9, Wis +7, Cha +10},
skills = {Perception +7},
damage-resistances = {cold; bludgeoning, piercing, slashing from nonmagical attacks that aren't silvered},
damage-immunities = {fire, poison},
condition-immunities = {charmed, poisoned},
senses = {darkvision 120 ft., passive perception 17},
languages = {Common, Infernal, telepathy 1,000 ft.},
challenge = 18,
]
\DndMonsterAction{Devil's Sight}
Magical darkness doesn't impede the amnizu's darkvision.

\DndMonsterAction{Magic Resistance}
The amnizu has advantage on saving throws against spells and other magical effects.

\DndMonsterAction{Spellcasting}
The amnizu casts one of the following spells, requiring no material components and using Intelligence as the spellcasting ability (spell save DC 19):
\begin{DndMonsterSpells}
  \DndInnateSpellLevel{\textit{command}}
  \DndInnateSpellLevel[1]{\textit{feeblemind}}
  \DndInnateSpellLevel[3]{\textit{dominate monster}}
\end{DndMonsterSpells}
\DndMonsterSection{Actions}
\DndMonsterAction{Multiattack}
The amnizu uses Blinding Rot or Forgetfulness, if available. It also makes two Taskmaster Whip attacks.

\DndMonsterAction{Taskmaster Whip}
\textsl{Melee Weapon Attack:} 11 to hit, reach 10 ft., one target. \textsl{Hit:} 9 (1d8 + 5) slashing damage plus 16 (3d10) force damage.

\DndMonsterAction{Blinding Rot}
The amnizu targets one or two creatures that it can see within 60 feet of it. Each target must succeed on a DC 19 Wisdom saving throw or take 26 (4d12) necrotic damage and be blinded until the start of the amnizu's next turn.

\DndMonsterAction{Forgetfulness (Recharge 6)}
The amnizu targets one creature it can see within 60 feet of it. That creature must succeed on a DC 18 Intelligence saving throw or take 26 (4d12) psychic damage and become stunned for 1 minute. A stunned creature repeats the saving throw at the end of each of its turns, ending the effect on itself on a success. If the target is stunned for the full minute, it forgets everything it sensed, experienced, and learned during the last 5 hours.

\DndMonsterSection{Reactions}
\DndMonsterAction{Instinctive Charm}
When a creature within 60 feet of the amnizu makes an attack roll against it, and another creature is within the attack's range, the attacker must make a DC 19 Wisdom saving throw. On a failed save, the attacker must target the creature that is closest to it, not including the amnizu or itself. If multiple creatures are closest, the attacker chooses which one to target. If the saving throw is successful, the attacker is immune to the amnizu's Instinctive Charm for 24 hours.

\end{DndMonster}



\begin{DndMonster}[float=t]{Anacreda}
\DndMonsterType{Large fey, neutral evil}
\DndMonsterBasics[
armor-class = {19 (natural armor)},
hit-points = {\DndDice{22d10 + 132}},
speed = {40 ft., fly 40 ft.}
]
\DndMonsterAbilityScores[
 str = 25,
 dex = 16,
 con = 22,
 int = 15,
 wis = 16,
 cha = 19
]
\DndMonsterDetails[
saving-throws = {Con +12, Cha +10},
skills = {Deception +10, Perception +9},
damage-immunities = {cold; fire; bludgeoning, piercing, slashing from nonmagical attacks},
senses = {darkvision 120 ft., truesight 10 ft., passive perception 19},
languages = {Common, Giant, Infernal, Sylvan},
challenge = 17,
]
\DndMonsterSection{Actions}
\DndMonsterAction{Multiattack}
Anacreda makes two Claw attacks. She can replace one of the attacks with a Crushing Hug attack.

\DndMonsterAction{Claw}
\textsl{Melee Weapon Attack:} 13 to hit, reach 5 ft., one target. \textsl{Hit:} 16 (2d8 + 7) slashing damage.

\DndMonsterAction{Crushing Hug}
\textsl{Melee Weapon Attack:} 13 to hit, reach 5 ft., one target. \textsl{Hit:} 51 (8d10 + 7) bludgeoning damage, and if the target is a Large or smaller creature, Anacreda grabs it with both arms. The target has the grappled condition (escape DC 18). Until the grapple ends, the target takes 73 (12d10 + 7) necrotic damage at the start of each of Anacreda's turns. Anacreda can't make attacks while grappling a creature in this way.

\DndMonsterAction{Spinewall}
Anacreda causes a 50-foot-long, 5-foot-wide row of bones to pierce the earth at a point within 120 feet of her. She chooses the direction of the line. The bones rise 10 feet into the air and form a solid wall. A 5-foot section of bone has an AC of 20 and 10 hit points.

\DndMonsterAction{Shattering Shriek (1/Day)}
Anacreda shrieks loud enough to shatter stone and bone alike. Each creature within a 50-foot-radius sphere centered on Anacreda must make a DC 18 Constitution saving throw. A creature takes 55 (10d10) thunder damage on a failed save, or half as much damage on a successful one. Unworked stone, exposed bone, and similar materials in the radius are shattered and turned into difficult terrain.

\end{DndMonster}



\begin{DndMonster}[float=t]{Anagwendol}
\DndMonsterType{Medium celestial, lawful good}
\DndMonsterBasics[
armor-class = {17 (natural armor)},
hit-points = {\DndDice{20d8 + 100}},
speed = {30 ft., fly 90 ft.}
]
\DndMonsterAbilityScores[
 str = 20,
 dex = 18,
 con = 20,
 int = 18,
 wis = 20,
 cha = 22
]
\DndMonsterDetails[
saving-throws = {Int +9, Wis +10, Cha +11},
skills = {Perception +10},
damage-resistances = {radiant; bludgeoning, piercing, slashing from nonmagical attacks},
condition-immunities = {charmed, exhaustion, frightened},
senses = {truesight 120 ft., passive perception 20},
languages = {all, telepathy 120 ft.},
challenge = 16,
]
\DndMonsterAction{Magic Resistance}
Anagwendol has advantage on saving throws against spells and other magical effects.

\DndMonsterSection{Actions}
\DndMonsterAction{Multiattack}
Anagwendol makes two attacks using her Defending Greatsword, Longbow, or a combination of the two.

\DndMonsterAction{Defending Greatsword}
\textsl{Melee Weapon Attack:} 10 to hit, reach 5 ft., one target. \textsl{Hit:} 19 (4d6 + 5) slashing damage plus 22 (5d8) radiant damage.

\DndMonsterAction{Longbow}
\textsl{Ranged Weapon Attack:} 9 to hit, range 150/600 ft., one target. \textsl{Hit:} 8 (1d8 + 4) piercing damage damage plus 22 (5d8) radiant damage.

\DndMonsterAction{Healing Touch (2/Day)}
Anagwendol touches another creature. The target magically regains 30 (6d8 + 3) hit points and is freed from any curse, disease, poison, blindness, or deafness.

\DndMonsterSection{Bonus Actions}
\DndMonsterAction{Defensive Swordwork}
Anagwendol can forgo attacking with her Defending Greatsword on this turn and gain a +3 to her armor class. She can still attack with her Longbow.

\end{DndMonster}



\begin{DndMonster}[float=t]{Animated Furniture}
\DndMonsterType{Large construct, unaligned}
\DndMonsterBasics[
armor-class = {17 (natural armor)},
hit-points = {\DndDice{17d10 + 85}},
speed = {30 ft.}
]
\DndMonsterAbilityScores[
 str = 22,
 dex = 9,
 con = 20,
 int = 3,
 wis = 11,
 cha = 1
]
\DndMonsterDetails[
damage-immunities = {poison; psychic; bludgeoning, piercing, slashing from nonmagical attacks that aren't adamantine},
condition-immunities = {charmed, exhaustion, frightened, paralyzed, petrified, poisoned},
senses = {darkvision 120 ft., passive perception 10},
languages = {understands the languages of its creator but can't speak},
challenge = 10,
]
\DndMonsterAction{Copy Warning}
This content has been copied from another creature, but the data replacement hasn't been run. Check details.

\DndMonsterAction{Immutable Form}
The golem is immune to any spell or effect that would alter its form.

\DndMonsterAction{Magic Resistance}
The golem has advantage on saving throws against spells and other magical effects.

\DndMonsterAction{Magic Weapons}
The golem's weapon attacks are magical.

\DndMonsterSection{Actions}
\DndMonsterAction{Multiattack}
The golem makes two slam attacks.

\DndMonsterAction{Slam}
\textsl{Melee Weapon Attack:} 10 to hit, reach 5 ft., one target. \textsl{Hit:} 19 (3d8 + 6) bludgeoning damage.

\DndMonsterAction{Slow (Recharge 5\textendash 6)}
The golem targets one or more creatures it can see within 10 feet of it. Each target must make a DC 17 Wisdom saving throw against this magic. On a failed save, a target can't use reactions, its speed is halved, and it can't make more than one attack on its turn. In addition, the target can take either an action or a bonus action on its turn, not both. These effects last for 1 minute. A target can repeat the saving throw at the end of each of its turns, ending the effect on itself on a success.

\end{DndMonster}



\begin{DndMonster}[float=t]{Archmage}
\DndMonsterType{Medium humanoid (any race), any alignment}
\DndMonsterBasics[
armor-class = {12 (15 with \textit{mage armor})},
hit-points = {\DndDice{18d8 + 18}},
speed = {30 ft.}
]
\DndMonsterAbilityScores[
 str = 10,
 dex = 14,
 con = 12,
 int = 20,
 wis = 15,
 cha = 16
]
\DndMonsterDetails[
saving-throws = {Int +9, Wis +6},
skills = {Arcana +13, History +13},
damage-resistances = {damage from spells; bludgeoning, piercing, slashing (from stoneskin)},
languages = {any six languages},
challenge = 12,
]
\DndMonsterAction{Magic Resistance}
The archmage has advantage on saving throws against spells and other magical effects.

\DndMonsterAction{Spellcasting}
The archmage is an 18th-level spellcaster. Its spellcasting ability is Intelligence (spell save DC 17, 9 to hit with spell attacks). The archmage can cast \textit{disguise self} and \textit{invisibility} at will and has the following wizard spells prepared:
\begin{DndMonsterSpells}
  \DndInnateSpellLevel{\textit{disguise self}, \textit{invisibility}}
  \DndMonsterSpellLevel{\textit{fire bolt}, \textit{light}, \textit{mage hand}, \textit{prestidigitation}, \textit{shocking grasp}}
   \DndMonsterSpellLevel[1][4]{\textit{detect magic}, \textit{identify}, \textit{mage armor}*, \textit{magic missile}}
   \DndMonsterSpellLevel[2][3]{\textit{detect thoughts}, \textit{mirror image}, \textit{misty step}}
   \DndMonsterSpellLevel[3][3]{\textit{counterspell}, \textit{fly}, \textit{lightning bolt}}
   \DndMonsterSpellLevel[4][3]{\textit{banishment}, \textit{fire shield}, \textit{stoneskin}*}
   \DndMonsterSpellLevel[5][3]{\textit{cone of cold}, \textit{scrying}, \textit{wall of force}}
   \DndMonsterSpellLevel[6][1]{\textit{globe of invulnerability}}
   \DndMonsterSpellLevel[7][1]{\textit{teleport}}
   \DndMonsterSpellLevel[8][1]{\textit{mind blank}*}
\end{DndMonsterSpells}
\DndMonsterSection{Actions}
\DndMonsterAction{Dagger}
\textsl{Melee or Ranged Weapon Attack:} 6 to hit, reach 5 ft. or range 20/60 ft., one target. \textsl{Hit:} 4 (1d4 + 2) piercing damage.

\end{DndMonster}



\begin{DndMonster}[float*=b,width=\textwidth + 8pt]{Asmodeus}
\begin{multicols}{2}
\DndMonsterType{Large fiend (devil), lawful evil}
\DndMonsterBasics[
armor-class = {22 (natural armor)},
hit-points = {\DndDice{50d10 + 450}},
speed = {30 ft., fly 120 ft. \textit{(hover)}}
]
\DndMonsterAbilityScores[
 str = 30,
 dex = 21,
 con = 28,
 int = 25,
 wis = 30,
 cha = 30
]
\DndMonsterDetails[
saving-throws = {Str +19, Con +18, Int +16, Cha +19},
skills = {Arcana +16, Deception +28, Insight +28, Intimidation +19, Perception +19, Persuasion +19, Religion +16},
damage-resistances = {cold},
damage-immunities = {fire; poison; bludgeoning, piercing, slashing from nonmagical attacks that aren't silvered},
condition-immunities = {blinded, charmed, frightened, paralyzed, poisoned, stunned},
senses = {darkvision 120 ft., truesight 60 ft., passive perception 29},
languages = {all, telepathy 120 ft.},
challenge = 30,
]
\DndMonsterAction{Antimagic Aura}
Asmodeus has advantage on saving throws against spells and other magical effects. Spells of level 4 or lower can't damage Asmodeus.

\DndMonsterAction{Fiendish Regeneration}
Asmodeus regains 40 hit points at the start of his turn. If he takes radiant damage, this trait doesn't function at the start of his next turn. Asmodeus dies only if he starts his turn with 0 hit points and doesn't regenerate. If Asmodeus dies, his body regenerates, returning to life 13 (2d12) days later.

\DndMonsterAction{Legendary Resistance (3/Day)}
If Asmodeus fails a saving throw, he can choose to succeed instead.

\DndMonsterAction{Submission Aura}
When a creature moves to be within 120 feet of Asmodeus, they must succeed on a DC 24 Wisdom saving throw or have the frightened condition until they leave the aura. While frightened, creatures kneel before Asmodeus, taking no actions unless he commands otherwise. If Asmodeus deals damage to a kneeling creature, the creature's frightened condition ends. A creature that succeeds the saving throw is immune to this effect for 1 hour.

\DndMonsterAction{Spellcasting}
Asmodeus casts one of the following spells, requiring no material components and using Wisdom as the spellcasting ability (spell save DC 27):
\begin{DndMonsterSpells}
  \DndInnateSpellLevel{\textit{Bane}, \textit{Bestow Curse}, \textit{Command}, \textit{Cure Wounds}, \textit{Darkness}, \textit{Lesser Restoration}, \textit{Remove Curse}}
  \DndInnateSpellLevel[1e]{\textit{Antimagic Field}, \textit{Earthquake}, \textit{Flame Strike}, \textit{Gate}, \textit{Heal}, \textit{Mass Heal}, \textit{Wish}}
\end{DndMonsterSpells}
\DndMonsterSection{Actions}
\DndMonsterAction{Multiattack}
Asmodeus makes four attacks using Ruby Rod, Chilling Gaze, or a combination of the two. He can replace one of the attacks with Hellish Smite.

\DndMonsterAction{Ruby Rod}
\textsl{Melee Weapon Attack:} 21 to hit, reach 5 ft., one target. \textsl{Hit:} 21 (2d8 + 12) necrotic damage, which also reduces the target's hit point maximum by the damage taken. This damage can't reduce a target's hit point maximum below 1. Any effect that removes a disease allows a creature's hit point maximum to return to normal.

\DndMonsterAction{Chilling Gaze}
\textsl{Ranged Spell Attack:} 16 to hit, range 60 ft., one target. \textsl{Hit:} 18 (2d10 + 7) cold damage, and the target has a-1 penalty to all attack rolls they make until the end of their next turn. This penalty can stack.

\DndMonsterAction{Hellish Smite}
Asmodeus targets a creature he can see within 300 feet of him, calling down a bolt of hellish lightning. The target must make a DC 24 Dexterity saving throw, taking 14 (4d6) lightning damage plus 14 (4d6) fire damage on a failed save, or half as much damage on a successful one.

\DndMonsterAction{Infernal Word: Stun (Recharge 4\textendash 6)}
Asmodeus targets a creature he can see within 120 feet of him. The target must make a DC 24 Intelligence saving throw, taking 45 (10d8) force damage on a failed save, or half as much damage on a successful one. A creature that failed the save is stunned for up to 1 minute. The stunned creature can repeat the saving throw at the end of each of its turns, ending the effect on itself on a success, but taking the damage again on a failed save.

\DndMonsterAction{Hell-Lord's Suggestion (Recharge 5\textendash 6)}
Asmodeus targets a creature he can see within 120 feet of him. The target must make a DC 24 Charisma saving throw, taking 44 (8d10) psychic damage on a failed save, or half as much damage on a successful one. A creature that failed the save has the charmed condition for up to 1 minute. While charmed, a creature is controlled by Asmodeus, performing only actions that he suggests. A charmed creature can repeat the saving throw at the end of each of its turns, ending the effect on itself on a success, but taking the damage again on a failed save.

\DndMonsterAction{Fires of the Nine Hells (Recharge 5\textendash 6)}
Asmodeus chooses a point he can see within 500 feet of him, which explodes into a roaring inferno. All creatures within a 60-foot-radius sphere centered on that point must make a DC 24 Dexterity saving throw, taking 70 (20d6) fire damage on a failed save, or half as much damage on a successful one.

\DndMonsterSection{Reactions}
\DndMonsterAction{Ruby Absorption}
As a reaction to being the sole target of a spell, Asmodeus can attempt to absorb the spell into the Ruby Rod.

Asmodeus rolls (1d20 + 2), and if the result is higher than (10 + spell level), the spell is absorbed into the rod and has no effect.

\DndMonsterSection{Legendary Actions}
Asmodeus can take 3 legendary actions, choosing from the options below. Only one legendary action can be used at a time and only at the end of another creature's turn. Asmodeus regains spent legendary actions at the start of its turn.

\begin{DndMonsterLegendaryActions}
\DndMonsterLegendaryAction{Channel Ruby Rod}{The Ruby Rod has 66 charges. Asmodeus may spend its charges to cast one of the following spells from it, with a spell save DC of 20: \textit{Lesser Restoration} (3 charges), \textit{Wall of Force} (7 charges), \textit{Greater Restoration} (12 charges), \textit{Chain Lightning} (15 charges), \textit{Prismatic Spray} (23 charges), \textit{Globe of Invulnerability} (33 charges)
}
\DndMonsterLegendaryAction{Strike (Costs 2 Actions)}{Asmodeus makes a Ruby Rod attack.
}
\DndMonsterLegendaryAction{Call Underling (Costs 3 Actions)}{Asmodeus summons an allied \textbf{pit fiend} in an unoccupied space that he can see.
}
\end{DndMonsterLegendaryActions}
\DndMonsterSection{Regional Effects}
The region containing Asmodeus's lair is corrupted by his presence, which creates the following effect:


\begin{description}
\item[The Devil's Greeting.]
Within 500 miles of his lair, an illusory version of Asmodeus (as if he has cast the \textit{Project Image} spell) appears. Asmodeus attempts to communicate with any characters and decipher their intentions. If they're hostile, Asmodeus pretends to prepare himself for combat, in an effort to make the character's waste their resources.

\end{description}

If Asmodeus dies, the effects fade over the course of 1d10 days.

\end{multicols}
\end{DndMonster}



\begin{DndMonster}[float=t]{Assassin}
\DndMonsterType{Medium humanoid (any race), any non-good alignment}
\DndMonsterBasics[
armor-class = {15 (studded leather armor)},
hit-points = {\DndDice{12d8 + 24}},
speed = {30 ft.}
]
\DndMonsterAbilityScores[
 str = 11,
 dex = 16,
 con = 14,
 int = 13,
 wis = 11,
 cha = 10
]
\DndMonsterDetails[
saving-throws = {Dex +6, Int +4},
skills = {Acrobatics +6, Deception +3, Perception +3, Stealth +9},
damage-resistances = {poison},
languages = {Thieves' cant plus any two languages},
challenge = 8,
]
\DndMonsterAction{Assassinate}
During its first turn, the assassin has advantage on attack rolls against any creature that hasn't taken a turn. Any hit the assassin scores against a Surprise creature is a critical hit.

\DndMonsterAction{Evasion}
If the assassin is subjected to an effect that allows it to make a Dexterity saving throw to take only half damage, the assassin instead takes no damage if it succeeds on the saving throw, and only half damage if it fails.

\DndMonsterAction{Sneak Attack (1/Turn)}
The assassin deals an extra 14 (4d6) damage when it hits a target with a weapon attack and has advantage on the attack roll, or when the target is within 5 feet of an ally of the assassin that isn't incapacitated and the assassin doesn't have disadvantage on the attack roll.

\DndMonsterSection{Actions}
\DndMonsterAction{Multiattack}
The assassin makes two shortsword attacks.

\DndMonsterAction{Shortsword}
\textsl{Melee Weapon Attack:} 6 to hit, reach 5 ft., one target. \textsl{Hit:} 6 (1d6 + 3) piercing damage, and the target must make a DC 15 Constitution saving throw, taking 24 (7d6) poison damage on a failed save, or half as much damage on a successful one.

\DndMonsterAction{Light Crossbow}
\textsl{Ranged Weapon Attack:} 6 to hit, range 80/320 ft., one target. \textsl{Hit:} 7 (1d8 + 3) piercing damage, and the target must make a DC 15 Constitution saving throw, taking 24 (7d6) poison damage on a failed save, or half as much damage on a successful one.

\end{DndMonster}



\begin{DndMonster}[float=t]{Avatar of Baalzebul}
\DndMonsterType{Medium fiend (devil), lawful evil}
\DndMonsterBasics[
armor-class = {14 (natural armor)},
hit-points = {\DndDice{30d8 + 150}},
speed = {20 ft., fly 30 ft.}
]
\DndMonsterAbilityScores[
 str = 21,
 dex = 13,
 con = 20,
 int = 17,
 wis = 18,
 cha = 20
]
\DndMonsterDetails[
saving-throws = {Str +11, Cha +11},
skills = {Deception +11, Insight +10, Intimidation +11, Persuasion +11},
damage-vulnerabilities = {radiant},
damage-resistances = {acid; cold; bludgeoning, piercing, slashing from nonmagical attacks that aren't silvered},
damage-immunities = {fire, poison},
condition-immunities = {charmed, poisoned},
senses = {darkvision 120 ft., passive perception 14},
languages = {all, telepathy 120 ft.},
challenge = 18,
]
\DndMonsterAction{Stench of the Slug}
Any creature that starts its turn within 10 feet of the avatar must succeed on a DC 19 Constitution saving throw or have the poisoned condition until the start of their next turn. On a successful saving throw, the creature is immune to this stench for 1 hour.

\DndMonsterSection{Actions}
\DndMonsterAction{Multiattack}
The avatar makes four Sickening Claw attacks. It can replace one of the attacks with Insect Gorge (if available).

\DndMonsterAction{Sickening Claw}
\textsl{Melee Weapon Attack:} 11 to hit, reach 5 ft., one target. \textsl{Hit:} 9 (1d8 + 5) slashing damage plus 13 (3d8) acid damage. The target's hit point maximum is reduced by the acid damage taken. The reduction lasts until the target finishes a long rest. The target dies if this effect reduces its hit point maximum to 0.

\DndMonsterAction{Insect Gorge (Recharge 4\textendash 6)}
The avatar disgorges a swarm of biting flies at a point it can see within 300 feet of itself. Each creature within a 20-foot-radius sphere centered on that point must make a DC 19 Constitution saving throw. A creature takes 55 (10d10) piercing damage on a failed save, or half as much damage on a successful one. The biting flies persist for 1 minute or until the avatar is slain. A creature must also make this saving throw when it enters the insects' area for the first time on a turn or ends its turn there.

\DndMonsterAction{Caustic Slimepool (1/Day)}
The avatar causes caustic slime to emanate from itself, covering a 15-foot-radius circle. Each creature within the slime must succeed on a DC 19 Dexterity saving throw or have the grappled condition (escape DC 17). A grappled creature takes 21 (6d6) acid damage at the start of its turn.

\end{DndMonster}



\begin{DndMonster}[float*=b,width=\textwidth + 8pt]{Awful Fisher}
\begin{multicols}{2}
\DndMonsterType{Gargantuan monstrosity, unaligned}
\DndMonsterBasics[
armor-class = {17 (natural armor)},
hit-points = {\DndDice{24d20 + 96}},
speed = {30 ft., climb 30 ft.}
]
\DndMonsterAbilityScores[
 str = 23,
 dex = 16,
 con = 18,
 int = 5,
 wis = 10,
 cha = 8
]
\DndMonsterDetails[
saving-throws = {Dex +9, Con +10},
skills = {Perception +6, Stealth +9, Survival +6},
senses = {blindsight 120 ft., darkvision 60 ft., passive perception 16},
challenge = 18,
]
\DndMonsterAction{Flammable Blood}
If the fisher drops below 200 hit points, it gains vulnerability to fire damage.

\DndMonsterAction{Legendary Resistance (3/Day)}
If the fisher fails a saving throw, it can choose to succeed instead.

\DndMonsterAction{Spider Climb}
The fisher can climb difficult surfaces, including upside down on ceilings, without needing to make an ability check.

\DndMonsterSection{Actions}
\DndMonsterAction{Multiattack}
The fisher makes two Claw attacks.

\DndMonsterAction{Claw}
\textsl{Melee Weapon Attack:} 12 to hit, reach 15 ft., one target. \textsl{Hit:} 20 (4d6 + 6) slashing damage.

\DndMonsterAction{Retract Filament}
One Huge or smaller creature grappled by the fisher's Adhesive Filament must make a DC 20 Strength saving throw. On a failed save, the target is pulled into an unoccupied space within 5 feet of the fisher, and the fisher makes one Claw attack against it with advantage. Anyone else who was attached to the filament is released. Until the grapple ends on the target, the fisher can't use Adhesive Filament.

\DndMonsterSection{Bonus Actions}
\DndMonsterAction{Adhesive Filament}
The fisher extends a sticky filament up to 90 feet in a straight line, which adheres to anything that touches it. A Huge or smaller creature the filament adheres to has the grappled condition (escape DC 20), and ability checks made to escape this grapple have disadvantage. The filament can be attacked (AC 16; 25 hit points; immunity to poison and psychic damage). A weapon that fails to sever it becomes stuck to it, requiring an action and a successful DC 20 Strength check to pull free. Destroying the filament deals no damage to the fisher. The filament crumbles away if the fisher takes this bonus action again.

\DndMonsterSection{Legendary Actions}
Awful Fisher can take 3 legendary actions, choosing from the options below. Only one legendary action can be used at a time and only at the end of another creature's turn. Awful Fisher regains spent legendary actions at the start of its turn.

\begin{DndMonsterLegendaryActions}
\DndMonsterLegendaryAction{Slash}{The fisher makes a Claw attack.
}
\DndMonsterLegendaryAction{Spit}{The fisher uses Adhesive Filament.
}
\DndMonsterLegendaryAction{Reel In}{The fisher uses Retract Filament.
}
\end{DndMonsterLegendaryActions}
\end{multicols}
\end{DndMonster}



\begin{DndMonster}[float=t]{Ayperobo Swarm}
\DndMonsterType{Medium fiend (devil), lawful evil}
\DndMonsterBasics[
armor-class = {20 (natural armor)},
hit-points = {\DndDice{20d8 + 40}},
speed = {5 ft., fly 60 ft. \textit{(hover)}}
]
\DndMonsterAbilityScores[
 str = 3,
 dex = 24,
 con = 14,
 int = 8,
 wis = 13,
 cha = 13
]
\DndMonsterDetails[
saving-throws = {Dex +11, Cha +5},
skills = {Intimidation +5, Perception +5, Stealth +11, Survival +5},
damage-resistances = {cold; bludgeoning, piercing, slashing from nonmagical attacks that aren't silvered},
damage-immunities = {fire, poison},
condition-immunities = {charmed, frightened, grappled, paralyzed, petrified, poisoned, prone, restrained, stunned},
senses = {blindsight 30 ft., darkvision 120 ft., passive perception 15},
languages = {Celestial, Draconic, Infernal, telepathy 120 ft.},
challenge = 12,
]
\DndMonsterAction{Devil's Sight}
Magical darkness doesn't impede the ayperobo swarm's darkvision.

\DndMonsterAction{Magic Resistance}
The ayperobo swarm has advantage on saving throws against spells and other magical effects.

\DndMonsterAction{Swarm}
The ayperobo swarm can occupy another creature's space and vice versa, and the swarm can move through any opening large enough for a Tiny creature. The ayperobo swarm can't regain hit points or gain temporary hit points.

\DndMonsterSection{Actions}
\DndMonsterAction{Multiattack}
The ayperobo swarm makes three Bite attacks.

\DndMonsterAction{Bite}
\textsl{Melee Weapon Attack:} 11 to hit, reach 0 ft., one target. \textsl{Hit:} 27 (8d4 + 7) piercing damage, or 17 (4d4 + 7) piercing damage if the swarm has half of its hit points (65) or fewer.

\DndMonsterAction{Burrow (Recharge 6)}
The ayperobo swarm makes a Bite attack. On a hit, the swarm burrows inside the creature, dealing an additional 35 (10d6) necrotic damage and taking control of the creature. At the start of each of its turns, the target may make a DC 17 Constitution saving throw, expelling the ayperobo swarm and taking 14 (4d6) piercing damage on a success.

While burrowed inside a creature, the ayperobo swarm has total cover against attacks targeting them, and no longer acts on their own initiative. Instead, they take total control over their host, gaining all of its abilities, attacks, and equipment.

They act during the host's turn, taking actions based on their host. If a host dies while the ayperobo swarm is burrowed, they leave the host at the end of the host's turn, rolling initiative if necessary. Creatures acting as hosts to the ayperobo swarm are fully aware of their actions and surroundings, but are incapable of physically operating their body.

\end{DndMonster}



\begin{DndMonster}[float=t]{Azer}
\DndMonsterType{Medium elemental, lawful neutral}
\DndMonsterBasics[
armor-class = {17 (natural armor)},
hit-points = {\DndDice{6d8 + 12}},
speed = {30 ft.}
]
\DndMonsterAbilityScores[
 str = 17,
 dex = 12,
 con = 15,
 int = 12,
 wis = 13,
 cha = 10
]
\DndMonsterDetails[
saving-throws = {Con +4},
damage-immunities = {fire, poison},
condition-immunities = {poisoned},
languages = {Ignan},
challenge = 2,
]
\DndMonsterAction{Heated Body}
A creature that touches the azer or hits it with a melee attack while within 5 feet of it takes 5 (1d10) fire damage.

\DndMonsterAction{Heated Weapons}
When the azer hits with a metal melee weapon, it deals an extra 3 (1d6) fire damage (included in the attack).

\DndMonsterAction{Illumination}
The azer sheds bright light in a 10-foot radius and dim light for an additional 10 feet.

\DndMonsterSection{Actions}
\DndMonsterAction{Warhammer}
\textsl{Melee Weapon Attack:} 5 to hit, reach 5 ft., one target. \textsl{Hit:} 7 (1d8 + 3) bludgeoning damage, or 8 (1d10 + 3) bludgeoning damage if used with two hands to make a melee attack, plus 3 (1d6) fire damage.

\end{DndMonster}



\begin{DndMonster}[float*=b,width=\textwidth + 8pt]{Baalzebul}
\begin{multicols}{2}
\DndMonsterType{Huge fiend (devil), lawful evil}
\DndMonsterBasics[
armor-class = {19 (natural armor)},
hit-points = {\DndDice{40d12 + 280}},
speed = {30 ft., fly 60 ft.}
]
\DndMonsterAbilityScores[
 str = 28,
 dex = 15,
 con = 25,
 int = 21,
 wis = 24,
 cha = 26
]
\DndMonsterDetails[
saving-throws = {Str +17, Con +15, Int +13, Cha +16},
skills = {Athletics +17, Deception +24, Insight +15, Intimidation +16, Persuasion +16},
damage-resistances = {acid; cold; bludgeoning, piercing, slashing from nonmagical attacks that aren't silvered},
damage-immunities = {fire, poison},
condition-immunities = {charmed, exhaustion, frightened, poisoned},
senses = {darkvision 120 ft., tremorsense 10 ft., passive perception 17},
languages = {all, telepathy 120 ft.},
challenge = 26,
]
\DndMonsterAction{Devil's Sight}
Magical darkness doesn't impede Baalzebul's darkvision.

\DndMonsterAction{Fiendish Regeneration}
Baalzebul regains 20 hit points at the start of his turn. If he takes radiant damage this trait doesn't function at the start of his next turn. Baalzebul dies only if he starts his turn with 0 hit points and doesn't regenerate.

\DndMonsterAction{Lord of Flies}
Insects don't attack Baalzebul, and he can issue orders to them.

\DndMonsterAction{Legendary Resistance (3/Day)}
If Baalzebul fails a saving throw, he can choose to succeed instead.

\DndMonsterAction{Magic Resistance}
Baalzebul has advantage on saving throws against spells and other magical effects.

\DndMonsterAction{Stench of the Slug (Slug Form Only)}
A creature that starts its turn within 10 feet of Baalzebul must succeed on a DC 21 Constitution saving throw or have the poisoned condition until the start of its next turn. On a successful save, a creature is immune to this stench for 1 hour.

\DndMonsterAction{Spellcasting (Winged Form Only)}
Baalzebul casts one of the following spells, requiring no material components and using Charisma as the spellcasting ability (spell save DC 24):
\begin{DndMonsterSpells}
  \DndInnateSpellLevel{\textit{Animate Dead}, \textit{Darkness}, \textit{Detect Evil and Good}, \textit{Detect Magic}, \textit{Dispel Evil and Good}, \textit{Hallow}, \textit{Hold Monster}, \textit{Suggestion}, \textit{Teleport}, \textit{True Seeing}}
  \DndInnateSpellLevel[1e]{\textit{Symbol} (pain or insanity), \textit{Wish}}
\end{DndMonsterSpells}
\DndMonsterSection{Actions}
\DndMonsterAction{Multiattack}
Baalzebul uses Heart of Pestilence (if available), then makes three Slam attacks.

\DndMonsterAction{Slam}
\textsl{Melee Weapon Attack:} 17 to hit, reach 10 ft., one target. \textsl{Hit:} 19 (3d6 + 9) bludgeoning damage plus 7 (2d6) necrotic damage.

\DndMonsterAction{Pungent Eruption (Slug Form Only, Recharge 4-6)}
Baalzebul causes refuse and caustic liquid to spill from the ground at a point he can see within 300 feet. Each creature in a 15-foot-radius sphere centered on that point must make a DC 21 Constitution saving throw. Targets take 36 (8d8) acid damage on a failed save, or half as much on a successful one. Creatures that fail the save have the poisoned condition for 1 minute.

\DndMonsterAction{Heart of Pestilence (Recharge 5\textendash 6)}
Baalzebul channels pestilence into his body, sickening creatures within 50 feet that can see him. Sickened creatures must make a DC 21 Charisma saving throw, followed by a DC 21 Constitution saving throw. Failing the Charisma save makes a creature have the frightened condition, while failing the Constitution save makes a creature weak. While weakened, a creature deals half damage on melee attacks. Both effects last for 1 minute or until a \textit{Greater Restoration} spell or similar is cast.

\DndMonsterSection{Bonus Actions}
\DndMonsterAction{Displace}
Baalzebul magically teleports, along with any equipment he is wearing and carrying, up to 120 feet to an unoccupied space he can see.

\DndMonsterSection{Legendary Actions}
Baalzebul can take 3 legendary actions, choosing from the options below. Only one legendary action can be used at a time and only at the end of another creature's turn. Baalzebul regains spent legendary actions at the start of its turn.

\begin{DndMonsterLegendaryActions}
\DndMonsterLegendaryAction{Pummel}{Baalzebul makes a Slam attack.
}
\DndMonsterLegendaryAction{Teleport}{Baalzebul uses Displace.
}
\DndMonsterLegendaryAction{Insect Gorge (Costs 2 Actions)}{Baalzebul disgorges a swarm of biting flies at a point he can see within 300 feet of himself. Each creature within a 20-foot-radius sphere centered on that point must make a DC 21 Constitution saving throw. A creature takes 44 (8d10) piercing damage on a failed save, or half as much damage on a successful one. The biting flies persist for 1 minute, or until Baalzebul uses this ability again. Creatures that enter the flies' area or end their turn inside it must repeat the saving throw.
}
\DndMonsterLegendaryAction{Call Underling (Costs 3 Actions)}{Baalzebul summons an allied \textbf{bone devil} in an unoccupied space that he can see.
}
\end{DndMonsterLegendaryActions}
\DndMonsterSection{Regional Effects}
The region containing Baalzebul's lair is corrupted by his presence, which creates the following effect:


\begin{description}
\item[Filth.]
Within 1 mile of the lair, stinking piles of filth accumulate as if from nowhere. The entire area is covered with garbage and feces and is treated as difficult terrain.

\end{description}

If Baalzebul dies, the effects fade over the course of 1d10 days.

\end{multicols}
\end{DndMonster}


\clearpage

\begin{DndMonster}[float=t]{Balor}
\DndMonsterType{Huge fiend (demon), chaotic evil}
\DndMonsterBasics[
armor-class = {19 (natural armor)},
hit-points = {\DndDice{21d12 + 126}},
speed = {40 ft., fly 80 ft.}
]
\DndMonsterAbilityScores[
 str = 26,
 dex = 15,
 con = 22,
 int = 20,
 wis = 16,
 cha = 22
]
\DndMonsterDetails[
saving-throws = {Str +14, Con +12, Wis +9, Cha +12},
damage-resistances = {cold; lightning; bludgeoning, piercing, slashing from nonmagical attacks},
damage-immunities = {fire, poison},
condition-immunities = {poisoned},
senses = {truesight 120 ft., passive perception 13},
languages = {Abyssal, telepathy 120 ft.},
challenge = 19,
]
\DndMonsterAction{Death Throes}
When the balor dies, it explodes, and each creature within 30 feet of it must make a DC 20 Dexterity saving throw, taking 70 (20d6) fire damage on a failed save, or half as much damage on a successful one. The explosion ignites flammable objects in that area that aren't being worn or carried, and it destroys the balor's weapons.

\DndMonsterAction{Fire Aura}
At the start of each of the balor's turns, each creature within 5 feet of it takes 10 (3d6) fire damage, and flammable objects in the aura that aren't being worn or carried ignite. A creature that touches the balor or hits it with a melee attack while within 5 feet of it takes 10 (3d6) fire damage.

\DndMonsterAction{Magic Resistance}
The balor has advantage on saving throws against spells and other magical effects.

\DndMonsterAction{Magic Weapons}
The balor's weapon attacks are magical.

\DndMonsterSection{Actions}
\DndMonsterAction{Multiattack}
The balor makes two attacks: one with its longsword and one with its whip.

\DndMonsterAction{Longsword}
\textsl{Melee Weapon Attack:} 14 to hit, reach 10 ft., one target. \textsl{Hit:} 21 (3d8 + 8) slashing damage plus 13 (3d8) lightning damage. If the balor scores a critical hit, it rolls damage dice three times, instead of twice.

\DndMonsterAction{Whip}
\textsl{Melee Weapon Attack:} 14 to hit, reach 30 ft., one target. \textsl{Hit:} 15 (2d6 + 8) slashing damage plus 10 (3d6) fire damage, and the target must succeed on a DC 20 Strength saving throw or be pulled up to 25 feet toward the balor.

\DndMonsterAction{Teleport}
The balor magically teleports, along with any equipment it is wearing or carrying, up to 120 feet to an unoccupied space it can see.

\end{DndMonster}



\begin{DndMonster}[float=t]{Bandit}
\DndMonsterType{Medium humanoid (any race), any non-lawful alignment}
\DndMonsterBasics[
armor-class = {12 (\textit{leather armor})},
hit-points = {\DndDice{2d8 + 2}},
speed = {30 ft.}
]
\DndMonsterAbilityScores[
 str = 11,
 dex = 12,
 con = 12,
 int = 10,
 wis = 10,
 cha = 10
]
\DndMonsterDetails[
languages = {any one language (usually Common)},
challenge = 1/8,
]
\DndMonsterSection{Actions}
\DndMonsterAction{Scimitar}
\textsl{Melee Weapon Attack:} 3 to hit, reach 5 ft., one target. \textsl{Hit:} 4 (1d6 + 1) slashing damage.

\DndMonsterAction{Light Crossbow}
\textsl{Ranged Weapon Attack:} 3 to hit, range 80/320 ft., one target. \textsl{Hit:} 5 (1d8 + 1) piercing damage.

\end{DndMonster}



\begin{DndMonster}[float=t]{Bandit Captain}
\DndMonsterType{Medium humanoid (any race), any non-lawful alignment}
\DndMonsterBasics[
armor-class = {15 (studded leather armor)},
hit-points = {\DndDice{10d8 + 20}},
speed = {30 ft.}
]
\DndMonsterAbilityScores[
 str = 15,
 dex = 16,
 con = 14,
 int = 14,
 wis = 11,
 cha = 14
]
\DndMonsterDetails[
saving-throws = {Str +4, Dex +5, Wis +2},
skills = {Athletics +4, Deception +4},
languages = {any two languages},
challenge = 2,
]
\DndMonsterSection{Actions}
\DndMonsterAction{Multiattack}
The captain makes three melee attacks: two with its scimitar and one with its dagger. Or the captain makes two ranged attacks with its daggers.

\DndMonsterAction{Scimitar}
\textsl{Melee Weapon Attack:} 5 to hit, reach 5 ft., one target. \textsl{Hit:} 6 (1d6 + 3) slashing damage.

\DndMonsterAction{Dagger}
\textsl{Melee or Ranged Weapon Attack:} 5 to hit, reach 5 ft. or range 20/60 ft., one target. \textsl{Hit:} 5 (1d4 + 3) piercing damage.

\DndMonsterSection{Reactions}
\DndMonsterAction{Parry}
The captain adds 2 to its AC against one melee attack that would hit it. To do so, the captain must see the attacker and be wielding a melee weapon.

\end{DndMonster}



\begin{DndMonster}[float*=b,width=\textwidth + 8pt]{Baphomet}
\begin{multicols}{2}
\DndMonsterType{Huge fiend (demon), chaotic evil}
\DndMonsterBasics[
armor-class = {22 (natural armor)},
hit-points = {\DndDice{22d12 + 176}},
speed = {40 ft.}
]
\DndMonsterAbilityScores[
 str = 30,
 dex = 14,
 con = 26,
 int = 18,
 wis = 24,
 cha = 16
]
\DndMonsterDetails[
saving-throws = {Dex +9, Con +15, Wis +14},
skills = {Intimidation +17, Perception +14},
damage-resistances = {cold, fire, lightning},
damage-immunities = {poison; bludgeoning, piercing, slashing that is nonmagical},
condition-immunities = {charmed, exhaustion, frightened, poisoned},
senses = {truesight 120 ft., passive perception 24},
languages = {all, telepathy 120 ft.},
challenge = 23,
]
\DndMonsterAction{Labyrinthine Recall}
Baphomet can perfectly recall any path he has traveled, and he is immune to the \textit{maze} spell.

\DndMonsterAction{Legendary Resistance (3/Day)}
If Baphomet fails a saving throw, he can choose to succeed instead.

\DndMonsterAction{Magic Resistance}
Baphomet has advantage on saving throws against spells and other magical effects.

\DndMonsterAction{Spellcasting}
Baphomet casts one of the following spells, requiring no material components and using Charisma as the spellcasting ability (spell save DC 18):
\begin{DndMonsterSpells}
  \DndInnateSpellLevel[1]{\textit{teleport}}
  \DndInnateSpellLevel[3e]{\textit{dispel magic}, \textit{dominate beast}, \textit{maze}, \textit{wall of stone}}
\end{DndMonsterSpells}
\DndMonsterSection{Actions}
\DndMonsterAction{Multiattack}
Baphomet makes one Bite attack, one Gore attack, and one Heartcleaver attack. He also uses Frightful Presence.

\DndMonsterAction{Bite}
\textsl{Melee Weapon Attack:} 17 to hit, reach 10 ft., one target. \textsl{Hit:} 19 (2d8 + 10) piercing damage.

\DndMonsterAction{Gore}
\textsl{Melee Weapon Attack:} 17 to hit, reach 10 ft., one target. \textsl{Hit:} 17 (2d6 + 10) piercing damage. If Baphomet moved at least 10 feet straight toward the target immediately before the hit, the target takes an extra 16 (3d10) piercing damage. If the target is a creature, it must succeed on a DC 25 Strength saving throw or be pushed up to 10 feet away and knocked prone.

\DndMonsterAction{Heartcleaver}
\textsl{Melee Weapon Attack:} 17 to hit, reach 15 ft., one target. \textsl{Hit:} 21 (2d10 + 10) force damage.

\DndMonsterAction{Frightful Presence}
Each creature of Baphomet's choice within 120 feet of him and aware of him must succeed on a DC 18 Wisdom saving throw or become frightened for 1 minute. A frightened creature can repeat the saving throw at the end of each of its turns, ending the effect on itself on a success. These later saves have disadvantage if Baphomet is within line of sight of the creature.

If a creature succeeds on any of these saves or the effect ends on it, the creature is immune to Baphomet's Frightful Presence for the next 24 hours.

\DndMonsterSection{Legendary Actions}
Baphomet can take 3 legendary actions, choosing from the options below. Only one legendary action can be used at a time and only at the end of another creature's turn. Baphomet regains spent legendary actions at the start of its turn.

\begin{DndMonsterLegendaryActions}
\DndMonsterLegendaryAction{Heartcleaver Attack}{Baphomet makes one Heartcleaver attack.
}
\DndMonsterLegendaryAction{Charge (Costs 2 Actions)}{Baphomet moves up to his speed without provoking opportunity attack, then makes a Gore attack.
}
\end{DndMonsterLegendaryActions}
\DndMonsterSection{Lair Actions}
On initiative count 20 (losing initiative ties), Baphomet can take one of the following lair actions; he can't take the same lair action two rounds in a row:


\begin{description}
\item[Illusory Room.]
Baphomet casts mirage arcane, affecting a room within the lair that is no larger in any dimension than 100 feet. The effect ends on the next initiative count 20. Charisma is Baphomet's spellcasting ability for this spell.

\item[Reverse Gravity.]
Baphomet chooses a room within the lair that is no larger in any dimension than 100 feet. Until the next initiative count 20, gravity is reversed within that room. Any creatures or objects in the room when this happens fall in the direction of the new pull of gravity, unless they have some means of remaining aloft. Baphomet can ignore the gravity reversal if he's in the room, although he likes to use this action to land on a ceiling to attack targets flying near it.

\item[Seal the Way.]
Baphomet seals one doorway or other entryway within the lair. The opening must be unoccupied. It is filled with solid stone for 1 minute or until Baphomet takes this lair action again.

\end{description}

\DndMonsterSection{Regional Effects}
The region containing Baphomet's lair is warped by his magic, creating one or more of the following effects:


\begin{description}
\item[Beguiling Realm.]
Within 6 miles of the lair, all Charisma (Persuasion) and Charisma (Performance) checks have disadvantage, and all Charisma (Deception) and Charisma (Intimidation) checks have advantage.

\item[Hedge Mazes.]
Plant life within 1 mile of the lair grows thick and forms walls of trees, hedges, and other flora in the form of small mazes.

\item[Panicked Beasts.]
Beasts within 1 mile of the lair become frightened and disoriented, as though constantly under threat of being hunted, and might lash out or panic even when no visible threat is nearby.

\end{description}

If Baphomet dies, these effects fade over the course of 1d10 days.

\end{multicols}
\end{DndMonster}



\begin{DndMonster}[float=t]{Barachiel}
\DndMonsterType{Medium celestial, lawful good}
\DndMonsterBasics[
armor-class = {18 (plate armor)},
hit-points = {\DndDice{16d8 + 64}},
speed = {30 ft., fly 50 ft.}
]
\DndMonsterAbilityScores[
 str = 18,
 dex = 12,
 con = 18,
 int = 11,
 wis = 14,
 cha = 20
]
\DndMonsterDetails[
saving-throws = {Con +9, Wis +7},
damage-resistances = {necrotic, poison},
condition-immunities = {exhaustion, frightened, poisoned},
languages = {Celestial, Common, Infernal},
challenge = 13,
]
\DndMonsterSection{Actions}
\DndMonsterAction{Multiattack}
Barachiel makes three Greatsword attacks.

\DndMonsterAction{Greatsword}
\textsl{Melee Weapon Attack:} 9 to hit, reach 5 ft., one target. \textsl{Hit:} 11 (2d6 + 4) slashing damage plus 22 (5d8) radiant damage.

\DndMonsterAction{Purifying Flames (Recharge 4\textendash 6)}
Barachiel wreathes his sword in white flame, then makes a Greatsword attack. On a hit, the target takes an additional 45 (7d12) fire damage plus 45 (7d12) radiant damage. The target must succeed on a DC 18 Charisma saving throw or have the stunned condition until the end of their next turn.

\DndMonsterSection{Reactions}
\DndMonsterAction{Parry}
Barachiel adds 5 to his AC against one attack. To do so, Barachiel must see the attacker and be wielding a melee weapon.

\end{DndMonster}



\begin{DndMonster}[float=t]{Barbed Devil}
\DndMonsterType{Medium fiend (devil), lawful evil}
\DndMonsterBasics[
armor-class = {15 (natural armor)},
hit-points = {\DndDice{13d8 + 52}},
speed = {30 ft.}
]
\DndMonsterAbilityScores[
 str = 16,
 dex = 17,
 con = 18,
 int = 12,
 wis = 14,
 cha = 14
]
\DndMonsterDetails[
saving-throws = {Str +6, Con +7, Wis +5, Cha +5},
skills = {Deception +5, Insight +5, Perception +8},
damage-resistances = {cold; bludgeoning, piercing, slashing from nonmagical attacks that aren't silvered},
damage-immunities = {fire, poison},
condition-immunities = {poisoned},
senses = {darkvision 120 ft., passive perception 18},
languages = {Infernal, telepathy 120 ft.},
challenge = 5,
]
\DndMonsterAction{Barbed Hide}
At the start of each of its turns, the barbed devil deals 5 (1d10) piercing damage to any creature grappling it.

\DndMonsterAction{Devil's Sight}
Magical darkness doesn't impede the devil's darkvision.

\DndMonsterAction{Magic Resistance}
The devil has advantage on saving throws against spells and other magical effects.

\DndMonsterSection{Actions}
\DndMonsterAction{Multiattack}
The devil makes three melee attacks: one with its tail and two with its claws. Alternatively, it can use Hurl Flame twice.

\DndMonsterAction{Claw}
\textsl{Melee Weapon Attack:} 6 to hit, reach 5 ft., one target. \textsl{Hit:} 6 (1d6 + 3) piercing damage.

\DndMonsterAction{Tail}
\textsl{Melee Weapon Attack:} 6 to hit, reach 5 ft., one target. \textsl{Hit:} 10 (2d6 + 3) piercing damage.

\DndMonsterAction{Hurl Flame}
\textsl{Ranged Spell Attack:} 5 to hit, range 150 ft., one target. \textsl{Hit:} 10 (3d6) fire damage. If the target is a flammable object that isn't being worn or carried, it also catches fire.

\end{DndMonster}



\begin{DndMonster}[float=t]{Barlgura}
\DndMonsterType{Large fiend (demon), chaotic evil}
\DndMonsterBasics[
armor-class = {15 (natural armor)},
hit-points = {\DndDice{8d10 + 24}},
speed = {40 ft., climb 40 ft.}
]
\DndMonsterAbilityScores[
 str = 18,
 dex = 15,
 con = 16,
 int = 7,
 wis = 14,
 cha = 9
]
\DndMonsterDetails[
saving-throws = {Dex +5, Con +6},
skills = {Perception +5, Stealth +5},
damage-resistances = {cold, fire, lightning},
damage-immunities = {poison},
condition-immunities = {poisoned},
senses = {blindsight 30 ft., darkvision 120 ft., passive perception 15},
languages = {Abyssal, telepathy 120 ft.},
challenge = 5,
]
\DndMonsterAction{Reckless}
At the start of its turn, the barlgura can gain advantage on all melee weapon attack rolls it makes during that turn, but attack rolls against it have advantage until the start of its next turn.

\DndMonsterAction{Running Leap}
The barlgura's long jump is up to 40 feet and its high jump is up to 20 feet when it has a running start.

\DndMonsterAction{Innate Spellcasting}
The barlgura's spellcasting ability is Wisdom (spell save DC 13). The barlgura can innately cast the following spells, requiring no material components:
\begin{DndMonsterSpells}
  \DndInnateSpellLevel[1e]{\textit{entangle}, \textit{phantasmal force}}
  \DndInnateSpellLevel[2e]{\textit{disguise self}, \textit{invisibility} (self only)}
\end{DndMonsterSpells}
\DndMonsterSection{Actions}
\DndMonsterAction{Multiattack}
The barlgura makes three attacks: one with its bite and two with its fists.

\DndMonsterAction{Bite}
\textsl{Melee Weapon Attack:} 7 to hit, reach 5 ft., one target. \textsl{Hit:} 11 (2d6 + 4) piercing damage.

\DndMonsterAction{Fist}
\textsl{Melee Weapon Attack:} 7 to hit, reach 5 ft., one target. \textsl{Hit:} 9 (1d10 + 4) bludgeoning damage.

\end{DndMonster}



\begin{DndMonster}[float=t]{Bearded Devil}
\DndMonsterType{Medium fiend (devil), lawful evil}
\DndMonsterBasics[
armor-class = {13 (natural armor)},
hit-points = {\DndDice{8d8 + 16}},
speed = {30 ft.}
]
\DndMonsterAbilityScores[
 str = 16,
 dex = 15,
 con = 15,
 int = 9,
 wis = 11,
 cha = 11
]
\DndMonsterDetails[
saving-throws = {Str +5, Con +4, Wis +2},
damage-resistances = {cold; bludgeoning, piercing, slashing from nonmagical attacks that aren't silvered},
damage-immunities = {fire, poison},
condition-immunities = {poisoned},
senses = {darkvision 120 ft., passive perception 10},
languages = {Infernal, telepathy 120 ft.},
challenge = 3,
]
\DndMonsterAction{Devil's Sight}
Magical darkness doesn't impede the devil's darkvision.

\DndMonsterAction{Magic Resistance}
The devil has advantage on saving throws against spells and other magical effects.

\DndMonsterAction{Steadfast}
The devil can't be frightened while it can see an allied creature within 30 feet of it.

\DndMonsterSection{Actions}
\DndMonsterAction{Multiattack}
The devil makes two attacks: one with its beard and one with its glaive.

\DndMonsterAction{Beard}
\textsl{Melee Weapon Attack:} 5 to hit, reach 5 ft., one creature. \textsl{Hit:} 6 (1d8 + 2) piercing damage, and the target must succeed on a DC 12 Constitution saving throw or be poisoned for 1 minute. While poisoned in this way, the target can't regain hit points. The target can repeat the saving throw at the end of each of its turns, ending the effect on itself on a success.

\DndMonsterAction{Glaive}
\textsl{Melee Weapon Attack:} 5 to hit, reach 10 ft., one target. \textsl{Hit:} 8 (1d10 + 3) slashing damage. If the target is a creature other than an undead or a construct, it must succeed on a DC 12 Constitution saving throw or lose 5 (1d10) hit points at the start of each of its turns due to an infernal wound. Each time the devil hits the wounded target with this attack, the damage dealt by the wound increases by 5 (1d10). Any creature can take an action to stanch the wound with a successful DC 12 Wisdom (Medicine) check. The wound also closes if the target receives magical healing.

\end{DndMonster}



\begin{DndMonster}[float*=b,width=\textwidth + 8pt]{Beholder}
\begin{multicols}{2}
\DndMonsterType{Large aberration, lawful evil}
\DndMonsterBasics[
armor-class = {18 (natural armor)},
hit-points = {\DndDice{19d10 + 76}},
speed = {0 ft., fly 20 ft. \textit{(hover)}}
]
\DndMonsterAbilityScores[
 str = 10,
 dex = 14,
 con = 18,
 int = 17,
 wis = 15,
 cha = 17
]
\DndMonsterDetails[
saving-throws = {Int +8, Wis +7, Cha +8},
skills = {Perception +12},
condition-immunities = {prone},
senses = {darkvision 120 ft., passive perception 22},
languages = {Deep Speech, Undercommon},
challenge = 13,
]
\DndMonsterAction{Antimagic Cone}
The beholder's central eye creates an area of antimagic, as in the \textit{antimagic field} spell, in a 150-foot cone. At the start of each of its turns, the beholder decides which way the cone faces and whether the cone is active. The area works against the beholder's own eye rays.

\DndMonsterSection{Actions}
\DndMonsterAction{Bite}
\textsl{Melee Weapon Attack:} 5 to hit, reach 5 ft., one target. \textsl{Hit:} 14 (4d6) piercing damage.

\DndMonsterAction{Eye Rays}
The beholder shoots three of the following magical eye rays at random (reroll duplicates), choosing one to three targets it can see within 120 feet of it:


\begin{description}
\item[1. Charm Ray.]
The targeted creature must succeed on a DC 16 Wisdom saving throw or be charmed by the beholder for 1 hour, or until the beholder harms the creature.
\item[2. Paralyzing Ray.]
The targeted creature must succeed on a DC 16 Constitution saving throw or be paralyzed for 1 minute. The target can repeat the saving throw at the end of each of its turns, ending the effect on itself on a success.
\item[3. Fear Ray.]
The targeted creature must succeed on a DC 16 Wisdom saving throw or be frightened for 1 minute. The target can repeat the saving throw at the end of each of its turns, ending the effect on itself on a success.
\item[4. Slowing Ray.]
The targeted creature must succeed on a DC 16 Dexterity saving throw. On a failed save, the target's speed is halved for 1 minute. In addition, the creature can't take reactions, and it can take either an action or a bonus action on its turn, not both. The creature can repeat the saving throw at the end of each of its turns, ending the effect on itself on a success.
\item[5. Enervation Ray.]
The targeted creature must make a DC 16 Constitution saving throw, taking 36 (8d8) necrotic damage on a failed save, or half as much damage on a successful one.
\item[6. Telekinetic Ray.]
If the target is a creature, it must succeed on a DC 16 Strength saving throw or the beholder moves it up to 30 feet in any direction. It is restrained by the ray's telekinetic grip until the start of the beholder's next turn or until the beholder is incapacitated.

If the target is an object weighing 300 pounds or less that isn't being worn or carried, it is moved up to 30 feet in any direction. The beholder can also exert fine control on objects with this ray, such as manipulating a simple tool or opening a door or a container.

\item[7. Sleep Ray.]
The targeted creature must succeed on a DC 16 Wisdom saving throw or fall asleep and remain unconscious for 1 minute. The target awakens if it takes damage or another creature takes an action to wake it. This ray has no effect on constructs and undead.
\item[8. Petrification Ray.]
The targeted creature must make a DC 16 Dexterity saving throw. On a failed save, the creature begins to turn to stone and is restrained. It must repeat the saving throw at the end of its next turn. On a success, the effect ends. On a failure, the creature is petrified until freed by the  \textit{greater restoration} spell or other magic.
\item[9. Disintegration Ray.]
If the target is a creature, it must succeed on a DC 16 Dexterity saving throw or take 45 (10d8) force damage. If this damage reduces the creature to 0 hit points, its body becomes a pile of fine gray dust.

If the target is a Large or smaller nonmagical object or creation of magical force, it is disintegrated without a saving throw. If the target is a Huge or larger object or creation of magical force, this ray disintegrates a 10-foot cube of it.

\item[10. Death Ray.]
The targeted creature must succeed on a DC 16 Dexterity saving throw or take 55 (10d10) necrotic damage. The target dies if the ray reduces it to 0 hit points.
\end{description}

\DndMonsterSection{Legendary Actions}
The beholder can take 3 legendary actions, using the Eye Ray option below. It can take only one legendary action at a time and only at the end of another creature's turn. The beholder regains spent legendary actions at the start of its turn.

\begin{DndMonsterLegendaryActions}
\DndMonsterLegendaryAction{Eye Ray}{The beholder uses one random eye ray.
}
\end{DndMonsterLegendaryActions}
\DndMonsterSection{Lair Actions}
When fighting inside its lair, a beholder can invoke the ambient magic to take lair actions. On initiative count 20 (losing initiative ties), the beholder can take one lair action to cause one of the following effects:


\begin{itemize}
\item A 50-foot-square area of ground within 120 feet of the beholder becomes slimy; that area is difficult terrain until initiative count 20 on the next round.
\item Walls within 120 feet of the beholder sprout grasping appendages until initiative count 20 on the round after next. Each creature of the beholder's choice that starts its turn within 10 feet of such a wall must succeed on a DC 15 Dexterity saving throw or be grappled. Escaping requires a successful DC 15 Strength (Athletics) or Dexterity (Acrobatics) check.
\item An eye opens on a solid surface within 60 feet of the beholder. One random eye ray of the beholder shoots from that eye at a target of the beholder's choice that it can see. The eye then closes and disappears.
\end{itemize}

The beholder can't repeat an effect until they have all been used, and it can't use the same effect two rounds in a row.

\DndMonsterSection{Regional Effects}
A region containing a beholder's lair is warped by the creature's unnatural presence, which creates one or more of the following effects:


\begin{itemize}
\item Creatures within 1 mile of the beholder's lair sometimes feel as if they're being watched when they aren't.
\item When the beholder sleeps, minor warps in reality occur within 1 mile of its lair and then vanish 24 hours later. Marks on cave walls might change subtly, an eerie trinket might appear where none existed before, harmless slime might coat a statue, and so on. These effects apply only to natural surfaces and to nonmagical objects that aren't on anyone's person.
\end{itemize}

If the beholder dies, these effects fade over the course of 1d10 days.

\end{multicols}
\end{DndMonster}



\begin{DndMonster}[float*=b,width=\textwidth + 8pt]{Bel}
\begin{multicols}{2}
\DndMonsterType{Large fiend (devil), lawful evil}
\DndMonsterBasics[
armor-class = {19 (natural armor)},
hit-points = {\DndDice{27d10 + 216}},
speed = {30 ft., fly 60 ft.}
]
\DndMonsterAbilityScores[
 str = 28,
 dex = 14,
 con = 26,
 int = 25,
 wis = 19,
 cha = 26
]
\DndMonsterDetails[
saving-throws = {Dex +10, Con +16, Wis +12},
skills = {Arcana +15, Deception +16, Insight +12, Persuasion +16},
damage-resistances = {cold; bludgeoning, piercing, slashing from nonmagical attacks that aren't silvered},
damage-immunities = {fire, poison},
condition-immunities = {poisoned},
senses = {truesight 120 ft., passive perception 14},
languages = {Common, Infernal, telepathy 120 ft.},
challenge = 25,
]
\DndMonsterAction{Fear Aura}
Any creature hostile to Bel that starts its turn within 20 feet of him must make a DC 23 Wisdom saving throw unless Bel is incapacitated. If the save fails, the creature has the frightened condition until the start of its next turn. If a creature's saving throw is successful, the creature is immune to Bel's Fear Aura for the next 24 hours.

\DndMonsterAction{Legendary Resistance (3/Day)}
If Bel fails a saving throw, he can choose to succeed instead.

\DndMonsterAction{Magic Resistance}
Bel has advantage on saving throws against spells and other magical effects.

\DndMonsterAction{Spellcasting}
Bel casts one of the following spells, requiring no material components and using Charisma as the spellcasting ability (spell save DC 24):
\begin{DndMonsterSpells}
  \DndInnateSpellLevel{\textit{Detect Magic}, \textit{Fireball}}
  \DndInnateSpellLevel[2e]{\textit{Dispel Magic}, \textit{Hold Monster}, \textit{Mirror Image}, \textit{Mislead}, \textit{Raise Dead}, \textit{Teleport}, \textit{Wall of Fire}}
  \DndInnateSpellLevel[1e]{\textit{Imprisonment}, \textit{Meteor Swarm}}
\end{DndMonsterSpells}
\DndMonsterSection{Actions}
\DndMonsterAction{Multiattack}
Bel attacks twice with his Greatsword and once with his Tail.

\DndMonsterAction{Greatsword}
\textsl{Melee Weapon Attack:} 17 to hit, reach 10 ft., one target. \textsl{Hit:} 23 (4d6 + 9) slashing damage plus 21 (6d6) fire damage. If the target is a flammable object that isn't being held or worn, it catches fire.

\DndMonsterAction{Tail}
\textsl{Melee Weapon Attack:} 17 to hit, reach 10 ft., one target. \textsl{Hit:} 25 (3d10 + 9) bludgeoning damage. If the target is a creature, it must succeed on a DC 23 Constitution saving throw or have the stunned condition until the end of its next turn.

\DndMonsterSection{Legendary Actions}
Bel can take 3 legendary actions, choosing from the options below. Only one legendary action can be used at a time and only at the end of another creature's turn. Bel regains spent legendary actions at the start of its turn.

\begin{DndMonsterLegendaryActions}
\DndMonsterLegendaryAction{Fireball}{Bel casts \textit{Fireball}.
}
\DndMonsterLegendaryAction{Tactical Edge (Costs 2 Actions)}{Roll a d6 for Bel. The number rolled on the die is subtracted from the next attack roll made against Bel or an ally of his choice within the next minute.
}
\DndMonsterLegendaryAction{Summon Ice Devil (Costs 3 Actions)}{Bel magically summons an \textbf{ice devil} with an ice spear (as described in the ice devil's entry in the Monster Manual). The ice devil appears in an unoccupied space within 60 feet of Bel, acts as Bel's ally, and can summon other devils if it has such power. The ice devil remains until Bel dies or until he dismisses it with an action.
}
\end{DndMonsterLegendaryActions}
\DndMonsterSection{Regional Effects}
The region containing Bel's lair is corrupted by his presence, which creates the following effect:


\begin{description}
\item[Lava.]
There are numerous rivers of lava within 3 miles of the lair. Anyone who enters the lava for the first time or starts their turn in it takes 33 (6d10) fire damage.

\end{description}

If Bel dies, the effects fade over the course of 1d10 days.

\end{multicols}
\end{DndMonster}



\begin{DndMonster}[float*=b,width=\textwidth + 8pt]{Belial}
\begin{multicols}{2}
\DndMonsterType{Medium fiend (devil), lawful evil}
\DndMonsterBasics[
armor-class = {21 (natural armor)},
hit-points = {\DndDice{40d8 + 240}},
speed = {30 ft., fly 60 ft.}
]
\DndMonsterAbilityScores[
 str = 25,
 dex = 17,
 con = 22,
 int = 22,
 wis = 22,
 cha = 29
]
\DndMonsterDetails[
saving-throws = {Str +15, Dex +11, Con +14, Wis +14},
skills = {Athletics +15, Deception +25, Insight +14, Intimidation +17, Perception +14, Persuasion +25},
damage-resistances = {cold; bludgeoning, piercing, slashing from nonmagical attacks that aren't silvered},
damage-immunities = {fire, poison},
condition-immunities = {charmed, poisoned},
senses = {darkvision 120 ft., passive perception 24},
languages = {Celestial, Common, Draconic, Infernal, telepathy 120 ft.},
challenge = 25,
]
\DndMonsterAction{Deep Wounds}
All damage dealt by Belial can only be restored through magical means.

\DndMonsterAction{Devil's Sight}
Magical darkness doesn't impede Belials darkvision.

\DndMonsterAction{Fear Aura}
When a creature starts their turn within 120 feet of Belial, they must succeed on a DC 22 Wisdom saving throw or have the frightened condition until they leave the aura. A creature that succeeds on the saving throw is immune to this effect for 1 hour.

\DndMonsterAction{Fiendish Regeneration}
Belial regains 20 hit points at the start of his turn. If he takes radiant damage this trait doesn't function at the start of his next turn. Belial dies only if he starts his turn with 0 hit points and doesn't regenerate.

\DndMonsterAction{Legendary Resistance (3/Day)}
If Belial fails a saving throw, he can choose to succeed instead.

\DndMonsterAction{Magic Resistance}
Belial has advantage on saving throws against spells and other magical effects.

\DndMonsterSection{Actions}
\DndMonsterAction{Multiattack}
Belial makes three attacks using his Ranseur of Torture.

\DndMonsterAction{Ranseur of Torture}
\textsl{Melee Weapon Attack:} 18 to hit, reach 10 ft., one target. \textsl{Hit:} 21 (2d10 + 10) piercing damage plus 5 (1d10) necrotic damage, and the target must make a DC 23 Constitution saving throw. On a failed save, the target's movement speed is reduced by 5 feet until the end of their next turn. This penalty can stack.

\DndMonsterAction{Walls of Phlegethos}
Belial creates a wall of fire on a solid surface within 120 feet of him. The wall is 20 feet high and 1 foot thick, and Belial can make the wall up to 60 feet long or create a ringed wall up to 20 feet in diameter. The wall is opaque and lasts until Belial dismisses it. When the wall appears, each creature within its area must make a DC 25 Dexterity saving throw, taking 36 (8d8) fire damage on a failed save, or half as much damage on a successful one. A creature that ends its turn inside the wall, or within 5 feet of it, must repeat the save.

\DndMonsterAction{Symbol of Insanity (Recharge 4\textendash 6)}
Belial draws a magic symbol at a point he can see within 120 feet of him. All creatures within a 20-foot-radius sphere centered on the symbol must make a DC 25 Intelligence saving throw. Targets take 45 (10d8) psychic damage on a failed save, or half as much damage on a successful one. Creatures that fail the save immediately use their reaction to make a melee attack against a random creature within their reach, then move their full movement in a random direction.

\DndMonsterSection{Reactions}
\DndMonsterAction{Pleasure in Pain}
As a reaction to taking damage, Belial immediately regenerates half the damage he took.

\DndMonsterSection{Legendary Actions}
Belial can take 3 legendary actions, choosing from the options below. Only one legendary action can be used at a time and only at the end of another creature's turn. Belial regains spent legendary actions at the start of its turn.

\begin{DndMonsterLegendaryActions}
\DndMonsterLegendaryAction{Ranseur}{Belial makes a Ranseur of Torture attack.
}
\DndMonsterLegendaryAction{Eruption (Costs 2 Actions)}{Belial causes molten stone to erupt from a point he can see within 90 feet of him. All creatures within a 20-foot-radius sphere centered on that point must make a DC 23 Dexterity saving throw. Targets take 28 (8d6) fire damage on a failed save, or half as much damage on a successful one.
}
\DndMonsterLegendaryAction{Call Underling (Costs 3 Actions)}{Belial summons an allied \textbf{horned devil} in an unoccupied space that he can see.
}
\end{DndMonsterLegendaryActions}
\DndMonsterSection{Regional Effects}
The region containing Belial's lair is corrupted by his presence, which creates the following effect:


\begin{description}
\item[Smoke-cloud.]
A vast fog rises 3 feet above ground within 1 mile of the lair, completely obscuring the ground. Any creature having the prone condition at the start of their turn must succeed on a DC 22 Constitution saving throw or take 7 (2d12) necrotic damage from choking on the noxious fumes.

\end{description}

If Belial dies, the effects fade over the course of 1d10 days.

\end{multicols}
\end{DndMonster}



\begin{DndMonster}[float=t]{Berserker}
\DndMonsterType{Medium humanoid (any race), any chaotic alignment}
\DndMonsterBasics[
armor-class = {13 (\textit{hide armor})},
hit-points = {\DndDice{9d8 + 27}},
speed = {30 ft.}
]
\DndMonsterAbilityScores[
 str = 16,
 dex = 12,
 con = 17,
 int = 9,
 wis = 11,
 cha = 9
]
\DndMonsterDetails[
languages = {any one language (usually Common)},
challenge = 2,
]
\DndMonsterAction{Reckless}
At the start of its turn, the berserker can gain advantage on all melee weapon attack rolls during that turn, but attack rolls against it have advantage until the start of its next turn.

\DndMonsterSection{Actions}
\DndMonsterAction{Greataxe}
\textsl{Melee Weapon Attack:} 5 to hit, reach 5 ft., one target. \textsl{Hit:} 9 (1d12 + 3) slashing damage.

\end{DndMonster}



\begin{DndMonster}[float=t]{Black Abishai}
\DndMonsterType{Medium fiend (devil), lawful evil}
\DndMonsterBasics[
armor-class = {15 (natural armor)},
hit-points = {\DndDice{9d8 + 18}},
speed = {30 ft., fly 40 ft.}
]
\DndMonsterAbilityScores[
 str = 14,
 dex = 17,
 con = 14,
 int = 13,
 wis = 16,
 cha = 11
]
\DndMonsterDetails[
saving-throws = {Dex +6, Wis +6},
skills = {Perception +6, Stealth +6},
damage-resistances = {cold; bludgeoning, piercing, slashing from nonmagical attacks that aren't silvered},
damage-immunities = {acid, fire, poison},
condition-immunities = {poisoned},
senses = {darkvision 120 ft., passive perception 16},
languages = {Draconic, Infernal, telepathy 120 ft.},
challenge = 7,
]
\DndMonsterAction{Devil's Sight}
Magical darkness doesn't impede the abishai's darkvision.

\DndMonsterAction{Magic Resistance}
The abishai has advantage on saving throws against spells and other magical effects.

\DndMonsterAction{Creeping Darkness (Recharge 6)}
The abishai casts \textit{darkness} at a point within 120 feet of it, requiring no spell components or concentration. Wisdom is its spellcasting ability for this spell. While the spell persists, the abishai can move the area of darkness up to 60 feet as a bonus action.
\begin{DndMonsterSpells}
\end{DndMonsterSpells}
\DndMonsterSection{Actions}
\DndMonsterAction{Multiattack}
The abishai makes one Bite attack and two Scimitar attacks.

\DndMonsterAction{Bite}
\textsl{Melee Weapon Attack:} 6 to hit, reach 5 ft., one target. \textsl{Hit:} 8 (1d10 + 3) piercing damage plus 9 (2d8) acid damage.

\DndMonsterAction{Scimitar}
\textsl{Melee Weapon Attack:} 6 to hit, reach 5 ft., one target. \textsl{Hit:} 6 (1d6 + 3) force damage.

\DndMonsterSection{Bonus Actions}
\DndMonsterAction{Shadow Stealth}
While in dim light or darkness, the abishai takes the Hide action.

\end{DndMonster}



\begin{DndMonster}[float=t]{Black Pudding}
\DndMonsterType{Large ooze, unaligned}
\DndMonsterBasics[
armor-class = {7},
hit-points = {\DndDice{10d10 + 30}},
speed = {20 ft., climb 20 ft.}
]
\DndMonsterAbilityScores[
 str = 16,
 dex = 5,
 con = 16,
 int = 1,
 wis = 6,
 cha = 1
]
\DndMonsterDetails[
damage-immunities = {acid, cold, lightning, slashing},
condition-immunities = {blinded, charmed, deafened, exhaustion, frightened, prone},
senses = {blindsight 60 ft. (blind beyond this radius), passive perception 8},
challenge = 4,
]
\DndMonsterAction{Amorphous}
The pudding can move through a space as narrow as 1 inch wide without squeezing.

\DndMonsterAction{Corrosive Form}
A creature that touches the pudding or hits it with a melee attack while within 5 feet of it takes 4 (1d8) acid damage. Any nonmagical weapon made of metal or wood that hits the pudding corrodes. After dealing damage, the weapon takes a permanent and cumulative −1 penalty to damage rolls. If its penalty drops to −5, the weapon is destroyed. Nonmagical ammunition made of metal or wood that hits the pudding is destroyed after dealing damage. The pudding can eat through 2-inch-thick, nonmagical wood or metal in 1 round.

\DndMonsterAction{Spider Climb}
The pudding can climb difficult surfaces, including upside down on ceilings, without needing to make an ability check.

\DndMonsterSection{Actions}
\DndMonsterAction{Pseudopod}
\textsl{Melee Weapon Attack:} 5 to hit, reach 5 ft., one target. \textsl{Hit:} 6 (1d6 + 3) bludgeoning damage plus 18 (4d8) acid damage. In addition, nonmagical armor worn by the target is partly dissolved and takes a permanent and cumulative −1 penalty to the AC it offers. The armor is destroyed if the penalty reduces its AC to 10.

\DndMonsterSection{Reactions}
\DndMonsterAction{Split}
When a pudding that is Medium or larger is subjected to lightning or slashing damage, it splits into two new puddings if it has at least 10 hit points. Each new pudding has hit points equal to half the original pudding's, rounded down. New puddings are one size smaller than the original pudding.

\end{DndMonster}



\begin{DndMonster}[float=t]{Blue Abishai}
\DndMonsterType{Medium fiend (devil, wizard), lawful evil}
\DndMonsterBasics[
armor-class = {19 (natural armor)},
hit-points = {\DndDice{27d8 + 81}},
speed = {30 ft., fly 50 ft.}
]
\DndMonsterAbilityScores[
 str = 15,
 dex = 14,
 con = 17,
 int = 22,
 wis = 23,
 cha = 18
]
\DndMonsterDetails[
saving-throws = {Int +12, Wis +12},
skills = {Arcana +12},
damage-resistances = {cold; bludgeoning, piercing, slashing from nonmagical attacks that aren't silvered},
damage-immunities = {fire, lightning, poison},
condition-immunities = {poisoned},
senses = {darkvision 120 ft., passive perception 16},
languages = {Draconic, Infernal, telepathy 120 ft.},
challenge = 17,
]
\DndMonsterAction{Devil's Sight}
Magical darkness doesn't impede the abishai's darkvision.

\DndMonsterAction{Magic Resistance}
The abishai has advantage on saving throws against spells and other magical effects.

\DndMonsterAction{Spellcasting}
The abishai casts one of the following spells, using Intelligence as the spellcasting ability (spell save DC 20):
\begin{DndMonsterSpells}
  \DndInnateSpellLevel{\textit{disguise self}, \textit{mage hand}, \textit{minor illusion}}
  \DndInnateSpellLevel[2e]{\textit{charm person}, \textit{dispel magic}, \textit{greater invisibility}, \textit{wall of force}}
\end{DndMonsterSpells}
\DndMonsterSection{Actions}
\DndMonsterAction{Multiattack}
The abishai makes three Bite or Lightning Strike attacks.

\DndMonsterAction{Bite}
\textsl{Melee Weapon Attack:} 8 to hit, reach 5 ft., one target. \textsl{Hit:} 13 (2d10 + 2) piercing damage plus 14 (4d6) lightning damage.

\DndMonsterAction{Lightning Strike}
\textsl{Ranged Spell Attack:} 12 to hit, range 120 ft., one target. \textsl{Hit:} 36 (8d8) lightning damage.

\DndMonsterSection{Bonus Actions}
\DndMonsterAction{Teleport}
The abishai teleports, along with any equipment it is wearing or carrying, up to 30 feet to an unoccupied space that it can see.

\end{DndMonster}



\begin{DndMonster}[float=t]{Bone Devil}
\DndMonsterType{Large fiend (devil), lawful evil}
\DndMonsterBasics[
armor-class = {19 (natural armor)},
hit-points = {\DndDice{15d10 + 60}},
speed = {40 ft., fly 40 ft.}
]
\DndMonsterAbilityScores[
 str = 18,
 dex = 16,
 con = 18,
 int = 13,
 wis = 14,
 cha = 16
]
\DndMonsterDetails[
saving-throws = {Int +5, Wis +6, Cha +7},
skills = {Deception +7, Insight +6},
damage-resistances = {cold; bludgeoning, piercing, slashing from nonmagical attacks that aren't silvered},
damage-immunities = {fire, poison},
condition-immunities = {poisoned},
senses = {darkvision 120 ft., passive perception 12},
languages = {Infernal, telepathy 120 ft.},
challenge = 9,
]
\DndMonsterAction{Devil's Sight}
Magical darkness doesn't impede the devil's darkvision.

\DndMonsterAction{Magic Resistance}
The devil has advantage on saving throws against spells and other magical effects.

\DndMonsterSection{Actions}
\DndMonsterAction{Multiattack}
The devil makes three attacks: two with its claws and one with its sting.

\DndMonsterAction{Claw}
\textsl{Melee Weapon Attack:} 8 to hit, reach 10 ft., one target. \textsl{Hit:} 8 (1d8 + 4) slashing damage.

\DndMonsterAction{Sting}
\textsl{Melee Weapon Attack:} 8 to hit, reach 10 ft., one target. \textsl{Hit:} 13 (2d8 + 4) piercing damage plus 17 (5d6) poison damage, and the target must succeed on a DC 14 Constitution saving throw or become poisoned for 1 minute. The target can repeat the saving throw at the end of each of its turns, ending the effect on itself on a success.

\end{DndMonster}



\begin{DndMonster}[float*=b,width=\textwidth + 8pt]{Brother Adramalech}
\begin{multicols}{2}
\DndMonsterType{Medium fiend (devil), lawful evil}
\DndMonsterBasics[
armor-class = {19 (natural armor)},
hit-points = {\DndDice{25d8 + 100}},
speed = {30 ft.}
]
\DndMonsterAbilityScores[
 str = 14,
 dex = 16,
 con = 18,
 int = 25,
 wis = 25,
 cha = 25
]
\DndMonsterDetails[
saving-throws = {Int +14, Wis +14, Cha +14},
skills = {Arcana +14, Deception +21, Insight +14, Intimidation +14, Persuasion +21, Religion +14},
damage-resistances = {cold; bludgeoning, piercing, slashing from nonmagical attacks that aren't silvered},
damage-immunities = {fire, poison},
condition-immunities = {charmed, poisoned},
senses = {darkvision 120 ft., passive perception 17},
languages = {Celestial, Common, Draconic, Infernal, telepathy 120 ft.},
challenge = 21,
]
\DndMonsterAction{Devil's Sight}
Magical darkness doesn't impede Adramalech's darkvision.

\DndMonsterAction{Magic Resistance}
Adramalech has advantage on saving throws against spells and other magical effects.

\DndMonsterAction{Unbreakable Bond}
If Adramalech is reduced to 0 hit points while Brother Morax still lives, Adramalech regenerates 50 hit points at the start of his next turn. This regeneration is interrupted if Brother Morax is reduced to 0 hit points before the start of Adramalech's turn.

\DndMonsterSection{Actions}
\DndMonsterAction{Multiattack}
Adramalech makes four attacks using Mental

\DndMonsterAction{Barrage}
He can replace one of the attacks with Brain Freeze or Cone of Madness (if available).

\DndMonsterAction{Mental Barrage}
\textsl{Melee or Ranged Weapon Attack:} 14 to hit, reach 5 ft. or range 60/120 ft., one target. \textsl{Hit:} 25 (4d8 + 7) psychic damage.

\DndMonsterAction{Brain Freeze}
Adramalech attempts to overload the brain of a creature he can see within 60 feet of him. The target must make a DC 22 Intelligence saving throw, taking 22 (4d10) psychic damage on a failed save, or half as much on a successful one. Creatures that fail the saving throw have the stunned condition for up to 1 minute. A stunned creature may repeat the saving throw at the end of each of their turns, ending the effect on a success.

\DndMonsterAction{Cone of Madness (Recharge 4\textendash 6)}
Adramalech unleashes visions in a 30-foot cone in front of him. All creatures in the cone must make a DC 22 Wisdom saving throw, taking 17 (5d6) psychic damage on a failed save, or half as much on a successful one. For up to 1 minute, creatures that failed the save can't take reactions and each turn must use their action to make a melee attack against the nearest creature. Creatures may repeat the saving throw at the end of each of their turns, ending the effect on a success.

\DndMonsterSection{Bonus Actions}
\DndMonsterAction{Conjure Effigy (Recharge 6)}
Adramalech targets a creature, reading its thoughts and creating an effigy of a loved one in an unoccupied space that he can see within 60 feet of Adramalech. The effigy has 55 (10d10) hit points, an AC of 10, and ability scores of 10. It occasionally shouts for help, but otherwise takes no actions. While the effigy is alive, the target creature has disadvantage on attacks targeting Adramalech and any time the effigy takes damage, the target creature also takes an equal amount as psychic damage. If the effigy is killed, the target creature gains a level of exhaustion. If Adramalech is killed while effigies exist, they harmlessly melt into a warm sludge.

\DndMonsterSection{Legendary Actions}
Brother Adramalech can take 3 legendary actions, choosing from the options below. Only one legendary action can be used at a time and only at the end of another creature's turn. Brother Adramalech regains spent legendary actions at the start of its turn.

\begin{DndMonsterLegendaryActions}
\DndMonsterLegendaryAction{Barrage}{Adramalech makes a Mental Barrage attack.
}
\DndMonsterLegendaryAction{Mental Blocker}{Adramalech fortifies the mind of a creature he can see within 60 feet of him, or himself. The target benefits as if the \textit{Greater Restoration} spell has been cast on them and is granted resistance to psychic damage until the end of Adramalech's next turn.
}
\end{DndMonsterLegendaryActions}
\end{multicols}
\end{DndMonster}



\begin{DndMonster}[float*=b,width=\textwidth + 8pt]{Brother Morax}
\begin{multicols}{2}
\DndMonsterType{Large fiend (devil), lawful evil}
\DndMonsterBasics[
armor-class = {21 (natural armor)},
hit-points = {\DndDice{25d10 + 200}},
speed = {40 ft.}
]
\DndMonsterAbilityScores[
 str = 24,
 dex = 22,
 con = 27,
 int = 14,
 wis = 16,
 cha = 18
]
\DndMonsterDetails[
saving-throws = {Str +14, Dex +13, Con +15},
skills = {Acrobatics +20, Athletics +21, Intimidation +11, Medicine +10, Stealth +13, Survival +10},
damage-resistances = {cold; bludgeoning, piercing, slashing from nonmagical attacks that aren't silvered},
damage-immunities = {fire, poison},
condition-immunities = {poisoned},
senses = {darkvision 120 ft., passive perception 13},
languages = {Celestial, Common, Draconic, Infernal, telepathy 120 ft.},
challenge = 21,
]
\DndMonsterAction{Brittle Spine}
For each Brittle Spine a creature has, the creature takes 3 (1d6) poison damage at the start of each of their turns, and when Morax activates Burrow or Deep Burrow. Spines can only be removed by a Regeneration spell or similar, or by killing Morax.

\DndMonsterAction{Devil's Sight}
Magical darkness doesn't impede Morax's darkvision.

\DndMonsterAction{Magic Resistance}
Morax has advantage on saving throws against spells and other magical effects.

\DndMonsterAction{Spined Hide}
Whenever Morax is hit with a melee attack, the attacker takes 7 (2d6) piercing damage and gains 1 Brittle Spine.

\DndMonsterAction{Unbreakable Bond}
If Morax is reduced to 0 hit points while Brother Adramalech still lives, Morax regenerates 50 hit points at the start of his next turn. This regeneration is interrupted if Brother Adramalech is reduced to 0 hit points before the start of Morax's turn.

\DndMonsterSection{Actions}
\DndMonsterAction{Multiattack}
Morax makes four attacks using his Claw, Spine Fling, or a combination of the two.

\DndMonsterAction{Claw}
\textsl{Melee Weapon Attack:} 14 to hit, reach 10 ft., one target. \textsl{Hit:} 17 (3d6 + 7) slashing damage, and the target gains 1 Brittle Spine.

\DndMonsterAction{Spine Fling}
\textsl{Ranged Weapon Attack:} 13 to hit, range 30/60 ft., one target. \textsl{Hit:} 12 (1d12 + 6) piercing damage, and the target gains 1 Brittle Spine.

\DndMonsterAction{Deep Burrow}
Morax immediately deals Brittle Spine ongoing damage to all creatures that have at least one Brittle Spine that he can see.

\DndMonsterAction{Bone Spikes (Recharge 6)}
Morax shoves an arm into the ground, rapidly growing bones that erupt at a point he can see within 100 feet of him. Each creature in a 60-foot-radius sphere centered on that point must make a DC 21 Dexterity saving throw, taking 28 (8d6) piercing damage on a failed save, or half as much damage on a successful one. Creatures that fail the save gain 2 Brittle Spines, while creatures that succeed only gain 1 Brittle Spine.

\DndMonsterSection{Legendary Actions}
Brother Morax can take 3 legendary actions, choosing from the options below. Only one legendary action can be used at a time and only at the end of another creature's turn. Brother Morax regains spent legendary actions at the start of its turn.

\begin{DndMonsterLegendaryActions}
\DndMonsterLegendaryAction{Claw}{Morax makes a Claw attack.
}
\DndMonsterLegendaryAction{Bone Splinters (Costs 2 Actions)}{Morax gestures at a creature he can see that has at least 1 Brittle Spine. That creature must succeed on a DC 21 Constitution saving throw or the number of Brittle Spines it has doubles.
}
\DndMonsterLegendaryAction{Burrow (Costs 2 Actions)}{Morax bristles the spines within a creature that he can see. The target immediately takes Brittle Spine ongoing damage.
}
\end{DndMonsterLegendaryActions}
\end{multicols}
\end{DndMonster}



\begin{DndMonster}[float*=b,width=\textwidth + 8pt]{Buer}
\begin{multicols}{2}
\DndMonsterType{Large fiend (devil), lawful evil}
\DndMonsterBasics[
armor-class = {21 (natural armor)},
hit-points = {\DndDice{30d10 + 150}},
speed = {25 ft.}
]
\DndMonsterAbilityScores[
 str = 23,
 dex = 22,
 con = 21,
 int = 16,
 wis = 18,
 cha = 17
]
\DndMonsterDetails[
saving-throws = {Str +15, Dex +13, Con +14},
skills = {Arcana +10, Athletics +22, Insight +11, Intimidation +17, Perception +11, Survival +11},
damage-resistances = {cold; bludgeoning, piercing, slashing from nonmagical attacks that aren't silvered},
damage-immunities = {fire, poison},
condition-immunities = {charmed, poisoned},
senses = {darkvision 120 ft., passive perception 21},
languages = {Celestial, Common, Draconic, Infernal, telepathy 120 ft.},
challenge = 21,
]
\DndMonsterAction{Contractually Obligated}
Buer has advantage to hit creatures under an infernal contract and when she hits such creatures she delivers an additional 9 (2d8) damage of a type the target is most vulnerable to. A creature bound under an infernal contract also makes all saving throws against effects originating from Buer with disadvantage.

\DndMonsterAction{Devil's Sight}
Magical darkness doesn't impede Buer's darkvision.

\DndMonsterAction{Magic Resistance}
Buer has advantage on saving throws against spells and other magical effects.

\DndMonsterSection{Actions}
\DndMonsterAction{Multiattack}
Buer makes three attacks using her Claw. She can replace two of the attacks with a Cleansing Bite attack.

\DndMonsterAction{Claw}
\textsl{Melee Weapon Attack:} 13 to hit, reach 10 ft., one target. \textsl{Hit:} 22 (3d10 + 6) slashing damage.

\DndMonsterAction{Cleansing Bite}
\textsl{Melee Weapon Attack:} 13 to hit, reach 5 ft., one target. \textsl{Hit:} 19 (2d12 + 6) bludgeoning damage, healing Buer for an amount equal to the damage dealt.

\DndMonsterAction{Spectral Reaping (Recharge 6)}
Buer focuses her ethereal form, dashing up to 50 feet in a straight line through the ethereal plane. Any creatures in the line must make a DC 21 Constitution saving throw, taking 42 (12d6) force damage on a failed save, or half as much damage on a successful one. Buer may reappear at any unoccupied space along the line. For 1 minute after using Spectral Reaping, Buer is empowered by ethereal essence. Her Cleansing Bite attack now deals force damage instead of bludgeoning.

\DndMonsterSection{Bonus Actions}
\DndMonsterAction{Defensive Trip (Recharge 5\textendash 6)}
Buer attempts to slow the actions of a creature she can see within 60 feet of her. The target must succeed on a DC 18 Wisdom saving throw or be slowed until the end of its next turn. A slowed creature's speed is halved, it takes a-2 penalty to AC and Dexterity saving throws, and it can't use its reactions. Regardless of the slowed creature's abilities or magic items, it can't make more than one melee or ranged attack during its turn.

\DndMonsterSection{Reactions}
\DndMonsterAction{Soul Exchange}
When Buer is subjected to a spell targeting only her, she can attempt to swap places with a creature she damaged on her last turn. The creature must make a DC 19 Wisdom saving throw. On a failed save, the creature and Buer swap positions, and the spell now targets the creature instead of Buer.

\DndMonsterSection{Legendary Actions}
Buer can take 3 legendary actions, choosing from the options below. Only one legendary action can be used at a time and only at the end of another creature's turn. Buer regains spent legendary actions at the start of its turn.

\begin{DndMonsterLegendaryActions}
\DndMonsterLegendaryAction{Teamwork}{Buer or one other devil within 60 feet of Buer makes a melee weapon attack against a target of Buer's choosing.
}
\DndMonsterLegendaryAction{Backdoor Clause}{Buer wreathes energy around a creature she can see within 60 feet of her. Until the end of Buer's next turn, the target is covered in fire or cold energy, Buer's choice. Each time the target is struck with a melee attack, the attacker takes 9 (2d8) fire or cold damage.
}
\DndMonsterLegendaryAction{Refresh Souls (Costs 2 Actions)}{Buer devours a Soul Coin, instantly refreshing one Recharge ability of any devil within 60 feet, including herself.
}
\end{DndMonsterLegendaryActions}
\end{multicols}
\end{DndMonster}



\begin{DndMonster}[float=t]{Bugbear}
\DndMonsterType{Medium humanoid (goblinoid), chaotic evil}
\DndMonsterBasics[
armor-class = {16 (\textit{hide armor})},
hit-points = {\DndDice{5d8 + 5}},
speed = {30 ft.}
]
\DndMonsterAbilityScores[
 str = 15,
 dex = 14,
 con = 13,
 int = 8,
 wis = 11,
 cha = 9
]
\DndMonsterDetails[
skills = {Stealth +6, Survival +2},
senses = {darkvision 60 ft., passive perception 10},
languages = {Common, Goblin},
challenge = 1,
]
\DndMonsterAction{Brute}
A melee weapon deals one extra die of its damage when the bugbear hits with it (included in the attack).

\DndMonsterAction{Surprise Attack}
If the bugbear surprises a creature and hits it with an attack during the first round of combat, the target takes an extra 7 (2d6) damage from the attack.

\DndMonsterSection{Actions}
\DndMonsterAction{Morningstar}
\textsl{Melee Weapon Attack:} 4 to hit, reach 5 ft., one target. \textsl{Hit:} 11 (2d8 + 2) piercing damage.

\DndMonsterAction{Javelin}
\textsl{Melee or Ranged Weapon Attack:} 4 to hit, reach 5 ft. or range 30/120 ft., one target. \textsl{Hit:} 9 (2d6 + 2) piercing damage in melee or 5 (1d6 + 2) piercing damage at range.

\end{DndMonster}


\clearpage

\begin{DndMonster}[float=t]{Cambion}
\DndMonsterType{Medium fiend, any evil alignment}
\DndMonsterBasics[
armor-class = {19 (\textit{scale mail})},
hit-points = {\DndDice{11d8 + 33}},
speed = {30 ft., fly 60 ft.}
]
\DndMonsterAbilityScores[
 str = 18,
 dex = 18,
 con = 16,
 int = 14,
 wis = 12,
 cha = 16
]
\DndMonsterDetails[
saving-throws = {Str +7, Con +6, Int +5, Cha +6},
skills = {Deception +6, Intimidation +6, Perception +4, Stealth +7},
damage-resistances = {cold; fire; lightning; poison; bludgeoning, piercing, slashing from nonmagical attacks},
senses = {darkvision 60 ft., passive perception 14},
languages = {Abyssal, Common, Infernal},
challenge = 5,
]
\DndMonsterAction{Fiendish Blessing}
The AC of the cambion includes its Charisma bonus.

\DndMonsterAction{Innate Spellcasting}
The cambion's spellcasting ability is Charisma (spell save DC 14). The cambion can innately cast the following spells, requiring no material components:
\begin{DndMonsterSpells}
  \DndInnateSpellLevel[1]{\textit{plane shift} (self only)}
  \DndInnateSpellLevel[3e]{\textit{alter self}, \textit{command}, \textit{detect magic}}
\end{DndMonsterSpells}
\DndMonsterSection{Actions}
\DndMonsterAction{Multiattack}
The cambion makes two melee attacks or uses its Fire Ray twice.

\DndMonsterAction{Spear}
\textsl{Melee or Ranged Weapon Attack:} 7 to hit, reach 5 ft. or range 20/60 ft., one target. \textsl{Hit:} 7 (1d6 + 4) piercing damage, or 8 (1d8 + 4) piercing damage if used with two hands to make a melee attack, plus 3 (1d6) fire damage.

\DndMonsterAction{Fire Ray}
\textsl{Ranged Spell Attack:} 7 to hit, range 120 ft., one target. \textsl{Hit:} 10 (3d6) fire damage.

\DndMonsterAction{Fiendish Charm}
One humanoid the cambion can see within 30 feet of it must succeed on a DC 14 Wisdom saving throw or be magically charmed for 1 day. The charmed target obeys the cambion's spoken commands. If the target suffers any harm from the cambion or another creature or receives a suicidal command from the cambion, the target can repeat the saving throw, ending the effect on itself on a success. If a target's saving throw is successful, or if the effect ends for it, the creature is immune to the cambion's Fiendish Charm for the next 24 hours.

\end{DndMonster}



\begin{DndMonster}[float=t]{Centaur}
\DndMonsterType{Large monstrosity, neutral good}
\DndMonsterBasics[
armor-class = {12},
hit-points = {\DndDice{6d10 + 12}},
speed = {50 ft.}
]
\DndMonsterAbilityScores[
 str = 18,
 dex = 14,
 con = 14,
 int = 9,
 wis = 13,
 cha = 11
]
\DndMonsterDetails[
skills = {Athletics +6, Perception +3, Survival +3},
languages = {Elvish, Sylvan},
challenge = 2,
]
\DndMonsterAction{Charge}
If the centaur moves at least 30 feet straight toward a target and then hits it with a pike attack on the same turn, the target takes an extra 10 (3d6) piercing damage.

\DndMonsterSection{Actions}
\DndMonsterAction{Multiattack}
The centaur makes two attacks: one with its pike and one with its hooves or two with its longbow.

\DndMonsterAction{Pike}
\textsl{Melee Weapon Attack:} 6 to hit, reach 10 ft., one target. \textsl{Hit:} 9 (1d10 + 4) piercing damage.

\DndMonsterAction{Hooves}
\textsl{Melee Weapon Attack:} 6 to hit, reach 5 ft., one target. \textsl{Hit:} 11 (2d6 + 4) bludgeoning damage.

\DndMonsterAction{Longbow}
\textsl{Ranged Weapon Attack:} 4 to hit, range 150/600 ft., one target. \textsl{Hit:} 6 (1d8 + 2) piercing damage.

\end{DndMonster}



\begin{DndMonster}[float=t]{Chain Devil}
\DndMonsterType{Medium fiend (devil), lawful evil}
\DndMonsterBasics[
armor-class = {16 (natural armor)},
hit-points = {\DndDice{10d8 + 40}},
speed = {30 ft.}
]
\DndMonsterAbilityScores[
 str = 18,
 dex = 15,
 con = 18,
 int = 11,
 wis = 12,
 cha = 14
]
\DndMonsterDetails[
saving-throws = {Con +7, Wis +4, Cha +5},
damage-resistances = {cold; bludgeoning, piercing, slashing from nonmagical attacks that aren't silvered},
damage-immunities = {fire, poison},
condition-immunities = {poisoned},
senses = {darkvision 120 ft., passive perception 11},
languages = {Infernal, telepathy 120 ft.},
challenge = 8,
]
\DndMonsterAction{Devil's Sight}
Magical darkness doesn't impede the devil's darkvision.

\DndMonsterAction{Magic Resistance}
The devil has advantage on saving throws against spells and other magical effects.

\DndMonsterSection{Actions}
\DndMonsterAction{Multiattack}
The devil makes two attacks with its chains.

\DndMonsterAction{Chain}
\textsl{Melee Weapon Attack:} 8 to hit, reach 10 ft., one target. \textsl{Hit:} 11 (2d6 + 4) slashing damage. The target is grappled (escape DC 14) if the devil isn't already grappling a creature. Until this grapple ends, the target is restrained and takes 7 (2d6) piercing damage at the start of each of its turns.

\DndMonsterAction{Animate Chains (Recharges after a Short or Long Rest)}
Up to four chains the devil can see within 60 feet of it magically sprout razor-edged barbs and animate under the devil's control, provided that the chains aren't being worn or carried.

Each animated chain is an object with AC 20, 20 hit points, resistance to piercing damage, and immunity to psychic and thunder damage. When the devil uses Multiattack on its turn, it can use each animated chain to make one additional chain attack. An animated chain can grapple one creature of its own but can't make attacks while grappling. An animated chain reverts to its inanimate state if reduced to 0 hit points or if the devil is incapacitated or dies.

\DndMonsterSection{Reactions}
\DndMonsterAction{Unnerving Mask}
When a creature the devil can see starts its turn within 30 feet of the devil, the devil can create the illusion that it looks like one of the creature's departed loved ones or bitter enemies. If the creature can see the devil, it must succeed on a DC 14 Wisdom saving throw or be frightened until the end of its turn.

\end{DndMonster}



\begin{DndMonster}[float=t]{Chasme}
\DndMonsterType{Large fiend (demon), chaotic evil}
\DndMonsterBasics[
armor-class = {15 (natural armor)},
hit-points = {\DndDice{13d10 + 13}},
speed = {20 ft., fly 60 ft.}
]
\DndMonsterAbilityScores[
 str = 15,
 dex = 15,
 con = 12,
 int = 11,
 wis = 14,
 cha = 10
]
\DndMonsterDetails[
saving-throws = {Dex +5, Wis +5},
skills = {Perception +5},
damage-resistances = {cold, fire, lightning},
damage-immunities = {poison},
condition-immunities = {poisoned},
senses = {blindsight 10 ft., darkvision 120 ft., passive perception 15},
languages = {Abyssal, telepathy 120 ft.},
challenge = 6,
]
\DndMonsterAction{Drone}
The chasme produces a horrid droning sound to which demons are immune. Any other creature that starts its turn with in 30 feet of the chasme must succeed on a DC 12 Constitution saving throw or fall unconscious for 10 minutes. A creature that can't hear the drone automatically succeeds on the save. The effect on the creature ends if it takes damage or if another creature takes an action to splash it with holy water. If a creature's saving throw is successful or the effect ends for it, it is immune to the drone for the next 24 hours.

\DndMonsterAction{Magic Resistance}
The chasme has advantage on saving throws against spells and other magical effects.

\DndMonsterAction{Spider Climb}
The chasme can climb difficult surfaces, including upside down on ceilings, without needing to make an ability check.

\DndMonsterSection{Actions}
\DndMonsterAction{Proboscis}
\textsl{Melee Weapon Attack:} 5 to hit, reach 5 ft., one creature. \textsl{Hit:} 16 (4d6 + 2) piercing damage plus 24 (7d6) necrotic damage, and the target's hit point maximum is reduced by an amount equal to the necrotic damage taken. If this effect reduces a creature's hit point maximum to 0, the creature dies. This reduction to a creature's hit point maximum lasts until the creature finishes a long rest or until it is affected by a spell like  \textit{greater restoration}.

\end{DndMonster}



\begin{DndMonster}[float=t]{Chuul}
\DndMonsterType{Large aberration, chaotic evil}
\DndMonsterBasics[
armor-class = {16 (natural armor)},
hit-points = {\DndDice{11d10 + 33}},
speed = {30 ft., swim 30 ft.}
]
\DndMonsterAbilityScores[
 str = 19,
 dex = 10,
 con = 16,
 int = 5,
 wis = 11,
 cha = 5
]
\DndMonsterDetails[
skills = {Perception +4},
damage-immunities = {poison},
condition-immunities = {poisoned},
senses = {darkvision 60 ft., passive perception 14},
languages = {understands Deep Speech but can't speak},
challenge = 4,
]
\DndMonsterAction{Amphibious}
The chuul can breathe air and water.

\DndMonsterAction{Sense Magic}
The chuul senses magic within 120 feet of it at will. This trait otherwise works like the \textit{detect magic} spell but isn't itself magical.

\DndMonsterSection{Actions}
\DndMonsterAction{Multiattack}
The chuul makes two pincer attacks. If the chuul is grappling a creature, the chuul can also use its tentacles once.

\DndMonsterAction{Pincer}
\textsl{Melee Weapon Attack:} 6 to hit, reach 10 ft., one target. \textsl{Hit:} 11 (2d6 + 4) bludgeoning damage. The target is grappled (escape DC 14) if it is a Large or smaller creature and the chuul doesn't have two other creatures grappled.

\DndMonsterAction{Tentacles}
One creature grappled by the chuul must succeed on a DC 13 Constitution saving throw or be poisoned for 1 minute. Until this poison ends, the target is paralyzed. The target can repeat the saving throw at the end of each of its turns, ending the effect on itself on a success.

\end{DndMonster}



\begin{DndMonster}[float=t]{Cloaker}
\DndMonsterType{Large aberration, chaotic neutral}
\DndMonsterBasics[
armor-class = {14 (natural armor)},
hit-points = {\DndDice{12d10 + 12}},
speed = {10 ft., fly 40 ft.}
]
\DndMonsterAbilityScores[
 str = 17,
 dex = 15,
 con = 12,
 int = 13,
 wis = 12,
 cha = 14
]
\DndMonsterDetails[
skills = {Stealth +5},
senses = {darkvision 60 ft., passive perception 11},
languages = {Deep Speech, Undercommon},
challenge = 8,
]
\DndMonsterAction{Damage Transfer}
While attached to a creature, the cloaker takes only half the damage dealt to it (rounded down). and that creature takes the other half.

\DndMonsterAction{False Appearance}
While the cloaker remains motionless without its underside exposed, it is indistinguishable from a dark leather cloak.

\DndMonsterAction{Light Sensitivity}
While in bright light, the cloaker has disadvantage on attack rolls and Wisdom (Perception) checks that rely on sight.

\DndMonsterSection{Actions}
\DndMonsterAction{Multiattack}
The cloaker makes two attacks: one with its bite and one with its tail.

\DndMonsterAction{Bite}
\textsl{Melee Weapon Attack:} 6 to hit, reach 5 ft., one creature. \textsl{Hit:} 10 (2d6 + 3) piercing damage, and if the target is Large or smaller, the cloaker attaches to it. If the cloaker has advantage against the target, the cloaker attaches to the target's head, and the target is blinded and unable to breathe while the cloaker is attached. While attached, the cloaker can make this attack only against the target and has advantage on the attack roll. The cloaker can detach itself by spending 5 feet of its movement. A creature, including the target, can take its action to detach the cloaker by succeeding on a DC 16 Strength check.

\DndMonsterAction{Tail}
\textsl{Melee Weapon Attack:} 6 to hit, reach 10 ft., one creature. \textsl{Hit:} 7 (1d8 + 3) slashing damage.

\DndMonsterAction{Moan}
Each creature within 60 feet of the cloaker that can hear its moan and that isn't an aberration must succeed on a DC 13 Wisdom saving throw or become frightened until the end of the cloaker's next turn. If a creature's saving throw is successful, the creature is immune to the cloaker's moan for the next 24 hours.

\DndMonsterAction{Phantasms (Recharges after a Short or Long Rest)}
The cloaker magically creates three illusory duplicates of itself if it isn't in bright light. The duplicates move with it and mimic its actions, shifting position so as to make it impossible to track which cloaker is the real one. If the cloaker is ever in an area of bright light, the duplicates disappear.

Whenever any creature targets the cloaker with an attack or a harmful spell while a duplicate remains, that creature rolls randomly to determine whether it targets the cloaker or one of the duplicates. A creature is unaffected by this magical effect if it can't see or if it relies on senses other than sight.

A duplicate has the cloaker's AC and uses its saving throws. If an attack hits a duplicate, or if a duplicate fails a saving throw against an effect that deals damage, the duplicate disappears.

\end{DndMonster}



\begin{DndMonster}[float=t]{Cloud Giant}
\DndMonsterType{Huge giant, neutral good (50\%) neutral evil (50\%)}
\DndMonsterBasics[
armor-class = {14 (natural armor)},
hit-points = {\DndDice{16d12 + 96}},
speed = {40 ft.}
]
\DndMonsterAbilityScores[
 str = 27,
 dex = 10,
 con = 22,
 int = 12,
 wis = 16,
 cha = 16
]
\DndMonsterDetails[
saving-throws = {Con +10, Wis +7, Cha +7},
skills = {Insight +7, Perception +7},
languages = {Common, Giant},
challenge = 9,
]
\DndMonsterAction{Keen Smell}
The giant has advantage on Wisdom (Perception) checks that rely on smell.

\DndMonsterAction{Innate Spellcasting}
The giant's innate spellcasting ability is Charisma. It can innately cast the following spells, requiring no material components:
\begin{DndMonsterSpells}
  \DndInnateSpellLevel{\textit{detect magic}, \textit{fog cloud}, \textit{light}}
  \DndInnateSpellLevel[3e]{\textit{feather fall}, \textit{fly}, \textit{misty step}, \textit{telekinesis}}
  \DndInnateSpellLevel[1e]{\textit{control weather}, \textit{gaseous form}}
\end{DndMonsterSpells}
\DndMonsterSection{Actions}
\DndMonsterAction{Multiattack}
The giant makes two morningstar attacks.

\DndMonsterAction{Morningstar}
\textsl{Melee Weapon Attack:} 12 to hit, reach 10 ft., one target. \textsl{Hit:} 21 (3d8 + 8) piercing damage.

\DndMonsterAction{Rock}
\textsl{Ranged Weapon Attack:} 12 to hit, range 60/240 ft., one target. \textsl{Hit:} 30 (4d10 + 8) bludgeoning damage.

\end{DndMonster}



\begin{DndMonster}[float=t]{Commoner}
\DndMonsterType{Medium humanoid (any race), any alignment}
\DndMonsterBasics[
armor-class = {10},
hit-points = {\DndDice{1d8}},
speed = {30 ft.}
]
\DndMonsterAbilityScores[
 str = 10,
 dex = 10,
 con = 10,
 int = 10,
 wis = 10,
 cha = 10
]
\DndMonsterDetails[
languages = {any one language (usually Common)},
challenge = 0,
]
\DndMonsterSection{Actions}
\DndMonsterAction{Club}
\textsl{Melee Weapon Attack:} 2 to hit, reach 5 ft., one target. \textsl{Hit:} 2 (1d4) bludgeoning damage.

\end{DndMonster}



\begin{DndMonster}[float=t]{Corruption Devil (Paeliryon)}
\DndMonsterType{Large fiend (devil), lawful evil}
\DndMonsterBasics[
armor-class = {16 (natural armor)},
hit-points = {\DndDice{19d10 + 114}},
speed = {20 ft., burrow 20 ft., fly 60 ft. \textit{(hover)}}
]
\DndMonsterAbilityScores[
 str = 21,
 dex = 14,
 con = 22,
 int = 19,
 wis = 15,
 cha = 14
]
\DndMonsterDetails[
saving-throws = {Con +11, Wis +7},
skills = {Arcana +9, Deception +7, Insight +7, Intimidation +7},
damage-resistances = {cold; bludgeoning, piercing, slashing from nonmagical attacks that aren't silvered},
damage-immunities = {fire, poison},
condition-immunities = {poisoned, stunned},
senses = {blindsight 30 ft., darkvision 120 ft., passive perception 12},
languages = {all, telepathy 120 ft.},
challenge = 14,
]
\DndMonsterAction{Devil's Sight}
Magical darkness doesn't impede the paeliryon's darkvision.

\DndMonsterAction{Keen Eyes}
The paeliryon scores a critical hit on a roll of 19 or 20.

\DndMonsterAction{Magic Resistance}
The paeliryon has advantage on saving throws against spells and other magical effects.

\DndMonsterAction{Spellcasting}
The paeliryon casts one of the following spells, requiring no material components and using Intelligence as the spellcasting ability (spell save DC 17):
\begin{DndMonsterSpells}
  \DndInnateSpellLevel{\textit{Bestow Curse}, \textit{Command}, \textit{Darkness}, \textit{Hold Monster}, \textit{Major Image}}
  \DndInnateSpellLevel[1e]{\textit{Antilife Shell}, \textit{Incendiary Cloud}, \textit{Scrying}, \textit{Weird}}
\end{DndMonsterSpells}
\DndMonsterSection{Actions}
\DndMonsterAction{Multiattack}
The paeliryon makes three Claw attacks. It can replace one of the attacks with Draining Grasp (if available).

\DndMonsterAction{Claw}
\textsl{Melee Weapon Attack:} 10 to hit, reach 10 ft., one target. \textsl{Hit:} 16 (2d10 + 5) slashing damage plus 10 (3d6) fire damage.

\DndMonsterAction{Draining Grasp (Recharge 4\textendash 6)}
The paeliryon makes a Claw attack. On a hit, the target's Charisma score is reduced by 2 (1d4). This reduction lasts until the target finishes a short or long rest. A creature whose Charisma is reduced to 3 or less because of this ability, has the stunned condition for 1 minute.

\end{DndMonster}



\begin{DndMonster}[float=t]{Couatl}
\DndMonsterType{Medium celestial, lawful good}
\DndMonsterBasics[
armor-class = {19 (natural armor)},
hit-points = {\DndDice{13d8 + 39}},
speed = {30 ft., fly 90 ft.}
]
\DndMonsterAbilityScores[
 str = 16,
 dex = 20,
 con = 17,
 int = 18,
 wis = 20,
 cha = 18
]
\DndMonsterDetails[
saving-throws = {Con +5, Wis +7, Cha +6},
damage-resistances = {radiant},
damage-immunities = {psychic; bludgeoning, piercing, slashing from nonmagical attacks},
senses = {truesight 120 ft., passive perception 15},
languages = {all, telepathy 120 ft.},
challenge = 4,
]
\DndMonsterAction{Magic Weapons}
The couatl's weapon attacks are magical.

\DndMonsterAction{Shielded Mind}
The couatl is immune to scrying and to any effect that would sense its emotions, read its thoughts, or detect its location.

\DndMonsterAction{Innate Spellcasting}
The couatl's spellcasting ability is Charisma (spell save DC 14). It can innately cast the following spells, requiring only verbal components:
\begin{DndMonsterSpells}
  \DndInnateSpellLevel{\textit{detect evil and good}, \textit{detect magic}, \textit{detect thoughts}}
  \DndInnateSpellLevel[3e]{\textit{bless}, \textit{create food and water}, \textit{cure wounds}, \textit{lesser restoration}, \textit{protection from poison}, \textit{sanctuary}, \textit{shield}}
  \DndInnateSpellLevel[1e]{\textit{dream}, \textit{greater restoration}, \textit{scrying}}
\end{DndMonsterSpells}
\DndMonsterSection{Actions}
\DndMonsterAction{Bite}
\textsl{Melee Weapon Attack:} 8 to hit, reach 5 ft., one creature. \textsl{Hit:} 8 (1d6 + 5) piercing damage, and the target must succeed on a DC 13 Constitution saving throw or be poisoned for 24 hours. Until this poison ends, the target is unconscious. Another creature can use an action to shake the target awake.

\DndMonsterAction{Constrict}
\textsl{Melee Weapon Attack:} 6 to hit, reach 10 ft., one Medium or smaller creature. \textsl{Hit:} 10 (2d6 + 3) bludgeoning damage, and the target is grappled (escape DC 15). Until this grapple ends, the target is restrained, and the couatl can't constrict another target.

\DndMonsterAction{Change Shape}
The couatl magically polymorphs into a humanoid or beast that has a challenge rating equal to or less than its own, or back into its true form. It reverts to its true form if it dies. Any equipment it is wearing or carrying is absorbed or borne by the new form (the couatl's choice).

In a new form, the couatl retains its game statistics and ability to speak, but its AC, movement modes, Strength, Dexterity, and other actions are replaced by those of the new form, and it gains any statistics and capabilities (except class features, legendary actions, and lair actions) that the new form has but that it lacks. If the new form has a bite attack, the couatl can use its bite in that form.

\end{DndMonster}



\begin{DndMonster}[float=t]{Cult Fanatic}
\DndMonsterType{Medium humanoid (any race), any non-good alignment}
\DndMonsterBasics[
armor-class = {13 (\textit{leather armor})},
hit-points = {\DndDice{6d8 + 6}},
speed = {30 ft.}
]
\DndMonsterAbilityScores[
 str = 11,
 dex = 14,
 con = 12,
 int = 10,
 wis = 13,
 cha = 14
]
\DndMonsterDetails[
skills = {Deception +4, Persuasion +4, Religion +2},
languages = {any one language (usually Common)},
challenge = 2,
]
\DndMonsterAction{Dark Devotion}
The fanatic has advantage on saving throws against being charmed or frightened.

\DndMonsterAction{Spellcasting}
The fanatic is a 4th-level spellcaster. Its spellcasting ability is Wisdom (spell save DC 11, 3 to hit with spell attacks). The fanatic has the following cleric spells prepared:
\begin{DndMonsterSpells}
  \DndMonsterSpellLevel{\textit{light}, \textit{sacred flame}, \textit{thaumaturgy}}
   \DndMonsterSpellLevel[1][4]{\textit{command}, \textit{inflict wounds}, \textit{shield of faith}}
   \DndMonsterSpellLevel[2][3]{\textit{hold person}, \textit{spiritual weapon}}
\end{DndMonsterSpells}
\DndMonsterSection{Actions}
\DndMonsterAction{Multiattack}
The fanatic makes two melee attacks.

\DndMonsterAction{Dagger}
\textsl{Melee or Ranged Weapon Attack:} 4 to hit, reach 5 ft. or range 20/60 ft., one creature. \textsl{Hit:} 4 (1d4 + 2) piercing damage.

\end{DndMonster}



\begin{DndMonster}[float=t]{Death Knight}
\DndMonsterType{Medium undead, chaotic evil}
\DndMonsterBasics[
armor-class = {20 (\textit{plate armor})},
hit-points = {\DndDice{19d8 + 95}},
speed = {30 ft.}
]
\DndMonsterAbilityScores[
 str = 20,
 dex = 11,
 con = 20,
 int = 12,
 wis = 16,
 cha = 18
]
\DndMonsterDetails[
saving-throws = {Dex +6, Wis +9, Cha +10},
damage-immunities = {necrotic, poison},
condition-immunities = {exhaustion, frightened, poisoned},
senses = {darkvision 120 ft., passive perception 13},
languages = {Abyssal, Common},
challenge = 17,
]
\DndMonsterAction{Magic Resistance}
The death knight has advantage on saving throws against spells and other magical effects.

\DndMonsterAction{Marshal Undead}
Unless the death knight is incapacitated, it and undead creatures of its choice within 60 feet of it have advantage on saving throws against features that turn undead.

\DndMonsterAction{Spellcasting}
The death knight is a 19th-level spellcaster. Its spellcasting ability is Charisma (spell save DC 18, 10 to hit with spell attacks). It has the following paladin spells prepared:
\begin{DndMonsterSpells}
   \DndMonsterSpellLevel[1][4]{\textit{command}, \textit{compelled duel}, \textit{searing smite}}
   \DndMonsterSpellLevel[2][3]{\textit{hold person}, \textit{magic weapon}}
   \DndMonsterSpellLevel[3][3]{\textit{dispel magic}, \textit{elemental weapon}}
   \DndMonsterSpellLevel[4][3]{\textit{banishment}, \textit{staggering smite}}
   \DndMonsterSpellLevel[5][2]{\textit{destructive wave} (necrotic)}
\end{DndMonsterSpells}
\DndMonsterSection{Actions}
\DndMonsterAction{Multiattack}
The death knight makes three longsword attacks.

\DndMonsterAction{Longsword}
\textsl{Melee Weapon Attack:} 11 to hit, reach 5 ft., one target. \textsl{Hit:} 9 (1d8 + 5) slashing damage, or 10 (1d10 + 5) slashing damage if used with two hands, plus 18 (4d8) necrotic damage.

\DndMonsterAction{Hellfire Orb (1/Day)}
The death knight hurls a magical ball of fire that explodes at a point it can see within 120 feet of it. Each creature in a 20-foot-radius sphere centered on that point must make a DC 18 Dexterity saving throw. The sphere spreads around corners. A creature takes 35 (10d6) fire damage and 35 (10d6) necrotic damage on a failed save, or half as much damage on a successful one.

\DndMonsterSection{Reactions}
\DndMonsterAction{Parry}
The death knight adds 6 to its AC against one melee attack that would hit it. To do so, the death knight must see the attacker and be wielding a melee weapon.

\end{DndMonster}



\begin{DndMonster}[float=t]{Demon Ichor}
\DndMonsterType{Large ooze, unaligned}
\DndMonsterBasics[
armor-class = {7},
hit-points = {\DndDice{10d10 + 30}},
speed = {20 ft., climb 20 ft.}
]
\DndMonsterAbilityScores[
 str = 16,
 dex = 5,
 con = 16,
 int = 1,
 wis = 6,
 cha = 1
]
\DndMonsterDetails[
damage-immunities = {acid, cold, lightning, slashing},
condition-immunities = {blinded, charmed, deafened, exhaustion, frightened, prone},
senses = {blindsight 60 ft. (blind beyond this radius), passive perception 8},
challenge = 4,
]
\DndMonsterAction{Copy Warning}
This content has been copied from another creature, but the data replacement hasn't been run. Check details.

\DndMonsterAction{Amorphous}
The pudding can move through a space as narrow as 1 inch wide without squeezing.

\DndMonsterAction{Corrosive Form}
A creature that touches the pudding or hits it with a melee attack while within 5 feet of it takes 4 (1d8) acid damage. Any nonmagical weapon made of metal or wood that hits the pudding corrodes. After dealing damage, the weapon takes a permanent and cumulative −1 penalty to damage rolls. If its penalty drops to −5, the weapon is destroyed. Nonmagical ammunition made of metal or wood that hits the pudding is destroyed after dealing damage. The pudding can eat through 2-inch-thick, nonmagical wood or metal in 1 round.

\DndMonsterAction{Spider Climb}
The pudding can climb difficult surfaces, including upside down on ceilings, without needing to make an ability check.

\DndMonsterSection{Actions}
\DndMonsterAction{Pseudopod}
\textsl{Melee Weapon Attack:} 5 to hit, reach 5 ft., one target. \textsl{Hit:} 6 (1d6 + 3) bludgeoning damage plus 18 (4d8) acid damage. In addition, nonmagical armor worn by the target is partly dissolved and takes a permanent and cumulative −1 penalty to the AC it offers. The armor is destroyed if the penalty reduces its AC to 10.

\DndMonsterSection{Reactions}
\DndMonsterAction{Split}
When a pudding that is Medium or larger is subjected to lightning or slashing damage, it splits into two new puddings if it has at least 10 hit points. Each new pudding has hit points equal to half the original pudding's, rounded down. New puddings are one size smaller than the original pudding.

\end{DndMonster}



\begin{DndMonster}[float=t]{Deva}
\DndMonsterType{Medium celestial, lawful good}
\DndMonsterBasics[
armor-class = {17 (natural armor)},
hit-points = {\DndDice{16d8 + 64}},
speed = {30 ft., fly 90 ft.}
]
\DndMonsterAbilityScores[
 str = 18,
 dex = 18,
 con = 18,
 int = 17,
 wis = 20,
 cha = 20
]
\DndMonsterDetails[
saving-throws = {Wis +9, Cha +9},
skills = {Insight +9, Perception +9},
damage-resistances = {radiant; bludgeoning, piercing, slashing from nonmagical attacks},
condition-immunities = {charmed, exhaustion, frightened},
senses = {darkvision 120 ft., passive perception 19},
languages = {all, telepathy 120 ft.},
challenge = 10,
]
\DndMonsterAction{Angelic Weapons}
The deva's weapon attacks are magical. When the deva hits with any weapon, the weapon deals an extra 4d8 radiant damage (included in the attack).

\DndMonsterAction{Magic Resistance}
The deva has advantage on saving throws against spells and other magical effects.

\DndMonsterAction{Innate Spellcasting}
The deva's spellcasting ability is Charisma (spell save DC 17). The deva can innately cast the following spells, requiring only verbal components:
\begin{DndMonsterSpells}
  \DndInnateSpellLevel{\textit{detect evil and good}}
  \DndInnateSpellLevel[1e]{\textit{commune}, \textit{raise dead}}
\end{DndMonsterSpells}
\DndMonsterSection{Actions}
\DndMonsterAction{Multiattack}
The deva makes two melee attacks.

\DndMonsterAction{Mace}
\textsl{Melee Weapon Attack:} 8 to hit, reach 5 ft., one target. \textsl{Hit:} 7 (1d6 + 4) bludgeoning damage plus 18 (4d8) radiant damage.

\DndMonsterAction{Healing Touch (3/Day)}
The deva touches another creature. The target magically regains 20 (4d8 + 2) hit points and is freed from any curse, disease, poison, blindness, or deafness.

\DndMonsterAction{Change Shape}
The deva magically polymorphs into a humanoid or beast that has a challenge rating equal to or less than its own, or back into its true form. It reverts to its true form if it dies. Any equipment it is wearing or carrying is absorbed or borne by the new form (the deva's choice).

In a new form, the deva retains its game statistics and ability to speak, but its AC, movement modes, Strength, Dexterity, and special senses are replaced by those of the new form, and it gains any statistics and capabilities (except class features, legendary actions, and lair actions) that the new form has but that it lacks.

\end{DndMonster}



\begin{DndMonster}[float=t]{Devorastus}
\DndMonsterType{Huge fiend (demon), unaligned}
\DndMonsterBasics[
armor-class = {14 (natural armor)},
hit-points = {\DndDice{22d12 + 132}},
speed = {30 ft.}
]
\DndMonsterAbilityScores[
 str = 24,
 dex = 14,
 con = 22,
 int = 8,
 wis = 6,
 cha = 10
]
\DndMonsterDetails[
damage-resistances = {cold; fire; poison; bludgeoning, piercing, slashing from nonmagical attacks that aren't silvered},
damage-immunities = {acid},
condition-immunities = {blinded, charmed, deafened, exhaustion, frightened, prone},
senses = {blindsight 60 ft. (blind beyond this radius), passive perception 8},
languages = {telepathy 120 ft.},
challenge = 20,
]
\DndMonsterAction{Legendary Resistance (3/Day)}
If Devorastus fails a saving throw, it can choose to succeed instead.

\DndMonsterAction{Ooze Cube}
Devorastus takes up its entire space. Other creatures can enter the space, but a creature that does so is subjected to Engulf and has disadvantage on the saving throw. Creatures inside the cube can be seen but have total cover. A creature within 5 feet of the cube can take an action to pull a creature or object out of the cube. Doing so requires a successful DC 21 Strength check, and the creature making the attempt takes 35 (10d6) acid damage. The cube can hold only one Huge creature or up to six Large or smaller creatures inside it at a time.

\DndMonsterAction{Sticky}
Devorastus can climb difficult surfaces, including upside down on ceilings, without needing to make an ability check.

\DndMonsterAction{Transparent}
Even when Devorastus is in plain sight, it takes a successful DC 15 Wisdom (Perception) check to spot Devorastus if it has neither moved nor attacked. A creature that tries to enter Devorastus' space while unaware of Devorastus is surprised by it.

\DndMonsterSection{Actions}
\DndMonsterAction{Multiattack}
Devorastus makes four attacks using its Pseudopod, Spit, or a combination of the two.

\DndMonsterAction{Pseudopod}
\textsl{Melee Weapon Attack:} 13 to hit, reach 20 ft., one target. \textsl{Hit:} 42 (12d6) acid damage.

\DndMonsterAction{Spit}
\textsl{Ranged Weapon Attack:} 8 to hit, range 100/200 ft., one target. \textsl{Hit:} 28 (8d6) acid damage. A target hit by this attack must make a DC 21 Constitution saving throw. On a failed save, the target has the prone condition.

\DndMonsterAction{Engulf}
Devorastus moves up to its speed. While doing so, it can enter Large or smaller creatures' spaces. Whenever Devorastus enters a creature's space, the creature must make a DC 21 Dexterity saving throw. On a successful save, the creature can choose to be pushed 5 feet back or to the side of Devorastus. A creature that chooses not to be pushed suffers the consequences of a failed saving throw. On a failed save, Devorastus enters the creature's space, the creature takes 35 (10d6) acid damage, and is engulfed. The engulfed creature can't breathe, has the restrained condition, and takes 42 (12d6) acid damage at the start of each of Devorastus' turns. When Devorastus moves, the engulfed creature moves with it. An engulfed creature can try to escape by making a DC 21 Strength check as an action. On a success, the creature escapes and enters a space of its choice within 5 feet of Devorastus.

\end{DndMonster}



\begin{DndMonster}[float*=b,width=\textwidth + 8pt]{Dispater}
\begin{multicols}{2}
\DndMonsterType{Large fiend (devil), lawful evil}
\DndMonsterBasics[
armor-class = {21 (natural armor)},
hit-points = {\DndDice{33d10 + 231}},
speed = {30 ft.}
]
\DndMonsterAbilityScores[
 str = 28,
 dex = 17,
 con = 25,
 int = 25,
 wis = 22,
 cha = 24
]
\DndMonsterDetails[
saving-throws = {Str +17, Dex +11, Con +15, Int +15},
skills = {Arcana +23, Insight +22, Intimidation +15, Perception +14, Persuasion +15, Stealth +11},
damage-resistances = {cold; bludgeoning, piercing, slashing from nonmagical attacks that aren't silvered},
damage-immunities = {fire, poison},
condition-immunities = {charmed, frightened, poisoned},
senses = {darkvision 120 ft., passive perception 24},
languages = {Celestial, Common, Draconic, Infernal, telepathy 120 ft.},
challenge = 27,
]
\DndMonsterAction{Devil's Sight}
Magical darkness doesn't impede Dispater's darkvision.

\DndMonsterAction{Fear Aura}
When a creature starts their turn within 120 feet of Dispater, they must succeed on a DC 22 Wisdom saving throw or have the frightened condition until they leave the aura. A creature that succeeds on the saving throw is immune to this effect for 1 hour.

\DndMonsterAction{Fiendish Regeneration}
Dispater regains 20 hit points at the start of his turn. If he takes radiant damage this trait doesn't function at the start of his next turn. Dispater dies only if he starts his turn with 0 hit points and is unable to regenerate.

\DndMonsterAction{Legendary Resistance (3/Day)}
If Dispater fails a saving throw, he can choose to succeed instead.

\DndMonsterAction{Magic Resistance}
Dispater has advantage on saving throws against spells and other magical effects.

\DndMonsterAction{Rust Metal}
Any nonmagical weapon or ammunition made of metal that hits Dispater corrodes. After its hit, the weapon or ammunition is destroyed.

\DndMonsterSection{Actions}
\DndMonsterAction{Multiattack}
Dispater makes four attacks using his Wrought- Iron Tower, Iron Column, or a combination of the two. He can replace one of the attacks with Superheat Metal (if available).

\DndMonsterAction{Wrought-Iron Tower}
\textsl{Melee Weapon Attack:} 20 to hit, reach 10 ft., one target. \textsl{Hit:} 21 (2d8 + 12) bludgeoning damage plus 7 (2d6) fire damage. Dispater may replace the bludgeoning damage with damage of a type that the target is vulnerable to instead.

\DndMonsterAction{Iron Column}
\textsl{Ranged Spell Attack:} 15 to hit, range 60 ft., one target. \textsl{Hit:} 21 (2d8 + 12) bludgeoning damage plus 7 (2d6) fire damage. Dispater may replace the bludgeoning damage with damage of a type that the target is vulnerable to instead.

\DndMonsterAction{Superheat Metal (Recharge 4\textendash 6)}
Dispater targets a metal object he can see within 90 feet of him. The metal glows white-hot, burning all that touches it. If the object was held by a creature, the creature must succeed on a DC 23 Dexterity saving throw to drop the object or take 45 (10d8) fire damage. If the object is attached to a creature, or they couldn't feasibly drop it, they automatically fail the save. The object returns to room temperature at the end of Dispater's turn.

\DndMonsterAction{Nail Spray (2/Day)}
Dispater conjures a storm of iron nails and launches them forward in a 90-foot-long, 5-foot-wide line. All creatures within the line must make a DC 23 Dexterity saving throw, taking 40 (16d4) piercing damage on a failed save, or half as much damage on a successful one. Nails launched from this ability stick into soft surfaces or fall to the ground after the attack finishes.

\DndMonsterSection{Legendary Actions}
Dispater can take 3 legendary actions, choosing from the options below. Only one legendary action can be used at a time and only at the end of another creature's turn. Dispater regains spent legendary actions at the start of its turn.

\begin{DndMonsterLegendaryActions}
\DndMonsterLegendaryAction{Staff}{Dispater makes a Wrought-Iron Tower attack.
}
\DndMonsterLegendaryAction{Flesh to Iron (Costs 2 Actions)}{\textsl{Melee Spell Attack:} 15 to hit, reach 5 ft., one target. \textsl{Hit:} 40 (6d10 + 7) force damage, and the target must make a DC 23 Constitution saving throw. On a failed save, the target's flesh begins to harden, and the creature's movement speed is halved for 1 minute. If the creature is harmed by Flesh to Iron again, while still partly iron, the rest of the creature also turns to iron, killing them.
}
\DndMonsterLegendaryAction{Call Underling (Costs 3 Actions)}{Dispater summons an allied \textbf{erinyes} in an unoccupied space that he can see.
}
\end{DndMonsterLegendaryActions}
\DndMonsterSection{Regional Effects}
The region containing Dispater's lair is corrupted by his presence, which creates the following effect:


\begin{description}
\item[Din.]
Within 6 miles of the lair, the sound of metal working is near inescapable. Any creature attempting a short or long rest may do so, but must make a DC 21 Constitution saving throw afterwards. Failure negates any benefits of the rest.

\end{description}

If Dispater dies, the effects fade over the course of 1d10 days.

\end{multicols}
\end{DndMonster}



\begin{DndMonster}[float=t]{Displacer Fiend}
\DndMonsterType{Large fiend, neutral evil}
\DndMonsterBasics[
armor-class = {16 (natural armor)},
hit-points = {\DndDice{18d10 + 36}},
speed = {40 ft.}
]
\DndMonsterAbilityScores[
 str = 17,
 dex = 19,
 con = 14,
 int = 8,
 wis = 13,
 cha = 10
]
\DndMonsterDetails[
saving-throws = {Str +7, Dex +8},
skills = {Acrobatics +8, Athletics +7, Perception +5, Stealth +12},
damage-resistances = {fire, poison},
condition-immunities = {charmed, frightened},
senses = {darkvision 120 ft., passive perception 15},
challenge = 9,
]
\DndMonsterAction{Avoidance}
If the displacer fiend is subjected to an effect that allows it to make a saving throw to take only half damage, it instead takes no damage if it succeeds on the saving throw, and only half damage if it fails.

\DndMonsterAction{Displacement}
The displacer fiend projects a magical illusion that makes it appear to be standing near its actual location, causing attack rolls against it to have disadvantage. If it is hit by an attack, this trait is disrupted until the end of its next turn. This trait is also disrupted while the displacer fiend is incapacitated or has a speed of 0.

\DndMonsterSection{Actions}
\DndMonsterAction{Multiattack}
The displacer fiend makes three Tentacle attacks. It can replace one of the attacks with a Bite attack.

\DndMonsterAction{Tentacle}
\textsl{Melee Weapon Attack:} 8 to hit, reach 5 ft., one target. \textsl{Hit:} 18 (4d6 + 4) piercing damage plus 7 (2d6) necrotic damage.

\DndMonsterAction{Bite}
\textsl{Melee Weapon Attack:} 8 to hit, reach 10 ft., one target. \textsl{Hit:} 23 (3d12 + 4) piercing damage, and the target's Dexterity score is reduced by 2 (1d4). The target is killed if this reduces its Dexterity to 0. The reduction lasts until the target finishes a short or long rest.

\DndMonsterSection{Bonus Actions}
\DndMonsterAction{Shadowhop}
While in dim light or darkness, the displacer fiend can teleport up to 60 feet, provided its destination is also in dim light or darkness.

\end{DndMonster}



\begin{DndMonster}[float=t]{Djinni}
\DndMonsterType{Large elemental, chaotic good}
\DndMonsterBasics[
armor-class = {17 (natural armor)},
hit-points = {\DndDice{14d10 + 84}},
speed = {30 ft., fly 90 ft.}
]
\DndMonsterAbilityScores[
 str = 21,
 dex = 15,
 con = 22,
 int = 15,
 wis = 16,
 cha = 20
]
\DndMonsterDetails[
saving-throws = {Dex +6, Wis +7, Cha +9},
damage-immunities = {lightning, thunder},
senses = {darkvision 120 ft., passive perception 13},
languages = {Auran},
challenge = 11,
]
\DndMonsterAction{Elemental Demise}
If the djinni dies, its body disintegrates into a warm breeze, leaving behind only equipment the djinni was wearing or carrying.

\DndMonsterAction{Innate Spellcasting}
The djinni's innate spellcasting ability is Charisma (spell save DC 17, 9 to hit with spell attacks). It can innately cast the following spells, requiring no material components:
\begin{DndMonsterSpells}
  \DndInnateSpellLevel{\textit{detect evil and good}, \textit{detect magic}, \textit{thunderwave}}
  \DndInnateSpellLevel[3e]{\textit{create food and water} (can create wine instead of water), \textit{tongues}, \textit{wind walk}}
  \DndInnateSpellLevel[1e]{\textit{conjure elemental} (\textbf{air elemental} only), \textit{creation}, \textit{gaseous form}, \textit{invisibility}, \textit{major image}, \textit{plane shift}}
\end{DndMonsterSpells}
\DndMonsterSection{Actions}
\DndMonsterAction{Multiattack}
The djinni makes three scimitar attacks.

\DndMonsterAction{Scimitar}
\textsl{Melee Weapon Attack:} 9 to hit, reach 5 ft., one target. \textsl{Hit:} 12 (2d6 + 5) slashing damage plus 3 (1d6) lightning or thunder damage (djinni's choice).

\DndMonsterAction{Create Whirlwind}
A 5-foot-radius, 30-foot-tall cylinder of swirling air magically forms on a point the djinni can see within 120 feet of it. The whirlwind lasts as long as the djinni maintains concentration (as if concentration on a spell). Any creature but the djinni that enters the whirlwind must succeed on a DC 18 Strength saving throw or be restrained by it. The djinni can move the whirlwind up to 60 feet as an action, and creatures restrained by the whirlwind move with it. The whirlwind ends if the djinni loses sight of it.

A creature can use its action to free a creature restrained by the whirlwind, including itself, by succeeding on a DC 18 Strength check. If the check succeeds, the creature is no longer restrained and moves to the nearest space outside the whirlwind.

\end{DndMonster}



\begin{DndMonster}[float=t]{Doppelganger}
\DndMonsterType{Medium monstrosity (shapechanger), neutral}
\DndMonsterBasics[
armor-class = {14},
hit-points = {\DndDice{8d8 + 16}},
speed = {30 ft.}
]
\DndMonsterAbilityScores[
 str = 11,
 dex = 18,
 con = 14,
 int = 11,
 wis = 12,
 cha = 14
]
\DndMonsterDetails[
skills = {Deception +6, Insight +3},
condition-immunities = {charmed},
senses = {darkvision 60 ft., passive perception 11},
languages = {Common},
challenge = 3,
]
\DndMonsterAction{Shapechanger}
The doppelganger can use its action to polymorph into a Small or Medium humanoid it has seen, or back into its true form. Its statistics, other than its size, are the same in each form. Any equipment it is wearing or carrying isn't transformed. It reverts to its true form if it dies.

\DndMonsterAction{Ambusher}
In the first round of a combat, the doppelganger has advantage on attack rolls against any creature it Surprise.

\DndMonsterAction{Surprise Attack}
If the doppelganger surprises a creature and hits it with an attack during the first round of combat, the target takes an extra 10 (3d6) damage from the attack.

\DndMonsterSection{Actions}
\DndMonsterAction{Multiattack}
The doppelganger makes two melee attacks.

\DndMonsterAction{Slam}
\textsl{Melee Weapon Attack:} 6 to hit, reach 5 ft., one target. \textsl{Hit:} 7 (1d6 + 4) bludgeoning damage.

\DndMonsterAction{Read Thoughts}
The doppelganger magically reads the surface thoughts of one creature within 60 feet of it. The effect can penetrate barriers, but 3 feet of wood or dirt, 2 feet of stone, 2 inches of metal, or a thin sheet of lead blocks it. While the target is in range, the doppelganger can continue reading its thoughts, as long as the doppelganger's concentration isn't broken (as if concentration on a spell). While reading the target's mind, the doppelganger has advantage on Wisdom (Insight) and Charisma (Deception, Intimidation, and Persuasion) checks against the target.

\end{DndMonster}



\begin{DndMonster}[float=t]{Dragon Turtle}
\DndMonsterType{Gargantuan dragon, neutral}
\DndMonsterBasics[
armor-class = {20 (natural armor)},
hit-points = {\DndDice{22d20 + 110}},
speed = {20 ft., swim 40 ft.}
]
\DndMonsterAbilityScores[
 str = 25,
 dex = 10,
 con = 20,
 int = 10,
 wis = 12,
 cha = 12
]
\DndMonsterDetails[
saving-throws = {Dex +6, Con +11, Wis +7},
damage-resistances = {fire},
senses = {darkvision 120 ft., passive perception 11},
languages = {Aquan, Draconic},
challenge = 17,
]
\DndMonsterAction{Amphibious}
The dragon turtle can breathe air and water.

\DndMonsterSection{Actions}
\DndMonsterAction{Multiattack}
The dragon turtle makes three attacks: one with its bite and two with its claws. It can make one tail attack in place of its two claw attacks.

\DndMonsterAction{Bite}
\textsl{Melee Weapon Attack:} 13 to hit, reach 15 ft., one target. \textsl{Hit:} 26 (3d12 + 7) piercing damage.

\DndMonsterAction{Claw}
\textsl{Melee Weapon Attack:} 13 to hit, reach 10 ft., one target. \textsl{Hit:} 16 (2d8 + 7) slashing damage.

\DndMonsterAction{Tail}
\textsl{Melee Weapon Attack:} 13 to hit, reach 15 ft., one target. \textsl{Hit:} 26 (3d12 + 7) bludgeoning damage. If the target is a creature, it must succeed on a DC 20 Strength saving throw or be pushed up to 10 feet away from the dragon turtle and knocked prone.

\DndMonsterAction{Steam Breath (Recharge 5\textendash 6)}
The dragon turtle exhales scalding steam in a 60-foot cone. Each creature in that area must make a DC 18 Constitution saving throw, taking 52 (15d6) fire damage on a failed save, or half as much damage on a successful one. Being underwater doesn't grant resistance against this damage.

\end{DndMonster}


\clearpage

\begin{DndMonster}[float=t]{Dretch}
\DndMonsterType{Small fiend (demon), chaotic evil}
\DndMonsterBasics[
armor-class = {11 (natural armor)},
hit-points = {\DndDice{4d6 + 4}},
speed = {20 ft.}
]
\DndMonsterAbilityScores[
 str = 11,
 dex = 11,
 con = 12,
 int = 5,
 wis = 8,
 cha = 3
]
\DndMonsterDetails[
damage-resistances = {cold, fire, lightning},
damage-immunities = {poison},
condition-immunities = {poisoned},
senses = {darkvision 60 ft., passive perception 9},
languages = {Abyssal, telepathy 60 ft. (works only with creatures that understand Abyssal)},
challenge = 1/4,
]
\DndMonsterSection{Actions}
\DndMonsterAction{Multiattack}
The dretch makes two attacks: one with its bite and one with its claws.

\DndMonsterAction{Bite}
\textsl{Melee Weapon Attack:} 2 to hit, reach 5 ft., one target. \textsl{Hit:} 3 (1d6) piercing damage.

\DndMonsterAction{Claws}
\textsl{Melee Weapon Attack:} 2 to hit, reach 5 ft., one target. \textsl{Hit:} 5 (2d4) slashing damage.

\DndMonsterAction{Fetid Cloud (1/Day)}
A 10-foot radius of disgusting green gas extends out from the dretch. The gas spreads around corners, and its area is lightly obscured. It lasts for 1 minute or until a strong wind disperses it. Any creature that starts its turn in that area must succeed on a DC 11 Constitution saving throw or be poisoned until the start of its next turn. While poisoned in this way, the target can take either an action or a bonus action on its turn, not both, and can't take reactions.

\end{DndMonster}



\begin{DndMonster}[float=t]{Drow Mage}
\DndMonsterType{Medium humanoid (elf), neutral evil}
\DndMonsterBasics[
armor-class = {12 (15 with \textit{mage armor})},
hit-points = {\DndDice{10d8}},
speed = {30 ft.}
]
\DndMonsterAbilityScores[
 str = 9,
 dex = 14,
 con = 10,
 int = 17,
 wis = 13,
 cha = 12
]
\DndMonsterDetails[
skills = {Arcana +6, Deception +4, Perception +4, Stealth +5},
senses = {darkvision 120 ft., passive perception 14},
languages = {Elvish, Undercommon},
challenge = 7,
]
\DndMonsterAction{Fey Ancestry}
The drow has advantage on saving throws against being charmed, and magic can't put the drow to sleep.

\DndMonsterAction{Sunlight Sensitivity}
While in sunlight, the drow has disadvantage on attack rolls, as well as on Wisdom (Perception) checks that rely on sight.

\DndMonsterAction{Innate Spellcasting}
The drow's spellcasting ability is Charisma (spell save DC 12). It can innately cast the following spells, requiring no material components:
\begin{DndMonsterSpells}
  \DndInnateSpellLevel{\textit{dancing lights}}
  \DndInnateSpellLevel[1e]{\textit{darkness}, \textit{faerie fire}, \textit{levitate} (self only)}
\end{DndMonsterSpells}
\DndMonsterAction{Spellcasting}
The drow is a 10th-level spellcaster. Its spellcasting ability is Intelligence (spell save DC 14, 6 to hit with spell attacks). The drow has the following wizard spells prepared:
\begin{DndMonsterSpells}
  \DndMonsterSpellLevel{\textit{mage hand}, \textit{minor illusion}, \textit{poison spray}, \textit{ray of frost}}
   \DndMonsterSpellLevel[1][4]{\textit{mage armor}, \textit{magic missile}, \textit{shield}, \textit{witch bolt}}
   \DndMonsterSpellLevel[2][3]{\textit{alter self}, \textit{misty step}, \textit{web}}
   \DndMonsterSpellLevel[3][3]{\textit{fly}, \textit{lightning bolt}}
   \DndMonsterSpellLevel[4][3]{\textit{Evard's black tentacles}, \textit{greater invisibility}}
   \DndMonsterSpellLevel[5][2]{\textit{cloudkill}}
\end{DndMonsterSpells}
\DndMonsterSection{Actions}
\DndMonsterAction{Staff}
\textsl{Melee Weapon Attack:} 2 to hit, reach 5 ft., one target. \textsl{Hit:} 2 (1d6 - 1) bludgeoning damage, or 3 (1d8 - 1) bludgeoning damage if used with two hands, plus 3 (1d6) poison damage.

\DndMonsterAction{Summon Demon (1/Day)}
The drow magically summons a \textbf{quasit}, or attempts to summon a \textbf{shadow demon} with a 50\% summoning chance chance of success. The summoned demon appears in an unoccupied space within 60 feet of its summoner, acts as an ally of its summoner, and can't summon other demons. It remains for 10 minutes, until it or its summoner dies, or until its summoner dismisses it as an action.

\end{DndMonster}



\begin{DndMonster}[float=t]{Druid}
\DndMonsterType{Medium humanoid (any race), any alignment}
\DndMonsterBasics[
armor-class = {11 (16 with \textit{barkskin})},
hit-points = {\DndDice{5d8 + 5}},
speed = {30 ft.}
]
\DndMonsterAbilityScores[
 str = 10,
 dex = 12,
 con = 13,
 int = 12,
 wis = 15,
 cha = 11
]
\DndMonsterDetails[
skills = {Medicine +4, Nature +3, Perception +4},
languages = {Druidic plus any two languages},
challenge = 2,
]
\DndMonsterAction{Spellcasting}
The druid is a 4th-level spellcaster. Its spellcasting ability is Wisdom (spell save DC 12, 4 to hit with spell attacks). It has the following druid spells prepared:
\begin{DndMonsterSpells}
  \DndMonsterSpellLevel{\textit{druidcraft}, \textit{produce flame}, \textit{shillelagh}}
   \DndMonsterSpellLevel[1][4]{\textit{entangle}, \textit{longstrider}, \textit{speak with animals}, \textit{thunderwave}}
   \DndMonsterSpellLevel[2][3]{\textit{animal messenger}, \textit{barkskin}}
\end{DndMonsterSpells}
\DndMonsterSection{Actions}
\DndMonsterAction{Quarterstaff}
\textsl{Melee Weapon Attack:} 2 to hit (4 to hit with shillelagh), reach 5 ft., one target. \textsl{Hit:} 3 (1d6) bludgeoning damage, 4 (1d8) bludgeoning damage if wielded with two hands, or 6 (1d8 + 2) bludgeoning damage with \textit{shillelagh}.

\end{DndMonster}



\begin{DndMonster}[float=t]{Duergar}
\DndMonsterType{Medium humanoid (dwarf), lawful evil}
\DndMonsterBasics[
armor-class = {16 (\textit{scale mail})},
hit-points = {\DndDice{4d8 + 4}},
speed = {25 ft.}
]
\DndMonsterAbilityScores[
 str = 14,
 dex = 11,
 con = 14,
 int = 11,
 wis = 10,
 cha = 9
]
\DndMonsterDetails[
damage-resistances = {poison},
senses = {darkvision 120 ft., passive perception 10},
languages = {Dwarvish, Undercommon},
challenge = 1,
]
\DndMonsterAction{Duergar Resilience}
The duergar has advantage on saving throws against poison, spells, and illusions, as well as to resist being charmed or paralyzed.

\DndMonsterAction{Sunlight Sensitivity}
While in sunlight, the duergar has disadvantage on attack rolls, as well as on Wisdom (Perception) checks that rely on sight.

\DndMonsterSection{Actions}
\DndMonsterAction{Enlarge (Recharges after a Short or Long Rest)}
For 1 minute, the duergar magically increases in size, along with anything it is wearing or carrying. While enlarged, the duergar is Large, doubles its damage dice on Strength-based weapon attacks (included in the attacks), and makes Strength checks and Strength saving throws with advantage. If the duergar lacks the room to become Large, it attains the maximum size possible in the space available.

\DndMonsterAction{War Pick}
\textsl{Melee Weapon Attack:} 4 to hit, reach 5 ft., one target. \textsl{Hit:} 6 (1d8 + 2) piercing damage, or 11 (2d8 + 2) piercing damage while enlarged.

\DndMonsterAction{Javelin}
\textsl{Melee or Ranged Weapon Attack:} 4 to hit, reach 5 ft. or range 30/120 ft., one target. \textsl{Hit:} 5 (1d6 + 2) piercing damage, or 9 (2d6 + 2) piercing damage while enlarged.

\DndMonsterAction{Invisibility (Recharges after a Short or Long Rest)}
The duergar magically turns invisible until it attacks, casts a spell, or uses its Enlarge, or until its concentration is broken, up to 1 hour (as if concentration on a spell). Any equipment the duergar wears or carries is invisible with it.

\end{DndMonster}



\begin{DndMonster}[float=t]{Efreeti}
\DndMonsterType{Large elemental, lawful evil}
\DndMonsterBasics[
armor-class = {17 (natural armor)},
hit-points = {\DndDice{16d10 + 112}},
speed = {40 ft., fly 60 ft.}
]
\DndMonsterAbilityScores[
 str = 22,
 dex = 12,
 con = 24,
 int = 16,
 wis = 15,
 cha = 16
]
\DndMonsterDetails[
saving-throws = {Int +7, Wis +6, Cha +7},
damage-immunities = {fire},
senses = {darkvision 120 ft., passive perception 12},
languages = {Ignan},
challenge = 11,
]
\DndMonsterAction{Elemental Demise}
If the efreeti dies, its body disintegrates in a flash of fire and puff of smoke, leaving behind only equipment the efreeti was wearing or carrying.

\DndMonsterAction{Innate Spellcasting}
The efreeti's innate spellcasting ability is Charisma (spell save DC 15, 7 to hit with spell attacks). It can innately cast the following spells, requiring no material components:
\begin{DndMonsterSpells}
  \DndInnateSpellLevel{\textit{detect magic}}
  \DndInnateSpellLevel[3e]{\textit{enlarge/reduce}, \textit{tongues}}
  \DndInnateSpellLevel[1e]{\textit{conjure elemental} (\textbf{fire elemental} only), \textit{gaseous form}, \textit{invisibility}, \textit{major image}, \textit{plane shift}, \textit{wall of fire}}
\end{DndMonsterSpells}
\DndMonsterSection{Actions}
\DndMonsterAction{Multiattack}
The efreeti makes two scimitar attacks or uses its Hurl Flame twice.

\DndMonsterAction{Scimitar}
\textsl{Melee Weapon Attack:} 10 to hit, reach 5 ft., one target. \textsl{Hit:} 13 (2d6 + 6) slashing damage plus 7 (2d6) fire damage.

\DndMonsterAction{Hurl Flame}
\textsl{Ranged Spell Attack:} 7 to hit, range 120 ft., one target. \textsl{Hit:} 17 (5d6) fire damage.

\end{DndMonster}



\begin{DndMonster}[float=t]{Ekengarik}
\DndMonsterType{Large fiend, lawful evil}
\DndMonsterBasics[
armor-class = {16 (natural armor)},
hit-points = {\DndDice{22d10 + 110}},
speed = {30 ft.}
]
\DndMonsterAbilityScores[
 str = 23,
 dex = 16,
 con = 21,
 int = 14,
 wis = 17,
 cha = 21
]
\DndMonsterDetails[
saving-throws = {Con +10, Cha +10},
skills = {Deception +10, Insight +8, Persuasion +10},
damage-resistances = {acid; lightning; bludgeoning, piercing, slashing from nonmagical attacks},
damage-immunities = {cold, fire, poison},
condition-immunities = {charmed, poisoned},
senses = {blindsight 30 ft., darkvision 120 ft., passive perception 13},
languages = {Common, Formian, telepathy 120 ft.},
challenge = 16,
]
\DndMonsterAction{Hive Mind}
All fiendish formians within 1 mile of Ekengarik can telepathically communicate with each other and Ekengarik.

\DndMonsterAction{Regeneration}
Ekengarik regenerates 10 hit points at the start of her turn, unless she took radiant damage in the last round.

\DndMonsterAction{Spellcasting}
Ekengarik casts one of the following spells, requiring no material components and using Charisma as the spellcasting ability (spell save DC 18):
\begin{DndMonsterSpells}
  \DndInnateSpellLevel{\textit{Detect Magic}, \textit{Dispel Magic}, \textit{Heroism}, \textit{Magic Missile}}
  \DndInnateSpellLevel[1e]{\textit{Evard's Black Tentacles}, \textit{Cone of Cold}, \textit{Confusion}, \textit{Dimension Door}, \textit{Geas}, \textit{Invisibility}, \textit{Prismatic Wall}, \textit{Slow}}
\end{DndMonsterSpells}
\DndMonsterSection{Actions}
\DndMonsterAction{Multiattack}
Ekengarik makes three Bite attacks. She can replace one of the attacks with an Acid Spray (if available) attack or a use of Spellcasting.

\DndMonsterAction{Bite}
\textsl{Melee Weapon Attack:} 11 to hit, reach 5 ft., one target. \textsl{Hit:} 15 (2d8 + 6) piercing damage and the target must make a DC 19 Constitution saving throw. On a failed save, the target's Strength score is reduced by 2 (1d4). The target dies if this reduces its Strength to 0. Otherwise, the reduction lasts until the target finishes a short or long rest.

\DndMonsterAction{Acid Spray (Recharge 5\textendash 6)}
Ekengarik spits acid at one creature within 60 feet of her, or two creatures within 60 feet of her and 5 feet of each other. Targets must make a DC 18 Dexterity saving throw, taking 33 (6d10) acid damage on a failed saving throw, or half as much on a successful one.

\DndMonsterSection{Reactions}
\DndMonsterAction{Hardened Carapace (Recharge 4\textendash 6)}
When hit with an attack, Ekengarik temporarily hardens her carapace, reducing her speed to 0 and increasing her AC by 5 until the start of her next turn.

\end{DndMonster}



\begin{DndMonster}[float*=b,width=\textwidth + 8pt]{Elder Brain}
\begin{multicols}{2}
\DndMonsterType{Large aberration (mind flayer), lawful evil}
\DndMonsterBasics[
armor-class = {10},
hit-points = {\DndDice{20d10 + 100}},
speed = {5 ft., swim 10 ft.}
]
\DndMonsterAbilityScores[
 str = 15,
 dex = 10,
 con = 20,
 int = 21,
 wis = 19,
 cha = 24
]
\DndMonsterDetails[
saving-throws = {Int +10, Wis +9, Cha +12},
skills = {Arcana +10, Deception +12, Insight +14, Intimidation +12, Persuasion +12},
senses = {blindsight 120 ft., passive perception 14},
languages = {understands Common, Deep Speech, and Undercommon but can't speak, telepathy 5 miles},
challenge = 14,
]
\DndMonsterAction{Creature Sense}
The elder brain is aware of creatures within 5 miles of it that have an Intelligence score of 4 or higher. It knows the distance and direction to each creature, as well as each one's Intelligence score, but can't sense anything else about it. A creature protected by a \textit{mind blank} spell, a \textit{nondetection} spell, or similar magic can't be perceived in this manner.

\DndMonsterAction{Legendary Resistance (3/Day)}
If the elder brain fails a saving throw, it can choose to succeed instead.

\DndMonsterAction{Magic Resistance}
The elder brain has advantage on saving throws against spells and other magical effects.

\DndMonsterAction{Telepathic Hub}
The elder brain can use its telepathy to initiate and maintain telepathic conversations with up to ten creatures at a time. The elder brain can let those creatures telepathically hear each other while connected in this way.

\DndMonsterAction{Spellcasting (Psionics)}
The elder brain casts one of the following spells, requiring no spell components and using Intelligence as the spellcasting ability (spell save DC 18):
\begin{DndMonsterSpells}
  \DndInnateSpellLevel{\textit{detect thoughts}, \textit{levitate}}
  \DndInnateSpellLevel[3]{\textit{modify memory}}
  \DndInnateSpellLevel[1e]{\textit{dominate monster}, \textit{plane shift} (self only)}
\end{DndMonsterSpells}
\DndMonsterSection{Actions}
\DndMonsterAction{Tentacle}
\textsl{Melee Weapon Attack:} 7 to hit, reach 30 ft., one target. \textsl{Hit:} 20 (4d8 + 2) bludgeoning damage. If the target is a Huge or smaller creature, it is grappled (escape DC 15) and takes 9 (1d8 + 5) psychic damage at the start of each of its turns until the grapple ends. The elder brain can have up to four targets grappled at a time.

\DndMonsterAction{Mind Blast (Recharge 5\textendash 6)}
Creatures of the elder brain's choice within 60 feet of it must succeed on a DC 18 Intelligence saving throw or take 32 (5d10 + 5) psychic damage and be stunned for 1 minute. A target can repeat the saving throw at the end of each of its turns, ending the effect on itself on a success.

\DndMonsterSection{Bonus Actions}
\DndMonsterAction{Psychic Link}
The elder brain targets one incapacitated creature it senses with its Creature Sense trait and establishes a psychic link with the target. Until the link ends, the elder brain can perceive everything the target senses. The target becomes aware that something is linked to its mind once it is no longer incapacitated, and the elder brain can terminate the link at any time (no action required). The target can use an action on its turn to attempt to break the link, doing so with a successful DC 18 Charisma saving throw. On a successful save, the target takes 10 (3d6) psychic damage. The link also ends if the target and the elder brain are more than 5 miles apart. The elder brain can form psychic links with up to ten creatures at a time.

\DndMonsterAction{Sense Thoughts}
The elder brain targets a creature with which it has a psychic link. The elder brain gains insight into the target's emotional state and foremost thoughts (including worries, loves, and hates).

\DndMonsterSection{Legendary Actions}
Elder Brain can take 3 legendary actions, choosing from the options below. Only one legendary action can be used at a time and only at the end of another creature's turn. Elder Brain regains spent legendary actions at the start of its turn.

\begin{DndMonsterLegendaryActions}
\DndMonsterLegendaryAction{Break Concentration}{The elder brain targets one creature within 120 feet of it with which it has a psychic link. The elder brain breaks the creature's concentration on a spell it has cast. The creature also takes 2 (1d4) psychic damage per level of the spell.
}
\DndMonsterLegendaryAction{Psychic Pulse}{The elder brain targets one creature within 120 feet of it with which it has a psychic link. The target and enemies of the elder brain within 30 feet of target take 10 (3d6) psychic damage.
}
\DndMonsterLegendaryAction{Sever Psychic Link}{The elder brain targets one creature within 120 feet of it with which it has a psychic link. The elder brain ends the link, causing the creature to have disadvantage on all ability checks, attack rolls, and saving throws until the end of the creature's next turn.
}
\DndMonsterLegendaryAction{Tentacle (Costs 2 Actions)}{The elder brain makes one Tentacle attack.
}
\end{DndMonsterLegendaryActions}
\DndMonsterSection{Lair Actions}
On initiative count 20 (losing initiative ties), an elder brain can take one of the following lair actions; the elder brain can't take the same lair action two rounds in a row:


\begin{description}
\item[Force Wall.]
The elder brain casts \textit{wall of force}.

\item[Psionic Anchor.]
The elder brain targets one creature it can sense within 120 feet of it and anchors it by sheer force of will. The target must make a DC 18 Charisma saving throw. On a failed save, its speed is reduced to 0, and it can't teleport. It can repeat the saving throw at the end of each of its turns, ending the effect on itself on a success.

\item[Psychic Inspiration.]
The elder brain targets one friendly creature it can sense within 120 feet of it. The target has a flash of inspiration and gains advantage on one attack roll, ability check, or saving throw it makes before the end of its next turn.

\end{description}

\DndMonsterSection{Regional Effects}
The territory within 5 miles of an elder brain is altered by the creature's psionic presence, which creates one or more of the following effects:


\begin{description}
\item[Paranoia.]
Creatures within 5 miles of an elder brain feel as if they are being followed, even when they're not.

\item[Psychic Whispers.]
Any creature with which the elder brain has formed a psychic link hears faint, incomprehensible whispers in the deepest recesses of its mind. This psychic detritus consists of the elder brain's stray thoughts commingled with those of other creatures to which it is linked.

\item[Telepathic Eavesdropping.]
The elder brain can overhear any telepathic conversation within 5 miles of it. The creature that initiated the telepathic conversation makes a DC 18 Wisdom saving throw when telepathic contact is first established. If the save is successful, the creature is aware that something is eavesdropping. The nature of the eavesdropper isn't revealed.

\end{description}

If the elder brain dies, these effects immediately end.

\end{multicols}
\end{DndMonster}



\begin{DndMonster}[float=t]{Elephant}
\DndMonsterType{Huge beast, unaligned}
\DndMonsterBasics[
armor-class = {12 (natural armor)},
hit-points = {\DndDice{8d12 + 24}},
speed = {40 ft.}
]
\DndMonsterAbilityScores[
 str = 22,
 dex = 9,
 con = 17,
 int = 3,
 wis = 11,
 cha = 6
]
\DndMonsterDetails[
challenge = 4,
]
\DndMonsterAction{Trampling Charge}
If the elephant moves at least 20 feet straight toward a creature and then hits it with a gore attack on the same turn, that target must succeed on a DC 12 Strength saving throw or be knocked prone. If the target is prone, the elephant can make one stomp attack against it as a bonus action.

\DndMonsterSection{Actions}
\DndMonsterAction{Gore}
\textsl{Melee Weapon Attack:} 8 to hit, reach 5 ft., one target. \textsl{Hit:} 19 (3d8 + 6) piercing damage.

\DndMonsterAction{Stomp}
\textsl{Melee Weapon Attack:} 8 to hit, reach 5 ft., one prone creature. \textsl{Hit:} 22 (3d10 + 6) bludgeoning damage.

\end{DndMonster}



\begin{DndMonster}[float*=b,width=\textwidth + 8pt]{Empyrean}
\begin{multicols}{2}
\DndMonsterType{Huge celestial (titan), chaotic good (75\%) neutral evil (25\%)}
\DndMonsterBasics[
armor-class = {22 (natural armor)},
hit-points = {\DndDice{19d12 + 190}},
speed = {50 ft., fly 50 ft., swim 50 ft.}
]
\DndMonsterAbilityScores[
 str = 30,
 dex = 21,
 con = 30,
 int = 21,
 wis = 22,
 cha = 27
]
\DndMonsterDetails[
saving-throws = {Str +17, Int +12, Wis +13, Cha +15},
skills = {Insight +13, Persuasion +15},
damage-immunities = {bludgeoning, piercing, slashing from nonmagical attacks},
senses = {truesight 120 ft., passive perception 16},
languages = {all},
challenge = 23,
]
\DndMonsterAction{Legendary Resistance (3/Day)}
If the empyrean fails a saving throw, it can choose to succeed instead.

\DndMonsterAction{Magic Resistance}
The empyrean has advantage on saving throws against spells and other magical effects.

\DndMonsterAction{Magic Weapons}
The empyrean's weapon attacks are magical.

\DndMonsterAction{Innate Spellcasting}
The empyrean's innate spellcasting ability is Charisma (spell save DC 23, 15 to hit with spell attacks). It can innately cast the following spells, requiring no material components:
\begin{DndMonsterSpells}
  \DndInnateSpellLevel{\textit{greater restoration}, \textit{pass without trace}, \textit{water breathing}, \textit{water walk}}
  \DndInnateSpellLevel[1e]{\textit{commune}, \textit{dispel evil and good}, \textit{earthquake}, \textit{fire storm}, \textit{plane shift} (self only)}
\end{DndMonsterSpells}
\DndMonsterSection{Actions}
\DndMonsterAction{Maul}
\textsl{Melee Weapon Attack:} 17 to hit, reach 10 ft., one target. \textsl{Hit:} 31 (6d6 + 10) bludgeoning damage. If the target is a creature, it must succeed on a DC 15 Constitution saving throw or be stunned until the end of the empyrean's next turn.

\DndMonsterAction{Bolt}
\textsl{Ranged Spell Attack:} 15 to hit, range 600 ft., one target. \textsl{Hit:} 24 (7d6) damage of one of the following types (empyrean's choice): acid, cold, fire, force, lightning, radiant, or thunder.

\DndMonsterSection{Legendary Actions}
Empyrean can take 3 legendary actions, choosing from the options below. Only one legendary action can be used at a time and only at the end of another creature's turn. Empyrean regains spent legendary actions at the start of its turn.

\begin{DndMonsterLegendaryActions}
\DndMonsterLegendaryAction{Attack}{The empyrean makes one attack.
}
\DndMonsterLegendaryAction{Bolster}{The empyrean bolsters all nonhostile creatures within 120 feet of it until the end of its next turn. Bolstered creatures can't be charmed or frightened, and they gain advantage on ability checks and saving throws until the end of the empyrean's next turn.
}
\DndMonsterLegendaryAction{Trembling Strike (Costs 2 Actions)}{The empyrean strikes the ground with its maul, triggering an earth tremor. All other creatures on the ground within 60 feet of the empyrean must succeed on a DC 25 Strength saving throw or be knocked prone.
}
\end{DndMonsterLegendaryActions}
\end{multicols}
\end{DndMonster}



\begin{DndMonster}[float*=b,width=\textwidth + 8pt]{Eriflamme}
\begin{multicols}{2}
\DndMonsterType{Gargantuan elemental, unaligned}
\DndMonsterBasics[
armor-class = {19 (natural armor)},
hit-points = {\DndDice{28d20 + 224}},
speed = {40 ft.}
]
\DndMonsterAbilityScores[
 str = 30,
 dex = 11,
 con = 26,
 int = 5,
 wis = 10,
 cha = 10
]
\DndMonsterDetails[
saving-throws = {Int +5, Wis +8, Cha +8},
skills = {Arcana +5, Athletics +18},
damage-immunities = {fire; poison; bludgeoning, piercing, slashing from nonmagical attacks},
condition-immunities = {charmed, frightened, paralyzed, poisoned},
senses = {blindsight 120 ft., passive perception 10},
challenge = 28,
]
\DndMonsterAction{Legendary Resistance (3/Day)}
If the Eriflamme fails a saving throw, it can choose to succeed instead.

\DndMonsterAction{Magic Resistance}
The Eriflamme has advantage on saving throws against spells and other magical effects.

\DndMonsterSection{Actions}
\DndMonsterAction{Multiattack}
The Eriflamme makes four Claw attacks. It can replace one of the attacks with a Bite attack.

\DndMonsterAction{Claw}
\textsl{Melee Weapon Attack:} 18 to hit, reach 15 ft., one target. \textsl{Hit:} 20 (3d6 + 10) slashing damage plus 7 (2d6) fire damage.

\DndMonsterAction{Bite}
\textsl{Melee Weapon Attack:} 18 to hit, reach 10 ft., one target. \textsl{Hit:} 24 (4d6 + 10) piercing damage plus 33 (6d10) fire damage.

\DndMonsterAction{Reform}
All the Eriflamme's fire elementals vanish and the Eriflamme reappears in one of their spaces. The Eriflamme's hit points are equivalent to the combined hit points of the fire elementals.

\DndMonsterSection{Reactions}
\DndMonsterAction{Split (2/Day)}
As a reaction to taking damage, the Eriflamme can split into fire elementals and reduce the damage to 0. The Eriflamme can form up to 10 fire elementals, splitting its remaining hit points between them. The fire elementals appear in spaces adjacent to or within the Eriflamme's old location, and all act on the Eriflamme's initiative. They gain a +5 bonus to attack rolls and have the Eriflamme's Reform action.

\DndMonsterSection{Legendary Actions}
Eriflamme can take 3 legendary actions, choosing from the options below. Only one legendary action can be used at a time and only at the end of another creature's turn. Eriflamme regains spent legendary actions at the start of its turn.

\begin{DndMonsterLegendaryActions}
\DndMonsterLegendaryAction{Slam}{The Eriflamme makes a Claw attack, which deals force damage instead of slashing damage if it hits.
}
\DndMonsterLegendaryAction{Fireglide}{The Eriflamme moves up to half its speed. It can fly during this movement.
}
\DndMonsterLegendaryAction{Chomp (Costs 2 Actions)}{The Eriflamme makes a Bite attack with advantage.
}
\end{DndMonsterLegendaryActions}
\end{multicols}
\end{DndMonster}



\begin{DndMonster}[float=t]{Erinyes}
\DndMonsterType{Medium fiend (devil), lawful evil}
\DndMonsterBasics[
armor-class = {18 (\textit{plate armor})},
hit-points = {\DndDice{18d8 + 72}},
speed = {30 ft., fly 60 ft.}
]
\DndMonsterAbilityScores[
 str = 18,
 dex = 16,
 con = 18,
 int = 14,
 wis = 14,
 cha = 18
]
\DndMonsterDetails[
saving-throws = {Dex +7, Con +8, Wis +6, Cha +8},
damage-resistances = {cold; bludgeoning, piercing, slashing from nonmagical attacks that aren't silvered},
damage-immunities = {fire, poison},
condition-immunities = {poisoned},
senses = {truesight 120 ft., passive perception 12},
languages = {Infernal, telepathy 120 ft.},
challenge = 12,
]
\DndMonsterAction{Hellish Weapons}
The erinyes's weapon attacks are magical and deal an extra 13 (3d8) poison damage on a hit (included in the attacks).

\DndMonsterAction{Magic Resistance}
The erinyes has advantage on saving throws against spells and other magical effects.

\DndMonsterSection{Actions}
\DndMonsterAction{Multiattack}
The erinyes makes three attacks.

\DndMonsterAction{Longsword}
\textsl{Melee Weapon Attack:} 8 to hit, reach 5 ft., one target. \textsl{Hit:} 8 (1d8 + 4) slashing damage, or 9 (1d10 + 4) slashing damage if used with two hands, plus 13 (3d8) poison damage.

\DndMonsterAction{Longbow}
\textsl{Ranged Weapon Attack:} 7 to hit, range 150/600 ft., one target. \textsl{Hit:} 7 (1d8 + 3) piercing damage plus 13 (3d8) poison damage, and the target must succeed on a DC 14 Constitution saving throw or be poisoned. The poison lasts until it is removed by the \textit{lesser restoration} spell or similar magic.

\DndMonsterSection{Reactions}
\DndMonsterAction{Parry}
The erinyes adds 4 to its AC against one melee attack that would hit it. To do so, the erinyes must see the attacker and be wielding a melee weapon.

\end{DndMonster}



\begin{DndMonster}[float=t]{Ettin}
\DndMonsterType{Large giant, chaotic evil}
\DndMonsterBasics[
armor-class = {12 (natural armor)},
hit-points = {\DndDice{10d10 + 30}},
speed = {40 ft.}
]
\DndMonsterAbilityScores[
 str = 21,
 dex = 8,
 con = 17,
 int = 6,
 wis = 10,
 cha = 8
]
\DndMonsterDetails[
skills = {Perception +4},
senses = {darkvision 60 ft., passive perception 14},
languages = {Giant, Orc},
challenge = 4,
]
\DndMonsterAction{Two Heads}
The ettin has advantage on Wisdom (Perception) checks and on saving throws against being blinded, charmed, deafened, frightened, stunned, and knocked unconscious.

\DndMonsterAction{Wakeful}
When one of the ettin's heads is asleep, its other head is awake.

\DndMonsterSection{Actions}
\DndMonsterAction{Multiattack}
The ettin makes two attacks: one with its battleaxe and one with its morningstar.

\DndMonsterAction{Battleaxe}
\textsl{Melee Weapon Attack:} 7 to hit, reach 5 ft., one target. \textsl{Hit:} 14 (2d8 + 5) slashing damage.

\DndMonsterAction{Morningstar}
\textsl{Melee Weapon Attack:} 7 to hit, reach 5 ft., one target. \textsl{Hit:} 14 (2d8 + 5) piercing damage.

\end{DndMonster}



\begin{DndMonster}[float=t]{Fiendish Formian}
\DndMonsterType{Medium fiend, lawful evil}
\DndMonsterBasics[
armor-class = {15 (natural armor)},
hit-points = {\DndDice{6d8 + 18}},
speed = {40 ft.}
]
\DndMonsterAbilityScores[
 str = 18,
 dex = 14,
 con = 17,
 int = 10,
 wis = 12,
 cha = 11
]
\DndMonsterDetails[
saving-throws = {Str +6},
skills = {Acrobatics +4, Stealth +4},
damage-resistances = {bludgeoning, piercing, slashing from nonmagical attacks that aren't silvered},
damage-immunities = {cold, fire, poison},
condition-immunities = {charmed, poisoned},
senses = {darkvision 60 ft., passive perception 11},
languages = {Formian},
challenge = 4,
]
\DndMonsterAction{Hive Mind}
All fiendish formians within 1 mile of Ekengarik can telepathically communicate with each other and Ekengarik.

\DndMonsterSection{Actions}
\DndMonsterAction{Multiattack}
The fiendish formian makes two Claw attacks.

\DndMonsterAction{Claw}
\textsl{Melee Weapon Attack:} 6 to hit, reach 5 ft., one target. \textsl{Hit:} 7 (1d6 + 4) slashing damage.

\DndMonsterAction{Bite}
\textsl{Melee Weapon Attack:} 6 to hit, reach 5 ft., one target. \textsl{Hit:} 8 (1d8 + 4) piercing damage and the target must make a DC 14 Constitution saving throw. On a failed save, the target's Strength score is reduced by 1. The target dies if this reduces its Strength to 0. Otherwise, the reduction lasts until the target finishes a short or long rest.

\end{DndMonster}



\begin{DndMonster}[float*=b,width=\textwidth + 8pt]{Fierna}
\begin{multicols}{2}
\DndMonsterType{Medium fiend (devil), lawful evil}
\DndMonsterBasics[
armor-class = {20 (natural armor)},
hit-points = {\DndDice{35d8 + 210}},
speed = {30 ft., fly 60 ft.}
]
\DndMonsterAbilityScores[
 str = 17,
 dex = 25,
 con = 22,
 int = 22,
 wis = 22,
 cha = 29
]
\DndMonsterDetails[
saving-throws = {Dex +15, Con +14, Wis +14, Cha +17},
skills = {Acrobatics +15, Deception +25, Insight +14, Intimidation +17, Perception +14, Persuasion +25},
damage-resistances = {bludgeoning, piercing, slashing from nonmagical attacks that aren't silvered},
damage-immunities = {cold, fire, poison},
condition-immunities = {charmed, poisoned},
senses = {darkvision 120 ft., passive perception 24},
languages = {Celestial, Common, Draconic, Infernal, telepathy 120 ft.},
challenge = 25,
]
\DndMonsterAction{Fear Aura}
When a creature starts their turn within 120 feet of Fierna, they must succeed on a DC 22 Wisdom saving throw or have the frightened condition until they leave the aura. A creature that succeeds on the saving throw is immune to this effect for 1 hour.

\DndMonsterAction{Fiendish Regeneration}
Fierna regains 20 hit points at the start of her turn. If she takes radiant damage this trait doesn't function at the start of her next turn. Fierna dies only if she starts her turn with 0 hit points and is unable to regenerate.

\DndMonsterAction{Legendary Resistance (3/Day)}
If Fierna fails a saving throw, she can choose to succeed instead.

\DndMonsterAction{Magic Resistance}
Fierna has advantage on saving throws against spells and other magical effects.

\DndMonsterAction{Princess of Hellfire}
Fire conjured by Fierna becomes Hellfire. It ignores resistances and immunities to fire damage, deals double damage to creatures vulnerable to fire or necrotic damage, can't be extinguished through any means, and is capable of melting stone or igniting inflammable objects.

\DndMonsterSection{Actions}
\DndMonsterAction{Multiattack}
Fierna makes three attacks using Flame Blade, Mote of Flame, or a combination of the two.

\DndMonsterAction{Flame Blade}
\textsl{Melee Spell Attack:} 20 to hit, reach 5 ft., one target. \textsl{Hit:} 23 (4d6 + 9) fire damage.

\DndMonsterAction{Mote of Flame}
Fierna targets a creature she can see within 90 feet of her. The target must make a DC 25 Dexterity saving throw, taking 22 (3d8 + 9) fire damage on a failed save, or half as much damage on a successful one.

\DndMonsterSection{Reactions}
\DndMonsterAction{Wreath of Flames}
When hit by a melee attack and Fierna can see her attacker, she covers herself in Hellfire, and the attacker takes 13 (3d8) fire damage.

\DndMonsterSection{Legendary Actions}
Fierna can take 3 legendary actions, choosing from the options below. Only one legendary action can be used at a time and only at the end of another creature's turn. Fierna regains spent legendary actions at the start of its turn.

\begin{DndMonsterLegendaryActions}
\DndMonsterLegendaryAction{Blade}{Fierna makes a Flame Blade attack.
}
\DndMonsterLegendaryAction{Conjure Hellfire (Costs 2 Actions)}{Fierna chooses a point she can see within 150 feet of her. All creatures in a 20-foot-radius sphere centered on that point must make a DC 25 Dexterity saving throw, taking 28 (8d6) fire damage on a failed save, or half as much damage on a successful one.
}
\DndMonsterLegendaryAction{Call Underling (Costs 3 Actions)}{Fierna summons an allied \textbf{spined devil} in an unoccupied space that she can see.
}
\end{DndMonsterLegendaryActions}
\DndMonsterSection{Regional Effects}
The region containing Fierna's lair is corrupted by her presence, which creates the following effect:


\begin{description}
\item[Smoke-cloud.]
A vast fog rises 3 feet above ground within 1 mile of the lair, completely obscuring the ground. Any creature having the prone condition at the start of their turn must succeed on a DC 22 Constitution saving throw or take 7 (2d12) necrotic damage from choking on the noxious fumes.

\end{description}

If Fierna dies, the effects fade over the course of 1d10 days.

\end{multicols}
\end{DndMonster}



\begin{DndMonster}[float=t]{Fire Elemental}
\DndMonsterType{Large elemental, neutral}
\DndMonsterBasics[
armor-class = {13},
hit-points = {\DndDice{12d10 + 36}},
speed = {50 ft.}
]
\DndMonsterAbilityScores[
 str = 10,
 dex = 17,
 con = 16,
 int = 6,
 wis = 10,
 cha = 7
]
\DndMonsterDetails[
damage-resistances = {bludgeoning, piercing, slashing from nonmagical attacks},
damage-immunities = {fire, poison},
condition-immunities = {exhaustion, grappled, paralyzed, petrified, poisoned, prone, restrained, unconscious},
senses = {darkvision 60 ft., passive perception 10},
languages = {Ignan},
challenge = 5,
]
\DndMonsterAction{Fire Form}
The elemental can move through a space as narrow as 1 inch wide without squeezing. A creature that touches the elemental or hits it with a melee attack while within 5 feet of it takes 5 (1d10) fire damage. In addition, the elemental can enter a hostile creature's space and stop there. The first time it enters a creature's space on a turn, that creature takes 5 (1d10) fire damage and catches fire; until someone takes an action to douse the fire, the creature takes 5 (1d10) fire damage at the start of each of its turns.

\DndMonsterAction{Illumination}
The elemental sheds bright light in a 30-foot radius and dim light in an additional 30 feet.

\DndMonsterAction{Water Susceptibility}
For every 5 feet the elemental moves in water, or for every gallon of water splashed on it, it takes 1 cold damage.

\DndMonsterSection{Actions}
\DndMonsterAction{Multiattack}
The elemental makes two touch attacks.

\DndMonsterAction{Touch}
\textsl{Melee Weapon Attack:} 6 to hit, reach 5 ft., one target. \textsl{Hit:} 10 (2d6 + 3) fire damage. If the target is a creature or a flammable object, it ignites. Until a creature takes an action to douse the fire, the target takes 5 (1d10) fire damage at the start of each of its turns.

\end{DndMonster}



\begin{DndMonster}[float=t]{Fire Giant}
\DndMonsterType{Huge giant, lawful evil}
\DndMonsterBasics[
armor-class = {18 (\textit{plate armor})},
hit-points = {\DndDice{13d12 + 78}},
speed = {30 ft.}
]
\DndMonsterAbilityScores[
 str = 25,
 dex = 9,
 con = 23,
 int = 10,
 wis = 14,
 cha = 13
]
\DndMonsterDetails[
saving-throws = {Dex +3, Con +10, Cha +5},
skills = {Athletics +11, Perception +6},
damage-immunities = {fire},
languages = {Giant},
challenge = 9,
]
\DndMonsterSection{Actions}
\DndMonsterAction{Multiattack}
The giant makes two greatsword attacks.

\DndMonsterAction{Greatsword}
\textsl{Melee Weapon Attack:} 11 to hit, reach 10 ft., one target. \textsl{Hit:} 28 (6d6 + 7) slashing damage.

\DndMonsterAction{Rock}
\textsl{Ranged Weapon Attack:} 11 to hit, range 60/240 ft., one target. \textsl{Hit:} 29 (4d10 + 7) bludgeoning damage.

\end{DndMonster}



\begin{DndMonster}[float*=b,width=\textwidth + 8pt]{Fire Kraken}
\begin{multicols}{2}
\DndMonsterType{Gargantuan elemental, chaotic neutral}
\DndMonsterBasics[
armor-class = {16 (natural armor)},
hit-points = {\DndDice{23d20 + 161}},
speed = {30 ft., swim 60 ft.}
]
\DndMonsterAbilityScores[
 str = 27,
 dex = 11,
 con = 24,
 int = 10,
 wis = 20,
 cha = 15
]
\DndMonsterDetails[
saving-throws = {Dex +7, Con +14},
skills = {Athletics +15, Nature +7, Survival +12},
damage-immunities = {fire; bludgeoning, piercing, slashing from nonmagical attacks},
condition-immunities = {frightened, paralyzed},
senses = {blindsight 120 ft., passive perception 15},
languages = {understands Abyssal, Celestial, Infernal, and Primordial but can't speak, telepathy 120 ft.},
challenge = 21,
]
\DndMonsterAction{Freedom of Movement}
The kraken ignores difficult terrain, and magical effects can't reduce its speed or cause it to be restrained. It can spend 5 feet of movement to escape from nonmagical restraints or grapples.

\DndMonsterAction{Legendary Resistance (3/Day)}
If the kraken fails a saving throw, it can choose to succeed instead.

\DndMonsterAction{Siege Monster}
The kraken deals double damage to objects and structures.

\DndMonsterSection{Actions}
\DndMonsterAction{Multiattack}
The kraken makes two attacks using its Tentacle, Fling, or a combination of the two.

\DndMonsterAction{Bite}
\textsl{Melee Weapon Attack:} 15 to hit, reach 10 ft., one target. \textsl{Hit:} 18 (3d6 + 8) piercing damage and 7 (2d6) fire damage. If the target is a Large or smaller creature grappled by the kraken, the creature is swallowed and the grapple ends. While swallowed, a creature has the blinded and restrained conditions, has total cover against attacks and other effects outside the kraken, and takes 42 (12d6) fire damage at the start of the kraken's turns. If the kraken takes 50 damage or more on a single turn from a creature inside it, it must succeed on a DC 23 Constitution saving throw or regurgitate all swallowed creatures, which fall into a space within 10 feet of the kraken and now have the prone condition. If the kraken dies, a swallowed creature is no longer restrained by it and can escape from the corpse using 15 feet of movement, exiting with the prone condition.

\DndMonsterAction{Tentacle}
\textsl{Melee Weapon Attack:} 15 to hit, reach 20 ft., one target. \textsl{Hit:} 13 (2d4 + 8) fire damage. If the target is a Huge or smaller creature, it has the grappled condition (escape DC 16). While grappled, the target also has the restrained condition. Any target that starts its turn grappled by the kraken takes 5 (2d4) fire damage. The kraken has five tentacles, each of which can grapple one target.

\DndMonsterAction{Fling}
One Large or smaller object or creature grappled by the kraken is thrown up to 60 feet in a random direction, and now has the prone condition. If a thrown target strikes a solid surface, the target takes 3 (1d6) bludgeoning damage for every 10 feet it was thrown. If the target is thrown at a creature, that creature must succeed on a DC 16 Dexterity saving throw or take the same damage and be knocked down with the prone condition.

\DndMonsterSection{Legendary Actions}
Fire Kraken can take 3 legendary actions, choosing from the options below. Only one legendary action can be used at a time and only at the end of another creature's turn. Fire Kraken regains spent legendary actions at the start of its turn.

\begin{DndMonsterLegendaryActions}
\DndMonsterLegendaryAction{Appendage}{The kraken makes a Fling attack.
}
\DndMonsterLegendaryAction{Fire Breath (Costs 3 Actions)}{The kraken exhales fire in a 60-foot cone. Each creature in that area must make a DC 20 Dexterity saving throw, taking 36 (8d8) fire damage on a failed save, or half as much damage on a successful one.
}
\end{DndMonsterLegendaryActions}
\end{multicols}
\end{DndMonster}



\begin{DndMonster}[float=t]{Fire Snake}
\DndMonsterType{Medium elemental, neutral evil}
\DndMonsterBasics[
armor-class = {14 (natural armor)},
hit-points = {\DndDice{5d8}},
speed = {30 ft.}
]
\DndMonsterAbilityScores[
 str = 12,
 dex = 14,
 con = 11,
 int = 7,
 wis = 10,
 cha = 8
]
\DndMonsterDetails[
damage-vulnerabilities = {cold},
damage-resistances = {bludgeoning, piercing, slashing from nonmagical attacks},
damage-immunities = {fire},
senses = {darkvision 60 ft., passive perception 10},
languages = {understands Ignan but can't speak},
challenge = 1,
]
\DndMonsterAction{Heated Body}
A creature that touches the snake or hits it with a melee attack while within 5 feet of it takes 3 (1d6) fire damage.

\DndMonsterSection{Actions}
\DndMonsterAction{Multiattack}
The snake makes two attacks: one with its bite and one with its tail.

\DndMonsterAction{Bite}
\textsl{Melee Weapon Attack:} 3 to hit, reach 5 ft., one target. \textsl{Hit:} 3 (1d4 + 1) piercing damage plus 3 (1d6) fire damage.

\DndMonsterAction{Tail}
\textsl{Melee Weapon Attack:} 3 to hit, reach 5 ft., one target. \textsl{Hit:} 3 (1d4 + 1) bludgeoning damage plus 3 (1d6) fire damage.

\end{DndMonster}



\begin{DndMonster}[float=t]{Flameskull}
\DndMonsterType{Tiny undead, neutral evil}
\DndMonsterBasics[
armor-class = {13},
hit-points = {\DndDice{9d4 + 18}},
speed = {0 ft., fly 40 ft. \textit{(hover)}}
]
\DndMonsterAbilityScores[
 str = 1,
 dex = 17,
 con = 14,
 int = 16,
 wis = 10,
 cha = 11
]
\DndMonsterDetails[
skills = {Arcana +5, Perception +2},
damage-resistances = {lightning, necrotic, piercing},
damage-immunities = {cold, fire, poison},
condition-immunities = {charmed, frightened, paralyzed, poisoned, prone},
senses = {darkvision 60 ft., passive perception 12},
languages = {Common},
challenge = 4,
]
\DndMonsterAction{Illumination}
The flameskull sheds either dim light in a 15-foot radius, or bright light in a 15-foot radius and dim light for an additional 15 feet. It can switch between the options as an action.

\DndMonsterAction{Magic Resistance}
The flameskull has advantage on saving throws against spells and other magical effects.

\DndMonsterAction{Rejuvenation}
If the flameskull is destroyed, it regains all its hit points in 1 hour unless holy water is sprinkled on its remains or a \textit{dispel magic} or \textit{remove curse} spell is cast on them.

\DndMonsterAction{Spellcasting}
The flameskull is a 5th-level spellcaster. Its spellcasting ability is Intelligence (spell save DC 13, 5 to hit with spell attacks). It requires no somatic or material components to cast its spells. The flameskull has the following wizard spells prepared:
\begin{DndMonsterSpells}
  \DndMonsterSpellLevel{\textit{mage hand}}
   \DndMonsterSpellLevel[1][3]{\textit{magic missile}, \textit{shield}}
   \DndMonsterSpellLevel[2][2]{\textit{blur}, \textit{flaming sphere}}
   \DndMonsterSpellLevel[3][1]{\textit{fireball}}
\end{DndMonsterSpells}
\DndMonsterSection{Actions}
\DndMonsterAction{Multiattack}
The flameskull uses Fire Ray twice.

\DndMonsterAction{Fire Ray}
\textsl{Ranged Spell Attack:} 5 to hit, range 30 ft., one target. \textsl{Hit:} 10 (3d6) fire damage.

\end{DndMonster}



\begin{DndMonster}[float=t]{Flesh Golem}
\DndMonsterType{Medium construct, neutral}
\DndMonsterBasics[
armor-class = {9},
hit-points = {\DndDice{11d8 + 44}},
speed = {30 ft.}
]
\DndMonsterAbilityScores[
 str = 19,
 dex = 9,
 con = 18,
 int = 6,
 wis = 10,
 cha = 5
]
\DndMonsterDetails[
damage-immunities = {lightning; poison; bludgeoning, piercing, slashing from nonmagical attacks that aren't adamantine},
condition-immunities = {charmed, exhaustion, frightened, paralyzed, petrified, poisoned},
senses = {darkvision 60 ft., passive perception 10},
languages = {understands the languages of its creator but can't speak},
challenge = 5,
]
\DndMonsterAction{Berserk}
Whenever the golem starts its turn with 40 hit points or fewer, roll a d6. On a 6, the golem goes berserk. On each of its turns while berserk, the golem attacks the nearest creature it can see. If no creature is near enough to move to and attack, the golem attacks an object, with preference for an object smaller than itself. Once the golem goes berserk, it continues to do so until it is destroyed or regains all its hit points.

The golem's creator, if within 60 feet of the berserk golem, can try to calm it by speaking firmly and persuasively. The golem must be able to hear its creator, who must take an action to make a DC 15 Charisma (Persuasion) check. If the check succeeds, the golem ceases being berserk. If it takes damage while still at 40 hit points or fewer, the golem might go berserk again.

\DndMonsterAction{Aversion of Fire}
If the golem takes fire damage, it has disadvantage on attack rolls and ability checks until the end of its next turn.

\DndMonsterAction{Immutable Form}
The golem is immune to any spell or effect that would alter its form.

\DndMonsterAction{Lightning Absorption}
Whenever the golem is subjected to lightning damage, it takes no damage and instead regains a number of hit points equal to the lightning damage dealt.

\DndMonsterAction{Magic Resistance}
The golem has advantage on saving throws against spells and other magical effects.

\DndMonsterAction{Magic Weapons}
The golem's weapon attacks are magical.

\DndMonsterSection{Actions}
\DndMonsterAction{Multiattack}
The golem makes two slam attacks.

\DndMonsterAction{Slam}
\textsl{Melee Weapon Attack:} 7 to hit, reach 5 ft., one target. \textsl{Hit:} 13 (2d8 + 4) bludgeoning damage.

\end{DndMonster}


\clearpage

\begin{DndMonster}[float=t]{Fomorian}
\DndMonsterType{Huge giant, chaotic evil}
\DndMonsterBasics[
armor-class = {14 (natural armor)},
hit-points = {\DndDice{13d12 + 65}},
speed = {30 ft.}
]
\DndMonsterAbilityScores[
 str = 23,
 dex = 10,
 con = 20,
 int = 9,
 wis = 14,
 cha = 6
]
\DndMonsterDetails[
skills = {Perception +8, Stealth +3},
senses = {darkvision 120 ft., passive perception 18},
languages = {Giant, Undercommon},
challenge = 8,
]
\DndMonsterSection{Actions}
\DndMonsterAction{Multiattack}
The fomorian attacks twice with its greatclub or makes one greatclub attack and uses Evil Eye once.

\DndMonsterAction{Greatclub}
\textsl{Melee Weapon Attack:} 9 to hit, reach 15 ft., one target. \textsl{Hit:} 19 (3d8 + 6) bludgeoning damage.

\DndMonsterAction{Evil Eye}
The fomorian magically forces a creature it can see within 60 feet of it to make a DC 14 Charisma saving throw. The creature takes 27 (6d8) psychic damage on a failed save, or half as much damage on a successful one.

\DndMonsterAction{Curse of the Evil Eye (Recharges after a Short or Long Rest)}
With a stare, the fomorian uses Evil Eye, but on a failed save, the creature is also cursed with magical deformities. While deformed, the creature has its speed halved and has disadvantage on ability checks, saving throws, and attacks based on Strength or Dexterity.

The transformed creature can repeat the saving throw whenever it finishes a long rest, ending the effect on a success.

\end{DndMonster}



\begin{DndMonster}[float=t]{Froghemoth}
\DndMonsterType{Huge monstrosity, unaligned}
\DndMonsterBasics[
armor-class = {14 (natural armor)},
hit-points = {\DndDice{14d12 + 70}},
speed = {30 ft., swim 30 ft.}
]
\DndMonsterAbilityScores[
 str = 23,
 dex = 13,
 con = 20,
 int = 2,
 wis = 12,
 cha = 5
]
\DndMonsterDetails[
saving-throws = {Con +9, Wis +5},
skills = {Perception +9, Stealth +5},
damage-resistances = {fire, lightning},
senses = {darkvision 60 ft., passive perception 19},
challenge = 10,
]
\DndMonsterAction{Amphibious}
The froghemoth can breathe air and water.

\DndMonsterAction{Shock Susceptibility}
If the froghemoth takes lightning damage, it suffers two effects until the end of its next turn: its speed is halved, and it has disadvantage on Dexterity saving throws.

\DndMonsterSection{Actions}
\DndMonsterAction{Multiattack}
The froghemoth makes one Bite attack and two Tentacle attacks, and it can use Tongue.

\DndMonsterAction{Bite}
\textsl{Melee Weapon Attack:} 10 to hit, reach 5 ft., one target. \textsl{Hit:} 22 (3d10 + 6) piercing damage, and the target is swallowed if it is a Medium or smaller creature. A swallowed creature is blinded and restrained, has total cover against attacks and other effects outside the froghemoth, and takes 10 (3d6) acid damage at the start of each of the froghemoth's turns.

The froghemoth's gullet can hold up to two creatures at a time. If the froghemoth takes 20 damage or more on a single turn from a creature inside it, the froghemoth must succeed on a DC 20 Constitution saving throw at the end of that turn or regurgitate all swallowed creatures, each of which falls prone in a space within 10 feet of the froghemoth. If the froghemoth dies, any swallowed creature is no longer restrained by it and can escape from the corpse using 10 feet of movement, exiting prone.

\DndMonsterAction{Tentacle}
\textsl{Melee Weapon Attack:} 10 to hit, reach 20 ft., one target. \textsl{Hit:} 19 (3d8 + 6) bludgeoning damage, and the target is grappled (escape DC 16) if it is a Huge or smaller creature. Until the grapple ends, the froghemoth can't use this tentacle on another target. The froghemoth has four tentacles.

\DndMonsterAction{Tongue}
The froghemoth targets one Medium or smaller creature that it can see within 20 feet of it. The target must make a DC 18 Strength saving throw. On a failed save, the target is pulled into an unoccupied space within 5 feet of the froghemoth.

\end{DndMonster}



\begin{DndMonster}[float=t]{Frost Giant}
\DndMonsterType{Huge giant, neutral evil}
\DndMonsterBasics[
armor-class = {15 (patchwork armor)},
hit-points = {\DndDice{12d12 + 60}},
speed = {40 ft.}
]
\DndMonsterAbilityScores[
 str = 23,
 dex = 9,
 con = 21,
 int = 9,
 wis = 10,
 cha = 12
]
\DndMonsterDetails[
saving-throws = {Con +8, Wis +3, Cha +4},
skills = {Athletics +9, Perception +3},
damage-immunities = {cold},
languages = {Giant},
challenge = 8,
]
\DndMonsterSection{Actions}
\DndMonsterAction{Multiattack}
The giant makes two greataxe attacks.

\DndMonsterAction{Greataxe}
\textsl{Melee Weapon Attack:} 9 to hit, reach 10 ft., one target. \textsl{Hit:} 25 (3d12 + 6) slashing damage.

\DndMonsterAction{Rock}
\textsl{Ranged Weapon Attack:} 9 to hit, range 60/240 ft., one target. \textsl{Hit:} 28 (4d10 + 6) bludgeoning damage.

\end{DndMonster}



\begin{DndMonster}[float=t]{Frozen Golem}
\DndMonsterType{Large construct, unaligned}
\DndMonsterBasics[
armor-class = {20 (natural armor)},
hit-points = {\DndDice{20d10 + 100}},
speed = {30 ft.}
]
\DndMonsterAbilityScores[
 str = 24,
 dex = 9,
 con = 20,
 int = 3,
 wis = 11,
 cha = 1
]
\DndMonsterDetails[
damage-immunities = {fire; poison; psychic; bludgeoning, piercing, slashing from nonmagical attacks that aren't adamantine},
condition-immunities = {charmed, exhaustion, frightened, paralyzed, petrified, poisoned},
senses = {darkvision 120 ft., passive perception 10},
languages = {understands the languages of its creator but can't speak},
challenge = 16,
]
\DndMonsterAction{Copy Warning}
This content has been copied from another creature, but the data replacement hasn't been run. Check details.

\DndMonsterAction{Fire Absorption}
Whenever the golem is subjected to fire damage, it takes no damage and instead regains a number of hit points equal to the fire damage dealt.

\DndMonsterAction{Immutable Form}
The golem is immune to any spell or effect that would alter its form.

\DndMonsterAction{Magic Resistance}
The golem has advantage on saving throws against spells and other magical effects.

\DndMonsterAction{Magic Weapons}
The golem's weapon attacks are magical.

\DndMonsterSection{Actions}
\DndMonsterAction{Multiattack}
The golem makes two melee attacks.

\DndMonsterAction{Slam}
\textsl{Melee Weapon Attack:} 13 to hit, reach 5 ft., one target. \textsl{Hit:} 20 (3d8 + 7) bludgeoning damage.

\DndMonsterAction{Sword}
\textsl{Melee Weapon Attack:} 13 to hit, reach 10 ft., one target. \textsl{Hit:} 23 (3d10 + 7) slashing damage.

\DndMonsterAction{Poison Breath (Recharge 5\textendash 6)}
The golem exhales poisonous gas in a 15-foot cone. Each creature in that area must make a DC 19 Constitution saving throw, taking 45 (10d8) poison damage on a failed save, or half as much damage on a successful one.

\end{DndMonster}



\begin{DndMonster}[float=t]{Giant Scorpion}
\DndMonsterType{Large beast, unaligned}
\DndMonsterBasics[
armor-class = {15 (natural armor)},
hit-points = {\DndDice{7d10 + 14}},
speed = {40 ft.}
]
\DndMonsterAbilityScores[
 str = 15,
 dex = 13,
 con = 15,
 int = 1,
 wis = 9,
 cha = 3
]
\DndMonsterDetails[
senses = {blindsight 60 ft., passive perception 9},
challenge = 3,
]
\DndMonsterSection{Actions}
\DndMonsterAction{Multiattack}
The scorpion makes three attacks: two with its claws and one with its sting.

\DndMonsterAction{Claw}
\textsl{Melee Weapon Attack:} 4 to hit, reach 5 ft., one target. \textsl{Hit:} 6 (1d8 + 2) bludgeoning damage, and the target is grappled (escape DC 12). The scorpion has two claws, each of which can grapple only one target.

\DndMonsterAction{Sting}
\textsl{Melee Weapon Attack:} 4 to hit, reach 5 ft., one creature. \textsl{Hit:} 7 (1d10 + 2) piercing damage, and the target must make a DC 12 Constitution saving throw, taking 22 (4d10) poison damage on a failed save, or half as much damage on a successful one.

\end{DndMonster}



\begin{DndMonster}[float=t]{Giant Spider}
\DndMonsterType{Large beast, unaligned}
\DndMonsterBasics[
armor-class = {14 (natural armor)},
hit-points = {\DndDice{4d10 + 4}},
speed = {30 ft., climb 30 ft.}
]
\DndMonsterAbilityScores[
 str = 14,
 dex = 16,
 con = 12,
 int = 2,
 wis = 11,
 cha = 4
]
\DndMonsterDetails[
skills = {Stealth +7},
senses = {blindsight 10 ft., darkvision 60 ft., passive perception 10},
challenge = 1,
]
\DndMonsterAction{Spider Climb}
The spider can climb difficult surfaces, including upside down on ceilings, without needing to make an ability check.

\DndMonsterAction{Web Sense}
While in contact with a web, the spider knows the exact location of any other creature in contact with the same web.

\DndMonsterAction{Web Walker}
The spider ignores movement restrictions caused by webbing.

\DndMonsterSection{Actions}
\DndMonsterAction{Bite}
\textsl{Melee Weapon Attack:} 5 to hit, reach 5 ft., one creature. \textsl{Hit:} 7 (1d8 + 3) piercing damage, and the target must make a DC 11 Constitution saving throw, taking 9 (2d8) poison damage on a failed save, or half as much damage on a successful one. If the poison damage reduces the target to 0 hit points, the target is stable but poisoned for 1 hour, even after regaining hit points, and is paralyzed while poisoned in this way.

\DndMonsterAction{Web (Recharge 5\textendash 6)}
\textsl{Ranged Weapon Attack:} 5 to hit, range 30/60 ft., one creature. \textsl{Hit:} The target is restrained by webbing. As an action, the restrained target can make a DC 12 Strength check, bursting the webbing on a success. The webbing can also be attacked and destroyed (AC 10; hp 5; vulnerability to fire damage; immunity to bludgeoning, poison, and psychic damage).

\end{DndMonster}



\begin{DndMonster}[float=t]{Giant Vulture}
\DndMonsterType{Large beast, neutral evil}
\DndMonsterBasics[
armor-class = {10},
hit-points = {\DndDice{3d10 + 6}},
speed = {10 ft., fly 60 ft.}
]
\DndMonsterAbilityScores[
 str = 15,
 dex = 10,
 con = 15,
 int = 6,
 wis = 12,
 cha = 7
]
\DndMonsterDetails[
skills = {Perception +3},
languages = {understands Common but can't speak},
challenge = 1,
]
\DndMonsterAction{Keen Sight and Smell}
The vulture has advantage on Wisdom (Perception) checks that rely on sight or smell.

\DndMonsterAction{Pack Tactics}
The vulture has advantage on an attack roll against a creature if at least one of the vulture's allies is within 5 feet of the creature and the ally isn't incapacitated.

\DndMonsterSection{Actions}
\DndMonsterAction{Multiattack}
The vulture makes two attacks: one with its beak and one with its talons.

\DndMonsterAction{Beak}
\textsl{Melee Weapon Attack:} 4 to hit, reach 5 ft., one target. \textsl{Hit:} 7 (2d4 + 2) piercing damage.

\DndMonsterAction{Talons}
\textsl{Melee Weapon Attack:} 4 to hit, reach 5 ft., one target. \textsl{Hit:} 9 (2d6 + 2) slashing damage.

\end{DndMonster}



\begin{DndMonster}[float=t]{Gibbering Mouther}
\DndMonsterType{Medium aberration, neutral}
\DndMonsterBasics[
armor-class = {9},
hit-points = {\DndDice{9d8 + 27}},
speed = {10 ft., swim 10 ft.}
]
\DndMonsterAbilityScores[
 str = 10,
 dex = 8,
 con = 16,
 int = 3,
 wis = 10,
 cha = 6
]
\DndMonsterDetails[
condition-immunities = {prone},
senses = {darkvision 60 ft., passive perception 10},
challenge = 2,
]
\DndMonsterAction{Aberrant Ground}
The ground in a 10-foot radius around the mouther is doughlike difficult terrain. Each creature that starts its turn in that area must succeed on a DC 10 Strength saving throw or have its speed reduced to 0 until the start of its next turn.

\DndMonsterAction{Gibbering}
The mouther babbles incoherently while it can see any creature and isn't incapacitated. Each creature that starts its turn within 20 feet of the mouther and can hear the gibbering must succeed on a DC 10 Wisdom saving throw. On a failure, the creature can't take reactions until the start of its next turn and rolls a d8 to determine what it does during its turn. On a 1 to 4, the creature does nothing. On a 5 or 6, the creature takes no action or bonus action and uses all its movement to move in a randomly determined direction. On a 7 or 8, the creature makes a melee attack against a randomly determined creature within its reach or does nothing if it can't make such an attack.

\DndMonsterSection{Actions}
\DndMonsterAction{Multiattack}
The gibbering mouther makes one bite attack and, if it can, uses its Blinding Spittle.

\DndMonsterAction{Bites}
\textsl{Melee Weapon Attack:} 2 to hit, reach 5 ft., one creature. \textsl{Hit:} 17 (5d6) piercing damage. If the target is Medium or smaller, it must succeed on a DC 10 Strength saving throw or be knocked prone. If the target is killed by this damage, it is absorbed into the mouther.

\DndMonsterAction{Blinding Spittle (Recharge 5\textendash 6)}
The mouther spits a chemical glob at a point it can see within 15 feet of it. The glob explodes in a blinding flash of light on impact. Each creature within 5 feet of the flash must succeed on a DC 13 Dexterity saving throw or be blinded until the end of the mouther's next turn.

\end{DndMonster}



\begin{DndMonster}[float=t]{Githyanki Knight}
\DndMonsterType{Medium humanoid (gith), lawful evil}
\DndMonsterBasics[
armor-class = {18 (\textit{plate armor})},
hit-points = {\DndDice{14d8 + 28}},
speed = {30 ft.}
]
\DndMonsterAbilityScores[
 str = 16,
 dex = 14,
 con = 15,
 int = 14,
 wis = 14,
 cha = 15
]
\DndMonsterDetails[
saving-throws = {Con +5, Int +5, Wis +5},
languages = {Gith},
challenge = 8,
]
\DndMonsterAction{Innate Spellcasting (Psionics)}
The githyanki's innate spellcasting ability is Intelligence (spell save DC 13, 5 to hit with spell attacks). It can innately cast the following spells, requiring no components:
\begin{DndMonsterSpells}
  \DndInnateSpellLevel{\textit{mage hand} (the hand is invisible)}
  \DndInnateSpellLevel[3e]{\textit{jump}, \textit{misty step}, \textit{nondetection} (self only), \textit{tongues}}
  \DndInnateSpellLevel[1e]{\textit{plane shift}, \textit{telekinesis}}
\end{DndMonsterSpells}
\DndMonsterSection{Actions}
\DndMonsterAction{Multiattack}
The githyanki makes two silver greatsword attacks.

\DndMonsterAction{Silver Greatsword}
\textsl{Melee Weapon Attack:} 9 to hit, reach 5 ft., one target. \textsl{Hit:} 13 (2d6 + 6) slashing damage plus 10 (3d6) psychic damage. This is a magic weapon attack. On a critical hit against a target in an astral body (as with the \textit{astral projection} spell), the githyanki can cut the silvery cord that tethers the target to its material body, instead of dealing damage.

\end{DndMonster}



\begin{DndMonster}[float=t]{Glabrezu}
\DndMonsterType{Large fiend (demon), chaotic evil}
\DndMonsterBasics[
armor-class = {17 (natural armor)},
hit-points = {\DndDice{15d10 + 75}},
speed = {40 ft.}
]
\DndMonsterAbilityScores[
 str = 20,
 dex = 15,
 con = 21,
 int = 19,
 wis = 17,
 cha = 16
]
\DndMonsterDetails[
saving-throws = {Str +9, Con +9, Wis +7, Cha +7},
damage-resistances = {cold; fire; lightning; bludgeoning, piercing, slashing from nonmagical attacks},
damage-immunities = {poison},
condition-immunities = {poisoned},
senses = {truesight 120 ft., passive perception 13},
languages = {Abyssal, telepathy 120 ft.},
challenge = 9,
]
\DndMonsterAction{Magic Resistance}
The glabrezu has advantage on saving throws against spells and other magical effects.

\DndMonsterAction{Innate Spellcasting}
The glabrezu's spellcasting ability is Intelligence (spell save DC 16). The glabrezu can innately cast the following spells, requiring no material components:
\begin{DndMonsterSpells}
  \DndInnateSpellLevel{\textit{darkness}, \textit{detect magic}, \textit{dispel magic}}
  \DndInnateSpellLevel[1e]{\textit{confusion}, \textit{fly}, \textit{power word stun}}
\end{DndMonsterSpells}
\DndMonsterSection{Actions}
\DndMonsterAction{Multiattack}
The glabrezu makes four attacks: two with its pincers and two with its fists. Alternatively, it makes two attacks with its pincers and casts one spell.

\DndMonsterAction{Pincer}
\textsl{Melee Weapon Attack:} 9 to hit, reach 10 ft., one target. \textsl{Hit:} 16 (2d10 + 5) bludgeoning damage. If the target is a Medium or smaller creature, it is grappled (escape DC 15). The glabrezu has two pincers, each of which can grapple only one target.

\DndMonsterAction{Fist}
\textsl{Melee Weapon Attack:} 9 to hit, reach 5 ft., one target. \textsl{Hit:} 7 (2d4 + 2) bludgeoning damage.

\end{DndMonster}



\begin{DndMonster}[float=t]{Gladiator}
\DndMonsterType{Medium humanoid (any race), any alignment}
\DndMonsterBasics[
armor-class = {16 (studded leather armor)},
hit-points = {\DndDice{15d8 + 45}},
speed = {30 ft.}
]
\DndMonsterAbilityScores[
 str = 18,
 dex = 15,
 con = 16,
 int = 10,
 wis = 12,
 cha = 15
]
\DndMonsterDetails[
saving-throws = {Str +7, Dex +5, Con +6},
skills = {Athletics +10, Intimidation +5},
languages = {any one language (usually Common)},
challenge = 5,
]
\DndMonsterAction{Brave}
The gladiator has advantage on saving throws against being frightened.

\DndMonsterAction{Brute}
A melee weapon deals one extra die of its damage when the gladiator hits with it (included in the attack).

\DndMonsterSection{Actions}
\DndMonsterAction{Multiattack}
The gladiator makes three melee attacks or two ranged attacks.

\DndMonsterAction{Spear}
\textsl{Melee or Ranged Weapon Attack:} 7 to hit, reach 5 ft. and range 20/60 ft., one target. \textsl{Hit:} 11 (2d6 + 4) piercing damage, or 13 (2d8 + 4) piercing damage if used with two hands to make a melee attack.

\DndMonsterAction{Shield Bash}
\textsl{Melee Weapon Attack:} 7 to hit, reach 5 ft., one creature. \textsl{Hit:} 9 (2d4 + 4) bludgeoning damage. If the target is a Medium or smaller creature, it must succeed on a DC 15 Strength saving throw or be knocked prone.

\DndMonsterSection{Reactions}
\DndMonsterAction{Parry}
The gladiator adds 3 to its AC against one melee attack that would hit it. To do so, the gladiator must see the attacker and be wielding a melee weapon.

\end{DndMonster}



\begin{DndMonster}[float*=b,width=\textwidth + 8pt]{Glasya}
\begin{multicols}{2}
\DndMonsterType{Medium fiend (devil), lawful evil}
\DndMonsterBasics[
armor-class = {21 (natural armor)},
hit-points = {\DndDice{40d8 + 200}},
speed = {40 ft., fly 80 ft.}
]
\DndMonsterAbilityScores[
 str = 22,
 dex = 28,
 con = 20,
 int = 21,
 wis = 25,
 cha = 28
]
\DndMonsterDetails[
saving-throws = {Dex +17, Int +13, Wis +15, Cha +17},
skills = {Deception +25, Intimidation +17, Perception +15, Persuasion +17, Stealth +25, Survival +15},
damage-resistances = {cold; necrotic; radiant; bludgeoning, piercing, slashing from nonmagical attacks that aren't silvered},
damage-immunities = {fire, poison},
condition-immunities = {charmed, poisoned},
senses = {blindsight 60 ft., darkvision 120 ft., passive perception 25},
languages = {Celestial, Common, Draconic, Infernal, telepathy 120 ft.},
challenge = 25,
]
\DndMonsterAction{Charm Aura}
When a creature moves to be within 120 feet of Glasya, they must succeed on a DC 23 Wisdom saving throw or have the charmed condition, treating Glasya as an ally, until they leave the aura. A creature that succeeds the saving throw is immune to this effect for 1 hour. Whenever a charmed creature takes damage, they may repeat the saving throw, ending the effect on a success.

\DndMonsterAction{Dominating Presence}
Glasya deals an additional 13 (2d12) psychic damage when attacking a charmed creature, and her attacks don't break charms.

\DndMonsterAction{Fiendish Regeneration}
Glasya regains 20 hit points at the start of her turn. If she takes radiant damage this trait doesn't function at the start of her next turn. Glasya dies only if she starts her turn with 0 hit points and is unable to regenerate. If Glasya is killed, her body slowly regenerates, returning to life 9 (2d8) weeks later.

\DndMonsterAction{Legendary Resistance (3/Day)}
If Glasya fails a saving throw, she can choose to succeed instead.

\DndMonsterAction{Magic Resistance}
Glasya has advantage on saving throws against spells and other magical effects.

\DndMonsterSection{Actions}
\DndMonsterAction{Multiattack}
Glasya makes four attacks using Scourge of Shadow, Necrotic Grasp, or a combination of the two. She can replace one of the attacks with Mesmerizing Gaze.

\DndMonsterAction{Scourge of Shadow}
\textsl{Melee Weapon Attack:} 20 to hit, reach 10 ft., one target. \textsl{Hit:} 19 (3d4 + 12) slashing damage, plus 5 (2d4) necrotic damage. This attack scores a critical on a roll of 18, 19, or 20.

\DndMonsterAction{Necrotic Grasp}
Glasya summons necrotic hands to grope and claw at a creature she can see within 90 feet of her. The target must make a DC 25 Constitution saving throw, taking 20 (2d10 + 9) necrotic damage on a failed save, or half as much damage on a successful one.

\DndMonsterAction{Mesmerizing Gaze}
Glasya focuses her gaze in a 60-foot cone in front of her. Each creature of her choice within the cone must make a DC 23 Charisma saving throw or have the charmed condition until the end of their next turn. While charmed, creatures automatically fail saving throws from abilities used by Glasya and move to be as close to her as possible.

\DndMonsterAction{Deathly Exclamation (1/Day)}
Glasya targets a creature she can see within 60 feet of her. The target must make a DC 25 Constitution saving throw, taking 61 (7d8 + 30) necrotic damage on a failed save, or half as much damage on a successful one. A Humanoid killed by this ability rises at the start of Glasya's next turn as a zombie permanently under her command.

\DndMonsterSection{Legendary Actions}
Glasya can take 3 legendary actions, choosing from the options below. Only one legendary action can be used at a time and only at the end of another creature's turn. Glasya regains spent legendary actions at the start of its turn.

\begin{DndMonsterLegendaryActions}
\DndMonsterLegendaryAction{Whip}{Glasya makes a Scourge of Shadow attack.
}
\DndMonsterLegendaryAction{Disorient (Costs 2 Actions)}{Glasya casts \textit{Confusion} (30-foot-radius, spell save DC 25), but only lasting until the end of Glasya's next turn.
}
\DndMonsterLegendaryAction{Call Underling (Costs 3 Actions)}{Glasya summons an allied \textbf{erinyes} in an unoccupied space that she can see.
}
\end{DndMonsterLegendaryActions}
\DndMonsterSection{Regional Effects}
The region containing Glasya's lair is corrupted by her presence, which creates the following effect:


\begin{description}
\item[Hag's Gaze.]
Every hour that a good creature travels within 1 mile of the lair, a blast of fire erupts from one of the eye sockets of the lair. If the target fails a DC 23 Dexterity saving throw, they take 42 (12d6) fire damage, half on a successful save. If the creature was flying and fails the saving throw, it falls. Only one such creature is targeted each hour.

\end{description}

If Glasya dies, the effects fade over the course of 1d10 days.

\end{multicols}
\end{DndMonster}



\begin{DndMonster}[float=t]{Gnoll}
\DndMonsterType{Medium humanoid (gnoll), chaotic evil}
\DndMonsterBasics[
armor-class = {15 (\textit{hide armor})},
hit-points = {\DndDice{5d8}},
speed = {30 ft.}
]
\DndMonsterAbilityScores[
 str = 14,
 dex = 12,
 con = 11,
 int = 6,
 wis = 10,
 cha = 7
]
\DndMonsterDetails[
senses = {darkvision 60 ft., passive perception 10},
languages = {Gnoll},
challenge = 1/2,
]
\DndMonsterAction{Rampage}
When the gnoll reduces a creature to 0 hit points with a melee attack on its turn, the gnoll can take a bonus action to move up to half its speed and make a bite attack.

\DndMonsterSection{Actions}
\DndMonsterAction{Bite}
\textsl{Melee Weapon Attack:} 4 to hit, reach 5 ft., one creature. \textsl{Hit:} 4 (1d4 + 2) piercing damage.

\DndMonsterAction{Spear}
\textsl{Melee or Ranged Weapon Attack:} 4 to hit, reach 5 ft. or range 20/60 ft., one target. \textsl{Hit:} 5 (1d6 + 2) piercing damage, or 6 (1d8 + 2) piercing damage if used with two hands to make a melee attack.

\DndMonsterAction{Longbow}
\textsl{Ranged Weapon Attack:} 3 to hit, range 150/600 ft., one target. \textsl{Hit:} 5 (1d8 + 1) piercing damage.

\end{DndMonster}



\begin{DndMonster}[float=t]{Goblin}
\DndMonsterType{Small humanoid (goblinoid), neutral evil}
\DndMonsterBasics[
armor-class = {15 (\textit{leather armor})},
hit-points = {\DndDice{2d6}},
speed = {30 ft.}
]
\DndMonsterAbilityScores[
 str = 8,
 dex = 14,
 con = 10,
 int = 10,
 wis = 8,
 cha = 8
]
\DndMonsterDetails[
skills = {Stealth +6},
senses = {darkvision 60 ft., passive perception 9},
languages = {Common, Goblin},
challenge = 1/4,
]
\DndMonsterAction{Nimble Escape}
The goblin can take the Disengage or Hide action as a bonus action on each of its turns.

\DndMonsterSection{Actions}
\DndMonsterAction{Scimitar}
\textsl{Melee Weapon Attack:} 4 to hit, reach 5 ft., one target. \textsl{Hit:} 5 (1d6 + 2) slashing damage.

\DndMonsterAction{Shortbow}
\textsl{Ranged Weapon Attack:} 4 to hit, range 80/320 ft., one target. \textsl{Hit:} 5 (1d6 + 2) piercing damage.

\end{DndMonster}



\begin{DndMonster}[float=t]{Goristro}
\DndMonsterType{Huge fiend (demon), chaotic evil}
\DndMonsterBasics[
armor-class = {19 (natural armor)},
hit-points = {\DndDice{23d12 + 161}},
speed = {40 ft.}
]
\DndMonsterAbilityScores[
 str = 25,
 dex = 11,
 con = 25,
 int = 6,
 wis = 13,
 cha = 14
]
\DndMonsterDetails[
saving-throws = {Str +13, Dex +6, Con +13, Wis +7},
skills = {Perception +7},
damage-resistances = {cold; fire; lightning; bludgeoning, piercing, slashing from nonmagical attacks},
damage-immunities = {poison},
condition-immunities = {poisoned},
senses = {darkvision 120 ft., passive perception 17},
languages = {Abyssal},
challenge = 17,
]
\DndMonsterAction{Charge}
If the goristro moves at least 15 feet straight toward a target and then hits it with a gore attack on the same turn, the target takes an extra 38 (7d10) piercing damage. If the target is a creature, it must succeed on a DC 21 Strength saving throw or be pushed up to 20 feet away and knocked prone.

\DndMonsterAction{Labyrinthine Recall}
The goristro can perfectly recall any path it has traveled.

\DndMonsterAction{Magic Resistance}
The goristro has advantage on saving throws against spells and other magical effects.

\DndMonsterAction{Siege Monster}
The goristro deals double damage to objects and structures.

\DndMonsterSection{Actions}
\DndMonsterAction{Multiattack}
The goristro makes three attacks: two with its fists and one with its hoof.

\DndMonsterAction{Fist}
\textsl{Melee Weapon Attack:} 13 to hit, reach 10 ft., one target. \textsl{Hit:} 20 (3d8 + 7) bludgeoning damage.

\DndMonsterAction{Hoof}
\textsl{Melee Weapon Attack:} 13 to hit, reach 5 ft., one target. \textsl{Hit:} 23 (3d10 + 7) bludgeoning damage. If the target is a creature, it must succeed on a DC 21 Strength saving throw or be knocked prone.

\DndMonsterAction{Gore}
\textsl{Melee Weapon Attack:} 13 to hit, reach 10 ft., one target. \textsl{Hit:} 45 (7d10 + 7) piercing damage.

\end{DndMonster}



\begin{DndMonster}[float=t]{Gray Render}
\DndMonsterType{Large monstrosity, chaotic neutral}
\DndMonsterBasics[
armor-class = {19 (natural armor)},
hit-points = {\DndDice{18d10 + 90}},
speed = {30 ft.}
]
\DndMonsterAbilityScores[
 str = 19,
 dex = 13,
 con = 20,
 int = 3,
 wis = 6,
 cha = 8
]
\DndMonsterDetails[
saving-throws = {Str +8, Con +9},
skills = {Perception +2},
senses = {darkvision 60 ft., passive perception 12},
challenge = 12,
]
\DndMonsterSection{Actions}
\DndMonsterAction{Multiattack}
The gray render makes one Bite attack and two Claw attacks.

\DndMonsterAction{Bite}
\textsl{Melee Weapon Attack:} 8 to hit, reach 5 ft., one target. \textsl{Hit:} 17 (2d12 + 4) piercing damage. If the target is Medium or smaller, the target must succeed on a DC 16 Strength saving throw or be knocked prone.

\DndMonsterAction{Claw}
\textsl{Melee Weapon Attack:} 8 to hit, reach 10 ft., one target. \textsl{Hit:} 13 (2d8 + 4) slashing damage, plus 10 (3d6) bludgeoning damage if the target is prone.

\DndMonsterSection{Reactions}
\DndMonsterAction{Bloody Rampage}
When the gray render takes damage, it makes one Claw attack against a random creature within its reach, other than its master.

\end{DndMonster}



\begin{DndMonster}[float*=b,width=\textwidth + 8pt]{Greater Tyrant Shadow}
\begin{multicols}{2}
\DndMonsterType{Gargantuan aberration, lawful evil}
\DndMonsterBasics[
armor-class = {18 (natural armor)},
hit-points = {\DndDice{24d20 + 120}},
speed = {30 ft.}
]
\DndMonsterAbilityScores[
 str = 25,
 dex = 15,
 con = 20,
 int = 24,
 wis = 16,
 cha = 22
]
\DndMonsterDetails[
saving-throws = {Int +14, Wis +10},
skills = {Arcana +14, Deception +20, Perception +10, Stealth +16},
damage-resistances = {fire},
damage-immunities = {necrotic; poison; psychic; bludgeoning, piercing, slashing from nonmagical attacks},
condition-immunities = {blinded, charmed, frightened, poisoned},
languages = {understands all, telepathy 300 ft.},
challenge = 22,
]
\DndMonsterAction{Empath}
The shadow uses telepathy to sense nearby creatures. The shadow can see a creature if that creature has thoughts and is within the range of the shadow's telepathy. The shadow can't see creatures immune to telepathy.

\DndMonsterAction{Undying Connection}
The shadow was originally spawned from a specific devil and, if the shadow is slain, it returns to life in 10 (1d20) days within 1 mile of that devil. The shadow is killed permanently only if its progenitor devil is slain first.

\DndMonsterSection{Actions}
\DndMonsterAction{Multiattack}
The shadow makes three Claw attacks. It can replace one of the attacks with Shadows of Death (if available).

\DndMonsterAction{Claw}
\textsl{Melee Weapon Attack:} 14 to hit, reach 5 ft., one target. \textsl{Hit:} 21 (4d6 + 7) necrotic damage, plus 10 (3d6) psychic damage.

\DndMonsterAction{Cloak of Shadows}
The shadow bends darkness around itself, and now has the invisible condition for 1 minute. The invisibility ends if the shadow is subjected to a \textit{Daylight} spell or similar.

\DndMonsterAction{Empathic Link}
The shadow attempts to form a link between itself and a creature it can see. The target must make a DC 22 Intelligence saving throw, taking 26 (4d12) psychic damage on a success. On a failed save, the creature and the shadow are linked. While linked, every time the shadow is damaged, the creature takes half of that damage as psychic damage. The creature may repeat the saving throw at the end of each of its turns, taking 26 (4d12) psychic damage on a failed save, or ending the effect on a successful one.

\DndMonsterAction{Shadows of Death (Recharge 5\textendash 6)}
The shadow discharges inky-black shadows in a 90-foot-radius sphere around itself. Creatures in the shadows must succeed on a DC 22 Wisdom saving throw or take 36 (8d8) necrotic damage and have the frightened condition for 1 minute. A frightened creature may repeat the saving throw at the end of each of its turns, ending the effect on a successful save, but taking the damage again on a failed save.

\DndMonsterSection{Bonus Actions}
\DndMonsterAction{Change Shape}
The shadow transforms into a Beast, Humanoid or Fiend that it has an Empathic Link with, or back into its true form. In a new form, the shadow retains its align ment, hit points, hit dice, telepathy, and Intelligence, Wisdom, and Charisma scores, as well as this action. Its statistics and capabilities are otherwise replaced by those of the new form, except any class features, lair actions, or legendary actions of that form.

\DndMonsterSection{Legendary Actions}
Greater Tyrant Shadow can take 3 legendary actions, choosing from the options below. Only one legendary action can be used at a time and only at the end of another creature's turn. Greater Tyrant Shadow regains spent legendary actions at the start of its turn.

\begin{DndMonsterLegendaryActions}
\DndMonsterLegendaryAction{Shadow Blink}{The shadow surrounds itself in black mist, then teleports up to 60 feet to an unoccupied space that it can see.
}
\DndMonsterLegendaryAction{Shapeshift}{The shadow uses its Change Shape ability.
}
\DndMonsterLegendaryAction{Mind Spike (Costs 2 Actions)}{\textsl{Ranged Spell Attack:} 14 to hit, range 300 ft., one target. \textsl{Hit:} 23 (3d10 + 7) psychic damage.
}
\end{DndMonsterLegendaryActions}
\DndMonsterSection{Lair Actions}
On initiative count 20 (losing initiative ties), the greater tyrant shadow can take one lair action to cause the following effect; it can't take the same lair action two rounds in a row:


\begin{description}
\item[Despair.]
One creature within 300 feet of the shadow must succeed on a DC 22 Charisma saving throw or be overcome with despair and self-doubt. The creature has the stunned condition for 1 minute. At the end of each of its turns, the creature can repeat the saving throw, ending the effect on a success.

\item[Renewal.]
The shadow recharges its Shadows of Death ability.

\end{description}

\end{multicols}
\end{DndMonster}



\begin{DndMonster}[float=t]{Green Abishai}
\DndMonsterType{Medium fiend (devil), lawful evil}
\DndMonsterBasics[
armor-class = {18 (natural armor)},
hit-points = {\DndDice{26d8 + 78}},
speed = {30 ft., fly 40 ft.}
]
\DndMonsterAbilityScores[
 str = 12,
 dex = 17,
 con = 16,
 int = 17,
 wis = 12,
 cha = 19
]
\DndMonsterDetails[
saving-throws = {Int +8, Cha +9},
skills = {Deception +9, Insight +6, Perception +6, Persuasion +9},
damage-resistances = {cold; bludgeoning, piercing, slashing from nonmagical attacks that aren't silvered},
damage-immunities = {fire, poison},
condition-immunities = {poisoned},
senses = {darkvision 120 ft., passive perception 16},
languages = {Draconic, Infernal, telepathy 120 ft.},
challenge = 15,
]
\DndMonsterAction{Devil's Sight}
Magical darkness doesn't impede the abishai's darkvision.

\DndMonsterAction{Magic Resistance}
The abishai has advantage on saving throws against spells and other magical effects.

\DndMonsterAction{Spellcasting}
The abishai casts one of the following spells, requiring no material components and using Charisma as the spellcasting ability (spell save DC 17):
\begin{DndMonsterSpells}
  \DndInnateSpellLevel{\textit{alter self}, \textit{major image}}
  \DndInnateSpellLevel[3e]{\textit{charm person}, \textit{detect thoughts}, \textit{fear}}
  \DndInnateSpellLevel[1e]{\textit{confusion}, \textit{dominate person}, \textit{mass suggestion}}
\end{DndMonsterSpells}
\DndMonsterSection{Actions}
\DndMonsterAction{Multiattack}
The abishai makes two Fiendish Claw attacks, or it makes one Fiendish Claw attack and uses Spellcasting.

\DndMonsterAction{Fiendish Claw}
\textsl{Melee Weapon Attack:} 8 to hit, reach 5 ft., one target. \textsl{Hit:} 12 (2d8 + 3) force damage. If the target is a creature, it must succeed on a DC 16 Constitution saving throw or take 16 (3d10) poison damage and become poisoned for 1 minute. The poisoned target can repeat the saving throw at the end of each of its turns, ending the effect on itself on a success.

\end{DndMonster}



\begin{DndMonster}[float=t]{Grell}
\DndMonsterType{Medium aberration, neutral evil}
\DndMonsterBasics[
armor-class = {12},
hit-points = {\DndDice{10d8 + 10}},
speed = {10 ft., fly 30 ft. \textit{(hover)}}
]
\DndMonsterAbilityScores[
 str = 15,
 dex = 14,
 con = 13,
 int = 12,
 wis = 11,
 cha = 9
]
\DndMonsterDetails[
skills = {Perception +4, Stealth +6},
damage-immunities = {lightning},
condition-immunities = {blinded, prone},
senses = {blindsight 60 ft. (blind beyond this radius), passive perception 14},
languages = {Grell},
challenge = 3,
]
\DndMonsterSection{Actions}
\DndMonsterAction{Multiattack}
The grell makes two attacks: one with its tentacles and one with its beak.

\DndMonsterAction{Tentacles}
\textsl{Melee Weapon Attack:} 4 to hit, reach 10 ft., one creature. \textsl{Hit:} 7 (1d10 + 2) piercing damage, and the target must succeed on a DC 11 Constitution saving throw or be poisoned for 1 minute. The poisoned target is paralyzed, and it can repeat the saving throw at the end of each of its turns, ending the effect on a success.

The target is also grappled (escape DC 15). If the target is Medium or smaller, it is also restrained until this grapple ends. While grappling the target, the grell has advantage on attack rolls against it and can 't use this attack against other targets. When the grell moves, any Medium or smaller target it is grappling moves with it.

\DndMonsterAction{Beak}
\textsl{Melee Weapon Attack:} 4 to hit, reach 5 ft., one target. \textsl{Hit:} 7 (2d4 + 2) piercing damage.

\end{DndMonster}



\begin{DndMonster}[float=t]{Grimlock}
\DndMonsterType{Medium humanoid (grimlock), neutral evil}
\DndMonsterBasics[
armor-class = {11},
hit-points = {\DndDice{2d8 + 2}},
speed = {30 ft.}
]
\DndMonsterAbilityScores[
 str = 16,
 dex = 12,
 con = 12,
 int = 9,
 wis = 8,
 cha = 6
]
\DndMonsterDetails[
skills = {Athletics +5, Perception +3, Stealth +3},
condition-immunities = {blinded},
senses = {blindsight 30 ft. or 10 ft. while deafened (blind beyond this radius), passive perception 13},
languages = {Undercommon},
challenge = 1/4,
]
\DndMonsterAction{Blind Senses}
The grimlock can't use its blindsight while deafened and unable to smell.

\DndMonsterAction{Keen Hearing and Smell}
The grimlock has advantage on Wisdom (Perception) checks that rely on hearing or smell.

\DndMonsterAction{Stone Camouflage}
The grimlock has advantage on Dexterity (Stealth) checks made to hide in rocky terrain.

\DndMonsterSection{Actions}
\DndMonsterAction{Spiked Bone Club}
\textsl{Melee Weapon Attack:} 5 to hit, reach 5 ft., one target. \textsl{Hit:} 5 (1d4 + 3) bludgeoning damage plus 2 (1d4) piercing damage.

\end{DndMonster}


\clearpage

\begin{DndMonster}[float=t]{Grotesque Tentacle}
\DndMonsterType{Large monstrosity, lawful good}
\DndMonsterBasics[
armor-class = {18 (natural armor)},
hit-points = {\DndDice{15d10 + 45}},
speed = {40 ft.}
]
\DndMonsterAbilityScores[
 str = 19,
 dex = 18,
 con = 16,
 int = 16,
 wis = 19,
 cha = 18
]
\DndMonsterDetails[
saving-throws = {Dex +8, Con +7, Int +7, Wis +8, Cha +8},
damage-immunities = {poison},
condition-immunities = {charmed, poisoned},
senses = {darkvision 60 ft., passive perception 14},
languages = {Celestial, Common},
challenge = 10,
]
\DndMonsterAction{Copy Warning}
This content has been copied from another creature, but the data replacement hasn't been run. Check details.

\DndMonsterAction{Rejuvenation}
If it dies, the naga returns to life in 1d6 days and regains all its hit points. Only a \textit{wish} spell can prevent this trait from functioning.

\DndMonsterAction{Spellcasting}
The naga is an 11th-level spellcaster. Its spellcasting ability is Wisdom (spell save DC 16, 8 to hit with spell attacks), and it needs only verbal components to cast its spells. It has the following cleric spells prepared:
\begin{DndMonsterSpells}
  \DndMonsterSpellLevel{\textit{mending}, \textit{sacred flame}, \textit{thaumaturgy}}
   \DndMonsterSpellLevel[1][4]{\textit{command}, \textit{cure wounds}, \textit{shield of faith}}
   \DndMonsterSpellLevel[2][3]{\textit{calm emotions}, \textit{hold person}}
   \DndMonsterSpellLevel[3][3]{\textit{bestow curse}, \textit{clairvoyance}}
   \DndMonsterSpellLevel[4][3]{\textit{banishment}, \textit{freedom of movement}}
   \DndMonsterSpellLevel[5][2]{\textit{flame strike}, \textit{geas}}
   \DndMonsterSpellLevel[6][1]{\textit{true seeing}}
\end{DndMonsterSpells}
\DndMonsterSection{Actions}
\DndMonsterAction{Bite}
\textsl{Melee Weapon Attack:} 8 to hit, reach 10 ft., one creature. \textsl{Hit:} 8 (1d8 + 4) piercing damage, and the target must make a DC 15 Constitution saving throw, taking 45 (10d8) poison damage on a failed save, or half as much damage on a successful one.

\DndMonsterAction{Spit Poison}
\textsl{Ranged Weapon Attack:} 8 to hit, range 15/30 ft., one creature. \textsl{Hit:} The target must make a DC 15 Constitution saving throw, taking 45 (10d8) poison damage on a failed save, or half as much damage on a successful one.

\end{DndMonster}



\begin{DndMonster}[float=t]{Guard}
\DndMonsterType{Medium humanoid (any race), any alignment}
\DndMonsterBasics[
armor-class = {16 (\textit{chain shirt})},
hit-points = {\DndDice{2d8 + 2}},
speed = {30 ft.}
]
\DndMonsterAbilityScores[
 str = 13,
 dex = 12,
 con = 12,
 int = 10,
 wis = 11,
 cha = 10
]
\DndMonsterDetails[
skills = {Perception +2},
languages = {any one language (usually Common)},
challenge = 1/8,
]
\DndMonsterSection{Actions}
\DndMonsterAction{Spear}
\textsl{Melee or Ranged Weapon Attack:} 3 to hit, reach 5 ft. or range 20/60 ft., one target. \textsl{Hit:} 4 (1d6 + 1) piercing damage, or 5 (1d8 + 1) piercing damage if used with two hands to make a melee attack.

\end{DndMonster}



\begin{DndMonster}[float=t]{Guardian Naga}
\DndMonsterType{Large monstrosity, lawful good}
\DndMonsterBasics[
armor-class = {18 (natural armor)},
hit-points = {\DndDice{15d10 + 45}},
speed = {40 ft.}
]
\DndMonsterAbilityScores[
 str = 19,
 dex = 18,
 con = 16,
 int = 16,
 wis = 19,
 cha = 18
]
\DndMonsterDetails[
saving-throws = {Dex +8, Con +7, Int +7, Wis +8, Cha +8},
damage-immunities = {poison},
condition-immunities = {charmed, poisoned},
senses = {darkvision 60 ft., passive perception 14},
languages = {Celestial, Common},
challenge = 10,
]
\DndMonsterAction{Rejuvenation}
If it dies, the naga returns to life in 1d6 days and regains all its hit points. Only a \textit{wish} spell can prevent this trait from functioning.

\DndMonsterAction{Spellcasting}
The naga is an 11th-level spellcaster. Its spellcasting ability is Wisdom (spell save DC 16, 8 to hit with spell attacks), and it needs only verbal components to cast its spells. It has the following cleric spells prepared:
\begin{DndMonsterSpells}
  \DndMonsterSpellLevel{\textit{mending}, \textit{sacred flame}, \textit{thaumaturgy}}
   \DndMonsterSpellLevel[1][4]{\textit{command}, \textit{cure wounds}, \textit{shield of faith}}
   \DndMonsterSpellLevel[2][3]{\textit{calm emotions}, \textit{hold person}}
   \DndMonsterSpellLevel[3][3]{\textit{bestow curse}, \textit{clairvoyance}}
   \DndMonsterSpellLevel[4][3]{\textit{banishment}, \textit{freedom of movement}}
   \DndMonsterSpellLevel[5][2]{\textit{flame strike}, \textit{geas}}
   \DndMonsterSpellLevel[6][1]{\textit{true seeing}}
\end{DndMonsterSpells}
\DndMonsterSection{Actions}
\DndMonsterAction{Bite}
\textsl{Melee Weapon Attack:} 8 to hit, reach 10 ft., one creature. \textsl{Hit:} 8 (1d8 + 4) piercing damage, and the target must make a DC 15 Constitution saving throw, taking 45 (10d8) poison damage on a failed save, or half as much damage on a successful one.

\DndMonsterAction{Spit Poison}
\textsl{Ranged Weapon Attack:} 8 to hit, range 15/30 ft., one creature. \textsl{Hit:} The target must make a DC 15 Constitution saving throw, taking 45 (10d8) poison damage on a failed save, or half as much damage on a successful one.

\end{DndMonster}



\begin{DndMonster}[float=t]{Halog}
\DndMonsterType{Small aberration, unaligned}
\DndMonsterBasics[
armor-class = {14 (natural armor)},
hit-points = {\DndDice{20d6 + 80}},
speed = {40 ft.}
]
\DndMonsterAbilityScores[
 str = 13,
 dex = 14,
 con = 18,
 int = 4,
 wis = 8,
 cha = 8
]
\DndMonsterDetails[
saving-throws = {Dex +5, Con +7},
skills = {Athletics +4, Perception +2, Survival +5},
condition-immunities = {poisoned},
senses = {darkvision 120 ft., passive perception 12},
challenge = 7,
]
\DndMonsterAction{Adaptive Fangs}
The halog magically alters its fangs to inflict more damage on the creatures it hits (included in the attack).

\DndMonsterAction{Magic Resistance}
The halog has advantage on saving throws against spells and other magical effects.

\DndMonsterAction{Pack Tactics}
The halog gains advantage on attacks targeting creatures adjacent to one or more of the halog's conscious allies.

\DndMonsterSection{Actions}
\DndMonsterAction{Multiattack}
The halog makes two Bite attacks.

\DndMonsterAction{Bite}
\textsl{Melee Weapon Attack:} 5 to hit, reach 5 ft., one target. \textsl{Hit:} 12 (3d6 + 2) piercing damage plus 10 (3d6) damage of a type that the target is most vulnerable to.

\DndMonsterSection{Reactions}
\DndMonsterAction{Adaptive Hide}
After taking damage, the halog is now invulnerable to the damage type of the damage it just took, and any future damage of that type instead heals the halog by the damage amount. These benefits last until the halog uses this reaction again.

\end{DndMonster}



\begin{DndMonster}[float=t]{Harpy}
\DndMonsterType{Medium monstrosity, chaotic evil}
\DndMonsterBasics[
armor-class = {11},
hit-points = {\DndDice{7d8 + 7}},
speed = {20 ft., fly 40 ft.}
]
\DndMonsterAbilityScores[
 str = 12,
 dex = 13,
 con = 12,
 int = 7,
 wis = 10,
 cha = 13
]
\DndMonsterDetails[
languages = {Common},
challenge = 1,
]
\DndMonsterSection{Actions}
\DndMonsterAction{Multiattack}
The harpy makes two attacks: one with its claws and one with its club.

\DndMonsterAction{Claws}
\textsl{Melee Weapon Attack:} 3 to hit, reach 5 ft., one target. \textsl{Hit:} 6 (2d4 + 1) slashing damage.

\DndMonsterAction{Club}
\textsl{Melee Weapon Attack:} 3 to hit, reach 5 ft., one target. \textsl{Hit:} 3 (1d4 + 1) bludgeoning damage.

\DndMonsterAction{Luring Song}
The harpy sings a magical melody. Every humanoid and giant within 300 feet of the harpy that can hear the song must succeed on a DC 11 Wisdom saving throw or be charmed until the song ends. The harpy must take a bonus action on its subsequent turns to continue singing. It can stop singing at any time. The song ends if the harpy is incapacitated.

While charmed by the harpy, a target is incapacitated and ignores the songs of other harpies. If the charmed target is more than 5 feet away from the harpy, the target must move on its turn toward the harpy by the most direct route. It doesn't avoid opportunity attacks, but before moving into damaging terrain, such as lava or a pit, and whenever it takes damage from a source other than the harpy, a target can repeat the saving throw. A creature can also repeat the saving throw at the end of each of its turns. If a creature's saving throw is successful, the effect ends on it.

A target that successfully saves is immune to this harpy's song for the next 24 hours.

\end{DndMonster}



\begin{DndMonster}[float=t]{Hell Hound}
\DndMonsterType{Medium fiend, lawful evil}
\DndMonsterBasics[
armor-class = {15 (natural armor)},
hit-points = {\DndDice{7d8 + 14}},
speed = {50 ft.}
]
\DndMonsterAbilityScores[
 str = 17,
 dex = 12,
 con = 14,
 int = 6,
 wis = 13,
 cha = 6
]
\DndMonsterDetails[
skills = {Perception +5},
damage-immunities = {fire},
senses = {darkvision 60 ft., passive perception 15},
languages = {understands Infernal but can't speak it},
challenge = 3,
]
\DndMonsterAction{Keen Hearing and Smell}
The hound has advantage on Wisdom (Perception) checks that rely on hearing or smell.

\DndMonsterAction{Pack Tactics}
The hound has advantage on an attack roll against a creature if at least one of the hound's allies is within 5 feet of the creature and the ally isn't incapacitated.

\DndMonsterSection{Actions}
\DndMonsterAction{Bite}
\textsl{Melee Weapon Attack:} 5 to hit, reach 5 ft., one target. \textsl{Hit:} 7 (1d8 + 3) piercing damage plus 7 (2d6) fire damage.

\DndMonsterAction{Fire Breath (Recharge 5\textendash 6)}
The hound exhales fire in a 15-foot cone. Each creature in that area must make a DC 12 Dexterity saving throw, taking 21 (6d6) fire damage on a failed save, or half as much damage on a successful one.

\end{DndMonster}



\begin{DndMonster}[float=t]{Hellcat (Bezekira)}
\DndMonsterType{Large fiend, neutral evil}
\DndMonsterBasics[
armor-class = {16 (natural armor)},
hit-points = {\DndDice{16d10 + 48}},
speed = {40 ft.}
]
\DndMonsterAbilityScores[
 str = 21,
 dex = 19,
 con = 17,
 int = 10,
 wis = 14,
 cha = 10
]
\DndMonsterDetails[
saving-throws = {Dex +8, Con +7},
skills = {Acrobatics +8, Athletics +9, Perception +6, Stealth +12},
damage-resistances = {bludgeoning, piercing, slashing from nonmagical attacks that aren't silvered; damage from spells},
damage-immunities = {fire},
senses = {blindsight 30 ft., darkvision 60 ft., passive perception 16},
languages = {understands Infernal, telepathy 120 ft.},
challenge = 10,
]
\DndMonsterAction{Darkness Dweller}
Attacks made against a bezekira in bright or dim light have disadvantage.

\DndMonsterAction{Devil's Sight}
Magical darkness doesn't impede the bezekira's darkvision.

\DndMonsterSection{Actions}
\DndMonsterAction{Multiattack}
The bezekira makes three Claw attacks. It can replace one of the attacks with a Bite attack.

\DndMonsterAction{Claw}
\textsl{Melee Weapon Attack:} 9 to hit, reach 5 ft., one target. \textsl{Hit:} 15 (3d6 + 5) slashing damage plus 7 (2d6) fire damage. The bezekira deals an extra 28 (8d6) fire damage when it hits a target with an attack and has advantage on the attack roll.

\DndMonsterAction{Bite}
\textsl{Melee Weapon Attack:} 9 to hit, reach 5 ft., one target. \textsl{Hit:} 16 (2d10 + 5) bludgeoning damage, grasping the target with its mouth. If the target is a Medium or smaller creature, it has the grappled condition (escape DC 17). The bezekira can't make a Bite attack while a creature is grappled in this way.

\end{DndMonster}



\begin{DndMonster}[float=t]{Hezrou}
\DndMonsterType{Large fiend (demon), chaotic evil}
\DndMonsterBasics[
armor-class = {16 (natural armor)},
hit-points = {\DndDice{13d10 + 65}},
speed = {30 ft.}
]
\DndMonsterAbilityScores[
 str = 19,
 dex = 17,
 con = 20,
 int = 5,
 wis = 12,
 cha = 13
]
\DndMonsterDetails[
saving-throws = {Str +7, Con +8, Wis +4},
damage-resistances = {cold; fire; lightning; bludgeoning, piercing, slashing from nonmagical attacks},
damage-immunities = {poison},
condition-immunities = {poisoned},
senses = {darkvision 120 ft., passive perception 11},
languages = {Abyssal, telepathy 120 ft.},
challenge = 8,
]
\DndMonsterAction{Magic Resistance}
The hezrou has advantage on saving throws against spells and other magical effects.

\DndMonsterAction{Stench}
Any creature that starts its turn within 10 feet of the hezrou must succeed on a DC 14 Constitution saving throw or be poisoned until the start of its next turn. On a successful saving throw, the creature is immune to the hezrou's stench for 24 hours.

\DndMonsterSection{Actions}
\DndMonsterAction{Multiattack}
The hezrou makes three attacks: one with its bite and two with its claws.

\DndMonsterAction{Bite}
\textsl{Melee Weapon Attack:} 7 to hit, reach 5 ft., one target. \textsl{Hit:} 15 (2d10 + 4) piercing damage.

\DndMonsterAction{Claws}
\textsl{Melee Weapon Attack:} 7 to hit, reach 5 ft., one target. \textsl{Hit:} 11 (2d6 + 4) slashing damage.

\end{DndMonster}



\begin{DndMonster}[float=t]{Hill Giant}
\DndMonsterType{Huge giant, chaotic evil}
\DndMonsterBasics[
armor-class = {13 (natural armor)},
hit-points = {\DndDice{10d12 + 40}},
speed = {40 ft.}
]
\DndMonsterAbilityScores[
 str = 21,
 dex = 8,
 con = 19,
 int = 5,
 wis = 9,
 cha = 6
]
\DndMonsterDetails[
skills = {Perception +2},
languages = {Giant},
challenge = 5,
]
\DndMonsterSection{Actions}
\DndMonsterAction{Multiattack}
The giant makes two greatclub attacks.

\DndMonsterAction{Greatclub}
\textsl{Melee Weapon Attack:} 8 to hit, reach 10 ft., one target. \textsl{Hit:} 18 (3d8 + 5) bludgeoning damage.

\DndMonsterAction{Rock}
\textsl{Ranged Weapon Attack:} 8 to hit, range 60/240 ft., one target. \textsl{Hit:} 21 (3d10 + 5) bludgeoning damage.

\end{DndMonster}



\begin{DndMonster}[float=t]{Hobgoblin}
\DndMonsterType{Medium humanoid (goblinoid), lawful evil}
\DndMonsterBasics[
armor-class = {18 (\textit{chain mail})},
hit-points = {\DndDice{2d8 + 2}},
speed = {30 ft.}
]
\DndMonsterAbilityScores[
 str = 13,
 dex = 12,
 con = 12,
 int = 10,
 wis = 10,
 cha = 9
]
\DndMonsterDetails[
senses = {darkvision 60 ft., passive perception 10},
languages = {Common, Goblin},
challenge = 1/2,
]
\DndMonsterAction{Martial Advantage}
Once per turn, the hobgoblin can deal an extra 7 (2d6) damage to a creature it hits with a weapon attack if that creature is within 5 feet of an ally of the hobgoblin that isn't incapacitated.

\DndMonsterSection{Actions}
\DndMonsterAction{Longsword}
\textsl{Melee Weapon Attack:} 3 to hit, reach 5 ft., one target. \textsl{Hit:} 5 (1d8 + 1) slashing damage, or 6 (1d10 + 1) slashing damage if used with two hands.

\DndMonsterAction{Longbow}
\textsl{Ranged Weapon Attack:} 3 to hit, range 150/600 ft., one target. \textsl{Hit:} 5 (1d8 + 1) piercing damage.

\end{DndMonster}



\begin{DndMonster}[float=t]{Hook Horror}
\DndMonsterType{Large monstrosity, neutral}
\DndMonsterBasics[
armor-class = {15 (natural armor)},
hit-points = {\DndDice{10d10 + 20}},
speed = {30 ft., climb 30 ft.}
]
\DndMonsterAbilityScores[
 str = 18,
 dex = 10,
 con = 15,
 int = 6,
 wis = 12,
 cha = 7
]
\DndMonsterDetails[
skills = {Perception +3},
senses = {blindsight 60 ft., darkvision 120 ft., passive perception 13},
languages = {Hook Horror},
challenge = 3,
]
\DndMonsterAction{Echolocation}
The hook horror can't use its blindsight while deafened.

\DndMonsterAction{Keen Hearing}
The hook horror has advantage on Wisdom (Perception) checks that rely on hearing.

\DndMonsterSection{Actions}
\DndMonsterAction{Multiattack}
The hook horror makes two hook attacks.

\DndMonsterAction{Hook}
\textsl{Melee Weapon Attack:} 6 to hit, reach 10 ft., one target. \textsl{Hit:} 11 (2d6 + 4) piercing damage.

\end{DndMonster}



\begin{DndMonster}[float=t]{Horned Devil}
\DndMonsterType{Large fiend (devil), lawful evil}
\DndMonsterBasics[
armor-class = {18 (natural armor)},
hit-points = {\DndDice{17d10 + 85}},
speed = {20 ft., fly 60 ft.}
]
\DndMonsterAbilityScores[
 str = 22,
 dex = 17,
 con = 21,
 int = 12,
 wis = 16,
 cha = 17
]
\DndMonsterDetails[
saving-throws = {Str +10, Dex +7, Wis +7, Cha +7},
damage-resistances = {cold; bludgeoning, piercing, slashing from nonmagical attacks that aren't silvered},
damage-immunities = {fire, poison},
condition-immunities = {poisoned},
senses = {darkvision 120 ft., passive perception 13},
languages = {Infernal, telepathy 120 ft.},
challenge = 11,
]
\DndMonsterAction{Devil's Sight}
Magical darkness doesn't impede the devil's darkvision.

\DndMonsterAction{Magic Resistance}
The devil has advantage on saving throws against spells and other magical effects.

\DndMonsterSection{Actions}
\DndMonsterAction{Multiattack}
The devil makes three melee attacks: two with its fork and one with its tail. It can use Hurl Flame in place of any melee attack.

\DndMonsterAction{Fork}
\textsl{Melee Weapon Attack:} 10 to hit, reach 10 ft., one target. \textsl{Hit:} 15 (2d8 + 6) piercing damage.

\DndMonsterAction{Tail}
\textsl{Melee Weapon Attack:} 10 to hit, reach 10 ft., one target. \textsl{Hit:} 10 (1d8 + 6) piercing damage. If the target is a creature other than an undead or a construct, it must succeed on a DC 17 Constitution saving throw or lose 10 (3d6) hit points at the start of each of its turns due to an infernal wound. Each time the devil hits the wounded target with this attack, the damage dealt by the wound increases by 10 (3d6). Any creature can take an action to stanch the wound with a successful DC 12 Wisdom (Medicine) check. The wound also closes if the target receives magical healing.

\DndMonsterAction{Hurl Flame}
\textsl{Ranged Spell Attack:} 7 to hit, range 150 ft., one target. \textsl{Hit:} 14 (4d6) fire damage. If the target is a flammable object that isn't being worn or carried, it also catches fire.

\end{DndMonster}



\begin{DndMonster}[float=t]{Ice Devil}
\DndMonsterType{Large fiend (devil), lawful evil}
\DndMonsterBasics[
armor-class = {18 (natural armor)},
hit-points = {\DndDice{19d10 + 76}},
speed = {40 ft.}
]
\DndMonsterAbilityScores[
 str = 21,
 dex = 14,
 con = 18,
 int = 18,
 wis = 15,
 cha = 18
]
\DndMonsterDetails[
saving-throws = {Dex +7, Con +9, Wis +7, Cha +9},
damage-resistances = {bludgeoning, piercing, slashing from nonmagical attacks that aren't silvered},
damage-immunities = {fire, poison, cold},
condition-immunities = {poisoned},
senses = {blindsight 60 ft., darkvision 120 ft., passive perception 12},
languages = {Infernal, telepathy 120 ft.},
challenge = 14,
]
\DndMonsterAction{Devil's Sight}
Magical darkness doesn't impede the devil's darkvision.

\DndMonsterAction{Magic Resistance}
The devil has advantage on saving throws against spells and other magical effects.

\DndMonsterSection{Actions}
\DndMonsterAction{Multiattack}
The devil makes three attacks: one with its bite, one with its claws, and one with its tail.

\DndMonsterAction{Bite}
\textsl{Melee Weapon Attack:} 10 to hit, reach 5 ft., one target. \textsl{Hit:} 12 (2d6 + 5) piercing damage plus 10 (3d6) cold damage.

\DndMonsterAction{Claws}
\textsl{Melee Weapon Attack:} 10 to hit, reach 5 ft., one target. \textsl{Hit:} 10 (2d4 + 5) slashing damage plus 10 (3d6) cold damage.

\DndMonsterAction{Tail}
\textsl{Melee Weapon Attack:} 10 to hit, reach 10 ft., one target. \textsl{Hit:} 12 (2d6 + 5) bludgeoning damage plus 10 (3d6) cold damage.

\DndMonsterAction{Wall of Ice (Recharge 6)}
The devil magically forms an opaque wall of ice on a solid surface it can see within 60 feet of it. The wall is 1 foot thick and up to 30 feet long and 10 feet high, or it's a hemispherical dome up to 20 feet in diameter.

When the wall appears, each creature in its space is pushed out of it by the shortest route. The creature chooses which side of the wall to end up on, unless the creature is incapacitated. The creature then makes a DC 17 Dexterity saving throw, taking 35 (10d6) cold damage on a failed save, or half as much damage on a successful one.

The wall lasts for 1 minute or until the devil is incapacitated or dies. The wall can be damaged and breached; each 10-foot section has AC 5, 30 hit points, vulnerability to fire damage, and immunity to acid, cold, necrotic, poison, and psychic damage. If a section is destroyed, it leaves behind a sheet of frigid air in the space the wall occupied. Whenever a creature finishes moving through the frigid air on a turn, willingly or otherwise, the creature must make a DC 17 Constitution saving throw, taking 17 (5d6) cold damage on a failed save, or half as much damage on a successful one. The frigid air dissipates when the rest of the wall vanishes.

\end{DndMonster}



\begin{DndMonster}[float=t]{Icy Simulacrum}
\DndMonsterType{Medium monstrosity (shapechanger), neutral}
\DndMonsterBasics[
armor-class = {14},
hit-points = {\DndDice{8d8 + 16}},
speed = {30 ft.}
]
\DndMonsterAbilityScores[
 str = 11,
 dex = 18,
 con = 14,
 int = 11,
 wis = 12,
 cha = 14
]
\DndMonsterDetails[
skills = {Deception +6, Insight +3},
condition-immunities = {charmed},
senses = {darkvision 60 ft., passive perception 11},
languages = {Common},
challenge = 3,
]
\DndMonsterAction{Copy Warning}
This content has been copied from another creature, but the data replacement hasn't been run. Check details.

\DndMonsterAction{Shapechanger}
The doppelganger can use its action to polymorph into a Small or Medium humanoid it has seen, or back into its true form. Its statistics, other than its size, are the same in each form. Any equipment it is wearing or carrying isn't transformed. It reverts to its true form if it dies.

\DndMonsterAction{Ambusher}
In the first round of a combat, the doppelganger has advantage on attack rolls against any creature it Surprise.

\DndMonsterAction{Surprise Attack}
If the doppelganger surprises a creature and hits it with an attack during the first round of combat, the target takes an extra 10 (3d6) damage from the attack.

\DndMonsterSection{Actions}
\DndMonsterAction{Multiattack}
The doppelganger makes two melee attacks.

\DndMonsterAction{Slam}
\textsl{Melee Weapon Attack:} 6 to hit, reach 5 ft., one target. \textsl{Hit:} 7 (1d6 + 4) bludgeoning damage.

\DndMonsterAction{Read Thoughts}
The doppelganger magically reads the surface thoughts of one creature within 60 feet of it. The effect can penetrate barriers, but 3 feet of wood or dirt, 2 feet of stone, 2 inches of metal, or a thin sheet of lead blocks it. While the target is in range, the doppelganger can continue reading its thoughts, as long as the doppelganger's concentration isn't broken (as if concentration on a spell). While reading the target's mind, the doppelganger has advantage on Wisdom (Insight) and Charisma (Deception, Intimidation, and Persuasion) checks against the target.

\end{DndMonster}



\begin{DndMonster}[float=t]{Imp}
\DndMonsterType{Tiny fiend (devil), lawful evil}
\DndMonsterBasics[
armor-class = {13},
hit-points = {\DndDice{3d4 + 3}},
speed = {20 ft., fly 40 ft.}
]
\DndMonsterAbilityScores[
 str = 6,
 dex = 17,
 con = 13,
 int = 11,
 wis = 12,
 cha = 14
]
\DndMonsterDetails[
skills = {Deception +4, Insight +3, Persuasion +4, Stealth +5},
damage-resistances = {cold; bludgeoning, piercing, slashing from nonmagical attacks not made with silvered weapons},
damage-immunities = {fire, poison},
condition-immunities = {poisoned},
senses = {darkvision 120 ft., passive perception 11},
languages = {Infernal, Common},
challenge = 1,
]
\DndMonsterAction{Shapechanger}
The imp can use its action to polymorph into a beast form that resembles a rat (speed 20 ft.), a raven (20 ft., fly 60 ft.), or a spider (20 ft., climb 20 ft.), or back into its true form. Its statistics are the same in each form, except for the speed changes noted. Any equipment it is wearing or carrying isn't transformed. It reverts to its true form if it dies.

\DndMonsterAction{Devil's Sight}
Magical darkness doesn't impede the imp's darkvision.

\DndMonsterAction{Magic Resistance}
The imp has advantage on saving throws against spells and other magical effects.

\DndMonsterSection{Actions}
\DndMonsterAction{Sting (Bite in Beast Form)}
\textsl{Melee Weapon Attack:} 5 to hit, reach 5 ft., one target. \textsl{Hit:} 5 (1d4 + 3) piercing damage, and the target must make a DC 11 Constitution saving throw, taking 10 (3d6) poison damage on a failed save, or half as much damage on a successful one.

\DndMonsterAction{Invisibility}
The imp magically turns invisible until it attacks, or until its concentration ends (as if concentration on a spell). Any equipment the imp wears or carries is invisible with it.

\end{DndMonster}



\begin{DndMonster}[float=t]{Incubus}
\DndMonsterType{Medium fiend (shapechanger), neutral evil}
\DndMonsterBasics[
armor-class = {15 (natural armor)},
hit-points = {\DndDice{12d8 + 12}},
speed = {30 ft., fly 60 ft.}
]
\DndMonsterAbilityScores[
 str = 8,
 dex = 17,
 con = 13,
 int = 15,
 wis = 12,
 cha = 20
]
\DndMonsterDetails[
skills = {Deception +9, Insight +5, Perception +5, Persuasion +9, Stealth +7},
damage-resistances = {cold; fire; lightning; poison; bludgeoning, piercing, slashing from nonmagical attacks},
senses = {darkvision 60 ft., passive perception 15},
languages = {Abyssal, Common, Infernal, telepathy 60 ft.},
challenge = 4,
]
\DndMonsterAction{Telepathic Bond}
The fiend ignores the range restriction on its telepathy when communicating with a creature it has charmed. The two don't even need to be on the same plane of existence.

\DndMonsterAction{Shapechanger}
The fiend can use its action to polymorph into a Small or Medium humanoid, or back into its true form. Without wings, the fiend loses its flying speed. Other than its size and speed, its statistics are the same in each form. Any equipment it is wearing or carrying isn't transformed. It reverts to its true form if it dies.

\DndMonsterSection{Actions}
\DndMonsterAction{Claw (Fiend Form Only)}
\textsl{Melee Weapon Attack:} 5 to hit, reach 5 ft., one target. \textsl{Hit:} 6 (1d6 + 3) slashing damage.

\DndMonsterAction{Charm}
One humanoid the fiend can see within 30 feet of it must succeed on a DC 15 Wisdom saving throw or be magically charmed for 1 day. The charmed target obeys the fiend's verbal or telepathic commands. If the target suffers any harm or receives a suicidal command, it can repeat the saving throw, ending the effect on a success. If the target successfully saves against the effect, or if the effect on it ends, the target is immune to this fiend's Charm for the next 24 hours.

The fiend can have only one target charmed at a time. If it charms another, the effect on the previous target ends.

\DndMonsterAction{Draining Kiss}
The fiend kisses a creature charmed by it or a willing creature. The target must make a DC 15 Constitution saving throw against this magic, taking 32 (5d10 + 5) psychic damage on a failed save, or half as much damage on a successful one. The target's hit point maximum is reduced by an amount equal to the damage taken. This reduction lasts until the target finishes a long rest. The target dies if this effect reduces its hit point maximum to 0.

\DndMonsterAction{Etherealness}
The fiend magically enters the Ethereal Plane from the Material Plane, or vice versa.

\end{DndMonster}



\begin{DndMonster}[float=t]{Intellect Devourer}
\DndMonsterType{Tiny aberration, lawful evil}
\DndMonsterBasics[
armor-class = {12},
hit-points = {\DndDice{6d4 + 6}},
speed = {40 ft.}
]
\DndMonsterAbilityScores[
 str = 6,
 dex = 14,
 con = 13,
 int = 12,
 wis = 11,
 cha = 10
]
\DndMonsterDetails[
skills = {Perception +2, Stealth +4},
damage-resistances = {bludgeoning, piercing, slashing from nonmagical attacks},
condition-immunities = {blinded},
senses = {blindsight 60 ft. (blind beyond this radius), passive perception 12},
languages = {understands Deep Speech but can't speak, telepathy 60 ft.},
challenge = 2,
]
\DndMonsterAction{Detect Sentience}
The intellect devourer can sense the presence and location of any creature within 300 feet of it that has an Intelligence of 3 or higher, regardless of interposing barriers, unless the creature is protected by a \textit{mind blank} spell.

\DndMonsterSection{Actions}
\DndMonsterAction{Multiattack}
The intellect devourer makes one attack with its claws and uses Devour Intellect.

\DndMonsterAction{Claws}
\textsl{Melee Weapon Attack:} 4 to hit, reach 5 ft., one target. \textsl{Hit:} 7 (2d4 + 2) slashing damage.

\DndMonsterAction{Devour Intellect}
The intellect devourer targets one creature it can see within 10 feet of it that has a brain. The target must succeed on a DC 12 Intelligence saving throw against this magic or take 11 (2d10) psychic damage. Also on a failure, roll 3d6: If the total equals or exceeds the target's Intelligence score, that score is reduced to 0. The target is stunned until it regains at least one point of Intelligence.

\DndMonsterAction{Body Thief}
The intellect devourer initiates an Intelligence contest with an incapacitated humanoid within 5 feet of it that isn't protected by \textit{protection from evil and good}. If it wins the contest, the intellect devourer magically consumes the target's brain, teleports into the target's skull, and takes control of the target's body. While inside a creature, the intellect devourer has total cover against attacks and other effects originating outside its host. The intellect devourer retains its Intelligence, Wisdom, and Charisma scores, as well as its understanding of Deep Speech, its telepathy, and its traits. It otherwise adopts the target's statistics. It knows everything the creature knew, including spells and languages.

If the host body dies, the intellect devourer must leave it. A \textit{protection from evil and good} spell cast on the body drives the intellect devourer out. The intellect devourer is also forced out if the target regains its devoured brain by means of a \textit{wish}. By spending 5 feet of its movement, the intellect devourer can voluntarily leave the body, teleporting to the nearest unoccupied space within 5 feet of it. The body then dies, unless its brain is restored within 1 round.

\end{DndMonster}



\begin{DndMonster}[float=t]{Iron Golem}
\DndMonsterType{Large construct, unaligned}
\DndMonsterBasics[
armor-class = {20 (natural armor)},
hit-points = {\DndDice{20d10 + 100}},
speed = {30 ft.}
]
\DndMonsterAbilityScores[
 str = 24,
 dex = 9,
 con = 20,
 int = 3,
 wis = 11,
 cha = 1
]
\DndMonsterDetails[
damage-immunities = {fire; poison; psychic; bludgeoning, piercing, slashing from nonmagical attacks that aren't adamantine},
condition-immunities = {charmed, exhaustion, frightened, paralyzed, petrified, poisoned},
senses = {darkvision 120 ft., passive perception 10},
languages = {understands the languages of its creator but can't speak},
challenge = 16,
]
\DndMonsterAction{Fire Absorption}
Whenever the golem is subjected to fire damage, it takes no damage and instead regains a number of hit points equal to the fire damage dealt.

\DndMonsterAction{Immutable Form}
The golem is immune to any spell or effect that would alter its form.

\DndMonsterAction{Magic Resistance}
The golem has advantage on saving throws against spells and other magical effects.

\DndMonsterAction{Magic Weapons}
The golem's weapon attacks are magical.

\DndMonsterSection{Actions}
\DndMonsterAction{Multiattack}
The golem makes two melee attacks.

\DndMonsterAction{Slam}
\textsl{Melee Weapon Attack:} 13 to hit, reach 5 ft., one target. \textsl{Hit:} 20 (3d8 + 7) bludgeoning damage.

\DndMonsterAction{Sword}
\textsl{Melee Weapon Attack:} 13 to hit, reach 10 ft., one target. \textsl{Hit:} 23 (3d10 + 7) slashing damage.

\DndMonsterAction{Poison Breath (Recharge 5\textendash 6)}
The golem exhales poisonous gas in a 15-foot cone. Each creature in that area must make a DC 19 Constitution saving throw, taking 45 (10d8) poison damage on a failed save, or half as much damage on a successful one.

\end{DndMonster}



\begin{DndMonster}[float=t]{Jenevere}
\DndMonsterType{Large celestial, lawful good}
\DndMonsterBasics[
armor-class = {20 (natural armor)},
hit-points = {\DndDice{16d10 + 128}},
speed = {40 ft., fly 120 ft.}
]
\DndMonsterAbilityScores[
 str = 18,
 dex = 22,
 con = 26,
 int = 18,
 wis = 20,
 cha = 24
]
\DndMonsterDetails[
saving-throws = {Int +10, Wis +11, Cha +13},
skills = {Perception +11},
damage-resistances = {radiant; bludgeoning, piercing, slashing from nonmagical attacks},
condition-immunities = {charmed, exhaustion, frightened},
senses = {truesight 120 ft., passive perception 21},
languages = {all, telepathy 120 ft.},
challenge = 17,
]
\DndMonsterAction{Magic Resistance}
Jenevere has advantage on saving throws against spells and other magical effects.

\DndMonsterAction{Spellcasting}
Jenevere casts one of the following spells, requiring no material components and using Charisma as the spellcasting ability (spell save DC 21):
\begin{DndMonsterSpells}
  \DndInnateSpellLevel{\textit{Bless}, \textit{Detect Evil and Good}}
  \DndInnateSpellLevel[1]{\textit{Foresight}}
  \DndInnateSpellLevel[2e]{\textit{Commune}, \textit{Heal}, \textit{Legend Lore}, \textit{Raise Dead}, \textit{Scrying}}
\end{DndMonsterSpells}
\DndMonsterSection{Actions}
\DndMonsterAction{Multiattack}
Jenevere makes three Radiant Touch attacks.

\DndMonsterAction{Radiant Touch}
\textsl{Melee Weapon Attack:} 12 to hit, reach 5 ft., one target. \textsl{Hit:} 33 (6d8 + 6) radiant damage.

\DndMonsterAction{Forgiveness (2/Day)}
Jenevere selects a target within 10 feet. The target is freed from any curse, disease, poison, blindness, or deafness. If the target is a Fiend, it must succeed on a DC 21 saving throw or have disadvantage on all attack rolls for 24 hours.

\end{DndMonster}



\begin{DndMonster}[float=t]{Knight}
\DndMonsterType{Medium humanoid (any race), any alignment}
\DndMonsterBasics[
armor-class = {18 (\textit{plate armor})},
hit-points = {\DndDice{8d8 + 16}},
speed = {30 ft.}
]
\DndMonsterAbilityScores[
 str = 16,
 dex = 11,
 con = 14,
 int = 11,
 wis = 11,
 cha = 15
]
\DndMonsterDetails[
saving-throws = {Con +4, Wis +2},
languages = {any one language (usually Common)},
challenge = 3,
]
\DndMonsterAction{Brave}
The knight has advantage on saving throws against being frightened.

\DndMonsterSection{Actions}
\DndMonsterAction{Multiattack}
The knight makes two melee attacks.

\DndMonsterAction{Greatsword}
\textsl{Melee Weapon Attack:} 5 to hit, reach 5 ft., one target. \textsl{Hit:} 10 (2d6 + 3) slashing damage.

\DndMonsterAction{Heavy Crossbow}
\textsl{Ranged Weapon Attack:} 2 to hit, range 100/400 ft., one target. \textsl{Hit:} 5 (1d10) piercing damage.

\DndMonsterAction{Leadership (Recharges after a Short or Long Rest)}
For 1 minute, the knight can utter a special command or warning whenever a nonhostile creature that it can see within 30 feet of it makes an attack roll or a saving throw. The creature can add a d4 to its roll provided it can hear and understand the knight. A creature can benefit from only one Leadership die at a time. This effect ends if the knight is incapacitated.

\DndMonsterSection{Reactions}
\DndMonsterAction{Parry}
The knight adds 2 to its AC against one melee attack that would hit it. To do so, the knight must see the attacker and be wielding a melee weapon.

\end{DndMonster}


\clearpage

\begin{DndMonster}[float=t]{Kobold}
\DndMonsterType{Small humanoid (kobold), lawful evil}
\DndMonsterBasics[
armor-class = {12},
hit-points = {\DndDice{2d6 - 2}},
speed = {30 ft.}
]
\DndMonsterAbilityScores[
 str = 7,
 dex = 15,
 con = 9,
 int = 8,
 wis = 7,
 cha = 8
]
\DndMonsterDetails[
senses = {darkvision 60 ft., passive perception 8},
languages = {Common, Draconic},
challenge = 1/8,
]
\DndMonsterAction{Sunlight Sensitivity}
While in sunlight, the kobold has disadvantage on attack rolls, as well as on Wisdom (Perception) checks that rely on sight.

\DndMonsterAction{Pack Tactics}
The kobold has advantage on an attack roll against a creature if at least one of the kobold's allies is within 5 feet of the creature and the ally isn't incapacitated.

\DndMonsterSection{Actions}
\DndMonsterAction{Dagger}
\textsl{Melee Weapon Attack:} 4 to hit, reach 5 ft., one target. \textsl{Hit:} 4 (1d4 + 2) piercing damage.

\DndMonsterAction{Sling}
\textsl{Ranged Weapon Attack:} 4 to hit, range 30/120 ft., one target. \textsl{Hit:} 4 (1d4 + 2) bludgeoning damage.

\end{DndMonster}



\begin{DndMonster}[float=t]{Koh Tam}
\DndMonsterType{Medium humanoid, lawful neutral}
\DndMonsterBasics[
armor-class = {17 (\textit{breastplate})},
hit-points = {\DndDice{22d8 + 44}},
speed = {30 ft.}
]
\DndMonsterAbilityScores[
 str = 16,
 dex = 12,
 con = 15,
 int = 19,
 wis = 20,
 cha = 15
]
\DndMonsterDetails[
saving-throws = {Wis +9, Cha +6},
skills = {Arcana +8, History +8, Insight +9, Religion +8},
senses = {darkvision 120 ft., passive perception 15},
languages = {Abyssal, Celestial, Common, Infernal},
challenge = 10,
]
\DndMonsterAction{Spellcasting}
Koh Tam casts one of the following spells, using Wisdom as the spellcasting ability (spell save DC 17):
\begin{DndMonsterSpells}
  \DndInnateSpellLevel{\textit{False Life}, \textit{Gentle Repose}, \textit{Lesser Restoration}, \textit{Protection from Evil and Good}, \textit{Speak with Dead}}
  \DndInnateSpellLevel[1e]{\textit{Antilife Shell}, \textit{Blight}, \textit{Death Ward}, \textit{Dispel Evil and Good}, \textit{Greater Restoration}, \textit{Raise Dead}, \textit{Regenerate}, \textit{Revivify}, \textit{Spiritual Weapon}}
\end{DndMonsterSpells}
\DndMonsterSection{Actions}
\DndMonsterAction{Arbiter's Touch}
\textsl{Melee Spell Attack:} 9 to hit, reach 5 ft., one target. \textsl{Hit:} 12 (2d6 + 5) radiant damage.

\DndMonsterAction{Judgement}
Koh Tam judges a creature within 60 feet of him. The target must succeed on a DC 17 Wisdom saving throw or take 50 (10d8 + 5) radiant damage. If the target had fewer hit points than their maximum prior to this attack, the creature takes 70 (10d12 + 5) radiant damage instead.

\DndMonsterAction{Arbiter's Light (Recharge 4\textendash 6)}
\textsl{Ranged Spell Attack:} 9 to hit, range 120 ft., one target. \textsl{Hit:} 70 (10d12 + 5) radiant damage, and the target has the blinded condition until the start of Koh Tam's next turn.

\end{DndMonster}



\begin{DndMonster}[float*=b,width=\textwidth + 8pt]{Kraken}
\begin{multicols}{2}
\DndMonsterType{Gargantuan monstrosity (titan), chaotic evil}
\DndMonsterBasics[
armor-class = {18 (natural armor)},
hit-points = {\DndDice{27d20 + 189}},
speed = {20 ft., swim 60 ft.}
]
\DndMonsterAbilityScores[
 str = 30,
 dex = 11,
 con = 25,
 int = 22,
 wis = 18,
 cha = 20
]
\DndMonsterDetails[
saving-throws = {Str +17, Dex +7, Con +14, Int +13, Wis +11},
damage-immunities = {lightning; bludgeoning, piercing, slashing from nonmagical attacks},
condition-immunities = {frightened, paralyzed},
senses = {truesight 120 ft., passive perception 14},
languages = {Abyssal, Celestial, Infernal, Primordial, telepathy 120 ft. but can't speak},
challenge = 23,
]
\DndMonsterAction{Amphibious}
The kraken can breathe air and water.

\DndMonsterAction{Freedom of Movement}
The kraken ignores difficult terrain, and magical effects can't reduce its speed or cause it to be restrained. It can spend 5 feet of movement to escape from nonmagical restraints or being grappled.

\DndMonsterAction{Siege Monster}
The kraken deals double damage to objects and structures.

\DndMonsterSection{Actions}
\DndMonsterAction{Multiattack}
The kraken makes three tentacle attacks, each of which it can replace with one use of Fling.

\DndMonsterAction{Bite}
\textsl{Melee Weapon Attack:} 17 to hit, reach 5 ft., one target. \textsl{Hit:} 23 (3d8 + 10) piercing damage. If the target is a Large or smaller creature grappled by the kraken, that creature is swallowed, and the grapple ends. While swallowed, the creature is blinded and restrained, it has total cover against attacks and other effects outside the kraken, and it takes 42 (12d6) acid damage at the start of each of the kraken's turns. If the kraken takes 50 damage or more on a single turn from a creature inside it, the kraken must succeed on a DC 25 Constitution saving throw at the end of that turn or regurgitate all swallowed creatures, which fall prone in a space within 10 feet of the kraken. If the kraken dies, a swallowed creature is no longer restrained by it and can escape from the corpse using 15 feet of movement, exiting prone.

\DndMonsterAction{Tentacle}
\textsl{Melee Weapon Attack:} 17 to hit, reach 30 ft., one target. \textsl{Hit:} 20 (3d6 + 10) bludgeoning damage, and the target is grappled (escape DC 18). Until this grapple ends, the target is restrained. The kraken has ten tentacles, each of which can grapple one target.

\DndMonsterAction{Fling}
One Large or smaller object held or creature grappled by the kraken is thrown up to 60 feet in a random direction and knocked prone. If a thrown target strikes a solid surface, the target takes 3 (1d6) bludgeoning damage for every 10 feet it was thrown. If the target is thrown at another creature, that creature must succeed on a DC 18 Dexterity saving throw or take the same damage and be knocked prone.

\DndMonsterAction{Lightning Storm}
The kraken magically creates three bolts of lightning, each of which can strike a target the kraken can see within 120 feet of it. A target must make a DC 23 Dexterity saving throw, taking 22 (4d10) lightning damage on a failed save, or half as much damage on a successful one.

\DndMonsterSection{Legendary Actions}
Kraken can take 3 legendary actions, choosing from the options below. Only one legendary action can be used at a time and only at the end of another creature's turn. Kraken regains spent legendary actions at the start of its turn.

\begin{DndMonsterLegendaryActions}
\DndMonsterLegendaryAction{Tentacle Attack or Fling}{The kraken makes one tentacle attack or uses its Fling.
}
\DndMonsterLegendaryAction{Lightning Storm (Costs 2 Actions)}{The kraken uses Lightning Storm.
}
\DndMonsterLegendaryAction{Ink Cloud (Costs 3 Actions)}{While underwater, the kraken expels an ink cloud in a 60-foot radius. The cloud spreads around corners, and that area is heavily obscured to creatures other than the kraken. Each creature other than the kraken that ends its turn there must succeed on a DC 23 Constitution saving throw, taking 16 (3d10) poison damage on a failed save, or half as much damage on a successful one. A strong current disperses the cloud, which otherwise disappears at the end of the kraken's next turn.
}
\end{DndMonsterLegendaryActions}
\DndMonsterSection{Lair Actions}
On initiative count 20 (losing initiative ties), the kraken takes a lair action to cause one of the following magical effects:


\begin{itemize}
\item A strong current moves through the kraken's lair. Each creature within 60 feet of the kraken must succeed on a DC 23 Strength saving throw or be pushed up to 60 feet away from the kraken. On a success, the creature is pushed 10 feet away from the kraken.
\item Creatures in the water within 60 feet of the kraken have vulnerability to lightning damage until initiative count 20 on the next round.
\item The water in the kraken's lair becomes electrically charged. All creatures within 120 feet of the kraken must succeed on a DC 23 Constitution saving throw, taking 10 (3d6) lightning damage on a failed save, or half as much damage on a successful one.
\end{itemize}

\DndMonsterSection{Regional Effects}
The region containing a kraken's lair is warped by the creature's blasphemous presence, creating the following magical effects:


\begin{itemize}
\item The kraken can alter the weather at will in a 6-mile radius centered on its lair. The effect is identical to the \textit{control weather} spell.
\item Water elementals coalesce within 6 miles of the lair. These elementals can't leave the water and have Intelligence and Charisma scores of 1 (-5).
\item Aquatic creatures within 6 miles of the lair that have an Intelligence score of 2 or lower are charmed by the kraken and aggressive toward intruders in the area.
\end{itemize}

When the kraken dies, all of these regional effects fade immediately.

\end{multicols}
\end{DndMonster}



\begin{DndMonster}[float=t]{Kuo-toa Archpriest}
\DndMonsterType{Medium humanoid (kuo-toa), neutral evil}
\DndMonsterBasics[
armor-class = {13 (natural armor)},
hit-points = {\DndDice{13d8 + 39}},
speed = {30 ft., swim 30 ft.}
]
\DndMonsterAbilityScores[
 str = 16,
 dex = 14,
 con = 16,
 int = 13,
 wis = 16,
 cha = 14
]
\DndMonsterDetails[
skills = {Perception +9, Religion +7},
senses = {darkvision 120 ft., passive perception 19},
languages = {Undercommon},
challenge = 6,
]
\DndMonsterAction{Amphibious}
The kuo-toa can breathe air and water.

\DndMonsterAction{Otherworldly Perception}
The kuo-toa can sense the presence of any creature within 30 feet of it that is invisible or on the Ethereal Plane. It can pinpoint such a creature that is moving.

\DndMonsterAction{Slippery}
The kuo-toa has advantage on ability checks and saving throws made to escape a grapple.

\DndMonsterAction{Sunlight Sensitivity}
While in sunlight, the kuo-toa has disadvantage on attack rolls, as well as on Wisdom (Perception) checks that rely on sight.

\DndMonsterAction{Spellcasting}
The kuo-toa is a 10th-level spellcaster. Its spellcasting ability is Wisdom (spell save DC 14, 6 to hit with spell attacks). The kuo-toa has the following cleric spells prepared:
\begin{DndMonsterSpells}
  \DndMonsterSpellLevel{\textit{guidance}, \textit{sacred flame}, \textit{thaumaturgy}}
   \DndMonsterSpellLevel[1][4]{\textit{detect magic}, \textit{sanctuary}, \textit{shield of faith}}
   \DndMonsterSpellLevel[2][3]{\textit{hold person}, \textit{spiritual weapon}}
   \DndMonsterSpellLevel[3][3]{\textit{spirit guardians}, \textit{tongues}}
   \DndMonsterSpellLevel[4][3]{\textit{control water}, \textit{divination}}
   \DndMonsterSpellLevel[5][2]{\textit{mass cure wounds}, \textit{scrying}}
\end{DndMonsterSpells}
\DndMonsterSection{Actions}
\DndMonsterAction{Multiattack}
The kuo-toa makes two melee attacks.

\DndMonsterAction{Scepter}
\textsl{Melee Weapon Attack:} 6 to hit, reach 5 ft., one target. \textsl{Hit:} 6 (1d6 + 3) bludgeoning damage plus 14 (4d6) lightning damage.

\DndMonsterAction{Unarmed Strike}
\textsl{Melee Weapon Attack:} 6 to hit, reach 5 ft., one target. \textsl{Hit:} 5 (1d4 + 3) bludgeoning damage.

\end{DndMonster}



\begin{DndMonster}[float=t]{Kuo-toa Monitor}
\DndMonsterType{Medium humanoid (kuo-toa), neutral evil}
\DndMonsterBasics[
armor-class = {13 (natural armor)},
hit-points = {\DndDice{10d8 + 20}},
speed = {30 ft., swim 30 ft.}
]
\DndMonsterAbilityScores[
 str = 14,
 dex = 10,
 con = 14,
 int = 12,
 wis = 14,
 cha = 11
]
\DndMonsterDetails[
skills = {Perception +6, Religion +5},
senses = {darkvision 120 ft., passive perception 16},
languages = {Undercommon},
challenge = 3,
]
\DndMonsterAction{Amphibious}
The kuo-toa can breathe air and water.

\DndMonsterAction{Otherworldly Perception}
The kuo-toa can sense the presence of any creature within 30 feet of it that is invisible or on the Ethereal Plane. It can pinpoint such a creature that is moving.

\DndMonsterAction{Slippery}
The kuo-toa has advantage on ability checks and saving throws made to escape a grapple.

\DndMonsterAction{Sunlight Sensitivity}
While in sunlight, the kuo-toa has disadvantage on attack rolls, as well as on Wisdom (Perception) checks that rely on sight.

\DndMonsterAction{Unarmored Defense}
The kuo-toa adds its Wisdom modifier to its armor class.

\DndMonsterSection{Actions}
\DndMonsterAction{Multiattack}
The kuo-toa makes one bite attack and two unarmed strikes.

\DndMonsterAction{Bite}
\textsl{Melee Weapon Attack:} 6 to hit, reach 5 ft., one target. \textsl{Hit:} 4 (1d4 + 2) piercing damage.

\DndMonsterAction{Unarmed Strike}
\textsl{Melee Weapon Attack:} 6 to hit, reach 5 ft., one target. \textsl{Hit:} 5 (1d6 + 2) bludgeoning damage plus 3 (1d6) lightning damage, and the target can't take reactions until the end of the kuo-toa's next turn.

\end{DndMonster}



\begin{DndMonster}[float=t]{Larva}
\DndMonsterType{Medium fiend, neutral evil}
\DndMonsterBasics[
armor-class = {9},
hit-points = {\DndDice{2d8}},
speed = {20 ft.}
]
\DndMonsterAbilityScores[
 str = 9,
 dex = 9,
 con = 10,
 int = 6,
 wis = 10,
 cha = 2
]
\DndMonsterDetails[
languages = {understands the languages it knew in life but can't speak},
challenge = 0,
]
\DndMonsterSection{Actions}
\DndMonsterAction{Bite}
\textsl{Melee Weapon Attack:} 1 to hit, reach 5 ft., one target. \textsl{Hit:} 1 (1d4 - 1) piercing damage.

\end{DndMonster}



\begin{DndMonster}[float=t]{Lemure}
\DndMonsterType{Medium fiend (devil), lawful evil}
\DndMonsterBasics[
armor-class = {7},
hit-points = {\DndDice{3d8}},
speed = {15 ft.}
]
\DndMonsterAbilityScores[
 str = 10,
 dex = 5,
 con = 11,
 int = 1,
 wis = 11,
 cha = 3
]
\DndMonsterDetails[
damage-resistances = {cold},
damage-immunities = {fire, poison},
condition-immunities = {charmed, frightened, poisoned},
senses = {darkvision 120 ft., passive perception 10},
languages = {understands Infernal but can't speak},
challenge = 0,
]
\DndMonsterAction{Devil's Sight}
Magical darkness doesn't impede the lemure's darkvision.

\DndMonsterAction{Hellish Rejuvenation}
A lemure that dies in the Nine Hells comes back to life with all its hit points in 1d10 days unless it is killed by a good-aligned creature with a \textit{bless} spell cast on that creature or its remains are sprinkled with holy water.

\DndMonsterSection{Actions}
\DndMonsterAction{Fist}
\textsl{Melee Weapon Attack:} 3 to hit, reach 5 ft., one target. \textsl{Hit:} 2 (1d4) bludgeoning damage

\end{DndMonster}



\begin{DndMonster}[float=t]{Lesser Tyrant Shadow}
\DndMonsterType{Medium aberration, lawful evil}
\DndMonsterBasics[
armor-class = {20 (natural armor)},
hit-points = {\DndDice{12d8 + 60}},
speed = {30 ft.}
]
\DndMonsterAbilityScores[
 str = 20,
 dex = 18,
 con = 20,
 int = 18,
 wis = 16,
 cha = 18
]
\DndMonsterDetails[
saving-throws = {Int +8, Wis +7},
skills = {Arcana +8, Deception +12, Perception +7, Stealth +12},
damage-resistances = {fire},
damage-immunities = {necrotic; poison; psychic; bludgeoning, piercing, slashing from nonmagical attacks},
condition-immunities = {blinded, charmed, frightened, poisoned},
languages = {understands all, telepathy 300 ft.},
challenge = 10,
]
\DndMonsterAction{Empath}
The shadow uses telepathy to sense nearby creatures. The shadow can see a creature if that creature has thoughts and is within the range of the shadow's telepathy. The shadow can't see creatures immune to telepathy.

\DndMonsterAction{Undying Connection}
The shadow was originally spawned from a specific devil and, if the shadow is slain, it returns to life in 10 (1d20) days within 1 mile of that devil. The shadow is killed permanently only if its progenitor devil is slain first.

\DndMonsterSection{Actions}
\DndMonsterAction{Multiattack}
The shadow makes two Claw attacks.

\DndMonsterAction{Claw}
\textsl{Melee Weapon Attack:} 9 to hit, reach 5 ft., one target. \textsl{Hit:} 12 (2d6 + 5) necrotic damage, plus 3 (1d6) psychic damage.

\DndMonsterAction{Cloak of Shadows}
The shadow bends darkness around itself, and now has the invisible condition for 1 minute. The invisibility ends if the shadow is subjected to a \textit{Daylight} spell or similar.

\DndMonsterAction{Empathic Link}
The shadow attempts to form a link between itself and a creature it can see. The target must make a DC 16 Intelligence saving throw, taking 6 (1d12) psychic damage on a success. On a failed save, the creature and the shadow are linked. While linked, every time the shadow is damaged, the creature takes half of that damage as psychic damage. The creature may repeat the saving throw at the end of each of its turns, taking 6 (1d12) psychic damage on a failed save, or ending the effect on a successful one.

\DndMonsterSection{Bonus Actions}
\DndMonsterAction{Change Shape}
The shadow transforms into a Beast, Humanoid or Fiend that it has an Empathic Link with, or back into its true form. In a new form, the shadow retains its align ment, hit points, hit dice, telepathy, and Intelligence, Wisdom, and Charisma scores, as well as this action. Its statistics and capabilities are otherwise replaced by those of the new form, except any class features, lair actions, or legendary actions of that form.

\end{DndMonster}



\begin{DndMonster}[float*=b,width=\textwidth + 8pt]{Levistus}
\begin{multicols}{2}
\DndMonsterType{Medium fiend (devil), lawful evil}
\DndMonsterBasics[
armor-class = {23 (natural armor)},
hit-points = {\DndDice{32d8 + 192}},
speed = {0 ft.}
]
\DndMonsterAbilityScores[
 str = 19,
 dex = 26,
 con = 22,
 int = 25,
 wis = 28,
 cha = 26
]
\DndMonsterDetails[
saving-throws = {Dex +16, Con +14, Int +15, Wis +17},
skills = {Acrobatics +24, Deception +16, Intimidation +16, Perception +25, Persuasion +16, Stealth +16},
damage-resistances = {bludgeoning, piercing, slashing from nonmagical attacks that aren't silvered},
damage-immunities = {cold, fire, poison},
condition-immunities = {charmed, poisoned},
senses = {darkvision 120 ft., passive perception 35},
languages = {Celestial, Common, Draconic, Infernal, telepathy 1,000 ft.},
challenge = 26,
]
\DndMonsterAction{Armor of Ice}
Levistus is always entombed in ice. The ice provides 100 temporary hit points. While Levistus has these temporary hit points, he has vulnerability to fire damage. If Levistus starts his turn with no temporary hit points from the ice, he regains 100 temporary hit points.

\DndMonsterAction{Devil's Sight}
Magical darkness doesn't impede Levistus's darkvision.

\DndMonsterAction{Legendary Resistance (3/Day)}
If Levistus fails a saving throw, he can choose to succeed instead.

\DndMonsterAction{Magic Resistance}
Levistus has advantage on saving throws against spells and other magical effects.

\DndMonsterAction{Sub-zero}
Cold damage dealt by Levistus ignores resistances and immunities.

\DndMonsterAction{Spellcasting}
Levistus casts one of the following spells, requiring no material components and using Intelligence as the spellcasting ability (spell save DC 23):
\begin{DndMonsterSpells}
  \DndInnateSpellLevel{\textit{Cone of Cold}, \textit{Darkness}, \textit{Ice Storm}, \textit{Wall of Ice}}
\end{DndMonsterSpells}
\DndMonsterSection{Actions}
\DndMonsterAction{Multiattack}
Levistus makes three Touch of Stygia or Ice Bolt attacks. He can replace one of the attacks with Blizzard (if available).

\DndMonsterAction{Touch of Stygia}
\textsl{Melee Spell Attack:} 15 to hit, reach 5 ft., one target. \textsl{Hit:} 20 (3d8 + 7) psychic damage, and the target must succeed on a DC 23 Intelligence saving throw or have the stunned condition for 1 minute. The target may repeat the saving throw at the end of its turns, ending the effect on a success. A creature that succeeds on the saving throw is immune to this effect for 1 hour.

\DndMonsterAction{Ice Bolt}
\textsl{Ranged Spell Attack:} 15 to hit, range 120 ft., one target. \textsl{Hit:} 23 (3d10 + 7) cold damage.

\DndMonsterAction{Blizzard (Recharge 4\textendash 6)}
Levistus chooses a point he can see within 300 feet of him. Each creature in a 20-foot-radius, 40-foot-high cylinder centered on that point must make a DC 23 Dexterity saving throw. Targets take 13 (3d8) force damage plus 13 (3d8) cold damage on a failed save, or half as much damage on a successful one.

\DndMonsterSection{Reactions}
\DndMonsterAction{Counterspell}
As a reaction to a creature he can see casting a spell, Levistus can attempt to prevent the casting. Levistus makes an Intelligence check (+7) against a DC of (10 + spell level). On a success, the creature's spell fails and has no effect.

\DndMonsterSection{Legendary Actions}
Levistus can take 3 legendary actions, choosing from the options below. Only one legendary action can be used at a time and only at the end of another creature's turn. Levistus regains spent legendary actions at the start of its turn.

\begin{DndMonsterLegendaryActions}
\DndMonsterLegendaryAction{Amnesia}{Levistus selects a stunned creature that he can see. The creature must succeed on a DC 23 Charisma saving throw or Levistus erases a specific memory from the target. The memory must be one that the target experienced in the last 24 hours and that lasted no more than 10 minutes. The memory can only be restored with a Remove Curse or a \textit{Greater Restoration} spell.
}
\DndMonsterLegendaryAction{Froststrike}{Levistus makes an Ice Bolt attack. On a hit, the target falls and has the prone condition.
}
\DndMonsterLegendaryAction{Call Underling (Costs 3 Actions)}{Levistus summons an allied \textbf{ice devil} in an unoccupied space that he can see.
}
\end{DndMonsterLegendaryActions}
\DndMonsterSection{Regional Effects}
The icy regions containing Levistus' lair is corrupted by his presence, which creates the following effect:


\begin{description}
\item[Frozen Landscape.]
Any creature that has the prone or unconscious conditions within 1 mile of the lair, takes 21 (6d6) cold damage at the start of their turn.

\end{description}

If Levistus dies, the effects fade over the course of 1d10 days.

\end{multicols}
\end{DndMonster}



\begin{DndMonster}[float*=b,width=\textwidth + 8pt]{Lich}
\begin{multicols}{2}
\DndMonsterType{Medium undead, any evil alignment}
\DndMonsterBasics[
armor-class = {17 (natural armor)},
hit-points = {\DndDice{18d8 + 54}},
speed = {30 ft.}
]
\DndMonsterAbilityScores[
 str = 11,
 dex = 16,
 con = 16,
 int = 20,
 wis = 14,
 cha = 16
]
\DndMonsterDetails[
saving-throws = {Con +10, Int +12, Wis +9},
skills = {Arcana +19, History +12, Insight +9, Perception +9},
damage-resistances = {cold, lightning, necrotic},
damage-immunities = {poison; bludgeoning, piercing, slashing from nonmagical attacks},
condition-immunities = {charmed, exhaustion, frightened, paralyzed, poisoned},
senses = {truesight 120 ft., passive perception 19},
languages = {Common plus up to five other languages},
challenge = 21,
]
\DndMonsterAction{Legendary Resistance (3/Day)}
If the lich fails a saving throw, it can choose to succeed instead.

\DndMonsterAction{Rejuvenation}
If it has a phylactery, a destroyed lich gains a new body in 1d10 days, regaining all its hit points and becoming active again. The new body appears within 5 feet of the phylactery.

\DndMonsterAction{Turn Resistance}
The lich has advantage on saving throws against any effect that turns undead.

\DndMonsterAction{Spellcasting}
The lich is an 18th-level spellcaster. Its spellcasting ability is Intelligence (spell save DC 20, 12 to hit with spell attacks). The lich has the following wizard spells prepared:
\begin{DndMonsterSpells}
  \DndMonsterSpellLevel{\textit{mage hand}, \textit{prestidigitation}, \textit{ray of frost}}
   \DndMonsterSpellLevel[1][4]{\textit{detect magic}, \textit{magic missile}, \textit{shield}, \textit{thunderwave}}
   \DndMonsterSpellLevel[2][3]{\textit{detect thoughts}, \textit{invisibility}, \textit{Melf's acid arrow}, \textit{mirror image}}
   \DndMonsterSpellLevel[3][3]{\textit{animate dead}, \textit{counterspell}, \textit{dispel magic}, \textit{fireball}}
   \DndMonsterSpellLevel[4][3]{\textit{blight}, \textit{dimension door}}
   \DndMonsterSpellLevel[5][3]{\textit{cloudkill}, \textit{scrying}}
   \DndMonsterSpellLevel[6][1]{\textit{disintegrate}, \textit{globe of invulnerability}}
   \DndMonsterSpellLevel[7][1]{\textit{finger of death}, \textit{plane shift}}
   \DndMonsterSpellLevel[8][1]{\textit{dominate monster}, \textit{power word stun}}
\end{DndMonsterSpells}
\DndMonsterSection{Actions}
\DndMonsterAction{Paralyzing Touch}
\textsl{Melee Spell Attack:} 12 to hit, reach 5 ft., one creature. \textsl{Hit:} 10 (3d6) cold damage. The target must succeed on a DC 18 Constitution saving throw or be paralyzed for 1 minute. The target can repeat the saving throw at the end of each of its turns, ending the effect on itself on a success.

\DndMonsterSection{Legendary Actions}
Lich can take 3 legendary actions, choosing from the options below. Only one legendary action can be used at a time and only at the end of another creature's turn. Lich regains spent legendary actions at the start of its turn.

\begin{DndMonsterLegendaryActions}
\DndMonsterLegendaryAction{Cantrip}{The lich casts a cantrip.
}
\DndMonsterLegendaryAction{Paralyzing Touch (Costs 2 Actions)}{The lich uses its Paralyzing Touch.
}
\DndMonsterLegendaryAction{Frightening Gaze (Costs 2 Actions)}{The lich fixes its gaze on one creature it can see within 10 feet of it. The target must succeed on a DC 18 Wisdom saving throw against this magic or become frightened for 1 minute. The frightened target can repeat the saving throw at the end of each of its turns, ending the effect on itself on a success. If a target's saving throw is successful or the effect ends for it, the target is immune to the lich's gaze for the next 24 hours.
}
\DndMonsterLegendaryAction{Disrupt Life (Costs 3 Actions)}{Each non-undead creature within 20 feet of the lich must make a DC 18 Constitution saving throw against this magic, taking 21 (6d6) necrotic damage on a failed save, or half as much damage on a successful one.
}
\end{DndMonsterLegendaryActions}
\DndMonsterSection{Lair Actions}
On initiative count 20 (losing initiative ties), the lich can take a lair action to cause one of the following magical effects; the lich can't use the same effect two rounds in a row:


\begin{itemize}
\item The lich rolls a d8 and regains a spell slot of that level or lower. If it has no spent spell slots of that level or lower, nothing happens.
\item The lich targets one creature it can see within 30 feet of it. A crackling cord of negative energy tethers the lich to the target. Whenever the lich takes damage, the target must make a DC 18 Constitution saving throw. On a failed save, the lich takes half the damage (rounded down), and the target takes the remaining damage. This tether lasts until initiative count 20 on the next round or until the lich or the target is no longer in the lich's lair.
\item The lich calls forth the spirits of creatures that died in its lair. These apparitions materialize and attack one creature that the lich can see within 60 feet of it. The target must succeed on a DC 18 Constitution saving throw, taking 52 (15d6) necrotic damage on a failed save, or half as much damage on a success. The apparitions then disappear.
\end{itemize}

\end{multicols}
\end{DndMonster}



\begin{DndMonster}[float=t]{Lion}
\DndMonsterType{Large beast, unaligned}
\DndMonsterBasics[
armor-class = {12},
hit-points = {\DndDice{4d10 + 4}},
speed = {50 ft.}
]
\DndMonsterAbilityScores[
 str = 17,
 dex = 15,
 con = 13,
 int = 3,
 wis = 12,
 cha = 8
]
\DndMonsterDetails[
skills = {Perception +3, Stealth +6},
challenge = 1,
]
\DndMonsterAction{Keen Smell}
The lion has advantage on Wisdom (Perception) checks that rely on smell.

\DndMonsterAction{Pack Tactics}
The lion has advantage on an attack roll against a creature if at least one of the lion's allies is within 5 feet of the creature and the ally isn't incapacitated.

\DndMonsterAction{Pounce}
If the lion moves at least 20 feet straight toward a creature and then hits it with a claw attack on the same turn, that target must succeed on a DC 13 Strength saving throw or be knocked prone. If the target is prone, the lion can make one bite attack against it as a bonus action.

\DndMonsterAction{Running Leap}
With a 10-foot running start, the lion can long jump up to 25 feet.

\DndMonsterSection{Actions}
\DndMonsterAction{Bite}
\textsl{Melee Weapon Attack:} 5 to hit, reach 5 ft., one target. \textsl{Hit:} 7 (1d8 + 3) piercing damage.

\DndMonsterAction{Claw}
\textsl{Melee Weapon Attack:} 5 to hit, reach 5 ft., one target. \textsl{Hit:} 6 (1d6 + 3) slashing damage.

\end{DndMonster}



\begin{DndMonster}[float=t]{Maelephant Nomad}
\DndMonsterType{Large fiend, lawful evil}
\DndMonsterBasics[
armor-class = {16 (\textit{chain mail})},
hit-points = {\DndDice{20d10 + 80}},
speed = {40 ft.}
]
\DndMonsterAbilityScores[
 str = 22,
 dex = 11,
 con = 18,
 int = 14,
 wis = 17,
 cha = 15
]
\DndMonsterDetails[
saving-throws = {Str +11, Con +9},
skills = {Athletics +16, Insight +8, Intimidation +7, Perception +8},
damage-resistances = {acid; fire; lightning; bludgeoning, piercing, slashing from nonmagical attacks},
condition-immunities = {frightened, poisoned},
senses = {blindsight 30 ft., darkvision 120 ft., passive perception 18},
languages = {Common, Infernal, Maelephant},
challenge = 14,
]
\DndMonsterAction{Magic Resistance}
The maelephant has advantage on saving throws against spells and other magical effects.

\DndMonsterAction{Regenerative}
At the start of its turn, if the maelephant has less than half its hit points, it regenerates 20 hit points. This trait doesn't function if the maelephant starts its turn with less than 1 hit point.

\DndMonsterSection{Actions}
\DndMonsterAction{Multiattack}
The maelephant makes two attacks using its Claw, Trunk, or a combination of the two.

\DndMonsterAction{Claw}
\textsl{Melee Weapon Attack:} 11 to hit, reach 5 ft., one target. \textsl{Hit:} 17 (2d10 + 6) slashing damage.

\DndMonsterAction{Trunk}
\textsl{Melee Weapon Attack:} 11 to hit, reach 10 ft., one target. \textsl{Hit:} 16 (3d6 + 6) piercing damage. If the target is a Medium or smaller creature, it has the grappled condition (escape DC 18). Until this grapple ends, the target has the restrained condition. While it is grappling a creature, the maelephant can't use its Trunk attack against other creatures.

\DndMonsterAction{Noxious Breath (Recharge 5\textendash 6)}
The maelephant releases a cloud of noxious gas. All creatures in a 30-foot cone in front of the maelephant must make a DC 16 Constitution saving throw. Targets take 42 (12d6) acid damage on a failed save, or half as much damage on a successful one. In addition, on a failed save, the creature's Intelligence and Charisma scores become 1. The creature can't cast spells, activate magic items, understand language, or communicate in any intelligible way. The creature can, however, identify its friends, follow them, and even protect them. A feebled creature can repeat the save every minute, ending the effect on a success.

\DndMonsterSection{Bonus Actions}
\DndMonsterAction{Reckless Charge}
The maelephant moves up to 40 feet towards a creature it can see and if it moves at least 20 feet in a straight line, it may make a Claw attack. On a hit, the target takes an additional 14 (4d6) force damage and has the prone condition.

\DndMonsterSection{Reactions}
\DndMonsterAction{Defensive Stance}
As a reaction to taking damage, the maelephant enters a defensive stance until the end of its next turn. While in a defensive stance, the maelephant moves at half speed and can't take bonus actions, but its AC increases by 5.

\end{DndMonster}



\begin{DndMonster}[float=t]{Mage}
\DndMonsterType{Medium humanoid (any race), any alignment}
\DndMonsterBasics[
armor-class = {12 (15 with \textit{mage armor})},
hit-points = {\DndDice{9d8}},
speed = {30 ft.}
]
\DndMonsterAbilityScores[
 str = 9,
 dex = 14,
 con = 11,
 int = 17,
 wis = 12,
 cha = 11
]
\DndMonsterDetails[
saving-throws = {Int +6, Wis +4},
skills = {Arcana +6, History +6},
languages = {any four languages},
challenge = 6,
]
\DndMonsterAction{Spellcasting}
The mage is a 9th-level spellcaster. Its spellcasting ability is Intelligence (spell save DC 14, 6 to hit with spell attacks). The mage has the following wizard spells prepared:
\begin{DndMonsterSpells}
  \DndMonsterSpellLevel{\textit{fire bolt}, \textit{light}, \textit{mage hand}, \textit{prestidigitation}}
   \DndMonsterSpellLevel[1][4]{\textit{detect magic}, \textit{mage armor}, \textit{magic missile}, \textit{shield}}
   \DndMonsterSpellLevel[2][3]{\textit{misty step}, \textit{suggestion}}
   \DndMonsterSpellLevel[3][3]{\textit{counterspell}, \textit{fireball}, \textit{fly}}
   \DndMonsterSpellLevel[4][3]{\textit{greater invisibility}, \textit{ice storm}}
   \DndMonsterSpellLevel[5][1]{\textit{cone of cold}}
\end{DndMonsterSpells}
\DndMonsterSection{Actions}
\DndMonsterAction{Dagger}
\textsl{Melee or Ranged Weapon Attack:} 5 to hit, reach 5 ft. or range 20/60 ft., one target. \textsl{Hit:} 4 (1d4 + 2) piercing damage.

\end{DndMonster}



\begin{DndMonster}[float*=b,width=\textwidth + 8pt]{Mammon}
\begin{multicols}{2}
\DndMonsterType{Huge fiend (devil), lawful evil}
\DndMonsterBasics[
armor-class = {18 (natural armor)},
hit-points = {\DndDice{32d12 + 256}},
speed = {40 ft.}
]
\DndMonsterAbilityScores[
 str = 28,
 dex = 17,
 con = 26,
 int = 22,
 wis = 25,
 cha = 25
]
\DndMonsterDetails[
saving-throws = {Dex +11, Con +16, Wis +15, Cha +15},
skills = {Arcana +14, Deception +15, Insight +15, Intimidation +23, Investigation +14, Persuasion +23},
damage-resistances = {acid; cold; bludgeoning, piercing, slashing from nonmagical attacks that aren't silvered},
damage-immunities = {fire, poison},
condition-immunities = {charmed, poisoned},
senses = {darkvision 120 ft., truesight 120 ft., passive perception 17},
languages = {Celestial, Common, Draconic, Infernal, telepathy 120 ft.},
challenge = 26,
]
\DndMonsterAction{Devil's Sight}
Magical darkness doesn't impede Mammon's darkvision.

\DndMonsterAction{Fiendish Regeneration}
Mammon regains 20 hit points at the start of his turn. If he takes radiant damage this trait doesn't function at the start of his next turn. Mammon only dies if he starts his turn with 0 hit points and is unable to regenerate.

\DndMonsterAction{Legendary Resistance (3/Day)}
If Mammon fails a saving throw, he can choose to succeed instead

\DndMonsterAction{Magic Resistance}
Mammon has advantage on saving throws against spells and other magical effects.

\DndMonsterSection{Actions}
\DndMonsterAction{Multiattack}
Mammon uses Fearful Gaze, then makes four attacks using Tail Swipe, Constrict, or a combination of the two.

\DndMonsterAction{Tail Swipe}
\textsl{Melee Weapon Attack:} 17 to hit, reach 10 ft., one target. \textsl{Hit:} 31 (4d10 + 9) bludgeoning damage. If the target is a Large or smaller creature, it has the grappled condition (escape DC 21). While grappled, the creature also has the restrained condition, and Mammon can't make Tail Swipe attacks.

\DndMonsterAction{Fearful Gaze}
Mammon focuses his gaze in a 90-foot cone in front of him. Each creature of his choice within the cone must make a DC 23 Wisdom saving throw, taking 19 (3d12) psychic damage on a failed save, or half as much damage on a successful one. Creatures that fail the save have the frightened condition until the end of their next turn.

\DndMonsterAction{Constrict}
Mammon tightens his grip around a creature he has grappled. The target takes 35 (4d12 + 9) bludgeoning damage.

\DndMonsterAction{Golden Gaze (Recharge 6)}
Mammon attempts to turn one creature he can see within 60 feet of him into solid gold. The target must make a DC 23 Constitution saving throw. On a failed save, the target's flesh begins to harden, and it has the restrained condition. A restrained creature must repeat the saving throw at the end of each turn. If it successfully saves three times, the effect ends. If it fails three times, it turns to gold and now has the petrified condition. Successes and failures don't need to be consecutive.

\DndMonsterSection{Reactions}
\DndMonsterAction{Pain Tax}
As a reaction to taking damage, Mammon causes the attacker's valuables to glow red hot. This deals 3 (1d6) fire damage to the attacker for every 100 gold (or equivalent in gems, jewelry, and other coinage) that it currently carries.

\DndMonsterSection{Legendary Actions}
Mammon can take 3 legendary actions, choosing from the options below. Only one legendary action can be used at a time and only at the end of another creature's turn. Mammon regains spent legendary actions at the start of its turn.

\begin{DndMonsterLegendaryActions}
\DndMonsterLegendaryAction{Change Form}{Mammon changes his form into that of a pit fiend, maintaining his statistics but losing his actions. He gains the actions and mobility of a pit fiend, including its weapon, and maintains his legendary actions and lair actions. Mammon's serpentine abilities can still recharge while he is in pit fiend form. If Mammon is already a pit fiend, this action changes him back into his serpentine form.
}
\DndMonsterLegendaryAction{Golden Touch (Costs 2 Actions)}{Mammon recharges Golden Gaze.
}
\DndMonsterLegendaryAction{Call Underling (Costs 3 Actions)}{Mammon summons an allied \textbf{horned devil} in an unoccupied space that he can see.
}
\end{DndMonsterLegendaryActions}
\DndMonsterSection{Regional Effects}
The region containing Mammon's lair is corrupted by his presence, which creates the following effect:


\begin{description}
\item[Rotting Vegetation.]
Within 3 miles of the lair, the plant life is even more sickly and decayed, and the water is polluted and toxic. A creature that isn't a Fiend that eats or drinks from these sources must succeed on a DC 19 Constitution saving throw or die.

\end{description}

If Mammon dies, the effects fade over the course of 1d10 days.

\end{multicols}
\end{DndMonster}



\begin{DndMonster}[float=t]{Mammoth}
\DndMonsterType{Huge beast, unaligned}
\DndMonsterBasics[
armor-class = {13 (natural armor)},
hit-points = {\DndDice{11d12 + 55}},
speed = {40 ft.}
]
\DndMonsterAbilityScores[
 str = 24,
 dex = 9,
 con = 21,
 int = 3,
 wis = 11,
 cha = 6
]
\DndMonsterDetails[
challenge = 6,
]
\DndMonsterAction{Trampling Charge}
If the mammoth moves at least 20 feet straight toward a creature and then hits it with a gore attack on the same turn, that target must succeed on a DC 18 Strength saving throw or be knocked prone. If the target is prone, the mammoth can make one stomp attack against it as a bonus action.

\DndMonsterSection{Actions}
\DndMonsterAction{Gore}
\textsl{Melee Weapon Attack:} 10 to hit, reach 10 ft., one target. \textsl{Hit:} 25 (4d8 + 7) piercing damage.

\DndMonsterAction{Stomp}
\textsl{Melee Weapon Attack:} 10 to hit, reach 5 ft., one prone creature. \textsl{Hit:} 29 (4d10 + 7) bludgeoning damage.

\end{DndMonster}



\begin{DndMonster}[float=t]{Manes}
\DndMonsterType{Small fiend (demon), chaotic evil}
\DndMonsterBasics[
armor-class = {9 (natural armor)},
hit-points = {\DndDice{2d6 + 2}},
speed = {20 ft.}
]
\DndMonsterAbilityScores[
 str = 10,
 dex = 9,
 con = 13,
 int = 3,
 wis = 8,
 cha = 4
]
\DndMonsterDetails[
damage-resistances = {cold, fire, lightning},
damage-immunities = {poison},
condition-immunities = {charmed, frightened, poisoned},
senses = {darkvision 60 ft., passive perception 9},
languages = {understands Abyssal but can't speak},
challenge = 1/8,
]
\DndMonsterSection{Actions}
\DndMonsterAction{Claws}
\textsl{Melee Weapon Attack:} 2 to hit, reach 5 ft., one target. \textsl{Hit:} 5 (2d4) slashing damage.

\end{DndMonster}



\begin{DndMonster}[float=t]{Marilith}
\DndMonsterType{Large fiend (demon), chaotic evil}
\DndMonsterBasics[
armor-class = {18 (natural armor)},
hit-points = {\DndDice{18d10 + 90}},
speed = {40 ft.}
]
\DndMonsterAbilityScores[
 str = 18,
 dex = 20,
 con = 20,
 int = 18,
 wis = 16,
 cha = 20
]
\DndMonsterDetails[
saving-throws = {Str +9, Con +10, Wis +8, Cha +10},
damage-resistances = {cold; fire; lightning; bludgeoning, piercing, slashing from nonmagical attacks},
damage-immunities = {poison},
condition-immunities = {poisoned},
senses = {truesight 120 ft., passive perception 13},
languages = {Abyssal, telepathy 120 ft.},
challenge = 16,
]
\DndMonsterAction{Magic Resistance}
The marilith has advantage on saving throws against spells and other magical effects.

\DndMonsterAction{Magic Weapons}
The marilith's weapon attacks are magical.

\DndMonsterAction{Reactive}
The marilith can take one reaction on every turn in combat.

\DndMonsterSection{Actions}
\DndMonsterAction{Multiattack}
The marilith can make seven attacks: six with its longswords and one with its tail.

\DndMonsterAction{Longsword}
\textsl{Melee Weapon Attack:} 9 to hit, reach 5 ft., one target. \textsl{Hit:} 13 (2d8 + 4) slashing damage.

\DndMonsterAction{Tail}
\textsl{Melee Weapon Attack:} 9 to hit, reach 10 ft., one creature. \textsl{Hit:} 15 (2d10 + 4) bludgeoning damage. If the target is Medium or smaller, it is grappled (escape DC 19). Until this grapple ends, the target is restrained, the marilith can automatically hit the target with its tail, and the marilith can't make tail attacks against other targets.

\DndMonsterAction{Teleport}
The marilith magically teleports, along with any equipment it is wearing or carrying, up to 120 feet to an unoccupied space it can see.

\DndMonsterSection{Reactions}
\DndMonsterAction{Parry}
The marilith adds 5 to its AC against one melee attack that would hit it. To do so, the marilith must see the attacker and be wielding a melee weapon.

\end{DndMonster}



\begin{DndMonster}[float=t]{Mastiff}
\DndMonsterType{Medium beast, unaligned}
\DndMonsterBasics[
armor-class = {12},
hit-points = {\DndDice{1d8 + 1}},
speed = {40 ft.}
]
\DndMonsterAbilityScores[
 str = 13,
 dex = 14,
 con = 12,
 int = 3,
 wis = 12,
 cha = 7
]
\DndMonsterDetails[
skills = {Perception +3},
challenge = 1/8,
]
\DndMonsterAction{Keen Hearing and Smell}
The mastiff has advantage on Wisdom (Perception) checks that rely on hearing or smell.

\DndMonsterSection{Actions}
\DndMonsterAction{Bite}
\textsl{Melee Weapon Attack:} 3 to hit, reach 5 ft., one target. \textsl{Hit:} 4 (1d6 + 1) piercing damage. If the target is a creature, it must succeed on a DC 11 Strength saving throw or be knocked prone.

\end{DndMonster}



\begin{DndMonster}[float=t]{Medusa}
\DndMonsterType{Medium monstrosity, lawful evil}
\DndMonsterBasics[
armor-class = {15 (natural armor)},
hit-points = {\DndDice{17d8 + 51}},
speed = {30 ft.}
]
\DndMonsterAbilityScores[
 str = 10,
 dex = 15,
 con = 16,
 int = 12,
 wis = 13,
 cha = 15
]
\DndMonsterDetails[
skills = {Deception +5, Insight +4, Perception +4, Stealth +5},
senses = {darkvision 60 ft., passive perception 14},
languages = {Common},
challenge = 6,
]
\DndMonsterAction{Petrifying Gaze}
When a creature that can see the medusa's eyes starts its turn within 30 feet of the medusa, the medusa can force it to make a DC 14 Constitution saving throw if the medusa isn't incapacitated and can see the creature. If the saving throw fails by 5 or more, the creature is instantly petrified. Otherwise, a creature that fails the save begins to turn to stone and is restrained. The restrained creature must repeat the saving throw at the end of its next turn, becoming petrified on a failure or ending the effect on a success. The petrification lasts until the creature is freed by the  \textit{greater restoration} spell or other magic.

Unless Surprise, a creature can avert its eyes to avoid the saving throw at the start of its turn. If the creature does so, it can't see the medusa until the start of its next turn, when it can avert its eyes again. If the creature looks at the medusa in the meantime, it must immediately make the save.

If the medusa sees itself reflected on a polished surface within 30 feet of it and in an area of bright light, the medusa is, due to its curse, affected by its own gaze.

\DndMonsterSection{Actions}
\DndMonsterAction{Multiattack}
The medusa makes either three melee attacks—one with its snake hair and two with its shortsword—or two ranged attacks with its longbow.

\DndMonsterAction{Snake Hair}
\textsl{Melee Weapon Attack:} 5 to hit, reach 5 ft., one creature. \textsl{Hit:} 4 (1d4 + 2) piercing damage plus 14 (4d6) poison damage.

\DndMonsterAction{Shortsword}
\textsl{Melee Weapon Attack:} 5 to hit, reach 5 ft., one target. \textsl{Hit:} 5 (1d6 + 2) piercing damage.

\DndMonsterAction{Longbow}
\textsl{Ranged Weapon Attack:} 5 to hit, range 150/600 ft., one target. \textsl{Hit:} 6 (1d8 + 2) piercing damage plus 7 (2d6) poison damage.

\end{DndMonster}



\begin{DndMonster}[float*=b,width=\textwidth + 8pt]{Mephistopheles}
\begin{multicols}{2}
\DndMonsterType{Large fiend (devil), lawful evil}
\DndMonsterBasics[
armor-class = {21 (natural armor)},
hit-points = {\DndDice{40d10 + 240}},
speed = {40 ft., fly 100 ft.}
]
\DndMonsterAbilityScores[
 str = 22,
 dex = 23,
 con = 22,
 int = 30,
 wis = 28,
 cha = 26
]
\DndMonsterDetails[
saving-throws = {Int +18, Wis +17, Cha +16},
skills = {Deception +24, Insight +25, Persuasion +24},
damage-resistances = {acid; cold; bludgeoning, piercing, slashing from nonmagical attacks that aren't silvered},
damage-immunities = {fire, poison},
condition-immunities = {charmed, exhaustion, frightened, poisoned},
senses = {darkvision 120 ft., passive perception 19},
languages = {all, telepathy 120 ft.},
challenge = 27,
]
\DndMonsterAction{Devil's Sight}
Magical darkness doesn't impede Mephistopheles' darkvision.

\DndMonsterAction{Fiendish Regeneration}
Mephistopheles regains 20 hit points at the start of his turn. If he takes radiant damage this trait doesn't function at the start of his next turn. Mephistopheles only dies if he starts his turn with 0 hit points and is unable to regenerate.

\DndMonsterAction{Flyby}
Mephistopheles doesn't provoke opportunity attacks when he flies out of an enemy's reach.

\DndMonsterAction{Hellfire Mastery}
Fire damage that Mephistopheles inflicts ignores resistances and immunities. Mephistopheles is immune to this effect.

\DndMonsterAction{Legendary Resistance (3/Day)}
If Mephistopheles fails a saving throw, he can choose to succeed instead.

\DndMonsterAction{Magic Resistance}
Mephistopheles has advantage on saving throws against spells and other magical effects.

\DndMonsterAction{Spellcasting}
Mephistopheles casts one of the following spells, requiring no material components and using Intelligence as the spellcasting ability (spell save DC 26):
\begin{DndMonsterSpells}
  \DndInnateSpellLevel{\textit{Darkness}, \textit{Detect Evil and Good}, \textit{Detect Magic}, \textit{Dispel Magic}, \textit{Geas} (duration 1 year), \textit{Greater Restoration}, \textit{Hallow}, \textit{Hold Monster}, \textit{Locate Creature}, \textit{Locate Object}, \textit{Major Image}, \textit{Resurrection}, \textit{Scrying}, \textit{Suggestion}, \textit{Wall of Fire}}
  \DndInnateSpellLevel[1e]{\textit{Meteor Swarm}, \textit{Symbol}, \textit{Wish}}
\end{DndMonsterSpells}
\DndMonsterSection{Actions}
\DndMonsterAction{Multiattack}
Mephistopheles makes three Ranseur attacks. He can replace one of the attacks with Hellfire Lash (if available).

\DndMonsterAction{Ranseur +3}
\textsl{Melee Weapon Attack:} 17 to hit, reach 10 ft., one target. \textsl{Hit:} 20 (2d10 + 9) piercing damage plus 18 (4d8) fire or cold damage (Mephistopheles' choice).

\DndMonsterAction{Hellfire Implosion (Recharge 4\textendash 6)}
A bright streak of Hellfire appears at a point that Mephistopheles can see within 150 feet of him. Each creature in a 20-foot-radius sphere centered on that point must make a DC 22 Dexterity saving throw. Targets take 31 (9d6) fire damage on a failed save, or half as much damage on a successful one.

\DndMonsterAction{Hellfire Lash (Recharge 5\textendash 6)}
Mephistopheles unleashes a 60-foot-long, 5-foot-wide lash of Hellfire that ignites a 5-foot- radius sphere around where it strikes. Any targets in the area of effect must make a DC 22 Dexterity saving throw, taking 36 (8d8) fire damage on a failed save, or half as much on a successful one.

\DndMonsterSection{Bonus Actions}
\DndMonsterAction{Ashen Teleport}
Mephistopheles' body and any equipment he is wearing or carrying turns to ash, then he magically teleports up to 120 feet to an unoccupied space he can see and reforms.

\DndMonsterSection{Legendary Actions}
Mephistopheles can take 3 legendary actions, choosing from the options below. Only one legendary action can be used at a time and only at the end of another creature's turn. Mephistopheles regains spent legendary actions at the start of its turn.

\begin{DndMonsterLegendaryActions}
\DndMonsterLegendaryAction{Hellfire Wings}{Mephistopheles uses his wings to generate a burst of heat. Each creature within 10 feet of him must succeed on a DC 22 Dexterity saving throw or take 9 (2d8) bludgeoning damage plus 9 (2d8) fire damage. Creatures that failed the save also have the prone condition.
}
\DndMonsterLegendaryAction{Hellfire Storm (Costs 2 Actions)}{Mephistopheles creates a rain of Hellfire at a point he can see within 150 feet. Each creature within a 20-foot-radius sphere centered on that point must make a DC 24 Dexterity saving throw. Targets take 38 (7d10) fire damage on a failed save, or half as much damage on a successful one.
}
\DndMonsterLegendaryAction{Call Underling (Costs 3 Actions)}{Mephistopheles summons an allied \textbf{horned devil} in an unoccupied space that he can see.
}
\end{DndMonsterLegendaryActions}
\DndMonsterSection{Regional Effects}
The region containing Mephistopheles' lair is corrupted by his presence, which creates the following effect:


\begin{description}
\item[Geysers.]
Within 1 mile of Mephistopheles' lair, every 100 feet of movement results in a flame geyser erupting. Creatures within 10 feet of a geyser when it erupts must succeed on a DC 22 Dexterity saving throw or take 10 (3d6) fire damage. The geyser continues erupting for 24 hours.

\end{description}

If Mephistopheles dies, the effects fade over the course of 1d10 days.

\end{multicols}
\end{DndMonster}


\clearpage

\begin{DndMonster}[float=t]{Merregon}
\DndMonsterType{Medium fiend (devil), lawful evil}
\DndMonsterBasics[
armor-class = {16 (natural armor)},
hit-points = {\DndDice{6d8 + 18}},
speed = {30 ft.}
]
\DndMonsterAbilityScores[
 str = 18,
 dex = 14,
 con = 17,
 int = 6,
 wis = 12,
 cha = 8
]
\DndMonsterDetails[
damage-resistances = {cold; bludgeoning, piercing, slashing from nonmagical attacks that aren't silvered},
damage-immunities = {fire, poison},
condition-immunities = {frightened, poisoned},
senses = {darkvision 60 ft., passive perception 11},
languages = {understands Infernal but can't speak, telepathy 120 ft.},
challenge = 4,
]
\DndMonsterAction{Devil's Sight}
Magical darkness doesn't impede the merregon's darkvision.

\DndMonsterAction{Magic Resistance}
The merregon has advantage on saving throws against spells and other magical effects.

\DndMonsterSection{Actions}
\DndMonsterAction{Multiattack}
The merregon makes three Halberd attacks.

\DndMonsterAction{Halberd}
\textsl{Melee Weapon Attack:} 6 to hit, reach 10 ft., one target. \textsl{Hit:} 9 (1d10 + 4) slashing damage.

\DndMonsterAction{Heavy Crossbow}
\textsl{Ranged Weapon Attack:} 4 to hit, range 100/400 ft., one target. \textsl{Hit:} 7 (1d10 + 2) piercing damage.

\DndMonsterSection{Reactions}
\DndMonsterAction{Loyal Bodyguard}
When another Fiend within 5 feet of the merregon is hit by an attack roll, the merregon causes itself to be hit instead.

\end{DndMonster}



\begin{DndMonster}[float=t]{Miasmorne}
\DndMonsterType{Huge monstrosity, unaligned}
\DndMonsterBasics[
armor-class = {19 (natural armor)},
hit-points = {\DndDice{20d12 + 100}},
speed = {20 ft., burrow 20 ft.}
]
\DndMonsterAbilityScores[
 str = 22,
 dex = 10,
 con = 21,
 int = 3,
 wis = 10,
 cha = 10
]
\DndMonsterDetails[
saving-throws = {Str +11, Con +10},
skills = {Arcana +1, Athletics +16, Survival +5},
damage-immunities = {acid; fire; bludgeoning, piercing, slashing from nonmagical attacks},
condition-immunities = {petrified},
senses = {tremorsense 60 ft., passive perception 10},
challenge = 16,
]
\DndMonsterAction{Acidic Attacks}
If a creature takes acid damage from the miasmorne and that creature is wearing nonmagical metal armor, that armor is destroyed.

\DndMonsterAction{Acidic Hide}
After a nonmagical metal weapon hits the miasmorne that weapon melts and is destroyed. Likewise, nonmagical metal ammunition that hits the miasmorne is destroyed.

\DndMonsterSection{Actions}
\DndMonsterAction{Multiattack}
The miasmorne makes three attacks using Bite, Flechette or a combination of the two.

\DndMonsterAction{Bite}
\textsl{Melee Weapon Attack:} 11 to hit, reach 5 ft., one target. \textsl{Hit:} 22 (3d10 + 6) bludgeoning damage plus 7 (2d6) acid damage. If the target is a Large or smaller creature, it has the grappled condition (escape DC 18). The miasmorne can't make Bite attacks while a creature is grappled in this way.

\DndMonsterAction{Flechette}
\textsl{Ranged Weapon Attack:} 11 to hit, range 60/120 ft., one target. \textsl{Hit:} 16 (3d6 + 6) acid damage.

\DndMonsterAction{Flechette Spray (Recharge 5\textendash 6)}
The miasmorne launches flechettes from its fins in a 90-foot cone in front of it. All creatures in the cone must make a DC 19 Dexterity saving throw, taking 21 (6d6) acid damage on a failed save, or half as much damage on a successful one.

\DndMonsterAction{Acidic Secretion (1/Day)}
The miasmorne secretes an acidic solution, then sprays it around itself. All creatures in a 20-foot-radius sphere, centered on the miasmorne, must make a DC 19 Dexterity saving throw. Targets take 26 (4d12) acid damage on a failed save, or half as much damage on a successful one. All nonmagical metals within the radius that aren't being carried, as well as nonmagical metal items carried or worn by creatures that failed the save, are destroyed.

\end{DndMonster}



\begin{DndMonster}[float=t]{Mind Flayer}
\DndMonsterType{Medium aberration, lawful evil}
\DndMonsterBasics[
armor-class = {15 (\textit{breastplate})},
hit-points = {\DndDice{13d8 + 13}},
speed = {30 ft.}
]
\DndMonsterAbilityScores[
 str = 11,
 dex = 12,
 con = 12,
 int = 19,
 wis = 17,
 cha = 17
]
\DndMonsterDetails[
saving-throws = {Int +7, Wis +6, Cha +6},
skills = {Arcana +7, Deception +6, Insight +6, Perception +6, Persuasion +6, Stealth +4},
senses = {darkvision 120 ft., passive perception 16},
languages = {Deep Speech, Undercommon, telepathy 120 ft.},
challenge = 7,
]
\DndMonsterAction{Magic Resistance}
The mind flayer has advantage on saving throws against spells and other magical effects.

\DndMonsterAction{Innate Spellcasting (Psionics)}
The mind flayer's innate spellcasting ability is Intelligence (spell save DC 15). It can innately cast the following spells, requiring no components:
\begin{DndMonsterSpells}
  \DndInnateSpellLevel{\textit{detect thoughts}, \textit{levitate}}
  \DndInnateSpellLevel[1e]{\textit{dominate monster}, \textit{plane shift} (self only)}
\end{DndMonsterSpells}
\DndMonsterSection{Actions}
\DndMonsterAction{Tentacles}
\textsl{Melee Weapon Attack:} 7 to hit, reach 5 ft., one creature. \textsl{Hit:} 15 (2d10 + 4) psychic damage. If the target is Medium or smaller, it is grappled (escape DC 15) and must succeed on a DC 15 Intelligence saving throw or be stunned until this grapple ends.

\DndMonsterAction{Extract Brain}
\textsl{Melee Weapon Attack:} 7 to hit, reach 5 ft., one incapacitated humanoid grappled by the mind flayer. \textsl{Hit:} The target takes 55 (10d10) piercing damage. If this damage reduces the target to 0 hit points, the mind flayer kills the target by extracting and devouring its brain.

\DndMonsterAction{Mind Blast (Recharge 5\textendash 6)}
The mind flayer magically emits psychic energy in a 60-foot cone. Each creature in that area must succeed on a DC 15 Intelligence saving throw or take 22 (4d8 + 4) psychic damage and be stunned for 1 minute. A creature can repeat the saving throw at the end of each of its turns, ending the effect on itself on a success.

\end{DndMonster}



\begin{DndMonster}[float*=b,width=\textwidth + 8pt]{Mummy Lord}
\begin{multicols}{2}
\DndMonsterType{Medium undead, lawful evil}
\DndMonsterBasics[
armor-class = {17 (natural armor)},
hit-points = {\DndDice{13d8 + 39}},
speed = {20 ft.}
]
\DndMonsterAbilityScores[
 str = 18,
 dex = 10,
 con = 17,
 int = 11,
 wis = 18,
 cha = 16
]
\DndMonsterDetails[
saving-throws = {Con +8, Int +5, Wis +9, Cha +8},
skills = {History +5, Religion +5},
damage-vulnerabilities = {fire},
damage-immunities = {necrotic; poison; bludgeoning, piercing, slashing from nonmagical attacks},
condition-immunities = {charmed, exhaustion, frightened, paralyzed, poisoned},
senses = {darkvision 60 ft., passive perception 14},
languages = {the languages it knew in life},
challenge = 15,
]
\DndMonsterAction{Magic Resistance}
The mummy lord has advantage on saving throws against spells and other magical effects.

\DndMonsterAction{Rejuvenation}
A destroyed mummy lord gains a new body in 24 hours if its heart is intact, regaining all its hit points and becoming active again. The new body appears within 5 feet of the mummy lord's heart.

\DndMonsterAction{Spellcasting}
The mummy lord is a 10th-level spellcaster. Its spellcasting ability is Wisdom (spell save DC 17, 9 to hit with spell attacks). The mummy lord has the following cleric spells prepared:
\begin{DndMonsterSpells}
  \DndMonsterSpellLevel{\textit{sacred flame}, \textit{thaumaturgy}}
   \DndMonsterSpellLevel[1][4]{\textit{command}, \textit{guiding bolt}, \textit{shield of faith}}
   \DndMonsterSpellLevel[2][3]{\textit{hold person}, \textit{silence}, \textit{spiritual weapon}}
   \DndMonsterSpellLevel[3][3]{\textit{animate dead}, \textit{dispel magic}}
   \DndMonsterSpellLevel[4][3]{\textit{divination}, \textit{guardian of faith}}
   \DndMonsterSpellLevel[5][2]{\textit{contagion}, \textit{insect plague}}
   \DndMonsterSpellLevel[6][1]{\textit{harm}}
\end{DndMonsterSpells}
\DndMonsterSection{Actions}
\DndMonsterAction{Multiattack}
The mummy can use its Dreadful Glare and makes one attack with its rotting fist.

\DndMonsterAction{Rotting Fist}
\textsl{Melee Weapon Attack:} 9 to hit, reach 5 ft., one target. \textsl{Hit:} 14 (3d6 + 4) bludgeoning damage plus 21 (6d6) necrotic damage. If the target is a creature, it must succeed on a DC 16 Constitution saving throw or be cursed with mummy rot. The cursed target can't regain hit points, and its hit point maximum decreases by 10 (3d6) for every 24 hours that elapse. If the curse reduces the target's hit point maximum to 0, the target dies, and its body turns to dust. The curse lasts until removed by the \textit{remove curse} spell or other magic.

\DndMonsterAction{Dreadful Glare}
The mummy lord targets one creature it can see within 60 feet of it. If the target can see the mummy lord, it must succeed on a DC 16 Wisdom saving throw against this magic or become frightened until the end of the mummy's next turn. If the target fails the saving throw by 5 or more, it is also paralyzed for the same duration. A target that succeeds on the saving throw is immune to the Dreadful Glare of all mummies and mummy lords for the next 24 hours.

\DndMonsterSection{Legendary Actions}
Mummy Lord can take 3 legendary actions, choosing from the options below. Only one legendary action can be used at a time and only at the end of another creature's turn. Mummy Lord regains spent legendary actions at the start of its turn.

\begin{DndMonsterLegendaryActions}
\DndMonsterLegendaryAction{Attack}{The mummy lord makes one attack with its rotting fist or uses its Dreadful Glare.
}
\DndMonsterLegendaryAction{Blinding Dust}{Blinding dust and sand swirls magically around the mummy lord. Each creature within 5 feet of the mummy lord must succeed on a DC 16 Constitution saving throw or be blinded until the end of the creature's next turn.
}
\DndMonsterLegendaryAction{Blasphemous Word (Costs 2 Actions)}{The mummy lord utters a blasphemous word. Each non-undead creature within 10 feet of the mummy lord that can hear the magical utterance must succeed on a DC 16 Constitution saving throw or be stunned until the end of the mummy lord's next turn.
}
\DndMonsterLegendaryAction{Channel Negative Energy (Costs 2 Actions)}{The mummy lord magically unleashes negative energy. Creatures within 60 feet of the mummy lord, including ones behind barriers and around corners, can't regain hit points until the end of the mummy lord's next turn.
}
\DndMonsterLegendaryAction{Whirlwind of Sand (Costs 2 Actions)}{The mummy lord magically transforms into a whirlwind of sand, moves up to 60 feet, and reverts to its normal form. While in whirlwind form, the mummy lord is immune to all damage, and it can't be grappled, petrified, knocked prone, restrained, or stunned. Equipment worn or carried by the mummy lord remain in its possession.
}
\end{DndMonsterLegendaryActions}
\DndMonsterSection{Lair Actions}
On initiative count 20 (losing initiative ties), the mummy lord takes a lair action to cause one of the following effects; the mummy lord can't use the same effect two rounds in a row:


\begin{itemize}
\item Each undead creature in the lair can pinpoint the location of each living creature within 120 feet of it until initiative count 20 on the next round.
\item Each undead in the lair has advantage on saving throws against effects that turn undead until initiative count 20 on the next round.
\item Until initiative count 20 on the next round, any non-undead creature that tries to cast a spell of 4th level or lower in the mummy lord's lair is wracked with pain. The creature can choose another action, but if it tries to cast the spell, it must make a DC 16 Constitution saving throw. On a failed save, it takes 1d6 necrotic damage per level of the spell, and the spell has no effect and is wasted.
\end{itemize}

\DndMonsterSection{Regional Effects}
A mummy lord's temple or tomb is warped in any of the following ways by the creature's dark presence:


\begin{itemize}
\item Food instantly molders and water instantly evaporates when brought into the lair. Other non magical drinks are spoiled - wine turning to vinegar, for instance.
\item \textit{Divination} spells cast within the lair by creatures other than the mummy lord have a 25 percent chance to provide misleading results, as determined by the DM. If a \textit{divination} spell already has a chance to fail or become unreliable when cast multiple times, that chance increases by 25 percent.
\item A creature that takes treasure from the lair is cursed until the treasure is returned. The cursed target has disadvantage on all saving throws. The curse lasts until removed by a \textit{remove curse} spell or other magic.
\end{itemize}

If the mummy lord is destroyed, these regional effects end immediately.

\end{multicols}
\end{DndMonster}



\begin{DndMonster}[float=t]{Nalfeshnee}
\DndMonsterType{Large fiend (demon), chaotic evil}
\DndMonsterBasics[
armor-class = {18 (natural armor)},
hit-points = {\DndDice{16d10 + 96}},
speed = {20 ft., fly 30 ft.}
]
\DndMonsterAbilityScores[
 str = 21,
 dex = 10,
 con = 22,
 int = 19,
 wis = 12,
 cha = 15
]
\DndMonsterDetails[
saving-throws = {Con +11, Int +9, Wis +6, Cha +7},
damage-resistances = {cold; fire; lightning; bludgeoning, piercing, slashing from nonmagical attacks},
damage-immunities = {poison},
condition-immunities = {poisoned},
senses = {truesight 120 ft., passive perception 11},
languages = {Abyssal, telepathy 120 ft.},
challenge = 13,
]
\DndMonsterAction{Magic Resistance}
The nalfeshnee has advantage on saving throws against spells and other magical effects.

\DndMonsterSection{Actions}
\DndMonsterAction{Multiattack}
The nalfeshnee uses Horror Nimbus if it can. It then makes three attacks: one with its bite and two with its claws.

\DndMonsterAction{Bite}
\textsl{Melee Weapon Attack:} 10 to hit, reach 5 ft., one target. \textsl{Hit:} 32 (5d10 + 5) piercing damage.

\DndMonsterAction{Claw}
\textsl{Melee Weapon Attack:} 10 to hit, reach 10 ft., one target. \textsl{Hit:} 15 (3d6 + 5) slashing damage.

\DndMonsterAction{Horror Nimbus (Recharge 5\textendash 6)}
The nalfeshnee magically emits scintillating, multicolored light. Each creature within 15 feet of the nalfeshnee that can see the light must succeed on a DC 15 Wisdom saving throw or be frightened for 1 minute. A creature can repeat the saving throw at the end of each of its turns, ending the effect on itself on a success. If a creature's saving throw is successful or the effect ends for it, the creature is immune to the nalfeshnee's Horror Nimbus for the next 24 hours.

\DndMonsterAction{Teleport}
The nalfeshnee magically teleports, along with any equipment it is wearing or carrying, up to 120 feet to an unoccupied space it can see.

\end{DndMonster}



\begin{DndMonster}[float=t]{Neogi}
\DndMonsterType{Small aberration, lawful evil}
\DndMonsterBasics[
armor-class = {15 (natural armor)},
hit-points = {\DndDice{6d6 + 12}},
speed = {30 ft., climb 30 ft.}
]
\DndMonsterAbilityScores[
 str = 6,
 dex = 16,
 con = 14,
 int = 13,
 wis = 12,
 cha = 15
]
\DndMonsterDetails[
skills = {Intimidation +4, Perception +3},
senses = {darkvision 60 ft., passive perception 13},
languages = {Common, Deep Speech, Undercommon},
challenge = 3,
]
\DndMonsterAction{Mental Fortitude}
The neogi has advantage on saving throws against being charmed or frightened, and magic can't put the neogi to sleep.

\DndMonsterAction{Spider Climb}
The neogi can climb difficult surfaces, including upside down on ceilings, without needing to make an ability check.

\DndMonsterSection{Actions}
\DndMonsterAction{Multiattack}
The neogi makes one Bite attack and two Claw attacks.

\DndMonsterAction{Bite}
\textsl{Melee Weapon Attack:} 5 to hit, reach 5 ft., one target. \textsl{Hit:} 6 (1d6 + 3) piercing damage plus 14 (4d6) poison damage, and the target must succeed on a DC 12 Constitution saving throw or become poisoned for 1 minute. A target can repeat the saving throw at the end of each of its turns, ending the effect on itself on a success.

\DndMonsterAction{Claw}
\textsl{Melee Weapon Attack:} 5 to hit, reach 5 ft., one target. \textsl{Hit:} 8 (2d4 + 3) slashing damage.

\DndMonsterSection{Bonus Actions}
\DndMonsterAction{Enslave (Recharges after a Short or Long Rest)}
The neogi targets one creature it can see within 30 feet of it. The target must succeed on a DC 14 Wisdom saving throw or be magically charmed by the neogi for 1 day, or until the neogi dies or is more than 1 mile from the target. The charmed target obeys the neogi's commands and can't take reactions, and the neogi and the target can communicate telepathically with each other at a distance of up to 1 mile. Whenever the charmed target takes damage, it can repeat the saving throw, ending the effect on itself on a success.

\end{DndMonster}



\begin{DndMonster}[float=t]{Neogi Master}
\DndMonsterType{Medium aberration (warlock), lawful evil}
\DndMonsterBasics[
armor-class = {15 (natural armor)},
hit-points = {\DndDice{11d8 + 22}},
speed = {30 ft., climb 30 ft.}
]
\DndMonsterAbilityScores[
 str = 6,
 dex = 16,
 con = 14,
 int = 16,
 wis = 12,
 cha = 18
]
\DndMonsterDetails[
saving-throws = {Wis +3},
skills = {Arcana +5, Deception +6, Intimidation +6, Perception +3, Persuasion +6},
senses = {darkvision 120 ft., passive perception 13},
languages = {Common, Deep Speech, Undercommon, telepathy 30 ft.},
challenge = 4,
]
\DndMonsterAction{Devil's Sight}
Magical darkness doesn't impede the neogi's darkvision.

\DndMonsterAction{Mental Fortitude}
The neogi has advantage on saving throws against being charmed or frightened, and magic can't put the neogi to sleep.

\DndMonsterAction{Spider Climb}
The neogi can climb difficult surfaces, including upside down on ceilings, without needing to make an ability check.

\DndMonsterAction{Spellcasting}
The neogi casts one of the following spells, using Charisma as the spellcasting ability (spell save DC 14):
\begin{DndMonsterSpells}
  \DndInnateSpellLevel{\textit{guidance}, \textit{mage hand}, \textit{minor illusion}, \textit{prestidigitation}}
  \DndInnateSpellLevel[1e]{\textit{dimension door}, \textit{hold person}, \textit{hunger of Hadar}}
\end{DndMonsterSpells}
\DndMonsterSection{Actions}
\DndMonsterAction{Multiattack}
The neogi makes one Bite attack and one Claw attack, or it makes two Tentacle of Hadar attacks.

\DndMonsterAction{Bite}
\textsl{Melee Weapon Attack:} 5 to hit, reach 5 ft., one target. \textsl{Hit:} 6 (1d6 + 3) piercing damage plus 14 (4d6) poison damage, and the target must succeed on a DC 12 Constitution saving throw or become poisoned for 1 minute. A target can repeat the saving throw at the end of each of its turns, ending the effect on itself on a success.

\DndMonsterAction{Claw}
\textsl{Melee Weapon Attack:} 5 to hit, reach 5 ft., one target. \textsl{Hit:} 8 (2d4 + 3) piercing damage.

\DndMonsterAction{Tentacle of Hadar}
\textsl{Ranged Spell Attack:} 6 to hit, range 120 ft., one target. \textsl{Hit:} 14 (3d6 + 4) necrotic damage, and the target can't take reactions until the end of the neogi's next turn, as a spectral tentacle clings to the target.

\DndMonsterSection{Bonus Actions}
\DndMonsterAction{Enslave (Recharges after a Short or Long Rest)}
The neogi targets one creature it can see within 30 feet of it. The target must succeed on a DC 14 Wisdom saving throw or be magically charmed by the neogi for 1 day, or until the neogi dies or is more than 1 mile from the target. The charmed target obeys the neogi's commands and can't take reactions, and the neogi and the target can communicate telepathically with each other at a distance of up to 1 mile. Whenever the charmed target takes damage, it can repeat the saving throw, ending the effect on itself on a success.

\end{DndMonster}



\begin{DndMonster}[float*=b,width=\textwidth + 8pt]{Night Hag}
\begin{multicols}{2}
\DndMonsterType{Medium fiend, neutral evil}
\DndMonsterBasics[
armor-class = {17 (natural armor)},
hit-points = {\DndDice{15d8 + 45}},
speed = {30 ft.}
]
\DndMonsterAbilityScores[
 str = 18,
 dex = 15,
 con = 16,
 int = 16,
 wis = 14,
 cha = 16
]
\DndMonsterDetails[
skills = {Deception +6, Insight +5, Perception +5, Stealth +5},
damage-resistances = {cold; fire; bludgeoning, piercing, slashing from nonmagical attacks that aren't silvered},
condition-immunities = {charmed},
senses = {darkvision 120 ft., passive perception 16},
languages = {Abyssal, Common, Infernal, Primordial},
challenge = 5,
]
\DndMonsterAction{Magic Resistance}
The hag has advantage on saving throws against spells and other magical effects.

\DndMonsterAction{Night Hag Items}
A night hag carries two very rare magic items that she must craft for herself. If either object is lost, the night hag will go to great lengths to retrieve it, as creating a new tool takes time and effort.

Heartstone: This lustrous black gem allows a night hag to become ethereal while it is in her possession. The touch of a heartstone also cures any disease. Crafting a heartstone takes 30 days.

Soul Bag: When an evil humanoid dies as a result of a night hag's Nightmare Haunting, the hag catches the soul in this black sack made of stitched flesh. A soul bag can hold only one evil soul at a time, and only the night hag who crafted the bag can catch a soul with it. Crafting a soul bag takes 7 days and a humanoid sacrifice (whose flesh is used to make the bag).

\DndMonsterAction{Innate Spellcasting}
The hag's innate spellcasting ability is Charisma (spell save DC 14, 6 to hit with spell attacks). She can innately cast the following spells, requiring no material components:
\begin{DndMonsterSpells}
  \DndInnateSpellLevel{\textit{detect magic}, \textit{magic missile}}
  \DndInnateSpellLevel[2e]{\textit{plane shift} (self only), \textit{ray of enfeeblement}, \textit{sleep}}
\end{DndMonsterSpells}
\DndMonsterSection{Actions}
\DndMonsterAction{Claws (Hag Form Only)}
\textsl{Melee Weapon Attack:} 7 to hit, reach 5 ft., one target. \textsl{Hit:} 13 (2d8 + 4) slashing damage.

\DndMonsterAction{Change Shape}
The hag magically polymorphs into a Small or Medium female humanoid, or back into her true form. Her statistics are the same in each form. Any equipment she is wearing or carrying isn't transformed. She reverts to her true form if she dies.

\DndMonsterAction{Etherealness}
The hag magically enters the Ethereal Plane from the Material Plane, or vice versa. To do so, the hag must have a heartstone in her possession.

\DndMonsterAction{Nightmare Haunting (1/Day)}
While on the Ethereal Plane, the hag magically touches a sleeping humanoid on the Material Plane. A \textit{protection from evil and good} spell cast on the target prevents this contact, as does a magic circle. As long as the contact persists, the target has dreadful visions. If these visions last for at least 1 hour, the target gains no benefit from its rest, and its hit point maximum is reduced by 5 (1d10). If this effect reduces the target's hit point maximum to 0, the target dies, and if the target was evil, its soul is trapped in the hag's soul bag. The reduction to the target's hit point maximum lasts until removed by the  \textit{greater restoration} spell or similar magic.

\end{multicols}
\end{DndMonster}



\begin{DndMonster}[float=t]{Nothic}
\DndMonsterType{Medium aberration, neutral evil}
\DndMonsterBasics[
armor-class = {15 (natural armor)},
hit-points = {\DndDice{6d8 + 18}},
speed = {30 ft.}
]
\DndMonsterAbilityScores[
 str = 14,
 dex = 16,
 con = 16,
 int = 13,
 wis = 10,
 cha = 8
]
\DndMonsterDetails[
skills = {Arcana +3, Insight +4, Perception +2, Stealth +5},
senses = {truesight 120 ft., passive perception 12},
languages = {Undercommon},
challenge = 2,
]
\DndMonsterAction{Keen Sight}
The nothic has advantage on Wisdom (Perception) checks that rely on sight.

\DndMonsterSection{Actions}
\DndMonsterAction{Multiattack}
The nothic makes two claw attacks.

\DndMonsterAction{Claw}
\textsl{Melee Weapon Attack:} 4 to hit, reach 5 ft., one target. \textsl{Hit:} 6 (1d6 + 3) slashing damage.

\DndMonsterAction{Rotting Gaze}
The nothic targets one creature it can see within 30 feet of it. The target must succeed on a DC 12 Constitution saving throw against this magic or take 10 (3d6) necrotic damage.

\DndMonsterAction{Weird Insight}
The nothic targets one creature it can see within 30 feet of it. The target must contest its Charisma (Deception) check against the nothic's Wisdom (Insight) check. If the nothic wins, it magically learns one fact or secret about the target. The target automatically wins if it is immune to being charmed.

\end{DndMonster}



\begin{DndMonster}[float=t]{Nupperibo}
\DndMonsterType{Medium fiend (devil), lawful evil}
\DndMonsterBasics[
armor-class = {13 (natural armor)},
hit-points = {\DndDice{2d8 + 2}},
speed = {20 ft.}
]
\DndMonsterAbilityScores[
 str = 16,
 dex = 11,
 con = 13,
 int = 3,
 wis = 8,
 cha = 1
]
\DndMonsterDetails[
skills = {Perception +1},
damage-resistances = {acid, cold},
damage-immunities = {fire, poison},
condition-immunities = {blinded, charmed, frightened, poisoned},
senses = {blindsight 20 ft. (blind beyond this radius), passive perception 11},
languages = {understands Infernal but can't speak},
challenge = 1/2,
]
\DndMonsterAction{Cloud of Vermin}
Any creature, other than a devil, that starts its turn within 20 feet of one or more nupperibos must succeed on a DC 11 Constitution saving throw or take 5 (2d4) acid damage. A creature within the areas of two or more nupperibos makes the saving throw with disadvantage.

\DndMonsterAction{Driven Tracker}
In the Nine Hells, the nupperibo can flawlessly track any creature that has taken damage from any nupperibo's Cloud of Vermin within the previous 24 hours.

\DndMonsterSection{Actions}
\DndMonsterAction{Bite}
\textsl{Melee Weapon Attack:} 5 to hit, reach 5 ft., one target. \textsl{Hit:} 6 (1d6 + 3) piercing damage.

\end{DndMonster}



\begin{DndMonster}[float=t]{Ogre}
\DndMonsterType{Large giant, chaotic evil}
\DndMonsterBasics[
armor-class = {11 (\textit{hide armor})},
hit-points = {\DndDice{7d10 + 21}},
speed = {40 ft.}
]
\DndMonsterAbilityScores[
 str = 19,
 dex = 8,
 con = 16,
 int = 5,
 wis = 7,
 cha = 7
]
\DndMonsterDetails[
senses = {darkvision 60 ft., passive perception 8},
languages = {Common, Giant},
challenge = 2,
]
\DndMonsterSection{Actions}
\DndMonsterAction{Greatclub}
\textsl{Melee Weapon Attack:} 6 to hit, reach 5 ft., one target. \textsl{Hit:} 13 (2d8 + 4) bludgeoning damage.

\DndMonsterAction{Javelin}
\textsl{Melee or Ranged Weapon Attack:} 6 to hit, reach 5 ft. or range 30/120 ft., one target. \textsl{Hit:} 11 (2d6 + 4) piercing damage.

\end{DndMonster}



\begin{DndMonster}[float=t]{Oneirovore}
\DndMonsterType{Large fiend, unaligned}
\DndMonsterBasics[
armor-class = {19 (natural armor)},
hit-points = {\DndDice{19d10 + 133}},
speed = {20 ft.}
]
\DndMonsterAbilityScores[
 str = 16,
 dex = 12,
 con = 25,
 int = 12,
 wis = 10,
 cha = 10
]
\DndMonsterDetails[
saving-throws = {Con +11, Int +5},
skills = {Insight +4, Nature +5, Perception +4, Survival +4},
damage-resistances = {fire; poison; bludgeoning, piercing, slashing from nonmagical attacks that aren't silvered},
damage-immunities = {psychic},
condition-immunities = {blinded, charmed, frightened},
senses = {truesight 60 ft., passive perception 14},
challenge = 11,
]
\DndMonsterSection{Actions}
\DndMonsterAction{Stomp}
\textsl{Melee Weapon Attack:} 7 to hit, reach 10 ft., one target. \textsl{Hit:} 16 (2d12 + 3) force damage.

\DndMonsterAction{Realize Nightmares (Recharge 4\textendash 6)}
The oneirovore manipulates the perception of nearby creatures, sending them into a panic. Creatures of the oneirovore's choice that it can see within 60 feet must make a DC 19 Charisma saving throw. On a failed save, the creature has the charmed condition, their vision distorted such that they view their allies as creatures of pure fear and hatred. While charmed, a creature will use their actions to attack their allies with intent to kill. A charmed creature may repeat the saving throw at the end of each turn, ending the effect on a success.

\DndMonsterAction{Warning Cry (1/Day)}
The oneirovore lets out an earsplitting shriek, calling for infernal reinforcements. Creatures that are not Fiends that are within 120 feet of the oneirovore must make a DC 19 Constitution saving throw, taking 55 (10d10) thunder damage on a failed save, or half as much damage on a successful one. Any lesser devils within a 1-mile-radius, centered on the oneirovore, are required by contract to come to its aid upon hearing its Warning Cry. Higher ranking devils may investigate if they desire, but the call holds no magical sway over them.

\end{DndMonster}



\begin{DndMonster}[float=t]{Oni}
\DndMonsterType{Large giant, lawful evil}
\DndMonsterBasics[
armor-class = {16 (\textit{chain mail})},
hit-points = {\DndDice{13d10 + 39}},
speed = {30 ft., fly 30 ft.}
]
\DndMonsterAbilityScores[
 str = 19,
 dex = 11,
 con = 16,
 int = 14,
 wis = 12,
 cha = 15
]
\DndMonsterDetails[
saving-throws = {Dex +3, Con +6, Wis +4, Cha +5},
skills = {Arcana +5, Deception +8, Perception +4},
senses = {darkvision 60 ft., passive perception 14},
languages = {Common, Giant},
challenge = 7,
]
\DndMonsterAction{Magic Weapons}
The oni's weapon attacks are magical.

\DndMonsterAction{Regeneration}
The oni regains 10 hit points at the start of its turn if it has at least 1 hit point.

\DndMonsterAction{Innate Spellcasting}
The oni's innate spellcasting ability is Charisma (spell save DC 13). The oni can innately cast the following spells, requiring no material components:
\begin{DndMonsterSpells}
  \DndInnateSpellLevel{\textit{darkness}, \textit{invisibility}}
  \DndInnateSpellLevel[1e]{\textit{charm person}, \textit{cone of cold}, \textit{gaseous form}, \textit{sleep}}
\end{DndMonsterSpells}
\DndMonsterSection{Actions}
\DndMonsterAction{Multiattack}
The oni makes two attacks, either with its claws or its glaive.

\DndMonsterAction{Claw (Oni Form Only)}
\textsl{Melee Weapon Attack:} 7 to hit, reach 5 ft., one target. \textsl{Hit:} 8 (1d8 + 4) slashing damage.

\DndMonsterAction{Glaive}
\textsl{Melee Weapon Attack:} 7 to hit, reach 10 ft., one target. \textsl{Hit:} 15 (2d10 + 4) slashing damage, or 9 (1d10 + 4) slashing damage in Small or Medium form.

\DndMonsterAction{Change Shape}
The oni magically polymorphs into a Small or Medium humanoid, into a Large giant, or back into its true form. Other than its size, its statistics are the same in each form. The only equipment that is transformed is its glaive, which shrinks so that it can be wielded in humanoid form. If the oni dies, it reverts to its true form, and its glaive reverts to its normal size.

\end{DndMonster}



\begin{DndMonster}[float=t]{Orc}
\DndMonsterType{Medium humanoid (orc), chaotic evil}
\DndMonsterBasics[
armor-class = {13 (\textit{hide armor})},
hit-points = {\DndDice{2d8 + 6}},
speed = {30 ft.}
]
\DndMonsterAbilityScores[
 str = 16,
 dex = 12,
 con = 16,
 int = 7,
 wis = 11,
 cha = 10
]
\DndMonsterDetails[
skills = {Intimidation +2},
senses = {darkvision 60 ft., passive perception 10},
languages = {Common, Orc},
challenge = 1/2,
]
\DndMonsterAction{Aggressive}
As a bonus action, the orc can move up to its speed toward a hostile creature that it can see.

\DndMonsterSection{Actions}
\DndMonsterAction{Greataxe}
\textsl{Melee Weapon Attack:} 5 to hit, reach 5 ft., one target. \textsl{Hit:} 9 (1d12 + 3) slashing damage.

\DndMonsterAction{Javelin}
\textsl{Melee or Ranged Weapon Attack:} 5 to hit, reach 5 ft. or range 30/120 ft., one target. \textsl{Hit:} 6 (1d6 + 3) piercing damage.

\end{DndMonster}



\begin{DndMonster}[float=t]{Orthon}
\DndMonsterType{Large fiend (devil), lawful evil}
\DndMonsterBasics[
armor-class = {17 (half plate armor)},
hit-points = {\DndDice{10d10 + 50}},
speed = {30 ft., climb 30 ft.}
]
\DndMonsterAbilityScores[
 str = 22,
 dex = 16,
 con = 21,
 int = 15,
 wis = 15,
 cha = 16
]
\DndMonsterDetails[
saving-throws = {Dex +7, Con +9, Wis +6},
skills = {Perception +10, Stealth +11, Survival +10},
damage-resistances = {cold; bludgeoning, piercing, slashing from nonmagical attacks that aren't silvered},
damage-immunities = {fire, poison},
condition-immunities = {charmed, exhaustion, poisoned},
senses = {darkvision 120 ft., truesight 30 ft., passive perception 20},
languages = {Common, Infernal, telepathy 120 ft.},
challenge = 10,
]
\DndMonsterAction{Magic Resistance}
The orthon has advantage on saving throws against spells and other magical effects.

\DndMonsterSection{Actions}
\DndMonsterAction{Infernal Dagger}
\textsl{Melee Weapon Attack:} 10 to hit, reach 5 ft., one target. \textsl{Hit:} 11 (2d4 + 6) force damage, and the target must make a DC 17 Constitution saving throw, taking 22 (4d10) poison damage on a failed save, or half as much damage on a successful one. On a failure, the target is poisoned for 1 minute. The poisoned target can repeat the saving throw at the end of each of its turns, ending the effect on itself on a success.

\DndMonsterAction{Brass Crossbow}
\textsl{Ranged Weapon Attack:} 7 to hit, range 100/400 ft., one target. \textsl{Hit:} 14 (2d10 + 3) force damage. The target also suffers one of the following effects of the orthon's choice; the orthon can't use the same effect two rounds in a row:


\begin{description}
\item[Acid.]
The target must make a DC 17 Constitution saving throw, taking an additional 17 (5d6) acid damage on a failed save, or half as much damage on a successful one.
\item[Blindness.]
The target takes 5 (1d10) radiant damage. In addition, the target and all other creatures within 20 feet of it must each make a successful DC 17 Dexterity saving throw or be blinded until the end of the orthon's next turn.
\item[Concussion.]
The target and each creature within 20 feet of it must make a DC 17 Constitution saving throw, taking 13 (2d12) thunder damage on a failed save, or half as much damage on a successful one.
\item[Entanglement.]
The target must make a successful DC 17 Dexterity saving throw or be restrained for 1 hour by strands of sticky webbing. The target can escape by taking an action to make a DC 17 Strength or Dexterity check and succeeding.
\item[Paralysis.]
The target takes 22 (4d10) lightning damage and must make a successful DC 17 Constitution saving throw or be paralyzed for 1 minute. The paralyzed target can repeat the saving throw at the end of each of its turns, ending the effect on itself on a success.
\item[Tracking.]
For the next 24 hours, the orthon knows the direction and distance to the target, as long as it's on the same plane of existence. If the target is on a different plane, the orthon knows which one, but not the exact location there.
\end{description}

\DndMonsterSection{Bonus Actions}
\DndMonsterAction{Invisibility Field (Recharge 4\textendash 6)}
The orthon becomes invisible. Any equipment it wears or carries is also invisible as long as the equipment is on its person. This invisibility ends immediately after it makes an attack roll or is hit by an attack roll.

\DndMonsterSection{Reactions}
\DndMonsterAction{Explosive Retribution}
In response to dropping to 15 hit points or fewer, the orthon explodes. All other creatures within 30 feet of it must each make a DC 17 Dexterity saving throw, taking 9 (2d8) fire damage plus 9 (2d8) thunder damage on a failed save, or half as much damage on a successful one. The orthon, its infernal dagger, and its brass crossbow are destroyed.

\end{DndMonster}



\begin{DndMonster}[float=t]{Pain Devil (Excruciarch)}
\DndMonsterType{Medium fiend (devil), lawful evil}
\DndMonsterBasics[
armor-class = {16 (natural armor)},
hit-points = {\DndDice{18d8 + 90}},
speed = {30 ft.}
]
\DndMonsterAbilityScores[
 str = 18,
 dex = 15,
 con = 20,
 int = 11,
 wis = 10,
 cha = 13
]
\DndMonsterDetails[
saving-throws = {Str +8, Con +9},
skills = {Deception +5, Insight +4, Intimidation +5, Perception +4},
damage-resistances = {cold; bludgeoning, piercing, slashing from nonmagical attacks that aren't silvered},
damage-immunities = {fire, poison},
condition-immunities = {poisoned},
senses = {darkvision 120 ft., passive perception 14},
languages = {Celestial, Common, Infernal, telepathy 120 ft.},
challenge = 12,
]
\DndMonsterAction{Aura of Torment}
Any creature hostile to the excruciarch that starts its turn within 20 feet of the excruciarch takes 10 (4d4) slashing damage, unless the excruciarch is incapacitated.

\DndMonsterAction{Devil's Sight}
Magical darkness doesn't impede the excruciarch's darkvision.

\DndMonsterAction{Magic Resistance}
The excruciarch has advantage on saving throws against spells and other magical effects.

\DndMonsterAction{Sadism}
The excruciarch gains a +1 bonus to attack and damage rolls for each creature it damaged on its previous turn.

\DndMonsterSection{Actions}
\DndMonsterAction{Multiattack}
The excruciarch makes two Scourge attacks. It can replace one of the attacks with Storm of Steel (if available).

\DndMonsterAction{Scourge}
\textsl{Melee Weapon Attack:} 8 to hit, reach 10 ft., one target. \textsl{Hit:} 13 (2d8 + 4) bludgeoning damage plus 13 (2d8 + 4) slashing damage.

\DndMonsterAction{Storm of Steel (Recharge 4\textendash 6)}
The excruciarch swings its scourge wildly around itself. All creatures within 15 feet of the excruciarch must make a DC 16 Dexterity saving throw. Targets take 31 (6d8 + 4) slashing damage on a failed save, or half as much damage on a successful one.

\DndMonsterSection{Reactions}
\DndMonsterAction{Vulnerable Gaze}
As a reaction to a creature resisting damage dealt by the excruciarch, it turns its gaze on that creature, negating any resistances and immunities that creature has until the start of the excruciarch's next turn.

\end{DndMonster}



\begin{DndMonster}[float=t]{Pentadrone}
\DndMonsterType{Large construct, lawful neutral}
\DndMonsterBasics[
armor-class = {16 (natural armor)},
hit-points = {\DndDice{5d10 + 5}},
speed = {40 ft.}
]
\DndMonsterAbilityScores[
 str = 15,
 dex = 14,
 con = 12,
 int = 10,
 wis = 10,
 cha = 13
]
\DndMonsterDetails[
skills = {Perception +4},
senses = {truesight 120 ft., passive perception 14},
languages = {Modron},
challenge = 2,
]
\DndMonsterAction{Axiomatic Mind}
The pentadrone can't be compelled to act in a manner contrary to its nature or its instructions.

\DndMonsterAction{Disintegration}
If the pentadrone dies, its body disintegrates into dust, leaving behind its weapons and anything else it was carrying.

\DndMonsterSection{Actions}
\DndMonsterAction{Multiattack}
The pentadrone makes five arm attacks.

\DndMonsterAction{Arm}
\textsl{Melee Weapon Attack:} 4 to hit, reach 5 ft., one target. \textsl{Hit:} 5 (1d6 + 2) bludgeoning damage.

\DndMonsterAction{Paralysis Gas (Recharge 5\textendash 6)}
The pentadrone exhales a 30-foot cone of gas. Each creature in that area must succeed on a DC 11 Constitution saving throw or be paralyzed for 1 minute. A creature can repeat the saving throw at the end of each of its turns, ending the effect on itself on a success.

\end{DndMonster}



\begin{DndMonster}[float*=b,width=\textwidth + 8pt]{Phoenix}
\begin{multicols}{2}
\DndMonsterType{Gargantuan elemental, neutral}
\DndMonsterBasics[
armor-class = {18},
hit-points = {\DndDice{10d20 + 70}},
speed = {20 ft., fly 120 ft.}
]
\DndMonsterAbilityScores[
 str = 19,
 dex = 26,
 con = 25,
 int = 2,
 wis = 21,
 cha = 18
]
\DndMonsterDetails[
saving-throws = {Wis +10, Cha +9},
damage-resistances = {bludgeoning, piercing, slashing from nonmagical attacks},
damage-immunities = {fire, poison},
condition-immunities = {exhaustion, grappled, paralyzed, petrified, poisoned, prone, restrained, stunned},
senses = {darkvision 60 ft., passive perception 15},
challenge = 16,
]
\DndMonsterAction{Fiery Death and Rebirth}
If the phoenix dies, it explodes. Each creature in 60-foot-radius sphere centered on the phoenix must make a DC 20 Dexterity saving throw, taking 22 (4d10) fire damage on a failed save, or half as much damage on a successful one. The fire ignites flammable objects in the area that aren't being worn or carried.

The explosion destroys the phoenix's body and leaves behind an egg-shaped cinder, which weighs 5 pounds. The cinder deals 21 (6d6) fire damage to any creature that touches it, though no more than once per round. The cinder is immune to all damage, and after 1d6 days, it hatches a new phoenix.

\DndMonsterAction{Fire Form}
The phoenix can move through a space as narrow as 1 inch wide without squeezing.

Any creature that touches the phoenix or hits it with a melee attack while within 5 feet of it takes 5 (1d10) fire damage. In addition, the phoenix can enter a hostile creature's space and stop there. The first time it enters a creature's space on a turn, that creature takes 5 (1d10) fire damage.

With a touch, the phoenix can also ignite flammable objects that aren't worn or carried (no action required).

\DndMonsterAction{Flyby}
The phoenix doesn't provoke opportunity attack when it flies out of an enemy's reach.

\DndMonsterAction{Illumination}
The phoenix sheds bright light in a 60-foot radius and dim light for an additional 30 feet.

\DndMonsterAction{Legendary Resistance (3/Day)}
If the phoenix fails a saving throw, it can choose to succeed instead.

\DndMonsterAction{Siege Monster}
The phoenix deals double damage to objects and structures.

\DndMonsterSection{Actions}
\DndMonsterAction{Multiattack}
The phoenix makes two attacks: one Beak attack and one Fiery Talons attack.

\DndMonsterAction{Beak}
\textsl{Melee Weapon Attack:} 13 to hit, reach 15 ft., one target. \textsl{Hit:} 15 (2d6 + 8) fire damage.

\DndMonsterAction{Fiery Talons}
\textsl{Melee Weapon Attack:} 13 to hit, reach 15 ft., one target. \textsl{Hit:} 17 (2d8 + 8) fire damage.

\DndMonsterSection{Legendary Actions}
Phoenix can take 3 legendary actions, choosing from the options below. Only one legendary action can be used at a time and only at the end of another creature's turn. Phoenix regains spent legendary actions at the start of its turn.

\begin{DndMonsterLegendaryActions}
\DndMonsterLegendaryAction{Move}{The phoenix moves up to its speed.
}
\DndMonsterLegendaryAction{Peck}{The phoenix makes one beak attack.
}
\DndMonsterLegendaryAction{Swoop (Costs 2 Actions)}{The phoenix moves up to its speed and makes one Fiery Talons attack.
}
\end{DndMonsterLegendaryActions}
\end{multicols}
\end{DndMonster}



\begin{DndMonster}[float=t]{Pit Fiend}
\DndMonsterType{Large fiend (devil), lawful evil}
\DndMonsterBasics[
armor-class = {19 (natural armor)},
hit-points = {\DndDice{24d10 + 168}},
speed = {30 ft., fly 60 ft.}
]
\DndMonsterAbilityScores[
 str = 26,
 dex = 14,
 con = 24,
 int = 22,
 wis = 18,
 cha = 24
]
\DndMonsterDetails[
saving-throws = {Dex +8, Con +13, Wis +10},
damage-resistances = {cold; bludgeoning, piercing, slashing from nonmagical attacks that aren't silvered},
damage-immunities = {fire, poison},
condition-immunities = {poisoned},
senses = {truesight 120 ft., passive perception 14},
languages = {Infernal, telepathy 120 ft.},
challenge = 20,
]
\DndMonsterAction{Fear Aura}
Any creature hostile to the pit fiend that starts its turn within 20 feet of the pit fiend must make a DC 21 Wisdom saving throw, unless the pit fiend is incapacitated. On a failed save, the creature is frightened until the start of its next turn. If a creature's saving throw is successful, the creature is immune to the pit fiend's Fear Aura for the next 24 hours.

\DndMonsterAction{Magic Resistance}
The pit fiend has advantage on saving throws against spells and other magical effects.

\DndMonsterAction{Magic Weapons}
The pit fiend's weapon attacks are magical.

\DndMonsterAction{Innate Spellcasting}
The pit fiend's spellcasting ability is Charisma (spell save DC 21). The pit fiend can innately cast the following spells, requiring no material components:
\begin{DndMonsterSpells}
  \DndInnateSpellLevel{\textit{detect magic}, \textit{fireball}}
  \DndInnateSpellLevel[3e]{\textit{hold monster}, \textit{wall of fire}}
\end{DndMonsterSpells}
\DndMonsterSection{Actions}
\DndMonsterAction{Multiattack}
The pit fiend makes four attacks: one with its bite, one with its claw, one with its mace, and one with its tail.

\DndMonsterAction{Bite}
\textsl{Melee Weapon Attack:} 14 to hit, reach 5 ft., one target. \textsl{Hit:} 22 (4d6 + 8) piercing damage. The target must succeed on a DC 21 Constitution saving throw or become poisoned. While poisoned in this way, the target can't regain hit points, and it takes 21 (6d6) poison damage at the start of each of its turns. The poisoned target can repeat the saving throw at the end of each of its turns, ending the effect on itself on a success.

\DndMonsterAction{Claw}
\textsl{Melee Weapon Attack:} 14 to hit, reach 10 ft., one target. \textsl{Hit:} 17 (2d8 + 8) slashing damage.

\DndMonsterAction{Mace}
\textsl{Melee Weapon Attack:} 14 to hit, reach 10 ft., one target. \textsl{Hit:} 15 (2d6 + 8) bludgeoning damage plus 21 (6d6) fire damage.

\DndMonsterAction{Tail}
\textsl{Melee Weapon Attack:} 14 to hit, reach 10 ft., one target. \textsl{Hit:} 24 (3d10 + 8) bludgeoning damage.

\end{DndMonster}



\begin{DndMonster}[float=t]{Planetar}
\DndMonsterType{Large celestial, lawful good}
\DndMonsterBasics[
armor-class = {19 (natural armor)},
hit-points = {\DndDice{16d10 + 112}},
speed = {40 ft., fly 120 ft.}
]
\DndMonsterAbilityScores[
 str = 24,
 dex = 20,
 con = 24,
 int = 19,
 wis = 22,
 cha = 25
]
\DndMonsterDetails[
saving-throws = {Con +12, Wis +11, Cha +12},
skills = {Perception +11},
damage-resistances = {radiant; bludgeoning, piercing, slashing from nonmagical attacks},
condition-immunities = {charmed, exhaustion, frightened},
senses = {truesight 120 ft., passive perception 21},
languages = {all, telepathy 120 ft.},
challenge = 16,
]
\DndMonsterAction{Angelic Weapons}
The planetar's weapon attacks are magical. When the planetar hits with any weapon, the weapon deals an extra 5d8 radiant damage (included in the attack).

\DndMonsterAction{Divine Awareness}
The planetar knows if it hears a lie.

\DndMonsterAction{Magic Resistance}
The planetar has advantage on saving throws against spells and other magical effects.

\DndMonsterAction{Innate Spellcasting}
The planetar's spellcasting ability is Charisma (spell save DC 20). The planetar can innately cast the following spells, requiring no material components:
\begin{DndMonsterSpells}
  \DndInnateSpellLevel{\textit{detect evil and good}, \textit{invisibility} (self only)}
  \DndInnateSpellLevel[3e]{\textit{blade barrier}, \textit{dispel evil and good}, \textit{flame strike}, \textit{raise dead}}
  \DndInnateSpellLevel[1e]{\textit{commune}, \textit{control weather}, \textit{insect plague}}
\end{DndMonsterSpells}
\DndMonsterSection{Actions}
\DndMonsterAction{Multiattack}
The planetar makes two melee attacks.

\DndMonsterAction{Greatsword}
\textsl{Melee Weapon Attack:} 12 to hit, reach 5 ft., one target. \textsl{Hit:} 21 (4d6 + 7) slashing damage plus 22 (5d8) radiant damage.

\DndMonsterAction{Healing Touch (4/Day)}
The planetar touches another creature. The target magically regains 30 (6d8 + 3) hit points and is freed from any curse, disease, poison, blindness, or deafness.

\end{DndMonster}


\clearpage

\begin{DndMonster}[float=t]{Polar Bear}
\DndMonsterType{Large beast, unaligned}
\DndMonsterBasics[
armor-class = {12 (natural armor)},
hit-points = {\DndDice{5d10 + 15}},
speed = {40 ft., swim 30 ft.}
]
\DndMonsterAbilityScores[
 str = 20,
 dex = 10,
 con = 16,
 int = 2,
 wis = 13,
 cha = 7
]
\DndMonsterDetails[
skills = {Perception +3},
challenge = 2,
]
\DndMonsterAction{Keen Smell}
The bear has advantage on Wisdom (Perception) checks that rely on smell.

\DndMonsterSection{Actions}
\DndMonsterAction{Multiattack}
The bear makes two attacks: one with its bite and one with its claws.

\DndMonsterAction{Bite}
\textsl{Melee Weapon Attack:} 7 to hit, reach 5 ft., one target. \textsl{Hit:} 9 (1d8 + 5) piercing damage.

\DndMonsterAction{Claws}
\textsl{Melee Weapon Attack:} 7 to hit, reach 5 ft., one target. \textsl{Hit:} 12 (2d6 + 5) slashing damage.

\end{DndMonster}



\begin{DndMonster}[float=t]{Priest}
\DndMonsterType{Medium humanoid (any race), any alignment}
\DndMonsterBasics[
armor-class = {13 (\textit{chain shirt})},
hit-points = {\DndDice{5d8 + 5}},
speed = {30 ft.}
]
\DndMonsterAbilityScores[
 str = 10,
 dex = 10,
 con = 12,
 int = 13,
 wis = 16,
 cha = 13
]
\DndMonsterDetails[
skills = {Medicine +7, Persuasion +3, Religion +5},
languages = {any two languages},
challenge = 2,
]
\DndMonsterAction{Divine Eminence}
As a bonus action, the priest can expend a spell slot to cause its melee weapon attacks to magically deal an extra 10 (3d6) radiant damage to a target on a hit. This benefit lasts until the end of the turn. If the priest expends a spell slot of 2nd level or higher, the extra damage increases by 1d6 for each level above 1st.

\DndMonsterAction{Spellcasting}
The priest is a 5th-level spellcaster. Its spellcasting ability is Wisdom (spell save DC 13, 5 to hit with spell attacks). The priest has the following cleric spells prepared:
\begin{DndMonsterSpells}
  \DndMonsterSpellLevel{\textit{light}, \textit{sacred flame}, \textit{thaumaturgy}}
   \DndMonsterSpellLevel[1][4]{\textit{cure wounds}, \textit{guiding bolt}, \textit{sanctuary}}
   \DndMonsterSpellLevel[2][3]{\textit{lesser restoration}, \textit{spiritual weapon}}
   \DndMonsterSpellLevel[3][2]{\textit{dispel magic}, \textit{spirit guardians}}
\end{DndMonsterSpells}
\DndMonsterSection{Actions}
\DndMonsterAction{Mace}
\textsl{Melee Weapon Attack:} 2 to hit, reach 5 ft., one target. \textsl{Hit:} 3 (1d6) bludgeoning damage.

\end{DndMonster}



\begin{DndMonster}[float=t]{Purple Worm}
\DndMonsterType{Gargantuan monstrosity, unaligned}
\DndMonsterBasics[
armor-class = {18 (natural armor)},
hit-points = {\DndDice{15d20 + 90}},
speed = {50 ft., burrow 30 ft.}
]
\DndMonsterAbilityScores[
 str = 28,
 dex = 7,
 con = 22,
 int = 1,
 wis = 8,
 cha = 4
]
\DndMonsterDetails[
saving-throws = {Con +11, Wis +4},
senses = {blindsight 30 ft., tremorsense 60 ft., passive perception 9},
challenge = 15,
]
\DndMonsterAction{Tunneler}
The worm can burrow through solid rock at half its burrow speed and leaves a 10-foot-diameter tunnel in its wake.

\DndMonsterSection{Actions}
\DndMonsterAction{Multiattack}
The worm makes two attacks: one with its bite and one with its stinger.

\DndMonsterAction{Bite}
\textsl{Melee Weapon Attack:} 14 to hit, reach 10 ft., one target. \textsl{Hit:} 22 (3d8 + 9) piercing damage. If the target is a Large or smaller creature, it must succeed on a DC 19 Dexterity saving throw or be swallowed by the worm. A swallowed creature is blinded and restrained, it has total cover against attacks and other effects outside the worm, and it takes 21 (6d6) acid damage at the start of each of the worm's turns.

If the worm takes 30 damage or more on a single turn from a creature inside it, the worm must succeed on a DC 21 Constitution saving throw at the end of that turn or regurgitate all swallowed creatures, which fall prone in a space within 10 feet of the worm. If the worm dies, a swallowed creature is no longer restrained by it and can escape from the corpse by using 20 feet of movement, exiting prone.

\DndMonsterAction{Tail Stinger}
\textsl{Melee Weapon Attack:} 14 to hit, reach 10 ft., one creature. \textsl{Hit:} 19 (3d6 + 9) piercing damage, and the target must make a DC 19 Constitution saving throw, taking 42 (12d6) poison damage on a failed save, or half as much damage on a successful one.

\end{DndMonster}



\begin{DndMonster}[float=t]{Quadrone}
\DndMonsterType{Medium construct, lawful neutral}
\DndMonsterBasics[
armor-class = {16 (natural armor)},
hit-points = {\DndDice{4d8 + 4}},
speed = {30 ft., fly 30 ft.}
]
\DndMonsterAbilityScores[
 str = 12,
 dex = 14,
 con = 12,
 int = 10,
 wis = 10,
 cha = 11
]
\DndMonsterDetails[
skills = {Perception +2},
senses = {truesight 120 ft., passive perception 12},
languages = {Modron},
challenge = 1,
]
\DndMonsterAction{Axiomatic Mind}
The quadrone can't be compelled to act in a manner contrary to its nature or its instructions.

\DndMonsterAction{Disintegration}
If the quadrone dies, its body disintegrates into dust, leaving behind its weapons and anything else it was carrying.

\DndMonsterSection{Actions}
\DndMonsterAction{Multiattack}
The quadrone makes two fist attacks or four shortbow attacks.

\DndMonsterAction{Fist}
\textsl{Melee Weapon Attack:} 3 to hit, reach 5 ft., one target. \textsl{Hit:} 3 (1d4 + 1) bludgeoning damage.

\DndMonsterAction{Shortbow}
\textsl{Ranged Weapon Attack:} 4 to hit, range 80/320 ft., one target. \textsl{Hit:} 5 (1d6 + 2) piercing damage.

\end{DndMonster}



\begin{DndMonster}[float=t]{Ramius}
\DndMonsterType{Medium humanoid, neutral good}
\DndMonsterBasics[
armor-class = {22 (plate armor)},
hit-points = {\DndDice{20d8 + 100}},
speed = {30 ft.}
]
\DndMonsterAbilityScores[
 str = 23,
 dex = 10,
 con = 21,
 int = 11,
 wis = 14,
 cha = 18
]
\DndMonsterDetails[
saving-throws = {Wis +7, Cha +9},
skills = {Animal Handling +12, Athletics +11, Insight +7, Religion +10},
damage-resistances = {bludgeoning, piercing, slashing from nonmagical attacks; damage from spells},
languages = {Abyssal, Celestial, Common, Infernal},
challenge = 14,
]
\DndMonsterAction{Aura of Protection}
Ramius and any ally that starts their turn within 30 feet of him have a +4 bonus to all saving throws and resistance to nonmagical damage.

\DndMonsterAction{Spellcasting}
Ramius casts one of the following spells, using Charisma as the spellcasting ability (spell save DC 17):
\begin{DndMonsterSpells}
  \DndInnateSpellLevel[1e]{\textit{Aura of Purity}, \textit{Compelled Duel}, \textit{Crusader's Mantle}, \textit{Death Ward}, \textit{Dispel Evil and Good}, \textit{Find Steed}, \textit{Lesser Restoration}, \textit{Remove Curse}, \textit{Revivify}, \textit{Shield of Faith}}
\end{DndMonsterSpells}
\DndMonsterSection{Actions}
\DndMonsterAction{Multiattack}
Ramius makes two Longsword attacks.

\DndMonsterAction{Longsword}
\textsl{Melee Weapon Attack:} 11 to hit, reach 5ft., one target. \textsl{Hit:} 10 (1d8 + 6) slashing damage or 11 (1d10 + 6) if wielded in two hands, plus 13 (3d8) radiant damage. If the target is an evil creature, it must succeed on a DC 17 Charisma saving throw or take an additional 13 (3d8) radiant damage.

\DndMonsterAction{Divine Wave (Recharge 6)}
Ramius strikes the ground and causes holy energy to radiate outwards. Each creature of his choice within 30 feet of him must make a DC 17 Constitution saving throw. Targets take 21 (6d6) thunder damage plus 21 (6d6) radiant damage on a failed save, or half as much damage on a successful one. Creatures that fail the save are knocked down (have the prone condition).

\DndMonsterSection{Reactions}
\DndMonsterAction{Healing Touch}
As a reaction to Ramius or an adjacent creature taking 50 points or more of damage, Ramius reaches out and channels divine energy, healing 13 (2d8 + 4) hit points.

\end{DndMonster}



\begin{DndMonster}[float=t]{Rat}
\DndMonsterType{Tiny beast, unaligned}
\DndMonsterBasics[
armor-class = {10},
hit-points = {\DndDice{1d4 - 1}},
speed = {20 ft.}
]
\DndMonsterAbilityScores[
 str = 2,
 dex = 11,
 con = 9,
 int = 2,
 wis = 10,
 cha = 4
]
\DndMonsterDetails[
senses = {darkvision 30 ft., passive perception 10},
challenge = 0,
]
\DndMonsterAction{Keen Smell}
The rat has advantage on Wisdom (Perception) checks that rely on smell.

\DndMonsterSection{Actions}
\DndMonsterAction{Bite}
\textsl{Melee Weapon Attack:} 0 to hit, reach 5 ft., one target. \textsl{Hit:} 1 piercing damage.

\end{DndMonster}



\begin{DndMonster}[float=t]{Red Abishai}
\DndMonsterType{Medium fiend (devil), lawful evil}
\DndMonsterBasics[
armor-class = {22 (natural armor)},
hit-points = {\DndDice{34d8 + 136}},
speed = {30 ft., fly 50 ft.}
]
\DndMonsterAbilityScores[
 str = 23,
 dex = 16,
 con = 19,
 int = 14,
 wis = 15,
 cha = 19
]
\DndMonsterDetails[
saving-throws = {Str +12, Con +10, Wis +8},
skills = {Intimidation +10, Perception +8},
damage-resistances = {cold; bludgeoning, piercing, slashing from nonmagical attacks that aren't silvered},
damage-immunities = {fire, poison},
condition-immunities = {frightened, poisoned},
senses = {darkvision 120 ft., passive perception 18},
languages = {Draconic, Infernal, telepathy 120 ft.},
challenge = 19,
]
\DndMonsterAction{Devil's Sight}
Magical darkness doesn't impede the abishai's darkvision.

\DndMonsterAction{Magic Resistance}
The abishai has advantage on saving throws against spells and other magical effects.

\DndMonsterSection{Actions}
\DndMonsterAction{Multiattack}
The abishai makes one Bite attack and one Claw attack, and it can use Frightful Presence or Incite Fanaticism.

\DndMonsterAction{Bite}
\textsl{Melee Weapon Attack:} 12 to hit, reach 5 ft., one target. \textsl{Hit:} 22 (3d10 + 6) piercing damage plus 38 (7d10) fire damage.

\DndMonsterAction{Claw}
\textsl{Melee Weapon Attack:} 12 to hit, reach 5 ft., one target. \textsl{Hit:} 17 (2d10 + 6) force damage plus 11 (2d10) fire damage.

\DndMonsterAction{Frightful Presence}
Each creature of the abishai's choice that is within 120 feet and aware of the abishai must succeed on a DC 18 Wisdom saving throw or become frightened of it for 1 minute. A creature can repeat the saving throw at the end of each of its turns, ending the effect on itself on a success. If a creature's saving throw is successful or the effect ends for it, the creature is immune to the abishai's Frightful Presence for the next 24 hours.

\DndMonsterAction{Incite Fanaticism}
The abishai chooses up to four other creatures within 60 feet of it that can see it. Until the start of the abishai's next turn, each of those creatures makes attack rolls with advantage and can't be frightened.

\DndMonsterAction{Power of the Dragon Queen}
The abishai targets one Dragon it can see within 120 feet of it. The Dragon must make a DC 18 Charisma saving throw. A chromatic dragon makes this save with disadvantage. On a successful save, the target is immune to the abishai's Power of the Dragon Queen for 1 hour. On a failed save, the target is charmed by the abishai for 1 hour. While charmed in this way, the target regards the abishai as a trusted friend to be heeded and protected. This effect ends if the abishai or its companions deal damage to the target.

\end{DndMonster}



\begin{DndMonster}[float=t]{Regular Orthoclath}
\DndMonsterType{Large construct, unaligned}
\DndMonsterBasics[
armor-class = {20 (natural armor)},
hit-points = {\DndDice{20d10 + 100}},
speed = {30 ft.}
]
\DndMonsterAbilityScores[
 str = 24,
 dex = 9,
 con = 20,
 int = 3,
 wis = 11,
 cha = 1
]
\DndMonsterDetails[
damage-immunities = {fire; poison; psychic; bludgeoning, piercing, slashing from nonmagical attacks that aren't adamantine},
condition-immunities = {charmed, exhaustion, frightened, paralyzed, petrified, poisoned},
senses = {darkvision 120 ft., passive perception 10},
languages = {understands the languages of its creator but can't speak},
challenge = 16,
]
\DndMonsterAction{Copy Warning}
This content has been copied from another creature, but the data replacement hasn't been run. Check details.

\DndMonsterAction{Fire Absorption}
Whenever the golem is subjected to fire damage, it takes no damage and instead regains a number of hit points equal to the fire damage dealt.

\DndMonsterAction{Immutable Form}
The golem is immune to any spell or effect that would alter its form.

\DndMonsterAction{Magic Resistance}
The golem has advantage on saving throws against spells and other magical effects.

\DndMonsterAction{Magic Weapons}
The golem's weapon attacks are magical.

\DndMonsterSection{Actions}
\DndMonsterAction{Multiattack}
The golem makes two melee attacks.

\DndMonsterAction{Slam}
\textsl{Melee Weapon Attack:} 13 to hit, reach 5 ft., one target. \textsl{Hit:} 20 (3d8 + 7) bludgeoning damage.

\DndMonsterAction{Sword}
\textsl{Melee Weapon Attack:} 13 to hit, reach 10 ft., one target. \textsl{Hit:} 23 (3d10 + 7) slashing damage.

\DndMonsterAction{Poison Breath (Recharge 5\textendash 6)}
The golem exhales poisonous gas in a 15-foot cone. Each creature in that area must make a DC 19 Constitution saving throw, taking 45 (10d8) poison damage on a failed save, or half as much damage on a successful one.

\end{DndMonster}



\begin{DndMonster}[float=t]{Remorhaz}
\DndMonsterType{Huge monstrosity, unaligned}
\DndMonsterBasics[
armor-class = {17 (natural armor)},
hit-points = {\DndDice{17d12 + 85}},
speed = {30 ft., burrow 20 ft.}
]
\DndMonsterAbilityScores[
 str = 24,
 dex = 13,
 con = 21,
 int = 4,
 wis = 10,
 cha = 5
]
\DndMonsterDetails[
damage-immunities = {cold, fire},
senses = {darkvision 60 ft., tremorsense 60 ft., passive perception 10},
challenge = 11,
]
\DndMonsterAction{Heated Body}
A creature that touches the remorhaz or hits it with a melee attack while within 5 feet of it takes 10 (3d6) fire damage.

\DndMonsterSection{Actions}
\DndMonsterAction{Bite}
\textsl{Melee Weapon Attack:} 11 to hit, reach 10 ft., one target. \textsl{Hit:} 40 (6d10 + 7) piercing damage plus 10 (3d6) fire damage. If the target is a creature, it is grappled (escape DC 17). Until this grapple ends, the target is restrained, and the remorhaz can't bite another target.

\DndMonsterAction{Swallow}
The remorhaz makes one bite attack against a Medium or smaller creature it is grappling. If the attack hits, that creature takes the bite's damage and is swallowed, and the grapple ends. While swallowed, the creature is blinded and restrained, it has total cover against attacks and other effects outside the remorhaz, and it takes 21 (6d6) acid damage at the start of each of the remorhaz's turns.

If the remorhaz takes 30 damage or more on a single turn from a creature inside it, the remorhaz must succeed on a DC 15 Constitution saving throw at the end of that turn or regurgitate all swallowed creatures, which fall prone in a space within 10 feet of the remorhaz. If the remorhaz dies, a swallowed creature is no longer restrained by it and can escape from the corpse using 15 feet of movement, exiting prone.

\end{DndMonster}



\begin{DndMonster}[float=t]{Rhinoceros}
\DndMonsterType{Large beast, unaligned}
\DndMonsterBasics[
armor-class = {11 (natural armor)},
hit-points = {\DndDice{6d10 + 12}},
speed = {40 ft.}
]
\DndMonsterAbilityScores[
 str = 21,
 dex = 8,
 con = 15,
 int = 2,
 wis = 12,
 cha = 6
]
\DndMonsterDetails[
challenge = 2,
]
\DndMonsterAction{Charge}
If the rhinoceros moves at least 20 feet straight toward a target and then hits it with a gore attack on the same turn, the target takes an extra 9 (2d8) bludgeoning damage. If the target is a creature, it must succeed on a DC 15 Strength saving throw or be knocked prone.

\DndMonsterSection{Actions}
\DndMonsterAction{Gore}
\textsl{Melee Weapon Attack:} 7 to hit, reach 5 ft., one target. \textsl{Hit:} 14 (2d8 + 5) bludgeoning damage.

\end{DndMonster}



\begin{DndMonster}[float=t]{Rimmon}
\DndMonsterType{Medium fiend (devil), lawful evil}
\DndMonsterBasics[
armor-class = {17 (natural armor)},
hit-points = {\DndDice{25d8 + 150}},
speed = {30 ft.}
]
\DndMonsterAbilityScores[
 str = 18,
 dex = 17,
 con = 22,
 int = 27,
 wis = 25,
 cha = 16
]
\DndMonsterDetails[
saving-throws = {Con +12, Int +14, Wis +13},
skills = {Arcana +20, History +20, Investigation +14, Perception +13, Sleight Of Hand +9},
damage-resistances = {cold; bludgeoning, piercing, slashing from nonmagical attacks that aren't silvered},
damage-immunities = {fire, poison},
condition-immunities = {charmed, poisoned},
senses = {darkvision 120 ft., passive perception 23},
languages = {Celestial, Common, Draconic, Infernal, telepathy 120 ft.},
challenge = 20,
]
\DndMonsterAction{Contractually Obligated}
Rimmon has advantage to hit creatures under an infernal contract and when he hits such creatures he delivers an additional 9 (2d8) damage of a type the target is most vulnerable to. A creature bound under an infernal contract also makes all saving throws against effects originating from Rimmon with disadvantage.

\DndMonsterAction{Devil's Sight}
Magical darkness doesn't impede Rimmon's darkvision.

\DndMonsterAction{Magic Resistance}
Rimmon has advantage on saving throws against spells and other magical effects.

\DndMonsterSection{Actions}
\DndMonsterAction{Multiattack}
Rimmon makes four Sharpened Tongue or Wordplay attacks.

\DndMonsterAction{Sharpened Tongue}
\textsl{Melee Weapon Attack:} 10 to hit, reach 10 ft., one target. \textsl{Hit:} 11 (3d4 + 4) slashing damage.

\DndMonsterAction{Wordplay}
\textsl{Ranged Spell Attack:} 14 to hit, range 60 ft., one target. \textsl{Hit:} 35 (6d8 + 8) psychic damage.

\DndMonsterAction{Connected Clauses (1/Day)}
Rimmon emits a burst of energy that arcs towards a creature he can see within 150 feet of him. Three bolts then leap from that target to as many as three other targets, each of which must be within 30 feet of the first target. A target must make a DC 22 Dexterity saving throw, taking 63 (14d8) damage on a failed save, or half as much damage on a successful one. The damage type inflicted is of a type the first target is most vulnerable to.

\end{DndMonster}



\begin{DndMonster}[float=t]{Rothé}
\DndMonsterType{Large beast (cattle), unaligned}
\DndMonsterBasics[
armor-class = {10},
hit-points = {\DndDice{2d10 + 4}},
speed = {30 ft.}
]
\DndMonsterAbilityScores[
 str = 18,
 dex = 10,
 con = 14,
 int = 2,
 wis = 10,
 cha = 4
]
\DndMonsterDetails[
senses = {darkvision 30 ft., passive perception 10},
challenge = 1/4,
]
\DndMonsterSection{Actions}
\DndMonsterAction{Gore}
\textsl{Melee Weapon Attack:} 6 to hit, reach 5 ft., one target. \textsl{Hit:} 7 (1d6 + 4) piercing damage. If the rothé moved at least 20 feet straight towards this target immediately before the hit, the target takes an extra 7 (2d6) piercing damage.

\end{DndMonster}



\begin{DndMonster}[float=t]{Salamander}
\DndMonsterType{Large elemental, neutral evil}
\DndMonsterBasics[
armor-class = {15 (natural armor)},
hit-points = {\DndDice{12d10 + 24}},
speed = {30 ft.}
]
\DndMonsterAbilityScores[
 str = 18,
 dex = 14,
 con = 15,
 int = 11,
 wis = 10,
 cha = 12
]
\DndMonsterDetails[
damage-vulnerabilities = {cold},
damage-resistances = {bludgeoning, piercing, slashing from nonmagical attacks},
damage-immunities = {fire},
senses = {darkvision 60 ft., passive perception 10},
languages = {Ignan},
challenge = 5,
]
\DndMonsterAction{Heated Body}
A creature that touches the salamander or hits it with a melee attack while within 5 feet of it takes 7 (2d6) fire damage.

\DndMonsterAction{Heated Weapons}
Any metal melee weapon the salamander wields deals an extra 3 (1d6) fire damage on a hit (included in the attack).

\DndMonsterSection{Actions}
\DndMonsterAction{Multiattack}
The salamander makes two attacks: one with its spear and one with its tail.

\DndMonsterAction{Spear}
\textsl{Melee or Ranged Weapon Attack:} 7 to hit, reach 5 ft. or range 20/60 ft., one target. \textsl{Hit:} 11 (2d6 + 4) piercing damage, or 13 (2d8 + 4) piercing damage if used with two hands to make a melee attack, plus 3 (1d6) fire damage.

\DndMonsterAction{Tail}
\textsl{Melee Weapon Attack:} 7 to hit, reach 10 ft., one target. \textsl{Hit:} 11 (2d6 + 4) bludgeoning damage plus 7 (2d6) fire damage, and the target is grappled (escape DC 14). Until this grapple ends, the target is restrained, the salamander can automatically hit the target with its tail, and the salamander can't make tail attacks against other targets.

\end{DndMonster}



\begin{DndMonster}[float*=b,width=\textwidth + 8pt]{Sarevok}
\begin{multicols}{2}
\DndMonsterType{Medium humanoid, neutral evil}
\DndMonsterBasics[
armor-class = {21 (plate armor of Bhaal)},
hit-points = {\DndDice{20d8 + 100}},
speed = {30 ft.}
]
\DndMonsterAbilityScores[
 str = 23,
 dex = 10,
 con = 21,
 int = 11,
 wis = 18,
 cha = 14
]
\DndMonsterDetails[
saving-throws = {Con +11, Wis +10},
skills = {History +6, Intimidation +14, Religion +6},
damage-immunities = {acid, necrotic, poison},
condition-immunities = {charmed, frightened, poisoned},
senses = {darkvision 60 ft., passive perception 14},
languages = {Abyssal, Common, Infernal},
challenge = 17,
]
\DndMonsterAction{Magic Resistance}
Sarevok has advantage on saving throws against spells and other magical effects.

\DndMonsterAction{Rejuvenation}
If Sarevok is killed, he gains a new body in 24 hours, regaining all his hit points. The new body appears on the altar of the temple of Bhaal beneath Baldur's Gate. This ability ceases to function if a cleric of good alignment casts \textit{Hallow} on the altar in the temple of Bhaal.

\DndMonsterAction{Spellcasting}
Sarevok casts one of the following spells, using Wisdom as the spellcasting ability (spell save DC 18):
\begin{DndMonsterSpells}
  \DndInnateSpellLevel{\textit{Resistance}, \textit{Thaumaturgy}}
  \DndInnateSpellLevel[1e]{\textit{Animate Dead}, \textit{Antimagic Field}, \textit{Astral Projection}, \textit{Blade Barrier}, \textit{Command}, \textit{Contagion}, \textit{Dispel Magic}, \textit{Divination}, \textit{Etherealness}, \textit{Fire Storm}, \textit{Harm}, \textit{Hold Person}, \textit{Silence}}
\end{DndMonsterSpells}
\DndMonsterSection{Actions}
\DndMonsterAction{Multiattack}
Sarevok makes three Longsword attacks. He can replace one of the attacks with Flames of Bhaal or Spellcasting.

\DndMonsterAction{Longsword}
\textsl{Melee Weapon Attack:} 12 to hit, reach 5ft., one target. \textsl{Hit:} 10 (1d8 + 6) slashing damage or 11 (1d10 + 6) if wielded in two hands. The target must succeed on a DC 18 Constitution saving throw or be cursed by Bhaal, preventing it from regaining hit points. The curse lasts until removed by the \textit{Remove Curse} spell or similar effect.

\DndMonsterAction{Flames of Bhaal}
Sarevok causes flames to engulf one creature that he can see within 60 feet. The target must succeed on a DC 18 Dexterity saving throw or take 18 (4d8) necrotic damage. This damage can't be restored except by a \textit{Lesser Restoration} spell or similar effect.

\DndMonsterSection{Legendary Actions}
Sarevok can take 3 legendary actions, choosing from the options below. Only one legendary action can be used at a time and only at the end of another creature's turn. Sarevok regains spent legendary actions at the start of its turn.

\begin{DndMonsterLegendaryActions}
\DndMonsterLegendaryAction{Slash}{Sarevok makes a Longsword attack.
}
\DndMonsterLegendaryAction{Unholy Flame}{Sarevok uses Flames of Bhaal.
}
\DndMonsterLegendaryAction{Channel Bhaal's Hate (Costs 2 Actions)}{Sarevok unleashes Bhaal's power. Creatures within 30 feet of Sarevok, including ones behind barriers and around corners, can't regain hit points until the end of Sarevok's next turn.
}
\end{DndMonsterLegendaryActions}
\end{multicols}
\end{DndMonster}



\begin{DndMonster}[float=t]{Scout}
\DndMonsterType{Medium humanoid (any race), any alignment}
\DndMonsterBasics[
armor-class = {13 (\textit{leather armor})},
hit-points = {\DndDice{3d8 + 3}},
speed = {30 ft.}
]
\DndMonsterAbilityScores[
 str = 11,
 dex = 14,
 con = 12,
 int = 11,
 wis = 13,
 cha = 11
]
\DndMonsterDetails[
skills = {Nature +4, Perception +5, Stealth +6, Survival +5},
languages = {any one language (usually Common)},
challenge = 1/2,
]
\DndMonsterAction{Keen Hearing and Sight}
The scout has advantage on Wisdom (Perception) checks that rely on hearing or sight.

\DndMonsterSection{Actions}
\DndMonsterAction{Multiattack}
The scout makes two melee attacks or two ranged attacks.

\DndMonsterAction{Shortsword}
\textsl{Melee Weapon Attack:} 4 to hit, reach 5 ft., one target. \textsl{Hit:} 5 (1d6 + 2) piercing damage.

\DndMonsterAction{Longbow}
\textsl{Ranged Weapon Attack:} 4 to hit, ranged 150/600 ft., one target. \textsl{Hit:} 6 (1d8 + 2) piercing damage.

\end{DndMonster}



\begin{DndMonster}[float*=b,width=\textwidth + 8pt]{Sea Hag}
\begin{multicols}{2}
\DndMonsterType{Medium fey, chaotic evil}
\DndMonsterBasics[
armor-class = {14 (natural armor)},
hit-points = {\DndDice{7d8 + 21}},
speed = {30 ft., swim 40 ft.}
]
\DndMonsterAbilityScores[
 str = 16,
 dex = 13,
 con = 16,
 int = 12,
 wis = 12,
 cha = 13
]
\DndMonsterDetails[
senses = {darkvision 60 ft., passive perception 11},
languages = {Aquan, Common, Giant},
challenge = 2,
]
\DndMonsterAction{Amphibious}
The hag can breathe air and water.

\DndMonsterAction{Horrific Appearance}
Any humanoid that starts its turn within 30 feet of the hag and can see the hag's true form must make a DC 11 Wisdom saving throw. On a failed save, the creature is frightened for 1 minute. A creature can repeat the saving throw at the end of each of its turns, with disadvantage if the hag is within line of sight, ending the effect on itself on a success. If a creature's saving throw is successful or the effect ends for it, the creature is immune to the hag's Horrific Appearance for the next 24 hours.

Unless the target is Surprise or the revelation of the hag's true form is sudden, the target can avert its eyes and avoid making the initial saving throw. Until the start of its next turn, a creature that averts its eyes has disadvantage on attack rolls against the hag.

\DndMonsterSection{Actions}
\DndMonsterAction{Claws}
\textsl{Melee Weapon Attack:} 5 to hit, reach 5 ft., one target. \textsl{Hit:} 10 (2d6 + 3) slashing damage.

\DndMonsterAction{Death Glare}
The hag targets one frightened creature she can see within 30 feet of her. If the target can see the hag, it must succeed on a DC 11 Wisdom saving throw against this magic or drop to 0 hit points.

\DndMonsterAction{Illusory Appearance}
The hag covers herself and anything she is wearing or carrying with a magical illusion that makes her look like an ugly creature of her general size and humanoid shape. The effect ends if the hag takes a bonus action to end it or if she dies.

The changes wrought by this effect fail to hold up to physical inspection. For example, the hag could appear to have no claws, but someone touching her hand might feel the claws. Otherwise, a creature must take an action to visually inspect the illusion and succeed on a DC 16 Intelligence (Investigation) check to discern that the hag is disguised.

\end{multicols}
\end{DndMonster}



\begin{DndMonster}[float=t]{Shadow Demon}
\DndMonsterType{Medium fiend (demon), chaotic evil}
\DndMonsterBasics[
armor-class = {13},
hit-points = {\DndDice{12d8 + 12}},
speed = {30 ft., fly 30 ft.}
]
\DndMonsterAbilityScores[
 str = 1,
 dex = 17,
 con = 12,
 int = 14,
 wis = 13,
 cha = 14
]
\DndMonsterDetails[
saving-throws = {Dex +5, Cha +4},
skills = {Stealth +7},
damage-vulnerabilities = {radiant},
damage-resistances = {acid; fire; necrotic; thunder; bludgeoning, piercing, slashing from nonmagical attacks},
damage-immunities = {cold, lightning, poison},
condition-immunities = {exhaustion, grappled, paralyzed, petrified, poisoned, prone, restrained},
senses = {darkvision 120 ft., passive perception 11},
languages = {Abyssal, telepathy 120 ft.},
challenge = 4,
]
\DndMonsterAction{Incorporeal Movement}
The demon can move through other creatures and objects as if they were difficult terrain. It takes 5 (1d10) force damage if it ends its turn inside an object.

\DndMonsterAction{Light Sensitivity}
While in bright light, the demon has disadvantage on attack rolls, as well as on Wisdom (Perception) checks that rely on sight.

\DndMonsterAction{Shadow Stealth}
While in dim light or darkness, the demon can take the Hide action as a bonus action.

\DndMonsterSection{Actions}
\DndMonsterAction{Claws}
\textsl{Melee Weapon Attack:} 5 to hit, reach 5 ft., one creature. \textsl{Hit:} 10 (2d6 + 3) psychic damage or, if the demon had advantage on the attack roll, 17 (4d6 + 3) psychic damage.

\end{DndMonster}



\begin{DndMonster}[float=t]{Shield Guardian}
\DndMonsterType{Large construct, unaligned}
\DndMonsterBasics[
armor-class = {17 (natural armor)},
hit-points = {\DndDice{15d10 + 60}},
speed = {30 ft.}
]
\DndMonsterAbilityScores[
 str = 18,
 dex = 8,
 con = 18,
 int = 7,
 wis = 10,
 cha = 3
]
\DndMonsterDetails[
damage-immunities = {poison},
condition-immunities = {charmed, exhaustion, frightened, paralyzed, poisoned},
senses = {blindsight 10 ft., darkvision 60 ft., passive perception 10},
languages = {understands commands given in any language but can't speak},
challenge = 7,
]
\DndMonsterAction{Bound}
The shield guardian is magically bound to an amulet. As long as the guardian and its amulet are on the same plane of existence, the amulet's wearer can telepathically call the guardian to travel to it, and the guardian knows the distance and direction to the amulet. If the guardian is within 60 feet of the amulet's wearer, half of any damage the wearer takes (rounded up) is transferred to the guardian.

\DndMonsterAction{Regeneration}
The shield guardian regains 10 hit points at the start of its turn if it has at least 1 hit point.

\DndMonsterAction{Spell Storing}
A spellcaster who wears the shield guardian's amulet can cause the guardian to store one spell of 4th level or lower. To do so, the wearer must cast the spell on the guardian. The spell has no effect but is stored within the guardian. When commanded to do so by the wearer or when a situation arises that was predefined by the spellcaster, the guardian casts the stored spell with any parameters set by the original caster, requiring no components. When the spell is cast or a new spell is stored, any previously stored spell is lost.

\DndMonsterSection{Actions}
\DndMonsterAction{Multiattack}
The guardian makes two fist attacks.

\DndMonsterAction{Fist}
\textsl{Melee Weapon Attack:} 7 to hit, reach 5 ft., one target. \textsl{Hit:} 11 (2d6 + 4) bludgeoning damage.

\DndMonsterSection{Reactions}
\DndMonsterAction{Shield}
When a creature makes an attack against the wearer of the guardian's amulet, the guardian grants a +2 bonus to the wearer's AC if the guardian is within 5 feet of the wearer.

\end{DndMonster}



\begin{DndMonster}[float=t]{Shredwing}
\DndMonsterType{Medium fiend, lawful evil}
\DndMonsterBasics[
armor-class = {21 (natural armor)},
hit-points = {\DndDice{20d8 + 80}},
speed = {0 ft., fly 60 ft. \textit{(hover)}}
]
\DndMonsterAbilityScores[
 str = 11,
 dex = 19,
 con = 19,
 int = 15,
 wis = 16,
 cha = 11
]
\DndMonsterDetails[
saving-throws = {Dex +8, Con +8},
skills = {Acrobatics +8, Athletics +4, Perception +11, Survival +7},
damage-resistances = {cold; fire; poison; bludgeoning, piercing, slashing from nonmagical attacks that aren't silvered},
condition-immunities = {blinded, charmed, poisoned},
senses = {darkvision 120 ft., passive perception 21},
challenge = 12,
]
\DndMonsterAction{Merged}
After using the Burrow ability, if the shredwing is merged with another creature, any damage dealt to the shredwing is split evenly between it and the merged creature.

\DndMonsterSection{Actions}
\DndMonsterAction{Multiattack}
The shredwing makes three Talon attacks.

\DndMonsterAction{Talon}
\textsl{Melee Weapon Attack:} 8 to hit, reach 5 ft., one target. \textsl{Hit:} 11 (2d6 + 4) slashing damage plus 7 (2d6) poison damage.

\DndMonsterAction{Burrow (Recharge 5\textendash 6)}
On a hit, the creature takes an additional 32 (5d12) slashing damage as the wings merge with it. While merged, the creature moves with the shredwing, takes 19 (3d12) necrotic damage at the start of each of its turn, and has the grappled condition. The shredwing can unmerge at any time, leaving the creature and dropping it if airborne. The only other way to end the merge is to kill the shredwing.

If the merged creature dies while merged, the shredwing permanently takes over its body. It gains the merged creature's hit dice, natural armor, possessions, and languages. If any of the merged creature's statistics are higher than the shredwing, it inherits those statistics as well. A creature permanently merged by the shredwing can only be resurrected through a \textit{Wish} spell or similar.

\DndMonsterSection{Bonus Actions}
\DndMonsterAction{Aggressive}
The shredwing can move up to its speed toward a hostile creature that it can see.

\end{DndMonster}



\begin{DndMonster}[float=t]{Shrieker}
\DndMonsterType{Medium plant, unaligned}
\DndMonsterBasics[
armor-class = {5},
hit-points = {\DndDice{3d8}},
speed = {0 ft.}
]
\DndMonsterAbilityScores[
 str = 1,
 dex = 1,
 con = 10,
 int = 1,
 wis = 3,
 cha = 1
]
\DndMonsterDetails[
condition-immunities = {blinded, deafened, frightened},
senses = {blindsight 30 ft. (blind beyond this radius), passive perception 6},
challenge = 0,
]
\DndMonsterAction{False Appearance}
While the shrieker remains motionless, it is indistinguishable from an ordinary fungus.

\DndMonsterSection{Reactions}
\DndMonsterAction{Shriek}
When bright light or a creature is within 30 feet of the shrieker, it emits a shriek audible within 300 feet of it. The shrieker continues to shriek until the disturbance moves out of range and for 1d4 of the shrieker's turns afterward.

\end{DndMonster}


\clearpage

\begin{DndMonster}[float=t]{Spined Devil}
\DndMonsterType{Small fiend (devil), lawful evil}
\DndMonsterBasics[
armor-class = {13 (natural armor)},
hit-points = {\DndDice{5d6 + 5}},
speed = {20 ft., fly 40 ft.}
]
\DndMonsterAbilityScores[
 str = 10,
 dex = 15,
 con = 12,
 int = 11,
 wis = 14,
 cha = 8
]
\DndMonsterDetails[
damage-resistances = {cold; bludgeoning, piercing, slashing from nonmagical attacks that aren't silvered},
damage-immunities = {fire, poison},
condition-immunities = {poisoned},
senses = {darkvision 120 ft., passive perception 12},
languages = {Infernal, telepathy 120 ft.},
challenge = 2,
]
\DndMonsterAction{Devil's Sight}
Magical darkness doesn't impede the devil's darkvision.

\DndMonsterAction{Flyby}
The devil doesn't provoke an opportunity attack when it flies out of an enemy's reach.

\DndMonsterAction{Limited Spines}
The devil has twelve tail spines. Used spines regrow by the time the devil finishes a long rest.

\DndMonsterAction{Magic Resistance}
The devil has advantage on saving throws against spells and other magical effects.

\DndMonsterSection{Actions}
\DndMonsterAction{Multiattack}
The devil makes two attacks: one with its bite and one with its fork or two with its tail spines.

\DndMonsterAction{Bite}
\textsl{Melee Weapon Attack:} 2 to hit, reach 5 ft., one target. \textsl{Hit:} 5 (2d4) slashing damage.

\DndMonsterAction{Fork}
\textsl{Melee Weapon Attack:} 2 to hit, reach 5 ft., one target. \textsl{Hit:} 3 (1d6) piercing damage.

\DndMonsterAction{Tail Spine}
\textsl{Ranged Weapon Attack:} 4 to hit, range 20/80 ft., one target. \textsl{Hit:} 4 (1d4 + 2) piercing damage plus 3 (1d6) fire damage.

\end{DndMonster}



\begin{DndMonster}[float=t]{Spy}
\DndMonsterType{Medium humanoid (any race), any alignment}
\DndMonsterBasics[
armor-class = {12},
hit-points = {\DndDice{6d8}},
speed = {30 ft.}
]
\DndMonsterAbilityScores[
 str = 10,
 dex = 15,
 con = 10,
 int = 12,
 wis = 14,
 cha = 16
]
\DndMonsterDetails[
skills = {Deception +5, Insight +4, Investigation +5, Perception +6, Persuasion +5, Sleight Of Hand +4, Stealth +4},
languages = {any two languages},
challenge = 1,
]
\DndMonsterAction{Cunning Action}
On each of its turns, the spy can use a bonus action to take the Dash, Disengage, or Hide action.

\DndMonsterAction{Sneak Attack (1/Turn)}
The spy deals an extra 7 (2d6) damage when it hits a target with a weapon attack and has advantage on the attack roll, or when the target is within 5 feet of an ally of the spy that isn't incapacitated and the spy doesn't have disadvantage on the attack roll.

\DndMonsterSection{Actions}
\DndMonsterAction{Multiattack}
The spy makes two melee attacks.

\DndMonsterAction{Shortsword}
\textsl{Melee Weapon Attack:} 4 to hit, reach 5 ft., one target. \textsl{Hit:} 5 (1d6 + 2) piercing damage.

\DndMonsterAction{Hand Crossbow}
\textsl{Ranged Weapon Attack:} 4 to hit, range 30/120 ft., one target. \textsl{Hit:} 5 (1d6 + 2) piercing damage.

\end{DndMonster}



\begin{DndMonster}[float=t]{Stone Giant}
\DndMonsterType{Huge giant, neutral}
\DndMonsterBasics[
armor-class = {17 (natural armor)},
hit-points = {\DndDice{11d12 + 55}},
speed = {40 ft.}
]
\DndMonsterAbilityScores[
 str = 23,
 dex = 15,
 con = 20,
 int = 10,
 wis = 12,
 cha = 9
]
\DndMonsterDetails[
saving-throws = {Dex +5, Con +8, Wis +4},
skills = {Athletics +12, Perception +4},
senses = {darkvision 60 ft., passive perception 14},
languages = {Giant},
challenge = 7,
]
\DndMonsterAction{Stone Camouflage}
The giant has advantage on Dexterity (Stealth) checks made to hide in rocky terrain.

\DndMonsterSection{Actions}
\DndMonsterAction{Multiattack}
The giant makes two greatclub attacks.

\DndMonsterAction{Greatclub}
\textsl{Melee Weapon Attack:} 9 to hit, reach 15 ft., one target. \textsl{Hit:} 19 (3d8 + 6) bludgeoning damage.

\DndMonsterAction{Rock}
\textsl{Ranged Weapon Attack:} 9 to hit, range 60/240 ft., one target. \textsl{Hit:} 28 (4d10 + 6) bludgeoning damage. If the target is a creature, it must succeed on a DC 17 Strength saving throw or be knocked prone.

\DndMonsterSection{Reactions}
\DndMonsterAction{Rock Catching}
If a rock or similar object is hurled at the giant, the giant can, with a successful DC 10 Dexterity saving throw, catch the missile and take no bludgeoning damage from it.

\end{DndMonster}



\begin{DndMonster}[float=t]{Stone Golem}
\DndMonsterType{Large construct, unaligned}
\DndMonsterBasics[
armor-class = {17 (natural armor)},
hit-points = {\DndDice{17d10 + 85}},
speed = {30 ft.}
]
\DndMonsterAbilityScores[
 str = 22,
 dex = 9,
 con = 20,
 int = 3,
 wis = 11,
 cha = 1
]
\DndMonsterDetails[
damage-immunities = {poison; psychic; bludgeoning, piercing, slashing from nonmagical attacks that aren't adamantine},
condition-immunities = {charmed, exhaustion, frightened, paralyzed, petrified, poisoned},
senses = {darkvision 120 ft., passive perception 10},
languages = {understands the languages of its creator but can't speak},
challenge = 10,
]
\DndMonsterAction{Immutable Form}
The golem is immune to any spell or effect that would alter its form.

\DndMonsterAction{Magic Resistance}
The golem has advantage on saving throws against spells and other magical effects.

\DndMonsterAction{Magic Weapons}
The golem's weapon attacks are magical.

\DndMonsterSection{Actions}
\DndMonsterAction{Multiattack}
The golem makes two slam attacks.

\DndMonsterAction{Slam}
\textsl{Melee Weapon Attack:} 10 to hit, reach 5 ft., one target. \textsl{Hit:} 19 (3d8 + 6) bludgeoning damage.

\DndMonsterAction{Slow (Recharge 5\textendash 6)}
The golem targets one or more creatures it can see within 10 feet of it. Each target must make a DC 17 Wisdom saving throw against this magic. On a failed save, a target can't use reactions, its speed is halved, and it can't make more than one attack on its turn. In addition, the target can take either an action or a bonus action on its turn, not both. These effects last for 1 minute. A target can repeat the saving throw at the end of each of its turns, ending the effect on itself on a success.

\end{DndMonster}



\begin{DndMonster}[float=t]{Storm Giant}
\DndMonsterType{Huge giant, chaotic good}
\DndMonsterBasics[
armor-class = {16 (\textit{scale mail})},
hit-points = {\DndDice{20d12 + 100}},
speed = {50 ft., swim 50 ft.}
]
\DndMonsterAbilityScores[
 str = 29,
 dex = 14,
 con = 20,
 int = 16,
 wis = 18,
 cha = 18
]
\DndMonsterDetails[
saving-throws = {Str +14, Con +10, Wis +9, Cha +9},
skills = {Arcana +8, Athletics +14, History +8, Perception +9},
damage-resistances = {cold},
damage-immunities = {lightning, thunder},
languages = {Common, Giant},
challenge = 13,
]
\DndMonsterAction{Amphibious}
The giant can breathe air and water.

\DndMonsterAction{Innate Spellcasting}
The giant's innate spellcasting ability is Charisma (spell save DC 17). It can innately cast the following spells, requiring no material components:
\begin{DndMonsterSpells}
  \DndInnateSpellLevel{\textit{detect magic}, \textit{feather fall}, \textit{levitate}, \textit{light}}
  \DndInnateSpellLevel[3e]{\textit{control weather}, \textit{water breathing}}
\end{DndMonsterSpells}
\DndMonsterSection{Actions}
\DndMonsterAction{Multiattack}
The giant makes two greatsword attacks.

\DndMonsterAction{Greatsword}
\textsl{Melee Weapon Attack:} 14 to hit, reach 10 ft., one target. \textsl{Hit:} 30 (6d6 + 9) slashing damage.

\DndMonsterAction{Rock}
\textsl{Ranged Weapon Attack:} 14 to hit, range 60/240 ft., one target. \textsl{Hit:} 35 (4d12 + 9) bludgeoning damage.

\DndMonsterAction{Lightning Strike (Recharge 5\textendash 6)}
The giant hurls a magical lightning bolt at a point it can see within 500 feet of it. Each creature within 10 feet of that point must make a DC 17 Dexterity saving throw, taking 54 (12d8) lightning damage on a failed save, or half as much damage on a successful one.

\end{DndMonster}



\begin{DndMonster}[float*=b,width=\textwidth + 8pt]{Styx Dragon}
\begin{multicols}{2}
\DndMonsterType{Gargantuan dragon, chaotic evil}
\DndMonsterBasics[
armor-class = {17 (natural armor)},
hit-points = {\DndDice{16d20 + 128}},
speed = {30 ft., burrow 20 ft., swim 60 ft.}
]
\DndMonsterAbilityScores[
 str = 27,
 dex = 14,
 con = 26,
 int = 19,
 wis = 20,
 cha = 19
]
\DndMonsterDetails[
saving-throws = {Dex +8, Con +14, Wis +11, Cha +10},
skills = {Athletics +21, Insight +11, Perception +11, Survival +17},
damage-resistances = {damage from spells},
damage-immunities = {cold, poison},
condition-immunities = {blinded, poisoned},
senses = {blindsight 120 ft., passive perception 21},
languages = {Draconic, Infernal},
challenge = 20,
]
\DndMonsterAction{Amphibious}
The dragon can breathe air and water.

\DndMonsterAction{Stygian Wasting}
Creatures afflicted with this disease are incapable of regaining hit points in any way until the disease is cured. Additionally, as their flesh starts to rot, a creature with this disease takes 36 (8d8) necrotic damage after each long rest they finish.

\DndMonsterAction{Styx Dweller}
The dragon is immune to all negative effects from the River Styx.

\DndMonsterSection{Actions}
\DndMonsterAction{Multiattack}
The dragon can use its Frightful Presence. It then makes three attacks using its Bite, Tail, or a combination of the two. It can replace one of the attacks with Constrict.

\DndMonsterAction{Bite}
\textsl{Melee Weapon Attack:} 14 to hit, reach 5 ft., one target. \textsl{Hit:} 19 (2d10 + 8) piercing damage plus 7 (2d6) poison damage, and the target must make a DC 19 Constitution saving throw. On a failed save, the target contracts Stygian Wasting.

\DndMonsterAction{Tail}
\textsl{Melee Weapon Attack:} 14 to hit, reach 10 ft., one target. \textsl{Hit:} 15 (2d6 + 8) bludgeoning damage. If the target is a Large or smaller creature, it has the grappled condition (escape DC 22). The dragon can't make a Tail attack while a creature is grappled in this way.

\DndMonsterAction{Constrict}
The dragon tightens its grip on a creature grappled by its tail. The creature takes 27 (3d12 + 8) bludgeoning damage, has the restrained condition until the start of the dragon's next turn, and must make a DC 19 Constitution saving throw. On a failed save, the target contracts Stygian Wasting.

\DndMonsterAction{Frightful Presence}
Each creature of the dragon's choice that is within 120 feet of the dragon must succeed on a DC 18 Wisdom saving throw or have the frightened condition for 1 minute. A creature can repeat the saving throw at the end of each of its turns, ending the effect on itself on a success. If a creature's saving throw is successful or the effect ends for it, the creature is immune to the dragon's Frightful Presence for the next 24 hours.

\DndMonsterAction{Stygian Breath (Recharge 5\textendash 6)}
The dragon exhales poisonous Styx water in a 60-foot-long, 10-foot-wide line. Each creature in that line must make a DC 19 Dexterity saving throw, taking 54 (12d8) poison damage on a failed save, or half as much on a successful one. Creatures that fail the save contract Stygian Wasting.

\DndMonsterSection{Legendary Actions}
Styx Dragon can take 3 legendary actions, choosing from the options below. Only one legendary action can be used at a time and only at the end of another creature's turn. Styx Dragon regains spent legendary actions at the start of its turn.

\begin{DndMonsterLegendaryActions}
\DndMonsterLegendaryAction{Detect}{The dragon makes a Wisdom (Perception) check.
}
\DndMonsterLegendaryAction{Tail Attack}{The dragon makes a Tail attack. If a creature is already grappled, the dragon may use Constrict instead.
}
\DndMonsterLegendaryAction{Death Throw (Costs 3 Actions)}{If the dragon has successfully restrained a creature with the Constrict ability, it may throw them up to 100 feet away. When it lands, the creature takes 21 (6d6) bludgeoning damage and must make a DC 18 Constitution saving throw. On a failed save, the creature has the unconscious condition for 10 minutes.
}
\end{DndMonsterLegendaryActions}
\end{multicols}
\end{DndMonster}



\begin{DndMonster}[float=t]{Succubus}
\DndMonsterType{Medium fiend (shapechanger), neutral evil}
\DndMonsterBasics[
armor-class = {15 (natural armor)},
hit-points = {\DndDice{12d8 + 12}},
speed = {30 ft., fly 60 ft.}
]
\DndMonsterAbilityScores[
 str = 8,
 dex = 17,
 con = 13,
 int = 15,
 wis = 12,
 cha = 20
]
\DndMonsterDetails[
skills = {Deception +9, Insight +5, Perception +5, Persuasion +9, Stealth +7},
damage-resistances = {cold; fire; lightning; poison; bludgeoning, piercing, slashing from nonmagical attacks},
senses = {darkvision 60 ft., passive perception 15},
languages = {Abyssal, Common, Infernal, telepathy 60 ft.},
challenge = 4,
]
\DndMonsterAction{Telepathic Bond}
The fiend ignores the range restriction on its telepathy when communicating with a creature it has charmed. The two don't even need to be on the same plane of existence.

\DndMonsterAction{Shapechanger}
The fiend can use its action to polymorph into a Small or Medium humanoid, or back into its true form. Without wings, the fiend loses its flying speed. Other than its size and speed, its statistics are the same in each form. Any equipment it is wearing or carrying isn't transformed. It reverts to its true form if it dies.

\DndMonsterSection{Actions}
\DndMonsterAction{Claw (Fiend Form Only)}
\textsl{Melee Weapon Attack:} 5 to hit, reach 5 ft., one target. \textsl{Hit:} 6 (1d6 + 3) slashing damage.

\DndMonsterAction{Charm}
One humanoid the fiend can see within 30 feet of it must succeed on a DC 15 Wisdom saving throw or be magically charmed for 1 day. The charmed target obeys the fiend's verbal or telepathic commands. If the target suffers any harm or receives a suicidal command, it can repeat the saving throw, ending the effect on a success. If the target successfully saves against the effect, or if the effect on it ends, the target is immune to this fiend's Charm for the next 24 hours.

The fiend can have only one target charmed at a time. If it charms another, the effect on the previous target ends.

\DndMonsterAction{Draining Kiss}
The fiend kisses a creature charmed by it or a willing creature. The target must make a DC 15 Constitution saving throw against this magic, taking 32 (5d10 + 5) psychic damage on a failed save, or half as much damage on a successful one. The target's hit point maximum is reduced by an amount equal to the damage taken. This reduction lasts until the target finishes a long rest. The target dies if this effect reduces its hit point maximum to 0.

\DndMonsterAction{Etherealness}
The fiend magically enters the Ethereal Plane from the Material Plane, or vice versa.

\end{DndMonster}



\begin{DndMonster}[float=t]{Thug}
\DndMonsterType{Medium humanoid (any race), any non-good alignment}
\DndMonsterBasics[
armor-class = {11 (\textit{leather armor})},
hit-points = {\DndDice{5d8 + 10}},
speed = {30 ft.}
]
\DndMonsterAbilityScores[
 str = 15,
 dex = 11,
 con = 14,
 int = 10,
 wis = 10,
 cha = 11
]
\DndMonsterDetails[
skills = {Intimidation +2},
languages = {any one language (usually Common)},
challenge = 1/2,
]
\DndMonsterAction{Pack Tactics}
The thug has advantage on an attack roll against a creature if at least one of the thug's allies is within 5 feet of the creature and the ally isn't incapacitated.

\DndMonsterSection{Actions}
\DndMonsterAction{Multiattack}
The thug makes two melee attacks.

\DndMonsterAction{Mace}
\textsl{Melee Weapon Attack:} 4 to hit, reach 5 ft., one creature. \textsl{Hit:} 5 (1d6 + 2) bludgeoning damage.

\DndMonsterAction{Heavy Crossbow}
\textsl{Ranged Weapon Attack:} 2 to hit, range 100/400 ft., one target. \textsl{Hit:} 5 (1d10) piercing damage.

\end{DndMonster}



\begin{DndMonster}[float=t]{Tiax}
\DndMonsterType{Small humanoid, chaotic neutral}
\DndMonsterBasics[
armor-class = {13 (16 with \textit{mage armor})},
hit-points = {\DndDice{20d6 + 40}},
speed = {25 ft.}
]
\DndMonsterAbilityScores[
 str = 11,
 dex = 16,
 con = 14,
 int = 15,
 wis = 19,
 cha = 18
]
\DndMonsterDetails[
saving-throws = {Wis +8, Cha +8},
skills = {Deception +8, Insight +8, Performance +8, Persuasion +8},
senses = {darkvision 60 ft., passive perception 14},
languages = {Celestial, Common, Gnomish},
challenge = 9,
]
\DndMonsterAction{Gnome Cunning}
Tiax has advantage on Intelligence, Wisdom, and Charisma saving throws against magic.

\DndMonsterAction{Special Equipment}
Tiax wears a \textit{Ring of Resistance} that grants him resistance to radiant damage.

\DndMonsterAction{Spellcasting}
Tiax casts one of the following spells, using Wisdom as the spellcasting ability (spell save DC 16):
\begin{DndMonsterSpells}
  \DndInnateSpellLevel{\textit{Bane}, \textit{Disguise Self}, \textit{Dispel Magic}, \textit{Mage Armor}, \textit{Mirror Image}, \textit{Thaumaturgy}}
  \DndInnateSpellLevel[1e]{\textit{Bestow Curse}, \textit{Blindness/Deafness}, \textit{Blink}, \textit{Charm Person}, \textit{Contagion}, \textit{Dimension Door}, \textit{Dominate Person}, \textit{Hold Person}, \textit{Pass Without Trace}, \textit{Polymorph}, \textit{Teleport}}
\end{DndMonsterSpells}
\DndMonsterSection{Actions}
\DndMonsterAction{Touch of Decay}
\textsl{Melee Spell Attack:} 8 to hit, reach 5 ft., one target. \textsl{Hit:} 37 (6d10 + 4) necrotic damage.

\DndMonsterAction{Unholy Firebolt}
Tiax causes unholy fire to flare up in a creature's space within 60 feet of him. The target must succeed on a DC 16 Dexterity saving throw or take 18 (4d8) necrotic damage plus 18 (4d8) poison damage.

\DndMonsterSection{Bonus Actions}
\DndMonsterAction{Heal (2/Day)}
Tiax heals 70 hit points.

\end{DndMonster}



\begin{DndMonster}[float=t]{Triceratops}
\DndMonsterType{Huge beast, unaligned}
\DndMonsterBasics[
armor-class = {13 (natural armor)},
hit-points = {\DndDice{10d12 + 30}},
speed = {50 ft.}
]
\DndMonsterAbilityScores[
 str = 22,
 dex = 9,
 con = 17,
 int = 2,
 wis = 11,
 cha = 5
]
\DndMonsterDetails[
challenge = 5,
]
\DndMonsterAction{Trampling Charge}
If the triceratops moves at least 20 feet straight toward a creature and then hits it with a gore attack on the same turn, that target must succeed on a DC 13 Strength saving throw or be knocked prone. If the target is prone, the triceratops can make one stomp attack against it as a bonus action.

\DndMonsterSection{Actions}
\DndMonsterAction{Gore}
\textsl{Melee Weapon Attack:} 9 to hit, reach 5 ft., one target. \textsl{Hit:} 24 (4d8 + 6) piercing damage.

\DndMonsterAction{Stomp}
\textsl{Melee Weapon Attack:} 9 to hit, reach 5 ft., one prone creature. \textsl{Hit:} 22 (3d10 + 6) bludgeoning damage

\end{DndMonster}



\begin{DndMonster}[float=t]{Tridrone}
\DndMonsterType{Medium construct, lawful neutral}
\DndMonsterBasics[
armor-class = {15 (natural armor)},
hit-points = {\DndDice{3d8 + 3}},
speed = {30 ft.}
]
\DndMonsterAbilityScores[
 str = 12,
 dex = 13,
 con = 12,
 int = 9,
 wis = 10,
 cha = 9
]
\DndMonsterDetails[
senses = {truesight 120 ft., passive perception 10},
languages = {Modron},
challenge = 1/2,
]
\DndMonsterAction{Axiomatic Mind}
The tridrone can't be compelled to act in a manner contrary to its nature or its instructions.

\DndMonsterAction{Disintegration}
If the tridrone dies, its body disintegrates into dust, leaving behind its weapons and anything else it was carrying.

\DndMonsterSection{Actions}
\DndMonsterAction{Multiattack}
The tridrone makes three fist attacks or three javelin attacks.

\DndMonsterAction{Fist}
\textsl{Melee Weapon Attack:} 3 to hit, reach 5 ft., one target. \textsl{Hit:} 3 (1d4 + 1) bludgeoning damage.

\DndMonsterAction{Javelin}
\textsl{Melee or Ranged Weapon Attack:} 3 to hit, reach 5 ft. or range 30/120 ft., one target. \textsl{Hit:} 4 (1d6 + 1) piercing damage.

\end{DndMonster}



\begin{DndMonster}[float=t]{Troll}
\DndMonsterType{Large giant, chaotic evil}
\DndMonsterBasics[
armor-class = {15 (natural armor)},
hit-points = {\DndDice{8d10 + 40}},
speed = {30 ft.}
]
\DndMonsterAbilityScores[
 str = 18,
 dex = 13,
 con = 20,
 int = 7,
 wis = 9,
 cha = 7
]
\DndMonsterDetails[
skills = {Perception +2},
senses = {darkvision 60 ft., passive perception 12},
languages = {Giant},
challenge = 5,
]
\DndMonsterAction{Keen Smell}
The troll has advantage on Wisdom (Perception) checks that rely on smell.

\DndMonsterAction{Regeneration}
The troll regains 10 hit points at the start of its turn. If the troll takes acid or fire damage, this trait doesn't function at the start of the troll's next turn. The troll dies only if it starts its turn with 0 hit points and doesn't regenerate.

\DndMonsterSection{Actions}
\DndMonsterAction{Multiattack}
The troll makes three attacks: one with its bite and two with its claws.

\DndMonsterAction{Bite}
\textsl{Melee Weapon Attack:} 7 to hit, reach 5 ft., one target. \textsl{Hit:} 7 (1d6 + 4) piercing damage.

\DndMonsterAction{Claw}
\textsl{Melee Weapon Attack:} 7 to hit, reach 5 ft., one target. \textsl{Hit:} 11 (2d6 + 4) slashing damage.

\end{DndMonster}



\begin{DndMonster}[float*=b,width=\textwidth + 8pt]{Tyrant Shadow}
\begin{multicols}{2}
\DndMonsterType{Huge aberration, lawful evil}
\DndMonsterBasics[
armor-class = {19 (natural armor)},
hit-points = {\DndDice{18d12 + 90}},
speed = {30 ft.}
]
\DndMonsterAbilityScores[
 str = 22,
 dex = 16,
 con = 20,
 int = 20,
 wis = 16,
 cha = 20
]
\DndMonsterDetails[
saving-throws = {Int +11, Wis +9},
skills = {Arcana +11, Deception +17, Perception +9, Stealth +15},
damage-resistances = {fire},
damage-immunities = {necrotic; poison; psychic; bludgeoning, piercing, slashing from nonmagical attacks},
condition-immunities = {blinded, charmed, frightened, poisoned},
languages = {understands all, telepathy 300 ft.},
challenge = 17,
]
\DndMonsterAction{Empath}
The shadow uses telepathy to sense nearby creatures. The shadow can see a creature if that creature has thoughts and is within the range of the shadow's telepathy. The shadow can't see creatures immune to telepathy.

\DndMonsterAction{Undying Connection}
The shadow was originally spawned from a specific devil and, if the shadow is slain, it returns to life in 10 (1d20) days within 1 mile of that devil. The shadow is killed permanently only if its progenitor devil is slain first.

\DndMonsterSection{Actions}
\DndMonsterAction{Multiattack}
The shadow makes three Claw attacks. It can replace one of the attacks with Shadows of Death (if available).

\DndMonsterAction{Claw}
\textsl{Melee Weapon Attack:} 12 to hit, reach 5 ft., one target. \textsl{Hit:} 16 (3d6 + 6) necrotic damage, plus 7 (2d6) psychic damage.

\DndMonsterAction{Cloak of Shadows}
The shadow bends darkness around itself, and now has the invisible condition for 1 minute. The invisibility ends if the shadow is subjected to a \textit{Daylight} spell or similar.

\DndMonsterAction{Empathic Link}
The shadow attempts to form a link between itself and a creature it can see. The target must make a DC 19 Intelligence saving throw, taking 13 (2d12) psychic damage on a success. On a failed save, the creature and the shadow are linked. While linked, every time the shadow is damaged, the creature takes half of that damage as psychic damage. The creature may repeat the saving throw at the end of each of its turns, taking 13 (2d12) psychic damage on a failed save, or ending the effect on a successful one.

\DndMonsterAction{Shadows of Death (Recharge 6)}
The shadow discharges inky-black shadows in a 90-foot-radius sphere around itself. Creatures in the shadows must succeed on a DC 19 Wisdom saving throw or take 27 (6d8) necrotic damage and have the frightened condition for 1 minute. A frightened creature may repeat the saving throw at the end of each of its turns, ending the effect on a succesful save, but taking the damage again on a failed save.

\DndMonsterSection{Bonus Actions}
\DndMonsterAction{Change Shape}
The shadow transforms into a Beast, Humanoid or Fiend that it has an Empathic Link with, or back into its true form. In a new form, the shadow retains its align ment, hit points, hit dice, telepathy, and Intelligence, Wisdom, and Charisma scores, as well as this action. Its statistics and capabilities are otherwise replaced by those of the new form, except any class features, lair actions, or legendary actions of that form.

\DndMonsterSection{Legendary Actions}
Tyrant Shadow can take 3 legendary actions, choosing from the options below. Only one legendary action can be used at a time and only at the end of another creature's turn. Tyrant Shadow regains spent legendary actions at the start of its turn.

\begin{DndMonsterLegendaryActions}
\DndMonsterLegendaryAction{Shadow Blink}{The shadow surrounds itself in black mist, then teleports up to 60 feet to an unoccupied space that it can see.
}
\DndMonsterLegendaryAction{Shapeshift}{The shadow uses its Change Shape ability.
}
\DndMonsterLegendaryAction{Mind Spike (Costs 2 Actions)}{\textsl{Ranged Spell Attack:} 11 to hit, range 300 ft., one target. \textsl{Hit:} 16 (2d10 + 5) psychic damage.
}
\end{DndMonsterLegendaryActions}
\end{multicols}
\end{DndMonster}



\begin{DndMonster}[float=t]{Umber Hulk}
\DndMonsterType{Large monstrosity, chaotic evil}
\DndMonsterBasics[
armor-class = {18 (natural armor)},
hit-points = {\DndDice{11d10 + 33}},
speed = {30 ft., burrow 20 ft.}
]
\DndMonsterAbilityScores[
 str = 20,
 dex = 13,
 con = 16,
 int = 9,
 wis = 10,
 cha = 10
]
\DndMonsterDetails[
senses = {darkvision 120 ft., tremorsense 60 ft., passive perception 10},
languages = {Umber Hulk},
challenge = 5,
]
\DndMonsterAction{Confusing Gaze}
When a creature starts its turn within 30 feet of the umber hulk and is able to see the umber hulk's eyes, the umber hulk can magically force it to make a DC 15 Charisma saving throw, unless the umber hulk is incapacitated.

On a failed saving throw, the creature can't take reactions until the start of its next turn and rolls a d8 to determine what it does during that turn. On a 1 to 4, the creature does nothing. On a 5 or 6, the creature takes no action but uses all its movement to move in a random direction. On a 7 or 8, the creature makes one melee attack against a random creature, or it does nothing if no creature is within reach.

Unless Surprise, a creature can avert its eyes to avoid the saving throw at the start of its turn. If the creature does so, it can't see the umber hulk until the start of its next turn, when it can avert its eyes again. If the creature looks at the umber hulk in the meantime, it must immediately make the save.

\DndMonsterAction{Tunneler}
The umber hulk can burrow through solid rock at half its burrowing speed and leaves a 5 foot-wide, 8-foot-high tunnel in its wake.

\DndMonsterSection{Actions}
\DndMonsterAction{Multiattack}
The umber hulk makes three attacks: two with its claws and one with its mandibles.

\DndMonsterAction{Claw}
\textsl{Melee Weapon Attack:} 8 to hit, reach 5 ft., one target. \textsl{Hit:} 9 (1d8 + 5) slashing damage.

\DndMonsterAction{Mandibles}
\textsl{Melee Weapon Attack:} 8 to hit, reach 5 ft., one target. \textsl{Hit:} 14 (2d8 + 5) slashing damage.

\end{DndMonster}



\begin{DndMonster}[float=t]{Veteran}
\DndMonsterType{Medium humanoid (any race), any alignment}
\DndMonsterBasics[
armor-class = {17 (\textit{splint armor})},
hit-points = {\DndDice{9d8 + 18}},
speed = {30 ft.}
]
\DndMonsterAbilityScores[
 str = 16,
 dex = 13,
 con = 14,
 int = 10,
 wis = 11,
 cha = 10
]
\DndMonsterDetails[
skills = {Athletics +5, Perception +2},
languages = {any one language (usually Common)},
challenge = 3,
]
\DndMonsterSection{Actions}
\DndMonsterAction{Multiattack}
The veteran makes two longsword attacks. If it has a shortsword drawn, it can also make a shortsword attack.

\DndMonsterAction{Longsword}
\textsl{Melee Weapon Attack:} 5 to hit, reach 5 ft., one target. \textsl{Hit:} 7 (1d8 + 3) slashing damage, or 8 (1d10 + 3) slashing damage if used with two hands.

\DndMonsterAction{Shortsword}
\textsl{Melee Weapon Attack:} 5 to hit, reach 5 ft., one target. \textsl{Hit:} 6 (1d6 + 3) piercing damage.

\DndMonsterAction{Heavy Crossbow}
\textsl{Ranged Weapon Attack:} 3 to hit, range 100/400 ft., one target. \textsl{Hit:} 6 (1d10 + 1) piercing damage.

\end{DndMonster}



\begin{DndMonster}[float=t]{Vorvolaka}
\DndMonsterType{Medium undead, neutral evil}
\DndMonsterBasics[
armor-class = {18 (natural armor)},
hit-points = {\DndDice{20d8 + 100}},
speed = {20 ft., fly 50 ft.}
]
\DndMonsterAbilityScores[
 str = 15,
 dex = 19,
 con = 21,
 int = 8,
 wis = 13,
 cha = 16
]
\DndMonsterDetails[
saving-throws = {Con +10, Wis +6},
skills = {Insight +6, Perception +11, Religion +4, Survival +6},
damage-resistances = {bludgeoning, piercing, slashing from nonmagical attacks},
damage-immunities = {necrotic},
condition-immunities = {charmed, exhaustion, paralyzed},
senses = {darkvision 120 ft., passive perception 21},
languages = {the languages it knew in life},
challenge = 14,
]
\DndMonsterAction{Fear Aura}
Any hostile creature that starts its turn within 20 feet of the vorvolaka must succeed on a DC 18 Wisdom saving throw or have frightened condition until the start of the creature's next turn. If a creature's saving throw is successful, they're immune to the vorvolaka's Fear Aura for the next 24 hours. This ability does not function if the vorvolaka is unconscious.

\DndMonsterAction{Magic Resistance}
The vorvolaka has advantage on saving throws against spells and other magical effects.

\DndMonsterAction{Regeneration}
The vorvolaka regains 15 hit points at the start of its turn if it has at least 1 hit point. If the vorvolaka takes radiant damage, this trait doesn't function at the start of the vorvolaka's next turn.

\DndMonsterSection{Actions}
\DndMonsterAction{Multiattack}
The vorvolaka makes two Talon attacks. It can replace one of the attacks with Rend Armor (if available).

\DndMonsterAction{Talon}
\textsl{Melee Weapon Attack:} 9 to hit, reach 5 ft., one target. \textsl{Hit:} 22 (4d8 + 4) piercing damage plus 14 (4d6) necrotic damage.

\DndMonsterAction{Rend Armor (Recharge 4\textendash 6)}
The vorvolaka makes a Talon attack against a creature. On a hit, the creature's AC is reduced by 2 for 1 minute. Creatures not wearing armor are immune to this effect, and the total reduction can't bring a creature below an AC of 10.

\end{DndMonster}



\begin{DndMonster}[float*=b,width=\textwidth + 8pt]{Waeloquay}
\begin{multicols}{2}
\DndMonsterType{Gargantuan elemental, neutral evil}
\DndMonsterBasics[
armor-class = {14 (natural armor)},
hit-points = {\DndDice{20d20 + 80}},
speed = {30 ft., swim 90 ft.}
]
\DndMonsterAbilityScores[
 str = 26,
 dex = 14,
 con = 18,
 int = 8,
 wis = 10,
 cha = 12
]
\DndMonsterDetails[
saving-throws = {Wis +6, Cha +7},
skills = {Athletics +14, Nature +5},
damage-resistances = {acid; bludgeoning, piercing, slashing from nonmagical attacks},
damage-immunities = {poison},
condition-immunities = {exhaustion, grappled, paralyzed, petrified, poisoned, prone, restrained, unconscious},
senses = {darkvision 60 ft., passive perception 10},
languages = {Aquan, Common},
challenge = 18,
]
\DndMonsterAction{Magic Resistance}
Waeloquay has advantage on saving throws against spells and other magical effects.

\DndMonsterAction{Legendary Resistance (3/Day)}
If Waeloquay fails a saving throw, he can choose to succeed instead.

\DndMonsterAction{Water Form}
Waeloquay can enter a hostile creature's space and stop there. It can move through a space as narrow as 1 inch wide without squeezing.

\DndMonsterSection{Actions}
\DndMonsterAction{Multiattack}
Waeloquay makes three Slam attacks.

\DndMonsterAction{Slam}
\textsl{Melee Weapon Attack:} 14 to hit, reach 15 ft., one target. \textsl{Hit:} 26 (4d8 + 8) bludgeoning damage.

\DndMonsterAction{Whelm (Recharge 4\textendash 6)}
Each creature in Waeloquay's space must make a DC 22 Strength saving throw. On a failed save, a target takes 21 (3d8 + 8) bludgeoning damage. If it is Large or smaller, it also has the grappled condition (escape DC 20). Until this grapple ends, the target is restrained and unable to breathe, unless it can breathe water. If the save is successful, the target is pushed out of Waeloquay's space. Waeloquay can grapple one Large creature or up to four Medium or smaller creatures at one time. At the start of each of Waeloquay's turns, each target grappled by it takes 21 (3d8 + 8) bludgeoning damage. A creature within 5 feet of Waeloquay can pull a creature or object out of it by taking an action to make a DC 22 Strength check and succeeding.

\DndMonsterSection{Legendary Actions}
Waeloquay can take 3 legendary actions, choosing from the options below. Only one legendary action can be used at a time and only at the end of another creature's turn. Waeloquay regains spent legendary actions at the start of its turn.

\begin{DndMonsterLegendaryActions}
\DndMonsterLegendaryAction{Flow}{Waeloquay moves up to its speed.
}
\DndMonsterLegendaryAction{Crush (Costs 2 Actions)}{One creature that Waeloquay is grappling takes 44 (8d8 + 8) bludgeoning damage.
}
\DndMonsterLegendaryAction{Pummel (Costs 2 Actions)}{Waeloquay makes a Slam attack with advantage.
}
\end{DndMonsterLegendaryActions}
\end{multicols}
\end{DndMonster}



\begin{DndMonster}[float=t]{War Devil}
\DndMonsterType{Huge fiend (devil), lawful evil}
\DndMonsterBasics[
armor-class = {11 (natural armor)},
hit-points = {\DndDice{19d12 + 114}},
speed = {40 ft., fly 60 ft.}
]
\DndMonsterAbilityScores[
 str = 24,
 dex = 9,
 con = 23,
 int = 7,
 wis = 13,
 cha = 9
]
\DndMonsterDetails[
saving-throws = {Str +12, Con +11},
skills = {Athletics +12, Insight +6, Intimidation +4, Perception +6},
damage-resistances = {cold},
damage-immunities = {fire; poison; bludgeoning, piercing, slashing from nonmagical attacks that aren't silvered},
condition-immunities = {frightened, poisoned},
senses = {darkvision 120 ft., passive perception 16},
languages = {Infernal, telepathy 120 ft.},
challenge = 16,
]
\DndMonsterAction{Devil's Sight}
Magical darkness doesn't impede the war devil's darkvision.

\DndMonsterAction{Fear Aura}
Any hostile creature that starts its turn within 20 feet of the war devil must succeed on a DC 18 Wisdom saving throw or have the frightened condition until the start of the creature's next turn. If a creature's saving throw is successful, the creature is immune to the war devil's Fear Aura for the next 24 hours. This ability doesn't function if the war devil is unconscious.

\DndMonsterAction{Magic Resistance}
The war devil has advantage on saving throws against spells and other magical effects.

\DndMonsterSection{Actions}
\DndMonsterAction{Multiattack}
The war devil makes two attacks using its Cold- Iron Ranseur, Claws, or a combination of the two.

\DndMonsterAction{Cold-Iron Ranseur}
\textsl{Melee Weapon Attack:} 12 to hit, reach 10 ft., one target. \textsl{Hit:} 23 (3d10 + 7) force damage plus 14 (4d6) cold damage.

\DndMonsterAction{Claws}
\textsl{Melee Weapon Attack:} 12 to hit, reach 5 ft., one target. \textsl{Hit:} 20 (3d8 + 7) slashing damage. If the target is a Large or smaller creature, it has the grappled condition (escape DC 20). While a creature is grappled, the war devil can't make Claws attacks.

\end{DndMonster}



\begin{DndMonster}[float=t]{White Abishai}
\DndMonsterType{Medium fiend (devil), lawful evil}
\DndMonsterBasics[
armor-class = {15 (natural armor)},
hit-points = {\DndDice{8d8 + 32}},
speed = {30 ft., fly 40 ft.}
]
\DndMonsterAbilityScores[
 str = 16,
 dex = 11,
 con = 18,
 int = 11,
 wis = 12,
 cha = 13
]
\DndMonsterDetails[
saving-throws = {Str +6, Con +7},
damage-resistances = {bludgeoning, piercing, slashing from nonmagical attacks that aren't silvered},
damage-immunities = {cold, fire, poison},
condition-immunities = {poisoned},
senses = {darkvision 120 ft., passive perception 11},
languages = {Draconic, Infernal, telepathy 120 ft.},
challenge = 6,
]
\DndMonsterAction{Devil's Sight}
Magical darkness doesn't impede the abishai's darkvision.

\DndMonsterAction{Magic Resistance}
The abishai has advantage on saving throws against spells and other magical effects.

\DndMonsterAction{Reckless}
At the start of its turn, the abishai can gain advantage on all melee weapon attack rolls during that turn, but attack rolls against it have advantage until the start of its next turn.

\DndMonsterSection{Actions}
\DndMonsterAction{Multiattack}
The abishai makes one Bite attack, one Claw attack, and one Longsword attack.

\DndMonsterAction{Bite}
\textsl{Melee Weapon Attack:} 6 to hit, reach 5 ft., one target. \textsl{Hit:} 5 (1d4 + 3) piercing damage plus 3 (1d6) cold damage.

\DndMonsterAction{Claw}
\textsl{Melee Weapon Attack:} 6 to hit, reach 5 ft., one target. \textsl{Hit:} 8 (1d10 + 3) slashing damage.

\DndMonsterAction{Longsword}
\textsl{Melee Weapon Attack:} 6 to hit, reach 5 ft., one target. \textsl{Hit:} 7 (1d8 + 3) force damage, or 8 (1d10 + 3) force damage if used with two hands.

\DndMonsterSection{Reactions}
\DndMonsterAction{Vicious Reprisal}
In response to taking damage, the abishai makes one Bite attack against a random creature within 5 feet of it. If no creature is within reach, the abishai moves up to half its speed toward an enemy it can see, without provoking opportunity attack.

\end{DndMonster}



\begin{DndMonster}[float=t]{Wolf}
\DndMonsterType{Medium beast, unaligned}
\DndMonsterBasics[
armor-class = {13 (natural armor)},
hit-points = {\DndDice{2d8 + 2}},
speed = {40 ft.}
]
\DndMonsterAbilityScores[
 str = 12,
 dex = 15,
 con = 12,
 int = 3,
 wis = 12,
 cha = 6
]
\DndMonsterDetails[
skills = {Perception +3, Stealth +4},
challenge = 1/4,
]
\DndMonsterAction{Keen Hearing and Smell}
The wolf has advantage on Wisdom (Perception) checks that rely on hearing or smell.

\DndMonsterAction{Pack Tactics}
The wolf has advantage on an attack roll against a creature if at least one of the wolf's allies is within 5 feet of the creature and the ally isn't incapacitated.

\DndMonsterSection{Actions}
\DndMonsterAction{Bite}
\textsl{Melee Weapon Attack:} 4 to hit, reach 5 ft., one target. \textsl{Hit:} 7 (2d4 + 2) piercing damage. If the target is a creature, it must succeed on a DC 11 Strength saving throw or be knocked prone.

\end{DndMonster}


\clearpage

\begin{DndMonster}[float=t]{Xorn}
\DndMonsterType{Medium elemental, neutral}
\DndMonsterBasics[
armor-class = {19 (natural armor)},
hit-points = {\DndDice{7d8 + 42}},
speed = {20 ft., burrow 20 ft.}
]
\DndMonsterAbilityScores[
 str = 17,
 dex = 10,
 con = 22,
 int = 11,
 wis = 10,
 cha = 11
]
\DndMonsterDetails[
skills = {Perception +6, Stealth +3},
damage-resistances = {piercing, slashing from nonmagical attacks that aren't adamantine},
senses = {darkvision 60 ft., tremorsense 60 ft., passive perception 16},
languages = {Terran},
challenge = 5,
]
\DndMonsterAction{Earth Glide}
The xorn can burrow through nonmagical, unworked earth and stone. While doing so, the xorn doesn't disturb the material it moves through.

\DndMonsterAction{Stone Camouflage}
The xorn has advantage on Dexterity (Stealth) checks made to hide in rocky terrain.

\DndMonsterAction{Treasure Sense}
The xorn can pinpoint, by scent, the location of precious metals and stones, such as coins and gems, within 60 feet of it.

\DndMonsterSection{Actions}
\DndMonsterAction{Multiattack}
The xorn makes three claw attacks and one bite attack.

\DndMonsterAction{Bite}
\textsl{Melee Weapon Attack:} 6 to hit, reach 5 ft., one target. \textsl{Hit:} 13 (3d6 + 3) piercing damage.

\DndMonsterAction{Claw}
\textsl{Melee Weapon Attack:} 6 to hit, reach 5 ft., one target. \textsl{Hit:} 6 (1d6 + 3) slashing damage.

\end{DndMonster}



\begin{DndMonster}[float=t]{Young Red Dragon}
\DndMonsterType{Large dragon, chaotic evil}
\DndMonsterBasics[
armor-class = {18 (natural armor)},
hit-points = {\DndDice{17d10 + 85}},
speed = {40 ft., climb 40 ft., fly 80 ft.}
]
\DndMonsterAbilityScores[
 str = 23,
 dex = 10,
 con = 21,
 int = 14,
 wis = 11,
 cha = 19
]
\DndMonsterDetails[
saving-throws = {Dex +4, Con +9, Wis +4, Cha +8},
skills = {Perception +8, Stealth +4},
damage-immunities = {fire},
senses = {blindsight 30 ft., darkvision 120 ft., passive perception 18},
languages = {Common, Draconic},
challenge = 10,
]
\DndMonsterSection{Actions}
\DndMonsterAction{Multiattack}
The dragon makes three attacks: one with its bite and two with its claws.

\DndMonsterAction{Bite}
\textsl{Melee Weapon Attack:} 10 to hit, reach 10 ft., one target. \textsl{Hit:} 17 (2d10 + 6) piercing damage plus 3 (1d6) fire damage.

\DndMonsterAction{Claw}
\textsl{Melee Weapon Attack:} 10 to hit, reach 5 ft., one target. \textsl{Hit:} 13 (2d6 + 6) slashing damage.

\DndMonsterAction{Fire Breath (Recharge 5\textendash 6)}
The dragon exhales fire in a 30-foot cone. Each creature in that area must make a DC 17 Dexterity saving throw, taking 56 (16d6) fire damage on a failed save, or half as much damage on a successful one.

\end{DndMonster}



\begin{DndMonster}[float=t]{Young Red Shadow Dragon}
\DndMonsterType{Large dragon, chaotic evil}
\DndMonsterBasics[
armor-class = {18 (natural armor)},
hit-points = {\DndDice{17d10 + 85}},
speed = {40 ft., climb 40 ft., fly 80 ft.}
]
\DndMonsterAbilityScores[
 str = 23,
 dex = 10,
 con = 21,
 int = 14,
 wis = 11,
 cha = 19
]
\DndMonsterDetails[
saving-throws = {Dex +5, Con +10, Wis +5, Cha +9},
skills = {Perception +10, Stealth +10},
damage-resistances = {necrotic},
damage-immunities = {fire},
senses = {blindsight 30 ft., darkvision 120 ft., passive perception 18},
languages = {Common, Draconic},
challenge = 13,
]
\DndMonsterAction{Living Shadow}
While in dim light or darkness, the dragon has resistance to damage that isn't force, psychic, or radiant.

\DndMonsterAction{Shadow Stealth}
While in dim light or darkness, the dragon can take the Hide action as a bonus action.

\DndMonsterAction{Sunlight Sensitivity}
While in sunlight, the dragon has disadvantage on attack rolls, as well as on Wisdom (Perception) checks that rely on sight.

\DndMonsterSection{Actions}
\DndMonsterAction{Multiattack}
The dragon makes three attacks: one with its bite and two with its claws.

\DndMonsterAction{Bite}
\textsl{Melee Weapon Attack:} 11 to hit, reach 10 ft., one target. \textsl{Hit:} 17 (2d10 + 6) piercing damage plus 3 (1d6) necrotic damage.

\DndMonsterAction{Claw}
\textsl{Melee Weapon Attack:} 11 to hit, reach 5 ft., one target. \textsl{Hit:} 13 (2d6 + 6) slashing damage.

\DndMonsterAction{Shadow Breath (Recharge 5\textendash 6)}
The dragon exhales shadowy fire in a 30-foot cone. Each creature in that area must make a DC 18 Dexterity saving throw, taking 56 (16d6) necrotic damage on a failed save, or half as much damage on a successful one. A humanoid reduced to 0 hit points by this damage dies, and an undead \textbf{shadow} rises from its corpse and acts immediately after the dragon in the initiative count. The shadow is under the dragon's control.

\end{DndMonster}



\begin{DndMonster}[float=t]{Zagum}
\DndMonsterType{Medium fiend (devil), lawful evil}
\DndMonsterBasics[
armor-class = {19 (natural armor)},
hit-points = {\DndDice{30d8 + 120}},
speed = {30 ft.}
]
\DndMonsterAbilityScores[
 str = 17,
 dex = 16,
 con = 18,
 int = 22,
 wis = 25,
 cha = 27
]
\DndMonsterDetails[
saving-throws = {Int +12, Wis +13, Cha +14},
skills = {Arcana +12, Deception +20, Insight +13, Intimidation +14, Performance +14, Persuasion +20},
damage-resistances = {cold; bludgeoning, piercing, slashing from nonmagical attacks that aren't silvered},
damage-immunities = {fire, poison, psychic},
condition-immunities = {charmed, poisoned},
senses = {darkvision 120 ft., passive perception 17},
languages = {Celestial, Common, Draconic, Infernal, telepathy 120 ft.},
challenge = 20,
]
\DndMonsterAction{Contractually Obligated}
Zagum has advantage to hit creatures under an infernal contract and when he hits such creatures he delivers an additional 9 (2d8) damage of a type the target is most vulnerable to. A creature bound under an infernal contract also makes all saving throws against effects originating from Zagum with disadvantage.

\DndMonsterAction{Devil's Sight}
Magical darkness doesn't impede Zagum's darkvision.

\DndMonsterAction{Magic Resistance}
Zagum has advantage on saving throws against spells and other magical effects.

\DndMonsterSection{Actions}
\DndMonsterAction{Multiattack}
Zagum makes three attacks with Disarming Tongue or Oration.

\DndMonsterAction{Disarming Tongue}
\textsl{Melee Weapon Attack:} 9 to hit, reach 10 ft., one target. \textsl{Hit:} 10 (3d4 + 3) slashing damage. A target hit by this attack has vulnerability to psychic damage until the start of their next turn.

\DndMonsterAction{Oration}
\textsl{Ranged Spell Attack:} 14 to hit, range 60 ft., one target. \textsl{Hit:} 35 (6d8 + 8) psychic damage.

\DndMonsterAction{Legal Jargon (Recharge 4\textendash 6)}
Zagum chooses a point he can see within 90 feet of him. All creatures within a 20-foot-radius sphere centered on that point must make a DC 22 Wisdom saving throw. Creatures that fail the save are confused for up to 1 minute and must roll a d10 at the start of each of their turns.

On a 1, a creature moves in a random direction then ends their turn. On a 2-6, a creature can't use actions this turn. On a 7-8, a creature must use their action to make a melee attack against a randomly determined creature within their reach that is not a Fiend. On a 9-10, a creature can move and act normally.

A confused creature may repeat the saving throw at the end of its turn, ending the effect on a success.

\end{DndMonster}



\begin{DndMonster}[float*=b,width=\textwidth + 8pt]{Zariel}
\begin{multicols}{2}
\DndMonsterType{Large fiend (devil), lawful evil}
\DndMonsterBasics[
armor-class = {21 (natural armor)},
hit-points = {\DndDice{29d10 + 261}},
speed = {50 ft., fly 150 ft.}
]
\DndMonsterAbilityScores[
 str = 27,
 dex = 24,
 con = 28,
 int = 26,
 wis = 27,
 cha = 30
]
\DndMonsterDetails[
saving-throws = {Int +16, Wis +16, Cha +18},
skills = {Intimidation +18, Perception +16},
damage-resistances = {cold; fire; radiant; bludgeoning, piercing, slashing from nonmagical attacks that aren't silvered},
damage-immunities = {necrotic, poison},
condition-immunities = {charmed, exhaustion, frightened, poisoned},
senses = {darkvision 120 ft., passive perception 26},
languages = {all, telepathy 120 ft.},
challenge = 26,
]
\DndMonsterAction{Devil's Sight}
Magical darkness doesn't impede Zariel's darkvision.

\DndMonsterAction{Legendary Resistance (3/Day)}
If Zariel fails a saving throw, she can choose to succeed instead.

\DndMonsterAction{Magic Resistance}
Zariel has advantage on saving throws against spells and other magical effects.

\DndMonsterAction{Regeneration}
Zariel regains 20 hit points at the start of her turn. If she takes radiant damage, this trait doesn't function at the start of her next turn. Zariel dies only if she starts her turn with 0 hit points and doesn't regenerate.

\DndMonsterAction{Spellcasting}
Zariel casts one of the following spells, requiring no material components and using Charisma as the spellcasting ability (spell save DC 26):
\begin{DndMonsterSpells}
  \DndInnateSpellLevel{\textit{Alter Self} (can become Medium when changing her appearance), \textit{Detect Evil and Good}, \textit{Fireball}, \textit{Invisibility} (self only), \textit{Major Image}, \textit{Wall of Fire}}
  \DndInnateSpellLevel[2e]{\textit{Blade Barrier}, \textit{Dispel Evil and Good}, \textit{Finger of Death}}
\end{DndMonsterSpells}
\DndMonsterSection{Actions}
\DndMonsterAction{Multiattack}
Zariel makes three Flail or Longsword attacks. She can replace one attack with a use of Horrid Touch, if available.

\DndMonsterAction{Flail}
\textsl{Melee Weapon Attack:} 16 to hit, reach 10 ft., one target. \textsl{Hit:} 17 (2d8 + 8) force damage plus 36 (8d8) fire damage.

\DndMonsterAction{Longsword}
\textsl{Melee Weapon Attack:} 16 to hit, reach 10 ft., one target. \textsl{Hit:} 17 (2d8 + 8) radiant damage or 19 (2d10 + 8) radiant damage when used with two hands, plus 36 (8d8) fire damage.

\DndMonsterAction{Horrid Touch (Recharge 5\textendash 6)}
Zariel touches one creature within 10 feet of her. The target must succeed on a DC 26 Constitution saving throw or take 44 (8d10) necrotic damage and have the poisoned condition for 1 minute. While poisoned in this way, the target has the blinded and deafened conditions. The target can repeat the saving throw at the end of each of its turns, ending the effect on itself on a success.

\DndMonsterAction{Teleport}
Zariel magically teleports, along with any equipment she is wearing or carrying, up to 120 feet to an unoccupied space she can see.

\DndMonsterSection{Legendary Actions}
Zariel can take 3 legendary actions, choosing from the options below. Only one legendary action can be used at a time and only at the end of another creature's turn. Zariel regains spent legendary actions at the start of its turn.

\begin{DndMonsterLegendaryActions}
\DndMonsterLegendaryAction{Teleport}{Zariel uses Teleport.
}
\DndMonsterLegendaryAction{Immolating Gaze (Costs 2 Actions)}{Zariel turns her magical gaze toward one creature she can see within 120 feet of her and commands it to burn. The target must succeed on a DC 26 Wisdom saving throw or take 22 (4d10) fire damage.
}
\end{DndMonsterLegendaryActions}
\DndMonsterSection{Regional Effects}
The region containing Zariel's lair is warped by her magic, which creates the following effect:


\begin{description}
\item[Hellscape.]
The area within 9 miles of the lair is filled with screaming voices and the stench of burning meat.

\end{description}

If Zariel dies, these effects fade over the course of 1d10 days.

\end{multicols}
\end{DndMonster}



\begin{DndMonster}[float=t]{Zythan}
\DndMonsterType{Medium humanoid, lawful neutral}
\DndMonsterBasics[
armor-class = {12 (15 with \textit{mage armor})},
hit-points = {\DndDice{18d8 + 36}},
speed = {30 ft.}
]
\DndMonsterAbilityScores[
 str = 10,
 dex = 14,
 con = 14,
 int = 20,
 wis = 15,
 cha = 16
]
\DndMonsterDetails[
saving-throws = {Int +10, Wis +7},
skills = {Arcana +15, History +15, Investigation +10, Religion +10},
damage-resistances = {bludgeoning, piercing, slashing from nonmagical attacks (from Stoneskin); damage from spells},
languages = {Celestial, Common, Draconic, Elvish, Infernal, Primordial},
challenge = 13,
]
\DndMonsterAction{Battlemage}
Being within 5 feet of a hostile creature doesn't impose disadvantage on Zythan's ranged attack rolls.

\DndMonsterAction{Magic Resistance}
Zythan has advantage on saving throws against spells and other magical effects.

\DndMonsterAction{Spellcasting}
Zythan casts one of the following spells, using Intelligence as the spellcasting ability (spell save DC 18).
\begin{DndMonsterSpells}
  \DndInnateSpellLevel{\textit{Arcane Eye}, \textit{Detect Magic}, \textit{Detect Thoughts}, \textit{Fly}, \textit{Light}, \textit{Locate Object}, \textit{Mage Armor}, \textit{Mage Hand}, \textit{Message}, \textit{Prestidigitation}}
  \DndInnateSpellLevel[1]{\textit{Cone of Cold}, \textit{Foresight}, \textit{Stoneskin}, \textit{Rary's telepathic bond}, \textit{True Seeing}}
\end{DndMonsterSpells}
\DndMonsterSection{Actions}
\DndMonsterAction{Multiattack}
Zythan makes three Dagger or Arcane Burst attacks.

\DndMonsterAction{Dagger}
\textsl{Melee or Ranged Weapon Attack:} 7 to hit, reach 5 ft. or range 20/60 ft., one target. \textsl{Hit:} 4 (1d4 + 2) piercing damage.

\DndMonsterAction{Arcane Burst}
\textsl{Ranged Spell Attack:} 10 to hit, range 120 ft., one target. \textsl{Hit:} 27 (4d10 + 5) radiant damage.

\DndMonsterAction{Spirit Fissure (Recharge 4\textendash 6)}
Zythan causes a disruption in arcane energies in a 300-foot line that is 5 feet wide. Each creature in the line must make a DC 18 Constitution saving throw, taking 65 (10d12) radiant damage on a failed save, or half as much damage on a successful one.

\DndMonsterSection{Reactions}
\DndMonsterAction{Portent (3/Day)}
As a reaction to a creature Zythan can see making an attack roll, a saving throw, or an ability check, Zythan rolls a d20 and chooses whether to use that roll in place of the original roll.

\end{DndMonster}


\end{document}
